\documentclass[12pt,a4paper]{report}

% =========================================================================
% 1. PACKAGES REQUIS
% =========================================================================
% Encodage & langue
\usepackage[utf8]{inputenc}
\usepackage[T1]{fontenc}
\usepackage[french]{babel}
\usepackage{longtable}
\usepackage{booktabs}
\usepackage{array}
\usepackage[table]{xcolor}
\usepackage{ragged2e}
\usepackage{needspace}
\usepackage{float}

\definecolor{TableHeaderBlue}{RGB}{32,78,121}
\definecolor{TableLightBlue}{RGB}{232,240,248}
\definecolor{TableVeryLight}{RGB}{248,251,253}


\setlength{\LTpre}{4pt}
\setlength{\LTpost}{6pt}
\renewcommand{\arraystretch}{1.25}
% Mise en page
\usepackage[a4paper,left=2cm,right=2cm,top=2cm,bottom=2cm]{geometry}
\usepackage{titlesec}
\usepackage{ragged2e}
\usepackage{needspace}
\usepackage{float}             % pour les figures/tableaux fixes
\usepackage[section]{placeins}  % empêche les flottants de dépasser leur section

% Figures et graphiques
\usepackage{graphicx}
\usepackage{tcolorbox}         % encadrés (chemin critique, etc.)

% Tableaux
\usepackage{array}
\usepackage{tabularx}
\usepackage{booktabs}
\usepackage{multirow}
\usepackage[table]{xcolor}     % couleurs dans les tableaux
\usepackage{pdflscape}

% Listes & citations
\usepackage{enumitem}
\usepackage{csquotes}
\usepackage{amsmath}
\usepackage{amssymb}
\usepackage[hypertexnames=false,bookmarks=false]{hyperref}
\usepackage{lmodern}
\usepackage{fix-cm}
\usepackage{fancyhdr}

% =========================================================================
% 2. CONFIGURATION DU DOCUMENT (À REMPLIR PAR LES ÉTUDIANTS)
% =========================================================================

% Informations sur les étudiants
\newcommand{\EtudiantUn}{Mohamed Ben Abdeladhim}
\newcommand{\EtudiantDeux}{Rahma Barouni}

% Informations sur le projet
\newcommand{\TitreProjet}{TalentPredict}
\newcommand{\LieuProjet}{AgileX Management}
\newcommand{\DateSoutenance}{\today} % ou remplacez \today par "6 mai 2026"

% Membres du Jury
\newcommand{\PresidentJury}{Pr. Nom Prénom}
\newcommand{\Rapporteur}{Pr. Nom Prénom}
\newcommand{\EncadreurAcademique}{Dr. Ben Aïcha Takwa}
\newcommand{\EncadrantProfessionnel}{Med Amine Gharbi}

% Chemins des logos
\newcommand{\LogoISIMA}{images/isima.png}
\newcommand{\LogoUniversite}{images/universit.png}


% =========================================================================
% 3. DÉFINITION DES COMMANDES ET STYLES
% =========================================================================
\usepackage{pgfplots}
\pgfplotsset{compat=1.18} % ou la version de ton compilateur

% Réglages de mise en page généraux
\setlength{\parindent}{0pt}
\setlength{\parskip}{0.4em}
\setlength{\tabcolsep}{4pt}
\setlength{\headheight}{28pt}
\renewcommand{\arraystretch}{1.05}
\raggedbottom
\setcounter{topnumber}{5}
\setcounter{bottomnumber}{5}
\setcounter{totalnumber}{10}
\renewcommand{\topfraction}{0.95}
\renewcommand{\bottomfraction}{0.95}
\renewcommand{\textfraction}{0.05}
\renewcommand{\floatpagefraction}{0.8}
\setlength{\emergencystretch}{3em}
\hbadness=10000
\vbadness=10000
\hfuzz=60pt
\vfuzz=600pt

% Réduction de l'espace autour des titres
\titleformat{\chapter}[display]
  {\normalfont\huge\bfseries}
  {Chapitre \thechapter}
  {0.4em}
  {}
\titlespacing*{\chapter}{0pt}{-14pt}{10pt}
\titlespacing*{\section}{0pt}{6pt}{6pt}

% Configuration de l'en-tête et du pied de page
\pagestyle{fancy}
\fancyhf{}
\fancyhead[L]{\ifnum\value{chapter}>0 Chapitre \thechapter\fi}
\fancyhead[R]{\TitreProjet}
\fancyfoot[C]{\thepage}
\renewcommand{\headrulewidth}{0.5pt}
\renewcommand{\footrulewidth}{0.5pt}
\fancypagestyle{plain}{%
  \fancyhf{}%
  \fancyhead[L]{\ifnum\value{chapter}>0 Chapitre \thechapter\fi}%
  \fancyhead[R]{\TitreProjet}%
  \fancyfoot[C]{\thepage}%
  \renewcommand{\headrulewidth}{0.5pt}%
  \renewcommand{\footrulewidth}{0.5pt}%
}

% Ligne horizontale personnalisée pour le titre
\newcommand{\HRule}{\rule{\linewidth}{0.4mm}}

% Réglages et types de colonnes pour les tableaux (ltablex)
\newcommand{\SafeRaggedRight}{\rightskip=0pt plus 2em\relax\spaceskip=.3333em\xspaceskip=.5em\relax}
\newcolumntype{L}{>{\SafeRaggedRight\arraybackslash}X}
\newcolumntype{C}[1]{>{\Centering\arraybackslash}p{#1}}
\newcolumntype{Y}{>{\SafeRaggedRight\arraybackslash}X}
\newlength{\TableWidth}
\setlength{\TableWidth}{0.92\textwidth}
\newcolumntype{P}[1]{>{\SafeRaggedRight\hspace{0pt}\arraybackslash}p{#1}}
\newcommand{\logoplaceholder}{\parbox[c][0.9cm][c]{1.0cm}{\centering\scriptsize Logo}}
\newcommand{\chapterfigure}[3][0.8\textwidth]{%
  \begin{figure}[H]
    \centering
    \IfFileExists{#2}{%
      \includegraphics[width=#1]{#2}%
    }{%
      \fbox{\parbox[c][5cm][c]{0.82\textwidth}{\centering\textit{Figure non disponible : #3}}}%
    }%
    \caption{#3}
  \end{figure}
}
\newcommand{\techlogo}[1]{%
  \IfFileExists{#1}{\includegraphics[width=1.45cm]{#1}}{\logoplaceholder}
}
\newcommand{\figureplaceholder}[2]{%
  \begin{figure}[H]
    \centering
    \fbox{\begin{minipage}[c][5cm][c]{0.9\textwidth}
      \centering
      \textit{Espace réservé pour l'insertion de la figure}
    \end{minipage}}
    \caption{#1}
    \label{#2}
  \end{figure}
}

% Macros de priorité
\newcommand{\MustFR}{\textbf{Must}}
\newcommand{\ShouldFR}{\textbf{Should}}
\newcommand{\CouldFR}{\textbf{Could}}

% Style de ligne "Epic" pour les tableaux
\newcommand{\EpicRow}[1]{\rowcolor{gray!15}\multicolumn{6}{l}{\textbf{#1}}\\}

\newcommand{\placeholderfigure}[2]{%
  \begin{figure}[H]
    \centering
    \IfFileExists{images/#1}{%
      \includegraphics[width=0.85\textwidth]{images/#1}%
    }{%
      \fbox{\parbox[c][6cm][c]{0.82\textwidth}{\centering Emplacement réservé pour la figure}}%
    }
    \caption{#2}
  \end{figure}
}


\usepackage{listings}


\lstdefinelanguage{json}{
    basicstyle=\ttfamily\small,
    sensitive=false,
    morestring=[b]",
    morecomment=[l]{//}
}
\lstdefinelanguage{yaml}{
    basicstyle=\ttfamily\small,
    sensitive=false,
    morecomment=[l]{\#},
    morestring=[b]',
    morestring=[b]"
}
\lstdefinelanguage{properties}{
    basicstyle=\ttfamily\small,
    sensitive=false,
    morecomment=[l]{\#},
    morecomment=[l]{!},
    morestring=[b]"
}



% Configuration globale pour le code (Prompt et JSON)
\lstset{
    basicstyle=\ttfamily\small,
    breaklines=true,
    frame=single,
    backgroundcolor=\color{gray!5},
    showstringspaces=false,
    captionpos=b,
    columns=fullflexible,
    literate=
    {à}{{\`a}}1 {â}{{\^a}}1 {ä}{{\"a}}1
    {ç}{{\c{c}}}1
    {é}{{\'e}}1 {è}{{\`e}}1 {ê}{{\^e}}1 {ë}{{\"e}}1
    {î}{{\^\i}}1 {ï}{{\"\i}}1
    {ô}{{\^o}}1 {ö}{{\"o}}1
    {ù}{{\`u}}1 {û}{{\^u}}1 {ü}{{\"u}}1
    {À}{{\`A}}1 {Â}{{\^A}}1 {Ä}{{\"A}}1
    {Ç}{{\c{C}}}1
    {É}{{\'E}}1 {È}{{\`E}}1 {Ê}{{\^E}}1 {Ë}{{\"E}}1
    {Î}{{\^I}}1 {Ï}{{\"I}}1
    {Ô}{{\^O}}1 {Ö}{{\"O}}1
    {Ù}{{\`U}}1 {Û}{{\^U}}1 {Ü}{{\"U}}1
}
% =========================================================================
% 4. DÉBUT DU DOCUMENT
% =========================================================================
\begin{document}

% -
% PAGE DE GARDE
% -


\begin{titlepage}
\newgeometry{top=0.35cm,bottom=0.35cm,left=2cm,right=2cm}
\thispagestyle{empty}
\pagestyle{empty}

% EN-TÊTE 
\noindent
\begin{tabular}{@{}p{0.42\textwidth}@{}c@{}p{0.42\textwidth}@{}}
\begin{minipage}[t]{0.42\textwidth}
  \small République Tunisienne\\
  \small Ministère de l'Enseignement Supérieur et\\
  \small de la Recherche Scientifique
\end{minipage}
&
\includegraphics[height=1.6cm]{\LogoISIMA}
&
\begin{minipage}[t]{0.42\textwidth}
  \raggedleft
  \small Université de Monastir\\
  \small Institut Supérieur d'Informatique de\\
  \small Mahdia
\end{minipage}
\end{tabular}

\vspace{0.18cm}
\noindent\rule{\textwidth}{0.8pt}

\begin{flushright}
  \textbf{\textit{MEMOIRE}}\\[1pt]
  \textbf{N° d'ordre :} LBC-BI-3-2026-032
\end{flushright}

\noindent\rule{\textwidth}{0.8pt}

%--- TITRE PRINCIPAL ---
\vspace{0.22cm}
\begin{center}
  {\fontsize{36}{42}\selectfont \textit{M\'{E}MOIRE}}
\end{center}

\vspace{0.22cm}

%--- CORPS ---
\begin{center}
  \textit{Présenté à}\\[0.2cm]
  {\large \textbf{Institut Supérieur d'Informatique de Mahdia}}\\[0.8cm]
  \textit{Effectué à}\\[0.25cm]
  \textbf{AgileX Management}\\[0.25cm]
  \textit{En vue de l'obtention du diplôme de}\\[1.2cm]
  \textit{Elaboré par :}\\[0.25cm]
  Mohamed Abdeladhim, Rahma Barouni
\end{center}

\vspace{0.5cm}
\noindent\rule{\textwidth}{2pt}

\begin{center}
  \vspace{0.08cm}
  \textbf{TalentPredict}
  \vspace{0.08cm}
\end{center}

\noindent\rule{\textwidth}{2pt}

\vspace{0.8cm}
\begin{center}
  \textit{\textbf{Soutenu le :}} 11 juin 2026\\[0.08cm]
  \textit{\textbf{devant le jury composé de :}}
\end{center}

\vspace{0.6cm}

%--- JURY ---
\begin{tabular}{ll}
  \textbf{Président}               & : Tarek M. Hamdani \\[1pt]
  \textbf{Rapporteur}              & : Emna Ben Ayed \\[1pt]
  \textbf{Encadrant Académique}    & : Ben Aichaa Takwa \\[1pt]
  \textbf{Encadrant Professionnel} & : Med Amine Gharbi \\
\end{tabular}

\vfill

\begin{center}
  Année universitaire 2025-2026
\end{center}
\end{titlepage}

% -------------------------------------------------------------------------
% DÉDICACE
% -------------------------------------------------------------------------
\chapter*{Dédicace}
\thispagestyle{empty}

\vspace*{1cm}

\begin{center}
\itshape
À celles et ceux qui nous ont soutenus sans réserve,\\
à nos familles pour leurs sacrifices,\\
à nos proches pour leur confiance et leur patience.

\vspace{0.4cm}

Pour leur encouragement constant, leur bienveillance et leur présence\\
dans les moments d’effort comme dans les instants de doute.

\vspace{0.6cm}

Ce travail leur est dédié.
\end{center}

\vspace{1cm}

\begin{flushright}
\large\itshape
\EtudiantUn \\
\EtudiantDeux
\end{flushright}

\clearpage

% -------------------------------------------------------------------------
% REMERCIEMENTS
% -------------------------------------------------------------------------
\chapter*{Remerciements} 
\thispagestyle{empty}

Nous tenons à exprimer notre profonde gratitude à toutes les personnes qui ont contribué, de près ou de loin, à la réalisation de ce projet.

Nos sincères remerciements s’adressent à notre encadrant pour son accompagnement rigoureux, ses conseils avisés et sa disponibilité tout au long de ce travail.

Nous remercions également les membres du jury pour l’intérêt qu’ils portent à ce mémoire et pour le temps consacré à son évaluation.

Notre reconnaissance va aussi à l’ensemble des enseignants et du personnel de l’Institut Supérieur d’Informatique de Mahdia (ISIMA) pour la qualité de la formation reçue.

Enfin, nous adressons une pensée particulière à nos familles et à nos proches pour leur soutien indéfectible et leurs encouragements constants.

\clearpage

% -------------------------------------------------------------------------
% RÉSUMÉ / ABSTRACT
% -------------------------------------------------------------------------
\chapter*{Résumé}
Dans un contexte de transformation numérique où la gestion des talents constitue un levier stratégique majeur, les méthodes classiques d'évaluation des compétences  fondées sur des CV déclaratifs, des entretiens ponctuels et des appréciations subjectives  montrent leurs limites : fragmentation des données, biais déclaratifs, impossibilité d'évaluer objectivement les compétences comportementales à grande échelle.
Ce mémoire de fin d'études, réalisé au sein de la société AgileX Management, présente la conception et le développement de TalentPredict AI, une plateforme intelligente d'évaluation multi-sources des compétences techniques et comportementales. La problématique centrale est la suivante : comment concevoir une solution capable d'analyser de manière structurée, fiable et sécurisée les profils des collaborateurs, tout en générant des recommandations de formation personnalisées exploitables par les décideurs RH ?
La solution repose sur une architecture hybride et modulaire articulant plusieurs couches complémentaires. L'évaluation des compétences techniques combine l'analyse du profil GitHub, le parsing du CV (150+ technologies), des QCM adaptatifs et des challenges de code évalués par CodeLlama. L'évaluation des compétences comportementales s'appuie sur le modèle psychométrique PCM, l'analyse sémantique locale via Ollama, et des scénarios d'évaluation générés par l'intelligence artificielle. Un moteur de recommandations transforme ensuite les lacunes identifiées en parcours de formation personnalisés, suivis via un tableau Kanban et validés par certification PDF.
Les technologies mises en œuvre incluent Angular, Spring Boot, FastAPI, n8n, Ollama (llama3.2), CodeLlama, PostgreSQL et Docker. L'inférence 100 \% locale garantit la conformité RGPD par conception. Le développement a suivi la méthodologie Agile Scrum en quatre sprints progressifs.
Les résultats obtenus valident la faisabilité d'une évaluation unifiée, automatisée et objectivée des talents, combinant preuves techniques observables et analyse comportementale. Ce travail ouvre des perspectives d'évolution vers l'intégration de référentiels de compétences standards (ROME, ESCO) et vers une mise à l'échelle cloud tout en préservant la souveraineté des données.

% -------------------------------------------------------------------------
% TABLE DES MATIÈRES ET LISTES
% -------------------------------------------------------------------------
\tableofcontents
\clearpage

\listoffigures
\clearpage

\listoftables
\clearpage

% -------------------------------------------------------------------------
% LISTE DES ABRÉVIATIONS
% -------------------------------------------------------------------------
\chapter*{Liste des Abréviations}
\thispagestyle{empty}

\begin{longtable}{|P{0.20\textwidth}|P{0.72\textwidth}|}
\hline
\rowcolor{TableHeaderBlue}
\color{white}\textbf{Acronyme} & \color{white}\textbf{Signification} \\
\hline
\endfirsthead

\hline
\rowcolor{TableHeaderBlue}
\color{white}\textbf{Acronyme} & \color{white}\textbf{Signification} \\
\hline
\endhead

\hline
\endfoot

\textbf{AI} & Intelligence Artificielle (\textit{Artificial Intelligence}) \\
\hline
\textbf{API} & Interface de Programmation Applicative \\
\hline
\textbf{ASGI} & Interface de Passerelle de Serveur Asynchrone (\textit{Asynchronous Server Gateway Interface}) \\
\hline
\textbf{BCrypt} & Algorithme de hachage de mots de passe adaptatif \\
\hline
\textbf{BI} & Intelligence d'Affaires (\textit{Business Intelligence}) \\
\hline
\textbf{CI/CD} & Intégration Continue / Déploiement Continu (\textit{Continuous Integration / Continuous Deployment}) \\
\hline
\textbf{CRF} & Champ Aléatoire Conditionnel (\textit{Conditional Random Field}) \\
\hline
\textbf{CRUD} & Créer, Lire, Mettre à jour, Supprimer (\textit{Create, Read, Update, Delete}) \\
\hline
\textbf{CSS} & Feuilles de Style en Cascade (\textit{Cascading Style Sheets}) \\
\hline
\textbf{CVE} & Vulnérabilités et Expositions Communes (\textit{Common Vulnerabilities and Exposures}) \\
\hline
\textbf{DISC} & Modèle psychométrique Dominance, Influence, Stabilité, Conformité \\
\hline
\textbf{DTO} & Objet de Transfert de Données (\textit{Data Transfer Object}) \\
\hline
\textbf{ERD} & Diagramme Entité-Relation (\textit{Entity-Relationship Diagram}) \\
\hline
\textbf{GPEC} & Gestion Prévisionnelle des Emplois et des Compétences \\
\hline
\textbf{HTML} & Langage de Balisage Hypertexte (\textit{HyperText Markup Language}) \\
\hline
\textbf{HTTP} & Protocole de Transfert Hypertexte (\textit{HyperText Transfer Protocol}) \\
\hline
\textbf{HTTPS} & Protocole de Transfert Hypertexte Sécurisé (\textit{HyperText Transfer Protocol Secure}) \\
\hline
\textbf{IaC} & Infrastructure en tant que Code (\textit{Infrastructure as Code}) \\
\hline
\textbf{IAM} & Gestion des Identités et des Accès (\textit{Identity and Access Management}) \\
\hline
\textbf{JPA} & Interface de Persistance Java (\textit{Java Persistence API}) \\
\hline
\textbf{JSON} & Notation Objet JavaScript \\
\hline
\textbf{JWT} & Jeton Web JSON \\
\hline
\textbf{KPI} & Indicateur Clé de Performance (\textit{Key Performance Indicator}) \\
\hline
\textbf{LLM} & Grand Modèle de Langage (\textit{Large Language Model}) \\
\hline
\textbf{LXP} & Plateforme d'Expérience d'Apprentissage (\textit{Learning Experience Platform}) \\
\hline
\textbf{MCQ / QCM} & Questionnaire à Choix Multiples (\textit{Multiple Choice Questions}) \\
\hline
\textbf{ML} & Apprentissage Automatique (\textit{Machine Learning}) \\
\hline
\textbf{MVC} & Modèle-Vue-Contrôleur (\textit{Model-View-Controller}) \\
\hline
\textbf{MVP} & Produit Minimum Viable (\textit{Minimum Viable Product}) \\
\hline
\textbf{NER} & Reconnaissance d'Entités Nommées (\textit{Named Entity Recognition}) \\
\hline
\textbf{NLP} & Traitement Automatique du Langage Naturel (\textit{Natural Language Processing}) \\
\hline
\textbf{ORM} & Mapping Objet-Relationnel (\textit{Object-Relational Mapping}) \\
\hline
\textbf{PCM} & Modèle de Communication de Processus (\textit{Process Communication Model}) \\
\hline
\textbf{PDF} & Format de Document Portable (\textit{Portable Document Format}) \\
\hline
\textbf{QA} & Assurance Qualité (\textit{Quality Assurance}) \\
\hline
\textbf{RBAC} & Contrôle d'Accès Basé sur les Rôles (\textit{Role-Based Access Control}) \\
\hline
\textbf{REST} & Transfert d'État Représentationnel (\textit{Representational State Transfer}) \\
\hline
\textbf{RGPD} & Règlement Général sur la Protection des Données \\
\hline
\textbf{RH} & Ressources Humaines \\
\hline
\textbf{SaaS} & Logiciel en tant que Service (\textit{Software as a Service}) \\
\hline
\textbf{SGBD} & Système de Gestion de Base de Données \\
\hline
\textbf{SIRH} & Système d'Information des Ressources Humaines \\
\hline
\textbf{SMTP} & Protocole Simple de Transfert de Courrier (\textit{Simple Mail Transfer Protocol}) \\
\hline
\textbf{SPA} & Application à Page Unique (\textit{Single Page Application}) \\
\hline
\textbf{SQL} & Langage de Requête Structuré (\textit{Structured Query Language}) \\
\hline
\textbf{SRP} & Principe de Responsabilité Unique (\textit{Single Responsibility Principle}) \\
\hline
\textbf{SSE} & Événements Envoyés par le Serveur (\textit{Server-Sent Events}) \\
\hline
\textbf{SSR} & Rendu Côté Serveur (\textit{Server-Side Rendering}) \\
\hline
\textbf{TLS} & Sécurité de la Couche de Transport (\textit{Transport Layer Security}) \\
\hline
\textbf{UML} & Langage de Modélisation Unifié (\textit{Unified Modeling Language}) \\
\hline
\textbf{UUID} & Identifiant Universel Unique (\textit{Universally Unique Identifier}) \\
\hline
\textbf{XHTML} & Langage de Balisage Hypertexte Extensible (\textit{Extensible HyperText Markup Language}) \\
\hline
\textbf{XSS} & Script Intersite (\textit{Cross-Site Scripting}) \\
\hline
\caption{Liste des abréviations utilisées dans le rapport}
\end{longtable}

\clearpage

% -------------------------------------------------------------------------
% CORPS DU DOCUMENT
% -------------------------------------------------------------------------
\pagenumbering{arabic}
\chapter*{Introduction Générale}
\addcontentsline{toc}{chapter}{Introduction Générale}

Dans un contexte professionnel marqué par la transformation numérique des processus métiers et par l’intégration croissante de l’intelligence artificielle dans les systèmes d’aide à la décision, la gestion des compétences constitue aujourd’hui un enjeu stratégique pour les organisations. Les entreprises doivent non seulement identifier les compétences disponibles en interne, mais aussi évaluer leur niveau de maîtrise, détecter les écarts par rapport aux besoins métier et proposer des parcours de développement adaptés. Or, les méthodes classiques d’évaluation des talents, souvent fondées sur des CV déclaratifs, des entretiens ponctuels ou des appréciations subjectives, montrent leurs limites lorsqu’il s’agit de produire une analyse objective, continue et exploitable à grande échelle.

C’est dans ce cadre que s’inscrit le présent projet de fin d’études, réalisé au sein de la société AgileX Management et intitulé TalentPredict AI. Ce projet vise à concevoir et développer une plateforme intelligente dédiée à l’évaluation automatisée des compétences humaines, en prenant en compte à la fois les compétences techniques et les compétences comportementales. La problématique centrale peut ainsi être formulée comme suit : comment mettre en place une solution logicielle capable d’analyser de manière structurée, fiable et sécurisée les profils des collaborateurs, tout en générant des recommandations de formation personnalisées et utiles à la prise de décision RH ?

La solution proposée repose sur une approche hybride combinant plusieurs sources d’information et plusieurs traitements complémentaires. L’évaluation des compétences techniques s’appuie sur l’analyse du CV, l’exploitation du profil GitHub, des tests QCM adaptatifs et des challenges de code. L’évaluation des compétences comportementales mobilise le modèle psychométrique PCM, l’analyse sémantique locale à l’aide d’Ollama, ainsi que des scénarios d’évaluation générés par l’intelligence artificielle. Ces traitements sont intégrés dans une architecture moderne articulée autour de n8n pour l’orchestration des workflows, FastAPI pour les services d’intelligence artificielle, Spring Boot pour le backend métier, Angular pour l’interface utilisateur et Docker pour la conteneurisation. Ce choix architectural vise à assurer la modularité de la plateforme, la portabilité de l’environnement, la souveraineté des données et une meilleure conformité aux exigences de confidentialité, notamment celles liées au RGPD.

Ce mémoire est structuré en six chapitres. Le premier chapitre présente le cadre général du projet, l’organisme d’accueil, la problématique métier, les objectifs, les besoins fonctionnels et non fonctionnels, ainsi que la méthodologie Agile Scrum adoptée. Le deuxième chapitre propose un état de l’art des concepts, technologies et solutions existantes liés à l’évaluation intelligente des compétences et au domaine de la HR Tech. Les chapitres suivants décrivent le déroulement du projet selon les sprints de développement : le Sprint 1 est consacré à la mise en place de l’architecture, de la sécurité et du système d’authentification ; le Sprint 2 porte sur le développement du moteur d’évaluation multi-sources des compétences ; le Sprint 3 traite du système de recommandation et d’upskilling piloté par l’intelligence artificielle ; enfin, le Sprint 4 présente les tableaux de bord analytiques, le pilotage stratégique et la mise en production de la plateforme.
\clearpage

% =========================================================
% Chapitre 1
% =========================================================

\chapter{Cadre général du projet}

\section{Introduction}
\phantomsection
\addcontentsline{toc}{section}{Introduction}
\phantomsection 

Ce chapitre présente le cadre général dans lequel s’inscrit le projet TalentPredict AI. Il introduit d’abord l’organisme d’accueil et le contexte métier ayant motivé la conception de la solution. Il précise ensuite la problématique, les objectifs, les besoins fonctionnels et non fonctionnels, ainsi que les acteurs impliqués. Enfin, il décrit la méthodologie Scrum adoptée, l’organisation de l’équipe projet et les principaux éléments de planification qui structurent le déroulement du travail.

\section{Présentation du contexte et de l’enjeu}

La compréhension du contexte organisationnel constitue une étape essentielle pour situer le projet dans son environnement réel. Cette section présente l’entreprise d’accueil, ses domaines d’expertise et les enjeux métiers auxquels le projet tente d’apporter une réponse. Elle permet également de montrer la cohérence entre les besoins de l’entreprise et la solution proposée.

\subsection{Présentation de l’organisme d’accueil}

L’organisme d’accueil joue un rôle important dans l’orientation fonctionnelle et technique du projet. Sa culture de travail, ses pratiques de développement et ses priorités stratégiques influencent directement les choix retenus durant la conception et la réalisation de la plateforme.

\subsubsection{Présentation générale d’AgileX}

AgilEX est une société spécialisée dans le développement de solutions logicielles sur mesure et dans l’accompagnement des entreprises dans leur transformation numérique. Elle intervient principalement dans la conception d’applications web et mobiles, l’intégration de systèmes, le conseil en architecture logicielle et le déploiement de solutions cloud.

Le nom AgileX reflète une orientation vers l’agilité, l’adaptation et l’excellence technique. L’entreprise privilégie des approches de développement itératives permettant de produire rapidement des incréments fonctionnels, de réduire les risques et d’intégrer progressivement les retours des parties prenantes.

La figure~\ref{fig:logo_orga} présente le logo de l’entreprise d’accueil. Elle permet d’identifier visuellement l’organisme dans lequel le projet a été réalisé.

\begin{figure}[H]
  \centering
  \includegraphics[width=0.3\textwidth]{images/logo_orga.png}
  \caption{Logo de l’organisme d’accueil AgileX }
  \label{fig:logo_orga}
\end{figure}


La culture de travail d’Agilx repose sur une organisation AgileX centrée sur la collaboration, la livraison continue de valeur et l’amélioration progressive des solutions développées. Cette culture est en adéquation avec le choix de la méthodologie Scrum retenue pour piloter le projet.

La figure~\ref{fig:struc_orga} illustre l’organisation structurelle de l’entreprise. Elle met en évidence les relations entre les pôles de travail et permet de mieux comprendre le contexte professionnel dans lequel s’inscrit le projet.

\begin{figure}[H]
  \centering
  \includegraphics[width=0.75\textwidth]{images/philosophie_orga.png}
  \caption{Organisation structurelle de l’entreprise AgileX}
  \label{fig:struc_orga}
\end{figure}


Le tableau~\ref{tab:fiche_signal_agilx} présente une fiche signalétique synthétique de l’entreprise. Il regroupe les informations principales permettant de situer l’organisme d’accueil du point de vue administratif, sectoriel et méthodologique.

\begin{longtable}{|P{0.32\textwidth}|P{0.58\textwidth}|}
\label{tab:fiche_signal_agilx}\\
\hline
\rowcolor{TableHeaderBlue}
\color{white}\color{white}\textbf{Élément} & \color{white}\color{white}\textbf{Information} \\
\hline
\endfirsthead

\hline
\rowcolor{TableHeaderBlue}
\color{white}\color{white}\textbf{Élément} & \color{white}\color{white}\textbf{Information} \\
\hline
\endhead

\hline
\endfoot

Raison sociale & AgileX \\
\hline
Secteur d’activité & Développement de logiciels et solutions digitales \\
\hline
Domaines d’intervention & Applications web et mobiles, architecture logicielle, intelligence artificielle, cloud et intégration de systèmes \\
\hline
Méthodologie de travail & Approche Agile, principalement inspirée de Scrum~\cite{scrum2020} \\
\hline
Encadrant professionnel & M. Mohamed Amine Gharbi, Manager SI \\
\hline
\caption{Fiche signalétique de l’entreprise AgileX}
\end{longtable}

\subsubsection{Portefeuille de services et domaines d’expertise}

L’offre de services d’AgileX s’articule autour de plusieurs domaines complémentaires. Cette diversité d’expertise a constitué un environnement favorable pour la réalisation de TalentPredict AI   , dans la mesure où le projet combine développement web, architecture logicielle, intelligence artificielle, orchestration de workflows et déploiement conteneurisé.

Le tableau~\ref{tab:portefeuille_agilx} synthétise les principaux pôles d’expertise de l’entreprise et précise leur contribution au projet. Cette présentation permet de relier les compétences de l’organisme d’accueil aux choix techniques et fonctionnels adoptés dans la solution.

\begin{longtable}{|P{0.24\textwidth}|P{0.34\textwidth}|P{0.32\textwidth}|}

\label{tab:portefeuille_agilx}\\
\hline
\rowcolor{TableHeaderBlue}
\color{white}\textbf{Pôle d’expertise} &
\color{white}\textbf{Description} &
\color{white}\textbf{Contribution à \textit{TalentPredict AI}} \\
\hline
\endfirsthead

\hline
\rowcolor{TableHeaderBlue}
\color{white}\textbf{Pôle d’expertise} &
\color{white}\textbf{Description} &
\color{white}\textbf{Contribution à \textit{TalentPredict AI}} \\
\hline
\endhead

\hline
\endfoot

Développement Full-Stack et mobilité &
Conception et développement d’applications web modernes, d’applications mobiles et de plateformes SaaS. &
Mise en œuvre du frontend Angular, des interfaces utilisateurs et des interactions avec le backend. \\
\hline

Ingénierie Cloud et DevOps &
Mise en place d’environnements cloud, conteneurisation, automatisation et amélioration de la portabilité des applications. &
Utilisation de Docker pour faciliter le déploiement, l’isolation des services et la reproductibilité de l’environnement. \\
\hline

Conseil en architecture logicielle &
Définition d’architectures robustes, évolutives et adaptées aux besoins métiers. &
Adoption d’une architecture modulaire séparant l’interface, la logique métier, les services IA et l’orchestration. \\
\hline

Intégration d’intelligence artificielle &
Intégration de capacités d’analyse sémantique, d’automatisation et d’aide à la décision dans les processus métiers. &
Utilisation d’Ollama, de FastAPI et de n8n pour l’analyse des compétences et la génération de recommandations. \\
\hline
\caption{Portefeuille de services et contribution au projet}
\end{longtable}

\subsection{Enjeu général du projet}

Dans les organisations modernes, la gestion des compétences ne se limite plus à la consultation des CV ou à l’organisation d’entretiens ponctuels. Elle nécessite une vision dynamique, objective et exploitable des compétences réellement disponibles au sein de l’entreprise. Cette exigence devient particulièrement importante dans les métiers du numérique, où les technologies évoluent rapidement et où les besoins en formation doivent être identifiés de manière continue.

Le projet \textit{TalentPredict AI} répond à cet enjeu en proposant une plateforme capable de croiser plusieurs sources d’information afin d’évaluer les compétences techniques et comportementales des collaborateurs. L’objectif n’est pas de remplacer l’expertise humaine des responsables RH, mais de fournir un outil d’aide à la décision capable de structurer les analyses, de réduire les biais déclaratifs et d’orienter les parcours de développement professionnel.

\section{Présentation du projet}

Après la présentation de l’organisme d’accueil, il est nécessaire de préciser la nature du projet, ses objectifs et les acteurs concernés. Cette section permet de délimiter le périmètre fonctionnel de la solution et d’expliciter la valeur attendue pour les utilisateurs finaux et les décideurs RH.

\subsection{Contexte et problématique métier}

Dans le secteur informatique, l’évaluation des compétences constitue un enjeu stratégique pour les entreprises qui souhaitent mieux connaître les capacités réelles de leurs collaborateurs . Cependant, les informations nécessaires à cette évaluation sont souvent réparties entre plusieurs sources hétérogènes, telles que les CV, les profils GitHub, les profils LinkedIn, les portfolios, les tests techniques, les feedbacks RH et les questionnaires comportementaux. Cette dispersion rend le processus d’évaluation plus long, plus complexe et moins facilement exploitable dans une démarche de formation continue.

Comme l’illustre la figure~\ref{fig:problematique-metier-talentpredict}, le problème principal réside dans la fragmentation des données et dans l’absence d’un mécanisme unifié permettant de croiser les informations disponibles. Les évaluations traditionnelles restent souvent dépendantes des déclarations du collaborateur, ce qui peut introduire des biais ou limiter la fiabilité des résultats. De plus, les compétences comportementales sont généralement plus difficiles à mesurer, car elles nécessitent une analyse plus fine du profil, des interactions, des expériences et des réponses à des situations professionnelles.

\begin{figure}[H]
    \centering
    \includegraphics[width=\textwidth]{images/contexte.png}
    \caption{Problématique métier liée à la fragmentation des données d’évaluation des compétences}
    \label{fig:problematique-metier-talentpredict}
\end{figure}

Ces limites peuvent conduire à des décisions RH moins précises, à des parcours de formation peu personnalisés et à une exploitation insuffisante du potentiel interne. Dans ce contexte, le besoin d’une plateforme intelligente devient essentiel afin de passer d’une évaluation dispersée et déclarative à une évaluation unifiée, multi-sources et exploitable.

La problématique centrale du projet peut donc être formulée comme suit : comment concevoir et développer une plateforme intelligente permettant d’évaluer de manière structurée les compétences techniques et comportementales des collaborateurs, tout en générant des recommandations de formation personnalisées, sécurisées et exploitables par l’entreprise ?

\subsection{Objectifs du projet}

L’objectif principal du projet est de concevoir et développer TalentPredict AI, une plateforme intelligente dédiée à l’évaluation et au développement des compétences. La solution vise à fournir une analyse plus complète du profil collaborateur en combinant des données déclaratives, des preuves techniques observables et des traitements d’intelligence artificielle. Elle permet ainsi de renforcer l’objectivité du processus d’évaluation et de produire des résultats plus utiles à la prise de décision RH.

Les objectifs spécifiques du projet sont les suivants :

\begin{itemize}[leftmargin=0.9cm]
    \item analyser les compétences techniques à partir du CV, du profil GitHub, de tests QCM et de challenges de code ;
    \item évaluer les compétences comportementales à travers le modèle PCM, l’analyse sémantique locale et des scénarios générés par l’intelligence artificielle ;
    \item proposer des recommandations de formation adaptées aux écarts de compétences identifiés ;
    \item offrir un tableau de bord permettant aux employés et aux responsables RH de suivre les résultats et l’évolution des compétences ;
    \item garantir un traitement sécurisé des données personnelles et professionnelles, notamment à travers une inférence locale et une architecture maîtrisée.
\end{itemize}

Ces objectifs permettent à la plateforme de ne pas se limiter à une simple notation des compétences. Ils orientent plutôt la solution vers une logique d’accompagnement, dans laquelle l’évaluation devient un point de départ pour construire des parcours de progression individualisés et alignés avec les besoins de l’entreprise.

\subsection{Acteurs concernés}

Le projet implique plusieurs catégories d’acteurs, ayant chacune des besoins, des responsabilités et des interactions spécifiques avec la plateforme. L’identification de ces acteurs est indispensable pour définir les fonctionnalités attendues, organiser les droits d’accès et assurer une communication claire entre les utilisateurs humains et les services automatisés.

Les principaux acteurs concernés sont l’employé, l’administrateur RH et les services automatisés intégrés à la plateforme. L’employé utilise la solution pour soumettre ses données, évaluer ses compétences, consulter ses résultats et suivre ses parcours de formation. L’administrateur RH supervise les profils, valide certaines informations, consulte les indicateurs globaux et exploite les rapports générés pour orienter les décisions de gestion des talents. Les services automatisés assurent l’analyse intelligente, l’orchestration des traitements, l’envoi des notifications et la communication en temps réel.

La figure~\ref{fig:acteurs-talentpredict} présente une vue synthétique des acteurs concernés et de leurs interactions avec \textit{TalentPredict AI}. Elle met en évidence le rôle central de la plateforme, qui coordonne les échanges entre les utilisateurs, les services d’intelligence artificielle, le service mail et le mécanisme de communication temps réel.

\begin{figure}[H]
    \centering
    \includegraphics[width=\textwidth]{images/acteurs.png}
    \caption{Acteurs concernés et interactions principales autour de la plateforme TalentPredict AI}
    \label{fig:acteurs-talentpredict}
\end{figure}

Cette organisation permet de mieux structurer les responsabilités de chaque acteur dans le système. Elle garantit également une séparation claire entre les actions humaines, comme la consultation des résultats ou la supervision RH, et les traitements automatisés, comme l’analyse des données, la génération des recommandations et l’envoi des notifications.

\section{Spécification des besoins}

La spécification des besoins permet de traduire la problématique métier en exigences concrètes. Elle distingue les besoins fonctionnels, qui décrivent les services attendus par les utilisateurs, et les besoins non fonctionnels, qui définissent les contraintes de qualité, de sécurité, de performance et d’évolutivité du système.

\subsection{Besoins fonctionnels}

Les besoins fonctionnels définissent les principales fonctionnalités attendues de la plateforme. Ils couvrent l’authentification, l’analyse des compétences, la recommandation de formation, le suivi des parcours et la restitution des résultats à travers des tableaux de bord.

\subsubsection{Module d’authentification et de gestion des utilisateurs}

Ce module constitue le point d’entrée sécurisé de la plateforme. Il permet aux utilisateurs de créer un compte, de se connecter, de gérer leur profil et d’accéder aux fonctionnalités correspondant à leur rôle. Il repose sur une authentification par jeton JWT et sur une gestion différenciée des rôles USER et ADMIN.

Les fonctionnalités attendues sont les suivantes :

\begin{itemize}[leftmargin=0.9cm]
    \item inscription avec adresse email et mot de passe ;
    \item connexion sécurisée avec génération d’un jeton JWT ;
    \item modification des informations du profil utilisateur ;
    \item gestion des rôles par l’administrateur ;
    \item récupération de mot de passe à travers un processus contrôlé.
\end{itemize}

\subsubsection{Module d’analyse des compétences comportementales}

Ce module vise à analyser les dimensions comportementales du collaborateur en combinant un modèle psychométrique, des données textuelles et une analyse sémantique locale. Le modèle PCM~\cite{kahler1996} est utilisé comme cadre de référence pour structurer les profils comportementaux, tandis qu’Ollama~\cite{ollama}, n8n~\cite{n8n} et FastAPI~\cite{fastapi} interviennent dans l’orchestration et l’analyse des données.

Les fonctionnalités attendues sont les suivantes :

\begin{itemize}[leftmargin=0.9cm]
    \item extraction sécurisée du texte des CV côté client ;
    \item analyse du profil GitHub afin d’identifier certains indicateurs comportementaux observables ;
    \item passation d’un test PCM ;
    \item orchestration d’un flux d’analyse à travers n8n ;
    \item calcul de scores liés à la communication, au leadership, à la collaboration, à la curiosité, à l’ownership et à la discipline ;
    \item identification du type de personnalité dominant ;
    \item génération de scénarios d’évaluation comportementale ;
    \item historisation des résultats afin de suivre l’évolution du profil.
\end{itemize}

\subsubsection{Module d’analyse des compétences techniques}

Ce module a pour objectif d’évaluer les compétences techniques du collaborateur à partir de plusieurs sources complémentaires. L’intérêt de cette approche multi-sources est de réduire la dépendance aux déclarations du CV et de confronter les compétences annoncées à des preuves techniques plus observables.

Les fonctionnalités attendues sont les suivantes :

\begin{itemize}[leftmargin=0.9cm]
    \item extraction et parsing du CV~\cite{celik2013,chien2020} ;
    \item analyse de l’activité GitHub~\cite{kalliamvakou2014} et des langages utilisés ;
    \item prise en compte éventuelle du profil LinkedIn ou d’un portfolio ;
    \item génération de QCM techniques adaptatifs à l’aide d’Ollama~\cite{ollama} ;
    \item génération et évaluation automatique de challenges de code avec Code Llama~\cite{roziere2023codellama} ;
    \item calcul d’un score de confiance par fusion multi-sources ;
    \item détection des incohérences entre compétences déclarées et compétences observées ;
    \item attribution d’un niveau de maîtrise par compétence.
\end{itemize}

\subsubsection{Module de recommandations et de développement des compétences}

Ce module transforme les résultats de l’évaluation en plans d’action exploitables. Il ne se limite pas à produire un score, mais propose des parcours de développement adaptés aux écarts de compétences identifiés et au profil du collaborateur.

Pour l’employé, le module permet de consulter des recommandations de formation, de suivre une roadmap d’apprentissage, de visualiser sa progression et de passer des mini-tests de validation. Pour l’administrateur RH, il permet de suivre les parcours, de valider ou refuser certaines demandes, de proposer des formations ciblées et d’observer les taux d’avancement.

\begin{itemize}[leftmargin=0.9cm]
    \item génération de recommandations personnalisées de formation~\cite{drachsler2009} ;
    \item construction d’une roadmap d’apprentissage ;
    \item suivi des statuts des formations ;
    \item intégration de mécanismes de progression et de gamification ;
    \item validation des acquis par mini-test ;
    \item génération d’un certificat interne lorsque les conditions sont satisfaites ;
    \item supervision des parcours par l’administrateur RH.
\end{itemize}

\subsubsection{Module de tableau de bord analytique}

Le tableau de bord constitue l’espace de restitution des résultats. Il permet de présenter les analyses sous une forme lisible, synthétique et exploitable. Les visualisations doivent être adaptées à la fois aux employés, qui consultent leurs propres résultats, et aux administrateurs RH, qui suivent les compétences à l’échelle de l’organisation.

Les fonctionnalités attendues sont les suivantes :

\begin{itemize}[leftmargin=0.9cm]
    \item visualisation des compétences à travers des radars, histogrammes et courbes d’évolution ;
    \item consultation du profil comportemental et des résultats techniques ;
    \item suivi de l’avancement des formations ;
    \item génération d’un Talent Passport individuel ;
    \item cartographie globale des compétences ;
    \item identification des écarts de compétences ;
    \item génération de rapports RH consolidés.
\end{itemize}

\subsection{Besoins non fonctionnels}

Les besoins non fonctionnels définissent les exigences de qualité qui encadrent le fonctionnement de la plateforme. Ils garantissent que la solution ne soit pas seulement fonctionnelle, mais également sécurisée, maintenable, performante et adaptée à une utilisation réelle en entreprise.

Les principales exigences non fonctionnelles sont résumées dans le tableau~\ref{tab:besoins_non_fonctionnels}. Elles couvrent les dimensions de performance, sécurité, disponibilité, maintenabilité, ergonomie, interopérabilité et confidentialité.

\begin{longtable}{|P{0.24\textwidth}|P{0.66\textwidth}|}
\label{tab:besoins_non_fonctionnels}\\
\hline
\rowcolor{TableHeaderBlue}
\color{white}\color{white}\textbf{Critère} & \color{white}\color{white}\textbf{Exigence} \\
\hline
\endfirsthead
\hline
\rowcolor{TableHeaderBlue}
\color{white}\color{white}\textbf{Critère} & \color{white}\color{white}\textbf{Exigence} \\
\hline
\endhead
\hline
\endfoot
Performance &
Les opérations courantes de l’API doivent répondre dans un délai raisonnable. Les traitements LLM, plus coûteux, doivent être accompagnés d’un retour visuel afin d’informer l’utilisateur de l’avancement. \\
\hline
Sécurité &
L’authentification doit reposer sur \texttt{JWT}, les mots de passe doivent être hachés avec \texttt{BCrypt}, et les accès doivent être contrôlés selon les rôles. Les protections contre les injections SQL, XSS et CSRF doivent être prises en compte. \\
\hline
Disponibilité et fiabilité &
La plateforme doit viser une disponibilité élevée, prévoir des mécanismes de reprise en cas d’erreur et limiter les interruptions lors des appels aux services externes ou aux modules IA. \\
\hline
Maintenabilité &
L’architecture doit être modulaire, documentée et organisée selon des conventions de codage claires pour Java, Python et Angular. \\
\hline
Évolutivité &
La solution doit pouvoir intégrer de nouveaux modèles IA, de nouvelles sources de données et de nouveaux modules fonctionnels sans remise en cause majeure de l’architecture. \\
\hline
Ergonomie et accessibilité &
L’interface doit être intuitive, responsive et adaptée à des utilisateurs non techniques. Les visualisations doivent être lisibles, cohérentes et accompagnées de légendes explicites. \\
\hline
Interopérabilité &
Les échanges entre services doivent reposer sur des API REST et des formats structurés, principalement JSON. Les appels aux modèles IA doivent être encapsulés pour faciliter un éventuel changement de fournisseur. \\
\hline
Confidentialité et conformité &
Le système doit limiter l’exposition des données sensibles, favoriser le traitement local lorsque cela est possible, recueillir le consentement explicite des utilisateurs et respecter les principes du RGPD~\cite{rgpd2016}. \\
\hline
\caption{Synthèse des besoins non fonctionnels du projet}
\end{longtable}
\subsection{Identification des utilisateurs et des rôles}

L’identification des utilisateurs permet de définir les responsabilités, les droits d’accès et les interactions attendues avec la plateforme. Dans le cadre de \textit{TalentPredict AI}, les acteurs ne sont pas uniquement humains : certains services automatisés jouent également un rôle central dans le fonctionnement global du système.

Le tableau~\ref{tab:acteurs_systeme} présente les principaux acteurs du système, leur rôle et leurs interactions clés.

\begin{longtable}{|P{0.18\textwidth}|P{0.20\textwidth}|P{0.27\textwidth}|P{0.25\textwidth}|}
\label{tab:acteurs_systeme}\\
\hline
\rowcolor{TableHeaderBlue}
\color{white}\textbf{Acteur} &
\color{white}\textbf{Rôle} &
\color{white}\textbf{Objectif principal} &
\color{white}\textbf{Interactions clés} \\
\hline
\endfirsthead

\hline
\rowcolor{TableHeaderBlue}
\color{white}\textbf{Acteur} &
\color{white}\textbf{Rôle} &
\color{white}\textbf{Objectif principal} &
\color{white}\textbf{Interactions clés} \\
\hline
\endhead

\hline
\endfoot

Employé &
Utilisateur principal &
Évaluer ses compétences, consulter ses résultats et suivre ses recommandations de formation. &
Soumission du CV, ajout du profil GitHub, passation des tests, consultation du tableau de bord. \\
\hline

Administrateur RH &
Administrateur métier &
Superviser les profils, analyser les résultats et suivre les parcours de développement. &
Validation des compétences, affectation des formations, consultation des rapports et indicateurs globaux. \\
\hline

Système IA &
Acteur automatisé d’analyse &
Analyser les données hétérogènes et produire des résultats exploitables par la plateforme. &
Inférence locale, analyse sémantique, génération de scénarios et recommandations. \\
\hline

Communication Hub &
Acteur automatisé temps réel &
Assurer la communication instantanée entre les traitements backend et l’interface utilisateur. &
Émission d’événements SSE et notification de l’état d’avancement des analyses. \\
\hline

Service Mail &
Acteur automatisé asynchrone &
Envoyer les notifications importantes et assurer la traçabilité de certaines actions. &
Envoi d’emails de confirmation, de récupération de mot de passe, de certificats et d’alertes. \\
\hline
\caption{Acteurs du système et interactions principales}
\end{longtable}

\section{Méthodologie de développement Scrum}

Le projet a été conduit selon une approche Agile afin de tenir compte de la complexité fonctionnelle et technique de la solution. La méthodologie Scrum a été retenue car elle permet une progression incrémentale, une adaptation continue aux retours des parties prenantes et une meilleure maîtrise des risques liés à l’intégration de technologies émergentes.

La figuree~\ref{fig:organisation-scrum} illustre de manière synthétique l'organisation globale de la méthode Scrum adoptée pour la réalisation de ce projet TalentPredict AI. Elle met en évidence la répartition des rôles, les artefacts et cérémonies mobilisés, ainsi que le découpage fonctionnel des quatre sprints.
\begin{figure}[H]
    \centering
    \includegraphics[width=\textwidth]{images/Organisation Scrum du projet.png}
    \caption{Organisation Scrum du projet}
    \label{fig:organisation-scrum}
\end{figure}
\subsection{Choix et justification de Scrum}

Le choix de Scrum se justifie par la nécessité de produire des incréments fonctionnels réguliers, de valider progressivement les modules développés et de maintenir une communication continue entre les membres de l’équipe. Dans le cadre de TalentPredict AI, cette méthodologie est particulièrement adaptée, car le projet combine des fonctionnalités classiques de développement logiciel avec des composants plus expérimentaux liés à l’intelligence artificielle.

Scrum a également permis d’organiser le travail autour de sprints successifs, chacun ayant un objectif clair et des livrables vérifiables. Cette organisation réduit le risque d’accumulation d’erreurs en fin de projet et favorise une amélioration progressive de la qualité technique.

\subsection{Principes fondamentaux appliqués}

Les principes Scrum appliqués durant le projet reposent sur l’itération, la collaboration, la priorisation et l’inspection continue. Ces principes ont guidé la planification, le développement, les revues et les ajustements réalisés au fil des sprints.

\begin{itemize}[leftmargin=0.9cm]
    \item organisation du travail en itérations courtes ;
    \item priorisation des fonctionnalités selon leur valeur métier ;
    \item validation régulière avec le Product Owner ;
    \item adaptation progressive du backlog selon les retours obtenus ;
    \item intégration continue des développements réalisés ;
    \item identification proactive des risques techniques.
\end{itemize}

\subsection{Rôles et artefacts Scrum}

La mise en œuvre de Scrum nécessite une clarification des rôles et des responsabilités. Dans ce projet, les rôles Scrum ont été adaptés au contexte académique et professionnel du PFE, tout en respectant la logique générale du cadre méthodologique.

Le tableau~\ref{tab:roles_scrum} présente les rôles identifiés au sein de l’équipe projet.

\begin{longtable}{|P{0.24\textwidth}|P{0.24\textwidth}|P{0.42\textwidth}|}
\hline
\rowcolor{TableHeaderBlue}
\color{white}\textbf{Rôle Scrum} &
\color{white}\textbf{Acteur} &
\color{white}\textbf{Responsabilités principales} \\
\hline
\endfirsthead

\hline
\rowcolor{TableHeaderBlue}
\color{white}\textbf{Rôle Scrum} &
\color{white}\textbf{Acteur} &
\color{white}\textbf{Responsabilités principales} \\
\hline
\endhead

\hline
\endfoot

Product Owner &
M. Mohamed Amine Gharbi &
Définition de la vision produit, priorisation du Product Backlog, clarification des besoins et validation des livrables. \\
\hline

Équipe de développement &
Rahma Barouni et Mohamed Abdeladhim &
Conception, développement, test et intégration des modules fonctionnels, techniques et IA de la plateforme. \\
\hline

Scrum Master  &
Mme. Takwa Ben Aicha &
Suivi méthodologique, validation de la cohérence scientifique et accompagnement dans la structuration du rapport. \\
\hline
\caption{Rôles Scrum identifiés dans le projet}
\label{tab:roles_scrum}\\
\end{longtable}

Les artefacts Scrum utilisés structurent le suivi du projet et assurent la traçabilité des fonctionnalités. Le tableau~\ref{tab:artefacts_scrum} résume les principaux artefacts mobilisés.

\begin{longtable}{|P{0.26\textwidth}|P{0.46\textwidth}|P{0.18\textwidth}|}
\hline
\rowcolor{TableHeaderBlue}
\color{white}\textbf{Artefact} &
\color{white}\textbf{Description} &
\color{white}\textbf{Fréquence d’utilisation} \\
\hline
\endfirsthead

\hline
\rowcolor{TableHeaderBlue}
\color{white}\textbf{Artefact} &
\color{white}\textbf{Description} &
\color{white}\textbf{Fréquence d’utilisation} \\
\hline
\endhead

\hline
\endfoot

Product Backlog &
Liste priorisée des besoins, fonctionnalités, user stories et améliorations à réaliser. &
Mise à jour continue \\
\hline

Sprint Backlog &
Sous-ensemble des user stories sélectionnées pour un sprint donné, avec les tâches associées. &
À chaque sprint \\
\hline

Incrément &
Version fonctionnelle, testée et potentiellement livrable de la plateforme. &
Fin de sprint \\
\hline

Definition of Ready &
Ensemble de critères indiquant qu’une user story est suffisamment claire pour être développée. &
Avant planification \\
\hline

Definition of Done &
Ensemble de critères permettant de considérer une fonctionnalité comme terminée. &
Fin de tâche ou de sprint \\
\hline
\caption{Artefacts Scrum utilisés}
\label{tab:artefacts_scrum}\\
\end{longtable}

\subsection{Organisation des sprints}

Le projet est organisé en quatre sprints principaux, chacun correspondant à une étape cohérente du développement. Cette structuration permet de construire progressivement la plateforme, depuis les fondations techniques jusqu’aux fonctionnalités avancées d’analyse, de recommandation et de pilotage.

Le tableau~\ref{tab:recap_sprints_ch1} présente la planification globale des sprints.

\begin{longtable}{|P{0.15\textwidth}|P{0.15\textwidth}|P{0.30\textwidth}|P{0.30\textwidth}|}
\hline
\rowcolor{TableHeaderBlue}
\color{white}\textbf{Sprint} &
\color{white}\textbf{Période} &
\color{white}\textbf{Objectif principal} &
\color{white}\textbf{Livrables clés} \\
\hline
\endfirsthead

\hline
\rowcolor{TableHeaderBlue}
\color{white}\textbf{Sprint} &
\color{white}\textbf{Période} &
\color{white}\textbf{Objectif principal} &
\color{white}\textbf{Livrables clés} \\
\hline
\endhead

\hline
\endfoot

Sprint 1 &
Semaines 1 à 4 &
Mise en place de l’architecture, de la sécurité et de l’authentification. &
Configuration Docker~\cite{docker}, backend sécurisé, authentification JWT, premières interfaces Angular. \\
\hline

Sprint 2 &
Semaines 5 à 8 &
Développement des moteurs d’analyse des compétences. &
Workflows n8n~\cite{n8n}, extraction de texte, analyse GitHub~\cite{kalliamvakou2014}, parsing CV et premiers scores. \\
\hline

Sprint 3 &
Semaines 9 à 12 &
Intégration de l’IA, des tests dynamiques et du moteur de recommandation. &
Ollama~\cite{ollama}, Code Llama~\cite{roziere2023codellama}, QCM adaptatifs, challenges de code et recommandations. \\
\hline

Sprint 4 &
Semaines 13 à 16 &
Finalisation des tableaux de bord, tests et déploiement. &
Dashboard Angular~\cite{angular}, rapports PDF, tests d’intégration, déploiement final et validation globale. \\
\hline
\caption{Récapitulatif de la planification des sprints}
\label{tab:recap_sprints_ch1}\\
\end{longtable}

\section{Constitution de l’équipe projet}

La constitution de l’équipe projet permet d’identifier clairement les responsabilités de chaque intervenant. Cette clarification facilite la coordination des tâches, la gestion des dépendances et la validation progressive des livrables.

Le tableau~\ref{tab:equipe_projet_ch1} présente les membres de l’équipe et leurs responsabilités principales dans la réalisation de \textit{TalentPredict AI}.

\begin{longtable}{|P{0.28\textwidth}|P{0.24\textwidth}|P{0.38\textwidth}|}
\hline
\rowcolor{TableHeaderBlue}
\color{white}\textbf{Membre} &
\color{white}\textbf{Rôle} &
\color{white}\textbf{Responsabilités principales} \\
\hline
\endfirsthead

\hline
\rowcolor{TableHeaderBlue}
\color{white}\textbf{Membre} &
\color{white}\textbf{Rôle} &
\color{white}\textbf{Responsabilités principales} \\
\hline
\endhead

\hline
\endfoot

M. Mohamed Amine Gharbi &
Product Owner et encadrant professionnel &
Définition des besoins, priorisation des fonctionnalités, validation des sprints et suivi professionnel du projet. \\
\hline

Takwa Ben Aicha &
Scrum Master &
Garantie du respect du cadre méthodologique Agile Scrum, animation des cérémonies (Sprint Planning, Daily Stand-up, Sprint Review, Retrospective), facilitation de la communication et élimination des obstacles entravant l'équipe. \\
\hline

Rahma Barouni &
Développeuse Full-Stack &
Développement du module des compétences comportementales, workflows n8n, analyse PCM, intégration d’Ollama et traitements via FastAPI~\cite{fastapi}. \\
\hline

Mohamed Abdeladhim &
Développeur Full-Stack &
Développement du module des compétences techniques, parsing de CV, intégration GitHub et LinkedIn, QCM adaptatifs, évaluation de code et persistance des résultats. \\
\hline
\caption{Constitution de l’équipe projet}
\label{tab:equipe_projet_ch1}\\
\end{longtable}

\section{Solution proposée : TalentPredict AI}

La solution proposée vise à répondre aux limites identifiées dans les processus classiques d’évaluation des compétences. Elle repose sur une approche hybride combinant des sources de données hétérogènes, des traitements d’intelligence artificielle et des mécanismes de restitution adaptés aux besoins des employés et des responsables RH.

\subsection{Description générale de la solution}

TalentPredict AI est une plateforme d’évaluation et de développement des compétences destinée à accompagner les entreprises dans l’analyse de leurs talents. Elle permet de collecter des données issues du CV, du profil GitHub, des tests techniques, des questionnaires comportementaux et des résultats d’évaluation, puis de les transformer en indicateurs exploitables.

La valeur ajoutée de la solution repose sur trois axes principaux. Le premier est l’analyse duale des compétences, qui combine les dimensions techniques et comportementales. Le deuxième est la validation multi-sources, qui permet de réduire la dépendance aux informations déclaratives. Le troisième est la souveraineté des données, rendue possible par l’utilisation de traitements locaux et d’une architecture contrôlée.

\subsection{Technologies de la plateforme}

La plateforme repose sur un ensemble de technologies complémentaires choisies pour leur maturité, leur compatibilité et leur adéquation avec les exigences du projet. Le tableau~\ref{tab:technologies_talentpredict_ch1} présente les principales technologies utilisées, ainsi que leur rôle dans l’architecture.

\begin{longtable}{|P{0.18\textwidth}|P{0.18\textwidth}|P{0.27\textwidth}|P{0.27\textwidth}|}
\hline
\rowcolor{TableHeaderBlue}
\color{white}\textbf{Technologie} &
\color{white}\textbf{Version indicative} &
\color{white}\textbf{Rôle dans le projet} &
\color{white}\textbf{Justification technique} \\
\hline
\endfirsthead

\hline
\rowcolor{TableHeaderBlue}
\color{white}\textbf{Technologie} &
\color{white}\textbf{Version indicative} &
\color{white}\textbf{Rôle dans le projet} &
\color{white}\textbf{Justification technique} \\
\hline
\endhead

\hline
\endfoot

Angular~\cite{angular} &
19+ &
Développement du frontend, des formulaires, des tableaux de bord et des visualisations. &
Architecture par composants, typage TypeScript, routage intégré et bonne intégration avec les interfaces interactives. \\
\hline

Spring Boot~\cite{springboot} / Java &
3.4.2 / 21 &
Backend métier, API REST, sécurité, persistance et génération de documents. &
Cadre robuste pour développer des services REST sécurisés, avec intégration native de Spring Security et JPA. \\
\hline

PostgreSQL~\cite{postgresql} &
15+ &
Stockage relationnel des utilisateurs, rôles, profils, compétences et résultats. &
SGBD fiable, conforme aux propriétés ACID et adapté à la gestion de données structurées et semi-structurées. \\
\hline

Python / FastAPI~\cite{fastapi} &
3.11 / 0.110+ &
Microservice IA pour l’analyse des CV, des profils, des QCM et des scénarios. &
Framework léger, asynchrone et adapté à la création rapide d’API pour les traitements IA. \\
\hline

n8n~\cite{n8n} &
latest &
Orchestration des workflows d’analyse et coordination des appels entre services. &
Permet de modéliser visuellement les flux, de modifier les workflows sans recompilation et de faciliter l’évolution des traitements. \\
\hline

Ollama~\cite{ollama} &
latest &
Exécution locale de modèles de langage pour l’analyse sémantique et la génération de recommandations. &
Permet une inférence locale, limitant l’externalisation des données sensibles et renforçant la souveraineté. \\
\hline

Code Llama~\cite{roziere2023codellama} &
latest &
Génération et évaluation de challenges de code. &
Modèle spécialisé dans le traitement du code, adapté aux tâches d’évaluation algorithmique. \\
\hline

Docker / Docker Compose~\cite{docker} &
24+ &
Conteneurisation et orchestration locale des services. &
Garantit la reproductibilité de l’environnement, l’isolation des composants et la portabilité du déploiement. \\
\hline

pdfjs-dist~\cite{pdfjs} &
5.6.205 &
Extraction de texte à partir de fichiers PDF côté client. &
Réduit l’exposition des CV bruts côté serveur et contribue à une logique de Privacy by Design. \\
\hline

Chart.js~\cite{chartjs} &
4+ &
Visualisation des scores, courbes d’évolution et radars de compétences. &
Bibliothèque légère, adaptée aux tableaux de bord et facilement intégrable avec Angular. \\
\hline
\caption{Technologies principales de TalentPredict AI}
\label{tab:technologies_talentpredict_ch1}\\
\end{longtable}

\subsection{Architecture générale de la solution}

L’architecture de TalentPredict AI est organisée autour de plusieurs composants spécialisés. Cette séparation des responsabilités permet de mieux maîtriser la complexité, de faciliter la maintenance et de préparer l’évolution future de la plateforme.

La figure présente l’architecture globale de la solution. Elle montre les interactions entre l’interface utilisateur, le backend métier, les services IA, l’orchestrateur de workflows, la base de données et les services de communication.
\begin{figure}[H]
    \centering
    \includegraphics[width=\textwidth]{images/archi Glob.png}
    \caption{Architecture globale de la plateforme TalentPredict AI}
    \label{fig:architecture-glob}
\end{figure}
\begin{figure}[H]
    \centering
    \includegraphics[width=\textwidth]{images/Architecture en couches2.png}
    \caption{Architecture globale de la plateforme TalentPredict AI}
    \label{fig:architecture-glob2}
\end{figure}

\subsection{Vue d’ensemble en couches}

L’architecture retenue adopte un modèle en couches afin de séparer la présentation, la logique métier, l’analyse intelligente, la persistance et le déploiement. Cette organisation améliore la lisibilité du système et permet de faire évoluer chaque couche de manière relativement indépendante.

La figure suivante présente une vue synthétique de l’architecture en couches. Elle permet de comprendre le rôle de chaque niveau et la circulation générale des données au sein de la plateforme.

\chapterfigure[\textwidth]{images/archi_couche.png}{Vue d’ensemble de l’architecture en couches de TalentPredict AI}

\section{Product Backlog}

Le Product Backlog regroupe les principales fonctionnalités attendues dans la première version de la plateforme. Il constitue une base de priorisation pour l’équipe Scrum et permet de transformer les besoins métiers en user stories exploitables dans les sprints.

Le tableau~\ref{tab:product_backlog_ch1} présente un extrait du Product Backlog initial, centré sur les fonctionnalités prioritaires du MVP.

\begin{longtable}{|P{0.10\textwidth}|P{0.15\textwidth}|P{0.15\textwidth}|P{0.18\textwidth}|P{0.22\textwidth}|P{0.10\textwidth}|}
\hline
\rowcolor{TableHeaderBlue}
\color{white}\textbf{ID} &
\color{white}\textbf{Epic} &
\color{white}\textbf{En tant que} &
\color{white}\textbf{Je souhaite} &
\color{white}\textbf{Afin de} &
\color{white}\textbf{Priorité} \\
\hline
\endfirsthead

\hline
\rowcolor{TableHeaderBlue}
\color{white}\textbf{ID} &
\color{white}\textbf{Epic} &
\color{white}\textbf{En tant que} &
\color{white}\textbf{Je souhaite} &
\color{white}\textbf{Afin de} &
\color{white}\textbf{Priorité} \\
\hline
\endhead

\hline
\endfoot

US-001 &
Authentification et sécurité &
Utilisateur / RH &
M’authentifier via un jeton JWT &
Accéder à l’application et à mes données de manière sécurisée &
Haute \\
\hline

US-002 &
Compétences comportementales &
Utilisateur &
Passer le test de personnalité PCM~\cite{kahler1996} &
Générer mon profil comportemental et identifier mon type dominant &
Haute \\
\hline

US-003 &
Compétences comportementales et IA &
Système IA &
Analyser sémantiquement les profils avec un modèle local &
Déduire des indicateurs comportementaux tout en préservant la confidentialité des données &
Haute \\
\hline

US-004 &
Analyse technique &
Utilisateur &
Transmettre mon CV et mon pseudonyme GitHub &
Permettre l’extraction automatique des données techniques &
Haute \\
\hline

US-005 &
Analyse technique &
Utilisateur &
Fournir l’URL de mon portfolio ou profil LinkedIn &
Enrichir l’évaluation par des signaux professionnels supplémentaires &
Haute \\
\hline

US-006 &
Scoring multi-sources &
Système d’évaluation &
Fusionner les résultats issus des différentes sources &
Attribuer un niveau de maîtrise et un score de confiance &
Haute \\
\hline

US-007 &
Recommandations &
Système de recommandation &
Croiser les compétences détectées avec les besoins du poste &
Proposer des parcours de formation pertinents &
Haute \\
\hline

US-008 &
Orchestration &
Développeur &
Modéliser les workflows d’analyse dans n8n &
Coordonner les appels entre frontend, backend et services IA &
Haute \\
\hline

US-009 &
Communication &
Administrateur RH &
Notifier les collaborateurs lors des étapes importantes &
Assurer un suivi opérationnel centralisé &
Haute \\
\hline
\caption{Product Backlog initial du projet}
\label{tab:product_backlog_ch1}\\
\end{longtable}

\section{Critères d’acceptation des user stories}

Les critères d’acceptation permettent de vérifier qu’une user story est suffisamment claire avant son développement et qu’elle peut être considérée comme terminée après réalisation. Ils contribuent à réduire les ambiguïtés, à améliorer la qualité des livrables et à faciliter la validation par le Product Owner.

\subsection{Definition of Ready}

Une user story est considérée comme prête lorsqu’elle peut être comprise, estimée et développée par l’équipe sans ambiguïté majeure. Dans ce projet, les critères de préparation suivants ont été appliqués :

\begin{itemize}[leftmargin=0.9cm]
    \item la user story est formulée selon le format « En tant que, je souhaite, afin de » ;
    \item les critères d’acceptation sont clairement définis ;
    \item les dépendances techniques sont identifiées ;
    \item les maquettes ou descriptions fonctionnelles nécessaires sont disponibles ;
    \item la priorité est définie dans le Product Backlog ;
    \item la story est validée par le Product Owner.
\end{itemize}

\subsection{Definition of Done}

Une fonctionnalité est considérée comme terminée lorsqu’elle respecte les exigences fonctionnelles, techniques et qualitatives définies par l’équipe. Ces critères assurent que l’incrément livré est stable, testable et potentiellement déployable.
\begin{itemize}[leftmargin=0.9cm]
    \item le code est implémenté, relu et intégré dans la branche principale ;
    \item les tests unitaires ou fonctionnels associés sont validés ;
    \item les API concernées sont documentées ;
    \item l’interface ne présente pas d’erreurs bloquantes ;
    \item la fonctionnalité respecte les règles de sécurité définies ;
    \item le livrable est validé par le Product Owner lors de la revue de sprint.
\end{itemize}
\section{Planification générale}

La planification générale permet de visualiser l’organisation globale du projet et les interactions attendues entre les acteurs et le système. Elle regroupe les représentations fonctionnelles et temporelles nécessaires à la compréhension du déroulement du projet.
\subsection{Diagramme des cas d’utilisation global}

Le diagramme des cas d’utilisation~\ref{fig:D-use-glob} offre une vue fonctionnelle d’ensemble de la plateforme. Il permet d’identifier les acteurs principaux, les services externes et les fonctionnalités majeures auxquelles chacun participe.

La figure~\ref{fig:D-use-glob} présente le diagramme global des cas d’utilisation de TalentPredict AI. Elle illustre les interactions entre l’employé, l’administrateur RH, les services IA, le service mail, le service SSE et les principales fonctionnalités de la plateforme.
\begin{figure}[H]
    \centering
    \includegraphics[width=1\textwidth]{images/D use case glob platfome.png}
    \caption{Diagramme global des cas d’utilisation de la plateforme TalentPredict AI}
    \label{fig:D-use-glob}
\end{figure}
\begin{figure}[H]
    \centering
    \includegraphics[width=1\textwidth]{images/use glob33.png}
    \caption{Diagramme global des cas d’utilisation de la plateforme TalentPredict AI}
    \label{fig:D-use-glob3}
\end{figure}

À noter que ce diagramme~\ref{fig:D-use-glob} se concentre uniquement sur les fonctionnalités de haut niveau. Les détails techniques internes, tels que les workflows d'intelligence artificielle ou les appels d'API spécifiques, ne sont pas représentés ici et seront détaillés dans les diagrammes d'architecture et de séquence.

\color{black}
\subsection{Diagramme de Gantt}

Le diagramme de Gantt permet de représenter la planification temporelle du projet. Il met en évidence l’enchaînement des sprints, les grandes phases de développement et la progression des livrables dans le temps.

La figure~\ref{fig:gantt-projet-ch1} l’emplacement prévu pour le diagramme de Gantt du projet. Elle devra représenter les quatre sprints principaux, leurs périodes respectives et les livrables associés.
\begin{figure}[H]
    \centering
    \includegraphics[width=\textwidth]{images/D gantt.png}
    \caption{Diagramme de Gantt du projet}
    \label{fig:gantt-projet-ch1}
\end{figure}
%\figureplaceholder{Diagramme de Gantt du projet}{fig:gantt-projet-ch1}

\section{Conclusion}
\phantomsection
\addcontentsline{toc}{section}{Conclusion}

Ce chapitre a permis de poser le cadre général du projet \textit{TalentPredict AI}. Il a présenté l’organisme d’accueil, la problématique métier, les objectifs, les besoins, les acteurs, la méthodologie Scrum et la planification générale. L’analyse menée montre que le projet répond à un besoin réel d’évaluation structurée, sécurisée et évolutive des compétences en entreprise. Les choix méthodologiques et architecturaux introduits dans ce chapitre constituent ainsi une base cohérente pour les chapitres suivants, qui détailleront l’état de l’art puis la réalisation progressive de la plateforme à travers les différents sprints.


\chapter{État de l'Art}

\textbf{Introduction}

Ce chapitre présente une revue systématique des travaux existants et des technologies fondatrices sur lesquels repose la plateforme TalentPredict AI. Le domaine de l'évaluation intelligente des compétences en entreprise se situe à l'intersection de plusieurs champs disciplinaires en évolution rapide : la gestion des ressources humaines augmentée par l'intelligence artificielle (HR Tech), le traitement automatique du langage naturel (NLP), les modèles psychométriques computationnels et les architectures logicielles distribuées orientées IA. L'objectif de ce chapitre n'est pas de dresser un inventaire exhaustif, mais d'analyser avec esprit critique les approches et solutions existantes afin d'identifier leurs limites structurelles et de justifier les choix conceptuels et techniques de TalentPredict AI.

Le chapitre s'organise en cinq sections. La première délimite le domaine et ses enjeux fondamentaux. La deuxième présente les concepts et technologies clés. La troisième analyse les travaux et solutions similaires. La quatrième propose une synthèse comparative. La cinquième positionne TalentPredict par rapport aux lacunes identifiées.

\section{Présentation du domaine}
L'optimisation du capital humain requiert aujourd'hui des outils capables d'évaluer le profil global des collaborateurs de manière objective et continue. Pour bien comprendre le positionnement de notre plateforme, nous explorerons dans cette section le contexte stratégique de la gestion des compétences internes, les dimensions complémentaires qui les composent, ainsi que les principaux verrous qu'un système d'évaluation automatisé par IA se doit de lever.
\subsection{La gestion des compétences internes : un enjeu stratégique majeur}

Dans une économie de la connaissance caractérisée par l'obsolescence rapide des compétences techniques et la raréfaction des profils spécialisés, les organisations ne peuvent plus se contenter d'une gestion réactive des talents fondée sur l'acquisition de ressources externes. La tendance de fond, documentée depuis les travaux précurseurs de Ulrich (1997) sur le HR Business Partner~\cite{ch2-ulrich1997}, est celle d'une gestion proactive et continue des compétences internes, désignée sous les termes d'Upskilling (montée en compétences sur le poste actuel) et de Reskilling (reconversion vers un nouveau métier).

Cette transition impose aux départements RH de disposer d'outils capables d'évaluer objectivement et en continu le niveau de compétence réel de chaque collaborateur  et non plus seulement leur niveau déclaré  afin d'identifier les écarts (skill gaps) entre les capacités actuelles et les exigences stratégiques de l'organisation. C'est précisément cette problématique que TalentPredict AI cherche à résoudre à travers un pipeline d'évaluation automatisé, multi-sources et piloté par l'intelligence artificielle.

\subsection{Les deux dimensions complémentaires de la compétence professionnelle}

La littérature en psychologie du travail et des organisations distingue classiquement deux grandes catégories de compétences, dont l'évaluation requiert des approches méthodologiques fondamentalement différentes.

Les compétences techniques (hard skills) désignent les savoirs opérationnels vérifiables et mesurables : maîtrise d'un langage de programmation, d'un framework, d'une plateforme cloud ou d'une méthodologie de développement logiciel. Leur évaluation peut s'appuyer sur des preuves objectives  résultats de tests techniques, analyse de code produit, historique de contributions réduisant ainsi la part de subjectivité.

Les compétences comportementales (soft skills) désignent les aptitudes transverses qui déterminent la manière dont un individu mobilise ses compétences techniques en contexte professionnel : capacité de communication, leadership, collaboration, gestion du stress, adaptabilité. Leur évaluation est intrinsèquement plus complexe car elle requiert l'observation de comportements en situation, difficilement capturables par un simple questionnaire auto-déclaratif~\cite{ch2-salgado1997}.

L'originalité de TalentPredict réside précisément dans l'intégration de ces deux dimensions au sein d'un même pipeline d'évaluation, permettant de produire un profil collaborateur à 360\textdegree{} combinant preuves objectives et analyse comportementale.

\subsection{Les défis de l'évaluation automatisée des compétences}

Trois défis fondamentaux caractérisent l'automatisation de l'évaluation des compétences en entreprise et constituent le fil conducteur de cette revue de l'état de l'art.

Le premier défi est celui du biais déclaratif : la quasi-totalité des systèmes d'évaluation existants reposent, en tout ou partie, sur ce que le collaborateur affirme de lui-même. Or, la psychologie sociale a établi de longue date que les auto-évaluations sont systématiquement biaisées par des effets de désirabilité sociale~\cite{ch2-donaldson2002} et de surestimation des compétences propres (Dunning-Kruger effect~\cite{ch2-kruger1999}).

Le second défi est celui de la fragmentation technique,comportementale : les solutions actuelles traitent séparément l'évaluation technique et l'évaluation comportementale, sans mécanisme de croisement permettant, par exemple, de vérifier la cohérence entre un profil psychométrique déclaré et des comportements observables dans des données asynchrones (historique de contributions, dynamiques de collaboration en équipe).

Le troisième défi est celui de la souveraineté des données : l'exploitation des grands modèles de langage (LLM) pour l'analyse sémantique des profils RH implique le traitement de données à caractère sensible (données psychologiques, données professionnelles personnelles), ce qui soulève des problématiques de conformité réglementaire, en particulier au regard du Règlement Général sur la Protection des Données (RGPD)~\cite{ch2-rgpd2016}.

\section{Concepts et technologies clés}
Pour répondre aux défis de la gestion des talents, notre plateforme s'appuie sur une convergence entre psychométrie éprouvée et intelligence artificielle de pointe. Cette section dresse un panorama des concepts clés du projet : elle débute par une analyse des méthodes d'évaluation traditionnelles pour ensuite présenter les ruptures technologiques (analyse sémantique, Edge AI, microservices) permettant d'automatiser et d'objectiver ce processus.
\subsection{Méthodes traditionnelles d'évaluation des compétences}

\textbf{L'évaluation managériale structurée:} L'échange d'évaluation (ou bilan de compétences face-à-face) constitue, depuis des décennies, l'outil central de la gestion des compétences en entreprise. Dans sa forme structurée  où les questions sont standardisées et les réponses cotées selon une grille prédéfinie  il présente une validité prédictive de la performance professionnelle significativement supérieure aux évaluations non structurées. McDaniel et al. (1994), dans une méta-analyse portant sur 114 études, ont établi une validité prédictive (r corrigé) de 0,51 pour les mises en situation structurées, contre 0,38 pour les échanges non structurés~\cite{ch2-mcdaniel1994}. Toutefois, ce format d'évaluation reste coûteux en temps RH, difficile à standardiser à grande échelle et sujet au biais de l'évaluateur (halo effect, biais de similarité).

\textbf{Les tests psychométriques standardisés:} Les outils psychométriques constituent le second pilier de l'évaluation comportementale traditionnelle. Parmi les modèles les plus utilisés en contexte organisationnel, on distingue :
\begin{itemize}[leftmargin=1.5cm]
  \item Le modèle des Big Five (OCEAN : Ouverture, Conscienciosité, Extraversion, Agréabilité, Névrosisme), issu de la psychologie différentielle et validé par plusieurs décennies de recherche empirique~\cite{ch2-costa1992}. Sa solidité scientifique est reconnue, mais son opérationnalisation en recommandations de formation concrètes reste limitée.
  \item Le modèle DISC (Dominance, Influence, Stabilité, Conformité), largement utilisé en développement organisationnel mais dont la validité scientifique est contestée ; plusieurs études ont mis en évidence une faible fidélité test-retest et une validité prédictive modeste~\cite{ch2-pittenger1993}.
  \item Le Process Communication Model (PCM), développé par le Dr Taibi Kahler dans les années 1970 et validé cliniquement dans le cadre du programme de sélection des astronautes de la NASA~\cite{ch2-kahler1996}. Le PCM structure la personnalité en six types fondamentaux (Empathique, Travaillomane, Persévérant, Rêveur, Rebelle, Promoteur), chacun associé à des canaux de communication préférentiels, des besoins psychologiques spécifiques et des comportements prévisibles sous stress. Sa richesse opérationnelle --- notamment sa dimension prescriptive sur la modalité de communication optimale --- en fait un cadre particulièrement adapté à la personnalisation des recommandations de formation. C'est ce modèle qui a été retenu dans TalentPredict AI pour piloter le module d'évaluation comportementale.
\end{itemize}

\textbf{Les centres d'évaluation (Assessment Centers):} Les assessment centers sont des dispositifs multi-méthodes combinant simulations de situations professionnelles, mises en situation collectives et tests psychométriques, évalués par plusieurs assesseurs formés. Ils présentent la validité prédictive la plus élevée parmi les outils d'évaluation des compétences managériales (r = 0,36 à 0,45 selon les méta-analyses)~\cite{ch2-thornton2006}. Toutefois, leur coût élevé (plusieurs milliers d'euros par collaborateur), leur durée (une à deux journées) et la difficulté de les déployer à l'échelle d'une organisation entière en font un outil réservé aux populations managériales, inapplicable à une évaluation systématique des compétences de l'ensemble du personnel.

\subsection{Traitement automatique du langage naturel appliqué aux profils professionnels}

Le NLP constitue la brique technologique centrale des systèmes modernes d'évaluation automatisée des compétences. Son application à l'analyse des documents professionnels (CV internes, dossiers de compétences, profils professionnels) a connu une évolution majeure avec l'avènement des architectures Transformer.

\textbf{Approches par règles et expressions régulières:} Les premières générations de parseurs de CV reposaient sur des règles heuristiques et des expressions régulières (regex) pour identifier les entités nommées (technologies, formations, expériences). Ces approches présentent l'avantage d'être déterministes, rapides et transparentes, mais souffrent de limites structurelles importantes : elles sont sensibles aux variations de formulation, ne capturent pas le contexte sémantique et nécessitent une maintenance manuelle intensive à mesure que le vocabulaire technique évolue. Les études de performance publiées font état de scores F1 typiquement compris entre 0,65 et 0,75 pour l'extraction de technologies depuis des CV~\cite{ch2-celik2013}.

\textbf{Approches par modèles Transformer (BERT, LLaMA)} L'avènement des modèles de langage pré-entraînés a profondément transformé les performances des systèmes d'extraction d'information. Chien \& Chia-Hui (2020) ont démontré que les architectures de type Transformer surpassent significativement les approches basées sur les champs aléatoires conditionnels (CRF) pour les tâches de reconnaissance d'entités nommées (NER) appliquées aux CV, avec un gain de performance de 12 à 18 points de F1~\cite{ch2-chien2020}. Des benchmarks plus récents sur des modèles de la famille LLaMA confirment des scores F1 supérieurs à 0,85 pour l'extraction de compétences techniques depuis des documents textuels non structurés, issues de la littérature~\cite{ch2-shi2023}. Dans TalentPredict AI, cette approche est mise en œuvre via le module FastAPI/Ollama qui complète l'extraction par règles d'un second niveau d'analyse sémantique LLM, permettant de capturer les technologies implicitement mentionnées ou formulées de manière non standard.

\textbf{Extraction d'indicateurs comportementaux depuis les données textuelles:} Au-delà de l'extraction de compétences techniques, le NLP permet d'inférer des signaux comportementaux depuis le texte d'un CV ou d'un profil professionnel : détection de termes évoquant le travail en équipe, le leadership ou l'adaptabilité, analyse de la progression de carrière, cohérence du parcours professionnel. Ces approches restent toutefois limitées par le caractère auto-déclaratif des sources analysées, qui reproduit le biais documenté en section 2.1.3.

La figure~\ref{fig:NLP-profile} synthétise l'application concrète de ces approches d'extraction au sein de TalentPredict AI. Elle illustre le pipeline NLP complet, détaillant le cheminement d'un document brut (CV ou profil professionnel) à travers les phases d'extraction, de nettoyage sémantique, de détection par modèles d'intelligence artificielle et de normalisation, jusqu'à la production d'un score quantifié.
\begin{figure}[H]
    \centering
    \includegraphics[width=\textwidth]{images/Pipeline NLP profile.png}
    \caption{Architecture en 6 étapes du pipeline NLP appliqué à l'extraction, la normalisation et le scoring des compétences professionnelles.}
    \label{fig:NLP-profile}
\end{figure}

\subsection{Analyse des dépôts de code source comme proxy de compétences}

L'exploitation des plateformes de versioning (GitHub, GitLab) comme source de données pour l'évaluation des compétences techniques constitue un paradigme émergent dans la recherche en ingénierie logicielle. Kalliamvakou et al. (2014) ont posé les jalons de cette approche en démontrant que les métriques d'activité GitHub (nombre de dépôts, fréquence de contribution, diversité des langages utilisés, participation aux pull requests) constituent des proxies valides  mais imparfaits  de l'expertise technique réelle d'un développeur~\cite{ch2-kalliamvakou2014}.

Les travaux de Marlow et al. (2013) ont par ailleurs établi que certaines dynamiques de collaboration observables dans les données GitHub fréquence des revues de code, vitesse de réponse aux issues, diversité des collaborateurs  présentent une corrélation modérée mais statistiquement significative avec des indicateurs de performance professionnelle évalués par les managers~\cite{ch2-marlow2013}. Ces résultats suggèrent que les données GitHub peuvent servir de proxy partiel pour certaines dimensions comportementales (ownership, collaboration, consistance) en complément des outils psychométriques classiques, ce qui est précisément le principe mis en œuvre dans le workflow d'analyse GitHub de TalentPredict.

La figure~\ref{fig:exploi-github} modélise ce workflow d'analyse spécifique à notre plateforme. Elle illustre concrètement comment les différentes métriques extraites via l'API GitHub (diversité des langages, volume de commits, régularité via la heat-map, et dynamiques de collaboration) sont transformées en indicateurs pondérés pour générer un score technique global et objectif.

\begin{figure}[H]
    \centering
    \includegraphics[width=\textwidth]{images/exploi github data en preuve tech.png}
    \caption{Exploitation des données
GitHub comme preuve
techniqu}
    \label{fig:exploi-github}
\end{figure}

Néanmoins, plusieurs limites à cette approche ont été identifiées dans la littérature : le biais en faveur des développeurs contribuant à des projets open source publics (défavorable aux collaborateurs travaillant exclusivement sur des dépôts privés en entreprise), la difficulté à distinguer quantité et qualité des contributions, et la variabilité des pratiques de versioning entre organisations~\cite{ch2-kalliamvakou2014}.

\subsection{Grands modèles de langage : architectures et performances comparées}
Le traitement sémantique avancé requis par TalentPredict s'appuie sur l'utilisation de Grands Modèles de Langage (LLM). Cette section dresse un état des lieux des architectures existantes, propose une analyse comparative de plusieurs modèles candidats, et justifie les choix technologiques retenus pour notre moteur d'inférence.
\textbf{Panorama des architectures LLM} Depuis la publication de l'architecture Transformer par Vaswani et al. (2017)~\cite{ch2-vaswani2017}, le champ des modèles de langage de grande taille a connu une évolution exponentielle. Deux grandes familles se distinguent pour les applications en entreprise :
\begin{itemize}[leftmargin=1.5cm]
  \item Les modèles propriétaires hébergés dans le cloud (GPT-4 d'OpenAI, Claude d'Anthropic, Gemini de Google) offrent des performances de pointe sur les tâches de compréhension sémantique complexe et de génération structurée, mais impliquent une dépendance à un fournisseur tiers, des coûts d'utilisation à l'appel potentiellement élevés, et surtout une externalisation du traitement des données vers des infrastructures non maîtrisées par l'organisation.
  \item Les modèles open-source déployables localement (famille LLaMA de Meta, Mistral AI, Phi-3 de Microsoft) permettent un déploiement on-premise ou dans une infrastructure privée, garantissant ainsi la souveraineté des données. Leurs performances, initialement inférieures aux modèles propriétaires, ont considérablement progressé avec les récentes générations.
\end{itemize}

\textbf{Comparaison quantitative et qualitative des modèles open-source candidats} Plusieurs modèles ont été évalués comme candidats potentiels pour le moteur d'inférence de TalentPredict AI. Le tableau suivant synthétise leurs caractéristiques selon les critères pertinents pour ce projet :

\begin{table}[H]
\centering
\small
\renewcommand{\arraystretch}{1.35}
\begin{tabularx}{\textwidth}{|>{\RaggedRight\arraybackslash}p{0.22\textwidth}|Y|Y|Y|Y|}
\hline
\rowcolor{TableHeaderBlue}
\color{white}\textbf{Critère} & \color{white}\textbf{LLaMA 3.2 (3B/8B)} & \color{white}\textbf{Mistral 7B} & \color{white}\textbf{Phi-3 Mini (3.8B)} & \color{white}\textbf{GPT-4 (cloud)} \\
\hline
Taille (paramètres) & 3B / 8B & 7B & 3.8B & \textasciitilde{}1T (estimé) \\
\hline
Déploiement local & Oui, via Ollama & Oui, via Ollama & Oui, via Ollama & Non, API cloud \\
\hline
Performance NLP générale (MMLU) & 63\% / 73\% & 64\% & 69\% & 86\% \\
\hline
Qualité génération JSON structuré & Bonne & Bonne & Moyenne & Excellente \\
\hline
Latence CPU (génération 200 tokens) & \textasciitilde{}8s / \textasciitilde{}18s & \textasciitilde{}15s & \textasciitilde{}6s & \textasciitilde{}3s (réseau) \\
\hline
Coût par requête & 0 (local) & 0 (local) & 0 (local) & \textasciitilde{}0,03€ \\
\hline
Conformité RGPD & Totale & Totale & Totale & Limitée : données cloud \\
\hline
Support CodeLlama (code) & Dédié & Partiel & Limité & Non disponible \\
\hline
Licence & Meta LLaMA & Apache 2.0 & MIT & Propriétaire \\
\hline
\end{tabularx}
\renewcommand{\arraystretch}{1.05}
\caption{Comparaison des modèles LLM candidats pour TalentPredict AI}
\end{table}

Sources : benchmarks MMLU publiés par les éditeurs ; mesures de latence estimées d'après la littérature sur du matériel standard (CPU 8 cœurs, 16 Go RAM)~\cite{ch2-meta2024,ch2-mistral2023,ch2-abdin2024}.

\textbf{Justification du choix de llama3.2 et CodeLlama} Au terme de cette analyse comparative, le modèle llama3.2 (variante 3B pour les analyses légères, 8B pour les synthèses comportementales complexes) a été retenu comme moteur d'inférence principal de TalentPredict AI, pour les raisons suivantes. Touvron et al. (2023) ont démontré que les architectures LLaMA offrent un rapport performance/taille remarquable pour les tâches de NLP standard, avec des performances compétitives par rapport aux modèles hébergés dans le cloud sur les tâches de classification et d'extraction sémantique~\cite{ch2-touvron2023}. Des évaluations ultérieures (Xu et al., 2024) ont confirmé la pertinence de llama3.2 quantifié en Q4 pour les tâches NLP appliquées aux ressources humaines, avec des latences acceptables sur des processeurs standards~\cite{ch2-xu2024}.

Le modèle CodeLlama, dérivé de LLaMA et spécialisé pour la compréhension et la génération de code, a été retenu pour le module de challenge de code en raison de ses performances supérieures aux modèles généralistes sur les tâches d'évaluation algorithmique. Mistral 7B a été écarté en raison d'une latence CPU plus élevée sans gain de performance significatif sur les tâches ciblées. Phi-3 Mini, bien que plus rapide, présentait une qualité de génération JSON structuré insuffisante pour les workflows d'orchestration.

\subsection{Edge AI et souveraineté des données}

Le concept d'Edge AI désigne le déploiement de modèles d'intelligence artificielle directement sur l'infrastructure locale de l'organisation, sans recours à des services cloud externes. Dans le contexte de l'évaluation des compétences RH, cette approche revêt une importance particulière au regard du cadre réglementaire européen.

L'Article 9 du RGPD classe les données psychologiques dans la catégorie des «~données à caractère sensible~», dont le traitement est soumis à des exigences renforcées~\cite{ch2-rgpd2016}. Les analyses de personnalité issues du modèle PCM, les scores comportementaux inférés depuis les profils professionnels et les résultats de tests psychométriques entrent directement dans cette catégorie. Leur transmission vers des infrastructures cloud tierces (serveurs hors UE, notamment pour les modèles américains) soulève des risques de non-conformité réglementaire que les DPO (Data Protection Officers) des grandes organisations ne peuvent pas accepter.

L'utilisation d'Ollama comme serveur d'inférence local dans TalentPredict AI résout structurellement ce problème : l'intégralité du traitement sémantique s'effectue sur l'infrastructure interne, les données ne quittent jamais le réseau privé Docker de la plateforme, et la conformité au principe de Privacy by Design (Article 25 du RGPD) est garantie par architecture.

La figure~\ref{fig:inference-local} illustre schématiquement cette rupture architecturale. Elle met en opposition l'approche cloud traditionnelle écartée en raison de ses risques inhérents de non conformité  avec l'approche d'inférence locale (via Ollama) retenue pour TalentPredict AI, garantissant ainsi une souveraineté et une confidentialité totales des données par conception.
\begin{figure}[H]
    \centering
    \includegraphics[width=\textwidth]{images/inference local ollama.png}
    \caption{Inférence locale avec Ol-
lama}
    \label{fig:inference-local}
\end{figure}

\subsection{Orchestration de workflows et architectures microservices}
L'intégration efficace de modèles d'intelligence artificielle au sein d'une application d'entreprise complexe exige une architecture logicielle robuste et modulaire. Cette section explore les principes d'orchestration et de conception en microservices qui sous-tendent l'infrastructure technique de TalentPredict, garantissant ainsi la flexibilité et la scalabilité de ses traitements.
\textbf{L'orchestration low-code:} Les plateformes d'orchestration visuelles de workflows (n8n, Zapier, Make, Apache Airflow) permettent de modéliser des séquences complexes d'appels à des services hétérogènes sans code impératif. van der Aalst (2013) a théorisé les bénéfices du découplage de la logique d'orchestration hors du noyau applicatif : meilleure maintenabilité, traçabilité des flux, et capacité à modifier les comportements (notamment les prompts LLM) sans recompilation du code source~\cite{ch2-vanderaalst2013}. Dans TalentPredict AI, n8n a été retenu pour orchestrer les workflows d'analyse comportementale (GitHub Analysis, CV Parser, PCM Test, Master Agent), permettant aux équipes RH ou IA de modifier les prompts directement depuis l'interface graphique sans intervention du développeur.

\textbf{Les architectures microservices pour l'IA:} Le pattern de découplage entre la logique métier (Spring Boot) et les traitements d'inférence IA (FastAPI/Python) s'inscrit dans les recommandations architecturales des systèmes MLOps modernes. Ce découplage permet une évolution indépendante des deux couches --- notamment le remplacement ou la mise à jour du modèle LLM sans impact sur le backend Java et facilite la gestion des ressources computationnelles spécifiques à l'inférence (GPU, RAM étendue). Les stratégies de timeout et de fallback heuristique, mises en œuvre dans TalentPredict pour gérer la latence variable des appels Ollama, constituent une bonne pratique documentée dans la littérature sur la résilience des systèmes distribués~\cite{ch2-newman2015}.

\textbf{Les systèmes de recommandation pédagogique:} Drachsler et al. (2009) ont établi, dans le contexte des environnements d'apprentissage enrichis par la technologie (Technology Enhanced Learning), que la contextualisation des recommandations de formation en fonction du profil précis de l'apprenant  plutôt qu'une approche générique par filtrage collaboratif  augmente le taux de réussite de 23 à 35\%~\cite{ch2-drachsler2009}. Ce résultat fonde scientifiquement le choix de TalentPredict de générer des recommandations prescriptives croisant le profil PCM du collaborateur avec ses lacunes techniques identifiées, plutôt que de se limiter à des suggestions génériques de cours.

\section{Travaux existants et solutions similaires}
Afin de positionner l'apport de notre plateforme par rapport à l'existant, cette section propose une analyse critique des solutions actuellement disponibles sur le marché. Celles-ci ont été classées en trois grandes catégories : les outils d'évaluation psychométrique, les plateformes de tests techniques dédiées aux développeurs, et enfin les systèmes d'apprentissage pilotés par l'IA (LXP).
\subsection{Solutions d'évaluation psychométrique}
La première catégorie d'outils étudiée concerne l'évaluation des compétences comportementales (soft skills). Cette sous-section analyse les plateformes psychométriques les plus représentatives du marché, qu'elles s'appuient sur des questionnaires auto-déclaratifs classiques ou sur l'inférence de traits de personnalité par l'analyse algorithmique
\textbf{PerformanSe} est une plateforme psychométrique française proposant des questionnaires basés sur le Big Five et d'autres modèles validés scientifiquement, destinés à l'évaluation des compétences comportementales des collaborateurs. Son principe repose sur des questionnaires auto-déclaratifs en ligne, avec restitution d'un rapport de personnalité et de recommandations de développement.

\textbf{Apports :} validation psychométrique rigoureuse des outils, restitution structurée du profil comportemental, utilisation en contexte RH français bien établie.

\textbf{Limites :} le biais déclaratif n'est pas atténué par des mécanismes de validation externe ; les recommandations de formation restent génériques et non corrélées à une analyse des compétences techniques du collaborateur ; aucune automatisation du cycle d'évaluation n'est proposée.

\textbf{Crystal Knows} prédit le type DISC d'un individu à partir de l'analyse NLP de son profil LinkedIn public. Son approche est intéressante car elle tente d'inférer des traits comportementaux depuis des données numériques observables plutôt que déclaratives.

\textbf{Apports :} réduction partielle du biais déclaratif en s'appuyant sur des données textuelles publiques.

\textbf{Limites :} la validité prédictive du modèle DISC est contestée dans la littérature scientifique~\cite{ch2-pittenger1993} ; l'accuracy du classificateur n'est pas publiée par l'éditeur ; la solution est totalement déconnectée de toute évaluation technique ; la dépendance aux données LinkedIn introduit un biais lié à la complétude et au style rédactionnel des profils.

\subsection{Systèmes d'évaluation technique par tests}
Les systèmes d’évaluation technique traditionnels se concentrent sur la mesure de l'expertise opérationnelle par le biais de challenges et de QCM. Nous examinons ici deux solutions majeures de ce domaine afin d'en dégager les apports et les limites.
\textbf{HackerRank} est l'une des plateformes de référence pour l'évaluation des compétences en programmation, proposant des challenges algorithmiques dans une interface de développement intégrée, avec correction automatique et scoring.

\textbf{Apports :} objectivité des scores (la solution produit ou non le résultat attendu), large bibliothèque de problèmes couvrant de nombreux langages, fiabilité du scoring technique.

\textbf{Limites :} les tests sont statiques et ne s'adaptent pas dynamiquement au profil du collaborateur ; aucune donnée comportementale n'est collectée ni analysée ; l'absence de croisement avec le profil GitHub ou le CV ne permet pas de détecter les incohérences entre compétences déclarées et compétences démontrées.

\textbf{TestGorilla} propose une approche combinant tests techniques standardisés et questionnaires comportementaux, mais sans lien organique entre les deux dimensions d'évaluation.

\textbf{Apports :} facilité de déploiement, variété des tests disponibles.

\textbf{Limites :} absence d'adaptation dynamique de la difficulté ; les tests comportementaux restent auto-déclaratifs sans validation par des données comportementales asynchrones ; aucun mécanisme de suivi de la progression dans le temps n'est intégré.

\subsection{Plateformes d'expérience d'apprentissage pilotées par IA}
Le troisième volet de notre étude de l'existant se concentre sur l'étape de prescription de formations. Nous analysons ici les solutions d'apprentissage modernes, dites LXP, afin de comprendre comment elles exploitent les algorithmes de recommandation, mais aussi pourquoi leurs suggestions manquent souvent de précision objective.
\textbf{Degreed} est une LXP (Learning Experience Platform) qui agrège des contenus de formation issus de multiples sources (MOOCs, articles, vidéos) et utilise des algorithmes de recommandation pour suggérer des parcours adaptés au profil déclaré du collaborateur.

\textbf{Apports :} large catalogue de ressources, interface utilisateur soignée, intégration avec les outils RH existants.

\textbf{Limites :} les recommandations sont basées sur des compétences déclarées et non validées par des preuves objectives ; l'utilisation de LLM cloud pour la personnalisation expose les données RH à des risques de conformité RGPD ; aucune évaluation des compétences techniques n'est intégrée en amont de la recommandation.

\textbf{360Learning} se distingue par son approche d'apprentissage collaboratif où les experts internes créent les contenus de formation.

\textbf{Apports :} valorisation de l'expertise interne, forte dimension sociale de l'apprentissage.

\textbf{Limites :} la recommandation de formations est faiblement automatisée et repose largement sur des décisions manuelles des équipes L\&D ; l'absence d'analyse automatique des écarts de compétences (skill gap analysis) conduit à des propositions de formation non nécessairement alignées sur les besoins réels du collaborateur.

\section{Analyse comparative}
Afin de situer clairement TalentPredict AI dans son écosystème fonctionnel et technologique, cette section propose de confronter notre solution aux principales plateformes d'évaluation et de formation actuellement disponibles sur le marché RH.

\subsection{Tableau comparatif multi-critères}

Le tableau~\ref{tab:comparatif-multi-carecteres} synthétise ce positionnement face à six solutions concurrentes reconnues, en les évaluant selon des critères déterminants pour notre problématique : la capacité d'évaluation multidimensionnelle, les mécanismes de validation croisée, la souveraineté des données (conformité RGPD) et le modèle économique.
\begin{table}[p]
\centering
\scriptsize
\renewcommand{\arraystretch}{1.28}
\setlength{\tabcolsep}{4pt}
\begin{tabularx}{\linewidth}{|>{\RaggedRight\arraybackslash}p{0.13\linewidth}|*{7}{>{\RaggedRight\arraybackslash}X|}}
\hline
\rowcolor{TableHeaderBlue}
\color{white}\textbf{Critère} & \color{white}\textbf{Performan\-Se} & \color{white}\textbf{Crystal Knows} & \color{white}\textbf{Hacker\-Rank} & \color{white}\textbf{Test\-Gorilla} & \color{white}\textbf{Degreed} & \color{white}\textbf{360\-Learning} & \color{white}\textbf{Talent\-Predict AI} \\
\hline
Évaluation technique & Non & Non & Oui, tests statiques & Oui, partielle & Non & Non & Oui, multi-sources \\
\hline
Évaluation comportementale & Oui, Big Five & Oui, DISC & Non & Oui, déclarative & Non & Non & Oui, PCM + IA locale \\
\hline
Validation croisée & Non & Non & Non & Non & Non & Non & Oui, GitHub -- CV -- Tests \\
\hline
Analyse GitHub & Non & Non & Non & Non & Non & Non & Oui \\
\hline
Recom\-mandation prescriptive & Partielle & Non & Non & Non & Oui, générique & Oui, générique & Oui, adaptée PCM \\
\hline
Souveraineté des données & Non, cloud & Non, cloud & Non, cloud & Non, cloud & Non, LLM cloud & Non, cloud & Oui, Edge AI / on-premise \\
\hline
Scoring multi-sources & Non & Non & Non & Non & Non & Non & Oui, 500+ alias \\
\hline
Suivi de progression & Non & Non & Non & Non & Oui & Oui & Oui, Kanban + gamification \\
\hline
Conformité RGPD données sensibles & Partielle & Non & Non & Non & Non & Non & Oui, totale \\
\hline
Coût d'exploitation & Licence élevée & Abonne\-ment & Licence & Licence & Licence élevée & Licence élevée & Open-source \\
\hline
\end{tabularx}
\renewcommand{\arraystretch}{1.05}
\caption{Analyse comparative multi-critères des solutions existantes et de TalentPredict AI}
\label{tab:comparatif-multi-carecteres}
\end{table}

Afin de synthétiser cette étude comparative, la figure~\ref{fig:positionnement-existant} propose une lecture graphique du marché sous la forme d'une matrice de positionnement. Cette cartographie permet de visualiser l'écart entre les solutions actuelles et TalentPredict AI, en croisant le degré d'objectivité de l'évaluation avec le niveau de souveraineté des données garanti par l'architecture.
\begin{figure}[H]
    \centering
    \includegraphics[width=\textwidth]{images/position existant.png}
    \caption{Matrice de positionnement concurrentiel de TalentPredict AI selon les axes de souveraineté et d'objectivité.}
    \label{fig:positionnement-existant}
\end{figure}
\subsection{Synthèse analytique des différenciateurs}

L'analyse comparative révèle trois fractures structurelles dans l'offre existante.

La première fracture est l'opposition irrésolue entre objectivité et complétude. Les solutions techniques (HackerRank, TestGorilla) offrent une bonne objectivité mais ignorent la dimension comportementale. Les solutions psychométriques (PerformanSe, Crystal Knows) couvrent la dimension comportementale mais de manière déclarative. Aucune solution n'intègre les deux dimensions dans un pipeline unifié avec validation croisée.

La deuxième fracture est le manque de validation factuelle des compétences déclarées. Toutes les solutions analysées acceptent sans vérification les compétences que le collaborateur ou son manager affirme maîtriser. TalentPredict introduit une rupture en confrontant systématiquement les déclarations à des preuves asynchrones (historique GitHub, résultats de tests génératifs).

La troisième fracture est l'incompatibilité entre puissance analytique LLM et souveraineté des données. Les solutions LXP modernes commencent à intégrer des LLM pour la personnalisation, mais via des APIs cloud qui exposent des données RH sensibles à des tiers. TalentPredict résout cette tension par l'inférence locale via Ollama.

\section{Positionnement du projet TalentPredict AI}
À la lumière de l'état de l'art et de l'analyse comparative présentés précédemment, cette section a pour objectif de définir le positionnement exact de TalentPredict AI. Pour ce faire, nous commencerons par formaliser les limites majeures des approches actuelles, justifiant ainsi les choix architecturaux et fonctionnels de notre solution.
\subsection{Lacunes identifiées dans la littérature et sur le marché}

L'analyse des travaux existants et des solutions disponibles permet d'identifier trois lacunes principales qui restent insuffisamment traitées.

\textbf{Lacune 1 : Absence de pipeline open-source unifié combinant évaluation technique et comportementale:} Aucune solution open-source ne propose, à ce jour, un système intégré permettant d'évaluer simultanément les compétences techniques d'un collaborateur (via analyse de code, parsing de CV, tests adaptatifs) et ses compétences comportementales (via modèle psychométrique et analyse sémantique), en croisant les résultats des deux dimensions pour produire un profil unifié et des recommandations prescriptives.

\textbf{Lacune 2 : Absence de validation croisée multi-sources avec déduplication sémantique:} Les systèmes d'évaluation existants traitent chaque source de données (CV, tests, LinkedIn) de manière indépendante. Aucun mécanisme de déduplication et de normalisation sémantique entre sources hétérogènes n'est proposé, conduisant à des évaluations fragmentées et potentiellement contradictoires. TalentPredict répond à cette lacune avec son algorithme de normalisation par dictionnaire d'alias (500+ variantes canoniques) et son moteur de scoring par triangulation multi-sources.

\textbf{Lacune 3 : Dépendance au cloud pour les solutions LLM appliquées aux RH:} Les solutions utilisant des LLM pour l'analyse des compétences comportementales externalisent toutes leur traitement vers des infrastructures cloud, compromettant la conformité RGPD pour les données psychologiques (Article 9). TalentPredict répond à cette lacune par un déploiement 100\% local via Ollama, garantissant la souveraineté totale des données.

\subsection{Contributions spécifiques de TalentPredict AI}

TalentPredict AI apporte trois contributions originales répondant directement aux lacunes identifiées.

La première contribution est une architecture hybride n8n/FastAPI pour l'Edge AI RH. L'orchestration des workflows comportementaux via n8n permet de séparer la logique de prompt engineering de la logique applicative, offrant une maintenabilité supérieure aux approches monolithiques tout en garantissant l'inférence locale. Cette architecture constitue, à notre connaissance, une approche originale dans le domaine de la HR Tech open-source.

La deuxième contribution est un algorithme de scoring multi-sources avec normalisation sémantique. Le pipeline d'évaluation technique de TalentPredict implémente un mécanisme de déduplication basé sur un dictionnaire de 500+ alias canoniques (par exemple : reactjs $\rightarrow$ React, k8s $\rightarrow$ Kubernetes, sklearn $\rightarrow$ Scikit-learn), couplé à un algorithme de nivellement par triangulation (Expert, Advanced, Intermediate, Beginner) selon le nombre de sources confirmant la compétence et la force du signal dans chacune. Ce mécanisme atténue structurellement le biais déclaratif en pondérant les sources objectives (GitHub, tests) davantage que les sources déclaratives (CV).

La troisième contribution est un pipeline PCM heuristique à reproductibilité totale couplé à une validation LLM sémantique. Le calcul des scores PCM est implémenté de manière entièrement algorithmique en JavaScript, garantissant une reproductibilité de 100\% (même entrée $\rightarrow$ même sortie, quelle que soit l'exécution), ce qu'un calcul LLM ne peut pas assurer en raison de la stochasticité inhérente à la génération de texte (accord inter-exécutions typiquement de 70 à 80\% pour les tâches de classification psychométrique~\cite{ch2-xu2024}). La dimension LLM intervient uniquement pour l'analyse sémantique des données textuelles et la synthèse narrative, où la variabilité est acceptable et même souhaitable.

Le positionnement scientifique et technique de TalentPredict AI repose sur trois piliers fondamentaux, ou contributions, qui structurent l'innovation du projet. Le schéma~\ref{fig:position-platfome} détaille ces contributions (C1, C2 et C3) en mettant en relation les objectifs méthodologiques avec la stack technologique et les mécanismes d'orchestration mis en œuvre.
\begin{figure}[H]
    \centering
    \includegraphics[width=\textwidth]{images/position platforme.png}
    \caption{Positionnement scienti-
fique et technique de Ta-
lentPredict AI}
    \label{fig:position-platfome}
\end{figure}
\subsection{Valeur ajoutée synthétique}
Afin de résumer la proposition de valeur de notre plateforme, le tableau~\ref{tab:valeur-ajoutee-synthetique} illustre comment TalentPredict AI apporte une réponse technologique et méthodologique directe à chacune des lacunes identifiées dans l'état de l'art.
\begingroup
\small
\renewcommand{\arraystretch}{1.35}
\setlength{\tabcolsep}{5pt}
\begin{longtable}{|>{\RaggedRight\arraybackslash}p{0.30\textwidth}|>{\RaggedRight\arraybackslash}p{0.62\textwidth}|}
\label{tab:valeur-ajoutee-synthetique}\\
\hline
\rowcolor{TableHeaderBlue}
\color{white}\textbf{Lacune identifiée} & \color{white}\textbf{Réponse de TalentPredict AI} \\
\hline
\endfirsthead
\hline
\rowcolor{TableHeaderBlue}
\color{white}\textbf{Lacune identifiée} & \color{white}\textbf{Réponse de TalentPredict AI} \\
\hline
\endhead
Fragmentation technique/comportementale & Pipeline unifié combinant l'évaluation technique (GitHub, CV, QCM et code) avec l'évaluation comportementale (PCM, Ollama et scénarios IA). \\
\hline
Biais déclaratif non contrôlé & Validation croisée multi-sources : les compétences déclarées sont confrontées à des preuves asynchrones et objectives. \\
\hline
Dépendance cloud / RGPD & Inférence locale via Ollama, sans externalisation des données RH sensibles. \\
\hline
Recommandations génériques & Recommandations prescriptives avec une modalité pédagogique adaptée au profil PCM du collaborateur. \\
\hline
Coûts de licence élevés & Stack open-source fondée sur n8n, Ollama, FastAPI, Spring Boot, PostgreSQL et Angular. \\
\hline
\caption{Synthèse des contributions de TalentPredict AI face aux lacunes identifiées}
\end{longtable}
\renewcommand{\arraystretch}{1.05}
\endgroup

\section{Conclusion}
\phantomsection
\addcontentsline{toc}{section}{Conclusion}

Ce chapitre établit le socle théorique de TalentPredict AI en démontrant que face aux limites des méthodes d'évaluation traditionnelles (biais déclaratif, coûts prohibitifs, fragmentation) et à l'incapacité des solutions actuelles du marché à garantir simultanément l'objectivité factuelle, l'analyse technique et comportementale, ainsi que la souveraineté des données, la conception d'une plateforme innovante s'impose. Pour combler ces lacunes, notre étude a directement dicté les choix architecturaux qui seront détaillés lors des quatre prochains sprints : l'exploitation des données GitHub comme indicateur comportemental, combinée à la puissance des LLMs via le déploiement local de llama3.2 sous Ollama pour assurer la conformité RGPD. Ces technologies s'articulent autour du modèle scientifique PCM, d'une orchestration des workflows par n8n, et d'un scoring multi-sources indispensable pour neutraliser efficacement les biais d'auto-évaluation.




\chapter{ Sprint 1 : Mise en place de l’Architecture, de la Sécurité et du Système d’Authentification}
\section*{Introduction}
\phantomsection
\addcontentsline{toc}{section}{Introduction}
\phantomsection 
\noindent
Le premier sprint de TalentPredict a constitué une phase essentielle de mise en place technique et sécuritaire. Avant l’intégration des modules d’intelligence artificielle, l’objectif était de construire une architecture stable, sécurisée et performante. Cette étape a reposé sur une infrastructure conteneurisée facilitant le déploiement, ainsi que sur l’implémentation des mécanismes de sécurité JWT et BCrypt, afin de préparer une intégration fluide des futurs modules prédictifs.
 \section{Objectifs du Sprint 1}
L’objectif principal de ce sprint est de livrer un incrément logiciel fonctionnel constituant la base de TalentPredict. Cette phase vise à mettre en place une infrastructure stable et structurée, capable d’intégrer les futurs modules métiers du projet. Elle garantit également une expérience utilisateur fluide dans un environnement sécurisé et optimisé pour les besoins de performance de l’application.

\subsubsection{1. Objectifs Fonctionnels et Interactifs}
\begin{itemize}
    \item \textbf{Système d'Authentification Complet (IAM) :}
    \begin{itemize}
        \item Mise en œuvre du cycle d’inscription (Register), consolidé par un système de validation d’adresse email afin de garantir l’authenticité des comptes créés.
        \item Déploiement d’une authentification sécurisée (Login), enrichie d’une gestion rigoureuse des erreurs de saisie pour optimiser l’expérience utilisateur et la sécurité des accès.
        \item Implémentation de la procédure de récupération de mot de passe (Forgot Password), s'appuyant sur l'émission de tokens temporaires pour assurer la confidentialité et l'intégrité des données.
    \end{itemize}
    
    \item \textbf{Gestion Dynamique des Profils :}
    \begin{itemize}
        \item Interface de consultation du profil utilisateur permettant de visualiser de manière structurée les données d’identité ainsi que les rôles associés.
        \item Formulaires d'édition optimisés pour la mise à jour des informations professionnelles, telles que l'intitulé du poste, la biographie et les réseaux sociaux.
    \end{itemize}
    
    \item \textbf{Interface Utilisateur (UX/UI) de Base :}
    \begin{itemize}
        \item Déploiement d’une page d'accueil responsive, intégrant les sections stratégiques : Hero, présentation des fonctionnalités et appels à l’action (CTA).
        \item Mise en place du Dashboard Shell, incluant une navigation latérale (Sidenav) ergonomique et une barre supérieure (Topbar) dotée du menu utilisateur.
        
    \end{itemize}
    
    \item \textbf{Autorisation Granulaire :}
    \begin{itemize}
        \item Développement de gardes de navigation (Route Guards) afin de sécuriser les accès et d'interdire systématiquement l'affichage des pages protégées aux utilisateurs non authentifiés.
        \item Adaptation dynamique de l’interface en fonction des privilèges de l'utilisateur, permettant notamment de restreindre la visibilité de certaines sections, telle que la gestion des administrateurs, aux seuls rôles autorisés.
    \end{itemize}
\end{itemize}

\subsubsection{2. Objectifs Techniques et Opérationnels}

\begin{itemize}
    \item \textbf{Standardisation de l'Architecture :}
    \begin{itemize}
        \item Initialisation du projet en \textit{Monolithe Modulaire} sous \texttt{Spring Boot 3.4.2}.
        \item Organisation rigoureuse des \textit{packages} (\texttt{Controller}, \texttt{Service}, \texttt{Repository}, \texttt{DTO}).
    \end{itemize}
    
    \item \textbf{Sécurité et Cryptographie :}
    \begin{itemize}
        \item Mise en place de \texttt{Spring Security} avec une politique de session \textit{stateless}.
        \item Implémentation du standard \texttt{JWT} (Génération, Signature HS256, Extraction et Validation).
        \item Hachage des secrets via l'algorithme \texttt{BCrypt}.
    \end{itemize}
    
    \item \textbf{Infrastructure et DevOps :}
    \begin{itemize}
        \item Orchestration complète via \texttt{Docker Compose} pour assurer la parité des environnements.
        \item Configuration sécurisée des variables d'environnement (\texttt{.env}).
    \end{itemize}
    
    \item \textbf{Robustesse et Qualité de Code :}
    \begin{itemize}
        \item Centralisation de la gestion des exceptions via un \texttt{GlobalExceptionHandler}.
        \item Intercepteur HTTP côté \texttt{Angular} pour l'injection automatique du \textit{token} dans les \textit{headers}.
        \item Persistance via \texttt{PostgreSQL}; initialisation automatique du schéma avec \texttt{JPA} et \texttt{Hibernate}.
    \end{itemize}
    
    \item \textbf{Exposition et Test de l'API :}
    \begin{itemize}
        \item Configuration de \texttt{Swagger} pour la documentation et les tests en temps réel.
        \item Mise en œuvre des politiques de \texttt{CORS} restrictives.
    \end{itemize}
\end{itemize}

\paragraph{Critères d’acceptation (Definition of Done)}

Pour valider chaque fonctionnalité de ce sprint, les critères suivants ont été appliqués :

\begin{itemize}
    \item \textbf{Fonctionnalité} : Conformité totale avec la \textit{User Story} et validation des tests unitaires.
    \item \textbf{Sécurité} : Endpoints protégés par JWT et mots de passe hachés via \textit{BCrypt}.
    \item \textbf{Qualité UI} : Interface responsive (mobile/desktop) et absence d'erreurs console.
    \item \textbf{Documentation} : Exposition des APIs via \textit{Swagger} et code documenté.
    \item \textbf{Déploiement} : Construction réussie de l'image \textit{Docker} et intégration au \textit{Compose}.
\end{itemize}



\section{Séance de Sprint Planning}
La séance de Sprint Planning a réuni l’équipe de développement, composée de Mohamed Ben Abdeladhim et Rahma Barouni, ainsi que le Product Owner, M. Med Amine Gharbi, et l’encadrante pédagogique, Mme Ben Aïcha Takwa. Cette réunion a permis de définir le périmètre fonctionnel et technique du premier sprint de TalentPredict.

Les points clés suivants ont été déterminés lors de cette séance :
\begin{itemize}
    \item \textbf{Durée :} La durée de cette itération a été fixée à quatre semaines. Ce délai a été jugé nécessaire pour mettre en place une infrastructure conteneurisée robuste et implémenter un système de sécurité complexe (\texttt{JWT}) garantissant l'intégrité des futures données d'évaluation.
    
    \item \textbf{Priorités Absolues :}
    \begin{itemize}
        \item Mise en place de l'architecture \textit{Monolithe Modulaire} sous \texttt{Spring Boot} et du schéma relationnel \texttt{PostgreSQL}.
        \item Développement du module \textit{Authentication \& Security} (Gestion des \textit{tokens} \texttt{JWT}, hachage \texttt{BCrypt} et contrôle d'accès \texttt{RBAC}).
        \item Implémentation du système d'Inscription et de Connexion pour les deux rôles principaux : Employé et Administrateur RH.
        \item Création de la page d’accueil et du \textit{Dashboard Shell} (\textit{Sidenav}, \textit{Topbar}) pour stabiliser l'interface utilisateur.
    \end{itemize}
    
    \item \textbf{Priorités Secondaires :}
    \begin{itemize}
        \item Configuration de la documentation interactive de l'API via \texttt{Swagger}.
        \item Centralisation de la gestion des erreurs (\textit{Global Exception Handling}) et configuration des politiques \texttt{CORS}.
        \item Mise en place du service de récupération de mot de passe et de gestion des profils utilisateurs.
    \end{itemize}
    
    \item \textbf{Prérequis :} La disponibilité de l'environnement de virtualisation \texttt{Docker} et l'initialisation du \textit{repository} \texttt{Git} sont essentielles pour permettre une collaboration fluide et un déploiement continu dès les premières phases du projet.
\end{itemize}

À la fin de cette réunion, les besoins ont été décomposés en User Stories (US) et estimés en jours-homme. La priorité a été accordée aux composants de sécurité et à l’infrastructure, essentiels à l’intégration des futurs modules d’intelligence artificielle.
\section{Backlog du Sprint }

Le \textit{Sprint Backlog} constitue le plan opérationnel détaillé de cette itération. Il résulte de la décomposition des \textit{User Stories} (besoins utilisateurs) en tâches techniques élémentaires, permettant ainsi une exécution structurée et un suivi précis de l'avancement.

Chaque tâche a fait l’objet d’une estimation en jours-homme (JH) afin d’évaluer la charge de travail et d’assurer une répartition équilibrée entre les membres du binôme de développement.
 Le tableau \ref{tab:backlog_sprint1} ci-dessous présente la feuille de route exhaustive de ce premier sprint :
\begin{longtable}{|c|>{\raggedright\arraybackslash}p{3cm}|>{\raggedright\arraybackslash}p{7cm}|c|>{\centering\arraybackslash}p{1.8cm}|}
    % --- First Page Header ---
    \hline
    \rowcolor{TableHeaderBlue}
    \color{white}\textbf{ID} & \color{white}\textbf{Module / US} & \color{white}\textbf{Tâches de Développement} & \color{white}\textbf{JH} & \color{white}\textbf{Priorité} \\
    \hline
    \endfirsthead

    % --- Subsequent Pages Header ---
    \hline
    \rowcolor{TableHeaderBlue}
    \color{white}\textbf{ID} & \color{white}\textbf{Module / US} & \color{white}\textbf{Tâches de Développement} & \color{white}\textbf{JH} & \color{white}\textbf{Priorité} \\
    \hline
    \endhead

    % --- Footer for pages before the last ---
    \hline
    \multicolumn{5}{|r|}{\textit{Suite à la page suivante...}} \\
    \hline
    \endfoot

    % --- Final Footer ---
    \hline
    \endlastfoot

    % --- Table Data ---
    1.1 & Auth / Enregistrement & 
    \begin{itemize}[leftmargin=*, nosep]
        \item Formulaire de \textit{Registration} (\texttt{Angular}).
        \item Service \textit{Backend} et hachage \texttt{BCrypt}.
        \item Validation et gestion des doublons.
    \end{itemize} & 4 & Haute \\
    \hline
    
    1.2 & Auth / Connexion JWT & 
    \begin{itemize}[leftmargin=*, nosep]
        \item Configuration \texttt{Spring Security} (\textit{Stateless}).
        \item Développement du \texttt{TokenProvider} (JWT).
        \item Composant \textit{Login} et stockage du \textit{token}.
    \end{itemize} & 5 & Haute \\
    \hline
    
    1.3 & Auth / Forget Password & Interface "Mot de passe oublié" et logique de génération de \textit{token} de récupération. & 3 & Haute \\
    \hline
    
    1.4 & Services / Notifications & 
    \begin{itemize}[leftmargin=*, nosep]
        \item Intégration \texttt{Spring Mail} (SMTP).
        \item \textit{Templates} d'emails et envoi asynchrone.
    \end{itemize} & 3 & Moyenne \\
    \hline
    
    1.5 & Profil / Différencié & 
    \begin{itemize}[leftmargin=*, nosep]
        \item \textbf{User :} Champs techniques (GitHub, CV).
        \item \textbf{Admin :} Champs RH.
        \item \textit{Endpoints} CRUD et affichage conditionnel.
    \end{itemize} & 5 & Haute \\
    \hline
    
    1.6 & Profil / Sécurité & Changement de MDP, option de désactivation et audit des connexions. & 3 & Moyenne \\
    \hline
    
    1.7 & Interface / Navigation & 
    \begin{itemize}[leftmargin=*, nosep]
        \item \textit{Design} Landing Page et \textit{Layout}.
        \item \textit{Route Guards} et Intercepteurs JWT.
        \item Système de Toasts/\textit{Snackbars}.
    \end{itemize} & 5 & Haute \\
    \hline
    
    1.8 & Infra / Environnement & 
    \begin{itemize}[leftmargin=*, nosep]
        \item Configuration \texttt{DOCKER Compose}.
        \item Setup \texttt{PostgreSQL} (JPA/Hibernate).
        \item Gestion des secrets via \texttt{.env}.
    \end{itemize} & 3 & Haute \\
\end{longtable}
\addtocounter{table}{-1} % Adjust counter for the manual caption below if needed
\begin{center}
    \small Table \refstepcounter{table}\thetable: Backlog du Sprint 1 \label{tab:backlog_sprint1}
\end{center}
\section{Démarrage du Sprint 1}
Après la phase de planification, l’équipe a entamé la réalisation technique du sprint afin de transformer les exigences du backlog en un incrément logiciel fonctionnel. Cette étape a porté sur les activités de développement, d’intégration et de test, avec une attention particulière accordée à la stabilité architecturale de TalentPredict avant l’intégration des modules d’intelligence artificielle. Cette section présente les travaux réalisés, les résultats obtenus ainsi que les défis techniques rencontrés durant le sprint.

\subsubsection{1. Activités réalisées}

L'équipe a procédé à la mise en place des fondations logicielles et infrastructurelles suivantes :

\begin{itemize}
    \item \textbf{Initialisation du projet :} Création de la structure modulaire du \textit{backend} (\texttt{Spring Boot 3.4.2}) et du \textit{frontend} (\texttt{Angular 17}) avec mise en place des standards de nommage et de codage.
    \item \textbf{Développement de la couche de sécurité :} Implémentation du système d'authentification \texttt{JWT}, incluant la signature des \textit{tokens} et la configuration de \texttt{Spring Security} pour le contrôle d'accès par rôle (\texttt{RBAC}).
    \item \textbf{Conception et déploiement de la base de données :} Modélisation des entités \texttt{User} et \texttt{Profile} différenciés, et exécution des scripts de création via \texttt{JPA}/\texttt{Hibernate} sous \texttt{PostgreSQL}.
    \item \textbf{Intégration des services de notification :} Configuration du serveur \texttt{SMTP} et développement du service asynchrone d'envoi d'emails pour l'enregistrement et la récupération de mot de passe.
    \item \textbf{Conception de l'interface utilisateur (UI) :} Développement de la \textit{Landing Page}, des formulaires d'authentification et de la structure globale du \textit{Dashboard} (\textit{Sidenav}, \textit{Topbar}).
    \item \textbf{Mise en place de l'infrastructure DevOps :} Rédaction et test des configurations \texttt{Docker} pour garantir un environnement de développement reproductible.
\end{itemize}

\subsubsection{2. Livrables produits}

À l'issue de cette itération, les livrables suivants ont été intégrés :

\begin{itemize}
    \item \textbf{Module d'Authentification Fonctionnel :} Un système complet de \textit{Login/Register} sécurisé par \textit{token}.
    \item \textbf{Profils Différenciés :} Des interfaces et des \textit{endpoints} spécifiques pour les données "User" (Employé) et "Admin" (RH).
    \item \textbf{Socle d'Infrastructure Conteneurisé :} Un fichier \texttt{docker-compose.yml} orchestrant l'application, le \textit{frontend} et la base de données.
    \item \textbf{Documentation de l'API :} Une interface \texttt{Swagger UI} exposant tous les \textit{endpoints} réalisés pour faciliter les tests.
    \item \textbf{Composants UI Réutilisables :} Une bibliothèque de composants \texttt{Angular} (Boutons, \textit{Cards}, \textit{Toasts}) basée sur \texttt{Angular Material}.
\end{itemize}

\subsubsection{3. Difficultés rencontrées et traitement}

L'exécution de ce sprint a été marquée par trois défis techniques majeurs :

\begin{itemize}
    \item \textbf{Difficulté 1 : Conflits de CORS}
    \begin{itemize}
        \item \textit{Description :} Le navigateur bloquait les requêtes entre le \textit{frontend} (port 4200) et le \textit{backend} (port 8001).
        \item \textit{Traitement :} Implémentation d'un \textit{bean} de configuration globale \texttt{CorsConfiguration} autorisant explicitement les \textit{headers} et méthodes requis.
    \end{itemize}
    \vspace{2mm}
    \item \textbf{Difficulté 2 : Gestion de l'authentification \textit{Stateless}}
    \begin{itemize}
        \item \textit{Description :} Difficulté à injecter le \textit{token} \texttt{JWT} de manière transparente dans toutes les requêtes \textit{frontend}.
        \item \textit{Traitement :} Création d'un \textit{Intercepteur HTTP} côté \texttt{Angular} pour automatiser l'ajout du \textit{header} \texttt{Authorization} à chaque appel API.
    \end{itemize}
    \vspace{2mm}
    \item \textbf{Difficulté 3 : Parité des environnements avec \texttt{Docker}}
    \begin{itemize}
        \item \textit{Description :} Des comportements différents observés entre le déploiement \texttt{Docker} et l'exécution locale.
        \item \textit{Traitement :} Standardisation des variables d'environnement via des fichiers \texttt{.env} et utilisation de \texttt{Docker Multi-stage builds} pour isoler les étapes de construction.
    \end{itemize}
\end{itemize}
\subsection{Analyse}

Cette phase d'analyse identifie les interactions spécifiques des deux types d'acteurs de TalentPredict et définit les garde-fous techniques nécessaires à la sécurisation de leurs espaces respectifs.

\subsubsection*{1. Identification des Acteurs}
\begin{itemize}
    \item \textbf{L'Employé (Utilisateur) :} C'est l'acteur qui utilise la plateforme pour ses auto-évaluations. Il doit pouvoir s'enregistrer, se connecter et enrichir son profil (liens \texttt{GitHub}/\texttt{LinkedIn}) qui servira de base à l'analyse de ses compétences techniques et comportementales.
    \item \textbf{Le RH (Administrateur) :} C'est l'acteur responsable du suivi et de la gouvernance. Il possède un accès privilégié lui permettant de visualiser l'ensemble des talents de l'entreprise. Dès ce sprint, il doit disposer d'un accès sécurisé à son interface de gestion pour préparer le suivi des futurs résultats.
    \item \textbf{Service Mail}:C'est un système externe automatisé qui interagit avec le backend. Son rôle est critique dans le flux de sécurité : il est chargé d'envoyer les emails transactionnels de manière asynchrone, notamment lors du processus de récupération de compte ("Mot de passe oublié") ou de la confirmation de sécurité, garantissant ainsi l'autonomie et la protection de l'utilisateur.
\end{itemize}

\subsubsection*{2. Clarification du besoin}
Le besoin pour ce premier sprint est de construire le socle de confiance entre l'Employé et le Admin RH :
\begin{itemize}
    \item \textbf{Séparation des espaces :} Garantir que l'Employé n'accède qu'à ses propres données, tandis que le RH accède à une vue globale consolidée.
    \item \textbf{Collecte de données structurées :} Permettre à l'Employé de fournir les vecteurs d'analyse (liens externes) nécessaires aux sprints suivants.
    \item \textbf{Système de confiance :} Mise en place d'une authentification \textit{Stateless} (\texttt{JWT}) pour assurer que chaque action sur la plateforme est tracée et autorisée selon le rôle (\texttt{Admin RH} ou \texttt{Employé}).
\end{itemize}

\subsubsection*{3. Scénarios de validation}
\begin{itemize}
    \item \textbf{Enregistrement de l'Employé :} Création d'un compte, hachage du mot de passe et redirection vers son \textit{dashboard} personnel.
    \item \textbf{Connexion sécurisée :} Saisie des identifiants (Email/Mot de passe) et réception d'un \textit{token} \texttt{JWT} valide pour l'accès aux services.
    \item \textbf{Détention de Rôle :} Tentative d'accès d'un Employé à la zone RH ; le système doit rejeter l'accès (Erreur \texttt{403 Forbidden}).
    \item \textbf{Récupération d'accès :} Envoi d'un email de réinitialisation de mot de passe en cas d'oubli, utilisable par les deux acteurs.
\end{itemize}

\subsubsection*{4. Contraintes techniques}
\begin{itemize}
    \item \textbf{Modèle de données évolutif :} Utilisation de \texttt{PostgreSQL} pour stocker les profils différenciés de l'Employé et du RH.
    \item \textbf{Standard Sécurité \texttt{JWT} :} Utilisation de \textit{tokens} auto-porteurs pour une authentification performante et décentralisée.
    \item \textbf{Interfaces \texttt{Angular Material} :} Développement de deux vues distinctes au sein du même \textit{Shell} applicatif, s'adaptant dynamiquement au rôle détecté dans le \textit{token}.
    \item \textbf{Infrastructure \texttt{Docker} :} Déploiement unifié du \textit{Backend} \texttt{Spring Boot} et du \textit{Frontend} \texttt{Angular} dans des conteneurs isolés et sécurisés.
\end{itemize}
\subsubsection*{5. Diagramme de cas d'utilisation}
Le premier sprint, dédié à la Gestion des Utilisateurs et de l’Authentification, pose les fon-
dations de la plateforme TalentPredict. Ce diagramme de cas d’utilisation \ref{fig:usecase_du_sprint1} synthétise les
fonctionnalités essentielles permettant de sécuriser et de personnaliser l’accès à l’application.

\begin{figure}[H]
    \centering
    \includegraphics[width=\textwidth, keepaspectratio]{images/usecasSprint1.png}
    \caption{Diagramme de cas d'utilisation - Sprint 1 }
    \label{fig:usecase_du_sprint1}
\end{figure}

Le tableau \ref{tab:legende_sprint1} suivant résume la sémantique visuelle et l'application concrète des différentes relations modélisées dans notre diagramme pour ce premier sprint.

\begin{table}[H]
    \centering
    \renewcommand{\arraystretch}{1.5} % Aère les lignes du tableau
    \begin{tabularx}{\textwidth}{@{} >{\raggedright\arraybackslash}p{3cm} >{\raggedright\arraybackslash}p{4.5cm} X @{}}
        \toprule
        \rowcolor{TableHeaderBlue}
        \color{white}\textbf{Type de Relation} & \color{white}\textbf{Représentation Graphique} & \color{white}\textbf{Application Concrète } \\
        \midrule
        \textbf{Généralisation} \newline \textit{\small{(Héritage)}} & 
        Trait plein continu + flèche triangulaire vide pointant vers l'acteur parent. & 
        Les acteurs Employé et Admin RH héritent de l'acteur de base Utilisateur et de tous ses accès par défaut (ex: \textit{Se connecter}). \\
        
        \textbf{Association} \newline \textit{\small{(Communication)}} & 
        Ligne continue simple. & 
        Lien direct entre l'Utilisateur et les fonctions de base, ou le Service Mail pour l'envoi de liens. \\
        
        \textbf{Inclusion} \newline \textit{\small{(«include»)}} & 
        Flèche en pointillés avec le mot-clé \texttt{«include»}, orientée vers le cas inclus. & 
        L'action \textit{Mot de passe oublié} déclenche obligatoirement \textit{Envoyer un lien de récupération}. \\
        
        \textbf{Extension} \newline \textit{\small{(«extend»)}} & 
        Flèche en pointillés avec le mot-clé \texttt{«extend»}, orientée vers le cas de base. & 
        L'action \textit{Modifier le profil} offre des sous-options supplémentaires comme \textit{Gérer les préférences} ou \textit{Changer le mot de passe}. \\
        \bottomrule
    \end{tabularx}
    \caption{Légende des associations du diagramme  cas d’utilisation - Sprint 1}
    \label{tab:legende_sprint1}
\end{table}
\subsection{Conception}
Cette phase de conception, réalisée après l’identification des acteurs et des besoins fonctionnels, vise à structurer l’architecture du système, modéliser les données et définir les interfaces de communication (API). Elle permet de mettre en place un socle technique robuste et sécurisé pour soutenir l’intégration des futurs modules d’intelligence artificielle de TalentPredict.

\subsubsection*{1. Objectifs de conception}
\begin{itemize}
    \item \textbf{\textit{Stratégie de Monolithe Modulaire } :} Afin d'éviter la complexité prématurée des microservices tout en préparant le futur, nous avons opté pour un monolithe modulaire. Chaque fonctionnalité (Auth, Évaluation, IA) est encapsulée dans son propre \textit{package}, communiquant via des services définis, respectant ainsi le principe de responsabilité unique (SRP).
    \item \textbf{\textit{Authentification sans état (Stateless) et Montée en charge (Scalability)} :} En choisissant \texttt{JWT}, nous éliminons le besoin de \texttt{HttpSession} côté serveur. Cela libère de la mémoire RAM et permet au \textit{backend} de monter en charge sans se soucier de la réplication des sessions.
    \item \textbf{Intégrité des données (\textit{Data Integrity}) :} Assurer que chaque utilisateur possède un profil unique dès sa création via des contraintes de clés étrangères au niveau de la base de données.
\end{itemize}

\subsubsection*{2. Architecture Technique }
L’architecture de TalentPredict repose sur une structure N-Tier (multicouche) moderne, conçue pour être modulaire, sécurisée et facilement déployable. \ref{fig:archi_globale_s1} expose l’organisation globale de l’infrastructure technique ainsi que l’articulation des flux de données entre les différents composants :
\begin{figure}[H]
    \centering
    \includegraphics[width=0.8\textwidth]{images/Diagramme d'Architecture.png}
    \vspace{2mm}
    \caption{Architecture globale et flux de données du Sprint 1}
    \label{fig:archi_globale_s1}
\end{figure}

\subsubsection*{Infrastructure et environnement \texttt{Docker}}
\begin{itemize}
    \item L’ensemble du système est orchestré via \texttt{Docker Compose}, garantissant la portabilité entre les environnements de développement et de production ainsi qu’une architecture multi-conteneurs.
    \item Un réseau \texttt{Docker} privé (\texttt{talentpredict-net}) assure une communication sécurisée et isolée entre les services, empêchant tout accès direct à la base de données depuis l’extérieur. L’accès aux données est exclusivement contrôlé par le \textit{backend}.
    \item La persistance des données est assurée par un \textit{Docker Volume} associé au conteneur \texttt{PostgreSQL}, garantissant la conservation des informations même après arrêt ou redémarrage des services.
\end{itemize}

\subsubsection*{Couche Client (\textit{Frontend Tier})}
\begin{itemize}
   \item L’interface utilisateur est développée avec \texttt{Angular 17} et déployée via \texttt{Nginx} sous forme d’application SPA (\textit{Single Page Application}).
    \item Un intercepteur HTTP \texttt{Angular} assure l’injection automatique du jeton \texttt{JWT} dans l’en-tête \texttt{Authorization}, permettant l’accès sécurisé aux ressources protégées.
\end{itemize}

\subsubsection*{Couche Application (\textit{Backend Tier})}
\begin{itemize}
    \item Le cœur métier est développé avec \texttt{Spring Boot 3} selon une architecture MVC enrichie.
\item Les contrôleurs REST exposent les \textit{endpoints} de l’API et utilisent des DTOs afin d’assurer une séparation claire des couches, tandis que la logique métier est encapsulée dans des services tels que \texttt{AuthService} et \texttt{UserService}.
\item La sécurité est renforcée par l’utilisation de \texttt{BCrypt} pour le hachage des mots de passe. L’accès à la base de données est géré via \texttt{Spring Data JPA}, réduisant les risques d’injection SQL grâce aux requêtes paramétrées.
\end{itemize}

\subsubsection*{Couche Données (\textit{Data Tier})}
La persistance des données est assurée par \texttt{PostgreSQL 16}, choisi pour sa robustesse, sa fiabilité et ses capacités avancées de gestion des transactions. Le modèle de données du Sprint 1 est principalement centré sur les entités liées aux utilisateurs, aux rôles et à l’authentification.
\subsubsection*{Services externes (SMTP)}
Le \textit{backend} communique avec un service \texttt{SMTP} externe (tel que \texttt{Gmail SMTP}) afin de gérer l’envoi d’e-mails transactionnels, notamment pour la réinitialisation des mots de passe. Ce mécanisme améliore la sécurité et l’autonomie des utilisateurs.

\subsubsection*{3. Modèle de données }
Le diagramme \ref{fig:schema_donnees_s1} présente le modèle de données du Sprint 1, couvrant les entités nécessaires à la gestion des utilisateurs, des profils, des rôles et de la sécurité d'authentification.

\begin{figure}[H]
    \centering
    % Remplacez 'chemin_schema_donnees.png' par le nom exact de votre image
    \includegraphics[width=1\textwidth]{images/ERDsprint1.png}
    \vspace{2mm}
    \caption{Diagramme entité-association et schéma de données du Sprint 1}
    \label{fig:schema_donnees_s1}
\end{figure}

Le modèle physique de données (ERD) est composé des entités suivantes :
\begin{itemize}
    \item \textbf{Table \texttt{ROLE}} : Stocke les rôles de sécurité applicatifs pour le contrôle d'accès (RBAC).
    \\ \textit{Points clés} : Unicité stricte du libellé via le champ \texttt{name} (UK) pour interdire les doublons sur \texttt{ROLE\_USER} et \texttt{ROLE\_ADMIN}.
    
    \item \textbf{Table \texttt{USER}} : Gère l'authentification globale et l'activation logique des comptes.
    \\ \textit{Points clés} : Identifiant de connexion unique (\texttt{email} UK), mot de passe sécurisé (\texttt{password\_hash} haché par BCrypt), contact (\texttt{phone}) et statut du compte (\texttt{is\_active} pour le blocage et la conformité RGPD).
    
    \item \textbf{Table \texttt{PROFILE}} : Contient les métadonnées professionnelles et les liens d'intégration (LinkedIn, GitHub, CV).
    \\ \textit{Points clés} : Clé étrangère unique \texttt{user\_id} (FK, UK) garantissant une relation de composition exclusive 1-to-1 (un profil par utilisateur).
    
    \item \textbf{Table \texttt{PASSWORD\_RESET\_TOKEN}} : Sécurise le processus de réinitialisation des mots de passe oubliés.
    \\ \textit{Points clés} : Jeton d'authentification unique (\texttt{token} UK) sous format UUID et sécurité temporelle via la date de fin de validité (\texttt{expiry\_date}).
\end{itemize}


Ce modèle garantit une structure normalisée, évolutive et adaptée aux exigences de sécurité ainsi qu’aux futurs modules d’analyse intelligente.
\subsubsection{4. Hachage des Mots de Passe}

Pour le stockage des mots de passe en base de données, nous avons retenu l'algorithme BCrypt avec un facteur de coût de 10, valeur par défaut recommandée par Spring Security. Ce choix s'appuie sur une analyse comparative de plusieurs standards de hachage, détaillée dans le Tableau \ref{tab:comparaison_hachage} ci-dessous:

\begin{table}[H]
    \centering
    \renewcommand{\arraystretch}{1.3} % Aère les lignes du tableau pour une meilleure lisibilité
    \begin{tabular}{|l|c|c|c|}
        \hline
        \rowcolor{TableHeaderBlue}
        \color{white}\textbf{Critère} & \color{white}\textbf{BCrypt (coût 10)} & \color{white}\textbf{MD5} & \color{white}\textbf{SHA-256} \\
        \hline
        Résistance force brute & Élevée & Faible & Faible \\
        \hline
        Sel automatique (\textit{Salt}) & Oui & Non & Non \\
        \hline
        Adaptatif & Oui & Non & Non \\
        \hline
        Temps de hachage & $\sim$100ms & $<$1ms & $<$1ms \\
        \hline
    \end{tabular}
    \caption{Comparaison des algorithmes de hachage de mots de passe}
    \label{tab:comparaison_hachage}
\end{table}

Cette implémentation garantit un haut niveau de sécurité pour trois raisons fondamentales :
\begin{itemize}
    \item \textbf{Protection contre la force brute :} Le facteur de coût 10 génère environ $2^{10}$ (soit 1024) itérations de hachage. Cela rend les attaques par dictionnaire et par force brute économiquement non viables en termes de temps de calcul pour un attaquant.
    \item \textbf{Résistance aux Rainbow Tables :} Un sel cryptographique (\textit{salt}) est généré automatiquement et de manière unique pour chaque mot de passe lors de sa création, éliminant ainsi les vulnérabilités liées aux tables d'empreintes précalculées.
    \item \textbf{Pérennité (Adaptabilité) :} L'aspect adaptatif de l'algorithme permet d'augmenter le facteur de coût dans les versions futures de la plateforme pour suivre l'évolution de la puissance de calcul matérielle, et ce, sans nécessiter de refonte de l'architecture de sécurité.
\end{itemize}

\subsubsection*{5. Contrats API }

Afin d’assurer l’interopérabilité entre le Frontend Angular et le Backend Spring Boot, des contrats d’interface clairs ont été définis. L’exemple suivant présente le service d’authentification, qui constitue le point d’accès sécurisé de l’application:
\vspace{2mm}
\textbf{Point d'accès :} \texttt{POST /api/auth/login}

\textbf{Entrée (\textit{Request Body}) :}
\begin{lstlisting}[language=json, caption={Structure de la requete de connexion}]
{
  "email": "user@talentpredict.exemple",
  "password": "securePassword123"
}
\end{lstlisting}

\textbf{Sortie (\textit{Response} - Success \texttt{200 OK}) :}
\begin{lstlisting}[language=json, caption={Structure de la reponse d'authentification (JWT)}]
{
  {
  "token": "eyJhbGciOiJIUzI1NiIsInR5cCI6IkpXVCJ9...",
  "type": "Bearer",
  "id": "550e8400-e29b-41d4-a716-446655440000",
  "email": "user@talentpredict.exemple",
  "role": "ROLE_USER",
  "nom": "Ben Ali",
  "prenom": "Wajih",
  "redirectUrl": "/dashboard",
  "emailVerified": true
}

}
\end{lstlisting}

\vspace{2mm}
\textbf{Précisions techniques :}
\begin{itemize}
    \item \textbf{\texttt{token}} : Jeton \texttt{JWT} signé par le serveur (\texttt{HS256}) permettant l'accès sécurisé aux ressources protégées.
    \item \textbf{\texttt{type}} : Identifie le schéma d'authentification utilisé (\textit{Bearer Authentication}).
    \item \textbf{\texttt{role}} : Information cruciale permettant au \textit{Frontend} \texttt{Angular} d'adapter dynamiquement l'affichage et de configurer les gardes de navigation (\textit{Route Guards}).
    \item \textbf{\texttt{nom} / \texttt{prenom}} : Données d'identité permettant une personnalisation immédiate de l'interface utilisateur dès la connexion.
    \item \textbf{\texttt{redirectUrl}} : \texttt{URL} de redirection suggérée par le serveur en fonction du profil ou du statut de l'utilisateur.
    \item \textbf{\texttt{emailVerified}} : Indicateur de sécurité permettant au \textit{Frontend} d'inviter l'utilisateur à vérifier son compte si nécessaire.
\end{itemize}


\subsubsection*{6. Gestion des Sessions et Sécurité JWT}
\subsubsection{Architecture d'Authentification \textit{Stateless} et \texttt{JWT}}

L’architecture ci-dessous \ref{fig:auth_flow} de TalentPredict adopte une authentification \textit{Stateless} basée sur les jetons \texttt{JWT} (\textit{JSON Web Token}), complétée par une stratégie de Rotation de Jeton (\textit{Token Rotation}) pour une sécurité maximale.
\begin{figure}[H]
    \centering
    \includegraphics[width=0.88\textwidth,height=0.48\textheight,keepaspectratio]{images/archdesecSprint1.png}
    \vspace{-2mm}
    \caption{Flux de la requête d'authentification et de sécurité via JWT}
    \label{fig:auth_flow}
\end{figure}
\vspace{-4mm}
\paragraph*{Structure et Double Jeton (\textit{Dual Token Strategy})}
Le système ne repose pas sur un seul jeton, mais sur un duo permettant d'équilibrer sécurité et expérience utilisateur :
\begin{itemize}[noitemsep, topsep=2pt]
    \item \textbf{\textit{Access Token} (12 heures) :} Un jeton à durée de vie très courte utilisé pour signer chaque requête \texttt{API}. Il contient les rôles RBAC (\textit{User}/\textit{Admin}).
    \item \textbf{\textit{Refresh Token} (7 jours) :} Un jeton persistant stocké de manière sécurisée, permettant de renouveler l'\textit{Access Token} sans obliger l'utilisateur à se reconnecter manuellement.
\end{itemize}
\paragraph*{Mécanisme \textit{Stateless} et Rotation}
Côté \textit{backend}, aucune session n’est conservée en mémoire vive. Pour chaque requête, le filtre \texttt{JwtAuthFilter} (illustré dans la Figure \ref{fig:auth_flow}) valide la signature (\texttt{HMAC-SHA256}). Nous avons implémenté la \textit{Token Rotation} : chaque demande de nouvel \textit{Access Token} invalide l'ancien \textit{Refresh Token} et en génère un nouveau, empêchant ainsi les attaques par rejeu (\textit{Replay Attacks}).
\paragraph*{Sécurité et RBAC}
Les rôles encodés dans le jeton contrôlent l’accès aux ressources. Toute tentative d’accès non autorisé est rejetée avec une erreur HTTP \texttt{403 Forbidden}. Le stockage côté client est géré de manière à minimiser les risques d'attaques \texttt{XSS}, tandis que les échanges sont protégés par le protocole \texttt{HTTPS}.
\subsubsection*{7. Diagramme de seqence }
\subsubsection{Cas 1 : Cycle de vie de l'Employé}
Ce cas modélise l'ensemble des interactions de l'employé avec la plateforme, depuis son inscription initiale jusqu'à la gestion sécurisée de son compte personnel.
\paragraph{Cas 1A : Inscription, Session et Récupération complèt :}
La figure~\ref{fig:seq_employe_public} représente le diagramme de séquence relatif à l'inscription, la session et la récupération complète du mot de passe pour l'employé.

\begin{figure}[H]
    \centering
    % Remplace le chemin ci-dessous par le nom réel de ton image
    \includegraphics[width=\textwidth]{images/diagramme_3A.png} 
    \caption{Diagramme de séquence : Inscription, Session et Récupération complète (Employé)}
    \label{fig:seq_employe_public}
\end{figure}

Ce diagramme illustre les processus publics accessibles à l'employé avant son authentification. Il se divise en trois phases principales :
\begin{itemize}
    \item \textbf{Phase d'inscription :} L'utilisateur soumet ses informations via l'interface Angular. Le contrôleur délègue la vérification à la couche de service pour s'assurer que l'email ou le numéro de téléphone n'existe pas déjà. Si les données sont valides, le mot de passe est haché à l'aide de BCrypt, le profil est inséré dans la base de données PostgreSQL, et un email de bienvenue est expédié.
    \item \textbf{Phase de récupération de mot de passe :} En cas d'oubli, l'employé demande une réinitialisation. Le système vérifie l'existence du compte, génère un jeton unique (token), l'enregistre en base, et envoie un lien sécurisé par email. Une fois le lien cliqué et le nouveau mot de passe soumis, le système valide le jeton et met à jour le mot de passe haché.
    \item \textbf{Phase de connexion et déconnexion :} Lors de l'authentification, les identifiants sont vérifiés (comparaison BCrypt). En cas de succès, un jeton JWT est généré et stocké localement côté frontend. La déconnexion consiste en la suppression de ce jeton du stockage local.
\end{itemize}

\vspace{0.5cm}
\paragraph{Cas 1B : Profil, Préférences et Mot de pass :}
La figure~\ref{fig:seq_employe_jwt} représente le diagramme de séquence détaillant la gestion du profil, des préférences et du mot de passe de l'employé sous authentification sécurisée.

\begin{figure}[H]
    \centering
    % Remplace le chemin ci-dessous par le nom réel de ton image
    \includegraphics[width=\textwidth]{images/diagramme_3B.png} 
    \caption{Diagramme de séquence : Profil, Préférences et Mot de passe (Employé - JWT)}
    \label{fig:seq_employe_jwt}
\end{figure}

Ce diagramme détaille les opérations sécurisées nécessitant une authentification préalable. Toutes les requêtes de ce flux incluent le jeton JWT dans l'en-tête HTTP.
\begin{itemize}
    \item \textbf{Interception de sécurité :} Avant d'atteindre le contrôleur, chaque requête traverse le \texttt{SecurityFilter} qui valide la signature et l'expiration du jeton JWT.
    \item \textbf{Gestion du Profil et Préférences :} Une fois le jeton validé, l'employé peut mettre à jour ses informations personnelles ou ses préférences de notification. Le contrôleur transmet les données au \texttt{ProfileService} qui exécute les requêtes \texttt{UPDATE} correspondantes en base de données.
    \item \textbf{Changement de mot de passe interne :} L'employé saisit son ancien et son nouveau mot de passe. Le service récupère le mot de passe actuel en base de données, le compare avec l'ancien mot de passe saisi via BCrypt. Si la correspondance est validée, le nouveau mot de passe est haché et enregistré.
\end{itemize}

\subsubsection{Cas 2 : Gouvernance et Accès de l'Administrateur RH}
Ce cas modélise les flux spécifiques à l'Administrateur RH, en mettant l'accent sur la vérification des privilèges d'accès et les opérations de gouvernance globale
\paragraph{Cas 2A :  Authentification, Gouvernance et Récupération complète :}
La figure~\ref{fig:seq_admin_public} représente le diagramme de séquence relatif à l'authentification, la gouvernance globale et la récupération du mot de passe pour l'administrateur RH.

\begin{figure}[H]
    \centering
    % Remplace le chemin ci-dessous par le nom réel de ton image
    \includegraphics[width=\textwidth]{images/diagramme_4A.png} 
    \caption{Diagramme de séquence : Authentification, Gouvernance et Récupération complète (Admin RH)}
    \label{fig:seq_admin_public}
\end{figure}

Ce diagramme modélise l'accès et les actions globales de l'Administrateur RH, en mettant l'accent sur les restrictions de rôles.
\begin{itemize}
    \item \textbf{Récupération et Connexion (Phases publiques) :} Le processus de mot de passe oublié est similaire à celui de l'employé. Cependant, lors de la phase de connexion, la requête SQL vérifie non seulement la validité des identifiants (via BCrypt), mais s'assure également que l'utilisateur possède explicitement le rôle \texttt{ROLE\_ADMIN}. Si ce rôle est absent, une exception \texttt{BadCredentialsException} est levée.
    \item \textbf{Gouvernance et Tableau de bord (Phase sécurisée) :} Pour accéder aux statistiques globales, la requête GET est interceptée par le \texttt{SecurityFilter} qui extrait le JWT et valide la présence du rôle Administrateur. En cas d'autorisation, l'\texttt{AdminService} exécute les requêtes d'agrégation (comme le comptage des utilisateurs) sur la base PostgreSQL et retourne un objet DTO formaté en JSON pour peupler le tableau de bord côté Angular.
\end{itemize}

\vspace{0.5cm}
\paragraph{Cas 2B :  Profil, Préférences et Mot de passe :}
La figure~\ref{fig:seq_admin_jwt} représente le diagramme de séquence illustrant la mise à jour des informations organisationnelles et des paramètres de sécurité de l'administrateur RH.

\begin{figure}[H]
    \centering
    % Remplace le chemin ci-dessous par le nom réel de ton image
    \includegraphics[width=0.8\textwidth]{images/diagramme_4B.png} 
    \caption{Diagramme de séquence : Profil, Préférences et Mot de passe (Admin RH - JWT)}
    \label{fig:seq_admin_jwt}
\end{figure}

Ce diagramme présente la gestion du compte personnel de l'Administrateur RH.
\begin{itemize}
    \item \textbf{Mécanisme :} Le flux est techniquement identique à celui de l'employé (Diagramme 3.B), reposant sur l'en-tête d'autorisation Bearer et la validation du JWT par le \texttt{SecurityFilter}.
    \item \textbf{Opérations :} Il couvre la mise à jour des informations organisationnelles de l'administrateur, la gestion de ses propres préférences de sécurité, ainsi que la procédure stricte de modification de son mot de passe nécessitant la vérification préalable de l'ancien mot de passe haché.
\end{itemize}

\subsubsection*{8. Diagramme de Classes }
Le diagramme de classes \ref{fig:class_diagram_sprint1} illustre l'architecture logicielle du module d'authentification et de gestion de profil, mis en place lors du premier sprint du projet TalentPredict. Il s'appuie sur une conception robuste en couches (N-tiers), caractéristique des applications modernes basées sur Spring Boot. Cette approche architecturale garantit une séparation stricte des responsabilités (\textit{Separation of Concerns}), facilitant ainsi la maintenance du code, l'évolutivité du système et l'application rigoureuse des principes de \textit{Security-by-Design}.
\begin{figure}[H]
    \centering
    % Remplacez 'chemin_diagramme_classes_sprint1.png' par le nom exact de votre image
    \includegraphics[width=\textwidth]{images/classSprint1.png.png}
    \vspace{2mm}
    \caption{Diagramme de classes : Architecture en couches du Sprint 1}
    \label{fig:class_diagram_sprint1}
\end{figure}

L'architecture est structurée autour de six composants ou couches principales, interagissant de manière unidirectionnelle pour traiter les requêtes des utilisateurs :

\begin{itemize}
    \item \textbf{Couche Modèle (Entités) :} Constitue le cœur du domaine de l'application. Elle regroupe les objets persistants mappés directement aux tables de la base de données (via JPA/Hibernate), tels que \texttt{User}, \texttt{Profile}, \texttt{PasswordResetToken} et \texttt{UserPrivacySettings}. Les énumérations, comme \texttt{Role} (USER, ADMIN), y sont également définies pour normaliser les données.
    
    \item \textbf{Couche Persistance (Repositories) :} Agit comme l'interface de communication avec la base de données. En étendant les interfaces de Spring Data JPA (comme \texttt{UserRepository} ou \texttt{ProfileRepository}), cette couche encapsule la logique d'accès aux données (opérations CRUD) et permet de récupérer ou manipuler les entités sans exposer les requêtes SQL complexes aux couches supérieures.
    
    \item \textbf{Couche DTOs (Data Transfer Objects) :} Joue un rôle de bouclier de protection et de formatage. Les objets tels que \texttt{LoginRequest}, \texttt{RegisterRequest} ou \texttt{AuthResponse} sont utilisés pour transporter strictement les données nécessaires entre le client (frontend Angular) et le serveur, évitant ainsi d'exposer directement la structure interne des entités de la base de données.
    
    \item \textbf{Couche Sécurité (Spring Security) :} Centralise les mécanismes de protection des ressources. Elle intègre les composants essentiels pour une authentification \textit{stateless}, notamment \texttt{JwtService} pour la génération et la validation des tokens JWT, et \texttt{JwtAuthenticationFilter} qui intercepte chaque requête entrante pour vérifier les habilitations de l'utilisateur avant d'autoriser l'accès aux API.
    
    \item \textbf{Couche Métier (Services) :} Contient la logique fonctionnelle et les règles métier de l'application. Les classes comme \texttt{AuthServiceImpl} et \texttt{ProfileService} orchestrent les opérations complexes (inscription, hachage des mots de passe, envoi d'e-mails de vérification, gestion des tokens de réinitialisation) en faisant le lien entre les contrôleurs et les repositories.
    
    \item \textbf{Couche API (Controllers) :} Représente le point d'entrée du backend (API REST). Les contrôleurs (\texttt{AuthController}, \texttt{ProfileController}) interceptent les requêtes HTTP (GET, POST, PUT), délèguent le traitement à la couche métier appropriée, puis formatent la réponse finale sous forme de \texttt{ResponseEntity} pour la renvoyer au client.
\end{itemize}




\subsection{Codage}

Cette sous-section détaille le passage de la conception à l'implémentation logicielle. Pour ce premier sprint, l'objectif principal était de bâtir un socle technologique résilient, sécurisé et prêt à monter en charge (\textit{Scalable}) pour accueillir les futurs modules d'intelligence artificielle.

\paragraph*{1. Écosystème Technologique (\textit{Stack} Technique)}
La sélection des technologies s'est portée sur des standards de l'industrie offrant un haut niveau de sécurité et de maintenabilité :

\begin{itemize}
    \item \textbf{Couche Métier et Sécurité (\textit{Backend}) :} Propulsée par \texttt{Java 17} et le \textit{framework} \texttt{Spring Boot 3.4.2}. La persistance des entités est orchestrée par \texttt{Spring Data JPA} (\texttt{Hibernate}). La sécurisation des \textit{endpoints} repose sur \texttt{Spring Security} couplé à la librairie \texttt{JJWT} pour la cryptographie des jetons (\texttt{HMAC-SHA256}).
    \item \textbf{Couche Présentation (\textit{Frontend}) :} Implémentée avec \texttt{Angular 17} (\texttt{TypeScript}). Nous avons privilégié une architecture réactive (\texttt{RxJS}) avec une interface utilisateur adaptative basée sur des composants UI modernes.
    \item \textbf{Persistance Relationnelle :} \texttt{PostgreSQL 16}, sélectionné pour sa robustesse transactionnelle (ACID) et son extensibilité (capacités futures de recherche vectorielle ou \texttt{JSONB} pour l'IA).
\end{itemize}

\subsubsection{2. Configuration du Service de Messagerie (Gmail SMTP)}

Pour assurer la fiabilité des communications transactionnelles (activation de compte, réinitialisation de mot de passe, alertes système), TalentPredict intègre le service \texttt{Gmail SMTP}. La configuration s'appuie sur le serveur \texttt{smtp.gmail.com} tilisant le protocole sécurisé \texttt{STARTTLS} sur le port \texttt{587}.

\begin{lstlisting}[language=properties, caption=Configuration SMTP dans le fichier application.properties]
spring.mail.host=smtp.gmail.com
spring.mail.port=587
spring.mail.username=${MAIL_USERNAME}
spring.mail.password=${MAIL_PASSWORD}
spring.mail.properties.mail.smtp.auth=true
spring.mail.properties.mail.smtp.starttls.enable=true
\end{lstlisting}
Les identifiants sont externalisés dans le fichier \texttt{.env} afin de sécuriser les secrets et d’éviter toute exposition lors du versionnage. L’envoi des e-mails est traité de manière asynchrone via l’annotation \texttt{@Async}, ce qui permet de ne pas bloquer le \textit{thread} principal et d’optimiser le temps de réponse de l’API. Par ailleurs, un \textit{App Password} Google est utilisé en remplacement du mot de passe principal du compte, conformément aux bonnes pratiques de sécurité.

\subsubsection*{3. Organisation Modulaire et \textit{Clean Architecture}}Le projet adopte une architecture de type Monolithe Modulaire en remplacement d’une approche monolithique classique. Le backend est structuré en domaines fonctionnels indépendants (ex. \texttt{modules.auth}), respectant strictement le principe de responsabilité unique (SRP).

\begin{itemize}
\item \textbf{\texttt{Controllers} :} Responsables de la gestion des requêtes HTTP et du routage.
\item \textbf{\texttt{Services} :} Contiennent la logique métier ainsi que l’orchestration des traitements.
\item \textbf{\texttt{Repositories} :} Assurent l’accès aux données via une couche d’abstraction.
\item \textbf{\texttt{DTOs} (\textit{Data Transfer Objects}) :} Structures immuables servant d’interface entre le frontend et le backend, limitant les risques liés au \textit{Mass Assignment}.
\end{itemize}

Côté Angular, l’architecture suit le pattern Core/Shared/Feature :
\begin{itemize}
\item Le dossier \texttt{core/} regroupe les éléments transverses tels que les intercepteurs HTTP (injection du \texttt{JWT}) et les \textit{guards} de sécurité.
\item Le dossier \texttt{modules/} contient les fonctionnalités applicatives isolées (ex. : \textit{Login}, \textit{Dashboard}).
\end{itemize}

\subsubsection*{4. Fonctionnalités Implémentées}
\begin{itemize}
    \item \textbf{Authentification \textit{Stateless} (\texttt{JWT}) :} Mise en place d'un flux de connexion où l'état de session n'est pas conservé sur le serveur. Les mots de passe sont hachés via \texttt{BCrypt} (avec ajout d'un "sel" cryptographique) pour contrer les attaques par dictionnaire.
    \item \textbf{Contrôle d'Accès Basé sur les Rôles (RBAC) :} Configuration de la sécurité bloquant l'accès aux routes d'administration (\texttt{/api/admin/**}) pour tout utilisateur ne possédant pas le \texttt{ROLE\_ADMIN} (renvoi d'une erreur HTTP \texttt{403}).
    \item \textbf{Collecte des Vecteurs IA (Profil) :} Développement des interfaces permettant aux employés de lier dynamiquement leurs URLs professionnelles (\texttt{GitHub}, \texttt{LinkedIn}, \texttt{Portfolio}) et de téléverser leur CV, données indispensables pour l'extraction de compétences prévue au Sprint 2.
\end{itemize}

\subsubsection*{5. Ingénierie DevOps et Conteneurisation}
Afin de garantir une parité absolue entre les environnements de développement et de production, une stratégie d’\textit{Infrastructure as Code} (IaC) a été adoptée dès l'initialisation du projet.

\subsubsection*{A. Orchestration de la Base de Données avec Docker}
Le SGBD PostgreSQL est instancié via un fichier \texttt{docker-compose.yml}. Ce dernier définit les limites de ressources, les réseaux isolés et surtout le montage des volumes locaux (\texttt{volumes}) pour assurer la persistance des données. Cette approche garantit qu'en cas de destruction ou de redémarrage du conteneur, aucune donnée utilisateur n'est perdue.

\subsubsection{B. Sécurisation des Configurations via le fichier \texttt{.env}}
Conformément aux meilleures pratiques de sécurité, aucun secret n'est stocké en dur dans le code source ou versionné sur Git. L'application consomme ses paramètres sensibles (identifiants de base de données, clés cryptographiques, configuration SMTP) via un fichier d'environnement local \texttt{.env}.

Voici un aperçu de la structure du fichier \texttt{.env} utilisé pour notre backend :

\begin{lstlisting}[language=bash, caption={Structure du fichier de configuration sécurisé (.env)}, label={lst:env_file}]
# --- DATABASE CONFIGURATION ---
DB_URL=jdbc:postgresql://localhost:5432/talentpredict_db
DB_USERNAME=tp_admin
DB_PASSWORD=******** # Secret masque

# --- SECURITY (JWT) ---
JWT_SECRET=4a2b6c8d9e0f1g2h3i4j5k6l7m8n9o0p1q2r3s4t5u6v7w8x9y0z
JWT_EXPIRATION=86400000  # 24 Heures

# --- MAIL SERVER (SMTP) ---
SMTP_HOST=smtp.gmail.com
SMTP_PORT=587
SMTP_USERNAME=noreply.talentpredict@gmail.com
SMTP_PASSWORD=app_password_security
\end{lstlisting}

Cette approche repose sur deux piliers fondamentaux de l'ingénierie logicielle :
\begin{itemize}
    \item \textbf{Séparation stricte du Code et de la Configuration :} Le code source reste identique quel que soit l'environnement (Développement, Staging, Production). Seul le fichier \texttt{.env} diffère, ce qui facilite grandement le déploiement continu.
    \item \textbf{Sécurité et Confidentialité :} L'inclusion du fichier \texttt{.env} dans le fichier \texttt{.gitignore} prévient toute fuite accidentelle d'identifiants sur des dépôts distants (GitHub/GitLab), protégeant ainsi l'infrastructure contre les accès non autorisés.
\end{itemize}

\subsubsection*{C. Sécurisation des Échanges via CORS}
Enfin, pour protéger l'API REST contre les attaques de type \textit{Cross-Site Request Forgery} (CSRF), un paramétrage strict du filtre CORS (\textit{Cross-Origin Resource Sharing}) a été appliqué au niveau de Spring Boot. Ce filtre n'autorise que les requêtes provenant de l'origine de confiance de notre application front-end (\texttt{http://localhost:4200}), bloquant systématiquement toute tentative d'interaction provenant de domaines non reconnus.

\subsubsection*{6. Assurance Qualité et Validation des Contrats API}
Avant toute connexion au client \texttt{Angular}, les APIs ont été éprouvées selon des méthodes d'ingénierie rigoureuses :

\begin{itemize}
    \item \textbf{Documentation OpenAPI (\texttt{Swagger 3}) :} \ref{fig:swagger_ui_s1} Intégration de \texttt{springdoc-openapi} pour exposer dynamiquement un portail interactif listant les schémas de requêtes/réponses et les codes HTTP attendus.

\begin{figure}[H]
    \centering
    % Remplacez 'chemin_swagger_ui.png' par le nom exact de votre image
    \includegraphics[width=0.8\textwidth]{images/swager.png}
    \vspace{2mm}
    \caption{Documentation interactive des APIs REST via Swagger UI}
    \label{fig:swagger_ui_s1}
\end{figure}
La documentation des API REST a été générée automatiquement grâce à l’intégration de \texttt{Springdoc OpenAPI} (\texttt{Swagger 3}). Cette interface permet de visualiser l’ensemble des \textit{endpoints} exposés par le backend Spring Boot, ainsi que leurs paramètres, les schémas \texttt{JSON} attendus et les codes de réponse HTTP associés.

\texttt{Swagger UI} facilite également le test des routes sécurisées en permettant l’injection manuelle du jeton \texttt{JWT} dans l’en-tête \texttt{Authorization}. Cette fonctionnalité améliore les phases de validation, de débogage et d’intégration entre le frontend Angular et le backend REST.

L’utilisation d’OpenAPI contribue par ailleurs à une meilleure maintenabilité du projet grâce à une documentation centralisée et automatiquement synchronisée avec le code source.

\begin{itemize}
\item \textbf{Tests automatisés (\texttt{Postman}) :} Mise en place de collections \texttt{Postman} intégrant des scripts de pré-requête (Pre-request scripts). Ces scripts permettent de récupérer automatiquement le jeton \texttt{JWT} généré lors de l’authentification (\texttt{/login}) et de l’injecter dans l’en-tête \texttt{Authorization: Bearer}, afin de tester les routes protégées dans des conditions proches d’un usage réel.
\end{itemize}

\end{itemize}
\subsection{Realisation}
ette section présente les interfaces majeures développées pour assurer l'interaction avec le socle de sécurité et de gestion des profils.

\subsubsection*{1. La Page d'Accueil (Landing Page)}
L’interface de \textit{Page d'Accueil} \ref{fig:ui_landing_s1} constitue la vitrine technologique de TalentPredict. Son rôle est de présenter la proposition de valeur de la plateforme : l'utilisation de l'intelligence artificielle pour l'évaluation prédictive des talents. Conçue avec une approche \textit{Mobile-First}, elle offre une navigation fluide vers les modules d'authentification. Elle utilise des composants visuels modernes (dégradés, animations légères) pour instaurer une image de marque innovante et professionnelle dès le premier contact avec l'utilisateur.

\begin{figure}[H]
    \centering
    \includegraphics[width=\textwidth]{images/landingpage.png}
    \vspace{2mm}
    \caption{Interface de la Landing Page de TalentPredict}
    \label{fig:ui_landing_s1}
\end{figure}

\subsubsection*{2. Interface d'Inscription (Register)}

L'interface d'inscription \ref{fig:ui_register_s1} est le point d'entrée pour les nouveaux employés. Elle repose sur un formulaire intelligent utilisant les \texttt{Angular Reactive Forms}, permettant une validation instantanée des données saisies (unicité de l'email, robustesse du mot de passe). Une fois le formulaire soumis, l'application communique avec l'\textit{endpoint} \texttt{/api/auth/register} du \textit{backend} \texttt{Spring Boot}, qui procède à la création sécurisée du compte et à l'initialisation automatique d'un profil vide, prêt à être enrichi par le candidat.

\begin{figure}[H]
    \centering
    \includegraphics[width=\textwidth]{images/signin.png}
    \vspace{2mm}
    \caption{Interface d'inscription des nouveaux employés}
    \label{fig:ui_register_s1}
\end{figure}

\subsubsection*{3. Interface de Connexion (Login)}

Cette interface \ref{fig:ui_login_s1} centralise l'accès sécurisé pour l'ensemble des utilisateurs (Employés et RH). Elle a été conçue pour être minimaliste et efficace, incluant des mécanismes de gestion d'erreurs clairs en cas d'identifiants incorrects. C'est ici que s'opère le flux critique d'authentification : en échange des identifiants valides, le serveur délivre un \textit{Token} \texttt{JWT}. Ce jeton est alors intercepté et stocké par le \textit{frontend} pour autoriser l'accès aux zones protégées de l'application selon le rôle de l'utilisateur (RBAC). Elle intègre également le lien vers le flux de récupération de compte en cas d'oubli du mot de passe.

\begin{figure}[H]
    \centering
    \includegraphics[width=\textwidth]{images/login.png}
    \vspace{2mm}
    \caption{Interface de connexion et gestion d'accès}
    \label{fig:ui_login_s1}
\end{figure}
\subsubsection*{4. Le Flux de Récupération de Compte (Self-Service Password Reset)}

Afin de garantir une autonomie totale des utilisateurs et une sécurité optimale, nous avons implémenté un flux de récupération de mot de passe complet et sécurisé. Ce mécanisme repose sur trois étapes clés :

\begin{itemize}
    \item \textbf{A. Demande de Récupération :} L'utilisateur saisit son adresse email sur une interface dédiée. Le système vérifie l'existence du compte et génère un jeton de réinitialisation unique (\textit{Token} \texttt{UUID}) à durée de vie limitée (7 heures). Cette étape prévient les tentatives d'usurpation grâce à un canal de communication chiffré.
    
    \begin{figure}[H]
        \centering
        \includegraphics[width=\textwidth]{images/forgetpassword.png}
        \caption{Interface de demande de récupération de compte}
        \label{fig:ui_forget_password}
    \end{figure}

    \item  \textbf{B. Email Transactionnel :} Un email professionnel est instantanément envoyé à l'utilisateur via un service \texttt{SMTP}. Cet email contient un lien sécurisé pointant vers l'application \textit{frontend}. Comme l'illustre la capture de la boîte de réception , le message fournit les instructions nécessaires pour poursuivre la démarche en toute confiance.

    \begin{figure}[H]
        \centering
        \includegraphics[width=\textwidth]{images/mail.png}
        \caption{Exemple d'email transactionnel reçu par l'utilisateur}
        \label{fig:ui_mail_reset}
    \end{figure}

    \item \textbf{C. Réinitialisation du Mot de Passe :} En cliquant sur le lien, l'utilisateur est redirigé vers une interface "Etape Finale" . Cette page valide le jeton présent dans l'URL et permet de définir un nouveau mot de passe robuste. L'interface bénéficie d'un design moderne et cohérent avec l'identité visuelle de TalentPredict.

    \begin{figure}[H]
        \centering
        \includegraphics[width=\textwidth]{images/modifierpassword.png}
        \caption{Interface de définition du nouveau mot de passe}
        \label{fig:ui_modify_password}
    \end{figure}
\end{itemize}
\subsubsection*{5. Contrôle d'Accès Dynamique et Navigation Différenciée}

L'une des réussites techniques majeures de ce sprint réside dans l'implémentation d'une navigation adaptative pilotée par le contexte d'authentification. Grâce à l'extraction des \textit{claims} (revendications) du jeton JWT (JSON Web Token), l'interface utilisateur (UI) se reconfigure dynamiquement en temps réel pour ne présenter que les fonctionnalités autorisées selon le rôle de l'utilisateur.
\begin{figure}[H]
    \centering
    \begin{minipage}{0.2\textwidth}
        \centering
        \includegraphics[width=\textwidth]{images/navbarUser.png} % Remplacez par le nom de l'Image 6
    \end{minipage}\hfill
    \begin{minipage}{0.2\textwidth}
        \centering
        \includegraphics[width=\textwidth]{images/navbarAdmin.png} % Remplacez par le nom de l'Image 7
    \end{minipage}
    \vspace{1mm}
    \caption{Interface de navigation côté administrateur et côté utilisateur}
    \label{fig:avigation côté administrateur et côté utilisateur}
\end{figure}

\subsubsection*{6.Gestion du Profil et Préférences de Sécurité}

Ces deux interfaces représentent le centre de contrôle de l'utilisateur au sein de la plateforme TalentPredict. Elles allient la collecte de données métier à la protection de la vie privée.

\subsubsection*{A. Interface "Mon Profil"}
Cette interface \ref{fig:ui_profile_user} est conçue comme un tableau de bord personnel permettant à l'employé de valoriser ses compétences. Elle centralise les vecteurs de données critiques qui seront analysés par le moteur d'IA lors du Sprint 2 :
\begin{figure}[H]
    \centering
    \includegraphics[width=0.7\textwidth]{images/profileUser.png}
    \caption{Interface de gestion du profil employé }
    \label{fig:ui_profile_user}
\end{figure}

\begin{itemize}
    \item \textbf{Données professionnelles :} Nombre de dépôts \texttt{GitHub}, lien vers le profil \texttt{LinkedIn} ,ien vers le profil \texttt{Portfolio} et téléversement du \texttt{CV} numérique.
    \item \textbf{Indicateurs d'engagement :} L'intégration d'une jauge de complétion et de conseils personnalisés sur les champs manquants garantit la qualité des données collectées pour le système.
\end{itemize}
\subsubsection*{B. Interface "Sécurité \& Confidentialité"}
Cette interface \ref{fig:ui_security_user} assure la confiance de l'utilisateur envers la plateforme en respectant les standards de \textit{Security by Design}. Elle offre une transparence totale sur l'utilisation du compte :
\begin{figure}[H]
    \centering
    \includegraphics[width=0.7\textwidth]{images/sécuritéUser.png}
    \caption{Interface de sécurité, historique de connexion et paramètres RGPD}
    \label{fig:ui_security_user}
\end{figure}

\begin{itemize}
    \item \textbf{Gestion des accès :} Fonctionnalités de modification sécurisée du mot de passe.
    \item \textbf{Traçabilité :} Historique de connexion détaillé incluant la date, l'adresse IP et le type d'appareil utilisé.
    \item \textbf{Conformité :} Un volet dédié au \textbf{RGPD} permet à l'utilisateur de gérer son consentement et ses données personnelles.
\end{itemize}


\subsubsection{7. Gestion du Profil et Sécurité Administrateur}

Dans l'optique de garantir une gestion rigoureuse des comptes à hauts privilèges, la plateforme \textit{TalentPredict} intègre des interfaces de gestion de profil spécifiquement configurées pour les utilisateurs possédant le rôle \texttt{ROLE\_ADMIN}.

\subsubsection{A. Profil Professionnel de l'Administrateur}
Cette interface \ref{fig:profil_admin} permet au gestionnaire RH de maintenir ses informations professionnelles et son expertise. Bien que la cohérence visuelle soit maintenue avec l'interface employé, cet espace est strictement isolé pour permettre une identification claire des responsables au sein de l'organisation :
\begin{figure}[H]
    \centering
    \includegraphics[width=0.7\textwidth]{images/profileAdmin.png}
    \caption{Interface de gestion du profil administrateur}
    \label{fig:profil_admin}
\end{figure}
\begin{itemize}
    \item \textbf{Informations Générales :} Mise à jour du titre professionnel, des années d'expérience et du niveau d'études.
    \item \textbf{Identité Numérique :} Intégration de liens vers les réseaux professionnels (LinkedIn, GitHub) et téléchargement de documents (CV, Photo de profil).
    \item \textbf{Bio \& Parcours :} Espace libre pour décrire la vision du recrutement et le parcours de l'administrateur.
\end{itemize}
\Needspace{18\baselineskip}

\subsubsection{B. Sécurité du Compte Privilégié}
Compte tenu des droits étendus de l'administrateur sur les données sensibles du personnel, l'interface de sécurité \ref{fig:security_admin} impose des contrôles rigoureux :
\begin{figure}[H]
    \centering
    \includegraphics[width=0.7\textwidth]{images/modifermdpAdmin.png}
    \caption{Interface de modification de mot de passe administrateur}
    \label{fig:security_admin}
\end{figure}
\begin{itemize}
    \item \textbf{Renouvellement des Identifiants :} Formulaire dédié imposant la saisie du mot de passe actuel pour toute modification.
    \item \textbf{Standards de Sécurité :} Le système exige un nouveau mot de passe robuste (minimum 8 caractères) pour prévenir les tentatives d'usurpation d'identité.
    \item \textbf{Confidentialité :} Accès rapide aux paramètres de visibilité des informations du compte.
\end{itemize}

\Needspace{18\baselineskip}

\subsubsection*{7. Processus de Déconnexion Sécurisée}

La clôture d'une session sur la plateforme \textit{TalentPredict} a été conçue pour allier confort d'utilisation et sécurité des données. Ce processus marque la fin du cycle d'activité de l'utilisateur et assure l'invalidité des accès après son départ.

\begin{figure}[H]
    \centering
    \includegraphics[width=0.8\textwidth]{images/logout.png}
    \caption{Interface de confirmation de déconnexion }
    \label{fig:logout_modal}
\end{figure}

\subsection{Phase de Test et Déploiement}

Cette dernière phase du Sprint 1 a pour objectif de s'assurer que les fonctionnalités développées répondent strictement aux exigences fixées, qu'elles sont sécurisées, et que le système est stable avant d'être déployé.

\subsubsection*{1. Stratégie d’Assurance Qualité (QA)}
Pour garantir la fiabilité de la plateforme TalentPredict, nous avons adopté une stratégie de test multi-niveaux :

\begin{itemize}
    \item \textbf{Validation des Règles Métier (\textit{Backend}) :} Utilisation de \texttt{JUnit 5} et \texttt{Mockito} pour tester unitairement les services critiques (notamment le hachage des mots de passe et la génération du jeton \texttt{JWT}).
    \item \textbf{Qualité du Code (Analyse Statique) :} Mise en place du \textit{plugin} Maven \texttt{Checkstyle} (via \texttt{checkstyle.xml}) pour imposer le respect des conventions de codage standard de Java.
    \item \textbf{Tests d'Intégration API :} Utilisation de \texttt{Postman} pour exécuter des requêtes sur tous les \textit{endpoints} REST, validant ainsi les codes de statut HTTP (\texttt{200 OK}, \texttt{201 Created}, \texttt{401 Unauthorized}, \texttt{403 Forbidden}, \texttt{409 Conflict}).
\end{itemize}
\subsubsection*{2. Synthèse des Tests d'Authentification Réalisés}

Pour valider le socle de sécurité du Sprint 1, une série de tests fonctionnels a été exécutée de bout en bout :

\begin{itemize}
    \item \textbf{Inscription (\texttt{/api/auth/register}) :} Validation de la création d'un utilisateur et de son profil associé (\texttt{201 Created}), du hachage sécurisé du mot de passe (\textbf{BCrypt}) et du rejet automatique des doublons d'adresse e-mail (\texttt{409 Conflict}).
    
    \item \textbf{Connexion (\texttt{/api/auth/login}) :} Validation de la vérification des identifiants, de la génération du jeton de session \textbf{JWT} signé, et du blocage des attaques par force brute via un limiteur de débit (\textbf{Bucket4j}) renvoyant un code d'erreur \texttt{429 Too Many Requests}.
    
    \item \textbf{Récupération de mot de passe :} Validation du cycle de réinitialisation sécurisée via la génération d'un jeton temporaire et à usage unique (\texttt{PasswordResetToken}), l'envoi du lien par courrier électronique, et la mise à jour effective du nouveau mot de passe haché.
    
    \item \textbf{Routes protégées et Rôles (RBAC) :} Validation de l'interception des requêtes HTTP par le filtre \texttt{JwtAuthFilter} et les Guards Angular. L'accès sans jeton valide renvoie systématiquement un code \texttt{401 Unauthorized}, tandis qu'une tentative d'accès d'un compte de rôle standard (\texttt{USER}) aux ressources réservées aux administrateurs (\texttt{/admin/**}) est bloquée avec un code de retour \texttt{403 Forbidden}.
\end{itemize}

\subsubsection*{3. Couverture de Code et Résultats des Tests}:

Afin de garantir la fiabilité des modules critiques, nous avons mesuré la couverture de code via le \textit{plugin} \texttt{JaCoCo} intégré à \texttt{Maven}. Les tests unitaires ont été rédigés avec \texttt{JUnit 5} et \texttt{Mockito}, en isolant la logique métier de la couche de persistance. Les résultats obtenus sur les modules de sécurité et d'authentification sont les suivants \ref{tab:coverage_s1} :

\begin{table}[H]
    \centering
    \renewcommand{\arraystretch}{1.5}
    \small
    \begin{tabular}{|l|c|}
        \hline
        \rowcolor{TableHeaderBlue}
        \color{white}\textbf{Module} & \color{white}\textbf{Couverture} \\
        \hline
        \texttt{AuthService} & $92\%$ \\
        \hline
        \texttt{JwtUtil} & $95\%$ \\
        \hline
        \texttt{ProfileService} & $82\%$ \\
        \hline
        \texttt{GlobalExceptionHandler} & $88\%$ \\
        \hline
        \rowcolor[HTML]{EAEAEA}
        \textbf{Total modules critiques} & \textbf{85\%} \\
        \hline
    \end{tabular}
    \vspace{2mm}
    \caption{Résultats de la couverture de code par module (Sprint 1)}
    \label{tab:coverage_s1}
\end{table}

Ces résultats confirment la robustesse du socle de sécurité livré lors de ce premier sprint.

\subsubsection*{4. Plan de Validation Fonctionnelle (Tests \textit{End-to-End})}
Le tableau \ref{tab:validation_sprint1} suivant résume les principaux scénarios de tests exécutés pour valider les flux de bout en bout de l'application :

\begin{table}[H]
    \centering
    \renewcommand{\arraystretch}{1.5}
    \small
    \begin{tabularx}{\textwidth}{|c|>{\raggedright\arraybackslash}p{3cm}|X|X|}
        \hline
        \rowcolor{TableHeaderBlue}
        \color{white}\textbf{ID} & \color{white}\textbf{Scénario Testé} & \color{white}\textbf{Pré-conditions \& Étapes} & \color{white}\textbf{Résultat Attendu} \\
        \hline
        TC-01 & Inscription d'un nouvel Employé & \textbf{Pré-cond. :} Utilisateur sur la \textit{Landing Page}. \newline \textbf{Étapes :} Saisir nom, email, mot de passe $\rightarrow$ Soumettre. & Redirection vers la page de \textit{login}. Code HTTP \texttt{201}. Création d'un \texttt{User} et d'un \texttt{Profile} vide en base. \\
        \hline
        TC-02 & Gestion des doublons (Inscription) & \textbf{Pré-cond. :} L'email "\texttt{test@tp.com}" existe déjà. \newline \textbf{Étapes :} Tenter une inscription avec cet email. & Le système bloque l'inscription. Affichage d'un message d'erreur. Code HTTP \texttt{409 Conflict}. \\
        \hline
        TC-03 & Connexion \& Contrôle d'Accès (RBAC) & \textbf{Pré-cond. :} Posséder un compte Employé valide. \newline \textbf{Étapes :} Se connecter $\rightarrow$ Tenter d'accéder à l'URL \texttt{/admin/dashboard}. & La connexion réussit (\textit{Token} \texttt{JWT} généré), mais l'accès au Panel RH est bloqué (Erreur \texttt{403 Forbidden}). \\
        \hline
        TC-04 & Enrichissement du profil (Vecteurs IA) & \textbf{Pré-cond. :} L'employé est connecté (\textit{Token} actif). \newline \textbf{Étapes :} Accéder au profil $\rightarrow$ Saisir lien \texttt{GitHub} $\rightarrow$ Enregistrer. & Mise à jour réussie en base de données. Message de succès côté interface (Code HTTP \texttt{200}). \\
        \hline
    \end{tabularx}
    \caption{Plan de validation fonctionnelle du Sprint 1}
    \label{tab:validation_sprint1}
\end{table}

\subsubsection*{5. Résultats obtenus et Corrections apportées (Débogage)} Au cours de l'exécution des tests, plusieurs défis techniques ont été identifiés et résolus, renforçant ainsi la robustesse du système :

\begin{itemize}
    \item \textbf{Erreur de Politique CORS (\textit{Cross-Origin Resource Sharing}) :}
    \begin{itemize}
        \item \textit{Problème :} Le navigateur bloquait les requêtes d'\texttt{Angular} (port 4200) vers \texttt{Spring Boot} (port 8081).
        \item \textit{Correction :} Implémentation d'une classe de configuration \texttt{WebMvcConfigurer} dans \texttt{Spring Boot} pour autoriser explicitement les requêtes, les en-têtes (notamment \texttt{Authorization}), et les méthodes HTTP provenant de l'interface \textit{front-end}.
    \end{itemize}
    \vspace{2mm}
    \item \textbf{Perte de contexte de sécurité (Filtre \texttt{JWT}) :}
    \begin{itemize}
        \item \textit{Problème :} Les rôles de l'utilisateur n'étaient pas reconnus après la connexion.
        \item \textit{Correction :} Ajout de l'extraction des \textit{Claims} (Rôles) directement depuis le \textit{Payload} du \textit{token} \texttt{JWT} dans la classe \texttt{JwtAuthFilter} et injection dans le \texttt{SecurityContextHolder} de \texttt{Spring Security}.
    \end{itemize}
    \vspace{2mm}
    \item \textbf{Problème de communication inter-conteneurs :}
    \begin{itemize}
        \item \textit{Problème :} Le \textit{backend} n'arrivait pas à se connecter à \texttt{PostgreSQL} lors du lancement \texttt{Docker}.
        \item \textit{Correction :} Ajustement de la configuration Docker pour pointer vers le nom du service (\texttt{postgres:5432}) au lieu de \texttt{localhost}.
    \end{itemize}
\end{itemize}

\subsubsection*{6. Renforcement de la Sécurité Applicative}
Afin de renforcer la sécurité de la plateforme \textbf{TalentPredict} face aux vulnérabilités web courantes, plusieurs mécanismes complémentaires ont été mis en place dès le Sprint 1 selon une approche \textit{Security by Design}.

\begin{itemize}
\item \textbf{Protection contre les injections SQL :} L’accès à la base de données repose sur \texttt{Spring Data JPA} et \texttt{Hibernate}, exploitant des requêtes paramétrées via ORM afin d’éviter toute exécution de requêtes SQL dynamiques non sécurisées.
\item \textbf{Validation des entrées utilisateur :} Les données sont validées à deux niveaux. Côté frontend, les \texttt{Reactive Forms} d’Angular assurent une validation immédiate des champs. Côté backend, les annotations de \textit{Bean Validation} (\texttt{@Valid}, \texttt{@NotNull}, \texttt{@Email}, \texttt{@Size}) garantissent l’intégrité des données reçues par les API REST.

\item \textbf{Contrôle d’accès :} Les routes sensibles sont protégées via \texttt{Spring Security} et un filtrage basé sur \texttt{JWT}, assurant l’authentification des utilisateurs et la vérification de leurs rôles avant l’accès aux ressources.

\item \textbf{Gestion des CORS :} Une politique CORS restrictive a été configurée afin de limiter les requêtes au seul frontend Angular officiel, réduisant les risques d’accès inter-domaines non autorisés.

\item \textbf{Hachage des mots de passe :} Les mots de passe sont stockés sous forme hachée à l’aide de \texttt{BCrypt} avec salage automatique, afin de se prémunir contre les attaques par dictionnaire et les \textit{rainbow tables}.

\item \textbf{Sécurisation des communications :} Les échanges entre le frontend, le backend et les services externes sont protégés via le protocole HTTPS/TLS, garantissant la confidentialité et l’intégrité des données.

\item \textbf{Journalisation :} Les opérations sensibles (authentifications, accès refusés, actions critiques) sont journalisées afin d’assurer la traçabilité des événements de sécurité.

\item \textbf{Protection contre les attaques par force brute :} Une limitation de débit a été mise en place via \texttt{Bucket4j}, restreignant les tentatives d’authentification à 5 requêtes par minute et par adresse IP sur les endpoints \texttt{/api/auth/**}. Au-delà de ce seuil, une réponse HTTP \texttt{429 Too Many Requests} est renvoyée.
\end{itemize}
\subsubsection*{7. Processus de Déploiement et Exécution}


Pour valider l’environnement cible et garantir la répétabilité de l'installation, l’application TalentPredict a été déployée en utilisant une approche d’\textit{Infrastructure as Code} (IaC). Cette méthode permet de décrire l'intégralité de l'infrastructure logicielle à travers des fichiers de configuration versionnés. Le déploiement s’effectue de manière automatisée via l'orchestrateur \texttt{Docker Compose}, selon le cycle suivant :

\begin{itemize}
    \item \textbf{Phase de Build (\textit{Backend} \& \textit{Frontend}) :} Compilation du code source Java via \texttt{Maven} (\texttt{mvn clean package}) et construction des \textit{bundles} \texttt{Angular}. Ces artefacts sont ensuite encapsulés dans des images \texttt{Docker} optimisées (utilisation de \textit{Multi-stage builds} pour réduire la taille des images).
    
    \item \textbf{Conteneurisation :} L’exécution de la commande \texttt{docker-compose up -d --build} automatise la création des conteneurs, l'allocation des ports réseau et le montage des volumes de persistance.
    
    \item \textbf{Orchestration et Dépendances :} \texttt{Docker} gère l'ordonnancement du démarrage. Le service \texttt{talentpredict-db} (\texttt{PostgreSQL 16}) est lancé en priorité. Grâce à une directive de \textit{Healthcheck}, le service \texttt{talentpredict-backend} (\texttt{Spring Boot 3.4.2}) attend que la base de données soit opérationnelle avant d'initier sa propre connexion. Enfin, le serveur \texttt{Nginx} servant l'interface \texttt{Angular} est déployé pour exposer l'application à l'utilisateur final.
\end{itemize}

\begin{figure}[H]
    \centering
    % Remplacez 'chemin_architecture_deploiement.png' par le nom exact de votre image
    \includegraphics[width=0.9\textwidth]{images/Deployment Architecture.png}
    \vspace{2mm}
    \caption{Architecture de déploiement Docker du Sprint 1}
    \label{fig:deploiement_docker_s1}
\end{figure}


Le schéma ci-dessus \ref{fig:deploiement_docker_s1} met en évidence la robustesse de l'infrastructure mise en place :

\begin{itemize}
    \item \textbf{Isolation Réseau :} L'utilisation d'un réseau privé \texttt{Docker} (\texttt{talentpredict-net}) assure qu'aucune brique sensible (\textit{Backend}, Base de données) n'est exposée directement sur Internet. Seul le serveur \texttt{Nginx} (Port 4200/443) est ouvert aux requêtes des utilisateurs.
    
    \item \textbf{Gestion de la Persistance :} Le volume \texttt{pg\_data} garantit la survie des données. En cas de mise à jour du conteneur \texttt{PostgreSQL}, les données restent intactes sur l'hôte, assurant une continuité de service.
    
    \item \textbf{Flux de Communication :} Le diagramme illustre le flux critique allant de la requête utilisateur (\textit{Browser}) vers les services métiers (\texttt{Spring Boot}), puis vers la persistance (\texttt{JDBC}). L'intégration d'un service \texttt{SMTP} externe complète ce dispositif pour les notifications critiques.
\end{itemize}


\section{Suivi et Pilotage du Sprint 1}
Le suivi et le pilotage du Sprint 1 s'inscrivent dans une démarche de gestion agile rigoureuse, visant à assurer la transparence de l'avancement et le respect des objectifs fixés. À travers l'analyse des indicateurs de vélocité et l'actualisation du backlog, ce contrôle continu permet de garantir l'alignement parfait entre les développements techniques et les exigences de la plateforme.
\subsection{Scrum Board}

Afin de garantir la transparence et l’inspection continue, essentielles dans la méthodologie \textit{Scrum}, nous avons utilisé des outils visuels pour gérer cette itération. Le sprint a duré trois semaines (soit 15 jours ouvrés), avec une charge de travail estimée à 38 Story Points.

Le suivi quotidien des tâches s’est matérialisé par un tableau de bord (\textit{Scrum Board}) reflétant l’état d’avancement des \textit{User Stories}. Ce management visuel a permis à l’équipe de développement d’identifier rapidement les goulots d’étranglement et de fluidifier le passage des tâches de la phase d'implémentation \textit{backend} vers la phase d'intégration et de test \textit{frontend}.

\begin{figure}[H]
    \centering
    % Remplacez 'scrum_board_sprint1' par le nom de votre fichier image
    \includegraphics[width=1\textwidth]{images/scrumboardSprint1.png}
    \vspace{2mm}
    \caption{État final du Scrum Board - Sprint 1}
    \label{fig:scrum_board_s1}
\end{figure}

\subsection{Le Burndown Chart}

Pour évaluer notre vélocité et s’assurer de l’atteinte de l’objectif dans le temps imparti, nous avons maintenu à jour le \textit{Burndown Chart} lors de nos \textit{Daily Stand-ups}.

\begin{figure}[H]
    \centering
    \begin{tikzpicture}
        \begin{axis}[
            width=16cm,
            height=8cm,
            title={\textbf{Burndown Chart - Sprint 1 }},
            xlabel={Jours ouvrés du Sprint (3 Semaines)},
            ylabel={Story Points restants},
            xmin=0, xmax=15,
            ymin=0, ymax=40,
            xtick={0,1,2,3,4,5,6,7,8,9,10,11,12,13,14,15},
            ytick={0,10,20,30,38,40},
            legend pos=north east,
            ymajorgrids=true,
            grid style=dashed,
            smooth
        ]
        
        % Ligne idéale (Tendance Idéale)
        \addplot[
            color=blue,
            dashed,
            mark=none,
            very thick
            ]
            coordinates {
            (0,38)(15,0)
            };
        \addlegendentry{Tendance Idéale}
        
        % Ligne réelle (Avancement Réel avec stagnation J7-J9)
        \addplot[
            color=red,
            mark=square*,
            very thick
            ]
            coordinates {
            (0,38)(1,36)(2,33)(3,30)(4,28)(5,25)(6,22)(7,22)(8,22)(9,20)(10,15)(11,10)(12,5)(13,2)(14,0)(15,0)
            };
        \addlegendentry{Avancement Réel}
        
        \end{axis}
    \end{tikzpicture}
    \caption{Burndown Chart du Sprint 1}
    \label{fig:burndown_sprint1}
\end{figure}

Le graphique montre la courbe de l’effort restant en \textit{Story Points} au fil des 15 jours du sprint. Nous avons rencontré quelques difficultés techniques notables :
\begin{itemize}
    \item \textbf{Stagnation (J7 à J9) :} Observable sur la courbe, cette période s'explique par les défis techniques liés à la configuration du filtre \texttt{JWT} dans \texttt{Spring Security} et la résolution des conflits de politique \texttt{CORS} entre l'interface \texttt{Angular} et le \textit{backend}.
    \item \textbf{Accélération finale :} La résolution de ces obstacles architecturaux a permis une forte accélération de la validation finale, clôturant ainsi le sprint dans les délais impartis.
\end{itemize}
\subsection{Indicateurs de Performance (KPIs) du Sprint 1}

Pour évaluer l'efficacité de cette première itération, nous avons défini un ensemble d'indicateurs de performance répartis selon deux axes : la gestion de projet Agile et la qualité technique.

\paragraph*{1. Indicateurs de Pilotage (Agilité)}
\begin{itemize}
    \item \textbf{Vélocité de l'Équipe :} L'objectif était de réaliser 38 \textit{Story Points}. Au terme du sprint, l'intégralité des points a été validée, soit un taux de complétion de 100\%.
    \item \textbf{Respect des Délais :} Le sprint a été clôturé en 15 jours, conformément au planning initial, malgré les défis techniques liés à l'intégration du module \texttt{Spring Security}.
    \item \textbf{Qualité du \textit{Backlog} :} Aucun \textit{Impediment} (blocage) majeur n'a nécessité le report d'une \textit{User Story} au sprint suivant.
\end{itemize}

\paragraph*{2. Indicateurs Techniques et Qualité}
Le tableau suivant résume les métriques de qualité logicielle obtenues :

\begin{table}[H]
    \centering
    \renewcommand{\arraystretch}{1.5}
    \small
    \begin{tabularx}{0.9\textwidth}{|X|c|c|c|}
        \hline
        \rowcolor{TableHeaderBlue}
        \color{white}\textbf{Indicateur} & \color{white}\textbf{Cible} & \color{white}\textbf{Résultat Obtenu} & \color{white}\textbf{État} \\
        \hline
        Couverture de code (Tests) & $> 80\%$ & $85\%$ & $\checkmark$ Validé \\
        \hline
        Temps de réponse moyen API & $< 200$ ms & $120$ ms & $\checkmark$ Validé \\
        \hline
        Failles de sécurité (\texttt{JWT}) & $0$ & $0$ & $\checkmark$ Validé \\
        \hline
        Disponibilité (\textit{Uptime} \texttt{Docker}) & $100\%$ & $100\%$ & $\checkmark$ Validé \\
        \hline
        Poids des images \texttt{Docker} & $< 500$ Mo & $320$ Mo & $\checkmark$ Validé \\
        \hline
    \end{tabularx}
    \caption{Tableau de bord des métriques techniques du Sprint 1}
    \label{tab:kpi_tech_s1}
\end{table}

\paragraph*{3. Analyse du Bilan}
Le bilan du Sprint 1 est extrêmement positif. La réussite des tests d'intégration et la stabilité de l'infrastructure \texttt{Docker} confirment que le socle de TalentPredict est sain. L'indicateur le plus significatif est la vélocité constante, qui démontre une bonne maîtrise de la \textit{stack} technologique (\texttt{Angular} / \texttt{Spring Boot}) par l'équipe de développement.
\section{Conclusion}
\addcontentsline{toc}{section}{Conclusion du Sprint 1}
Le premier \textit{sprint} du projet TalentPredict s'achève sur un bilan très positif. Cette itération, consacrée aux fondations techniques et à la gestion de l'identité, a permis de mettre en place les bases essentielles au bon fonctionnement de la plateforme. 

En appliquant rigoureusement la méthodologie \textit{Scrum}, nous sommes passés d'un concept théorique à une application fonctionnelle, sécurisée et conteneurisée. Les principaux objectifs ont été atteints avec succès : l'architecture a été stabilisée avec la mise en place d'une structure N-tiers robuste basée sur \texttt{Spring Boot} et \texttt{Angular} ; la sécurité a été assurée grâce à l'implémentation de l'authentification et de l'autorisation via des jetons \texttt{JWT} ; et la modélisation des données a permis de représenter les principaux éléments du profil professionnel, incluant le CV, \texttt{GitHub}, \texttt{LinkedIn} et le \textit{portfolio}.
Ce \textit{sprint} a non seulement validé la solidité du socle technique, mais a également démontré notre capacité à surmonter les défis d'implémentation, notamment à travers des tests d'\texttt{API} systématiques avec \texttt{Postman} et l'utilisation d'un environnement \texttt{DevOps} basé sur \texttt{Docker}. 
Cependant, un profil utilisateur reste une donnée statique tant qu'il n'est pas interprété par une couche d'intelligence. C'est pourquoi le Sprint 2 constitue une étape clé du projet, avec l'introduction du Moteur d'Évaluation Multidimensionnelle, visant à passer de la simple collecte de données à la génération de valeur, grâce à l'intégration d'une analyse automatique des compétences basée sur l'IA et d'une évaluation psychométrique (PCM), ouvrant ainsi la voie vers le cœur prédictif de TalentPredict.



\chapter{Sprint 2  Évaluation Multi-Source des Compétences}

\section{Introduction}
\phantomsection
\addcontentsline{toc}{section}{Introduction}

Ce troisième chapitre documente la réalisation du Sprint 2, dont l'objectif est de livrer le module complet d'évaluation des compétences techniques et comportementales de TalentPredict. Contrairement aux fondations posées lors du Sprint 1 (sécurité, authentification, architecture de base), ce sprint se concentre sur le cœur métier de la plateforme : l'analyse intelligente multi-source des talents.
Le module se décompose en deux sous-systèmes complémentaires développés en parallèle par les deux membres du binôme : le module Compétences comportementales  et le module Compétences techniques . Ces deux modules partagent une architecture commune mais reposent sur des pipelines technologiques distincts, reflétant la nature différente des données analysées.
La structure de ce chapitre suit la méthodologie Agile Scrum adoptée par l'équipe. Nous présentons successivement : l'objectif du sprint et ses livrables attendus , la séance de planification et les décisions d'architecture , le backlog estimé , puis le déroulement technique du sprint organisé selon les phases d'analyse, conception, codage et validation . 



\section{Objectifs du Sprint 2}

Avant de décrire le déroulement du sprint, il est essentiel de définir précisément le périmètre fonctionnel visé et les critères de validation qui guideront l’équipe tout au long des trois semaines de développement.

L’objectif majeur de ce sprint est de livrer un moteur d’évaluation multi-sources complet et autonome. L'enjeu est de faire évoluer la plateforme TalentPredict : d'une simple base de données déclarative, elle devient un véritable système de diagnostic intelligent capable de profiler un collaborateur sous toutes ses coutures en agrégeant de multiples sources d’analyse complémentaires pour produire une cartographie unifiée (compétences techniques et comportementales).

\paragraph*{1. Objectifs fonctionnels et interactifs :}
\begin{itemize}[leftmargin=1cm]
    \item \textbf{Évaluation comportementale :} Déploiement d'un workflow automatisé intégrant le test heuristique PCM, l'analyse sémantique du CV et de GitHub, ainsi que la résolution de scénarios immersifs générés par l'IA.
    \item \textbf{Évaluation technique :} Mise en œuvre d'un parcours de validation factuel combinant l'extraction de la stack technique (CV, LinkedIn), l'analyse de l'historique GitHub, et des tests interactifs (QCM adaptatif et challenge de code).
    \item \textbf{Job Matching :} Intégration d'un algorithme de correspondance permettant de calculer un score de compatibilité entre le profil consolidé du collaborateur et plus de 20 profils métiers standards.
\end{itemize}

\paragraph*{2. Objectifs techniques et opérationnels :}
\begin{itemize}[leftmargin=1cm]
    \item \textbf{Orchestration asynchrone :} Configuration et déploiement de workflows n8n (Master Agent, CV Parser, GitHub Analysis) pour piloter les flux de données lourds sans surcharger le backend central.
    \item \textbf{Inférence IA locale :} Implémentation du microservice Python (FastAPI) communiquant avec l'écosystème local Ollama (Llama 3 / CodeLlama) pour la génération, la correction et l'évaluation des scénarios et des défis algorithmiques.
    \item \textbf{Persistance unifiée :} Conception et migration du schéma de base de données PostgreSQL via Spring Boot pour stocker les entités Skill (en distinguant les énumérations TypeSkill.TECH et TypeSkill.SOFT).
\end{itemize}

\paragraph*{Critères d’acceptation (Definition of Done) :}
La validation du pipeline d'évaluation repose sur les critères stricts suivants :
\begin{itemize}[leftmargin=1cm]
    \item \textbf{Cohérence multi-sources :} Le  Master Agent doit agréger avec succès les scores comportementaux (0--100) issus de GitHub, du CV, du PCM et du scénario immersif pour générer un rapport unifié cohérent.
    \item \textbf{Exactitude de l'extraction technique :} Le \textit{parser} doit être capable d'identifier avec précision un catalogue de plus de 150 technologies avec leurs années d'expérience associées à partir d'un fichier brut.
    \item \textbf{Évaluation algorithmique :} Le challenge de code généré par CodeLlama doit être évalué de manière autonome selon 4 critères précis (syntaxe, logique, complexité, sécurité) pour en déduire un niveau (Beginner à Expert).
    \item \textbf{Intégrité des données :} Chaque nouvelle compétence détectée et validée par le pipeline doit être persistée instantanément et correctement catégorisée en base de données.
\end{itemize}

\vspace{0.5cm}
Le Sprint 2 vise à livrer treize composants techniques répartis sur les deux modules parallèles d'évaluation. Le tableau suivant formalise ces livrables avec leurs critères de validation respectifs :

\begin{table}[H]
\centering
\small
\renewcommand{\arraystretch}{1.3}
\begin{tabularx}{\textwidth}{|c|>{\bfseries}P{0.35\textwidth}|Y|}
\hline
\rowcolor{TableHeaderBlue}
\color{white}\textbf{\#} & \color{white}\textbf{Livrable technique} & \color{white}\textbf{Critère de validation détaillé} \\
\hline
1 & Analyse GitHub comportementale & Génération de 4 scores comportementaux (0--100). \\
\hline
2 & Analyse CV comportementale & Calcul du score CV et extraction des indices sémantiques. \\
\hline
3 & Test heuristique PCM & Identification des 6 dimensions cognitives et du type dominant via 18 questions. \\
\hline
4 & Synthèse Master Agent & Agrégation réussie produisant un rapport pondéré unifié. \\
\hline
5 & Scénario IA comportemental & Évaluation sémantique produisant 4 scores psychologiques. \\
\hline
6 & Analyse GitHub technique & Extraction avérée des langages, frameworks et du niveau d'engagement. \\
\hline
7 & Parsing CV technique & Identification précise sur une base de 150+ technologies et années d’expérience. \\
\hline
8 & Analyse LinkedIn & Établissement du positionnement professionnel et score /10. \\
\hline
9 & QCM technique adaptatif & Génération dynamique des questions et scoring pondéré validé. \\
\hline
10 & Challenge de code (CodeLlama) & Génération et évaluation algorithmique fonctionnelles (basées sur 4 critères). \\
\hline
11 & Scoring multi-sources & Classification de la maîtrise technique de Beginner à Expert. \\
\hline
12 & Moteur de Job matching & Calcul du score de compatibilité (/100) face à 20 profils métiers cibles. \\
\hline
13 & Persistance des compétences & Stockage relationnel réussi de l'entité \texttt{Skill} en base de données PostgreSQL. \\
\hline
\end{tabularx}
\renewcommand{\arraystretch}{1.0}
\caption{Livrables attendus et validation du Sprint 2}
\end{table}
La figure~\ref{fig:pip-glob-sprint2} illustre le pipeline global d'analyse implémenté durant cette phase. Elle retrace le flux complet des données, depuis l'ingestion des sources hétérogènes (CV, API GitHub, résultats de QCM, challenges de code et analyse comportementale PCM) jusqu'à leur agrégation en un score d'évaluation unifié.
\begin{figure}[H]
  \centering
  \includegraphics[width=0.9\textwidth]{images/pipline glob sprint2.png}
  \caption{Pipeline global d’analyse multi-sources}
  \label{fig:pip-glob-sprint2}
\end{figure}

\section{Séance du Sprint Planning}
La séance de Sprint Planning s'est tenue  en présence de l'ensemble de l'équipe projet. Cette réunion structurante a permis de traduire les exigences fonctionnelles en choix techniques concrets et de définir les modalités de collaboration entre les deux développeurs. Les décisions prises durant cette séance conditionnent l'architecture finale du module
Participants : Barouni Rahma, Abdeladhim Mohamed  , M. Mohamed Amine Gharbi (Product Owner/ encadrant professionnel et Takwa ben Aicha ( encadrente Academique )).

\textbf{Durée du sprint :} quatre semaines, correspondant aux semaines 5 à 8 du diagramme de Gantt du projet.

 

\textbf{Éléments discutés et décisions clés}

\textbf{Module compétences comportementales  :}
\begin{itemize}[leftmargin=1cm]
  \item n8n est retenu comme orchestrateur des workflows d'analyse GitHub, CV et PCM afin de dissocier la logique d'orchestration du backend applicatif et de permettre la modification des prompts sans recompilation.
  \item Ollama llama3.2:latest est choisi comme moteur LLM local pour l'ensemble du module compétences comportementales, ce qui favorise une conformité RGPD par conception.
  \item Le scoring PCM est implémenté en JavaScript pur, sans appel LLM, afin de garantir la reproductibilité et de réduire la latence d'exécution.
  \item Une séparation claire des responsabilités est établie : n8n gère l'analyse de profil statique, tandis que FastAPI prend en charge l'évaluation comportementale en temps réel via les scénarios IA.
\end{itemize}

\textbf{Module compétences techniques  :}
\begin{itemize}[leftmargin=1cm]
  \item Une architecture  :  Ollama \texttt{llama3.2} en mode local avec pipeline direct.
  \item \texttt{pdfplumber} est retenu pour l'extraction PDF côté serveur, car il offre une meilleure fiabilité que PyPDF2 sur les documents PDF vectoriels.
  \item \texttt{CodeLlama} est utilisé spécifiquement pour les challenges de code, grâce à de meilleures performances sur la génération et l'évaluation de code que llama3.2.
  \item Le service TalentPredictAiProxyService de Spring Boot assure le pont entre l'API Java et le microservice Python FastAPI avec un délai d'attente configurable de 120 secondes.
\end{itemize}

\textbf{Priorités retenues :}
\begin{enumerate}[leftmargin=1cm]
  \item Workflows n8n : GitHub Analysis, CV Parser, PCM Test et Master Agent.
  \item Scénario comportemental IA dans FastAPI.
  \item CV Parser et GitHub Analyzer pour les compétences techniques.
  \item QCM adaptatif et challenge de code avec FastAPI et CodeLlama.
  \item Persistance des compétences via Spring Boot et PostgreSQL.

\end{enumerate}

\textbf{Hypothèses et contraintes :}
\begin{itemize}[leftmargin=1cm]
  \item \textbf{Latence Ollama :} les délais d'attente sont configurés avec des fallbacks heuristiques afin de garantir une réponse exploitable même en cas de lenteur du modèle local.
  \item \textbf{Format JSON instable :} les réponses produites par Ollama font l'objet d'un nettoyage systématique avant leur exploitation par les services applicatifs.
  \item \textbf{GitHub rate limit :} un GITHUB\_TOKEN est requis en production, car l'API GitHub est limitée à 60 requêtes par heure sans authentification.
  \item \textbf{PDF scannés :} un avertissement est retourné et une analyse partielle est réalisée lorsque le texte du document n'est pas extractible.
\end{itemize}
Ces décisions d'architecture ont été formalisées dans le Sprint Backlog présenté en section suivante.

\section{Sprint Backlog Estimé}
Sur la base des décisions prises lors du Sprint Planning, l'équipe a décomposé l'objectif global en User Stories estimées. Le backlog ci-dessous reflète la priorisation métier validée par le Product Owner et l'effort estimé en jours-homme selon la méthode du Planning Poker. La charge totale de 37 jours-homme (20j pour Compétences techniques + 17j pour Compétences comportementales) correspond à un sprint de 3 semaines avec deux développeurs à temps plein.

Le tableau \ref{tab:backlog-hard-skills} présente le backlog du module Compétences techniques.

\begin{table}[H]
\centering
\small
\setlength{\tabcolsep}{3pt}
\begin{tabular}{|P{1.0cm}|P{2.4cm}|P{5.8cm}|C{0.7cm}|P{1.5cm}|}
\hline
\rowcolor{TableHeaderBlue}
\color{white}\textbf{ID} & \color{white}\textbf{Module / US} & \color{white}\textbf{Tâches de développement} & \color{white}\textbf{JH} & \color{white}\textbf{Priorité} \\
\hline
HS-01 & Analyse GitHub technique & Récupération des repos via l'API REST ; détection des langages, frameworks et niveau d'activité. & 3 & Haute \\
\hline
HS-02 & Parsing CV PDF/DOCX & Extraction de 150+ technologies par regex via pdfplumber et python-docx ; estimation des années d'expérience. & 2 & Haute \\
\hline
HS-03 & Analyse LinkedIn & Analyse via Ollama du positionnement, arc de carrière et cohérence CV/GitHub. & 2 & Haute \\
\hline
HS-04 & Génération QCM adaptatif & Moteur de génération de 6 à 12 questions par compétence et niveau via Ollama. & 3 & Haute \\
\hline
HS-05 & Scoring QCM pondéré & Évaluation selon trois dimensions : exactitude $\times$ confiance $\times$ temps de réponse. & 2 & Haute \\
\hline
HS-06 & Challenge de code CodeLlama & Génération fix\_bugs / complete / refactor ; évaluation sur 4 critères avec pénalité par indice. & 3 & Moyenne \\
\hline
HS-07 & Scoring multi-sources & Déduplication de 500+ alias ; nivellement Beginner $\rightarrow$ Expert ; job match sur 20+ profils métiers. & 3 & Haute \\
\hline
HS-08 & Persistance compétences TECH & Entité Skill (TypeSkill.TECH) avec endpoints Spring Boot pour sauvegarde et validation ADMIN. & 2 & Haute \\
\hline
\multicolumn{3}{|r|}{\textbf{Total}} & \textbf{20} & \\
\hline
\end{tabular}
\caption{Sprint Backlog estimé du module Compétences techniques}\label{tab:backlog-hard-skills}
\end{table}

Le tableau \ref{tab:backlog-soft-skills} détaille le backlog du module Compétences comportementales.

\begin{table}[H]
\centering
\small
\setlength{\tabcolsep}{3pt}
\begin{tabular}{|P{1.0cm}|P{2.5cm}|P{5.7cm}|C{0.7cm}|P{1.5cm}|}
\hline
\rowcolor{TableHeaderBlue}
\color{white}\textbf{ID} & \color{white}\textbf{Module / US} & \color{white}\textbf{Tâches de développement} & \color{white}\textbf{JH} & \color{white}\textbf{Priorité} \\
\hline
SS-01 & Analyse GitHub comportementale & Workflow n8n : scoring heuristique JS (repos, stars, forks) ; analyse Ollama des 4 dimensions comportementales. & 3 & Haute \\
\hline
SS-02 & Extraction compétences comportementales CV & Workflow n8n CV Parser : envoi du texte brut à Ollama ; extraction des indices comportementaux et score CV global. & 2 & Haute \\
\hline
SS-03 & Test PCM (18 questions) & Workflow n8n : calcul 100\% heuristique JS des 6 dimensions PCM, sans appel LLM. & 2 & Haute \\
\hline
SS-04 & Master Agent --- Synthèse pondérée & Orchestrateur n8n : appels parallèles aux 3 sous-workflows ; synthèse Ollama CV 40\% / PCM 30\% / GitHub 30\%. & 2 & Haute \\
\hline
SS-05 & Génération de scénario IA & FastAPI + Ollama : génération d'un dilemme professionnel contextualisé par rôle et niveau. & 2 & Moyenne \\
\hline
SS-06 & Évaluation de la réponse scénario & FastAPI + Ollama : évaluation sur 4 dimensions psychologiques (empathie, assertivité, pragmatisme, clarté). & 2 & Moyenne \\
\hline
SS-07 & Entretien vocal IA multi-tours & FastAPI + Ollama : entretien de 5 tours avec pivot automatique vers la dimension la plus faible (score < 55/100). & 3 & Moyenne \\
\hline
SS-08 & Persistance compétences SOFT & Sauvegarde de l'entité \texttt{Skill} (\texttt{TypeSkill.SOFT}) depuis les résultats du Master Agent et scores PCM via Spring Boot. & 1 & Haute \\
\hline
\multicolumn{3}{|r|}{\textbf{Total}} & \textbf{17} & \\
\hline
\end{tabular}
\caption{Sprint Backlog estimé du module Compétences comportementales}\label{tab:backlog-soft-skills}
\end{table}

\section{Démarrage du Sprint 2}

Le début du Sprint 2 a été caractérisé par une phase d’initialisation technique approfondie en préparation des développements relatifs au moteur d'évaluation multi-sources. L’équipe s’est aussitôt focalisée sur l'orchestration des flux de données via n8n et la configuration de l'intelligence artificielle locale (Ollama), tout en établissant les ponts de communication avec le backend Spring Boot. Cette phase essentielle a posé les fondations nécessaires pour l’analyse croisée des compétences techniques et comportementales, garantissant ainsi un passage en douceur de la modélisation théorique à la mise en œuvre effective des pipelines d'évaluation.

\subsubsection*{1. Activités réalisées}

Durant cette itération, l’équipe a mené de front des développements backend complexes et des intégrations poussées d’intelligence artificielle :
\begin{itemize}[leftmargin=1cm]
    \item \textbf{Ingénierie des flux (n8n) :} Conception et automatisation des workflows d'extraction et d'analyse pour décharger le backend des traitements asynchrones lourds.
    \item \textbf{Développement IA (FastAPI \& Ollama) :} Implémentation du microservice Python pour le parsing sémantique des CV, la génération de scénarios immersifs et l'évaluation du code source.
    \item \textbf{Développement Backend (Spring Boot) :} Création des endpoints REST pour la collecte des données brutes (CV, GitHub) et implémentation du stockage relationnel des entités de compétences.
    \item \textbf{Développement Frontend (Angular) :} Création des interfaces de passation de tests (QCM, questionnaire PCM, IDE de code) et des formulaires de soumission de profil.
    \item \textbf{Intégration d'API tierces :} Connexion sécurisée à l'API GitHub pour récupérer l'historique d'activité et les statistiques de collaboration des employés.
\end{itemize}

\subsubsection*{2. Livrables produits}

Les livrables suivants ont été validés et intégrés à la branche principale du projet :
\begin{itemize}[leftmargin=1cm]
    \item \textbf{Module d'évaluation comportementale :} Workflow complet de détermination des 6 dimensions cognitives.
    \item \textbf{Module d'évaluation technique :} Pipeline d'analyse factuelle du niveau de maîtrise (Beginner à Expert).
    \item \textbf{Service d'Intelligence Artificielle :} API locale propulsée de manière autonome par Ollama et CodeLlama.
    \item \textbf{Moteur de Job Matching :} Algorithme de calcul de compatibilité basé sur plus de 20 profils métiers.
    \item \textbf{Interfaces Utilisateur :} Écrans interactifs de passation de tests dynamiques et adaptatifs.
\end{itemize}

\subsubsection*{3. Difficultés rencontrées et traitement}

L’exécution de ce sprint a été marquée par trois défis techniques majeurs liés à l'intelligence artificielle et au traitement des données :
\begin{itemize}[leftmargin=1cm]
    \item \textbf{Difficulté 1 : Hétérogénéité des formats de CV}
    \begin{itemize}[label=--]
        \item \textit{Problème :} L'extraction de texte brut à partir de fichiers PDF générés par différents outils entraînait des pertes de données sémantiques ou des erreurs d'encodage.
        \item \textit{Traitement :} Implémentation de la bibliothèque \texttt{pdfplumber} dans le microservice FastAPI, couplée à des expressions régulières avancées (Regex) pour nettoyer et normaliser le texte avant son ingestion par le LLM.
    \end{itemize}
    \item \textbf{Difficulté 2 : Limites de requêtes de l'API GitHub (\textit{Rate Limits})}
    \begin{itemize}[label=--]
        \item \textit{Problème :} Les requêtes multiples pour extraire l'historique détaillé des commits provoquaient des blocages temporaires imposés par GitHub.
        \item \textit{Traitement :} Optimisation des appels via n8n en ciblant uniquement les métadonnées essentielles et mise en place d'un système de mise en cache temporaire.
    \end{itemize}
    \item \textbf{Difficulté 3 : Temps d'inférence de l'IA locale}
    \begin{itemize}[label=--]
        \item \textit{Problème :} La génération asynchrone des scénarios professionnels complexes par Ollama causait des timeouts côté frontend (Erreur 503).
        \item \textit{Traitement :} Transition vers une architecture événementielle en s'appuyant sur les Server-Sent Events (SSE) et le Communication Hub pour libérer la connexion HTTP et notifier l'utilisateur dès la fin du traitement de l'IA.
    \end{itemize}
\end{itemize}

\subsection{Analyse}

La phase d’analyse vise à définir le périmètre fonctionnel du module d'évaluation et à identifier avec précision les interactions entre le système et les utilisateurs. Ce sprint introduit une forte automatisation grâce à l'Intelligence Artificielle et à l'orchestration des flux.

\subsubsection*{1. Identification des Acteurs}

Le système interagit avec plusieurs acteurs principaux, parmi lesquels se trouvent des entités humaines et des acteurs automatisés :
\begin{itemize}[leftmargin=1cm]
    \item \textbf{Utilisateur (Employé) :} Constitue le cœur du processus de collecte de données. Il soumet ses preuves documentaires (CV, identifiant GitHub) et s'engage activement dans l'évaluation en répondant au test PCM, au QCM adaptatif, en résolvant le challenge de code, puis en interagissant avec le scénario immersif IA.
    \item \textbf{Administrateur (RH) :} Supervise le bon déroulement des évaluations de son équipe. Il consulte les résultats consolidés, valide formellement les compétences détectées par la machine et s'appuie sur le système pour générer les rapports PDF d'aide à la décision.
    \item \textbf{Moteurs IA et Orchestrateurs (« Systèmes ») :} Agissent de manière autonome en arrière-plan pour traiter la complexité algorithmique :
    \begin{itemize}[label=--]
        \item \textit{Ollama \& CodeLlama} : Réalisent l’analyse sémantique textuelle, la synthèse comportementale, ainsi que la génération et la correction intelligente des scénarios et du code.
        \item \textit{Microservice FastAPI} : Exécute le \textit{parsing} des CV, pré-traite les données brutes, et assure l'interface directe avec les modèles LLM.
        \item \textit{n8n} : Ordonnance les appels asynchrones entre le backend, les API externes et les services d'IA, garantissant la bonne chronologie de l'évaluation.
    \end{itemize}
    \item \textbf{Sources externes (GitHub API) :} Agit comme un acteur de confiance tiers, fournissant des données factuelles (dépôts, commits, langages, \textit{pull requests}) pour objectiver les déclarations initiales du collaborateur.
\end{itemize}

\subsubsection{Spécification des besoins fonctionnels}

\textbf{BF-SS-01 : Analyse de l'activité GitHub (compétences comportementales)}\par
Le système calcule un score heuristique JavaScript (repos 40\%, stars 40\%, forks 20\%), puis soumet les données à Ollama pour évaluer quatre dimensions comportementales : curiosité, collaboration, ownership et consistance. Un fallback heuristique est retourné en cas d'échec, avec des timeouts de 10 s (GitHub API) et 20 s (Ollama).

\textbf{BF-SS-02 : Extraction compétences comportementales depuis CV}\par
Le texte brut extrait via pdfjs-dist est analysé par Ollama pour détecter les indices comportementaux (travail en équipe, leadership, adaptabilité), le niveau d'éducation et les années d'expérience, produisant un score global de 0 à 100. Le prompt exclut explicitement les technologies pour éviter toute contamination avec le module compétences techniques (timeout : 15 s).

\textbf{BF-SS-03 : Test de personnalité PCM}\par
Les scores des six dimensions PCM (communication, discipline, curiosité, collaboration, ownership, leadership) sont calculés algorithmiquement en JavaScript pur, chaque dimension correspondant à la moyenne de trois questions sur une échelle de 0 à 5, sans appel LLM, garantissant une reproductibilité totale en moins de 100 ms.
\begin{figure}[H]
  \centering
  \includegraphics[width=0.9\textwidth]{images/pipline pcm sprint2.png}
  \caption{Pipeline d’évaluation PCM}
  \label{fig:pip-PCM-sprint2}
\end{figure}
\textbf{BF-SS-04 : Synthèse Master Agent}\par
Le Master Agent orchestre en parallèle via Promise.all les trois sous-workflows (GitHub 30\%, CV 40\%, PCM 30\%), puis soumet les résultats agrégés à Ollama pour produire un profil unifié (score global, type de personnalité, forces, recommandations). En cas d'indisponibilité GitHub, la pondération bascule dynamiquement vers CV 60\% / PCM 40\% (timeout global : 30 s).

\textbf{BF-SS-05 : Génération et évaluation de scénarios comportementaux IA}\par
Ollama génère un dilemme professionnel parmi cinq types (deadline, conflit, feedback client, décision, intégration) contextualisé par rôle et niveau (température 0.7). La réponse libre du employé est ensuite évaluée sur quatre dimensions psychologiques : empathie, assertivité, pragmatisme et clarté, avec une température de 0.2 pour garantir la reproductibilité, produisant scores, forces, axes d'amélioration et profil culture-add.

\textbf{BF-HS-01 : Analyse GitHub compétences techniques}par
Le système récupère les dépôts publics via l'API REST GitHub, calcule les statistiques de langages et détecte les frameworks depuis les topics et descriptions, avant enrichissement par Ollama. Un  GITHUB\_TOKEN  est requis en production (limite : 60 req/h sans token ; timeout : 30 s). Les champs  language\_stats et detected\_frameworks alimentent ensuite le moteur de scoring.

\textbf{BF-HS-02 : Extraction de technologies depuis CV}\par
\texttt{pdfplumber} (PDF vectoriel) ou python-docx (DOCX) extrait le texte, puis 150+ technologies sont détectées par expressions régulières organisées par catégorie (langages, frameworks, cloud, DevOps, IA/ML, bases de données). Les années d'expérience sont estimées via patterns linguistiques. En cas de PDF scanné, une analyse partielle est poursuivie avec avertissement.
\begin{figure}[H]
  \centering
  \includegraphics[width=0.9\textwidth]{images/pipline cv sprint2.png}
  \caption{Pipeline d’analyse du CV}
  \label{fig:pip-cv-tech-sprint2}
\end{figure}
\textbf{BF-HS-03 : Analyse LinkedIn}\par
Ollama analyse le profil LinkedIn (URL ou texte) pour extraire le rôle actuel, l'arc de carrière, les certifications et la cohérence avec le CV et GitHub. Un score de complétude de 0 à 10 est attribué, permettant de détecter les incohérences entre les déclarations du employé et son activité réelle.

\textbf{BF-HS-04 : Test QCM technique adaptatif}\par
Ollama génère dynamiquement 6 à 12 questions par compétence et niveau (Junior / Senior / Lead) à température 0.4, avec fallback sur une banque prédéfinie balisée source: fallback. Le scoring combine trois dimensions : exactitude $\times$ confiance $\times$ temps, avec un multiplicateur 1.2 pour haute confiance correcte, une pénalité de surconfiance pour haute confiance incorrecte, et un bonus de 0.05 pour bonne réponse difficile en moins de 15 secondes.
\begin{figure}[H]
  \centering
  \includegraphics[width=0.9\textwidth]{images/QCM.png}
  \caption{Génération de QCM adaptatifs par IA}
  \label{fig:pip-QCM-sprint2}
\end{figure}
\textbf{BF-HS-05 : Challenge de code }\par
CodeLlama génère un défi de type fix\_bugs, complete ou refactor et évalue la soumission selon quatre critères pondérés : correctness (40 pts), code quality (30 pts), efficiency (20 pts) et readability (10 pts). Une pénalité de 10 points par indice utilisé est appliquée, soit : score final = total brut $-$ (\texttt{hints\_used} $\times$ 10).
\begin{figure}[H]
  \centering
  \includegraphics[width=0.9\textwidth]{images/code challennge.png}
  \caption{Évaluation automatiquedes challenges de code}
  \label{fig:pip-code-challenge-sprint2}
\end{figure}
\textbf{BF-HS-06 : Scoring multi-sources et job matching}\par
Les compétences issues des trois sources~\ref{fig:Fusion-multi-scores-sprint2} sont consolidées via 500+ alias normalisés dans skill\_extractor.py, puis nivelées par triangulation : Expert ($\geq$ 3 sources, score $\geq$ 90), Advanced ($\geq$ 2 sources, 70--89), Intermediate (1 source forte, 40--69), Beginner ($\sim$25). Le job matching compare le profil avec 20+ profils métiers selon la formule :
\[
  \texttt{job\_match\_score} = \left(\frac{\texttt{earned\_weight}}{\texttt{total\_weight}}\right) \times 85 + \texttt{bonus\_extra\_skills},
\]
ajustée par l'\texttt{experience\_score}.
\begin{figure}[H]
  \centering
  \includegraphics[width=0.9\textwidth]{images/Fusion multi sources.png}
  \caption{Fusion multi-sources des
scores}
  \label{fig:Fusion-multi-scores-sprint2}
\end{figure}

\subsubsection{Objectifs non fonctionnels}

\begin{table}[H]
\centering
\renewcommand{\arraystretch}{1.3}
\begin{tabularx}{\TableWidth}{|P{2.6cm}|Y|}
\hline
\rowcolor{TableHeaderBlue}
\color{white}\textbf{Critère} & \color{white}\textbf{Exigence} \\
\hline
Performance & Analyse complète compétences comportementales en moins de 15 secondes et APIs REST courantes en moins de 2 secondes \\
\hline
Confidentialité & Traitement 100\% local via Ollama pour les données RH sensibles \\
\hline
Fiabilité & Fallbacks systématiques pour chaque appel LLM avec timeouts configurés et réponses heuristiques par défaut \\
\hline
Scalabilité & Architecture conteneurisée Docker autorisant le déploiement multi-instances \\
\hline
Sécurité & Authentification JWT Bearer Token sur les endpoints Spring Boot et RBAC USER/ADMIN \\
\hline
\end{tabularx}
\caption{Objectifs non fonctionnels du Sprint 2}
\label{tab:objectifs-non-fonctionnels}
\end{table}

\subsubsection{Contraintes techniques}
L'analyse préliminaire des technologies retenues a révélé quatre contraintes majeures susceptibles d'impacter la fiabilité et les performances du pipeline d'évaluation. Ces contraintes ont été anticipées dès la phase de planification afin de concevoir des mécanismes de mitigation adaptés, notamment des fallbacks heuristiques et des timeouts explicites, avant toute décision d'architecture.
\begin{itemize}[leftmargin=1cm]
  \item \textbf{Latence Ollama :} le modèle llama3.2 peut dépasser 20 secondes, d'où des timeouts explicites et des fallbacks heuristiques immédiats.
  \item \textbf{Format JSON des modèles locaux :} les réponses Ollama contiennent fréquemment des artefacts Markdown ; un nettoyage systématique est donc appliqué dans les workflows n8n et les services Python.
  \item \textbf{PCM sans LLM :} le scoring des 18 questions est entièrement algorithmique afin d'assurer une forte reproductibilité et une exécution inférieure à 100 ms.
  \item \textbf{CodeLlama :} ce modèle étant plus lourd, le timeout des challenges de code est porté à 600 secondes.

  \item \textbf{GitHub rate limit :} en production, GITHUB\_TOKEN est requis pour éviter la limite de 60 requêtes par heure.
  \item \textbf{PDF scannés :} pdfplumber ne peut extraire le texte d'un PDF image, ce qui conduit à une analyse partielle accompagnée d'un avertissement.
\end{itemize}

%\figureplaceholder{Diagramme de contexte unifié compétences comportementales + compétences techniques}{fig:contexte-unifie}



\subsubsection{Cas d'utilisation}
Afin de formaliser les besoins fonctionnels de la plateforme TalentPredict AI et d'illustrer les interactions entre les différents acteurs et le système, cette section détaille la modélisation des cas d'utilisation. Nous présenterons d'abord une vision globale du processus d'évaluation, avant de décomposer ces interactions à travers les deux modules clés du Sprint 2 : l'évaluation des compétences comportementales et celle des compétences techniques.

\textbf{1. Diagramme de cas d'utilisation global (Sprint 2)}

Le diagramme global~\ref{fig:usecase-global-sprint2} illustre l'ensemble des interactions fonctionnelles entre les acteurs principaux  l'Employé et l'Administrateur RH  et le système TalentPredict AI dans le cadre du Sprint 2. Il regroupe les fonctionnalités liées à l'évaluation des compétences techniques et comportementales. L'employé soumet ses données (CV, profil GitHub), passe ses évaluations (test PCM, QCM adaptatif, scénario immersif, challenge de code) et consulte les recommandations générées. L'Administrateur RH supervise les évaluations, valide les résultats et génère des rapports. Les services externes   Service IA, Service GitHub et Service de communication temps réel   interviennent respectivement pour l'analyse intelligente, la récupération des données de dépôts, et la diffusion des mises à jour en temps réel.

\begin{figure}[H]
  \centering
  \includegraphics[width=1\textwidth]{images/D use case sprint2.png}
  \caption{Diagramme de cas d'utilisation : Passation et consultation de l'évaluation globale}
  \label{fig:usecase-global-sprint2}
\end{figure}

\begin{table}[H]
\centering
\small
\renewcommand{\arraystretch}{1.25}
\begin{tabularx}{\textwidth}{|P{0.26\textwidth}|Y|}
\hline
\rowcolor{TableHeaderBlue}
\color{white}\textbf{Élément} & \color{white}\textbf{Description} \\
\hline
Nom du cas & Passation et consultation de l'évaluation globale \\
\hline
Acteur principal & Employé \\
\hline
Acteurs secondaires & Administrateur RH, Service IA, Service GitHub, Service de communication temps réel \\
\hline
Pré-conditions & L'employé possède un compte actif et s'est authentifié sur TalentPredict. \\
\hline
Scénario nominal & \begin{minipage}[t]{\linewidth}
\begin{enumerate}[leftmargin=*,nosep]
  \item L'employé s'authentifie puis accède au module d'évaluation.
  \item Il soumet son CV et renseigne son nom d'utilisateur GitHub pour déclencher les analyses comportementale et technique.
  \item Le système évalue les compétences comportementales (CV sémantique, collaboration GitHub, test PCM, scénario immersif) et les compétences techniques (CV technique, activité GitHub, QCM adaptatif, challenge de code).
  \item Les résultats d'évaluation sont produits et des recommandations de formation sont générées par le Service IA.
  \item L'employé consulte les recommandations et le score de correspondance au poste.
  \item Le Service de communication temps réel diffuse les mises à jour ; l'Administrateur RH supervise, valide les résultats et génère les rapports.
\end{enumerate}
\end{minipage} \\
\hline
Scénario alternatif & Si un module est indisponible (ex.\ : timeout du Service IA ou profil GitHub introuvable), l'évaluation partielle est sauvegardée et le Service de communication temps réel prévient l'employé de la reprise future. \\
\hline
Post-conditions & Le profil complet du collaborateur est généré (compétences techniques et comportementales), validé par l'Administrateur RH, et prêt pour l'intégration au plan de formation. \\
\hline
\end{tabularx}
\renewcommand{\arraystretch}{1.05}
\caption{Fiche descriptive du cas d'utilisation : Passation et consultation de l'évaluation globale}
\end{table}

\textbf{2. Diagramme de cas d'utilisation : Module Compétences comportementales}

Ce diagramme~\ref{fig:usecase-soft-skills-sprint2} présente les fonctionnalités du module d'évaluation comportementale. L'employé soumet son CV et renseigne son nom d'utilisateur GitHub, ce qui déclenche respectivement l'extraction sémantique du CV et l'analyse de sa dynamique de collaboration. Il passe également le test PCM et répond à un scénario métier immersif. L'ensemble de ces sous-évaluations constitue le cas central \og Évaluer les compétences comportementales \fg{}. L'employé consulte ensuite son profil et son rapport comportemental. L'Administrateur RH valide le profil et génère le rapport Soft Skills. Le Service de communication temps réel diffuse les notifications de manière optionnelle via une relation \og extend \fg{}.

\begin{figure}[H]
  \centering
  \includegraphics[width=0.9\textwidth]{images/D use case soft.png}
  \caption{Diagramme de cas d'utilisation : Évaluation et restitution des Compétences comportementales}
  \label{fig:usecase-soft-skills-sprint2}
\end{figure}

\begin{table}[H]
\centering
\small
\renewcommand{\arraystretch}{1.25}
\begin{tabularx}{\textwidth}{|P{0.26\textwidth}|Y|}
\hline
\rowcolor{TableHeaderBlue}
\color{white}\textbf{Élément} & \color{white}\textbf{Description} \\
\hline
Nom du cas & Évaluation et restitution du profil comportemental \\
\hline
Acteur principal & Employé \\
\hline
Acteurs secondaires & Administrateur RH, Service IA, Service GitHub, Service de communication temps réel \\
\hline
Pré-conditions & L'employé s'est authentifié et a accédé au module \og Compétences comportementales \fg{} de son espace d'évaluation. \\
\hline
Scénario nominal & \begin{minipage}[t]{\linewidth}
\begin{enumerate}[leftmargin=*,nosep]
  \item L'employé soumet son CV ; le système en extrait le ton et la sémantique via le Service IA.
  \item L'employé renseigne son nom d'utilisateur GitHub ; le Service GitHub et le Service IA analysent sa dynamique de collaboration.
  \item L'employé passe le test PCM ; les scores des six dimensions sont calculés par le Service IA.
  \item L'employé répond à un scénario métier immersif généré et évalué par le Service IA.
  \item Le système consolide les quatre sources et produit le profil comportemental avec le type de personnalité dominant.
  \item L'employé consulte son profil et son rapport comportemental. L'Administrateur RH valide le profil et génère le rapport Soft Skills.
\end{enumerate}
\end{minipage} \\
\hline
Scénario alternatif & Si l'employé ne termine pas le scénario professionnel, le système sauvegarde l'état actuel et le Service de communication temps réel programme une relance automatique. \\
\hline
Post-conditions & La cartographie de personnalité de l'employé est finalisée, permettant de prescrire les meilleurs leviers de communication et d'apprentissage. \\
\hline
\end{tabularx}
\renewcommand{\arraystretch}{1.05}
\caption{Fiche descriptive du cas d'utilisation : Évaluation et restitution du profil comportemental}
\end{table}

\textbf{3. Diagramme de cas d'utilisation : Module Compétences techniques}

Ce diagramme~\ref{fig:usecase-hard-skills-sprint2} présente les fonctionnalités du module d'évaluation technique. L'employé soumet son CV technique et renseigne son nom d'utilisateur GitHub, déclenchant respectivement l'extraction des compétences techniques du CV et l'analyse de l'activité GitHub. Il passe également un QCM technique adaptatif et réalise un challenge de code, tous deux supervisés par le Service IA. Le cas central \og Évaluer les compétences techniques \fg{} inclut l'ensemble de ces analyses ainsi que le calcul automatique du niveau technique. L'employé consulte ensuite son niveau, son score de correspondance au poste et son rapport technique. L'Administrateur RH valide le niveau certifié et génère le rapport Hard Skills.

\begin{figure}[H]
  \centering
  \includegraphics[width=0.9\textwidth]{images/D uses case hard.png}
  \caption{Diagramme de cas d'utilisation : Évaluation multi-sources des Compétences techniques}
  \label{fig:usecase-hard-skills-sprint2}
\end{figure}

\begin{table}[H]
\centering
\small
\renewcommand{\arraystretch}{1.25}
\begin{tabularx}{\textwidth}{|P{0.26\textwidth}|Y|}
\hline
\rowcolor{TableHeaderBlue}
\color{white}\textbf{Élément} & \color{white}\textbf{Description} \\
\hline
Nom du cas & Évaluation multi-sources et certification des Compétences techniques \\
\hline
Acteur principal & Employé \\
\hline
Acteurs secondaires & Administrateur RH, Service IA, Service GitHub, Service de communication temps réel \\
\hline
Pré-conditions & L'employé s'est authentifié et a accédé au module \og Compétences techniques \fg{}. \\
\hline
Scénario nominal & \begin{minipage}[t]{\linewidth}
\begin{enumerate}[leftmargin=*,nosep]
  \item L'employé soumet son CV technique ; le Service IA en extrait les compétences techniques.
  \item L'employé renseigne son nom d'utilisateur GitHub ; le Service GitHub et le Service IA analysent son activité et sa stack technique.
  \item L'employé passe un QCM technique adaptatif généré par le Service IA selon ses compétences détectées.
  \item L'employé réalise un challenge de code évalué par le Service IA sur quatre critères.
  \item Le Service IA calcule le niveau technique par triangulation multi-sources (\textit{Beginner} à \textit{Expert}).
  \item L'employé consulte son niveau technique, son score de correspondance au poste et son rapport technique.
  \item L'Administrateur RH valide le niveau certifié et génère le rapport Hard Skills.
\end{enumerate}
\end{minipage} \\
\hline
Scénario alternatif & Si le nom d'utilisateur GitHub est erroné ou introuvable, le système en avertit l'employé via le Service de communication temps réel et base le calcul du niveau uniquement sur le CV, le QCM et le challenge de code. \\
\hline
Post-conditions & Les compétences techniques sont mesurées, classées et le système est prêt à recommander des formations ciblées pour combler les lacunes identifiées (\textit{Upskilling}). \\
\hline
\end{tabularx}
\renewcommand{\arraystretch}{1.05}
\caption{Fiche descriptive du cas d'utilisation : Évaluation multi-sources et certification des Compétences techniques}
\end{table}

\subsubsection{Sources d'évaluation et dimensions mesurées}
L'approche multi-sources retenue croise données objectives, évaluations auto-déclarées et tests interactifs afin de produire un profil employé complet. Chaque source apporte un signal complémentaire irréductible aux autres, comme détaillé dans les tableaux des deux modules ci-dessous.

\textbf{Module compétences comportementales}

\begin{table}[H]
\centering
\renewcommand{\arraystretch}{1.3}
\begin{tabularx}{\TableWidth}{|P{2.4cm}|P{3.2cm}|Y|}
\hline
\rowcolor{TableHeaderBlue}
\color{white}\textbf{Source} & \color{white}\textbf{Technologie} & \color{white}\textbf{Compétences comportementales évaluées} \\
\hline
GitHub & n8n + GitHub API + Ollama & Curiosité, collaboration, ownership et consistance \\
\hline
CV (texte brut) & n8n + Ollama & Indices de compétences comportementales, années d'expérience et niveau d'éducation \\
\hline
Test PCM (18 questions) & n8n (JavaScript heuristique, sans LLM) & Communication, discipline, curiosité, collaboration, ownership et leadership \\
\hline
Scénario IA & FastAPI + Ollama & Empathie, assertivité, pragmatisme et clarté de communication \\
\hline
\end{tabularx}
\caption{Sources d'évaluation du module compétences comportementales}
\label{tab:sources-soft-skills}
\end{table}

\textbf{Module compétences techniques}

\begin{table}[H]
\centering
\small
\renewcommand{\arraystretch}{1.3}
\begin{tabularx}{\textwidth}{|P{2.8cm}|P{4.0cm}|Y|}
\hline
\rowcolor{TableHeaderBlue}
\color{white}\textbf{Source} & \color{white}\textbf{Technologie} & \color{white}\textbf{Ce qui est évalué} \\
\hline
GitHub & GitHub API REST + Ollama & Langages, frameworks détectés et niveau d'activité \\
\hline
CV (PDF/DOCX) & pdfplumber + regex sur 150+ mots-clés & Technologies mentionnées et années d'expérience \\
\hline
LinkedIn & Agent Ollama & Rôle actuel, arc de carrière, certifications et cohérence CV/GitHub \\
\hline
Test QCM & Ollama (llama3.2) & Maîtrise technique par compétence selon exactitude, confiance et temps \\
\hline
Challenge de code & CodeLlama & Correctness (40), code quality (30), efficiency (20) et readability (10) \\
\hline
\end{tabularx}
\caption{Sources d'évaluation du module compétences techniques}
\label{tab:sources-hard-skills}
\end{table}



\subsubsection{4.5.2 Conception}

Cette phase a pour but de traduire les exigences fonctionnelles et les cas d’utilisation identifiés précédemment en solutions techniques concrètes. L’enjeu majeur de ce deuxième sprint réside dans la conception d’une architecture capable d’ingérer et de croiser des données hétérogènes (fichiers PDF, données d'API tierces, réponses dynamiques) via une intelligence artificielle locale, tout en garantissant la fiabilité et la vélocité de l'évaluation globale.

\paragraph*{1. Objectifs de conception}
\begin{itemize}[leftmargin=1cm]
    \item \textbf{Découplage et Modularité :} Isoler les traitements cognitifs lourds (parsing sémantique, inférence LLM) dans un microservice dédié (FastAPI) et confier l'orchestration des flux de données à n8n afin d'éviter la saturation du cœur transactionnel (Spring Boot).
    \item \textbf{Fiabilité Multi-Sources :} Concevoir des algorithmes capables de corroborer de manière robuste les données déclaratives (CV, profil LinkedIn) avec des preuves factuelles (historique de commits GitHub, résultats aux QCM et aux défis de code).
    \item \textbf{Asynchronisme et Réactivité :} Mettre en œuvre des flux d'exécution hautement asynchrones, couplés au Communication Hub (SSE), pour garantir que les temps d'inférence de l'IA (Ollama) n'impactent pas la fluidité et l'expérience utilisateur sur le frontend.
\end{itemize}

\paragraph*{2. Architecture globale du module}
L'architecture repose sur quatre  couches aux responsabilités strictement séparées : Angular collecte les données et restitue les résultats ; n8n orchestre les workflows compétences comportementales via GitHub API et Ollama ; FastAPI gère le module compétences techniques et les scénarios comportementaux avec llama3.2 et CodeLlama ; Spring Boot assure la persistance PostgreSQL, la sécurité JWT et le proxy vers FastAPI ; l'API GitHub constitue la source externe commune aux deux modules.
\begin{figure}[H]
  \centering
  \includegraphics[width=0.9\textwidth]{images/archiglob.png}
  \caption{Architecture globale du module d'évaluation des compétences }
  \label{fig:architecture-globale}
\end{figure}

%\figureplaceholder{Architecture globale du module d'évaluation des compétences (soft + compétences techniques)}{fig:architecture-globale}

L'architecture repose sur un modèle hybride orchestré autour de quatre couches technologiques :

\begin{table}[H]
\centering
\small
\renewcommand{\arraystretch}{1.25}
\begin{tabularx}{\TableWidth}{|P{2.8cm}|P{3.4cm}|Y|}
\hline
\rowcolor{TableHeaderBlue}
\color{white}\textbf{Couche} & \color{white}\textbf{Technologie} & \color{white}\textbf{Rôle} \\
\hline
Orchestration IA & n8n (workflows visuels) & Analyse GitHub soft, extraction CV, calcul PCM et synthèse finale \\
\hline
API IA comportementale & FastAPI + Ollama & Scénarios IA, analyse GitHub/CV/LinkedIn compétences techniques, QCM et code challenge \\
\hline
Backend métier & Spring Boot (Java 21) & Persistance PostgreSQL, sécurité JWT, API REST et rapports PDF \\
\hline
Frontend & Angular 18~\cite{angular}+ & Interface employé/admin/RH, visualisations et restitution des résultats \\
\hline
\end{tabularx}
\caption{Couches technologiques de l'architecture hybride}
\label{tab:couches-architecture-hybride}
\end{table}
L'architecture retenue repose sur un modèle hybride structuré autour de deux pipelines parallèles.





%\figureplaceholder{Diagramme architecture globale --- Vue complète des couches Frontend, n8n, FastAPI, Spring Boot, PostgreSQL, Ollama/CodeLlama avec flux de données}{fig:architecture-globale-couches-flux}

\subsubsection{Architecture n8n  Quatre workflows}



Le pipeline compétences comportementales s'articule autour du Master Agent n8n, qui déclenche
en parallèle via Promise.all() trois sous-workflows : CV Parser (analyse sémantique Ollama),
GitHub Analysis (score heuristique JavaScript + évaluation Ollama) et PCM Test (calcul algorithmique
pur sans LLM). Les résultats sont agrégés selon la pondération CV~40\%, PCM~30\%, GitHub~30\%,
avec basculement automatique vers CV~60\%/PCM~40\% si GitHub est indisponible. Ollama produit
ensuite le profil final (score global, type de personnalité, forces, recommandations). FastAPI
complète l'évaluation via un scénario IA sur quatre dimensions psychologiques en temps réel.
\begin{figure}[H]
  \centering
  \includegraphics[width=0.9\textwidth]{images/piplinesoft1.png}
   \caption{Pipeline d'évaluation des compétences comportementales}
  \label{fig:pipeline-soft-skills}
\end{figure}
Ce pipeline~\ref{fig:pipeline-soft-skills} combine les résultats du CV, du test PCM, de l'analyse GitHub et du scénario
comportemental IA afin de produire un profil soft skills unifié. Une étape de validation des
données en amont garantit la présence du CV, de l'identifiant GitHub et des réponses PCM avant
le lancement des workflows. Les scores agrégés sont ensuite sauvegardés en base de données et
restitués à l'interface utilisateur.



\color{black}
\subsubsection{Architecture FastAPI Structure du code}

talent\_agent.run\_agent() orchestre en parallèle trois sources :
l'analyseur GitHub (API GitHub REST), l'analyseur CV (extraction de plus de 150~technologies
via pdfplumber et expressions régulières) et un agent Ollama pour l'analyse LinkedIn.
Les PDF scannés déclenchent une analyse partielle avec avertissement transmis au frontend.
Le module de fusion et déduplication des compétences normalise les résultats via plus de
500~alias canoniques\,; le moteur de calcul du niveau de maîtrise attribue un niveau par
triangulation multi-sources (Débutant, Intermédiaire, Avancé, Expert)\,; le moteur de
correspondance au poste compare le profil avec plus de 20~profils métiers. Des tests interactifs
optionnels (QCM adaptatif Ollama, challenge de code CodeLlama) enrichissent le profil final
sans bloquer la génération des résultats. L'ensemble est persisté en PostgreSQL via Spring Boot
sous l'entité Skill de type TECH.

\begin{figure}[H]
  \centering
  \includegraphics[width=0.9\textwidth]{images/piplinehard1.png}
  \caption{Pipeline d'évaluation des compétences techniques}
  \label{fig:pipeline-hard-skills}
\end{figure}
Ce pipeline~\ref{fig:pipeline-hard-skills} plusieurs sources d'information (GitHub, CV, LinkedIn), applique une fusion
et une déduplication des compétences, calcule un niveau de maîtrise par triangulation et produit
un profil final enrichi par des tests interactifs optionnels. Les sorties principales sont les
compétences détectées avec leur niveau estimé, le score d'expérience, le score de correspondance
au poste cible et les recommandations associées.

\color{black}
\subsubsection{Modèles de données}

\textbf{Modèle entité-relation Spring Boot / PostgreSQL unifié}

Le modèle de persistance PostgreSQL s'articule autour d'une entité centrale Utilisateur reliée à huit tables satellites par des relations many-to-one. Competence enregistre chaque compétence détectée (TECHNIQUE ou COMPORTEMENTALE) avec son niveau (DEBUTANT, INTERMEDIAIRE, AVANCE, EXPERT), sa source (CV, GITHUB, PCM, IA\_PYTHON), un score de confiance (0-100) et son statut de validation administrative via le booléen validee.

\texttt{RésultatTest} enregistre les résultats détaillés des différents types de tests (QCM, CODE, SCENARIO, PCM) avec leur score et leur statut (EN\_ATTENTE, TERMINE, ECHOUE). Prédiction conserve la synthèse globale du profil candidat avec les scores décomposés en trois axes : score\_global, score\_technique etscore\_comportemental, ainsi que le type de personnalité et le score d'adéquation emploi.

\texttt{ProfilComportemental} dédie une table séparée aux données de personnalité (MBTI, leadership, collaboration, adaptabilité, communication) pour une meilleure traçabilité. Recommandationstocke les suggestions de formation personnalisées avec leur priorité (BASSE, MOYENNE, HAUTE) et leur type (MONTEE\_EN\_COMPETENCES, RECONVERSION). 

\texttt{ParcoursFormation} suit la progression des parcours d'upskilling avec un statut (ACTIF, TERMINE, EN\_PAUSE) et un pourcentage de progression. \texttt{PasseportCompétence} historise les générations du passeport de compétences (PDF, JSON, XML) avec un contrôle de version.

Les entités \texttt{Invitation} et \texttt{Équipe} gèrent les interactions collaboratives entre utilisateurs. Cette architecture complète permet de reconstruire intégralement le profil employé depuis son identifiant unique, avec l'historique complet de ses apprentissages, ses évaluations, ses recommandations et ses parcours de formation.

\begin{figure}[H]
  \centering
  \includegraphics[width=0.85\textwidth]{images/modele entitee relation.png}
  \caption{Diagramme entité-relation du modèle de persistance unifié de TalentPredict AI}
  \label{fig:entite-relation-skill}
\end{figure}

\textbf{Traçabilité et intégrité des données}

Toutes les tables incluent les champs date\_creation et date\_modification pour assurer une traçabilité complète. Les clés étrangères sont explicitement nommées (id\_utilisateur, id\_competence, etc.) et les contraintes d'unicité sont appliquées sur nom\_utilisateur et email. Le mot de passe est stocké sous forme de hash\_mot\_de\_passe pour garantir la sécurité des données sensibles. Les relations many-to-one garantissent l'intégrité référentielle entre l'entité centrale Utilisateur et ses données associées.



\color{black}

\subsubsection{Diagrammes de séquence}
Après avoir défini la cartographie fonctionnelle, il convient de modéliser la dimension dynamique et chronologique du système. Cette section présente les diagrammes de séquence qui détaillent l'enchaînement des messages et l'orchestration asynchrone des microservices lors des deux grands processus de la plateforme : l'évaluation comportementale et l'évaluation technique.
\textbf{1. Diagramme de Séquence : Évaluation des Compétences Comportementales}
Ce diagramme de séquence modélise le flux d'exécution asynchrone régissant l'évaluation comportementale complète du collaborateur. L'architecture met en exergue l'orchestration des microservices via n8n. Le processus débute par la collecte d'informations non structurées : le texte brut du CV est extrait sémantiquement via FastAPI, et l'historique de collaboration (messages de commit, dynamiques d'équipe sur les Pull Requests) est récupéré via l'API GitHub.
Parallèlement, le collaborateur répond au questionnaire heuristique PCM. L'orchestrateur n8n agrège ces trois vecteurs de données textuelles et les soumet à l'IA locale (Ollama) pour générer un scénario immersif hyper-personnalisé. Les réactions de l'employé face à cette mise en situation sont ensuite analysées pour déduire les six dimensions cognitives de son profil. La persistance de ces métriques déclenche un événement de propagation (Server-Sent Events) par l'intermédiaire du Communication Hub, assurant ainsi une notification fluide sur l'interface utilisateur.

\begin{figure}[H]
  \centering
  \includegraphics[width=1\textwidth]{images/sequencesoft2.png}
  \caption{ Diagramme de sequence d'Évaluation des Compétences Comportementales}
  \label{fig:seq_soft}
\end{figure}

\textbf{2. Diagramme de Séquence : Évaluation des compétences technique} 
Ce second diagramme de séquence formalise le paradigme d'évaluation multi-sources implémenté pour objectiver l'expertise technique. Le flux d'exécution débute par la délégation de l'extraction documentaire (CV) au microservice Python (FastAPI). En parallèle, le noyau Spring Boot interroge l'API GitHub pour corroborer les compétences déclarées avec des traces d'activité réelles.

L'agrégation de ces vecteurs d'information permet au système de solliciter l'IA (CodeLlama/Ollama) pour instancier un environnement d'évaluation adaptatif en deux phases : un QCM sur-mesure pour valider les connaissances théoriques, suivi d'un challenge de code pour éprouver la logique pratique. La correction automatisée de cette double soumission par le LLM permet le calcul d'un Score de Confiance global. Ce triptyque de validation (Déclaratif, Factuel, Évaluatif) est persistant en base de données, et l'utilisateur est instantanément notifié de son nouveau niveau d'expertise via le Communication Hub.
\begin{figure}[H]
  \centering
  \includegraphics[width=0.9\textwidth]{images/sequencehard2.png}
  \caption{Diagramme de sequence d'evaluation des competences technique}
  \label{fig:seq_hard}
\end{figure}

\subsubsection{Diagramme de classes  Module scénarios IA }

Le microservice FastAPI s'articule autour de cinq classes Pydantic principales : ScenarioRequest (entrées rôle/niveau), ScenarioResponse (titre, description, compétences testées), EvaluationResult (quatre scores psychologiques, forces, axes d'amélioration, profil culture-add), EmployéeAnalysis (données consolidées GitHub/CV/LinkedIn) et Skill (attributs name, level, score, sources), partagée entre les deux modules. Toutes les entrées sont validées automatiquement par Pydantic avant traitement.
\begin{figure}[H]
  \centering
  \includegraphics[width=0.85\textwidth]{images/classes_fastapi.png}
  \caption{Diagramme de classes du module scénarios IA et compétences techniques (FastAPI) }
  \label{fig:classes-fastapi}
\end{figure}
%\figureplaceholder{Diagramme de classes du module scénarios IA et compétences techniques (FastAPI)}{fig:classes-fastapi}


\subsection{Codage}
Cette section documente la phase d'implémentation effective du sprint. Elle présente les technologies déployées, l'organisation du code source, et les bonnes pratiques appliquées pour garantir la maintenabilité et la lisibilité du système. Chaque composant clé est détaillé avec des extraits de code commentés illustrant les choi-x techniques retenus.
\subsubsection{Stack technologique }

{\footnotesize
\setlength{\tabcolsep}{3pt}
\begin{table}[H]
\centering
\footnotesize
\setlength{\tabcolsep}{3pt}
\begin{tabular}{|P{2.3cm}|P{2.2cm}|P{1.2cm}|P{4.3cm}|}
\hline
\rowcolor{TableHeaderBlue}
\color{white}\textbf{Composant} & \color{white}\textbf{Technologie} & \color{white}\textbf{Version} & \color{white}\textbf{Justification} \\
\hline
Orchestration workflows & n8n & latest & Interface visuelle, modification des prompts sans recompilation \\
\hline
LLM local généraliste & Ollama (llama3.2) & latest & Déploiement local et conformité RGPD \\
\hline
LLM spécialisé code & CodeLlama & latest & Meilleures performances pour la génération et l'évaluation de code \\
\hline
API IA & FastAPI & 0.110+ & Support asynchrone et auto-documentation OpenAPI \\
\hline
Extraction PDF client & pdfjs-dist & 5.6.205 & Extraction navigateur pour le module compétences comportementales \\
\hline
Extraction PDF serveur & pdfplumber & 0.10+ & Parsing précis pour le module compétences techniques \\
\hline
Backend métier & Spring Boot + Java 21 & 3.x & API REST, persistance JPA et sécurité JWT \\
\hline
Base de données & PostgreSQL & 15+ & Stockage \texttt{Skill}, \texttt{User}, \texttt{Prediction} et \texttt{TestResult} \\
\hline
Conteneurisation & Docker Compose & 24+ & Isolation des services et reproductibilité \\
\hline
\end{tabular}
\caption{Stack technologique du Sprint 2}\label{tab:stack-technologique}
\end{table}
}

\subsubsection{Implémentation des workflows n8n}

\paragraph{Workflow 1 : GitHub Analysis}
L'objectif de ce workflow est d'analyser l'activité open source du employé afin de déduire quatre compétences comportementales objectivables : curiosité, collaboration, ownership et consistance.
\begin{figure}[H]
  \centering
  \includegraphics[width=0.9\textwidth]{images/workflow-github.png}
  \caption{Capture du workflow GitHub Analysis dans n8n}
  \label{fig:workflow-github-analysis}
\end{figure}
%\figureplaceholder{Capture du workflow GitHub Analysis dans n8n}{fig:workflow-github-analysis}

\textbf{Nœud \og Compute GitHub Score \fg{}}

Le score GitHub est calculé de manière purement heuristique en JavaScript, sans appel LLM, afin de garantir rapidité et reproductibilité. Le calcul repose sur trois indicateurs pondérés selon le tableau ci-dessous.

\begin{table}[H]
\centering
\small
\renewcommand{\arraystretch}{1.25}
\begin{tabularx}{\TableWidth}{|P{3.0cm}|Y|P{2.1cm}|}
\hline
\rowcolor{TableHeaderBlue}
\color{white}\textbf{Indicateur} & \color{white}\textbf{Formule} & \color{white}\textbf{Poids maximal} \\
\hline
Nombre de dépôts & \texttt{Math.min((repoCount / 20) $\times$ 40, 40)} & 40 pts \\
\hline
Nombre d'étoiles & \texttt{Math.min((totalStars / 50) $\times$ 40, 40)} & 40 pts \\
\hline
Nombre de forks & \texttt{Math.min((totalForks / 20) $\times$ 20, 20)} & 20 pts \\
\hline
\textbf{Score final} & \texttt{Math.round(repoScore + starScore + forkScore)} & \textbf{100 pts} \\
\hline
\end{tabularx}
\caption{Calcul heuristique du score GitHub}
\label{tab:compute-github-score}
\end{table}

%\figureplaceholder{Capture du code JavaScript du nœud Compute GitHub Score}{fig:compute-github-score}

Le nœud \og Analyze Compétences comportementales \fg{} soumet ensuite les données à Ollama à travers un prompt structuré centré sur les variables de contribution et de diversité technique.


\begin{lstlisting}
 Tu es un expert RH chargé d'analyser le profil GitHub d'un développeur. À partir des données suivantes : \{repo\_count\} dépôts, \{total\_stars\} étoiles, \{total\_forks\} forks et langages utilisés : \{languages\}, évalue uniquement les quatre compétences comportementales objectivement mesurables via GitHub. Retourne exclusivement un JSON valide contenant : curiosity (0--100), collaboration (0--100), ownership (0--100) et consistency (0--100).
\end{lstlisting}


%\figureplaceholder{Capture du nœud Ollama --- prompt Analyze Compétences comportementales}{fig:prompt-github-soft}
%\figureplaceholder{Exemple de sortie JSON du workflow GitHub Analysis}{fig:output-github-analysis}

\paragraph{Workflow 2 : CV Parser}
Ce workflow a pour objectif d'extraire les indices de compétences comportementales depuis le texte brut du CV via une analyse sémantique Ollama.

%\figureplaceholder{Capture du workflow CV Parser dans n8n}{fig:workflow-cv-parser}
\begin{figure}[H]
  \centering
  \includegraphics[width=0.85\textwidth]{images/workflow_cv.png}
  \caption{Capture du workflow CV Parser dans n8n}
  \label{fig:workflow-cv-parser}
\end{figure}

\begin{lstlisting}
Tu es un expert RH analysant un CV. Extrais uniquement les signaux de compétences comportementales depuis ce texte. Retourne uniquement un JSON valide avec : behavioral_skills_hints, experience_years, education_level et cv_score.
\end{lstlisting}


%\figureplaceholder{Capture du nœud Ollama --- prompt CV Parser}{fig:prompt-cv-parser}
%\figureplaceholder{Exemple de sortie JSON du workflow CV Parser}{fig:output-cv-parser}

\paragraph{Workflow 3 : PCM Test}
L'objectif est de calculer les scores PCM sur six dimensions à partir de 18 réponses auto-évaluées, sans recours à un LLM.

%\figureplaceholder{Capture du workflow PCM Test dans n8n}{fig:workflow-pcm}
\begin{figure}[H]
  \centering
  \includegraphics[width=0.8\textwidth]{images/workflow_pcm.png}
  \caption{Capture du workflow PCM Test dans n8n}
  \label{fig:workflow-pcm}
\end{figure}
\textbf{Nœud \og Score PCM Answers \fg{}}

Le scoring PCM est entièrement algorithmique. Chaque dimension est calculée comme la moyenne arrondie à une décimale de ses trois questions associées, puis le score global correspond à la moyenne des six dimensions.

\begin{table}[H]
\centering
\small
\renewcommand{\arraystretch}{1.25}
\begin{tabularx}{\TableWidth}{|P{2.6cm}|P{2.2cm}|Y|}
\hline
\rowcolor{TableHeaderBlue}
\color{white}\textbf{Dimension} & \color{white}\textbf{Questions} & \color{white}\textbf{Formule} \\
\hline
Communication & q1, q2, q3 & \texttt{Math.round((q1+q2+q3) / 3 $\times$ 10) / 10} \\
\hline
Discipline & q4, q5, q6 & \texttt{Math.round((q4+q5+q6) / 3 $\times$ 10) / 10} \\
\hline
Curiosité & q7, q8, q9 & \texttt{Math.round((q7+q8+q9) / 3 $\times$ 10) / 10} \\
\hline
Collaboration & q10, q11, q12 & \texttt{Math.round((q10+q11+q12) / 3 $\times$ 10) / 10} \\
\hline
Ownership & q13, q14, q15 & \texttt{Math.round((q13+q14+q15) / 3 $\times$ 10) / 10} \\
\hline
Leadership & q16, q17, q18 & \texttt{Math.round((q16+q17+q18) / 3 $\times$ 10) / 10} \\
\hline
\textbf{Score global} & Toutes & Moyenne des six dimensions, arrondie à une décimale \\
\hline
\end{tabularx}
\caption{Calcul algorithmique des scores PCM}
\label{tab:score-pcm-answers}
\end{table}

%\figureplaceholder{Capture du code JavaScript du nœud Compute PCM Scores}{fig:compute-pcm-scores}
%\figureplaceholder{Exemple de sortie JSON du workflow PCM Test}{fig:output-pcm}

\paragraph{Workflow 4 : Master Compétences comportementales Agent}
Ce workflow orchestre les trois sous-workflows en parallèle puis produit une synthèse pondérée via Ollama.

%\figureplaceholder{Capture du workflow Master Compétences comportementales Agent dans n8n}{fig:workflow-master-agent}
\begin{figure}[H]
  \centering
  \includegraphics[width=0.9\textwidth]{images/workflow_master.png}
  \caption{Capture du workflow Master Compétences comportementales Agent dans n8n}
  \label{fig:workflow-master-agent}
\end{figure}

\textbf{Nœud \og Prepare \& Call Sub-Workflows \fg{}}

L'orchestration des trois sous-workflows est réalisée en parallèle via \texttt{Promise.all}, ce qui réduit la latence globale du Master Agent en évitant les appels séquentiels. Le tableau ci-dessous détaille les paramètres transmis à chaque sous-workflow.

\begin{table}[H]
\centering
\small
\renewcommand{\arraystretch}{1.25}
\begin{tabularx}{\TableWidth}{|P{3.0cm}|P{3.2cm}|Y|}
\hline
\rowcolor{TableHeaderBlue}
\color{white}\textbf{Sous-workflow} & \color{white}\textbf{Endpoint appelé} & \color{white}\textbf{Paramètres transmis} \\
\hline
CV Parser & \texttt{behavioral-}\newline\texttt{skills-}\newline\texttt{extract} & \texttt{extracted\_cv\_text}, \texttt{full\_name}, \texttt{email} \\
\hline
PCM Test & \texttt{pcm-test} & \texttt{full\_name}, \texttt{q1} à \texttt{q18} \\
\hline
GitHub Analysis & \texttt{github-analysis} & \texttt{github\_username} \\
\hline
\end{tabularx}
\caption{Paramètres transmis aux sous-workflows du Master Agent}
\label{tab:prepare-call-subworkflows}
\end{table}

Les trois appels sont déclenchés simultanément via \texttt{await Promise.all([...])} et leurs résultats sont ensuite agrégés avant d'être soumis au nœud d'analyse Ollama pour la synthèse pondérée finale.

%\figureplaceholder{Capture du code JavaScript du nœud Prepare \& Call Sub-Workflows}{fig:prepare-call-subworkflows}

Le nœud \og Master AI Analysis \fg{} envoie les résultats agrégés à Ollama afin de produire un profil comportemental unifié avec score global, type de personnalité, compétences consolidées, forces et recommandations.

Prompt Ollama clé Master Agent synthèse :

\begin{lstlisting}
You are an expert in organizational psychology. Synthesize this multi-source behavioral data.
Data sources and weights:
- GitHub (30\%): behavioral signals from open-source activity
- CV (40\%): behavioral signals from professional background
- PCM (30\%): self-assessed personality dimensions
Return ONLY valid JSON:
{"overall_score": <0-100>, "personality_type": "<French>", 
 "merged_behavioral_skills": {...}, "key_strengths": [...], 
 "training_recommendations": "..."}
\end{lstlisting}
Output exemple :
\begin{lstlisting}
{
    "overall_score": 72,
    "personality_type": "Innovateur",
    "skills": [
        {
            "name": "communication",
            "score": 80,
            "source": "pcm"
        },
        {
            "name": "problem_solving",
            "score": 90,
            "source": "github"
        },
        {
            "name": "adaptability",
            "score": 70,
            "source": "cv"
        },
        {
            "name": "leadership",
            "score": 60,
            "source": "pcm"
        },
        {
            "name": "teamwork",
            "score": 50,
            "source": "github"
        }
    ],
    "merged_behavioral_skills": {
        "communication": 80,
        "leadership": 60,
        "teamwork": 50,
        "adaptability": 70,
        "problem_solving": 90
    },
    "source_data": {
        "cv_score": 75,
        "pcm_score": 85,
        "github_score": 92
    },
    "summary": "",
    "key_strengths": [],
    "training_recommendations": ""
}
\end{lstlisting}

%\figureplaceholder{Capture du nœud Ollama --- prompt Master AI Analysis}{fig:prompt-master-agent}

%\figureplaceholder{Exemple de sortie JSON du Master Agent}{fig:output-master-agent}


\subsubsection{Implémentation FastAPI}

Le microservice FastAPI centralise l'ensemble de la logique d'évaluation comportementale et technique en deux catégories de composants.

Les services à appel LLM sollicitent Ollama avec des températures calibrées selon le degré de créativité requis : \texttt{generate\_soft\_skills\_scenario} (température 0.7) produit un dilemme professionnel contextualisé par rôle et niveau, tandis que \texttt{evaluate\_scenario\_response} (température 0.2) évalue la réponse libre du employé sur quatre dimensions psychologiques. \texttt{analyze\_github\_profile} (température 0.3) enrichit les statistiques de dépôts publics par une synthèse sémantique Ollama.

Les services à logique purement algorithmique n'effectuent aucun appel LLM, garantissant des temps de réponse inférieurs à 200 ms et une reproductibilité totale : analyze\_cv extrait 150+ technologies par expressions régulières depuis un PDF ou DOCX, extract\_skills normalise les compétences issues des trois sources via un dictionnaire de 500+ alias, \texttt{\_per\_question\_weight} calcule le poids par réponse selon la formule exactitude $\times$ confiance $\times$ temps, \_determine\_level attribue un niveau Beginner $\rightarrow$ Expert par triangulation multi-sources, et match\_job produit un score de correspondance pour 20+ profils métiers.


\subsubsection{Bonnes pratiques appliquées}

Quatre bonnes pratiques transversales ont guidé l'implémentation des deux modules afin de garantir performance, fiabilité et maintenabilité du pipeline.

\textbf{1. Séparation calcul heuristique / analyse LLM}

Le LLM n'est sollicité que lorsque l'analyse sémantique est strictement nécessaire. Toute métrique déterministe est traitée sans appel Ollama : le score GitHub et les dimensions PCM sont calculés en JavaScript pur ($< 100$ ms), la détection de 150+ technologies par regex et la déduplication de 500+ alias en Python ($< 200$ ms), et le scoring QCM par algorithme pondéré ($< 100$ ms). Seules l'analyse sémantique GitHub, la synthèse Master Agent et l'évaluation de scénario sollicitent Ollama (5--25 s). Cette séparation a réduit de 60\% le nombre d'appels LLM par rapport à une approche tout-Ollama.

\textbf{2. Nettoyage systématique des réponses JSON Ollama}

Le modèle \texttt{llama3.2} produit fréquemment des réponses enveloppées dans des blocs Markdown ou contenant des artefacts textuels. Un mécanisme de nettoyage en quatre étapes est appliqué systématiquement dans tous les workflows n8n et services Python : suppression des balises Markdown, tentative de parsing direct via JSON.parse(), extraction du premier objet JSON valide par regex \texttt{\textbackslash\{[\textbackslash s\textbackslash S]*\textbackslash\}}, puis retour d'un fallback sécurisé avec scores à 0 si le parsing demeure impossible. Ce mécanisme garantit qu'aucun appel Ollama ne bloque le pipeline.

\textbf{3. Gestion des timeouts et fallbacks}

Chaque appel externe est encapsulé dans un mécanisme de protection à deux niveaux : un timeout explicite adapté à la complexité de la tâche et un fallback heuristique immédiat en cas d'échec. Les valeurs retenues vont de 10 s (GitHub API) à 600 s (CodeLlama), chaque échec déclenchant automatiquement une réponse de substitution exploitable  scénario prédéfini, banque de questions fallback ou starter\_code générique. La fonction call\_ollama\_json centralise cette logique dans FastAPI en interceptant toute exception httpx.TimeoutException ou erreur réseau.

\textbf{4. Déduplication des compétences  Dictionnaire d'alias}

La fusion de sources hétérogènes (GitHub, CV, LinkedIn) génère des doublons lorsqu'une même technologie est désignée par des noms différents. Le dictionnaire \_ALIASES dans skill\_extractor.py normalise 500+ variantes vers leur forme canonique --- par exemple reactjs $\rightarrow$ React, k8s $\rightarrow$ Kubernetes, tf $\rightarrow$ TensorFlow, postgres $\rightarrow$ \texttt{PostgreSQL}. La fonction \_normalize\_skill\_name convertit le nom brut en minuscules, supprime tirets et underscores, puis recherche la correspondance dans \_ALIASES, retournant à défaut la forme title() du nom original. Cette normalisation est critique pour éviter qu'une compétence présente dans plusieurs sources ne soit comptabilisée avec des scores fragmentés.


\subsection{Réalisation}

Cette section présente les interfaces développées lors du Sprint 2, illustrant la concrétisation des pipelines d'évaluation multi-source (compétences techniques et comportementales).

\subsubsection{Module d'Évaluation Comportementale (\textit{Compétences comportementales})}

\textbf{1. Extraction de données (GitHub \& CV)}

L'interface permet au employé de soumettre son identifiant GitHub et son CV. Le système traite ces entrées pour extraire instantanément de premières tendances comportementales basées sur l'historique d'engagement et la sémantique du texte.
\begin{figure}[H]
  \centering
  \includegraphics[width=0.6\textwidth]{images/soft1.png}
  \caption{Interface de soumission des profils et extraction CV/GitHub}
  \label{fig:soft-submission-cv-github}
\end{figure}
%\figureplaceholder{Interface de soumission des profils et extraction CV/GitHub}{fig:soft-submission-cv-github}

\textbf{2. Test Process Communication Model (PCM)}

L'employe complète un questionnaire de personnalité de 18 questions. L'interface restitue immédiatement le profil psychologique dominant grâce à un calcul purement algorithmique.
\begin{figure}[H]
  \centering
  \includegraphics[width=0.5\textwidth]{images/softPCM.png}
  \caption{Interface du test PCM et restitution du profil dominant}
  \label{fig:soft-pcm-test-profile}
\end{figure}
%\figureplaceholder{Interface du test PCM et restitution du profil dominant}{fig:soft-pcm-test-profile}

\textbf{3. Paramétrage du Scénario (Rôle et Niveau)}

Pour garantir la pertinence de la mise en situation, le employé définit son rôle métier cible et son niveau de séniorité. Ces paramètres permettent au modèle de calibrer dynamiquement le contexte du test.
\begin{figure}[H]
  \centering
  \includegraphics[width=0.65\textwidth]{images/SoftScenarioGenerate.png}
  \caption{Paramétrage contextuel de l'évaluation IA}
  \label{fig:soft-scenario-parameters}
\end{figure}
%\figureplaceholder{Paramétrage contextuel de l'évaluation IA}{fig:soft-scenario-parameters}

\textbf{4. Évaluation par Scénario IA}

L'agent conversationnel soumet un dilemme professionnel sur mesure. Le système analyse la réponse en texte libre du employé pour évaluer quatre dimensions clés : empathie, assertivité, pragmatisme et clarté.
\begin{figure}[H]
  \centering
  \includegraphics[width=0.7\textwidth]{images/softScenario.png}
  \caption{Interaction avec le scénario comportemental généré par IA}
  \label{fig:soft-ai-scenario-interaction}
\end{figure}
%\figureplaceholder{Interaction avec le scénario comportemental généré par IA}{fig:soft-ai-scenario-interaction}

\textbf{5. Résultats Comportementaux}

Le module consolide l'ensemble des scores Compétences comportementales et génère un tableau de bord synthétique, offrant un diagnostic sous forme de graphique radar et des recommandations détaillées.
\begin{figure}[H]
  \centering
  \includegraphics[width=0.4\textwidth]{images/softResult1.png}\hfill
  \includegraphics[width=0.45\textwidth]{images/softResult2.png}
  \caption{Restitution globale des résultats Compétences comportementales}
  \label{fig:soft-skills-results-global}
\end{figure}
%\figureplaceholder{Restitution globale des résultats Compétences comportementales}{fig:soft-skills-results-global}

\subsubsection{Module d'Évaluation Technique (\textit{Compétences techniques})}

\textbf{1. Collecte Technique (GitHub, LinkedIn, CV)}

Le employé renseigne ses liens professionnels. Un pipeline d'extraction robuste croise ces trois sources (profils publics et document parsé) pour identifier exhaustivement les technologies manipulées.
\begin{figure}[H]
  \centering
  \includegraphics[width=0.65\textwidth]{images/hard1.png}
  \caption{Formulaire de collecte des sources techniques}
  \label{fig:hard-technical-sources-collection}
\end{figure}
%\figureplaceholder{Formulaire de collecte des sources techniques}{fig:hard-technical-sources-collection}

\textbf{2. Cartographie et Job Matching}

Le système restitue le dictionnaire des compétences extraites, classées par catégorie. Un algorithme calcule instantanément le score d'adéquation (\textit{Job Matching}) du employé face à des profils métiers standards.
\begin{figure}[H]
  \centering
  \includegraphics[width=0.75\textwidth]{images/hard2.png}
  \caption{Cartographie des compétences et taux de Job Matching}
  \label{fig:hard-skills-job-matching}
\end{figure}
%\figureplaceholder{Cartographie des compétences et taux de Job Matching}{fig:hard-skills-job-matching}

\textbf{3. QCM Technique Adaptatif}

Afin de valider le niveau de maîtrise déclaré, la plateforme génère dynamiquement un questionnaire à choix multiples. La difficulté s'adapte aux technologies extraites précédemment.
\begin{figure}[H]
  \centering
  \includegraphics[width=0.75\textwidth]{images/hard3.png}
  \caption{Interface du QCM technique adaptatif}
  \label{fig:hard-adaptive-qcm}
\end{figure}
%\figureplaceholder{Interface du QCM technique adaptatif}{fig:hard-adaptive-qcm}
\textbf{4. Résultats de l'Évaluation Technique}
À l'issue du QCM adaptatif et des épreuves techniques, le système calcule et affiche un score de performance global. L'interface détaille la justesse des réponses par technologie et attribue un niveau de maîtrise consolidé, venant ainsi valider les compétences initialement extraites.

\begin{figure}[H]
  \centering
  \includegraphics[width=0.7\textwidth]{images/hardres1.png}
  \caption{Résultats détaillés de l'évaluation technique et score final}
  \label{fig:hard-skills-results-global}
\end{figure}
\subsubsection{Bilan Consolidé}
\textbf{1. Profil de Compétences et Carrière (Global)}

L'aboutissement de ce sprint est le tableau de bord global. Il fusionne les évaluations techniques et comportementales, fournissant aux admin/RH une vision unifiée du potentiel du employé et des opportunités d'évolution (\textit{Upskilling}).
\begin{figure}[H]
  \centering
  \includegraphics[width=0.9\textwidth]{images/res_tech_soft.png}
  \caption{Tableau de bord global : profil de compétences unifié compétences techniques et comportementales}
  \label{fig:global-hard-soft-dashboard}
\end{figure}
%\figureplaceholder{Tableau de bord global : profil de compétences unifié compétences techniques et comportementales}{fig:global-hard-soft-dashboard}

\subsection{Tests et Déploiement}
La validation du sprint repose sur une stratégie de tests multi-niveaux : tests unitaires des workflows n8n, tests d'intégration des endpoints FastAPI, tests de robustesse des fallbacks en cas de défaillance LLM, et validation fonctionnelle end-to-end. Cette section présente les résultats de ces campagnes de tests ainsi que le processus de déploiement containerisé
\subsubsection{Tests réalisés}

\textbf{Tests des workflows n8n (module compétences comportementales)}
Chaque workflow a été validé directement depuis l'interface n8n en mode « Test Workflow », permettant de vérifier en temps réel la validité du JSON retourné, le bon enchaînement des nœuds et le déclenchement correct des mécanismes de fallback. Le tableau ci-dessous synthétise les cas testés et les résultats obtenus pour les quatre workflows du module compétences comportementales.
\begin{table}[H]
\centering
\small
\renewcommand{\arraystretch}{1.25}
\begin{tabularx}{\TableWidth}{|P{2.3cm}|P{2.5cm}|X|X|}
\hline
\rowcolor{TableHeaderBlue}
\color{white}\textbf{Workflow} & \color{white}\textbf{Webhook} & \color{white}\textbf{Cas testés} & \color{white}\textbf{Résultat} \\
\hline
GitHub Analysis & POST /webhook/\newline github-analysis & Nom d'utilisateur valide, vide et inexistant & Succès avec fallback si vide \\
\hline
CV Parser & POST /webhook/\newline behavioral-\newline skills-extract & CV riche et CV vide inférieur à 50 caractères & JSON structuré avec cv\_score = 0 si vide \\
\hline
PCM Test & POST /webhook/\newline pcm-test & 18 réponses numériques & Six scores, top strength et faiblesse principale \\
\hline
Master Agent & POST /webhook/\newline master-agent & Payload complet et GitHub indisponible & Synthèse produite avec pondération dynamique \\
\hline
\end{tabularx}
\caption{Tests des workflows n8n}
\label{tab:tests-n8n}
\end{table}

\textbf{Tests FastAPI   routes compétences comportementales}
Les endpoints comportementaux du module compétences comportementales ont été testés via la documentation interactive Swagger UI . Les deux routes principales ont été soumises à des cas représentatifs couvrant la génération de scénarios et l'évaluation de réponses libres. Le tableau ci-dessous présente les résultats obtenus.
\begin{table}[H]
\centering
\small
\renewcommand{\arraystretch}{1.25}
\begin{tabularx}{\textwidth}{|P{4.8cm}|Y|Y|}
\hline
\rowcolor{TableHeaderBlue}
\color{white}\textbf{Endpoint} & \color{white}\textbf{Cas testé} & \color{white}\textbf{Résultat} \\
\hline
POST \texttt{/api/test/}\newline\texttt{scenario/generate} & Role = SWE, Level = Junior & JSON avec \texttt{scenario\_title} et \texttt{skills\_tested} \\
\hline
POST \texttt{/api/test/}\newline\texttt{scenario/evaluate} & Réponse employée libre & Quatre scores psychologiques sur 100 \\
\hline
\end{tabularx}
\caption{Tests FastAPI pour les routes compétences comportementales}
\label{tab:tests-fastapi-soft}
\end{table}

\textbf{Tests FastAPI   routes compétences techniques}
L'ensemble des routes du module compétences techniques a été validé via Swagger UI en soumettant des payloads réels incluant un profil GitHub actif, un CV PDF vectoriel et une URL LinkedIn. Les tests couvrent le pipeline complet d'analyse employé ainsi que les deux types d'évaluation interactive : QCM adaptatif et challenge de code. Le tableau ci-dessous détaille les cas testés et les résultats observés.
\begin{table}[H]
\centering
\small
\renewcommand{\arraystretch}{1.25}
\begin{tabularx}{\TableWidth}{|P{3.2cm}|P{1.7cm}|X|X|}
\hline
\rowcolor{TableHeaderBlue}
\color{white}\textbf{Endpoint} & \color{white}\textbf{Méthode} & \color{white}\textbf{Cas testé} & \color{white}\textbf{Résultat} \\
\hline
/analyze-employée & POST (multipart) & GitHub + CV PDF + LinkedIn & JSON complet avec skills[] et job\_match \\
\hline
/api/test/generate & POST & Skill = Python, Level = Intermediate, n = 8 & Huit questions MCQ valides \\
\hline
/api/test/evaluate & POST & Huit réponses avec confiance déclarée & score réel, scores par compétence et précision de confiance \\
\hline
/api/test/code-challenge & POST & Skill = Python, Level = Junior & Type, starter\_code, hints[] et language \\
\hline
/api/test/code-submit & POST & Code soumis avec un indice utilisé & Score sur 100 avec détail des quatre critères \\
\hline
\end{tabularx}
\caption{Tests FastAPI pour les routes compétences techniques}
\label{tab:tests-fastapi-hard}
\end{table}

\textbf{Tests Spring Boot  API Skills}
La couche de persistance exposée par Spring Boot a été testée via la documentation Swagger UI du backend Java accessible à http://localhost:8080/swagger-ui.html, en s'authentifiant avec un token JWT valide. Les quatre opérations CRUD de l'entité Skill ont été vérifiées pour les deux rôles du système : USER pour la création et la consultation, et ADMIN pour la validation et la suppression. Le tableau ci-dessous résume les résultats obtenus.
\begin{table}[H]
\centering
\small
\renewcommand{\arraystretch}{1.25}
\begin{tabularx}{\textwidth}{|P{4.8cm}|Y|P{2.8cm}|}
\hline
\rowcolor{TableHeaderBlue}
\color{white}\textbf{Endpoint} & \color{white}\textbf{Test} & \color{white}\textbf{Résultat} \\
\hline
POST \texttt{/api/skills/}\newline\texttt{accounts/\{id\}} & Création d'une skill TECH Python niveau 3 & UUID retourné \\
\hline
GET /api/skills/\newline accounts/\{id\}/\newline type/TECH & Liste des skills techniques & Liste filtrée par type \\
\hline
PUT \texttt{/api/skills/}\newline\texttt{\{id\}/valider} & Validation d'une skill par ADMIN & validee = true \\
\hline
DELETE /api/skills/\{id\} & Suppression d'une skill & Code HTTP 204 \\
\hline
\end{tabularx}
\caption{Tests Spring Boot pour l'API Skills}
\label{tab:tests-spring-skills}
\end{table}

\textbf{Tests de robustesse   fallbacks}
La robustesse du pipeline face aux défaillances des services externes constitue une exigence critique du module. Des tests de panne simulée ont été conduits en arrêtant successivement Ollama, CodeLlama et en soumettant des PDF scannés sans couche texte extractible. Le tableau ci-dessous confirme que chaque scénario de défaillance déclenche bien le mécanisme de fallback prévu sans interruption du service.
\begin{table}[H]
\centering
\small
\renewcommand{\arraystretch}{1.25}
\begin{tabularx}{\TableWidth}{|X|P{2.3cm}|Y|P{1.2cm}|}
\hline
\rowcolor{TableHeaderBlue}
\color{white}\textbf{Cas} & \color{white}\textbf{Module} & \color{white}\textbf{Comportement} & \color{white}\textbf{Résultat} \\
\hline
Ollama hors ligne & GitHub Analysis n8n & Le nœud JavaScript retourne des scores nuls avec parse\_error = true & OK \\
\hline
Ollama hors ligne & Scénario FastAPI & Retour d'un scénario de secours prédéfini & OK \\
\hline
Ollama hors ligne & QCM FastAPI & Questions de fallback prédéfinies par catégorie & OK \\
\hline
CodeLlama hors ligne & Code Challenge & \texttt{starter\_code} générique de secours & OK \\

\hline
PDF scanné & CV Parser hard & Warning frontend et analyse partielle sans CV & OK \\
\hline
JSON malformé Ollama & Tous modules & Extraction du premier objet JSON valide et fallback sécurisé & OK \\
\hline
\end{tabularx}
\caption{Tests de robustesse et mécanismes de fallback}
\label{tab:tests-fallbacks}
\end{table}

\textbf{Tests de scoring multi-sources}
La cohérence de l'algorithme de nivellement a été vérifiée en construisant des jeux de données synthétiques couvrant les quatre niveaux de maîtrise possibles, de Beginner à Expert, avec des combinaisons variables de sources. Ces tests permettent de s'assurer que la logique de triangulation multi-sources produit des scores stables et conformes aux seuils définis dans \texttt{scoring\_engine.py}. Le tableau ci-dessous présente les cas de validation retenus et les résultats obtenus.
\begin{table}[H]
\centering
\small
\renewcommand{\arraystretch}{1.25}
\begin{tabularx}{\TableWidth}{|Y|Y|P{1.8cm}|}
\hline
\rowcolor{TableHeaderBlue}
\color{white}\textbf{Cas} & \color{white}\textbf{Résultat attendu} & \color{white}\textbf{Résultat obtenu} \\
\hline
Skill présente dans GitHub, CV et portfolio & Expert avec score supérieur ou égal à 90 & OK \\
\hline
Skill présente dans GitHub et CV uniquement & Advanced avec score entre 70 et 89 & OK \\
\hline
Skill présente uniquement dans le CV avec une mention & Beginner avec score proche de 25 & OK \\
\hline
Job match Python Developer avec 8 compétences sur 14 & Score proche de 68/100 & OK \\
\hline
\end{tabularx}
\caption{Validation du scoring multi-sources}
\label{tab:tests-scoring-multisources}
\end{table}


\subsubsection{Déploiement}

La plateforme est déployée via \textit{Docker Compose} en six services conteneurisés et interconnectés. Le tableau ci-dessous détaille la configuration de chaque service.

\begin{table}[H]
\centering
\scriptsize
\renewcommand{\arraystretch}{1.2}
\begin{tabularx}{\TableWidth}{|P{2.4cm}|P{2.5cm}|P{1.2cm}|Y|}
\hline
\rowcolor{TableHeaderBlue}
\color{white}\textbf{Service} & \color{white}\textbf{Image / Build} & \color{white}\textbf{Port exposé} & \color{white}\textbf{Variables d'environnement clés} \\
\hline
n8n & \texttt{n8nio/n8n} & 5678 & Volume persistant \texttt{n8n\_data} \\
\hline
Ollama & \texttt{ollama/ollama} & 11434 & Volume persistant \texttt{ollama\_data} \\
\hline
\texttt{talentpredict-}\newline\texttt{ai} (FastAPI) & \texttt{./talentpredict-}\newline\texttt{ai} & 8000 & \texttt{ollama}, \texttt{GITHUB\_TOKEN}, \texttt{llama3.2}, \texttt{CODE\_CHALLENGE\_MODEL=codellama} \\
\hline
\texttt{backend} (Spring Boot) & \texttt{./BackEnd} & 8080 & \texttt{SPRING\_DATASOURCE\_URL}, \texttt{TALENTPREDICT\_AI\_BASE\_URL}, timeout 120 s \\
\hline
\texttt{postgres} & \texttt{postgres:15} & 5432 & \texttt{POSTGRES\_DB=talentpredict}, volume persistant \\
\hline
\texttt{frontend}\newline (Angular) & \texttt{./FrontEnd} & 4200 & --- \\
\hline
\end{tabularx}
\caption{Configuration des services Docker Compose}
\label{tab:docker-compose-services}
\end{table}

Les services communiquent via le réseau Docker interne, ce qui évite toute exposition directe des composants sensibles (PostgreSQL, Ollama) sur l'interface publique. Le timeout configurable de 120 secondes entre Spring Boot et FastAPI absorbe la latence inhérente aux appels CodeLlama les plus lourds.


%\begin{table}[H]
%\centering
%%\label{tab:urls-acces}
%\small
%\renewcommand{\arraystretch}{1.25}
%\begin{tabularx}{\TableWidth}{|P{3.4cm}|Y|}
%\hline
%\textbf{Service} & \textbf{URL} \\
%\hline
%n8n (workflows visuels) & \url{http://localhost:5678} %\\
%\hline
%FastAPI (Swagger UI) & \url{http://localhost:8000/docs} \\
%\hline
%Spring Boot (Swagger UI) & \url{http://localhost:8080/swagger-ui.html} \\
%\hline
%Frontend Angular & \url{http://localhost:4200} \\
%\hline
%Ollama & \url{http://localhost:11434} \\
%\hline
%PostgreSQL & \texttt{localhost:5432} \\
%\hline
%\end{tabularx}
%
%\caption{URLs d'accès après déploiement}
%\end{table}


\section{Suivi et Pilotage du Sprint 2}

\subsection{Scrum Board}

Pour garantir une gestion transparente, ce sprint de trois semaines, estimé à 37 Story Points, a été piloté via un Scrum Board. Ce management visuel a permis à l'équipe d'identifier immédiatement les goulots d'étranglement et de fluidifier le passage des tâches liées à l'IA vers la phase d'intégration finale.
\begin{figure}[H]
  \centering
  \includegraphics[width=0.9\textwidth]{images/Scrum Board.png}
  \caption{État d'avancement du Scrum Board  Sprint 2}
  \label{fig:scrum-board-sprint2}
\end{figure}
%\figureplaceholder{État d'avancement du Scrum Board --- Sprint 2}{fig:scrum-board-sprint2}

\subsection{Le Burndown Chart}

Afin d'évaluer notre vélocité, le Burndown Chart a été mis à jour quotidiennement. Le graphique révèle un léger ralentissement entre le Jour 7 et le Jour 10, dû aux défis d'intégration du microservice IA, notamment l'instabilité des réponses JSON. La résolution de ces obstacles a provoqué une franche accélération, permettant d'achever le sprint dans les délais impartis.
\begin{figure}[H]
  \centering
  \includegraphics[width=0.9\textwidth]{images/Burndown Chart2.png}
  \caption{Burndown Chart : vélocité de l'équipe face aux défis techniques  Sprint 2}
  \label{fig:burndown-chart-sprint2}
\end{figure}
%\figureplaceholder{Burndown Chart : vélocité de l'équipe face aux défis techniques --- Sprint 2}{fig:burndown-chart-sprint2}
\subsection{Indicateurs de Performance (KPIs) du Sprint 2}

Afin de valider cette itération dédiée à l'évaluation multi-source (\textit{compétences techniques et comportementales}), nous avons mesuré l'efficacité de nos pipelines de collecte et de nos modèles d'IA.
\paragraph*{1. Pilotage et Agilité}
Malgré la complexité de l'intégration entre l'orchestrateur n8n et le microservice FastAPI, l'équipe a maintenu un rythme constant. La totalité des 37 Story Points planifiés a été livrée, sans aucun report (Carry-over), validant ainsi la précision de nos estimations initiales.
\paragraph*{2. Métriques Techniques}
Le tableau ci-dessous résume les performances clés atteintes à la fin de ce sprint, confirmant la viabilité et la rapidité de l'architecture d'évaluation.

\begin{table}[H]
    \centering
    \renewcommand{\arraystretch}{1.5}
    \small
    \begin{tabularx}{0.95\textwidth}{|X|c|c|c|}
        \hline
        \rowcolor{TableHeaderBlue}
        \color{white}\textbf{Indicateur} & \color{white}\textbf{Cible} & \color{white}\textbf{Résultat Obtenu} & \color{white}\textbf{État} \\
        \hline
        \textit{Story Points} complétés & 37 & 37 & Validé \\
        \hline
        Temps d'exécution workflow \texttt{n8n} & $< 10$ s & $\sim 6.2$ s & Validé \\
        \hline
        Précision de l'extraction CV (150+ technos) & $> 85\%$ & $91\%$ & Validé \\
        \hline
        Temps de réponse IA (Évaluation QCM/Code) & $< 8$ s & $\sim 5.5$ s & Validé \\
        \hline
        Taux de succès des tests d'intégration & $> 90\%$ & $100\%$ & Validé \\
        \hline
    \end{tabularx}
    \caption{Indicateurs de Performance du Sprint 2}
    \label{tab:kpi-tech-s2}
\end{table}

\paragraph*{3. Bilan de l'itération}
Le bilan du Sprint 2 est hautement satisfaisant. La triangulation des données (CV, GitHub, LinkedIn) via les workflows n8n couplée à l'analyse sémantique du modèle LLM s'est avérée robuste. Les temps de réponse obtenus respectent largement les contraintes métiers, garantissant une expérience employé fluide lors des évaluations.

\section{ Conclusion}

Dans ce chapitre, nous avons présenté la conception et la réalisation du moteur d'évaluation multi-sources, pilier diagnostique de TalentPredict. Grâce à une modélisation UML rigoureuse, nous avons structuré l'orchestration complexe des données (CV, GitHub, PCM) vers l'IA locale (Ollama). Le développement d'interfaces Angular interactives, couplées à une architecture backend hybride (Spring Boot, FastAPI, n8n), a permis de concrétiser l'enjeu du sprint : objectiver avec précision les compétences techniques et comportementales des collaborateurs.
Par ailleurs, l'intégration du Communication Hub (SSE) a garanti une expérience utilisateur fluide malgré les délais d'inférence du modèle LLM. L'application stricte de la méthode Agile Scrum a prouvé la fiabilité de ce pipeline d'évaluation, aboutissant à un calcul pertinent du score de \textit{Job Matching}.
Avec la validation de ce deuxième sprint, la plateforme dispose désormais d'une cartographie exhaustive et certifiée des talents. Le chapitre suivant s’attachera à présenter le troisième sprint, dédié à l'exploitation de ces résultats à travers le module d'Upskilling et de gouvernance RH.

\chapter{Sprint 3 : Moteur d'IA et Gestion de l'Upskilling}

\section*{Introduction}
\phantomsection
\addcontentsline{toc}{section}{Introduction}
\phantomsection 
Suite à la mise en place des mécanismes d’évaluation technique et comportementale lors du sprint précédent, cette troisième itération vise à transformer les résultats d’analyse en un véritable dispositif de développement des compétences. L’identification des écarts de compétences n’a de valeur que si elle est accompagnée de recommandations et d’actions de remédiation adaptées.

Ce chapitre traite de la conception et de l’implémentation du Moteur d’Upskilling Intelligent, intégré au Module de Gouvernance RH. Ce composant constitue le cœur d’apprentissage de la plateforme TalentPredict et permet de structurer un cycle complet de montée en compétences, allant de la recommandation de parcours de formation basée sur l’intelligence artificielle jusqu’à la certification des acquis, tout en assurant une supervision centralisée par les administrateurs RH.
\section{Objectifs du Sprint 3}

L’objectif principal de ce sprint est de livrer un module fonctionnel dédié à l’Upskilling et à la Gouvernance RH. Il s’agit de faire évoluer TalentPredict d’un simple outil d’évaluation vers une plateforme complète de développement des compétences, capable de recommander, suivre et certifier la montée en compétence des collaborateurs de manière automatisée et intelligente.


\textbf{1. Objectifs fonctionnels et interactifs:}
\begin{itemize}
    \item \textbf{Upskilling Personnalisé :} Déploiement d'un moteur de recommandation de formations ciblé, s'appuyant sur l'analyse précise des écarts de compétences (\textit{Skill Gaps}).
    \item \textbf{Tracking \& Progression :} Mise en œuvre d’un Dashboard
    global incluant une Roadmap dynamique et un tableau \textit{Kanban} interactif pour le suivi des apprentissages des utilisateurs.
    \item \textbf{Validation \& Gamification :} Intégration d’un module de mini-tests générés par l’IA et d’un Leaderboard pour stimuler l'engagement.
    \clearpage
    \item \textbf{Espace Gouvernance RH :} Conception d'une interface d'administration complète permettant aux responsables de gérer (approuver, suggérer ou refuser) les parcours de formation.
\end{itemize}
 \textbf{2. Objectifs techniques et opérationnels  :}
\begin{itemize}
    \item \textbf{Intégration IA :} Optimisation des flux d'inférence via le microservice Python pour assurer une génération fluide et structurée des contenus pédagogiques.
    \item \textbf{Communication Synchrone :} Implémentation d'un système de notifications en temps réel pour garantir un suivi réactif des flux de validation \textit{workflows} .
    \item \textbf{Reporting \& Certification :} Développement d'un moteur permettant la génération dynamique et automatisée de certificats au format PDF à la volée.
    \item \textbf{Sécurité \& Workflow :} Renforcement de la sécurité via un modèle de contrôle d'accès basé sur les rôles (\textbf{RBAC}) , sécurisant ainsi les processus de décision administrative.
    \end{itemize}

\paragraph{Critères d’acceptation (Definition of Done)}

La validation des modules d'Upskilling et de Gouvernance repose sur les critères suivants :

\begin{itemize}
    \item \textbf{Cohérence IA} : Les recommandations de formation et les mini-tests générés par le microservice Python doivent être cohérents avec le profil de l'utilisateur.
    \item \textbf{Réactivité (SSE)} : Les notifications de validation RH doivent être reçues instantanément par l'utilisateur sans latence perceptible.
    \item \textbf{Conformité PDF} : Les certificats générés doivent être conformes au modèle visuel et inclure les métadonnées de l'utilisateur sans erreur de rendu.
    \item \textbf{Sécurité RBAC} : Les routes d'administration (approbation/refus) doivent être inaccessibles aux utilisateurs ne possédant pas le rôle \textit{ADMIN}.
    \item \textbf{Intégrité Workflow} : Le passage d'une formation dans le tableau Kanban doit mettre à jour l'état de la base de données en temps réel.
\end{itemize}



\section{Séance du Sprint planning }

La séance de \textit{Sprint Planning} a réuni l’équipe de développement, composée de Mohamed Ben Abdeladhim et Rahma Barouni, ainsi que le Product Owner Mr Med Amine Gharbi, et notre encadrante pédagogique  Mme Ben Aïcha Takwa. Cette réunion a permis de définir le périmètre fonctionnel et technique de ce troisième sprint.

Les points clés suivants ont été déterminés lors de cette séance :

\begin{itemize}
    \item \textbf{Durée :} La durée de cette itération a été fixée à un mois, un délai nécessaire pour stabiliser les algorithmes de recommandation et assurer l'intégration fluide entre le backend Spring Boot et le microservice d'IA.
    
    \item \textbf{Priorités Absolues :}
    \begin{itemize}
        \item Mise en place du \textbf{Tableau de Bord d'Upskilling} pour les utilisateurs et l'espace de Gouvernance RH.
        \item Développement du moteur de Roadmap IA et du catalogue de formations.
        \item Implémentation du module de Mini-Tests et du système de Génération Automatisée des Certificats.
        \item Gestion du cycle de vie des formations (Approbation, Suggestion, Refus par les RH).
    \end{itemize}
    
    \item \textbf{Priorités Secondaires :}
    \begin{itemize}
        \item Intégration des mécaniques de \textbf{Gamification} (Leaderboard et Badges).
        \item Optimisation des performances du moteur d'IA pour la génération de questions.
        \item Mise à jour du pipeline CI/CD via Docker et GitLab pour inclure les nouveaux endpoints.
    \end{itemize}
    
    \item \textbf{Prérequis :} La stabilité des API livrées lors des sprints précédents est essentielle, notamment le module de diagnostic technique et l'évaluation comportementale , qui servent de données d'entrée au moteur de recommandation.
\end{itemize}

À l’issue de cette réunion, les besoins ont été décomposés en User Stories (US) et estimés en jours-homme. Une priorité élevée a été accordée aux livrables essentiels afin de garantir la continuité entre l’évaluation des talents et leur développement en compétences.
\section{Backlog du Sprint 3}
Le Sprint Backlog représente le plan d’action opérationnel de cette troisième itération. Il regroupe les User Stories (US) issues du Product Backlog afin de répondre aux objectifs liés à l’Upskilling définis précédemment.
Pour assurer une planification rigoureuse, chaque tâche est structurée par Epic (macro-fonctionnalités) et estimée en Story Points selon la suite de Fibonacci. Cette approche permet d’évaluer l’effort de développement tout en garantissant une vélocité cohérente et maîtrisée pour l’équipe.
\begin{table}[H]
    \centering
    \renewcommand{\arraystretch}{1.5}
    \small
    \begin{tabularx}{\textwidth}{|c|>{\raggedright\arraybackslash}p{2.5cm}|>{\raggedright\arraybackslash}X|c|>{\centering\arraybackslash}p{2cm}|}
        \hline
        \rowcolor[HTML]{EFEFEF} 
        \rowcolor{TableHeaderBlue}
        \color{white}\textbf{ID} & \color{white}\textbf{Module / US} & \color{white}\textbf{Tâches de Développement} & \color{white}\textbf{JH} & \color{white}\textbf{Priorité} \\
        \hline
        
        US 3.1 & Interface \textit{Upskilling} & Création du \textit{Dashboard} Global et du tableau \textit{Kanban} interactif. & 3 & Haute \\
        \hline
        
        US 3.2 & Moteur de \textit{Roadmap} & Intégration de l'API IA (\texttt{Python}) pour la génération de parcours basés sur le \textit{Skill Gap}. & 4 & \textbf{Critique} \\
        \hline
        
        US 3.3 & Catalogue Formations & Développement du service de recommandation et de la base de données des cours. & 2 & Haute \\
        \hline
        
        US 3.4 & Mini-Tests IA & Création du moteur de génération de QCM et du service de correction automatique. & 4 & Critique \\
        \hline
        
        US 3.5 & \textit{Gamification} & Implémentation du système de points, badges et du \textit{Leaderboard} temps réel. & 2 & Moyenne \\
        \hline
        
        US 3.6 & Gouvernance RH & Développement du workflow d'approbation (Approuver/Suggérer/Refuser) et espace Admin. & 3 & Haute \\
        \hline
        
        US 3.7 & Certification & Mise en place du moteur PDFBox pour la génération dynamique des certificats. & 2 & Haute \\
        \hline
        
        TS 3.8 & Assurance Qualité & Tests unitaires, tests d'intégration (\texttt{Spring}/\texttt{Python}) et tests de montée en charge. & 2 & \textbf{Critique} \\
        \hline
        
    \end{tabularx}
    \vspace{2mm}
    \caption{Backlog du Sprint 3}
    \label{tab:planification_sprint3}
\end{table}

\section{Démarrage du Sprint 3}
Le début du Sprint 3 a été marqué par une phase d’initialisation technique visant à préparer les développements liés à l’Upskilling. L’équipe s’est concentrée sur la mise en place des canaux de communication entre le backend Spring Boot et le microservice Python, ainsi que sur l’intégration des services de notifications en temps réel.

Cette étape a permis d’établir les fondations nécessaires à l’implémentation des algorithmes de recommandation et du système de gouvernance RH, assurant une transition fluide entre la phase de planification et la phase de réalisation des modules.
\subsubsection*{1. Activités réalisées}

Durant cette itération, l'équipe a mené de front des développements \textit{backend} complexes et des intégrations d'intelligence artificielle :

\begin{itemize}
    \item \textbf{Ingénierie des Prompts :} Conception de \textit{prompts} système pour forcer l'IA à générer des \textit{roadmaps} et des tests au format JSON strict.
    \item \textbf{Développement Backend (\texttt{Spring Boot}) :} Implémentation des services métier pour la gestion du cycle de vie des formations et le calcul du \textit{Readiness Score}.
    \item \textbf{Développement Frontend (\texttt{Angular}) :} Création de composants dynamiques pour le tableau \textit{Kanban} et le module interactif de mini-tests.
    \item \textbf{Mise en place de la Gouvernance :} Développement de l'espace HR avec des \textit{dashboards} de suivi et des actions de validation (Approbation/Refus).
    \item \textbf{Intégration du moteur PDF :} Configuration du service de génération de certificats à la volée.
\end{itemize}

\subsubsection*{2. Livrables produits}

Les livrables suivants ont été validés et intégrés à la branche principale du projet :

\begin{itemize}
    \item \textbf{Module "\textit{Upskilling Engine}" :} Moteur fonctionnel de recommandation et de génération de \textit{roadmaps}.
    \item \textbf{Module "\textit{Evaluation Validation}" :} Interface de passage de tests et correction automatique.
    \item \textbf{Module "\textit{HR Governance}" :} Console d'administration pour la supervision des talents.
    \item \textbf{Service de Certification :} Générateur de documents PDF officiels.
    \item \textbf{Documentation API :} Mise à jour de \texttt{Swagger} pour inclure les \textit{endpoints} de formation.
\end{itemize}

\subsubsection*{3. Difficultés rencontrées et traitement}

L'exécution de ce sprint a été marquée par trois défis techniques majeurs :

\begin{itemize}
    \item \textbf{Difficulté 1 : Instabilité des formats JSON de l’IA}
    \begin{itemize}
        \item \textbf{Problème :} Les réponses des modèles LLM présentaient parfois des structures incohérentes, compromettant le \textit{parsing} des \textit{roadmaps}.
        \item \textbf{Traitement :} Renforcement du \textit{system prompting} et intégration de validateurs de schémas rigoureux côté \textit{backend} pour garantir l'intégrité des données.
    \end{itemize}
    \vspace{2mm}
    
    \item \textbf{Difficulté 2 : Conflits de dépendances PDF}
    \begin{itemize}
        \item \textbf{Problème :} L'intégration de la génération de certificats a révélé des incompatibilités entre les bibliothèques de rendu.
        \item \textbf{Traitement :} Isolation des dépendances critiques et migration vers \texttt{OpenHTMLtoPDF}, assurant ainsi une stabilité du moteur de génération et un rendu style fidèle.
    \end{itemize}
    \vspace{2mm}
    
    \item \textbf{Difficulté 3 : Latence des appels API externes}
    \begin{itemize}
        \item \textbf{Problème :} Les délais de génération des tests par l’IA impactaient la fluidité de l'interface.
        \item \textbf{Traitement :} Implémentation d'une communication asynchrone via les \textit{Server-Sent Events} (\texttt{SSE}), permettant d'informer l'utilisateur en temps réel de l'état d'avancement des processus en tâche de fond.
    \end{itemize}
\end{itemize}

\subsection{Analyse}

La phase d’analyse vise à définir le périmètre fonctionnel du module de Formation.
et identifier avec précision les interactions entre le système et les utilisateurs. Contrairement au
Le premier sprint du module introduit une grande automatisation grâce à Intelligence Artificielle.

\subsubsection*{1. Identification des Acteurs}

Le système interagit avec trois acteurs principaux, parmi lesquels se trouve un acteur automatisé :

\begin{itemize}
    \item \textbf{Utilisateur (Employé) :} Constitue le cœur du module. Il réalise les évaluations, examine ses compétences et suit sa progression d'apprentissage via la roadmap interactive. Il soumet ses preuves de réussite (mini-tests ou certificats) pour valider sa montée en compétences.
    
    \item \textbf{Administrateur (RH) :} garantit la politique de formation de l’entreprise.
Le superviseur de l’équipe dirige les membres. Le superviseur a le droit d’approuver ou de rejeter chaque demande. L’inscription aux cours est obligatoire.
    
    \item \textbf{Moteur IA (\enquote{System}) :} Agit de manière autonome en arrière-plan. Il est sollicité pour générer des recommandations ciblées suite à un \textit{Assessment} et pour créer dynamiquement les questions des mini-tests de validation.
\end{itemize}


\subsubsection*{2. Clarification du besoin}

La finalité de ce sprint est de répondre à la problématique de la remédiation intelligente. Le système vise à automatiser le cycle d’Upskilling en réduisant l’écart entre les compétences identifiées et les besoins du marché.
L’objectif est de proposer une solution de GPEC (Gestion Prévisionnelle des Emplois et des Compétences) dynamique, dans laquelle chaque recommandation repose sur une analyse de données, et chaque progression est validée par une preuve formelle telle qu’un certificat.

\subsubsection*{3. Scénarios de validation}

Pour garantir la robustesse fonctionnelle, trois parcours types ont été modélisés :

\begin{itemize}
    \item \textbf{Parcours Nominal (Succès) :} \\
    Identification du besoin $\rightarrow$ Requête de formation $\rightarrow$ Approbation RH $\rightarrow$ Suivi du module $\rightarrow$ Réussite au test $\rightarrow$ Certification.
    
    \item \textbf{Parcours Alternatif (Ajustement RH) :} \\
    Requête utilisateur $\rightarrow$ Évaluation de la pertinence par le RH $\rightarrow$ Recommandation d'une ressource alternative $\rightarrow$ Mise à jour de la \textit{Roadmap}.
    
    \item \textbf{Parcours d'Échec (Réévaluation) :} \\
    Tentative de validation infructueuse $\rightarrow$ Analyse des erreurs par le système $\rightarrow$ Prescription de révision $\rightarrow$ Réactivation du test après un délai de carence.
\end{itemize}

\subsubsection*{4. Contraintes techniques}

L'implémentation de ces processus repose sur le respect de contraintes strictes :

\begin{itemize}
    \item \textbf{Persistance et Intégrité :} Garantie de la cohérence des états de formation dans la base de données \texttt{PostgreSQL} lors des transitions du \textit{Kanban}.
    \item \textbf{Interfaçage Asynchrone :} Gestion des temps de réponse des modèles LLM via une architecture événementielle pour maintenir une interface fluide.
    \item \textbf{Sécurité et RBAC :} Implémentation de contrôles d'accès granulaires pour les \textit{endpoints} de gouvernance, limitant les actions sensibles aux seuls profils RH.
    \item \textbf{Standardisation de la Certification :} Respect des normes de rendu documentaire pour la génération de certificats PDF à haute fidélité visuelle.
\end{itemize}


\subsubsection*{5. Diagramme de Cas d’Utilisation }
Le troisième sprint occupe une place centrale en se focalisant sur la gestion de l'apprentissage et la remédiation. La modélisation fonctionnelle de cette itération est illustrée par la figure \ref{fig:usecase_sprint3}. Ce diagramme de cas d'utilisation met en évidence les interactions directes entre les différents acteurs, qu'ils soient humains (employés, administrateurs) ou automatisés (service IA, services de notification), et les processus clés du système.
\begin{figure}[H]
    \centering
    \includegraphics[width=0.95\textwidth]{images/usecaseSprint3.png}

    \label{fig:usecase_sprint3}
\end{figure}
Le tableau \ref{tab:relations_sprint3} suivant résume la sémantique visuelle et l’application concrète des différentes
relations modélisées dans notre diagramme pour ce troisième sprint.

\begin{table}[H]
    \centering
    \label{tab:relations_sprint3}
\begin{tabular}{p{3.5cm} p{4.5cm} p{7cm}}
    \toprule
    \textbf{Type de Relation} & \textbf{Représentation Graphique} & \textbf{Application Concrète} \\
    \midrule

    \textbf{Association} (Communication) & Ligne continue simple & Relie l'acteur \textbf{Employé} au cas \textit{Consulter le Tableau de Bord d'Upskilling}. Indique que cet utilisateur interagit directement avec cette fonctionnalité. \\ \addlinespace
    
    \textbf{Inclusion} (\texttt{<<include>>}) & Flèche en pointillés pointant vers la fonction requise & \textit{Effectuer une demande de formation} inclut obligatoirement \textit{Calculer le nouveau Readiness Score}. Le calcul est indispensable à l'exécution de la demande. \\ \addlinespace
    
    \textbf{Extension} (\texttt{<<extend>>}) & Flèche en pointillés pointant vers la fonction de base & \textit{Suggérer une formation alternative} étend \textit{Calculer le nouveau Readiness Score}. Le système peut optionnellement proposer une alternative lors du calcul du score. \\
    \bottomrule
\end{tabular}
\caption{Types de relations dans le diagramme de cas d'utilisation}

\end{table}

\subsubsection* {Cas 1 : Processus de Planification de l'Apprentissage :}
Ce sous-module se concentre spécifiquement sur le processus de planification de l'apprentissage. La figure \ref{fig:usecase_planification} détaille les cas d'utilisation liés à la gestion du parcours, de la génération d'une roadmap personnalisée via le Service IA jusqu'au suivi quotidien via le système Kanban.
\begin{figure}[H]
    \centering
    \includegraphics[width=0.95\textwidth]{images/cas1Sprint3.png}
    \caption{Diagramme de cas d'utilisation du processus de planification de l'apprentissage}
    \label{fig:usecase_planification}
\end{figure}
Le tableau \ref{tab:relations_planification} résume les différents types de relations UML identifiés dans ce module d'apprentissage, en illustrant chaque concept par un exemple tiré du diagramme.

\begin{table}[H]
    \centering
    \label{tab:relations_planification}
\begin{tabular}{p{3.5cm} p{4.5cm} p{7cm}}
    \toprule
    \textbf{Type de Relation} & \textbf{Représentation Graphique} & \textbf{Application Concrète} \\
    \midrule

    \textbf{Association} (Communication) & Ligne continue simple & Relie l'acteur \textbf{Employé} au cas \textit{Gérer son parcours d'apprentissage}. Cela indique que l'utilisateur est l'initiateur de la gestion de son parcours. \\ \addlinespace
    
    \textbf{Inclusion} (\texttt{<<include>>}) & Flèche en pointillés pointant vers la fonction requise & Le cas \textit{Générer une roadmap personnalisée} inclut obligatoirement le cas \textit{Spécifier le profil métier cible}. L'IA ne peut générer la roadmap sans connaître l'objectif du profil visé. \\ \addlinespace
    
    \textbf{Extension} (\texttt{<<extend>>}) & Flèche en pointillés pointant vers la fonction de base & Le cas \textit{Ajouter un cours} étend le cas \textit{Piloter l'apprentissage via Kanban}. Cela signifie que lors du pilotage de son apprentissage, l'employé a la possibilité (optionnelle) d'ajouter un nouveau cours. \\
    \bottomrule
\end{tabular}
    \caption{Types de relations dans le module de planification de l'apprentissage}
\end{table}
\subsubsection*{Cas 2 : Processus de Gouvernance et Supervision RH :}
La supervision stratégique des compétences est assurée par le module de gouvernance RH, illustré par la figure \ref{fig:usecase_gouvernance}. Ce diagramme expose les fonctionnalités de haut niveau destinées à l'Administrateur RH, notamment la gestion des demandes de formation et l'analyse de l'adéquation des profils, tout en précisant les systèmes externes sollicités comme le Service de Notification.
\begin{figure}[H]
    \centering
    \includegraphics[width=0.95\textwidth]{images/usecaseSprint3.png}
    \caption{Diagramme de cas d'utilisation du module de gouvernance et de supervision RH}
    \label{fig:usecase_gouvernance}
\end{figure}
Pour assurer la cohérence métier des actions de supervision, plusieurs types de relations UML lient les cas d'utilisation entre eux. Le tableau \ref{tab:relations_gouvernance} explicite ces liens, en se basant sur les processus de traitement des demandes et d'analyse de montée en compétences.
\begin{table}[H]
    \centering
    \label{tab:relations_gouvernance}
\begin{tabular}{p{3.5cm} p{4.5cm} p{7cm}}
    \toprule
    \textbf{Type de Relation} & \textbf{Représentation Graphique} & \textbf{Application Concrète} \\
    \midrule

    \textbf{Association} (Communication) & Ligne continue simple & Relie l'acteur \textbf{Administrateur RH} au cas \textit{Suivre les indicateurs d'apprentissage de l'équipe}. Indique le droit de regard de l'acteur sur ces données. \\ \addlinespace
    
    \textbf{Inclusion} (\texttt{<<include>>}) & Flèche en pointillés pointant vers la fonction requise & Le cas \textit{Refuser la demande} inclut obligatoirement \textit{Saisir un motif de refus}. Cette étape est nécessaire pour finaliser le rejet d'une demande. \\ \addlinespace
    
    \textbf{Extension} (\texttt{<<extend>>}) & Flèche en pointillés pointant vers la fonction de base & Les cas \textit{Approuver la demande} ou \textit{Refuser la demande} étendent le cas générique \textit{Traiter les demandes de formation}. Le traitement est le cas de base qui se spécialise selon la décision prise. \\
    \bottomrule
\end{tabular}
    \caption{Types de relations dans le module de gouvernance RH}
\end{table}
\subsubsection*{Cas 3 : Processus d'Évaluation et de Certification}
Ce dernier sous-module se concentre sur l'aboutissement du parcours d'apprentissage de l'employé : l'évaluation finale et la certification. La figure \ref{fig:usecase_eval_certif} détaille les cas d'utilisation relatifs au passage du test de validation, à la préparation technique de cette évaluation par le Service IA, ainsi qu'aux actions d'audit menées par l'Administrateur RH pour statuer sur la réussite de l'apprenant.
\begin{figure}[H]
    \centering
    \includegraphics[width=0.95\textwidth]{images/case3Sprint3.png}
    \caption{Diagramme de cas d'utilisation : processus d'évaluation et de certification}
    \label{fig:usecase_eval_certif}
\end{figure}
Le tableau \ref{tab:relations_eval_certif} présente les différents types de relations UML structurant ce processus d'évaluation. Il met notamment en évidence les dépendances d'inclusion pour la génération automatique des certificats, ainsi que les relations d'extension qui traduisent les décisions conditionnelles de l'Administrateur RH (approbation de la certification ou exigence d'une réévaluation).
\begin{table}[H]
    \centering
    \label{tab:relations_eval_certif}
\begin{tabular}{p{3.5cm} p{4.5cm} p{7cm}}
    \toprule
    \textbf{Type de Relation} & \textbf{Représentation Graphique} & \textbf{Application Concrète} \\
    \midrule

    \textbf{Association} (Communication) & Ligne continue simple & Relie l'acteur \textbf{Employé} au cas \textit{Passer l'évaluation de validation}. L'acteur externe \textbf{Service IA} est, quant à lui, associé à \textit{Préparer l'évaluation de validation}. \\ \addlinespace
    
    \textbf{Inclusion} (\texttt{<<include>>}) & Flèche en pointillés pointant vers la fonction requise & Le cas \textit{Approuver la certification} inclut obligatoirement \textit{Générer le certificat PDF}. La validation RH déclenche de facto la création du document certifiant. \\ \addlinespace
    
    \textbf{Extension} (\texttt{<<extend>>}) & Flèche en pointillés pointant vers la fonction de base & Les cas \textit{Approuver la certification} et \textit{Refuser et exiger une réévaluation} étendent le cas de base \textit{Auditer les résultats du test}. Selon le résultat de l'audit RH, l'une de ces actions est déclenchée. \\
    \bottomrule
\end{tabular}
 \caption{Types de relations dans le processus d'évaluation et de certification}
\end{table}

\subsection{Conception }

Cette phase a pour but de traduire les exigences fonctionnelles et les cas d'utilisation identifiés précédemment en solutions techniques concrètes. L'enjeu majeur de ce troisième sprint réside dans la conception d'un système capable d'orchestrer une intelligence artificielle générative tout en préservant la robustesse transactionnelle requise pour la certification et la gouvernance RH.

\paragraph*{1. Objectifs de conception}
\begin{itemize}
    \item \textbf{Interopérabilité :} Concevoir une communication fluide et tolérante aux pannes entre le cœur transactionnel (Java) et le service cognitif (Python).
    \item \textbf{Traçabilité :} Assurer que chaque décision de l'IA (\textit{Roadmap} générée) et chaque décision RH (Approbation/Refus) soient historisées en base de données.
    \item \textbf{Réactivité :} Mettre en place des mécanismes asynchrones (Notifications) pour que les décisions de gouvernance ne bloquent pas l'expérience utilisateur.
\end{itemize}

\subsubsection*{2. Architecture Technique}

\begin{itemize}
    \item \textbf{Client (\texttt{Angular}) :} Responsable du rendu interactif du tableau \textit{Kanban} (\textit{Drag \& Drop}), du chronomètre pour les mini-tests, et de la visualisation des recommandations.
    \item \textbf{\textit{Backend Core} (\texttt{Spring Boot}) :} Joue le rôle de chef d'orchestre. Il sécurise les accès via \texttt{Spring Security} (JWT + RBAC), calcule les \textit{Skill Gaps}, sauvegarde les scores, et utilise \texttt{PDFBox} pour la génération de certificats côté serveur.
    \item \textbf{Microservice IA (\texttt{Python}/\texttt{FastAPI}) :} Isolé du \textit{backend} principal, il prend en charge la construction des \textit{prompts} et la communication avec l'API du LLM en forçant des sorties structurées (JSON).
    \item \textbf{Couche de Communication (SSE) :} Utilisation des \textit{Server-Sent Events} pour "pousser" les notifications de l'Admin vers l'Utilisateur en temps réel, évitant ainsi de surcharger le réseau avec du \textit{Long Polling}.
\end{itemize}
Afin d’assurer une séparation claire entre la logique métier \texttt{Java} et les traitements d’intelligence artificielle, une architecture microservices basée sur \texttt{FastAPI} et \texttt{Ollama} a été mise en place. 

Le diagramme de la Figure \ref{fig:archi_ia_fastapi} présente l’architecture détaillée du \textit{pipeline} d’IA, de la sollicitation par l’utilisateur jusqu’à la restitution des prédictions. Ce flux a été conçu pour maximiser la réactivité et la fiabilité du système.
\begin{figure}[H]
    \centering
    % Remplacez 'chemin_architecture_ia_fastapi.png' par le nom exact de votre image
     \includegraphics[width=0.75\textwidth]{images/ArchitectureFastAPI.png}
    \vspace{2mm}
    \caption{Architecture interne du microservice IA FastAPI}
    \label{fig:archi_ia_fastapi}
\end{figure}



\paragraph*{A. Orchestration et Flux Temps Réel (SSE)}
Contrairement aux architectures classiques où l’utilisateur doit attendre la fin d’un calcul long, TalentPredict implémente un flux de \textit{Streaming} via \textit{Server-Sent Events} (\texttt{SSE}).
\begin{itemize}
    \item Lorsqu'une analyse est lancée, le \textit{Backend} \texttt{Spring Boot} établit une connexion persistante avec le \textit{Frontend} \texttt{Angular}.
    \item Dès que les premiers \textit{tokens} sont générés par le modèle \texttt{Llama 3.2 8B}, ils sont transmis progressivement à l'interface, offrant une expérience utilisateur fluide et dynamique.
\end{itemize}

\paragraph*{B. Microservice FastAPI : Performance et Validation}
Le microservice \texttt{Python} agit comme un filtre d’intelligence et de validation :
\begin{itemize}
    \item \textbf{Validation Stricte (\texttt{Pydantic}) :} Chaque donnée entrante est validée par des schémas \texttt{Pydantic} pour éviter les erreurs de traitement. De même, la réponse de l'IA est rigoureusement vérifiée pour s'assurer qu'elle respecte le format \texttt{JSON} attendu par le \textit{backend} \texttt{Java}.
    \item \textbf{Traitement Asynchrone :} Grâce au moteur \texttt{AsyncIO} de \texttt{FastAPI}, les appels vers \texttt{Ollama} sont non-bloquants, permettant de gérer des files d'attente d'analyses sans dégrader les performances globales.
\end{itemize}

\paragraph*{C. Ingénierie de Prompt et Robustesse}
La "logique métier" de l'IA est encapsulée dans le \textit{Prompt Engine} :
\begin{itemize}
    \item \textbf{Contraintes de Système :} Le modèle est contraint par des \textit{System Prompts} qui définissent son rôle d'expert RH et lui imposent des règles de sortie déterministes.
    \item \textbf{Gestion des Erreurs (\textit{Retry Logic}) :} Le \textit{pipeline} intègre un module de gestion des exceptions. En cas de latence excessive du moteur \texttt{Ollama} ou de réponse malformée, le système déclenche automatiquement une logique de ré-essai (\textit{Retry}), garantissant la continuité du service.
\end{itemize}

\paragraph*{D. Sécurité et Inférence Locale}
La sécurité est assurée à deux niveaux :
\begin{itemize}
    \item \textbf{Authentification Inter-Services :} Les échanges entre \texttt{Spring Boot} et \texttt{FastAPI} sont sécurisés par des jetons \texttt{JWT} internes, empêchant toute sollicitation non autorisée du moteur d'IA.
    \item \textbf{Souveraineté des Données :} L'utilisation d'\texttt{Ollama} permet une inférence locale optimisée. Les données sensibles des collaborateurs restent confinées dans le réseau privé \texttt{Docker}, garantissant une conformité totale avec le RGPD.
\end{itemize}
\subsubsection*{3. Spécification des API du Microservice IA}

Le microservice IA expose plusieurs \textit{endpoints} REST permettant l’interaction avec le moteur de génération :

\begin{itemize}
    \item \textbf{\texttt{GET /api/formations/recommendations}} : Liste les formations suggérées en fonction des lacunes identifiées.
    \item \textbf{\texttt{GET /api/predictions/roadmap}} : Récupère la trajectoire de carrière générée par l'IA.
    \item \textbf{\texttt{POST /api/predictions/generate}} : Déclenche une nouvelle analyse du \textit{Gap} par le moteur \texttt{Llama 3}.
\end{itemize}

Chaque \textit{endpoint} retourne une réponse strictement formatée en \texttt{JSON} afin d’assurer la compatibilité avec le \textit{backend} \texttt{Spring Boot} et le \textit{frontend} \texttt{Angular}.
\subsubsection*{4. Validation et Robustesse des Réponses IA}

L'un des défis majeurs lors de l'intégration de modèles de langage (LLM) comme \texttt{Llama 3} est de transformer une réponse textuelle probabiliste en une donnée structurée et fiable pour une application d'entreprise. Pour garantir cette robustesse, TalentPredict implémente une stratégie de validation à plusieurs niveaux.

\paragraph*{A. Le Forçage de Structure (\textit{JSON Mode} \& Pydantic)}
Pour éviter les "hallucinations" de format, nous utilisons des schémas \texttt{Pydantic} qui définissent contractuellement la structure attendue (champs obligatoires, types de données, énumérations).
\begin{itemize}
    \item Le microservice \texttt{FastAPI} impose au modèle \texttt{Llama 3} un \texttt{JSON Schema} strict.
    \item Si la réponse de l'IA ne respecte pas ce schéma (ex: champ manquant ou type erroné), la donnée est rejetée avant même d'atteindre le \textit{backend} \texttt{Java}, évitant ainsi des erreurs de désérialisation en cascade.
\end{itemize}

\paragraph*{B. Déterminisme et \textit{Prompt Engineering}}
La robustesse passe également par la réduction de la "créativité" non désirée du modèle :
\begin{itemize}
    \item \textbf{Température basse :} Nous avons configuré la température du modèle à une valeur proche de 0 pour favoriser des réponses constantes et déterministes.
    \item \textbf{\textit{System Prompting} :} Des instructions système explicites interdisent au modèle d'ajouter des commentaires textuels (ex: "Voici votre analyse :") en dehors du bloc \texttt{JSON}, garantissant une sortie "pure" et directement \textit{parsable}.
\end{itemize}

\paragraph*{C. Mécanismes de Résilience (\textit{Retry \& Fallback})}
Malgré ces précautions, les modèles d'IA peuvent parfois échouer ou subir des latences. Nous avons donc mis en place :
\begin{itemize}
    \item \textbf{Stratégie de Re-essai (\textit{Retry Logic}) :} En cas de réponse invalide, le système effectue automatiquement jusqu'à trois tentatives avec un ajustement léger du \textit{prompt} pour corriger l'erreur.
    \item \textbf{Gestion des \textit{Timeouts} :} Pour garantir la disponibilité de l'application, un délai d'attente maximum est défini. Si l'IA ne répond pas dans le temps imparti, le système bascule sur un mode dégradé (affichage des dernières données connues) et alerte l'administrateur.
\end{itemize}

\paragraph*{D. Validation Métier (\textit{Sanity Checks})}
Une fois le \texttt{JSON} extrait, le \textit{backend} \texttt{Java} effectue des vérifications métier finales :
\begin{itemize}
    \item Validation de la cohérence des scores (de 0 à 100\%).
    \item Vérification de l'existence des compétences référencées. Cela assure que l'intelligence artificielle reste un outil d'aide à la décision fiable et non une source d'erreurs logiques.
\end{itemize}

\subsubsection*{5. Gestion Globale des Erreurs du Système IA}

Le système implémente une gestion centralisée des erreurs afin de garantir la stabilité globale de l’application. Les principaux scénarios couverts sont :

\begin{itemize}
    \item \textbf{\textit{Timeout} des requêtes IA} : Dépassement du temps d'attente autorisé lors de la génération.
    \item \textbf{Réponses \texttt{JSON} invalides} : Échec de la validation structurelle par les schémas \texttt{Pydantic}.
    \item \textbf{Indisponibilité du modèle \texttt{LLM}} : Déconnexion ou surcharge du moteur \texttt{Ollama}.
    \item \textbf{Erreurs réseau} : Problèmes de communication entre les conteneurs \texttt{Spring Boot} et \texttt{FastAPI}.
\end{itemize}

Dans ces cas, le système applique une stratégie de \textit{fallback} permettant de retourner une réponse contrôlée ou les dernières données valides disponibles, évitant ainsi un blocage de l'interface utilisateur.

\subsubsection*{6. Modèle de données : }
La figure \ref{fig:erd_talentpredict} présente le diagramme entité-association (ERD) détaillant l'architecture de notre base de données.

\begin{figure}[H]
    \centering
    % Remplacer 'images/erd.png' par le chemin exact de votre image
    \includegraphics[width=1\textwidth]{images/ERD3.png}
    \caption{Diagramme Entité-Association (ERD) du système}
    \label{fig:erd_talentpredict}
\end{figure}

Ce schéma conceptuel met en évidence la centralité des profils utilisateurs et la manière dont les différentes entités s'articulent pour gérer l'évaluation, l'apprentissage et les recommandations. Les tables constituant cette architecture sont décrites ci-dessous :

\begin{description}
    \item[\textbf{USERS} :] C'est la table centrale du système qui gère les comptes. Sa clé primaire (PK) est \texttt{id}. Elle ne contient pas de clé étrangère puisqu'elle est la référence principale. Elle utilise un \texttt{ENUM} pour le champ \texttt{role} avec les valeurs \texttt{USER} ou \texttt{ADMIN}.
    
    \item[\textbf{FORMATIONS} :] Cette table contient le catalogue des cours ou modules d'apprentissage. Sa PK est \texttt{id}. Elle ne possède ni clé étrangère (FK) ni d'énumération spécifique, se concentrant sur les détails descriptifs de la formation (titre, fournisseur, durée, url).
    
    \item[\textbf{PROFILES} :] Elle stocke les informations professionnelles détaillées des utilisateurs, comme les liens vers Github ou LinkedIn et un résumé généré par l'IA. La PK est \texttt{id} et elle contient une FK \texttt{user\_id} qui la relie directement à la table USERS.
    
    \item[\textbf{USER\_FORMATIONS} :] Il s'agit d'une table de jointure qui permet de suivre la progression d'un utilisateur dans un cours spécifique. Elle utilise des PK/FK composées : \texttt{user\_id} (référençant USERS) et \texttt{formation\_id} (référençant FORMATIONS). Elle intègre un \texttt{ENUM} sur le champ \texttt{status} avec les valeurs \texttt{EN COURS} ou \texttt{TERMINE}.
    
    \item[\textbf{USER\_NOTIFICATIONS} :] Cette table gère l'historique et l'état de lecture des alertes envoyées. Sa PK est \texttt{id} et elle utilise la FK \texttt{user\_id} pour associer la notification au bon utilisateur.
    
    \item[\textbf{LEARNING\_PATHS} :] Elle définit des parcours d'apprentissage personnalisés ou suggérés. La PK est \texttt{id} et la FK est \texttt{user\_id} pour lier le parcours à son propriétaire.
    
    \item[\textbf{RECOMMENDATIONS} :] Cette table enregistre les suggestions faites au système pour un utilisateur. La PK est \texttt{id} et la FK est \texttt{user\_id}. Le champ \texttt{type} agit comme un \texttt{ENUM} implicite ou explicite, limitant les choix à des éléments comme \texttt{FORMATION} ou \texttt{POSTE}.
    
    \item[\textbf{SKILLS} :] Elle répertorie le profil de compétences de l'utilisateur. La PK est \texttt{id} et la FK est \texttt{user\_id}. Cette table est riche en contraintes de type \texttt{ENUM} : le type de compétence (\texttt{TECH}, \texttt{SOFT}), le \texttt{level} (1-5), et la source de l'information (\texttt{CV}, \texttt{GITHUB}, \texttt{IA}).
    
    \item[\textbf{PREDICTIONS} :] Cette table stocke les résultats d'analyses algorithmiques ou IA sur le profil de l'utilisateur. La PK est \texttt{id} et la FK est \texttt{user\_id}. Le champ \texttt{type} utilise un \texttt{ENUM} indiquant la nature de la prédiction, avec des valeurs comme \texttt{CHURN} ou \texttt{UPSKILLING}.
\end{description}
\subsubsection*{7. Architecture du moteur de recommandation}
La figure~\ref{fig:arch_moteur_rec} présente l'architecture du système de recommandation du Sprint 3, illustrant le flux de transformation des résultats d'évaluation en parcours de formation personnalisés.

\begin{figure}[htbp]
    \centering
    \includegraphics[width=\textwidth]{images/archmoteurREC.png}
    \caption{Architecture du processus de recommandation}
    \label{fig:arch_moteur_rec}
\end{figure}

Comme l'illustre la figure~\ref{fig:arch_moteur_rec}, le système reçoit en entrée les scores issus des évaluations techniques et comportementales réalisées lors du Sprint 2. Ces résultats sont consolidés dans le profil de l'employé sous la forme de forces et de faiblesses. Une analyse des lacunes est ensuite effectuée par le Service IA (s'appuyant sur le modèle Llama 3.2 via Ollama), qui génère deux types de recommandations distinctes : des formations techniques ciblées pour combler les déficits de compétences (\textit{Hard Skills}), et des formations comportementales pour renforcer les aptitudes interpersonnelles (\textit{Soft Skills}) identifiées comme insuffisantes.
\subsubsection*{8. Algorithme de Calcul du Readiness Score: }
Dans notre architecture, le \textit{Readiness Score} (Taux de Préparation) est 
l'indicateur central qui permet de quantifier l'adéquation technique d'un candidat 
par rapport aux exigences d'un poste cible. Afin de garantir un calcul robuste, 
transparent et vérifiable, nous avons implémenté l'algorithme suivant dans le 
backend Java (\texttt{CareerController.java}) :

\begin{equation}
R_{brut} = \left( \frac{\sum_{i=1}^{n} Niveau_i}{n \times 10} \right) \times 100
\end{equation}

\noindent Définition des variables :
\begin{itemize}
    \item $R_{brut}$ : le Readiness Score théorique exprimé en pourcentage.
    \item $Niveau_i$ : le niveau de maîtrise actuel de l'utilisateur pour la 
    compétence $i$, évalué sur une échelle de 1 à 10.
    \item $n$ : le nombre total de compétences évaluées dans la Skill Gap Analysis.
\end{itemize}

\noindent Pour éviter les valeurs aberrantes, le résultat brut est soumis à une 
fonction de bornage (\textit{clamping}) avant sa persistance en base de données :

\begin{equation}
Readiness\_Score = \max\!\left(10,\ \min\!\left(95,\ \text{arrondi}(R_{brut})\right)\right)
\end{equation}

\noindent Cette borne garantit qu'aucun collaborateur ne se voit attribuer un score 
décourageant de 0\% ni une perfection irréaliste de 100\%.


\subsubsection*{9. Contrats API (Exemple de Conception): }
Afin de garantir la séparation entre le \textit{Frontend} et le \textit{Backend}, des contrats d'interface stricts ont été définis. Voici un exemple illustrant la requête et la réponse pour la génération asynchrone d'un mini-test :

\textbf{Point d'accès :} \texttt{POST /api/assessment/mini-test/generate}

\textbf{Requête (\textit{Payload}) :}
\begin{lstlisting}[language=json, caption={Payload de la requête de génération de test}]
{
  "candidateId": "123e4567-e89b-12d3-a456-426614174000",
  "targetSkill": "Spring Boot",
  "currentLevel": 4,
  "difficulty": "INTERMEDIATE"
}
\end{lstlisting}

\textbf{Réponse (200 OK - \texttt{TestDto}) :}
\begin{lstlisting}[language=json, caption={Structure de réponse d'un mini-test généré}]
{
  "sessionId": "eval-9048-abcd",
  "timeLimitSeconds": 600,
  "totalQuestions": 5,
  "questions": [
    {
      "questionId": "q-101",
      "text": "Quelle annotation Spring Boot expose des endpoints REST ?",
      "options": [
        { "key": "A", "value": "@Controller" },
        { "key": "B", "value": "@RestController" },
        { "key": "C", "value": "@ApiComponent" },
        { "key": "D", "value": "@RestApi" }
      ]
    }
  ]
}

\end{lstlisting}


\subsubsection*{10. Diagramme de Séquence}

\subsubsection*{ Cas 1: Planification de l'Apprentissage}

Le processus technique de planification de l'apprentissage est illustré en détail par le diagramme de séquence de la figure~\ref{fig:planification_apprentissage}.

\begin{figure}[htbp]
    \centering
    \includegraphics[width=\textwidth]{images/seq1Sprint3.png}
    \caption{Diagramme de séquence : planification de l'apprentissage}
    \label{fig:planification_apprentissage}
\end{figure}

Comme le montre la figure~\ref{fig:planification_apprentissage}, ce mécanisme se structure autour de trois phases clés :
\begin{itemize}
    \item \textbf{Consultation du Dashboard :} L'employé accède à son espace de montée en compétences via le Frontend (Angular). Après la vérification et validation du jeton JWT par le filtre de sécurité, le contrôleur sollicite le service métier. Ce dernier extrait les compétences actuelles de l'utilisateur et évalue ses écarts directement depuis la base de données PostgreSQL.
    \item \textbf{Génération de la Roadmap (IA) :} Suite à la sélection d'un poste cible par l'employé, le service métier transmet les données au Service IA, qui interroge le Service LLM afin de concevoir un parcours de formation personnalisé. Après la validation de la conformité du format JSON reçu, la roadmap structurée est persistée en base de données puis retournée à l'interface pour affichage.
    \item \textbf{Initialisation du Kanban et Gamification :} L'utilisateur a la possibilité d'intégrer ce plan à son tableau de suivi Kanban. Les tâches y sont automatiquement initialisées avec le statut \texttt{TO\_DO}. De plus, s'il s'agit du premier cours ajouté, le système déclenche l'envoi d'une notification de badge débloqué pour valoriser l'engagement de l'apprenant.
\end{itemize}

\subsubsection*{Cas 2: Gouvernance et Supervision RH }

Le processus de validation et de suivi des demandes par les administrateurs est détaillé dans le diagramme de séquence de la figure~\ref{fig:gouvernance_supervision}.

\begin{figure}[H]
    \centering
    \includegraphics[width=\textwidth]{images/seq2Sprint3.png}
    \caption{Diagramme de séquence : gouvernance et supervision RH}
    \label{fig:gouvernance_supervision}
\end{figure}

Comme l'illustre la figure~\ref{fig:gouvernance_supervision}, ce processus de supervision s'articule autour de trois étapes majeures :
\begin{itemize}
    \item \textbf{Accès sécurisé au Dashboard :} L'Administrateur RH se connecte à l'interface de gouvernance (Frontend Angular). Le filtre de sécurité s'assure de la validité du jeton JWT et vérifie le rôle administrateur (\texttt{401 Unauthorized} / \texttt{403 Forbidden} en cas d'échec). Une fois autorisé, le système charge les demandes en attente et les indicateurs d'équipe depuis PostgreSQL.
    \item \textbf{Analyse et décision :} L'administrateur étudie une demande spécifique. Il peut l'approuver, la refuser en justifiant un motif, ou suggérer une formation alternative. La décision est soumise au backend qui effectue une validation métier (gestion des erreurs de type \texttt{400 Bad Request}).
    \item \textbf{Mise à jour et Notification temps réel :} Selon l'action (Approuver, Suggérer, Refuser), le service métier met à jour le statut de la demande en base de données et ajuste les accès aux cours si nécessaire. Le processus se termine par la publication d'un événement via le service de notification en temps réel, poussant la décision directement sur le navigateur de l'employé concerné, tandis que l'interface administrateur confirme le traitement avec un statut \texttt{200 OK}.
\end{itemize}

\subsubsection*{Cas 3 : Évaluation de formation}
L'étape d'évaluation des compétences acquises et la gestion des résultats sont modélisées par le diagramme de séquence de la figure~\ref{fig:evaluation_formation}.

\begin{figure}[htbp]
    \centering
    \includegraphics[width=\textwidth]{images/seq3Sprint3.png}
    \caption{Diagramme de séquence : évaluation de la formation}
    \label{fig:evaluation_formation}
\end{figure}

Comme détaillé sur la figure~\ref{fig:evaluation_formation}, le déroulement de l'évaluation se divise en trois phases principales :
\begin{itemize}
    \item \textbf{Vérification et Génération du Test (IA) :} Lors du lancement du test, le backend s'assure d'abord que l'employé a bien achevé le cours correspondant (renvoyant une erreur \texttt{403 Forbidden} dans le cas contraire). Si l'accès est valide, le système délègue la création d'un QCM adapté au Microservice IA. Une fois le format JSON validé, la session de test est enregistrée et le test chronométré est présenté à l'utilisateur.
    \item \textbf{Soumission et Calcul du Score :} Une fois le test terminé, les réponses sont soumises via l'interface. Le service métier évalue les réponses, vérifie le temps imparti, calcule le score final et actualise l'indicateur de préparation de l'employé en base de données.
    \item \textbf{Gestion des Résultats et Remédiation :} Le système applique une logique conditionnelle basée sur un seuil de réussite :
    \begin{itemize}
        \item Si le score est suffisant, une certification en attente est automatiquement créée.
        \item Si le score est insuffisant, le système interroge à nouveau le Microservice IA pour concevoir et sauvegarder un plan de révision ciblé. 
    \end{itemize}
    Dans les deux cas, le bilan est finalement renvoyé au Frontend pour affichage.
\end{itemize}
\subsubsection*{Cas 4: Validation RH et Certification}
La dernière étape du processus, concernant la validation par les ressources humaines et la délivrance du certificat, est modélisée par le diagramme de séquence de la figure~\ref{fig:validation_certification}.

\begin{figure}[H]
    \centering
    \includegraphics[width=\textwidth]{images/seq4Sprint3.png}
    \caption{Diagramme de séquence : validation RH et certification}
    \label{fig:validation_certification}
\end{figure}

Comme illustré sur la figure~\ref{fig:validation_certification}, ce processus final se déroule en deux grandes phases :
\begin{itemize}
    \item \textbf{Approbation et Génération du PDF :} L'Administrateur RH approuve une certification via le Frontend (Angular). Après validation sécurisée de ses droits d'accès par le filtre de sécurité, le service métier actualise le statut en base de données et déclenche la génération du certificat officiel. En cas de succès de la génération, le chemin d'accès au fichier est sauvegardé.
    \item \textbf{Notification et Téléchargement :} Une fois le document généré et persisté, le service de notification alerte l'employé. Ce dernier peut alors demander le téléchargement de son certificat. Le système effectue une ultime vérification de sécurité pour s'assurer de la validité de sa session (jeton JWT) et de son droit d'accès exclusif (appartenance du certificat) avant de lui transmettre le flux du fichier PDF.
\end{itemize}
\subsubsection*{11. Diagramme de Classes}
La structure logicielle implémentée pour répondre aux objectifs de cette phase est illustrée par le diagramme de classes de la figure~\ref{fig:architecture_sprint3}. 

\begin{figure}[H]
    \centering
    \includegraphics[width=\textwidth]{classSprint3.png}
    \caption{Diagramme de classes de l'architecture du Sprint 3}
    \label{fig:architecture_sprint3}
\end{figure}

Comme le démontre la figure~\ref{fig:architecture_sprint3}, l'application s'appuie sur une architecture multi-tiers (en couches) garantissant une séparation stricte des responsabilités :
\begin{itemize}
    \item \textbf{Couche de Présentation (Controllers Layer) :} Cette couche expose les points d'entrée de l'API REST (\texttt{RoadmapController}, \texttt{GovernanceController}, \texttt{EvaluationController}). Elle réceptionne les requêtes HTTP envoyées par le Frontend Angular, extrait les paramètres et retourne les réponses HTTP correspondantes (\texttt{200 OK}, \texttt{400 Bad Request}, etc.) sans embarquer de logique métier.
    \item \textbf{Couche Métier (Services Layer) :} Elle centralise la logique métier, les algorithmes de décision et les règles de gestion du système. Cette couche orchestre également les interactions avec les clients externes, notamment les microservices d'intelligence artificielle (\texttt{RoadmapAiClient}, \texttt{EvaluationAiClient}), la génération de documents PDF (\texttt{PdfExportService}) et les flux de notification temps réel (\texttt{SseNotificationService}).
    \item \textbf{Couche d'Accès aux Données (Persistence Layer) :} Composée d'interfaces de type Repository, elle assure l'abstraction complète des interactions avec la base de données PostgreSQL à l'aide de Spring Data JPA, réalisant les opérations de persistance sans manipulation directe de requêtes SQL.
    \item \textbf{Entités du Domaine (Domain Entities / ORM) :} Elle contient les classes d'objets modélisant les concepts réels de l'application (\texttt{User}, \texttt{Roadmap}, \texttt{TrainingRequest}, \texttt{Evaluation}, \texttt{Certificate}). Grâce au mapping objet-relationnel, ces entités structurent directement les tables de la base de données PostgreSQL ainsi que leurs cardinalités.
\end{itemize}
\subsection{Codage }
Consacrée à la phase d'implémentation spécifique au Sprint 3, cette section met en lumière les aspects les plus complexes du développement. Les fondations infrastructurelles de la plateforme étant désormais opérationnelles, notre travail de codage s'est concentré sur des enjeux à haute valeur ajoutée : le pilotage des modèles d'IA générative, la gestion des échanges entre les différents microservices, ainsi que le rendu dynamique et sécurisé des documents de certification.

\paragraph*{1. Orchestration de l'Intelligence Artificielle (Architecture Distribuée)}
Pour répondre au besoin de génération de contenu dynamique (\textit{Roadmaps} et Mini-tests), nous avons mis en place un écosystème \texttt{Python} découplé du cœur transactionnel Java :

\begin{itemize}
    \item \textbf{\texttt{FastAPI} \& \texttt{Uvicorn} :} Sélectionné comme socle de notre API cognitive pour ses hautes performances basées sur l'\textit{Event Loop} ASGI (\texttt{Starlette}). Il expose les \textit{webhooks} de génération tout en générant nativement les contrats d'interface (OpenAPI/Swagger).
    \item \textbf{\texttt{Ollama Python SDK } (\texttt{AsyncIO}) :}Librairie utilisée pour instancier des clients non-bloquants interagissant avec le modèle Llama 3 hébergé localement. Elle orchestre les échanges asynchrones avec le moteur d'IA en exploitant le mode de forçage JSON (JSON Schema). Cela garantit que les analyses de compétences et les prédictions produites par l'IA soient déterministes, structurées et directement parsables par les services du backend Java.
    \item \textbf{\texttt{Pydantic} :} Déployée de manière systématique pour imposer un typage fort en Python. Elle assure la validation sémantique et syntaxique des \textit{Data Models} à la volée, agissant comme un bouclier contre les erreurs de désérialisation avant l'envoi des données vers le \textit{Backend} Java.
    \item \textbf{\texttt{Spring RestTemplate} :} Côté Java, implémentation d'un client HTTP \textit{thread-safe} configuré avec des stratégies de \textit{Timeout} agressives et des gestionnaires d'exceptions (\texttt{ResponseErrorHandler}) pour assurer une résilience réseau (prévention de la famine des \textit{threads}) lors des appels au microservice IA.
\end{itemize}

\paragraph*{2. Workflow de Gouvernance et Rendu Documentaire (Full-Stack)}
Pour transformer la plateforme en un outil certifiant et interactif, nous avons implémenté des mécanismes avancés de manipulation du DOM et de génération à la volée :

\begin{itemize}
    \item \textbf{\texttt{OpenHTMLToPDF} \& \texttt{Thymeleaf} :} Ce duo constitue notre moteur de \textit{Server-Side Rendering} (SSR) documentaire. \texttt{Thymeleaf} résout dynamiquement l'arbre DOM en y injectant le contexte métier (identifiants cryptographiques, \textit{Readiness Score}). \texttt{OpenHTMLtoPDF} compile ensuite ce DOM XHTML en un flux d'octets PDF vectoriel, inviolable et prêt au téléchargement.
    \item \textbf{\texttt{Spring SseEmitter} \& \texttt{RxJS} :} Implémentation du \textit{pattern Publisher/Subscriber}. Le \textit{backend} Java maintient une connexion HTTP unidirectionnelle (\textit{Server-Sent Events}) pour publier les décisions de l'Admin. Côté Angular, \texttt{RxJS} écoute ce flux via l'objet \texttt{EventSource}. Les flux sont gérés de manière réactive avec un nettoyage strict des souscriptions (\textit{unsubscribes}) pour éviter toute fuite de mémoire (\textit{Memory Leak}).
    \item \textbf{\texttt{Angular CDK} (\textit{Drag and Drop}) :} Intégration des composants de bas niveau de Google pour manipuler l'arbre DOM virtuel. Cela permet au candidat de gérer les statuts de sa formation sur un tableau \textit{Kanban} interactif garantissant une fluidité à 60 FPS.
    \item \textbf{\texttt{Ngx-Toastr} :} Couplée aux \textit{Observables} \texttt{RxJS}, cette librairie consomme les flux asynchrones SSE pour afficher des alertes visuelles élégantes et éphémères (ex. "Certification approuvée"), améliorant considérablement l'UX globale.
\end{itemize}

\paragraph*{3. Outillage DevOps et Qualité du Code (\textit{Clean Code})}
Afin de garantir la robustesse de ces nouvelles fonctionnalités critiques, nous avons maintenu un niveau d'exigence industriel sur l'outillage et les tests :

\begin{itemize}
    \item \textbf{\texttt{Postman} :} Utilisé intensivement pour la création de collections d'API. Il a permis de simuler (\textit{mocker}) les retours du modèle LLM et de tester les flux binaires (\textit{Headers} PDF) de manière isolée avant l'intégration du client Angular.
    \item \textbf{\texttt{Java Bean Validation} (\texttt{Hibernate Validator}) :} Mécanisme déployé systématiquement par programmation orientée aspect (AOP). Les annotations (\texttt{@Valid}, \texttt{@NotNull}) appliquées aux objets de transfert (DTOs) assurent l'assainissement (\textit{Sanitization}) des requêtes entrantes, bloquant les \textit{payloads} malformés dès l'entrée du contrôleur.
    \item \textbf{\texttt{GitLab CI/CD} \& Stratégie GitFlow :} L'historique du code source est protégé par une stratégie de branches stricte. Les \textit{Merge Requests} ciblent systématiquement la branche de pré-production (\texttt{main}), déclenchant les pipelines d'intégration continue pour auditer le code et prévenir toute régression architecturale.
\end{itemize}










\subsection{Réalisation}
Cette partie détaille les résultats techniques du sprint, notamment les interfaces utilisateur développées avec \textbf{Angular}. L’objectif était d’offrir une expérience immersive, gamifiée et intuitive, tout en reflétant visuellement la complexité des algorithmes d’IA et des workflows de validation en arrière-plan.
\Needspace{20\baselineskip}
\subsection*{1. Tableau de Bord Global pour l'Upskilling}

Le Tableau de Bord Global pour l’Upskilling débute par l’identification des manques de compétences. Il mesure l’employabilité d’un collaborateur grâce à un indicateur \textit{« Prêt pour le poste »} (\textit{Readiness Score}) qqui permet d’analyser précisément les écarts de compétences techniques et comportementales, mettant ainsi en évidence le besoin de formation.
L’interface montre aussi en temps réel le niveau de l’utilisateur (LVL) et ses points d'expérience (XP),  intégrant ainsi une dimension gamifiée pour inciter l’apprenant à s’y engager davantage.
\begin{figure}[H]
    \centering
    \includegraphics[width=\textwidth]{images/dashboredFormation.png} 
    \caption{Tableau de Bord Global pour l'Upskilling}
    \label{fig:dashboard_upskilling}
\end{figure}
\Needspace{20\baselineskip}
\subsection*{2. La Sélection (Catalogue de Recommandations Intelligentes)}

Afin de mettre en œuvre la \textit{Roadmap}, le système suggère une sélection de cours adaptés proposés par plusieurs prestataires extérieurs (\textit{Coursera}, \textbf{Udemy}, \textbf{LinkedIn Learning}). 

L’innovation principale de cette interface est son aspect prédictif : chaque carte de formation affiche l’impact estimé sur l’employabilité du candidat (ex: \textit{+5\% Readiness}).L’utilisateur se voit ainsi permettre de filtrer les cours par priorité et, en cliquant sur Ajouter au plan, de formuler une demande d’inscription.
\begin{figure}[H]
    \centering
    \includegraphics[width=\textwidth]{images/recommandation.png}
    \caption{ La Sélection (Catalogue de Recommandations Intelligentes)}
    \label{fig:course_recommendations}
\end{figure}
\subsection*{3. La Stratégie (Roadmap générée par l'IA)}

À partir du diagnostic, le système propose une solution macroscopique. Page d’apparence Parcours Stratégique transcrit les résultats du microservice IA dans une feuille de route actionnable.

Cette dernière propose des phases d’apprentissage précises  (ex: \textit{Phase 1 — Foundations \& Time Blocking}) sur mesure pour combler directement les faiblesses identifiées lors de l’évaluation initiale.
\begin{figure}[H]
    \centering
    \includegraphics[width=\textwidth]{images/roadmap.png}
    \caption{La Stratégie (Roadmap)}
    \label{fig:ai_roadmap}
\end{figure}
\subsection* {4. L'Exécution (Gestion Kanban des Formations)}
Lorsque les cours sont choisis dans le catalogue, le collaborateur passe à l’action. L’interface sous forme de tableau \texttt{Kanban} permet à l’usager de piloter visuellement le cycle de vie de ses formations (de l'état \textit{Proposées} à \textit{Terminées}). 

C'est l'espace de travail quotidien de l'apprenant pour suivre, d'une manière fluide et interactive, sa progression, ayant une visibilité claire sur les tâches en cours et les objectifs atteints.
\begin{figure}[H]
    \centering
    \includegraphics[width=\textwidth]{images/mescourse.png}
    \caption{L'Exécution (Gestion Kanban des Formations)}
    \label{fig:kanban_execution}
\end{figure}
\subsection*{4. La Validation (Interface du Mini-Test IA)}

Une fois la formation complétée à 100\%, l'utilisateur peut passer un \textbf{Mini-Test de certification} généré dynamiquement par l'IA. Ce test apparaît directement sur le tableau \texttt{Kanban} dans une fenêtre contextuelle. Ce questionnaire à choix multiples de quatre options est créé en fonction du contenu pédagogique (\textit{ex: Angular Material}). Cela garantit que l'évaluation soit pertinente et unique. Les réponses sont vérifiées de manière déterministe par Spring Boot afin de maximiser la performance. L'obtention d'un score $\geq 70\%$ automatise le passage au statut \texttt{TERMINEE}, l'attribution des XP et la mise à jour du \textit{Readiness Score}.
\begin{figure}[H]
    \centering
    \includegraphics[width=\textwidth]{images/minitest.png}
    \caption{Interface du Mini-Test}
    \label{fig:minitest_interface}
\end{figure}
\subsection*{5. La Gouvernance (Espace Administration RH)}

L'exécution de la formation nécessite un encadrement managérial rigoureux. L'interface d'administration centralise les KPI et permet aux responsables RH de superviser les demandes issues du tableau Kanban des collaborateurs. 

LLes administrateurs peuvent valider ou rejeter les inscriptions. Ils vérifient les certificats soumis et ajoutent des notes, ce que l'on appelle des (\textit{adminNote}), Cela permet d'assurer que les formations individuelles sont alignées sur les objectifs globaux de l'entreprise.

\begin{figure}[H]
    \centering
    \includegraphics[width=\textwidth]{images/adminFormation.png}
    \caption{La Gouvernance (Espace Administration RH)}
    \label{fig:rh_governance}
\end{figure}

\subsection*{6. La Rétention (Leaderboard et Gamification)}

Le parcours se termine par la reconnaissance des efforts fournis. Une fois que l'administration a validé les formations, le collaborateur obtient des badges spéciaux, comme  \textit{Premier Certif} / \textit{Compétences comportementales Pro} et accumule de l'XP .

L'interface de classement (\textit{Leaderboard}) stimule l'engagement continu et la saine compétition en permettant aux apprenants de se situer par rapport à leurs pairs, renforçant ainsi la culture de l'apprentissage au sein de l'organisation.
D'un point de vue technique, le calcul des points XP est effectué côté backend à chaque transition d'état dans le Kanban (\texttt{TO\_DO} $\rightarrow$ \texttt{IN\_PROGRESS} $\rightarrow$ \texttt{DONE}), garantissant ainsi l'intégrité du classement. Le Leaderboard est mis à jour en temps réel via le même canal SSE utilisé pour les notifications de gouvernance.
\begin{figure}[H]
    \centering
    \includegraphics[width=\textwidth]{images/classement.png}
    \caption{ La Rétention (Leaderboard et Gamification)}
    \label{fig:gamification_leaderboard}
\end{figure}

\subsection{Phase de Test et Déploiement }

Cette section décrit la méthodologie appliquée pour valider la robustesse de la solution et garantir que l'intégration des nouvelles bibliothèques complexes (moteurs PDF, connecteurs IA, flux asynchrones) ne compromet pas la stabilité du cœur transactionnel de l'application.

\subsubsection*{1. Stratégie d'Assurance Qualité (QA)}

Afin de garantir une couverture optimale, la validation du code a été découpée en deux niveaux de tests automatisés :

\begin{itemize}
    \item \textbf{Tests Unitaires (\texttt{JUnit 5} \& \texttt{Mockito}) :} L'objectif principal a été d'isoler la logique métier sans dépendre de la base de données. Des tests ont été rédigés pour le \texttt{PdfExportService} en simulant (\textit{mocking}) le contexte pour vérifier que la conversion HTML-vers-PDF ne lève pas d'exception. De même, l'algorithme de calcul du \textit{Readiness Score} a été testé avec différents jeux de réponses pour s'assurer de sa fiabilité mathématique.
    \item \textbf{Tests d'Intégration d'API (\texttt{Postman}) :} L'effort s'est concentré sur l'architecture distribuée. Des collections \texttt{Postman} ont été créées pour vérifier le dialogue HTTP entre le \textit{Backend} Java et le microservice \texttt{Python} (\texttt{FastAPI}). Nous avons volontairement simulé des latences réseau pour valider que le \texttt{RestTemplate} Java appliquait correctement ses \textit{Timeouts} sans saturer la mémoire du serveur.
\end{itemize}

\subsubsection*{2. Plan de Validation Fonctionnelle (Tests \textit{End-to-End})}

Pour valider l'expérience utilisateur, un plan de tests de bout en bout (E2E) a été exécuté depuis l'interface \texttt{Angular}, ciblant les trois cas d'utilisation majeurs du sprint :

\begin{table}[H]
    \centering
    \renewcommand{\arraystretch}{1.5}
    \small
    \begin{tabularx}{\textwidth}{|c|>{\raggedright\arraybackslash}p{3cm}|>{\raggedright\arraybackslash}X|>{\raggedright\arraybackslash}X|}
        \hline
        \rowcolor{TableHeaderBlue}
        \color{white}\textbf{ID} & \color{white}\textbf{Scénario Testé} & \color{white}\textbf{Pré-conditions \& Étapes} & \color{white}\textbf{Résultat Attendu} \\
        \hline
        TEST-01 & Génération de la \textit{Roadmap} par l'IA & \textbf{Pré-cond :} Profil utilisateur connecté avec un \textit{Skill Gap} existant. \newline \textbf{Étapes :} 1. Accès espace \textit{Upskilling}. 2. Choix poste cible. 3. Clic sur "Générer". & Le \textit{Backend} interroge \texttt{FastAPI} avec succès. Le JSON est désérialisé sans erreur. Le \textit{Kanban} est alimenté dynamiquement. \\
        \hline
        TEST-02 & Soumission du test et calcul du Score & \textbf{Pré-cond :} Test lancé et chronomètre actif. \newline \textbf{Étapes :} 1. Sélectionner les réponses. 2. Soumission. & Le \textit{Backend} corrige le test. Le \textit{Readiness Score} est actualisé. Une demande de certification \texttt{PENDING\_APPROVAL} est créée. \\
        \hline
        TEST-03 & Approbation RH et Notifications (SSE) & \textbf{Pré-cond :} Admin et utilisateur connectés simultanément. \newline \textbf{Étapes :} 1. Choix de la demande par l'Admin. 2. Clic sur "Approuver". & La transaction SQL passe à \texttt{APPROVED}. L'utilisateur reçoit une alerte visuelle (\texttt{Ngx-Toastr}) instantanément via \texttt{SSE}. \\
        \hline
        TEST-04 & Génération documentaire (Moteur PDF) & \textbf{Pré-cond :} Certification approuvée. \newline \textbf{Étapes :} 1. Clic sur "Télécharger le diplôme". & Pas d'erreur 500. Téléchargement d'un \texttt{.pdf} valide incluant dynamiquement le nom et le score via \texttt{Thymeleaf}. \\
        \hline
    \end{tabularx}
    \caption{Plan de validation fonctionnelle du Sprint 3}
    \label{tab:validation_sprint3}
\end{table}

\subsubsection*{3. Résultats obtenus et Corrections apportées (Débogage)}

Durant cette phase d'intégration, nous avons été confrontés à des erreurs techniques complexes :

\begin{itemize}
    \item \textbf{Erreur 422 (\textit{Unprocessable Content}) :} 
    \textit{Problème :} \texttt{FastAPI} rejetait le \textit{payload} JSON du Java.
    \textit{Correction :} Alignement des DTOs Java avec le modèle \texttt{Pydantic} (\texttt{CareerPredictionBody}) côté \texttt{Python}.
    \item \textbf{Erreur 500 (Conflit de dépendances) :} 
    \textit{Problème :} Conflit entre \texttt{PDFBox 3.x} et \texttt{OpenHTMLtoPDF}.
    \textit{Correction :} Ajustement des versions dans le \texttt{pom.xml} pour restaurer la compatibilité d'architecture.
    \item \textbf{Erreur 500 (Rendu XHTML) :} 
    \textit{Problème :} Erreur fatale lors du \textit{parsing} du certificat.
    \textit{Correction :} Réécriture du \textit{template} en XHTML strict (balises fermées) et correction du CSS pour le moteur de rendu.
\end{itemize}

\subsubsection*{3. Processus de Déploiement et Exécution}

L'intégration du microservice IA a nécessité une évolution de l'infrastructure via la conteneurisation :

\begin{itemize}
    \item \textbf{Conteneurisation (\texttt{Docker}) :} Création de \texttt{Dockerfiles} isolés pour le \textit{Backend} (JDK 21), le \textit{Frontend} (\texttt{Nginx}) et le microservice IA (\texttt{Python 3.11}).
    \item \textbf{Orchestration Locale (\texttt{Docker Compose}) :} Le fichier \texttt{docker-compose.yml} gère le démarrage ordonné : \texttt{PostgreSQL} $\rightarrow$ \texttt{FastAPI} $\rightarrow$ \texttt{Spring Boot} $\rightarrow$ \texttt{Angular}.
    \item \textbf{Sécurisation des Secrets :} LLes variables sensibles (Clé API inter-services, secrets \texttt{JWT}, identifiants \texttt{PostgreSQL}) sont injectées au \textit{runtime} via un fichier \texttt{.env} protégé. Cela permet de sécuriser la communication entre le \textit{backend} \texttt{Java} et le microservice IA sans exposer de clés privées dans le code source.
\end{itemize}
\section{Suivi et Pilotage du Sprint}
\subsection{ Scrum Board}
Afin de garantir la transparence et l’inspection continue, qui sont essentielles dans la méthodologie Scrum, nous avons utilisé deux outils visuels importants pour gérer cette itération. Le sprint a duré 3 semaines (soit 15 jours ouvrés), avec une charge de travail estimée à 34 Story Points.

Le suivi quotidien des tâches s'est matérialisé par un tableau de bord reflétant l'état d'avancement des \textit{User Stories}. Ce management visuel a permis à l'équipe de développement d'identifier rapidement les goulots d'étranglement et de fluidifier le passage des tâches de la phase de conception vers la phase de test.

\begin{figure}[H]
    \centering
    % Assurez-vous que l'image scrum_board.png est bien dans votre dossier d'images
    \includegraphics[width=0.9\textwidth]{images/scrum_bored.png}
    \caption{Scrum board du sprint 3}
    \label{fig:scrum_board}
\end{figure}

\subsection{ Le Burndown Chart}

Pour évaluer notre vélocité et s'assurer de l'atteinte de l'objectif dans le temps imparti, nous avons maintenu à jour le \textit{Burndown Chart} lors de nos \textit{Daily Stand-ups}. 

Le graphique montre la courbe de l’effort restant en  \textit{Story Points} au fil des 15 jours du sprint. Nous avons rencontré quelques difficultés techniques entre Jour 9 et le Jour 11 s'expliquent par les défis techniques rencontrés lors de la configuration du réseau interne Docker pour le microservice IA (FastAPI). La résolution de ces obstacles a permis une accélération de la validation finale, clôturant le sprint dans les délais.

\begin{figure}[H]
    \centering
    \begin{tikzpicture}
        \begin{axis}[
            width=0.9\textwidth,
            height=8cm,
            title={\textbf{Burndown Chart - Sprint 3 (Upskilling \& Formation)}},
            xlabel={Jours ouvrés du Sprint (3 Semaines)},
            ylabel={Story Points restants},
            xmin=0, xmax=15,
            ymin=0, ymax=40,
            xtick={0,1,2,3,4,5,6,7,8,9,10,11,12,13,14,15},
            ytick={0,10,20,30,34,40},
            legend pos=north east,
            ymajorgrids=true,
            xmajorgrids=true,
            grid style=dashed,
            axis background/.style={fill=gray!5},
            thick
        ]

        % La ligne idéale (l'objectif linéaire)
        \addplot[
            color=gray,
            dashed,
            very thick,
            mark=none
            ]
            coordinates {
            (0,34)(15,0)
            };
            \addlegendentry{Tendance Idéale}

        % La ligne réelle (l'avancement de l'équipe avec les blocages)
        \addplot[
            color=blue!70!black,
            very thick,
            mark=square*,
            mark options={fill=blue!70!black, scale=0.9}
            ]
            coordinates {
            (0,34)     % Jour 0 : Lancement
            (1,34)     % Jour 1 : Analyse
            (2,32)     % Jour 2 : Premières validations
            (3,29)     % Jour 3
            (4,29)     % Jour 4 : Focus sur le code long
            (5,26)     % Jour 5 : Fin de semaine 1
            (6,23)     % Jour 6 : Début semaine 2
            (7,21)     % Jour 7
            (8,21)     % Jour 8
            (9,18)     % Jour 9 : Début des problèmes Docker/IA
            (10,18)    % Jour 10 : Recherche de solution
            (11,18)    % Jour 11 : Blocage résolu
            (12,10)    % Jour 12 : Validation massive des tickets
            (13,6)     % Jour 13 : Gamification terminée
            (14,3)     % Jour 14 : Tests finaux
            (15,0)     % Jour 15 : Démo et clôture
            };
            \addlegendentry{Avancement Réel}

        \end{axis}
    \end{tikzpicture}
    \caption{Burndown Chart du sprint 3}
    \label{fig:burndown_chart}
\end{figure}

\subsection{Indicateurs de Performance (KPIs) du Sprint 3}

Afin d'évaluer la réussite de cette itération focalisée sur l'intelligence RH et l'\textit{upskilling}, nous avons mis en place un système de suivi basé sur des indicateurs précis.

\paragraph*{1. Indicateurs de Pilotage (Agilité)}
Le pilotage de ce sprint a été marqué par une forte densité technique (intégration de l'IA et des \texttt{WebSockets}).

\begin{itemize}
    \item \textbf{Vélocité Réalisée :} Sur les 34 \textit{Story Points} planifiés, la totalité a été livrée dans les délais, démontrant une maturité croissante de l'équipe sur la \textit{stack} technologique.
    \item \textbf{Densité de Développement :} La livraison de 12 nouveaux \textit{endpoints} \texttt{API} complexes souligne l'efficacité de la collaboration entre les modules \textit{Backend}, \textit{Frontend} et IA.
    \item \textbf{Qualité des \textit{User Stories} :} Le taux de validation en fin de sprint par le \textit{Product Owner} a atteint 100\%, sans aucun report de tâche (\textit{Carry-over}).
\end{itemize}

\paragraph*{2. Indicateurs Techniques et Qualité}
Le tableau ci-dessous synthétise les performances de la plateforme à l'issue du Sprint 3, couvrant les aspects logiciels, l'intelligence artificielle et l'infrastructure \texttt{Docker}.

\begin{table}[H]
    \centering
    \renewcommand{\arraystretch}{1.5}
    \small
    \begin{tabularx}{0.95\textwidth}{|X|c|c|c|}
        \hline
        \rowcolor{TableHeaderBlue}
        \color{white}\textbf{Indicateur} & \color{white}\textbf{Cible} & \color{white}\textbf{Résultat Obtenu} & \color{white}\textbf{État} \\
        \hline
        \textit{Story Points} complétés & 34 & 34 & $\checkmark$ Validé \\
        \hline
        Disponibilité (\textit{Uptime} \texttt{Docker}) & $100\%$ & $100\%$ & $\checkmark$ Validé \\
        \hline
        Poids des images \texttt{Docker} & $< 600$ Mo & $\sim 410$ Mo & $\checkmark$ Validé \\
        \hline
        Précision du \textit{Gap Analysis} (IA) & $> 90\%$ & $94\%$ & $\checkmark$ Validé \\
        \hline
        Fiabilité du \textit{Matching} Formation & $> 85\%$ & $89\%$ & $\checkmark$ Validé \\
        \hline
        Taux de réussite des tests E2E & $> 90\%$ & $100\%$ & $\checkmark$ Validé \\
        \hline
        Temps de génération \textit{Roadmap} IA & $< 5$ s & $\sim 3.2$ s & $\checkmark$ Validé \\
        \hline
        Temps de génération PDF (Passeport) & $< 3$ s & $\sim 1.8$ s & $\checkmark$ Validé \\
        \hline
        Latence \texttt{SSE} (Notifications temps réel) & $< 500$ ms & $\sim 180$ ms & $\checkmark$ Validé \\
        \hline
    \end{tabularx}
    \caption{Indicateurs de Performance du Sprint 3}
    \label{tab:kpi_tech_s3}
\end{table}

\paragraph*{3. Analyse du Bilan}
Le bilan du Sprint 3 est très positif, tant sur le plan fonctionnel que technique. L'intégration réussie de l'analyse des écarts (\textit{Skill-Gap Analysis}) sur les deux dimensions (\textit{compétences techniques et comportementales}) constitue la plus grande valeur ajoutée de cette itération.

La stabilité des infrastructures \texttt{Docker}, combinée à une latence de notification très faible (180 ms), prouve que la plateforme est capable de supporter des interactions en temps réel complexes. Enfin, la rapidité de génération des documents officiels (\texttt{PDF}) et des \textit{roadmaps} d'apprentissage confirme que l'architecture est optimisée pour une utilisation intensive par les départements RH.


\section{Conclusion }

Dans ce chapitre, nous avons présenté la conception, la réalisation et la validation du module  de\textit{Formation} et d'\textit{Upskilling},qui est au cœur de la plateforme TalentPredict. Nous avons utilisé une modélisation UML rigoureuse pour structurer les flux de données complexes entre les collaborateurs, l’administration RH et l’Intelligence Artificielle.
Le passage de la théorie à la pratique s'est matérialisé par le développement d'interfaces \textit{Angular} dynamiques et ergonomiques, soutenues par une architecture \textit{backend} hybride (\textit{Spring Boot} et \textit{FastAPI}). Nous avons réussi à concrétiser la promesse initiale du projet : transformer les lacunes de compétences (\textit{Skill Gaps}) en opportunités d'apprentissage personnalisées grâce aux prédictions génératives du LLM.

Par ailleurs, l'intégration de mécanismes de gamification (niveaux, badges) et d'un pilotage visuel (\textit{Kanban}) favorise l'engagement de l'apprenant, tandis que le \textit{workflow} d'approbation et les notifications en temps réel garantissent une gouvernance RH maîtrisée. Enfin, l'application stricte de la méthodologie Agile \textit{Scrum} et l'exécution des tests fonctionnels ont prouvé la fiabilité et la robustesse de ce système.

Avec la validation de ce troisième sprint, le cycle de vie de la gestion des talents est désormais opérationnel de bout en bout. Le chapitre suivant s'attachera à présenter le quatrième sprint.
\chapter{ Sprint 4 :Pilotage, Management \& Mise en Production}
\section*{Introduction}
\phantomsection
\addcontentsline{toc}{section}{Introduction}
\phantomsection 
Ce chapitre est la dernière étape du projet \textbf{TalentPredict}. Nous avons développé les moteurs d'évaluation et l'intelligence artificielle lors des étapes précédentes. Maintenant, nous nous concentrons sur la façon dont nous allons utiliser les données pour aider les ressources humaines. Nous voulons créer des outils pour les aider à prendre des décisions et à communiquer facilement. Nous allons également consolider tous les modules que nous avons créés, générer le \texttt{Talent Passport} et mettre la plateforme en ligne grâce à \textbf{Docker} et l'implémentation d'une chaîne
 \textbf{CI/CD} (Intégration et Déploiement Continus). Cette étape finale assure une qualité de code constante, une automatisation des tests et une mise en production sécurisée, livrant ainsi une plateforme solide, adaptable et prête à l’emploi.
\section{Objectifs du Sprint 4}
% Remarque : Utilisez \section{} au lieu de \subsection{} si "4.1" correspond au premier niveau de titre dans votre chapitre.

L’objectif principal de ce quatrième et dernier sprint est de renforcer les différents modules développés précédemment pour obtenir une version finale et stable de la plateforme \textbf{TalentPredict} . Ce sprint se concentre sur deux axes : la finition des interfaces de pilotage stratégique et l'industrialisation de la plateforme. 

    \subsubsection*{2. Objectifs fonctionnels et interactifs :}
    \begin{itemize}
        \item \textbf{Conception d'un tableau de bord décisionnel (\textit{Executive Dashboard}) :} Fournir aux administrateurs RH une vue centralisée et en temps réel des indicateurs clés de performance (KPIs) stratégiques pour le pilotage des talents.
        \item \textbf{Développement de l'espace collaborateur (\textit{User Dashboard}) :} Créer une interface personnalisée et ergonomique, en adoptant le design \textit{Bento Grid} . Cela inclut un radar de compétences dynamique et la présentation du persona IA.
        \item \textbf{Mise en place de la gestion administrative (Module RH) :} Développer un module complet (CRUD) pour gérer les comptes employés. Cela permettra de consulter leurs profils publics (liens sociaux, résultats détaillés, informations biographiques) tout en sécurisant les accès grâce à un contrôle strict basé sur les rôles (RBAC).
        \item \textbf{Création du \textit{Communication Hub} :} Centraliser les échanges avec un système de notifications en temps réel SSE(\textit{Server-Sent Events}) et permettre aux employés d’interagir avec les recommandations de l’IA via une boucle de rétroaction.
        \item \textbf{Génération de rapports professionnels :} 
        Produire, côté serveur, des documents structurés et fidèles. Cela inclut le \textit{Talent Passport} pour le bilan de compétences individuel des employés et le \textit{Rapport Global RH }  pour la synthèse analytique et stratégique de l'ensemble du capital humain..
    \end{itemize}

    \subsubsection*{2. Objectifs techniques et opérationnels (DevOps) :}
    \begin{itemize}
        \item \textbf{Conteneurisation de l'architecture :} Mettre en place un déploiement complet et isolé des différents services (Frontend, Backend, IA, Base de données) via \texttt{Docker} et \texttt{Docker Compose}.
        \item \textbf{Automatisation CI/CD :} Configurer un pipeline d'intégration et de déploiement continus sous \texttt{GitLab CI} pour garantir que les tests automatisés soient exécutés à chaque évolution du code.
        \item \textbf{Stabilisation et Performance :} Optimiser les temps de réponse des services asynchrones et assurer une cohérence totale des données entre les modules d'évaluation, de formation et de dashboarding.
    \end{itemize}

\subsubsection*{Critères d'acceptation (\textit{Definition of Done})}
Pour considérer ce sprint comme achevé, l'équipe a défini les critères de validation suivants :
\begin{itemize}
    \item Tous les tests unitaires du backend doivent s'exécuter avec succès.
    \item La génération du génération de rapports professionnels doit s'effectuer correctement, en respectant la mise en page spécifiée pour chaque utilisateur.
    \item Les notifications asynchrones doivent être réceptionnées en temps réel sans perte de connexion.
    \item  Le système doit montrer une stabilité satisfaisante dans l’environnement de pré-production  (\textit{staging}).
\end{itemize}
\section{Séance de Sprint Planning}

La séance de Sprint Planning a réuni l'équipe de développement, composée de \textit{Mohamed Ben Abdeladhim} et \textit{Rahma Barouni}, ainsi que le Product Owner,Mr \textit{Med Amine Gharbi}, et notre encadrante pédagogique, Mme \textit{Ben Aïcha Takwa}, afin de définir le périmètre de cette dernière itération de développement. Les points clés suivants ont été déterminés lors de cette réunion.

\begin{itemize}
    \item \textbf{Durée :} La durée de cette itération a été fixée à un mois, ce qui nous permet de disposer du temps nécessaire pour achever les développements complexes, réaliser les tests finaux et assurer le déploiement.
    \item \textbf{Priorités absolues :} Le développement des tableaux de bord pour les ressources humaines et les utilisateurs, le centre de communication, la génération automatisée des rapports professionnels et la mise en place du module de gestion des employés permettent aux ressources humaines d'avoir un contrôle total.
    \item \textbf{Priorités secondaires :} Optimisation du moteur d’IA pour améliorer ses performances.
 et la configuration du pipeline CI/CD via \texttt{Docker},et \texttt{GitLab}.
    \item \textbf{Prérequis :} Un prérequis essentiel pour cette itération est la stabilité confirmée des API livrées lors des sprints précédents, notamment le moteur d’IA, le module d’évaluation et le module de recommendations.
\end{itemize}

À la fin de cette réunion, les besoins ont été décomposés en \textit{User Stories} (US) et estimés en jours-homme, en donnant la priorité maximale aux livrables critiques pour la démonstration finale de \textbf{TalentPredict}.

\section{Backlog du Sprint}
% Utilisez \subsubsection{} selon la hiérarchie de vos titres.
\Needspace{18\baselineskip}

Pour mettre en œuvre les priorités définies lors de la planification du sprint et faciliter la répartition des tâches entre les membres de l'équipe (Mohamed et Rahma), chaque \textit{User Story} a été décomposée en tâches techniques spécifiques. Le tableau présente ce backlog enrichi, en incluant les responsabilités et les critères de validation associés.

\vspace{-0.3em}
\begin{table}[H]
    \centering
    \renewcommand{\arraystretch}{1}
    \scriptsize
    \begin{tabularx}{\textwidth}{|c|>{\raggedright\arraybackslash}p{3cm}|>{\raggedright\arraybackslash}X|c|>{\raggedright\arraybackslash}X|c|}
        \hline
        \rowcolor{TableHeaderBlue}
        \color{white}\textbf{ID} & \color{white}\textbf{Épic} & \color{white}\textbf{User Story} & \color{white}\textbf{Priorité} & \color{white}\textbf{Critères d'acceptation} & \color{white}\textbf{Estim. (SP)} \\
        \hline
        
        % US 4.1
        US4.1 & Système de Reporting 
        & En tant qu'employé, je veux exporter mon Passeport Talent en PDF. 
        & Moyenne 
        & \textbullet{} PDF fidèle au design, incluant scores IA et radar. 
        & 8 \\
        \hline
        
        % US 4.2
        US4.2 & \textit{Executive Analytics} 
        & En tant que RH, je veux visualiser les KPIs stratégiques sur un \textit{Dashboard} Admin. 
        & Haute 
        & \textbullet{} \textit{Dashboard Bento} fonctionnel avec graphiques \texttt{Chart.js}. 
        & 8 \\
        \hline
        
        % US 4.3
        US4.3 & Expérience Utilisateur 
        & En tant qu'employé, je veux voir mon radar de compétences et mon \textit{persona} IA. 
        & Haute
        & \textbullet{} Visualisation dynamique des compétences et résultats PCM. 
        & 5 \\
        \hline
        
        % US 4.4
        US4.4 & \textit{Communication Hub} 
        & En tant qu'utilisateur, je veux recevoir des notifications système en temps réel. 
        & Haute 
        & \textbullet{} Réception via WebSockets sans rafraîchissement de page. 
        & 5 \\
        \hline
        
        % US 4.5
        US4.5 & Pilotage Stratégique 
        & En tant qu'admin, je veux télécharger un rapport analytique global de l'entreprise. 
        & Moyenne 
        & \textbullet{} PDF consolidé avec statistiques par département. 
        & 5 \\
        \hline
        
        % US 4.6
        US4.6 & Gouvernance RH 
        & En tant qu'admin, je veux gérer (CRUD) et activer/désactiver les comptes employés. 
        & Haute 
        & \textbullet{} Changement de statut instantané et accès aux profils publics. 
        & 5 \\
        \hline
        
        % US 4.7
        US4.7 & Infrastructure \& DevOps 
        & En tant que Dev, je veux conteneuriser l'appli et automatiser le déploiement. 
        & Haute 
        & \textbullet{} Projet déployable via \texttt{Docker Compose} et pipeline \texttt{GitLab CI}. 
        & 13 \\
        \hline
        
    \end{tabularx}
    \vspace{1mm}
    \caption{Baklog du Sprint 4}
    \label{tab:backlog_sprint4}
\end{table}
\section{Démarrage du Sprint 4}
% Vous pouvez adapter le niveau de titre (\section, \subsection) selon la structure exacte de votre chapitre.
Une fois que l’équipe a validé le périmètre du sprint et a réparti les tâches entre ses membres, la phase de réalisation a réellement commencé. Cette phase a débuté par une analyse approfondie et détaillée des spécifications techniques. Ensuite, l’équipe a conçu l’architecture complète des nouveaux modules, ce qui est une étape cruciale. L’objectif principal de cette étape est de garantir que les solutions implémentées s’intègrent parfaitement avec le code produit lors des itérations précédentes. Les solutions doivent s’intégrer de manière harmonieuse avec le code existant pour éviter tout problème.

\subsubsection*{1. Activités réalisées}% La numérotation 3.4.1 se fera automatiquement selon la structure de votre document.

Durant cette itération, nous avons mené de front plusieurs activités critiques :

\begin{itemize}
    \item \textbf{Intégration du Moteur de Reporting :} Développement de la logique \textit{Backend} pour transformer les \textit{templates} HTML en documents PDF officiels.
    \item \textbf{Conception des Dashboards Bento :} Développement de l'interface réactive et intégration de la bibliothèque \texttt{Chart.js} pour la visualisation des compétences.
    \item \textbf{Configuration du Hub de Communication :}Implémentation du flux Server-Sent Events (SSE) pour les notifications en temps réel.
    \item \textbf{Travaux DevOps :} Écriture des fichiers de configuration Docker et mise en place du pipeline d'intégration continue sous \texttt{GitLab CI}.
    \item \textbf{Synchronisation des Données :} \textit{Mapping} des prédictions issues du service Python vers les objets Java (DTO) pour une restitution fluide.
\end{itemize}

\subsubsection*{2. Livrables produits}

À l'issue du sprint, les livrables suivants ont été finalisés :

\begin{itemize}
    \item \textbf{Modules Frontend :} Dashboards Admin/User, Centre de Notifications et Gestion des employés.
    \item \textbf{Services Backend :} Reporting Service,SSE Notification Hub et API de gestion RH.
    \item \textbf{Documents de Restitution :} \textit{Templates} PDF du Passeport Talent et du Rapport Global RH.
    \item \textbf{Artefacts DevOps :} Fichiers \texttt{Dockerfile} et \texttt{docker-compose.yml} opérationnels.
    \item \textbf{Suite de Tests :} Rapports de tests unitaires et d'intégration validés par la CI/CD.
\end{itemize}

\subsubsection{3. Difficultés rencontrées et traitement}

Le développement a présenté certains défis techniques que nous avons dû résoudre :

\begin{itemize}
    \item \textbf{Difficulté 1 : Rendu des polices dans le PDF}
    \begin{itemize}
        \item \textbf{Problème :} Les caractères spéciaux et les icônes ne s'affichaient pas correctement dans les rapports exportés.
        \item \textbf{Traitement :} Intégration de polices vectorielles TrueType (TTF) directement dans le moteur de rendu \texttt{OpenHTMLToPDF}.
    \end{itemize}
    \vspace{2mm} % Espacement léger pour aérer la liste
    
    \item \textbf{Difficulté 2 : Latence des notifications via Docker}
    \begin{itemize}
        \item \textbf{Problème :} Latence et stabilité du flux SSE via Docker.
        \item \textbf{Traitement :} Implémentation d'un mécanisme de 'Keep-Alive' (envoi périodique de commentaires vides) côté SseEmitter pour maintenir la connexion active.
    \end{itemize}
    \vspace{2mm}
    
    \item \textbf{Difficulté 3 : Cohérence des données IA}
    \begin{itemize}
        \item \textbf{Problème :} Décalage entre les types de données envoyés par \texttt{FastAPI} (Python) et ceux attendus par \texttt{Spring Boot} (Java).
        \item \textbf{Traitement :} Mise en place d'une couche de validation stricte via des DTOs communs pour garantir l'intégrité des analyses.
    \end{itemize}
\end{itemize}

\subsection{Analyse}


Cette phase a pour objectif de définir précisément le périmètre fonctionnel et technique du sprint, en identifiant les interactions entre les utilisateurs et les nouveaux modules de pilotage et de reporting.
\subsubsection*{1. Identification des Acteurs}
% J'utilise \subsubsection pour correspondre au niveau 4.4.2

Dans le cadre de ce sprint, les rôles se spécialisent autour de l'exploitation et de la valorisation des données :

\begin{itemize}
     
     \item \textbf{L'Administrateur RH (Acteur Principal) :} Responsable de la gestion stratégique, il utilise les outils décisionnels, gère le vivier de talents et valide les parcours de formation.
    
    \item \textbf{L'Employé (\textit{Acteur Principal}) :}tUtilisateur final qui consulte ses indicateurs de performance, reçoit des notifications et exporte son bilan de compétences officiel.

    \item \textbf{Le Système IA (Acteur Système) :} Service externe automatisé qui fournit les analyses prédictives nécessaires au fonctionnement des dashboards.
    \item \textbf{Le Service SSE (Acteur Système) :} Composant technique responsable du maintien de la connexion persistante et de la diffusion instantanée des notifications asynchrones vers les utilisateurs.
\end{itemize}

\subsubsection*{2. Clarification du besoin}

L'enjeu de ce sprint est de transformer la donnée brute en information décisionnelle. Le besoin se décline en trois axes :

\begin{itemize}
    \item \textbf{Visualisation :} Rendre les résultats de l'IA et des tests intelligibles via des graphiques dynamiques.
    \item \textbf{Interactivité :} Assurer une communication fluide et instantanée entre la direction et les employés.
    \item \textbf{Formalisation :} Produire des documents officiels (PDF) garantissant la traçabilité des compétences acquises.
\end{itemize}

\subsubsection*{3. Scénarios de validation}

Pour valider la réussite des développements, nous avons défini des scénarios clés :

\begin{itemize}
    \item \textbf{Scénario de Pilotage :} L'Admin RH consulte le \textit{dashboard} global ; le système doit afficher les KPIs agrégés après avoir sollicité le service IA.
    \item \textbf{Scénario de Communication :} L'Admin envoie un "\textit{Broadcast}" ; tous les employés connectés doivent recevoir l'alerte instantanément via SSE (Server-Sent Events).
    \item \textbf{Scénario d'Export :} L'Employé génère son "\textit{Talent Passport}" ; le système doit produire un PDF fidèle au design \textit{Bento}, incluant ses dernières analyses IA.
\end{itemize}

\subsubsection*{4. Contraintes techniques}

L'implémentation de ces besoins doit respecter plusieurs contraintes :

\begin{itemize}
    \item \textbf{Temps réel :} Utilisation du protocole \texttt{ SSE (Server-Sent Events)} pour garantir une latence minimale des notifications.
    \item \textbf{Performance :} La génération de rapports PDF complexes est gérée en dehors du cycle de vie principal (via des flux de sortie optimisés) pour ne pas bloquer les requêtes utilisateurs.
    \item \textbf{Sécurité (RBAC) :} Chaque action (export, consultation, gestion) doit être validée par le module de sécurité \texttt{JWT} selon le rôle de l'utilisateur.
\end{itemize}

\subsubsection{5. Diagramme de Cas d’Utilisation du Sprint 4}
Le diagramme présenté à la Figure~\ref{fig:usecase_sprint4} illustre les fonctionnalités clés développées au cours du Sprint 4, en mettant l'accent sur les outils de pilotage et de reporting.

\begin{figure}[htbp]
    \centering
    \includegraphics[width=0.85\textwidth]{images/usecaseSprint4.png}
    \caption{Diagramme de cas d'utilisation du Sprint 4 : pilotage et reporting}
    \label{fig:usecase_sprint4}
\end{figure}

Afin de faciliter la lecture de ce modèle, la Table~\ref{tab:legend_relations} détaille la sémantique et l'application concrète des différentes relations représentées graphiquement.

\begin{table}[H]
    \centering
    \caption{Légende des types de relations du diagramme de cas d'utilisation}
    \label{tab:legend_relations}
    \small
    \begin{tabular}{lp{5cm}p{6cm}}
        \toprule
        \textbf{Type de Relation} & \textbf{Représentation Graphique} & \textbf{Application Concrète} \\
        \midrule
        \textbf{Association} & Ligne continue simple & Indique le lien direct entre un acteur et une fonctionnalité, comme l'Administrateur RH qui accède à la cartographie des compétences. \\
        \addlinespace
        \textbf{Extension (\texttt{«extend»})} & Flèche en pointillés avec l'étiquette \texttt{«extend»} & Désigne une action optionnelle ou conditionnelle. Par exemple, l'export du passeport de compétences est une option facultative lors de la consultation du tableau de bord personnel. \\
        \addlinespace
        \textbf{Déclenchement Système} & Ligne continue vers un acteur \texttt{«System»} & Représente l'interaction automatisée entre un cas d'utilisation et un service tiers, comme la génération d'analyses par le Service IA. \\
        \bottomrule
    \end{tabular}
\end{table}
\subsubsection*{Cas 1 : Pilotage Collaborateur (User Dashboard)}
Le diagramme illustré par la Figure~\ref{fig:usecase_cas1_sprint4} présente l'organisation fonctionnelle du passeport de compétences et du tableau de bord personnel pour l'utilisateur[cite: 1].

\begin{figure}[htbp]
    \centering
    \includegraphics[width=0.85\textwidth]{images/cas1Sprint4.png}
    \caption{Diagramme de cas d'utilisation : Module Tableau de bord utilisateur et passeport de compétences}
    \label{fig:usecase_cas1_sprint4}
\end{figure}

Afin de faciliter la compréhension des interactions du module "Tableau de bord utilisateur", la Table~\ref{tab:legend_relations_cas1} détaille la sémantique des relations utilisées dans la modélisation[cite: 1].

\begin{table}[H]
    \centering
    \caption{Légende des types de relations du diagramme (Module Tableau de bord)}
    \label{tab:legend_relations_cas1}
    \small
    \begin{tabular}{lp{4.5cm}p{6.5cm}}
        \toprule
        \textbf{Type de Relation} & \textbf{Représentation Graphique} & \textbf{Application Concrète} \\
        \midrule
        \textbf{Association} & Ligne continue simple & Indique l'accès de l'**Employé** au cas d'utilisation principal "Consulter le tableau de bord personnel"[cite: 1]. \\
        \addlinespace
        \textbf{Inclusion (\texttt{«include»})} & Flèche en pointillés avec l'étiquette \texttt{«include»} & Représente une étape obligatoire. Ici, la consultation du tableau de bord nécessite impérativement de **S'authentifier**[cite: 1]. \\
        \addlinespace
        \textbf{Extension (\texttt{«extend»})} & Flèche en pointillés avec l'étiquette \texttt{«extend»} & Désigne des fonctionnalités optionnelles activables depuis le tableau de bord, comme le **téléchargement du passeport de compétences**[cite: 1]. \\
        \addlinespace
        \textbf{Interaction Système} & Ligne continue vers \texttt{«System»} & Indique la sollicitation du **Service IA** pour générer des analyses ou des recommandations de carrière[cite: 1]. \\
        \bottomrule
    \end{tabular}
\end{table}



\subsubsection*{CAS 2: Pilotage Stratégique (Executive Analytics)}

Le diagramme illustré par la Figure~\ref{fig:usecase_cas2_sprint4} présente les fonctionnalités de pilotage stratégique et d'analyses RH destinées à l'administration[cite: 1].

\begin{figure}[H]
    \centering
    \includegraphics[width=0.85\textwidth]{images/cas2Sprint4.png}
    \caption{Diagramme de cas d'utilisation : pilotage stratégique et analytics RH}
    \label{fig:usecase_cas2_sprint4}
\end{figure}

Pour une meilleure compréhension des interactions modélisées dans ce module, la Table~\ref{tab:legend_relations_cas2} définit les types de relations utilisés[cite: 1].

\begin{table}[H]
    \centering
    \caption{Légende des types de relations du diagramme (Pilotage et analyses RH)}
    \label{tab:legend_relations_cas2}
    \small
    \begin{tabular}{lp{4.5cm}p{6.5cm}}
        \toprule
        \textbf{Type de Relation} & \textbf{Représentation Graphique} & \textbf{Application Concrète} \\
        \midrule
        \textbf{Association} & Ligne continue simple & Relie l'**Administrateur RH** au cas d'utilisation central "Consulter le tableau de bord exécutif RH"[cite: 1]. \\
        \addlinespace
        \textbf{Inclusion (\texttt{«include»})} & Flèche en pointillés avec l'étiquette \texttt{«include»} & Indique que la consultation du tableau de bord nécessite obligatoirement de **S'authentifier**[cite: 1]. \\
        \addlinespace
        \textbf{Extension (\texttt{«extend»})} & Flèche en pointillés avec l'étiquette \texttt{«extend»} & Représente des analyses optionnelles comme l'**identification des talents** ou la **visualisation de la cartographie des compétences**[cite: 1]. \\
        \addlinespace
        \textbf{Interaction Système} & Ligne continue vers \texttt{«System»} & Marque la participation du **Service IA** pour la génération d'analyses globales automatisées[cite: 1]. \\
        \bottomrule
    \end{tabular}
\end{table}

\subsubsection*{Cas 3 : Hub de Communication (Real-time Hub)}
Le diagramme illustré par la Figure~\ref{fig:usecase_cas3_sprint4} présente l'organisation fonctionnelle du module de communication et de gestion des notifications en temps réel.

\begin{figure}[H]
    \centering
    \includegraphics[width=0.85\textwidth]{images/cas3Sprint4.png}
    \caption{Diagramme de cas d'utilisation : communication et notifications temps réel}
    \label{fig:usecase_cas3_sprint4}
\end{figure}

Pour une meilleure compréhension des flux et des dépendances de ce module, la Table~\ref{tab:legend_relations_cas3} détaille la sémantique des relations utilisées.

\begin{table}[H]
    \centering
    \caption{Légende des types de relations du diagramme (Communication et notifications)}
    \label{tab:legend_relations_cas3}
    \small
    \begin{tabular}{lp{4.5cm}p{6.5cm}}
        \toprule
        \textbf{Type de Relation} & \textbf{Représentation Graphique} & \textbf{Application Concrète} \\
        \midrule
        \textbf{Association} & Ligne continue simple & Relie l'**Employé** et l'**Administrateur RH** à leurs actions respectives (ex: Consulter le centre de notifications, Diffuser une annonce). \\
        \addlinespace
        \textbf{Inclusion (\texttt{«include»})} & Flèche en pointillés avec l'étiquette \texttt{«include»} & Indique une dépendance obligatoire. Par exemple, la consultation du centre requiert de **S'authentifier**, tandis que la diffusion d'une annonce entraîne obligatoirement l'action **Recevoir des alertes instantanées**. \\
        \addlinespace
        \textbf{Extension (\texttt{«extend»})} & Flèche en pointillés avec l'étiquette \texttt{«extend»} & Représente des actions complémentaires et facultatives lors de la navigation, comme **Filtrer les notifications** ou **Marquer une notification comme lue**. \\
        \addlinespace
        \textbf{Interaction Système} & Ligne continue vers \texttt{«System»} & Indique la connexion directe avec le **Service de notification temps réel** pour expédier et distribuer les alertes de manière automatisée. \\
        \bottomrule
    \end{tabular}
\end{table}
\subsubsection{Cas 4: Gouvernance \& Administration RH}
Le diagramme exposé à la Figure~\ref{fig:usecase_cas4_sprint4} détaille les fonctionnalités dédiées à la gouvernance et à l'administration du système par le responsable RH[cite: 1].

\begin{figure}[H]
    \centering
    \includegraphics[width=1\textwidth]{images/cas4Sprint4.png}
    \caption{Diagramme de cas d'utilisation : Gouvernance et administration RH}
    \label{fig:usecase_cas4_sprint4}
\end{figure}

La Table~\ref{tab:legend_relations_cas4} explicite les types de relations structurant ce module d'administration[cite: 1].

\begin{table}[H]
    \centering
    \caption{Légende des types de relations du diagramme (Gouvernance et administration)}
    \label{tab:legend_relations_cas4}
    \small
    \begin{tabular}{lp{4.5cm}p{6.5cm}}
        \toprule
        \textbf{Type de Relation} & \textbf{Représentation Graphique} & \textbf{Application Concrète} \\
        \midrule
        \textbf{Association} & Ligne continue simple & Relie l'**Administrateur RH** à la fonctionnalité racine "Gérer le vivier d'employés"[cite: 1]. \\
        \addlinespace
        \textbf{Inclusion (\texttt{«include»})} & Flèche en pointillés avec l'étiquette \texttt{«include»} & Signifie que la gestion du vivier impose impérativement de **S'authentifier** au préalable[cite: 1]. \\
        \addlinespace
        \textbf{Extension (\texttt{«extend»})} & Flèche en pointillés avec l'étiquette \texttt{«extend»} & Indique des actions spécifiques déclenchées selon les besoins, telles que **Gérer les rôles (RBAC)**, **Supprimer un compte** ou **Consulter le profil public détaillé**[cite: 1]. \\
        \bottomrule
    \end{tabular}
\end{table}


\subsection{Conception}

% La numérotation 3.2 s'adaptera automatiquement

Cette phase définit les choix techniques et architecturaux permettant de répondre aux besoins analysés précédemment. L'objectif est de garantir la robustesse du système et la fluidité des échanges de données.

\subsubsection*{1. Objectif de conception}

L'objectif principal de la conception pour ce sprint est la \textit{scalabilité} et l'intégration. Il s'agit de s'assurer que :

\begin{itemize}
    \item Le moteur de \textit{reporting} peut générer des fichiers lourds sans impacter les performances.
    \item Le système de notifications peut gérer des centaines de connexions simultanées via WebSockets.
    \item L'architecture permet une séparation nette entre la logique métier (Java) et l'intelligence artificielle (Python).
\end{itemize}
\subsubsection*{2. Architecture Technique}
Le système repose sur une architecture N-Tier (Multicouche), renforcée par une approche Micro-services hybride.
\begin{itemize}
    \item \textbf{Frontend :} Framework \texttt{Angular 17}, gérant l'état et la visualisation \textit{Bento Grid}.
    \item \textbf{Backend (Cœur) :} \texttt{Spring Boot 3}, agissant comme un chef d'orchestre (\textit{Gateway}) pour la sécurité et les données.
    \item \textbf{Service IA :} \texttt{FastAPI} (Python), dédié exclusivement aux calculs intensifs et aux modèles prédictifs.
    \item \textbf{Communication Temps Réel :} Protocole \texttt{SEE} via HTTP pour la gestion des notifications.
\end{itemize}

\subsubsection*{3. Algorithme de Détection \textit{HiPo} (\textit{High Potential})}

Pour automatiser l'identification des talents clés, nous avons implémenté un algorithme \textit{"Data-Driven"} basé sur la \textbf{Matrice 9-Box}, croisant les deux axes fondamentaux de l'évaluation RH :

\begin{itemize}
    \item \textbf{Performance ($S_{perf}$) :} Calculée via la moyenne arithmétique globale (sur 100) de l'ensemble des tests techniques passés par le candidat sur la plateforme.
    \item \textbf{Potentiel ($S_{pot}$) :} Évalué via la normalisation des scores comportementaux issus du test de personnalité \textit{PCM} (\textit{Process Communication Model}) ou, à défaut, des prédictions de l'IA.
\end{itemize}

L'algorithme génère un score \textit{HiPo} global pondéré selon la formule suivante :

\begin{equation}
    Score_{HiPo} = (S_{perf} \times 0.6) + (S_{pot} \times 0.4)
\end{equation}

Enfin, des seuils mathématiques placent automatiquement le collaborateur dans l'un des 9 quadrants de la matrice. Par exemple, un profil respectant la condition :

\begin{equation}
    (S_{perf} \ge 80\%) \cap (S_{pot} \ge 80\%)
\end{equation}

est immédiatement labellisé \textbf{"High Potential"} (\textit{Top Right}), permettant aux décideurs d'identifier les futurs leaders de manière objective et sans biais cognitifs. Cette classification automatisée facilite le pilotage stratégique du capital humain au sein de \textbf{TalentPredict}.

\subsubsection*{4. Modèle de données }
% --- Introduction ---
La gestion optimisée des ressources humaines et le développement des compétences constituent des piliers stratégiques pour la performance d'une organisation[cite: 1]. Ce modèle de \textbf{Gestion du Capital Humain} centralise les données relatives aux collaborateurs, à leurs compétences, et à leur évolution professionnelle[cite: 1]. L'objectif est d'aligner les aspirations individuelles avec les besoins stratégiques de l'entreprise via une analyse de données précise[cite: 1].

% --- Figure ---
\begin{figure}[h!]
    \centering
    \includegraphics[width=1.0\textwidth]{images/Erd4.png} 
    \caption{Diagramme Entité-Relation du Système de Gestion du Capital Humain}
    \label{fig:erd_diagram}
\end{figure}

Le schéma présenté dans la Figure \ref{fig:erd_diagram} s'articule autour des entités suivantes[cite: 1] :

\begin{itemize}
    \item \textbf{UTILISATEUR \& EMPLOYE} : Gèrent les accès sécurisés (\texttt{email}, \texttt{role}) et les informations identitaires du personnel (\texttt{nom}, \texttt{poste\_actuel})[cite: 1].
    \item \textbf{COMPETENCE \& CERTIFICATION} : Suivent le bagage technique des employés via les tables de jointure \texttt{EMPLOYE\_COMPETENCE} et \texttt{EMPLOYE\_CERTIFICATION}, incluant le \texttt{niveau\_maitrise} et la \texttt{date\_obtention}[cite: 1].
    \item \textbf{HISTORIQUE\_TEST \& PROFIL\_PSYCHOLOGIQUE} : Archivent les performances aux évaluations et les traits de personnalité stockés en format \texttt{json} pour identifier les hauts potentiels[cite: 1].
    \item \textbf{RECOMMANDATION\_CARRIERE} : Propose des évolutions vers un \texttt{POSTE\_CIBLE} en calculant un \texttt{score\_adequation} basé sur le profil de l'employé[cite: 1].
    \item \textbf{INSCRIPTION\_FORMATION} : Supervise la montée en compétences en suivant le \texttt{statut} et la \texttt{progression\_pourcentage} des modules suivis[cite: 1].
    \item \textbf{NOTIFICATION} : Assure la communication système via des alertes et relances entre les utilisateurs[cite: 1].
\end{itemize}

\subsubsection*{5. Contrats API (Exemple : Dashboard)}

Pour garantir l'interopérabilité, nous utilisons des DTO (\textit{Data Transfer Objects}). Exemple de réponse JSON pour le \textit{Dashboard} Admin :

\begin{lstlisting}[language=json, caption={Exemple de structure JSON pour le Dashboard Admin}, label={lst:json_dashboard}]
{
  "totalEmployees": 150,
  "highPotentialCount": 12,
  "skillsDistribution": { 
      "Java": 85, 
      "Python": 70, 
      "BehavioralSkills": 90 
  },
  "recentActivities": [...]
}
\end{lstlisting}

\subsubsection{6. Diagrammes de Séquence}

Cette section détaille la chorégraphie des interactions entre les différents composants de la plateforme à travers quatre scénarios clés du Sprint 4.

\subsubsection{Cas n°1 : Pilotage Collaborateur \& Passeport Talent}
Pour assurer la transparence du parcours professionnel de l'utilisateur, le diagramme de séquence de la \autoref{fig:seq1Sprint4} présente le processus de consultation du tableau de bord et les étapes de génération automatisée du passeport talent au format PDF.


\begin{figure}[H]
    \centering
    \includegraphics[width=0.9\textwidth]{images/seq1Sprint4.png}
    \caption{Diagramme de séquence : génération du passeport talent}
    \label{fig:seq1Sprint4}
\end{figure}
Le diagramme de la \autoref{fig:seq1Sprint4} illustre le flux opérationnel structuré en deux étapes majeures :
\subsubsection{Phase 1 : Consultation du Tableau de Bord}
Cette phase assure la visualisation sécurisée des indicateurs de performance de l'employé.
\begin{itemize}
    \item \textbf{Sécurité JWT :} Le système valide systématiquement le jeton d'accès via \textit{Spring Security} avant de traiter la requête (erreur 401 en cas d'échec).
    \item \textbf{Intelligence Artificielle :} Une analyse prédictive est déclenchée auprès du \textit{microservice IA} (si non cachée) pour enrichir le profil de l'utilisateur avec des recommandations personnalisées.
\end{itemize}

\subsubsection{Phase 2 : Exportation du Passeport Talent}
Cette étape gère la production d'un justificatif officiel au format numérique.
\begin{itemize}
    \item \textbf{Orchestration métier :} Après vérification des droits d'accès (erreur 403 si non autorisé), le service récupère l'intégralité du profil consolidé depuis la base \textit{PostgreSQL}.
    \item \textbf{Génération PDF :} Le service technique convertit dynamiquement le modèle HTML en un flux PDF sécurisé pour permettre le téléchargement final par l'utilisateur.
\end{itemize}

\subsubsection{Cas n°2 : Pilotage Stratégique \& Audit RH}
Afin de garantir une gouvernance optimale des données au sein de la plateforme TalentPredict, le diagramme de séquence de la \autoref{fig:seq2Sprint4} détaille les mécanismes de pilotage stratégique et les procédures d'audit réservés à l'Administrateur RH.
\begin{figure}[H]
    \centering
    \includegraphics[width=0.85\textwidth]{images/seq2Sprint4.png}
    \caption{Diagramme de séquence : pilotage stratégique et audit RH}
    \label{fig:seq2Sprint4}
\end{figure}

Le diagramme de la \autoref{fig:seq2Sprint4} détaille les interactions techniques structurées en deux phases majeures :

\subsubsection{Phase 1 : Génération des KPIs de Business Intelligence}
Cette phase permet la centralisation et la visualisation des indicateurs de performance globaux.
\begin{itemize}
    \item \textbf{Contrôle RBAC (Role-Based Access Control) :} En plus de la validation du jeton JWT, le filtre de sécurité vérifie de manière stricte le rôle de l'utilisateur. Tout accès non-administrateur est immédiatement rejeté avec une erreur 403 (Droits insuffisants).
    \item \textbf{Détection des hauts potentiels :} Le \textit{Service d'analyse} agrège les indicateurs de compétences depuis \textit{PostgreSQL} et sollicite optionnellement le \textit{Service IA} pour identifier les talents émergents et analyser les tendances globales, avant de retourner un DTO de synthèse JSON à l'interface.
\end{itemize}

\subsubsection{Phase 2 : Génération du Rapport d'Audit RH}
Cette étape prend en charge l'exportation massive et sécurisée des données de l'entreprise.
\begin{itemize}
    \item \textbf{Compilation du rapport :} Le \textit{Service d'export global} extrait les statistiques consolidées et orchestre la génération du document d'audit complet au format PDF (gérant une erreur 500 en cas d'échec technique).
    \item \textbf{Journalisation de sécurité :} Avant de délivrer le flux du fichier, le système exécute une action de traçabilité en enregistrant l'export dans les logs d'audit de la base \textit{PostgreSQL}, assurant ainsi une gouvernance stricte des données RH.
\end{itemize}
\subsubsection{Cas n°3 : Communication Hub \& Server-Sent Events (SSE)}
Afin d'assurer une réactivité immédiate aux événements du système, le diagramme de séquence de la \autoref{fig:seq3Sprint4} expose l'architecture de communication unidirectionnelle permettant l'envoi de notifications en temps réel aux utilisateurs.
\begin{figure}[H]
    \centering
    \includegraphics[width=0.65\textwidth]{images/seq3Sprint4.png}
    \caption{Diagramme de séquence : flux de notification temps réel}
    \label{fig:seq3Sprint4}
\end{figure}

Le diagramme de la \autoref{fig:seq3Sprint4} détaille la gestion du cycle de vie d'une connexion persistante, structurée en quatre phases :

\subsubsection{Phase 1 : Initialisation et Handshake}
\begin{itemize}
    \item \textbf{Établissement du flux :} L'interface Angular initie une connexion \textit{Server-Sent Events} (SSE) en transmettant le jeton JWT pour validation.
    \item \textbf{Persistance :} Une fois authentifié, le backend crée un émetteur SSE dédié et enregistre la session comme active pour l'utilisateur.
\end{itemize}

\subsubsection{Phase 2 : Maintien de la connexion (Keep-Alive)}
\begin{itemize}
    \item \textbf{Heartbeat :} Un signal de maintien est envoyé toutes les 30 secondes pour prévenir les coupures de connexion par les proxys réseau ou les conteneurs Docker.
\end{itemize}

\subsubsection{Phase 3 : Déclenchement et envoi d'alerte}
\begin{itemize}
    \item \textbf{Diffusion en ligne :} Lorsqu'un événement métier survient (ex: validation de formation par l'Administrateur), l'alerte est immédiatement transmise via le flux SSE si l'utilisateur est en ligne.
    \item \textbf{Stockage hors ligne :} Si aucune session active n'est trouvée, la notification est persistée dans le \textit{Dépôt des notifications} pour une consultation ultérieure.
\end{itemize}

\subsubsection{Phase 4 : Gestion de la déconnexion et expiration}
\begin{itemize}
    \item \textbf{Nettoyage des ressources :} Que ce soit par une fermeture volontaire de l'onglet ou par l'expiration du délai (Timeout), le système détecte la rupture, détruit l'émetteur SSE et libère la session active.
    \item \textbf{Résilience :} Le frontend intègre un mécanisme de reconnexion automatique en cas de perte de flux non sollicitée.
\end{itemize}
\subsubsection{Cas n°4 : Gouvernance \& Administration RH}
 Afin de permettre une supervision fine des collaborateurs, le diagramme de séquence de la \autoref{fig:seq4Sprint4} illustre les processus de consultation des profils publics du vivier ainsi que les actions administratives de gestion d'état des comptes.
 \begin{figure}[H]
    \centering
    \includegraphics[width=0.65\textwidth]{images/seq4Sprint4.png}
    \caption{Diagramme de séquence : gestion du vivier et consultation des profils publics}
    \label{fig:seq4Sprint4}
\end{figure}

Le diagramme de la \autoref{fig:seq4Sprint4} se divise en deux phases opérationnelles distinctes :

\subsubsection{Phase 1 : Consultation du Vivier et Profil Public}
\begin{itemize}
    \item \textbf{Accès sécurisé :} La consultation de la liste des employés et des détails de leur profil public (liens sociaux, biographie) est conditionnée par une validation stricte du jeton JWT et des droits RBAC.
    \item \textbf{Récupération de données :} Le \textit{Service des employés} agrège les informations depuis \textit{PostgreSQL} pour fournir un DTO complet à l'interface.
\end{itemize}

\subsubsection{Phase 2 : Gestion de l'état (Désactivation du compte)}
\begin{itemize}
    \item \textbf{Action administrative :} L'administrateur peut désactiver un compte, déclenchant une mise à jour du statut en base de données.
    \item \textbf{Révocation des accès :} En cas de succès, le système invalide immédiatement les droits futurs de l'employé et révoque toutes les sessions actives pour garantir la sécurité.
    \item \textbf{Notification :} Une confirmation de l'opération est envoyée au système, et le changement de statut est notifié visuellement à l'administrateur.
\end{itemize}
\subsubsection{7. Diagramme de Class }
Le diagramme présenté en Figure \ref{fig:architecture_talentpredict} illustre l'architecture technique du système \textit{TalentPredict AI}. Cette structure est conçue selon une approche modulaire, séparant clairement les responsabilités fonctionnelles des entités de données centrales. 

L'architecture repose sur une séparation en couches (Controllers, Services, Repositories) et intègre une communication asynchrone pour les notifications ainsi qu'une interaction avec un microservice externe spécialisé en IA.

\begin{figure}[htbp]
    \centering
    \includegraphics[width=1.0\textwidth]{images/classSprint4.png}
    \caption{Diagramme de classes technique du système TalentPredict AI.}
    \label{fig:architecture_talentpredict}
\end{figure}

\subsection{Description du diagramme}
Le diagramme met en évidence trois piliers fondamentaux :
\begin{itemize}
    \item \textbf{Le modèle de données central :} Il regroupe les entités métiers persistantes (ex: \texttt{User}, \texttt{Skill}, \texttt{Prediction}) qui forment le socle d'information partagé entre les différents modules.
    \item \textbf{La modularité fonctionnelle :} Le système est segmenté en quatre modules principaux (RH/Apprentissage, Analyses, Export, Communication), garantissant une maintenance aisée et une scalabilité du code.
    \item \textbf{L'interopérabilité :} La présence du \texttt{AiMicroserviceClient} démontre l'intégration d'un composant externe en Python/FastAPI, permettant au backend principal de bénéficier de capacités de prédiction avancées via des échanges REST.
\end{itemize}

\subsubsection*{8. Architecture Temps Réel via \texttt{SSE} (\textit{Server-Sent Events})}
Afin de garantir une communication fluide et instantanée, \textbf{TalentPredict} intègre une architecture temps réel basée sur le protocole \texttt{SSE}. Ce choix technologique, plus léger et robuste que les \texttt{WebSockets} pour un flux unidirectionnel (\textit{Serveur vers Client}), est parfaitement adapté à la diffusion de notifications et d'alertes sans rechargement de page.

Le \textit{backend} \texttt{Spring Boot} agit comme un émetteur d’événements grâce à l’\texttt{API} \texttt{SseEmitter}. Après authentification, les utilisateurs établissent une connexion persistante pour recevoir des mises à jour asynchrones.

\paragraph*{Composants de l'architecture :}
L’architecture repose sur les modules suivants :
\begin{itemize}
    \item \textbf{\texttt{NotificationSseService} :} Gère le cycle de vie des connexions (\texttt{SseEmitter}) et le registre des utilisateurs connectés via une structure \texttt{ConcurrentHashMap} sécurisée.
    \item \textbf{\texttt{NotificationController} :} Expose l’\textit{endpoint} de souscription \texttt{api}/\texttt{notifications}/\texttt{subscribe} permettant au client d'ouvrir le canal de communication.
    \item \textbf{Système Événementiel :} Centralise les notifications métiers (nouvelle recommandation IA, succès d’un test, validation RH, génération du \textit{Talent Passport}).
    \item \textbf{Angular \textit{EventSource Client} :} Utilise l’\texttt{API} native du navigateur pour maintenir la connexion ouverte et rafraîchir dynamiquement les composants \texttt{Angular} via \texttt{RxJS} à la réception d’un événement.
\end{itemize}

\begin{figure}[H]
    \centering
    % Remplacez 'nom_du_fichier' par le nom réel de votre image (ex: flux_sse.png)
    \includegraphics[width=0.85\textwidth]{images/DiagrammedeComposants.png}
    \caption{Architecture événementielle et gestion du registre des connexions  \texttt{SSE}}
    \label{fig:flux_communication_sse}
\end{figure}

Le schéma ci-dessus illustre le flux de communication temps réel au sein de \textbf{TalentPredict}. Le \texttt{NotificationController} réceptionne les demandes de souscription et délègue au \texttt{NotificationSseService} la création d'un canal persistant. 

Ce dernier stocke l'instance de la connexion dans une \texttt{ConcurrentHashMap} (mémoire vive), garantissant un accès \textit{thread-safe} et performant. Lors du déclenchement d'un événement métier, le système identifie le destinataire et "pousse" l'information instantanément vers le client \texttt{Angular}.

\paragraph*{Avantages techniques retenus :}
\begin{itemize}
    \item \textbf{Reconnexion Automatique :} Le protocole \texttt{SSE} gère nativement la reprise de connexion en cas de coupure réseau.
    \item \textbf{Efficacité énergétique :} Consomme moins de ressources côté client et serveur qu’une connexion \texttt{WebSocket} bidirectionnelle.
    \item \textbf{Compatibilité Proxy/Firewall :} Contrairement aux \texttt{WebSockets}, le \texttt{SSE} utilise le protocole \texttt{HTTP} standard, évitant ainsi les problèmes de blocage par les pare-feux d'entreprise ou les \textit{proxys} \texttt{Docker}.
\end{itemize}






\subsection{ Codage}
% Adaptez le niveau (section, subsection, subsubsection) selon votre plan.

Cette sous-section détaille l'étape de codage spécifique au quatrième sprint. L'infrastructure globale de l'application étant déjà en place, nos efforts de développement se sont concentrés sur l'intégration de bibliothèques spécialisées pour le rendu documentaire et la communication asynchrone.

\paragraph*{1. Implémentation du Moteur de Reporting (Backend)}
Pour répondre au besoin de dématérialisation (\textit{Passeport Talent} et Rapports RH), nous avons mis en place un service de génération de documents côté serveur (\textit{Server-Side Rendering}) :

\begin{itemize}
    \item \textbf{\texttt{OpenHTMLToPDF} (v1.0.10) :} Il s'agit de la bibliothèque centrale de ce sprint. Elle a été intégrée pour \textit{parser} nos gabarits dynamiques (HTML5/CSS3) et les transformer en documents PDF professionnels. Ce choix a été motivé par sa grande fidélité de rendu des standards Web récents.
    \item \textbf{\texttt{Apache PDFBox} :} Utilisée en complément pour la manipulation de bas niveau des flux générés (gestion des métadonnées du fichier et sécurisation du PDF).
    \item \textbf{\texttt{Lombok} :} Déployée systématiquement sur nos nouveaux objets de transfert (DTOs) de \textit{reporting} pour réduire le code \textit{boilerplate} et garantir la clarté des classes métiers.
    \item \textbf{Thymeleaf :} Choisi comme moteur de gabarits (Templating Engine) côté serveur pour injecter dynamiquement les métadonnées (nom de l'employé, scores IA, radar) dans les templates HTML avant leur transformation finale en documents PDF.
\end{itemize}
\subsubsection*{2. Génération Dynamique des Rapports PDF}

La génération documentaire constitue l’une des fonctionnalités majeures du Sprint 4. Le \textit{backend} \texttt{Spring Boot} génère dynamiquement des documents PDF professionnels à partir des données analytiques produites par les modules IA.

Deux types de documents clés ont été implémentés :
\begin{itemize}
    \item \textbf{\textit{Talent Passport} :} Une synthèse individuelle exhaustive regroupant les compétences extraites, les scores techniques, le profil comportemental PCM et les recommandations de carrière de l'IA.
    \item \textbf{Rapport RH Global :} Une vue consolidée destinée à la direction, présentant les statistiques et les tendances macroscopiques du capital humain.
\end{itemize}

\paragraph*{Mécanisme de production et persistance}
Le moteur de génération repose sur une architecture orientée \textit{Templates} (\texttt{Thymeleaf}) permettant de produire des documents cohérents, structurés et respectant la charte graphique de la plateforme. Une attention particulière a été portée à la persistance : les fichiers générés sont automatiquement stockés dans un volume \texttt{Docker} dédié (\texttt{app\_uploads}), assurant leur disponibilité et leur archivage même après le redémarrage des services applicatifs.


\subsubsection*{3. Hub de Communication et Tableaux de Bord }
Pour transformer la plateforme en un outil de pilotage interactif, nous avons implémenté des mécanismes de visualisation et de mise à jour en temps réel :

\begin{itemize}
    \item \textbf{\texttt{Spring SseEmitter} \& \texttt{RxJS} :} Ce duo technique a permis l'intégration du protocole SSE (Server-Sent Events). Côté serveur, \texttt{Spring Boot} gère la diffusion des messages. Côté client (\texttt{Angular}), \texttt{RxJS} orchestre les flux de données asynchrones, permettant de "pousser" les notifications aux utilisateurs sans aucun rafraîchissement manuel de l'interface.
    \item \textbf{\texttt{Ngx-Toastr} :} Consomme les flux SSE pour afficher des alertes visuelles élégantes et non intrusives (ex. "Nouvelle formation approuvée").
    \item \textbf{\texttt{Chart.js} :} Bibliothèque de \textit{data-visualisation} intégrée pour générer dynamiquement les radars de compétences de l'employé et les graphiques de tendances macroscopiques de l'\textit{Executive Dashboard}.
    \item \textbf{\texttt{Lucide-Angular} \& \texttt{FontAwesome} :} Packs vectoriels utilisés pour standardiser l'iconographie des nouveaux modules (boutons d'export, indicateurs d'état).
\end{itemize}

\subsubsection*{4. Outillage DevOps et Validation API}
Afin de garantir la stabilité de ces nouvelles fonctionnalités complexes avant leur déploiement, nous avons maintenu une rigueur stricte sur l'outillage :

\begin{itemize}
    \item \textbf{\texttt{Postman} :} Utilisé pour documenter et tester unitairement les nouveaux points d'accès REST complexes, notamment les flux binaires (téléchargement de PDF) et les retours de l'IA, avant leur intégration dans le \textit{frontend}.
    \item \textbf{Communication Inter-services :} Mise en œuvre de clients REST robustes (Spring RestTemplate) configurés avec des mécanismes de Timeout pour assurer une communication fluide et résiliente entre le Backend Java et le micro-service d'Intelligence Artificielle (FastAPI).
    \item \textbf{\texttt{Docker} \& \texttt{Docker Compose} :}la configuration du flux persistant pour le SSE a nécessité la mise à jour de nos fichiers de conteneurisation pour garantir un déploiement fluide de ce nouvel écosystème.
    \item \textbf{\texttt{GitLab CI/CD} :} Le pipeline d'intégration continue a été maintenu pour exécuter automatiquement les tests sur ces nouveaux modules à chaque validation de code (\textit{Commit}).
\end{itemize}
\subsection{Realisation}
% Ajustez le numéro de section au besoin (ex: remplacez \subsection par \section si c'est le titre principal 4.X)

L'aboutissement de ce quatrième sprint se matérialise par la livraison d'interfaces interactives et de modules de restitution documentaire. Le \textit{Frontend}, développé en \texttt{Angular 17}, communique de manière asynchrone avec nos différents microservices pour offrir une expérience utilisateur fluide et en temps réel.

\subsubsection*{1. Pilotage Stratégique : Executive Dashboard et Gestion RH}
Le cœur de ce sprint consistait à fournir aux administrateurs et décideurs RH une vue macroscopique et exploitable du capital humain de l'entreprise.


    \subsubsection*{ L'\textit{Executive Dashboard} (Vue d'ensemble) :} Cette interface centralise les indicateurs de performance (KPIs) en temps réel. Grâce à l'intégration de la bibliothèque \texttt{Chart.js}, l'administrateur peut visualiser les tendances d'évolution des compétences, le suivi de l'\textit{Onboarding}, ainsi que les alertes critiques (ex. : taux de couverture des tests insuffisant). L'IA intervient également ici pour mettre en évidence les "Talents à haut potentiel" (HiPo).


\begin{figure}[H]
    \centering
    \includegraphics[width=0.9\textwidth]{images/dashboardAdmin.png} % Remplacez par le nom de l'Image 1
    \vspace{2mm}
    \caption{\textit{Executive Dashboard} : Vue macroscopique et indicateurs clés}
    \label{fig:executive_dashboard}
\end{figure}

    \subsubsection*{Gestion des Employés et Profil Public :} Ce module permet une administration complète (CRUD) des collaborateurs. La liste interactive intègre des filtres dynamiques (statuts, rôles, départements). Un panneau latéral (\textit{Slide-over}) offre un aperçu rapide du profil sélectionné, affichant son volume de formations, ses prédictions IA et permettant de vérifier la validité de ses liens sociaux (LinkedIn, GitHub) sans quitter la page principale.


% --- Figure 1 : Aperçu du profil public ---
\begin{figure}[H]
    \centering
    \includegraphics[width=1\textwidth]{images/publicprofileUser.png}
    \caption{Interface de l'aperçu rapide du profil collaborateur}
    \label{fig:profil_public}
\end{figure}

% --- Figure 2 : Gestion des collaborateurs ---
\begin{figure}[H]
    \centering
    \includegraphics[width=1\textwidth]{images/gestiondeEmployee.png}
    \caption{Interface de gestion globale des effectifs (Vue Administrateur)}
    \label{fig:gestion_employes}
\end{figure}

\subsubsection*{2. Espace Collaborateur et Profiling IA}
Du côté de l'employé, l'interface a été enrichie pour restituer les analyses prédictives générées par notre moteur \texttt{FastAPI}.

    \subsubsection*{Tableau de bord de l'Employé :} Le collaborateur accède à son espace personnel ("\textit{Welcome back}") où figure son "\textit{AI Persona}" (ex. : Profil Stratégique). L'interface affiche une matrice de distribution des compétences (\textit{Soft \& Compétences techniques}), son inertie d'apprentissage (\textit{Momentum}), et surtout, les parcours de carrière recommandés (\textit{Career Pathways}) avec un pourcentage de correspondance calculé par l'intelligence artificielle.


\begin{figure}[H]
    \centering
    \includegraphics[width=0.9\textwidth]{images/dashboardUser.png} % Remplacez par le nom de l'Image 3
    \vspace{2mm}
    \caption{Tableau de bord collaborateur avec restitution des analyses IA}
    \label{fig:dashboard_employe}
\end{figure}

\subsubsection*{3. Interaction Temps Réel : Communication Hub et Notifications}
Pour transformer la plateforme en un écosystème vivant, nous avons implémenté des mécanismes de communication instantanée basés sur les \texttt{SEE}.


    \subsubsection*{Le \textit{Hub} de Communication :} Réservé aux administrateurs, cet espace permet de composer et d'envoyer des messages ciblés (Informations, Succès, Alertes) à des collaborateurs spécifiques, avec une option de relai par e-mail.

\begin{figure}[H]
    \centering
    \includegraphics[width=0.85\textwidth]{images/communicationHUB.png}
    \caption{Interface du Hub de communication pour l'envoi de messages}
    \label{fig:communication_hub}
\end{figure}
    \subsubsection*{Centre de Notifications (SSE) :} Côté utilisateur, un volet de notifications interactif capte les événements poussés par le serveur via le protocole \texttt{SSE} (ex. : "Demande de formation approuvée", "Nouveau test assigné"). L'état non lu/lu est géré dynamiquement sans aucun rafraîchissement de la page.


\begin{figure}[H]
    \centering
    \includegraphics[width=0.6\textwidth]{images/NoificationUser.png}
    \caption{Volet de notifications en temps réel pour l'utilisateur}
    \label{fig:notifications_user}
\end{figure}

\subsubsection*{4. Restitution Documentaire : Le Moteur de Reporting (PDF)}
La dématérialisation et l'exportation sécurisée des données ont été assurées par la bibliothèque \texttt{OpenHTMLToPDF} côté serveur.

    \subsubsection*{Le Passeport Talent:} Destiné à l'employé, ce document généré à la volée compile ses scores, son analyse IA détaillée, les recommandations de développement (\textit{Soft \& Tech}) et son historique de formation. La mise en page respecte une charte graphique professionnelle, idéale pour des entretiens annuels.

    \begin{figure}[H]
    \centering
    \includegraphics[width=0.85\textwidth]{images/talentPassport.png}
    \caption{Passeport Talent : Document PDF généré dynamiquement pour l'employé}
    \label{fig:talent_passport}
    \end{figure}
    \subsubsection*{Le Rapport Analytique RH :} Export destiné à la direction, synthétisant les données globales du personnel (département, nombre de tests, statut d'activité) sous forme de tableau paginé.


-
\begin{figure}[H]
    \centering
    \includegraphics[width=0.65\textwidth]{images/rapportRH.png}
    \caption{Rapport RH : Synthèse décisionnelle générée par le backend}
    \label{fig:rapport_rh}
\end{figure}
\subsection{Test et Déploiement }
% Ajustez la numérotation si nécessaire

Cette phase ultime constitue le point d'orgue de notre cycle de développement. Elle garantit que la plateforme \textbf{TalentPredict} est non seulement fonctionnelle d'un point de vue métier, mais aussi stable, sécurisée et prête à être exploitée dans un environnement de production à haute disponibilité.

\subsubsection*{1. Stratégie de Test et Assurance Qualité (QA)}

Afin de garantir une qualité logicielle irréprochable et de prévenir les régressions lors des mises à jour, nous avons mis en place une pyramide de tests rigoureuse :

\begin{itemize}
    \item \textbf{Tests Unitaires (Backend) :} En nous appuyant sur \texttt{JUnit 5} et \texttt{Mockito}, nous avons isolé et testé les composants critiques. Un effort particulier a été porté sur l'\texttt{ExportService} (vérification de la bonne agrégation des données avant injection dans le PDF) et le \texttt{DashboardService} (validation algorithmique du calcul des indicateurs globaux).
    \item \textbf{Tests d'Intégration et Validation API :} Nous avons utilisé \texttt{Postman} pour valider les contrats d'interface. Nous avons systématiquement testé la communication inter-services, notamment les requêtes asynchrones entre le \textit{Backend} \texttt{Spring Boot} et le moteur IA \texttt{FastAPI}, en simulant des scénarios de réponses complexes.
    \item \textbf{Tests de Recette Fonctionnelle et Temps Réel :} Les \textit{User Stories} ont fait l'objet d'une validation manuelle de bout en bout. Nous avons soumis le \textit{Communication Hub} (SSE) à des tests de charge simulée pour vérifier qu'aucune perte de notification n'intervenait lors de sessions concurrentes.
\end{itemize}

Afin de matérialiser notre démarche d'assurance qualité, le tableau ci-dessous présente un extrait des scénarios de tests majeurs exécutés lors de la phase de validation de ce quatrième sprint :

\begin{table}[H]
    \centering
    \renewcommand{\arraystretch}{1}
    \small
    \begin{tabularx}{\textwidth}{|c|>{\raggedright\arraybackslash}p{3cm}|>{\raggedright\arraybackslash}X|>{\raggedright\arraybackslash}X|}
        \hline
        \rowcolor{TableHeaderBlue}
        \color{white}\textbf{ID} & \color{white}\textbf{Scénario Testé} & \color{white}\textbf{Pré-conditions \& Étapes} & \color{white}\textbf{Résultat Attendu} \\
        \hline
        
        % TC-01
        TC-01 & \textbf{Génération PDF (Passeport Talent)} \newline \textit{(Test Unitaire / Backend)} 
        & \textbf{Pré-requis :} Employé avec des compétences évaluées en base. \newline
          \textbf{Action :} Appel de la méthode \texttt{generatePassport(userId)} de l'\texttt{ExportService}.
        & Génération réussie du flux binaire (\texttt{byte[]}). Le document s'ouvre sans erreur et les scores correspondent aux données de la base PostgreSQL. \\
        \hline
        
        % TC-02
        TC-02 & \textbf{Prédictions IA (API FastAPI)} \newline \textit{(Test d'Intégration)} 
        & \textbf{Pré-requis :} Microservice IA démarré. \newline
          \textbf{Action :} Envoi d'une requête \texttt{POST /predict/persona} avec un \textit{payload} JSON depuis le \textit{backend}.
        & Retour HTTP 200 (OK). Le \textit{payload} de réponse contient un "Persona" valide et des scores prédictifs correctement formatés. \\
        \hline
        
        % TC-03
        TC-03 & \textbf{Notification Temps Réel (SSE)} \newline \textit{(Test Fonctionnel de bout en bout)} 
        & \textbf{Pré-requis :} Employé et Manager RH connectés . \newline
          \textbf{Action :} Le manager clique sur "Approuver la formation".
        & L'employé reçoit instantanément une notification visuelle sur son interface Angular, sans aucun rafraîchissement manuel. \\
        \hline
        
        % TC-04
        TC-04 & \textbf{Restitution Visuelle (\textit{Executive Dashboard})} \newline \textit{(Test de Recette UI)} 
        & \textbf{Pré-requis :} Authentification réussie avec le rôle ADMIN. \newline
          \textbf{Action :} Navigation vers la vue macroscopique RH.
        & Les graphiques (\texttt{Chart.js}) se chargent sans erreur console. Les KPI globaux correspondent aux agrégations de la base de données. \\
        \hline
        
        % TC-05
        TC-05 & \textbf{Contrôle d'Accès (RBAC)} \newline \textit{(Test de Sécurité)} 
        & \textbf{Pré-requis :} Utilisateur authentifié avec le rôle EMPLOYEE. \newline
          \textbf{Action :} Tentative de téléchargement du rapport global RH (\texttt{GET /api/reporting/hr-global}).
        & Le filtre JWT intercepte l'appel. La requête est rejetée avec un code HTTP 403 (Forbidden). L'accès est bloqué. \\
        \hline
        
    \end{tabularx}
    \vspace{2mm}
    \caption{Synthèse des cas de test majeurs (Sprint 4)}
    \label{tab:test_cases_sprint4}
\end{table}

\subsubsection*{2.Tests des Flux Temps Réel (SSE)}

Afin de garantir la fiabilité des communications asynchrones au sein de \textbf{TalentPredict}, plusieurs scénarios de validation ont été réalisés sur l’architecture temps réel basée sur le protocole \texttt{SSE} (\textit{Server-Sent Events}).

Ces tests ont permis de vérifier plusieurs aspects critiques du système :

\begin{itemize}
    \item \textbf{Réception instantanée :} Les événements générés côté serveur, tels que les recommandations IA ou les validations RH, sont affichés sur le \textit{Dashboard} utilisateur avec une latence moyenne inférieure à 200 ms.
    
    \item \textbf{Synchronisation multi-session :} Un même utilisateur connecté simultanément sur plusieurs onglets ou appareils reçoit les notifications de manière cohérente et synchronisée en temps réel.
    
    \item \textbf{Résilience réseau :} Des tests ont été effectués afin d’évaluer la capacité du client \texttt{Angular} à rétablir automatiquement la connexion \texttt{SSE} après une interruption volontaire du réseau. Les résultats ont confirmé la continuité du flux événementiel sans perte de données ni désynchronisation des interfaces.
    
    \item \textbf{Sécurité et Confidentialité :} Des validations ont été menées afin de garantir que chaque flux \texttt{SSE} reste strictement associé à l’identité authentifiée de l’utilisateur via \texttt{JWT}. Cette approche empêche tout accès non autorisé aux notifications privées et renforce la confidentialité des échanges temps réel.
\end{itemize}

L’ensemble des expérimentations réalisées a confirmé la stabilité, la légèreté et la robustesse de l’architecture \texttt{SSE}, même dans des scénarios simulant plusieurs dizaines de connexions simultanées.

\subsubsection*{3. Architecture de Conteneurisation (Docker)}

Pour résoudre les problématiques de disparité d'environnements et faciliter le déploiement, nous avons intégralement conteneurisé l'écosystème :

\begin{itemize}
    \item \textbf{Dockerisation \textit{Multi-Stage} :} L'écriture de fichiers \texttt{Dockerfile} optimisés a permis de séparer les environnements de compilation des environnements d'exécution, réduisant drastiquement la taille des images finales.
    \item \textbf{Orchestration via Docker Compose :} La configuration du fichier \texttt{docker-compose.yml} orchestre le démarrage simultané de nos quatre micro-environnements au sein d'un réseau privé isolé (\texttt{talentpredict-net}) :
    \begin{itemize}
        \item \texttt{tp-frontend} : Application \texttt{Angular} servie par \texttt{Nginx}.
        \item \texttt{tp-backend} : Serveur \texttt{Spring Boot} exposant l'API REST.
        \item \texttt{tp-ai-service} : Moteur analytique \texttt{FastAPI}.
        \item \texttt{tp-database} : Instance \texttt{PostgreSQL} couplée à des volumes persistants (\texttt{tp-db-data}).
    \end{itemize}
\end{itemize}
\subsubsection*{4.Architecture Globale de Déploiement}

\begin{figure}[H]
    \centering
    % Remplacez 'chemin_architecture_globale.png' par le nom exact de votre image
    \includegraphics[width=0.95\textwidth]{images/ArchitectureGlobalede Déploiement.png}
    \vspace{2mm}
    \caption{Architecture globale de déploiement conteneurisé de la plateforme TalentPredict}
    \label{fig:archi_deploiement_globale}
\end{figure}

L'infrastructure de \textit{TalentPredict} est entièrement conteneurisée via \textbf{Docker Compose}, garantissant isolation et portabilité. L'écosystème repose sur un réseau privé \texttt{talentpredict-net} orchestrant cinq services clés :
\begin{itemize}
    \item \textbf{Frontend} : Application Angular servie par \textbf{Nginx} (Port 80).
    \item \textbf{Backend} : Serveur \textbf{Spring Boot 3.4} gérant la logique métier et les flux \textbf{SSE}.
    \item \textbf{IA Service} : Micro-service \textbf{FastAPI} pilotant l'intelligence artificielle.
    \item \textbf{Ollama Engine} : Inférence locale du modèle \textbf{Llama 3.2} pour la souveraineté des données.
    \item \textbf{Base de données} : Instance \textbf{PostgreSQL 16} avec volumes persistants (\texttt{pg\_data}).
\end{itemize}
Cette architecture modulaire assure une maintenance simplifiée et une \textbf{scalabilité horizontale} optimale, permettant de faire évoluer chaque composant indépendamment selon les besoins de charge.

\subsubsection*{5. Persistance des Données avec Docker Volumes}

Afin de garantir la conservation durable des données applicatives indépendamment du cycle de vie des conteneurs \texttt{Docker}, l’infrastructure de \textbf{TalentPredict} repose sur l’utilisation de \textit{Docker Volumes} dédiés.

Deux volumes principaux ont été configurés dans l’architecture de déploiement :

\begin{itemize}
    \item \textbf{\texttt{pg\_data}} : Volume associé au conteneur \texttt{PostgreSQL} permettant de persister l’ensemble des données relationnelles de la plateforme, notamment les comptes utilisateurs, les profils collaborateurs, les résultats d’évaluation, les notifications et les données analytiques RH.
    
    \item \textbf{\texttt{app\_uploads}} : Volume utilisé pour le stockage des fichiers générés ou téléversés par les utilisateurs, tels que les CV, les photos de profil, les documents PDF et les certificats produits automatiquement par la plateforme.
\end{itemize}

Cette approche permet de dissocier les données persistantes des conteneurs applicatifs eux-mêmes. Ainsi, même en cas de suppression, redémarrage ou mise à jour des conteneurs \texttt{Docker}, les données restent conservées et réutilisables lors du redéploiement du système.

L’utilisation des \textit{Docker Volumes} améliore également :
\begin{itemize}
    \item la fiabilité de l’infrastructure ;
    \item la sécurité des données ;
    \item la maintenabilité des environnements ;
    \item ainsi que la préparation à une future migration vers une infrastructure \textit{Cloud} ou \texttt{Kubernetes}.
\end{itemize}

\subsubsection{6.Stratégie de Supervision et Observabilité}

La pérennité de \textbf{TalentPredict} repose sur une surveillance proactive de ses composants conteneurisés. Pour assurer cette stabilité, nous avons déployé un dispositif d'observabilité articulé autour de deux axes : les métriques applicatives et la santé des conteneurs.

\paragraph*{Mise en œuvre du \textit{Monitoring} Applicatif}
Grâce à l’intégration de \texttt{Spring Boot Actuator}, le \textit{backend} de l'application devient "transparent". Il expose en continu des indicateurs critiques tels que l’utilisation de la mémoire \texttt{JVM}, la charge \texttt{CPU} et la disponibilité des ressources de persistance. Ces données permettent d'anticiper d'éventuelles dégradations de performance et facilitent le diagnostic technique à travers une centralisation des \textit{logs} applicatifs.

\paragraph*{Résilience et Auto-réparation \texttt{Docker}}
La robustesse de l'infrastructure est garantie par des mécanismes de \texttt{Healthcheck} intégrés à l'orchestration \texttt{Docker}. Chaque service (\textit{Backend}, IA, \texttt{PostgreSQL}) est soumis à des tests de santé périodiques. 

En cas d’anomalie ou d’indisponibilité d’un composant, \texttt{Docker} déclenche automatiquement un cycle de redémarrage. Cette capacité d’auto-réparation assure une continuité de service optimale et réduit drastiquement le temps d’indisponibilité (\textit{Downtime}) sans intervention manuelle, renforçant ainsi la fiabilité globale de la plateforme.


\subsubsection*{7. Pipeline CI/CD (Intégration et Déploiement Continus)}

Pour \textbf{TalentPredict}, nous avons mis en œuvre une architecture de \textit{pipeline} sophistiquée permettant de gérer indépendamment le cycle de vie du \textit{Backend} (\texttt{Spring Boot}) et du \textit{Frontend} (\texttt{Angular}).

\paragraph*{ Architecture \textit{Parent-Child} (Monorepo)}
Le fichier racine \texttt{.gitlab-ci.yml} agit comme un chef d'orchestre. Il ne déclenche les \textit{pipelines} spécifiques que si des modifications sont détectées dans les répertoires respectifs (\texttt{BackEnd/**} ou \texttt{FrontEnd/**}). Cette approche optimise les ressources et réduit le temps global d'intégration.

\paragraph*{Flux de Travail et Automatisation}
Le \textit{pipeline} du \textit{backend} automatise l'intégralité de la chaîne de production, de la validation syntaxique (\textit{Validate}) à la compilation (\textit{Build}). Une étape de test cruciale est intégrée, instanciant dynamiquement un service \texttt{PostgreSQL 16-alpine} pour valider la persistance des données. S'ensuivent une analyse de qualité via \texttt{Checkstyle}, la génération du \textit{package} exécutable, et la construction d'une image \texttt{Docker} poussée sur le registre privé de \texttt{GitLab}. Le déploiement final est protégé par une action manuelle, garantissant un contrôle humain avant la mise à jour sur le serveur.

\vspace{2mm}
\begin{lstlisting}[language=yaml, caption=Extrait du fichier de configuration (\texttt{BackEnd/.gitlab-ci.yml})]
test:
  stage: test
  services:
    - name: postgres:16-alpine
      alias: postgres
  variables:
    SPRING_DATASOURCE_URL: "jdbc:postgresql://postgres:5432/talentpredict_test"
  script:
    - cd BackEnd
    - mvn test -B
\end{lstlisting}
\vspace{2mm}

\paragraph*{ Bénéfices pour le Projet}
Cette industrialisation offre une fiabilité totale grâce au cache \texttt{Maven} et à l'archivage des rapports \texttt{JUnit}. La sécurité est renforcée par l'utilisation de variables d'environnement masquées, assurant une conformité aux standards \textbf{DevOps} professionnels.
La figure suivante présente la chaîne d'intégration et de déploiement continus (CI/CD) mise en place pour le projet TalentPredict AI via GitLab CI/CD, couvrant les pipelines frontend et backend ainsi que l'environnement de production conteneurisé :

\begin{figure}[H]
    \centering
    % Remplacez 'nom_de_votre_image' par le nom réel de votre fichier image
    \includegraphics[width=0.8\textwidth]{images/piplineCICD.png}
    \caption{Chaîne CI/CD du projet TalentPredict AI}
    \label{fig:pipeline-cicd}
\end{figure}

Comme illustré dans la figure \ref{fig:pipeline-cicd}, le pipeline CI/CD est déclenché automatiquement à chaque modification du code source. Le pipeline frontend exécute successivement l'installation des dépendances, l'analyse statique du code, la compilation Angular, les tests unitaires et la création de l'image Docker. Le pipeline backend suit une logique similaire avec Maven, en intégrant les tests JUnit 5, le contrôle qualité Checkstyle et l'empaquetage JAR avant la conteneurisation. Les services déployés communiquent via un réseau Docker privé, garantissant l'isolation et la sécurité des échanges entre les composants de la plateforme.

% Fin de document prématurée supprimée pour inclure la bibliographie active


\subsubsection*{8. Sécurisation du Pipeline DevOps (DevSecOps)}

Le déploiement automatisé via \texttt{GitLab CI/CD} introduit des vecteurs de risques (fuite de clés, librairies vulnérables). \textbf{TalentPredict} intègre des mesures de sécurité à chaque étape du \textit{pipeline} pour garantir l'intégrité de la plateforme.

\paragraph*{Gestion des Secrets et Variables d'Environnement}
Le principe fondamental de "\textit{No Secrets in Git}" est strictement appliqué.
\begin{itemize}
    \item \textbf{Variables \texttt{GitLab CI/CD} :} Les identifiants sensibles (mot de passe DB, clé secrète \texttt{JWT}, identifiants du registre \texttt{Docker}) sont stockés de manière chiffrée dans les paramètres \texttt{GitLab}. Ils ne sont jamais présents dans le code source.
    \item \textbf{Injection Dynamique :} Ces variables sont injectées sous forme de variables d'environnement au moment du \textit{build}, assurant qu'elles ne sont accessibles qu'au processus autorisé.
\end{itemize}

\paragraph*{Analyse de Sécurité des Dépendances}
Le \textit{pipeline} intègre des vérifications automatiques pour prévenir l'utilisation de librairies obsolètes ou malveillantes :
\begin{itemize}
    \item \textbf{\texttt{Maven Dependency Check} :} Lors de l'étape de Test, \texttt{Maven} vérifie si les dépendances \texttt{Java} (\texttt{pom.xml}) contiennent des vulnérabilités connues (CVE).
    \item \textbf{Audit \texttt{npm} :} Une vérification similaire est effectuée sur le \textit{frontend} \texttt{Angular} pour s'assurer de l'intégrité du catalogue de modules \texttt{Node.js}.
\end{itemize}

\paragraph*{Isolation et Sécurité du Registre Docker}
\begin{itemize}
    \item \textbf{Registre Privé :} Nous utilisons le registre privé de \texttt{GitLab} pour stocker les images de nos conteneurs. Seul le \textit{pipeline} de déploiement et le serveur de production (via des identifiants \texttt{deploy\_token}) peuvent y accéder.
    \item \textbf{Principe du Moindre Privilège :} Les conteneurs en production sont configurés pour ne pas fonctionner avec des privilèges "\texttt{root}", limitant ainsi l'impact potentiel en cas de compromission d'un service.
\end{itemize}

\paragraph*{Contrôle d'Accès et Traçabilité}
\begin{itemize}
    \item \textbf{\textit{Protected Branches} :} Seuls les "\textit{Maintainers}" du projet peuvent déclencher un déploiement sur la branche \texttt{main} (Production).
    \item \textbf{\textit{Audit Log} :} Chaque étape du \textit{pipeline} est tracée dans \texttt{GitLab}, permettant de savoir exactement qui a déclenché un déploiement, avec quel code et à quel moment.
\end{itemize}

\subsubsection{Bilan du Déploiement}

Le déploiement final a permis de confirmer l'extrême résilience de l'architecture choisie. Grâce au couplage de l'orchestration Docker et des pipelines CI/CD, la plateforme \textbf{TalentPredict} peut désormais être mise à jour en quelques minutes de manière transparente. De plus, cette architecture modulaire garantit la \textit{scalabilité} future du système, permettant par exemple de dupliquer indépendamment le moteur IA en cas de forte montée en charge.


\section{Suivi et Pilotage du Sprint}
\subsection{Le Scrum Board}

Le \textit{Scrum Board} a constitué le pilier de notre gestion visuelle et opérationnelle tout au long de ce sprint. Structuré sous forme de tableau \textit{Kanban}, il a agi comme un véritable radiateur d'information, offrant une transparence totale sur le cycle de vie de chaque \textit{User Story}, de sa phase de conception jusqu'à sa validation finale. Cet outil a grandement facilité la coordination quotidienne des développements au sein de l'équipe technique, en synergie avec Mohamed Ben Abdeladhim et Rahma Barouni, permettant de répartir efficacement la charge de travail et d'identifier instantanément les points de blocage. Par ailleurs, ce formalisme a offert à notre \textit{Product Owner}, Med Amine Gharbi, un cadre clair et structuré pour effectuer la recette fonctionnelle continue et certifier que chaque module respectait scrupuleusement la \textit{Definition of Done} (DoD).

% Insertion de l'image du Scrum Board
\begin{figure}[H]
    \centering
    % Remplacez 'nom_image_scrum_board.png' par le nom exact de votre fichier
    \includegraphics[width=0.9\textwidth]{images/scrumboardSpring.png} 
    \vspace{2mm}
    \caption{Tableau de suivi visuel (Scrum Board) du Sprint 4}
    \label{fig:scrum_board_sprint4}
\end{figure}

\subsection{Le Burndown Chart}

En complément du suivi quotidien, le \textit{Burndown Chart} a représenté notre principal outil de pilotage prédictif. Ce graphique nous a permis de mesurer avec précision la vélocité de l'équipe en confrontant l'effort global restant à fournir avec la trajectoire de complétion idéale sur les quatre semaines de ce mois de mai 2026. L'analyse itérative de cette courbe s'est révélée cruciale pour la réussite de cette dernière phase : elle nous a permis de détecter de manière proactive les variations de rythme, notamment lors des défis techniques liés à l'intégration asynchrone, d'ajuster notre capacité d'exécution en temps réel, et de sécuriser ainsi la livraison de l'intégralité du périmètre fonctionnel avant la clôture du projet.

\begin{figure}[H]
    \centering
    \begin{tikzpicture}
    \begin{axis}[
        title={Burndown Chart - Sprint 4 (TalentPredict)},
        xlabel={Jours du Sprint},
        ylabel={Effort Restant (Story Points)},
        xmin=1, xmax=20,
        ymin=0, ymax=45,
        xtick={1,4,8,12,16,20},
        ytick={0,10,20,30,40},
        legend pos=north east,
        xmajorgrids=true,
        ymajorgrids=true,
        grid style=dashed,
        width=12cm,
        height=7cm
    ]

    % Ligne de l'effort IDÉAL (Ligne droite de 40 à 0)
    \addplot[
        color=blue,
        domain=1:20,
        dashed,
        line width=1pt
    ] {40 - (40/19)*(x-1)};
    \addlegendentry{Effort Idéal}

    % Ligne de l'effort RÉEL (Données fictives réalistes pour le Sprint 4)
    \addplot[
        color=red,
        mark=square*,
        line width=1.2pt
    ] coordinates {
        (1,40) (3,38) (5,35) (8,30) (10,28) (12,22) (15,12) (18,5) (20,0)
    };
    \addlegendentry{Effort Réel}

    \end{axis}
    \end{tikzpicture}
    \caption{Suivi de l'avancement du Sprint 4}
\end{figure} 

\subsection{Indicateurs de Performance du Sprint 4}

Ce dernier sprint a été le garant de l'industrialisation et de la mise en service. Nous avons mesuré la performance à travers trois prismes : l'agilité, la robustesse \textbf{DevOps} et la valeur métier des tableaux de bord.

\paragraph*{1. Indicateurs de Pilotage (Agilité)}
Le sprint s'est achevé avec une vélocité de 28 \textit{Story Points}. Bien que le volume soit inférieur aux sprints précédents, la complexité résidait dans l'automatisation et l'interfaçage des systèmes.
\begin{itemize}
    \item \textbf{Taux de complétion :} 100\%. Toutes les tâches de déploiement et de pilotage ont été terminées.
    \item \textbf{Qualité des \textit{User Stories} :} Les \textit{dashboards} RH ont été validés avec un taux d'acceptation immédiat, grâce à la précision des données agrégées venant du moteur d'IA.
\end{itemize}

\paragraph*{2. Indicateurs DevOps et Infrastructure}
L'industrialisation via \texttt{GitLab CI/CD} a permis d'atteindre des niveaux de performance remarquables :
\begin{itemize}
    \item \textbf{Vitesse d'Intégration :} Le \textit{build} du \textit{backend} s'effectue en moyenne en 1 min 30 s, tandis que le \textit{bundle} \textit{frontend} nécessite 2 min 45 s.
    \item \textbf{Fiabilité du Déploiement :} L'utilisation de \texttt{docker-compose} en mode \textit{detached} assure des mises à jour transparentes avec une interruption de service quasi-nulle.
    \item \textbf{Optimisation des Images :} La taille finale des images \texttt{Docker} a été réduite de 25\% par rapport aux premiers tests, optimisant ainsi l'espace disque sur le serveur.
\end{itemize}

\paragraph*{3. Indicateurs de Pilotage RH (\textit{Dashboards})}
\begin{itemize}
    \item \textbf{Temps de réponse du \textit{Dashboard} :} L'agrégation des données statistiques (plusieurs milliers d'entrées potentielles) se fait en moins de 300 ms, garantissant une fluidité de navigation pour l'administrateur.
    \item \textbf{Cohérence des données :} 100\% de concordance entre les résultats des tests individuels et les graphiques de tendances globales affichés.
\end{itemize}

\subsubsection*{4. Tableau de Synthèse Final (KPIs Sprint 4)}
\begin{table}[H]
    \centering
    \renewcommand{\arraystretch}{1.5}
    \small
    \begin{tabularx}{0.95\textwidth}{|X|c|c|c|}
        \hline
        \rowcolor{TableHeaderBlue}
        \color{white}\textbf{Indicateur} & \color{white}\textbf{Cible} & \color{white}\textbf{Résultat Obtenu} & \color{white}\textbf{État} \\
        \hline
        \textit{Story Points} complétés & 28 & 28 & $\checkmark$ Validé \\
        \hline
        Succès des \textit{Pipelines} \texttt{GitLab} & $100\%$ & $100\%$ & $\checkmark$ Validé \\
        \hline
        Temps moyen de \textit{Build} (Global) & $< 6$ min & 4 min 15 s & $\checkmark$ Validé \\
        \hline
        \textit{Uptime} Production (\texttt{Docker}) & $100\%$ & $100\%$ & $\checkmark$ Validé \\
        \hline
        Latence agrégation \textit{Dashboards} & $< 500$ ms & $\sim 280$ ms & $\checkmark$ Validé \\
        \hline
        Audit Sécurité (CORS/\texttt{JWT}/SSL) & Conforme & $100\%$ & $\checkmark$ Validé \\
        \hline
    \end{tabularx}
    \caption{Tableau de bord des métriques globales du Sprint 4}
    \label{tab:kpi_tech_s4}
\end{table}

\paragraph*{5. Analyse du Bilan}
Le bilan du Sprint 4 confirme la viabilité industrielle de \textbf{TalentPredict}. L'automatisation du cycle de vie logiciel (CI/CD) couplée à une infrastructure conteneurisée robuste permet d'envisager sereinement une mise en production réelle. Les tableaux de bord fournissent enfin une vision macroscopique des talents, transformant les données IA individuelles en un véritable outil de pilotage stratégique pour l'entreprise.

\section*{Conclusion}
\phantomsection
\addcontentsline{toc}{section}{Conclusion}
% Utilisez \section*{Conclusion} si vous ne souhaitez pas qu'elle soit numérotée

Ce quatrième et dernier chapitre s'est attaché à décrire avec précision l'ultime itération de notre cycle de développement, marquant l'aboutissement technique de la plateforme TalentPredict. En nous appuyant sur une conception architecturale solide et modulaire, nous avons exposé l'implémentation de fonctionnalités avancées et critiques, telles que le moteur de reporting documentaire et le système de notifications asynchrones en temps réel.

L'intégration d'une démarche DevOps rigoureuse a constitué un axe majeur de cette phase finale. La mise en place d'une stratégie de tests multiniveau, couplée à la conteneurisation via Docker et à l'automatisation des déploiements (CI/CD), a permis de transformer un code en développement en un produit logiciel fiable, sécurisé et prêt pour la production. De plus, le pilotage agile de ce sprint, soutenu par des artefacts Scrum pertinents, a garanti une gestion optimale de l'effort d'équipe et le respect scrupuleux des délais impartis.

La réussite de cette dernière itération clôture ainsi la phase de réalisation technique de notre système, ouvrant naturellement la voie au bilan final et à la conclusion générale de ce projet de fin d'études.


\clearpage
\chapter*{Conclusion Générale}
\addcontentsline{toc}{chapter}{Conclusion Générale}

Le présent mémoire de fin d’études, réalisé au sein de la société AgileX Management dans le cadre de la formation dispensée à l’Institut Supérieur d’Informatique de Mahdia, a porté sur la conception et le développement de TalentPredict AI, une plateforme intelligente dédiée à l’évaluation et au développement des compétences en entreprise.

Ce projet s’inscrit dans un contexte où les organisations cherchent à mieux identifier, mesurer et valoriser les compétences de leurs collaborateurs. Les méthodes classiques d’évaluation, souvent fondées sur des CV déclaratifs, des entretiens ponctuels ou des appréciations subjectives, présentent plusieurs limites. Elles ne permettent pas toujours de croiser les informations disponibles, de valider objectivement les compétences techniques ni d’analyser de manière structurée les compétences comportementales. Face à ces limites, TalentPredict AIpropose une approche intégrée, multi-sources et orientée aide à la décision RH.

La réalisation du projet a été conduite selon une démarche Agile Scrum, organisée en quatre sprints complémentaires. Le premier sprint a permis de mettre en place les fondations techniques de la plateforme, notamment l’architecture applicative, le système d’authentification sécurisé par JWT, la gestion différenciée des profils employé et administrateur RH, ainsi que l’environnement conteneurisé avec Docker. Le deuxième sprint a constitué le cœur fonctionnel de l’évaluation des compétences, à travers l’analyse du CV, l’exploitation du profil GitHub, l’évaluation technique par QCM et challenge de code, ainsi que l’évaluation comportementale fondée sur le modèle PCM et enrichie par des traitements d’analyse sémantique. Le troisième sprint a transformé les résultats d’évaluation en recommandations concrètes, en intégrant un moteur d’upskilling, des parcours personnalisés, un suivi de progression et des mécanismes de validation. Enfin, le quatrième sprint a permis de finaliser la plateforme à travers les tableaux de bord analytiques, la génération du Talent Passport, la communication en temps réel et la préparation du déploiement.

Sur le plan technique, le projet a permis de mobiliser plusieurs technologies complémentaires. Le backend métier a été développé avec Spring Boot, tandis que les services d’intelligence artificielle ont été organisés autour de FastAPI, Ollama et n8n. L’interface utilisateur a été réalisée avec Angular, et la persistance des données a été assurée par PostgreSQL. Cette architecture modulaire facilite la séparation des responsabilités, améliore la maintenabilité du système et prépare l’intégration future de nouveaux modules d’analyse ou de recommandation.

La valeur ajoutée de TalentPredict AI réside principalement dans sa capacité à dépasser une évaluation purement déclarative des compétences. La plateforme confronte les informations fournies par le collaborateur à des preuves techniques observables, telles que le profil GitHub, les tests et les challenges de code. Elle prend également en compte les compétences comportementales à travers un cadre psychométrique structuré et une analyse sémantique locale. Cette combinaison permet de produire une vision plus complète du profil collaborateur, tout en générant des recommandations de formation adaptées aux écarts identifiés.

Le projet met également l’accent sur la souveraineté et la confidentialité des données. Le recours à une inférence locale via Ollama limite l’externalisation des données sensibles et renforce la conformité avec les exigences de protection des données personnelles. Cette orientation est particulièrement importante dans un contexte RH, où les informations traitées peuvent concerner les profils professionnels, les résultats d’évaluation, les préférences d’apprentissage et certaines dimensions comportementales.

Malgré les résultats obtenus, certaines limites doivent être prises en considération. L’évaluation sémantique pourrait être améliorée par l’entraînement ou l’adaptation de modèles spécialisés sur des corpus RH francophones. Les performances de la plateforme devraient également être testées sur un volume plus important d’utilisateurs et de profils afin de mieux évaluer sa robustesse en conditions réelles. Par ailleurs, l’intégration avec des référentiels de compétences reconnus, tels que ROME ou ESCO, permettrait d’améliorer la normalisation des résultats et la pertinence des recommandations.

Plusieurs perspectives peuvent être envisagées pour faire évoluer la solution. Il serait pertinent d’intégrer un module de feedback périodique entre collaborateurs afin d’enrichir l’évaluation comportementale par des observations issues du contexte de travail réel. La plateforme pourrait également être connectée à des systèmes d’information RH existants afin d’automatiser davantage la gestion des parcours de formation. À moyen terme, l’évolution vers une architecture plus scalable, basée sur une orchestration avancée des services, pourrait faciliter le passage d’un prototype académique à une solution exploitable à plus grande échelle.

En définitive, ce projet de fin d’études a permis de concevoir une plateforme complète, cohérente et évolutive répondant à une problématique réelle de gestion des talents. Il a également constitué une expérience formatrice, mobilisant des compétences en génie logiciel, intelligence artificielle, DevOps, conception d’interfaces et gestion de projet Agile. TalentPredict AI représente ainsi une base solide pour le développement futur d’un outil RH intelligent, capable d’accompagner les entreprises dans l’évaluation, la valorisation et le développement continu des compétences.

\color{red}
% =========================================================
% Suggestions de figures à ajouter pour améliorer le rapport
% =========================================================

\section*{Suggestions de figures à intégrer dans le rapport}
\phantomsection
\addcontentsline{toc}{section}{Suggestions de figures à intégrer dans le rapport}

Afin d’améliorer la lisibilité du rapport et de renforcer la valeur explicative des chapitres, plusieurs figures peuvent être ajoutées. Ces figures ont pour objectif de synthétiser les choix méthodologiques, fonctionnels, techniques et expérimentaux du projet \textit{TalentPredict AI}. Le tableau~\ref{tab:suggestions_figures_rapport} présente les principales figures recommandées, leur emplacement conseillé et leur objectif.

\begin{longtable}{|P{0.24\textwidth}|P{0.28\textwidth}|P{0.40\textwidth}|}
\hline
\rowcolor{TableHeaderBlue}
\color{white}\textbf{Emplacement recommandé} &
\color{white}\textbf{Titre proposé de la figure} &
\color{white}\textbf{Objectif de la figure} \\
\hline
\endfirsthead

\hline
\rowcolor{TableHeaderBlue}
\color{white}\textbf{Emplacement recommandé} &
\color{white}\textbf{Titre proposé de la figure} &
\color{white}\textbf{Objectif de la figure} \\
\hline
\endhead

\hline
\endfoot

Chapitre 1 -- Méthodologie Scrum &
Organisation Scrum du projet &
Présenter les rôles Scrum, les artefacts, les cérémonies et leur application dans le projet. \\
\hline

Chapitre 1 -- Solution proposée &
Architecture en couches &
Présenter les couches de la solution : présentation, API métier, orchestration IA, inférence locale, persistance et déploiement. \\
\hline


Chapitre 2 -- Concepts et technologies clés &
Pipeline NLP appliqué à l’analyse des profils &
Expliquer le passage du CV ou du profil professionnel vers l’extraction des compétences, la normalisation et le scoring. \\
\hline

Chapitre 2 -- Analyse GitHub &
Exploitation des données GitHub comme preuve technique &
Illustrer les signaux extraits de GitHub : langages, dépôts, commits, régularité, contribution et activité. \\
\hline

Chapitre 2 -- LLM et souveraineté des données &
Inférence locale avec Ollama &
Montrer la différence entre une approche cloud et une approche locale, en insistant sur la confidentialité et la conformité RGPD. \\
\hline

Chapitre 2 -- Analyse comparative &
Matrice de positionnement des solutions existantes &
Comparer \textit{TalentPredict AI} aux solutions existantes selon deux axes : objectivité de l’évaluation et souveraineté des données. \\
\hline

Chapitre 2 -- Contributions &
Positionnement scientifique et technique de \textit{TalentPredict AI} &
Synthétiser les trois contributions : validation multi-sources, évaluation technique et comportementale, inférence locale. \\
\hline

Sprint 1 -- Objectif du sprint &
Objectifs techniques et fonctionnels du Sprint 1 &
Présenter les objectifs du sprint : architecture, sécurité, authentification, profils et environnement Docker. \\
\hline

Sprint 1 -- Conception &
Architecture de sécurité du Sprint 1 &
Illustrer le fonctionnement de l’authentification JWT, du hachage BCrypt, des rôles et des routes protégées. \\
\hline


Sprint 1 -- Conception &
Modèle de données utilisateur et profil &
Montrer les entités principales liées à l’utilisateur, au rôle, au profil, au token et aux informations professionnelles. \\
\hline

Sprint 1 -- Déploiement &
Infrastructure Docker du Sprint 1 &
Présenter les conteneurs Angular, Spring Boot, PostgreSQL et leurs interactions réseau. \\
\hline

Sprint 1 -- Tests &
Stratégie de test de l’authentification &
Synthétiser les tests réalisés : inscription, connexion, récupération de mot de passe, routes protégées et accès par rôle. \\
\hline

Sprint 2 -- Objectif du sprint &
Pipeline global d’analyse multi-sources &
Montrer le flux complet depuis le CV, GitHub, QCM et challenge de code vers le scoring technique et comportemental. \\
\hline

Sprint 2 -- Analyse technique &
Pipeline d’analyse du CV &
Illustrer les étapes : chargement du fichier, extraction du texte, nettoyage, détection des compétences, normalisation et stockage. \\
\hline

Sprint 2 -- Analyse GitHub &
Workflow d’analyse GitHub &
Présenter les étapes d’appel à l’API GitHub, récupération des données, extraction des signaux techniques et calcul des scores. \\
\hline

Sprint 2 -- Analyse comportementale &
Pipeline d’évaluation PCM &
Montrer les étapes : questionnaire PCM, calcul heuristique, analyse sémantique, profil dominant et scores comportementaux. \\
\hline

Sprint 2 -- Orchestration &
Workflow n8n d’analyse des compétences &
Présenter comment n8n orchestre les appels entre les services, les prompts IA, FastAPI et le backend métier. \\
\hline

Sprint 2 -- Scoring &
Fusion multi-sources des scores &
Illustrer la triangulation entre CV, GitHub, QCM, challenge de code et analyse IA pour produire un score de confiance. \\
\hline

Sprint 3 -- Objectif du sprint &
Architecture du moteur de recommandation &
Présenter comment les résultats d’évaluation alimentent le moteur de recommandation de formations. \\
\hline

Sprint 2 -- IA générative &
Génération de QCM adaptatifs par IA &
Montrer le processus : profil de l’utilisateur, compétence cible, prompt, génération du QCM et validation des réponses. \\
\hline

Sprint 2 -- Challenge de code &
Évaluation automatique des challenges de code &
Illustrer le flux : génération du problème, soumission du code, évaluation, feedback et scoring. \\
\hline

Sprint 3 -- Upskilling &
Roadmap personnalisée de formation &
Présenter la construction du parcours de formation à partir des lacunes détectées et du profil comportemental. \\
\hline


Sprint 4 -- Objectif du sprint &
Pilotage stratégique des compétences &
Présenter la transformation des résultats individuels en indicateurs RH globaux. \\
\hline

Sprint 4 -- Tableaux de bord &
Tableau de bord employé &
Montrer les principales zones du dashboard employé : scores, recommandations, progression et Talent Passport. \\
\hline

Sprint 4 -- Tableaux de bord &
Tableau de bord administrateur RH &
Illustrer les vues RH : cartographie des compétences, écarts, formations, indicateurs globaux et rapports. \\
\hline

Sprint 4 -- Reporting &
Génération du Talent Passport &
Présenter le processus de génération du document synthétique à partir des résultats d’évaluation. \\
\hline

Sprint 4 -- DevOps &
Chaîne CI/CD et mise en production &
Montrer les étapes : commit, build, tests, image Docker, déploiement et supervision. \\
\hline

Sprint 4 -- Tests &
Stratégie globale de tests et validation &
Synthétiser les tests unitaires, tests d’intégration, tests API, tests UI et validation fonctionnelle. \\
\hline
\caption{Suggestions de figures pour améliorer le rapport}
\label{tab:suggestions_figures_rapport}\\
\end{longtable}
\color{black}

%\chapter*{Références}
\renewcommand{\bibname}{Références}
\begin{thebibliography}{99}
\bibitem{ch2-ulrich1997} Ulrich, D. (1997). \emph{Human Resource Champions: The Next Agenda for Adding Value and Delivering Results}. Boston, MA: Harvard Business School Press.
\bibitem{ch2-salgado1997} Salgado, J. F. (1997). The Five Factor Model of Personality and Job Performance in the European Community. \emph{Journal of Applied Psychology}, 82(1), 30--43.
\bibitem{ch2-donaldson2002} Donaldson, S. I., \& Grant-Vallone, E. J. (2002). Understanding Self-Report Bias in Organizational Behavior Research. \emph{Journal of Business and Psychology}, 17(2), 245--260.
\bibitem{ch2-kruger1999} Kruger, J., \& Dunning, D. (1999). Unskilled and Unaware of It: How Difficulties in Recognizing One's Own Incompetence Lead to Inflated Self-Assessments. \emph{Journal of Personality and Social Psychology}, 77(6), 1121--1134.
\bibitem{ch2-rgpd2016} Règlement (UE) 2016/679 du Parlement européen et du Conseil du 27 avril 2016 (RGPD), Articles 9 et 25. \emph{Journal officiel de l'Union européenne}, L 119, 4 mai 2016.
\bibitem{ch2-mcdaniel1994} McDaniel, M. A., Whetzel, D. L., Schmidt, F. L., \& Maurer, S. D. (1994). The Validity of Employment Interviews: A Comprehensive Review and Meta-Analysis. \emph{Journal of Applied Psychology}, 79(4), 599--616.
\bibitem{ch2-costa1992} Costa, P. T., \& McCrae, R. R. (1992). \emph{Revised NEO Personality Inventory (NEO-PI-R) and NEO Five-Factor Inventory (NEO-FFI): Professional Manual}. Odessa, FL: Psychological Assessment Resources.
\bibitem{ch2-pittenger1993} Pittenger, D. J. (1993). The Utility of the Myers-Briggs Type Indicator. \emph{Review of Educational Research}, 63(4), 467--488.
\bibitem{ch2-kahler1996} Kahler, T. (1996). \emph{Process Communication Model}. Little Rock, AR: Kahler Communications.
\bibitem{ch2-thornton2006} Thornton, G. C., \& Rupp, D. E. (2006). \emph{Assessment Centers in Human Resource Management: Strategies for Prediction, Diagnosis, and Development}. Mahwah, NJ: Lawrence Erlbaum Associates.
\bibitem{ch2-celik2013} Celik, D., \& Elçi, A. (2013). A SWRL-Based Ontology for Resume Information Extraction. In \emph{Proceedings of the 2013 International Conference on Information Technology Based Higher Education and Training (ITHET)}. IEEE.
\bibitem{ch2-chien2020} Chien, H.-J., \& Chia-Hui, C. (2020). Resume Information Extraction with Cascaded Hybrid Model. In \emph{Proceedings of the 58th Annual Meeting of the Association for Computational Linguistics (ACL 2020)}, pp. 5559--5569.
\bibitem{ch2-shi2023} Shi, W., et al. (2023). Large Language Models Are Not Robust Multiple Choice Selectors. \emph{arXiv preprint arXiv:2309.03882}.
\bibitem{ch2-kalliamvakou2014} Kalliamvakou, E., Gousios, G., Blincoe, K., Singer, L., German, D. M., \& Damian, D. (2014). The Promises and Perils of Mining GitHub. In \emph{Proceedings of the 11th Working Conference on Mining Software Repositories (MSR 2014)}, pp. 92--101. ACM.
\bibitem{ch2-marlow2013} Marlow, J., Dabbish, L., \& Herbsleb, J. (2013). Impression Formation in Online Peer Production: Activity Traces and Personal Profiles in GitHub. In \emph{Proceedings of the 2013 Conference on Computer Supported Cooperative Work (CSCW 2013)}, pp. 117--128. ACM.
\bibitem{ch2-vaswani2017} Vaswani, A., Shazeer, N., Parmar, N., Uszkoreit, J., Jones, L., Gomez, A. N., Kaiser, L., \& Polosukhin, I. (2017). Attention Is All You Need. In \emph{Advances in Neural Information Processing Systems (NeurIPS 2017)}, vol. 30.
\bibitem{ch2-meta2024} Meta AI. (2024). Llama 3 Model Card. Meta Platforms. Disponible à : \url{https://llama.meta.com/llama3}
\bibitem{ch2-mistral2023} Mistral AI. (2023). Mistral 7B Technical Report. Disponible à : \url{https://arxiv.org/abs/2310.06825}
\bibitem{ch2-abdin2024} Abdin, M., et al. (2024). Phi-3 Technical Report: A Highly Capable Language Model Locally on Your Phone. \emph{arXiv preprint arXiv:2404.14219}.
\bibitem{ch2-touvron2023} Touvron, H., et al. (2023). LLaMA: Open and Efficient Foundation Language Models. \emph{arXiv preprint arXiv:2302.13971}.
\bibitem{ch2-xu2024} Xu, Y., et al. (2024). Benchmarking LLaMA 3 on HR NLP Tasks. \emph{arXiv preprint}.
\bibitem{ch2-vanderaalst2013} van der Aalst, W. M. P. (2013). Business Process Management: A Comprehensive Survey. \emph{ISRN Software Engineering}, vol. 2013, Article ID 507984.
\bibitem{ch2-newman2015} Newman, S. (2015). \emph{Building Microservices: Designing Fine-Grained Systems}. Sebastopol, CA: O'Reilly Media.
\bibitem{ch2-drachsler2009} Drachsler, H., Hummel, H., \& Koper, R. (2009). Factors Affecting the Robustness of Recommender Systems in Technology Enhanced Learning Environments. \emph{Journal of Digital Information}, 10(2).

\bibitem{kahler1996} Kahler, T. (1996). \emph{Process Communication Model}. Little Rock, AR: Kahler Communications.
\bibitem{touvron2023llama} Touvron, H., et al. (2023). LLaMA: Open and Efficient Foundation Language Models. \emph{arXiv:2302.13971}.
\bibitem{roziere2023codellama} Rozière, B., et al. (2023). Code Llama: Open Foundation Models for Code. \emph{arXiv:2308.12950}.
\bibitem{chien2020} Chien, H.-J., \& Chang, C.-H. (2020). Resume Information Extraction with Cascaded Hybrid Model. In \emph{Proceedings of ACL 2020}, pp. 5559--5569.
\bibitem{kalliamvakou2014} Kalliamvakou, E., Gousios, G., Blincoe, K., Singer, L., German, D. M., \& Damian, D. (2014). The Promises and Perils of Mining GitHub. In \emph{MSR 2014}, pp. 92--101. ACM.
\bibitem{celik2013} Celik, I., \& Elçi, A. (2013). An ontology-based information extraction system for resume parsing.
\bibitem{rgpd2016} Règlement (UE) 2016/679 du Parlement européen et du Conseil (RGPD), Article 9.
\bibitem{donaldson2002} Donaldson, S. I., \& Grant-Vallone, E. J. (2002). Understanding Self-Report Bias in Organizational Behavior Research. \emph{Journal of Business and Psychology}, 17(2), 245--260.
\bibitem{kruger1999dunning} Kruger, J., \& Dunning, D. (1999). Unskilled and Unaware of It. \emph{Journal of Personality and Social Psychology}, 77(6), 1121--1134.
\bibitem{yao2023react} Yao, S., et al. (2023). ReAct: Synergizing Reasoning and Acting in Language Models. \emph{ICLR 2023}.
\bibitem{drachsler2009} Drachsler, H., Hummel, H., \& Koper, R. (2009). Factors Affecting the Robustness of Recommender Systems. \emph{Journal of Digital Information}, 10(2).
\bibitem{angular} Google. (2024). Angular Documentation. \url{https://angular.dev}
\bibitem{springboot} VMware Tanzu. (2025). Spring Boot Reference Guide. \url{https://spring.io/projects/spring-boot}
\bibitem{fastapi} Ramírez, S. (2024). FastAPI Documentation. \url{https://fastapi.tiangolo.com}
\bibitem{n8n} n8n GmbH. (2025). n8n Documentation. \url{https://docs.n8n.io}
\bibitem{ollama} Ollama. (2024). Ollama Documentation. \url{https://ollama.com}
\bibitem{docker} Docker Inc. (2024). Docker Documentation. \url{https://docs.docker.com}
\bibitem{eks} Amazon Web Services. (2025). Amazon Elastic Kubernetes Service User Guide. \url{https://docs.aws.amazon.com/eks/}
\bibitem{postgresql} The PostgreSQL Global Development Group. (2024). PostgreSQL Documentation. \url{https://www.postgresql.org/docs/}
\bibitem{pdfjs} Mozilla Foundation. (2025). PDF.js Documentation. \url{https://mozilla.github.io/pdf.js/}
\bibitem{chartjs} Chart.js Contributors. (2025). Chart.js Documentation. \url{https://www.chartjs.org}
\bibitem{pdfplumber} jsvine. (2024). pdfplumber Documentation. \url{https://github.com/jsvine/pdfplumber}
\bibitem{scrum2020} Schwaber, K., \& Sutherland, J. (2020). \emph{The Scrum Guide}. \url{https://scrumguides.org}
\bibitem{devlin2019bert} Devlin, J., Chang, M.-W., Lee, K., \& Toutanova, K. (2019). BERT: Pre-training of Deep Bidirectional Transformers for Language Understanding. In \emph{NAACL-HLT 2019}.
\bibitem{salgado1997} Salgado, J. F. (1997). The Five Factor Model of Personality and Job Performance. \emph{Journal of Applied Psychology}, 82(1), 30--43.
\bibitem{ulrich1997} Ulrich, D. (1997). \emph{Human Resource Champions}. Harvard Business School Press.
\end{thebibliography}
\end{document}
\documentclass[12pt,a4paper]{report}

% =========================================================================
% 1. PACKAGES REQUIS
% =========================================================================
% Encodage & langue
\usepackage[utf8]{inputenc}
\usepackage[T1]{fontenc}
\usepackage[french]{babel}
\usepackage{longtable}
\usepackage{booktabs}
\usepackage{array}
\usepackage[table]{xcolor}
\usepackage{ragged2e}
\usepackage{needspace}
\usepackage{float}

\definecolor{TableHeaderBlue}{RGB}{32,78,121}
\definecolor{TableLightBlue}{RGB}{232,240,248}
\definecolor{TableVeryLight}{RGB}{248,251,253}


\setlength{\LTpre}{4pt}
\setlength{\LTpost}{6pt}
\renewcommand{\arraystretch}{1.25}
% Mise en page
\usepackage[a4paper,left=2cm,right=2cm,top=2cm,bottom=2cm]{geometry}
\usepackage{titlesec}
\usepackage{ragged2e}
\usepackage{needspace}
\usepackage{float}             % pour les figures/tableaux fixes
\usepackage[section]{placeins}  % empêche les flottants de dépasser leur section

% Figures et graphiques
\usepackage{graphicx}
\usepackage{tcolorbox}         % encadrés (chemin critique, etc.)

% Tableaux
\usepackage{array}
\usepackage{tabularx}
\usepackage{booktabs}
\usepackage{multirow}
\usepackage[table]{xcolor}     % couleurs dans les tableaux
\usepackage{pdflscape}

% Listes & citations
\usepackage{enumitem}
\usepackage{csquotes}
\usepackage{amsmath}
\usepackage{amssymb}
\usepackage[hypertexnames=false,bookmarks=false]{hyperref}
\usepackage{lmodern}
\usepackage{fix-cm}
\usepackage{fancyhdr}

% =========================================================================
% 2. CONFIGURATION DU DOCUMENT (À REMPLIR PAR LES ÉTUDIANTS)
% =========================================================================

% Informations sur les étudiants
\newcommand{\EtudiantUn}{Mohamed Ben Abdeladhim}
\newcommand{\EtudiantDeux}{Rahma Barouni}

% Informations sur le projet
\newcommand{\TitreProjet}{TalentPredict}
\newcommand{\LieuProjet}{AgileX Management}
\newcommand{\DateSoutenance}{\today} % ou remplacez \today par "6 mai 2026"

% Membres du Jury
\newcommand{\PresidentJury}{Pr. Nom Prénom}
\newcommand{\Rapporteur}{Pr. Nom Prénom}
\newcommand{\EncadreurAcademique}{Dr. Ben Aïcha Takwa}
\newcommand{\EncadrantProfessionnel}{Med Amine Gharbi}

% Chemins des logos
\newcommand{\LogoISIMA}{images/isima.png}
\newcommand{\LogoUniversite}{images/universit.png}


% =========================================================================
% 3. DÉFINITION DES COMMANDES ET STYLES
% =========================================================================
\usepackage{pgfplots}
\pgfplotsset{compat=1.18} % ou la version de ton compilateur

% Réglages de mise en page généraux
\setlength{\parindent}{0pt}
\setlength{\parskip}{0.4em}
\setlength{\tabcolsep}{4pt}
\setlength{\headheight}{28pt}
\renewcommand{\arraystretch}{1.05}
\raggedbottom
\setcounter{topnumber}{5}
\setcounter{bottomnumber}{5}
\setcounter{totalnumber}{10}
\renewcommand{\topfraction}{0.95}
\renewcommand{\bottomfraction}{0.95}
\renewcommand{\textfraction}{0.05}
\renewcommand{\floatpagefraction}{0.8}
\setlength{\emergencystretch}{3em}
\hbadness=10000
\vbadness=10000
\hfuzz=60pt
\vfuzz=600pt

% Réduction de l'espace autour des titres
\titleformat{\chapter}[display]
  {\normalfont\huge\bfseries}
  {Chapitre \thechapter}
  {0.4em}
  {}
\titlespacing*{\chapter}{0pt}{-14pt}{10pt}
\titlespacing*{\section}{0pt}{6pt}{6pt}

% Configuration de l'en-tête et du pied de page
\pagestyle{fancy}
\fancyhf{}
\fancyhead[L]{\ifnum\value{chapter}>0 Chapitre \thechapter\fi}
\fancyhead[R]{\TitreProjet}
\fancyfoot[C]{\thepage}
\renewcommand{\headrulewidth}{0.5pt}
\renewcommand{\footrulewidth}{0.5pt}
\fancypagestyle{plain}{%
  \fancyhf{}%
  \fancyhead[L]{\ifnum\value{chapter}>0 Chapitre \thechapter\fi}%
  \fancyhead[R]{\TitreProjet}%
  \fancyfoot[C]{\thepage}%
  \renewcommand{\headrulewidth}{0.5pt}%
  \renewcommand{\footrulewidth}{0.5pt}%
}

% Ligne horizontale personnalisée pour le titre
\newcommand{\HRule}{\rule{\linewidth}{0.4mm}}

% Réglages et types de colonnes pour les tableaux (ltablex)
\newcommand{\SafeRaggedRight}{\rightskip=0pt plus 2em\relax\spaceskip=.3333em\xspaceskip=.5em\relax}
\newcolumntype{L}{>{\SafeRaggedRight\arraybackslash}X}
\newcolumntype{C}[1]{>{\Centering\arraybackslash}p{#1}}
\newcolumntype{Y}{>{\SafeRaggedRight\arraybackslash}X}
\newlength{\TableWidth}
\setlength{\TableWidth}{0.92\textwidth}
\newcolumntype{P}[1]{>{\SafeRaggedRight\hspace{0pt}\arraybackslash}p{#1}}
\newcommand{\logoplaceholder}{\parbox[c][0.9cm][c]{1.0cm}{\centering\scriptsize Logo}}
\newcommand{\chapterfigure}[3][0.8\textwidth]{%
  \begin{figure}[H]
    \centering
    \IfFileExists{#2}{%
      \includegraphics[width=#1]{#2}%
    }{%
      \fbox{\parbox[c][5cm][c]{0.82\textwidth}{\centering\textit{Figure non disponible : #3}}}%
    }%
    \caption{#3}
  \end{figure}
}
\newcommand{\techlogo}[1]{%
  \IfFileExists{#1}{\includegraphics[width=1.45cm]{#1}}{\logoplaceholder}
}
\newcommand{\figureplaceholder}[2]{%
  \begin{figure}[H]
    \centering
    \fbox{\begin{minipage}[c][5cm][c]{0.9\textwidth}
      \centering
      \textit{Espace réservé pour l'insertion de la figure}
    \end{minipage}}
    \caption{#1}
    \label{#2}
  \end{figure}
}

% Macros de priorité
\newcommand{\MustFR}{\textbf{Must}}
\newcommand{\ShouldFR}{\textbf{Should}}
\newcommand{\CouldFR}{\textbf{Could}}

% Style de ligne "Epic" pour les tableaux
\newcommand{\EpicRow}[1]{\rowcolor{gray!15}\multicolumn{6}{l}{\textbf{#1}}\\}

\newcommand{\placeholderfigure}[2]{%
  \begin{figure}[H]
    \centering
    \IfFileExists{images/#1}{%
      \includegraphics[width=0.85\textwidth]{images/#1}%
    }{%
      \fbox{\parbox[c][6cm][c]{0.82\textwidth}{\centering Emplacement réservé pour la figure}}%
    }
    \caption{#2}
  \end{figure}
}


\usepackage{listings}


\lstdefinelanguage{json}{
    basicstyle=\ttfamily\small,
    sensitive=false,
    morestring=[b]",
    morecomment=[l]{//}
}
\lstdefinelanguage{yaml}{
    basicstyle=\ttfamily\small,
    sensitive=false,
    morecomment=[l]{\#},
    morestring=[b]',
    morestring=[b]"
}
\lstdefinelanguage{properties}{
    basicstyle=\ttfamily\small,
    sensitive=false,
    morecomment=[l]{\#},
    morecomment=[l]{!},
    morestring=[b]"
}



% Configuration globale pour le code (Prompt et JSON)
\lstset{
    basicstyle=\ttfamily\small,
    breaklines=true,
    frame=single,
    backgroundcolor=\color{gray!5},
    showstringspaces=false,
    captionpos=b,
    columns=fullflexible,
    literate=
    {à}{{\`a}}1 {â}{{\^a}}1 {ä}{{\"a}}1
    {ç}{{\c{c}}}1
    {é}{{\'e}}1 {è}{{\`e}}1 {ê}{{\^e}}1 {ë}{{\"e}}1
    {î}{{\^\i}}1 {ï}{{\"\i}}1
    {ô}{{\^o}}1 {ö}{{\"o}}1
    {ù}{{\`u}}1 {û}{{\^u}}1 {ü}{{\"u}}1
    {À}{{\`A}}1 {Â}{{\^A}}1 {Ä}{{\"A}}1
    {Ç}{{\c{C}}}1
    {É}{{\'E}}1 {È}{{\`E}}1 {Ê}{{\^E}}1 {Ë}{{\"E}}1
    {Î}{{\^I}}1 {Ï}{{\"I}}1
    {Ô}{{\^O}}1 {Ö}{{\"O}}1
    {Ù}{{\`U}}1 {Û}{{\^U}}1 {Ü}{{\"U}}1
}
% =========================================================================
% 4. DÉBUT DU DOCUMENT
% =========================================================================
\begin{document}

% -
% PAGE DE GARDE
% -

\begin{titlepage}
\thispagestyle{empty}

%--- EN-TÊTE ---
\noindent
\begin{tabular}{@{}p{0.42\textwidth}@{}c@{}p{0.42\textwidth}@{}}
\begin{minipage}[t]{0.42\textwidth}
\small République Tunisienne\\
\small Ministère de l'Enseignement Supérieur et\\
\small de la Recherche Scientifique
\end{minipage}
&
\begin{minipage}[t]{0.16\textwidth}
\centering
\includegraphics[height=1.8cm]{\LogoISIMA}
\end{minipage}
&
\begin{minipage}[t]{0.42\textwidth}
\raggedleft
\small Université de Monastir\\
\small Institut Supérieur d'Informatique de\\
\small Mahdia
\end{minipage}
\end{tabular}

\vspace{0.3cm}
\noindent\rule{\textwidth}{0.8pt}

\begin{flushright}
  \textbf{\textit{MEMOIRE}}\\[1pt]
  \textbf{N° d'ordre :} LBC-BI-3-2026-032
\end{flushright}

\noindent\rule{\textwidth}{0.8pt}

%--- TITRE PRINCIPAL ---
\vspace{0.22cm}
\begin{center}
  {\fontsize{36}{42}\selectfont \textit{M\'{E}MOIRE}}
\end{center}

\vspace{0.22cm}

%--- CORPS ---
\begin{center}
  \textit{Présenté à}\\[0.2cm]
  {\large \textbf{Institut Supérieur d'Informatique de Mahdia}}\\[0.8cm]
  \textit{Effectué à}\\[0.25cm]
  \textbf{AgileX Management}\\[0.25cm]
  \textit{En vue de l'obtention du diplôme de}\\[1.2cm]
  \textit{Elaboré par :}\\[0.25cm]
  Mohamed Abdeladhim, Rahma Barouni
\end{center}

\vspace{0.5cm}
\noindent\rule{\textwidth}{2pt}

\begin{center}
  \vspace{0.08cm}
  \textbf{TalentPredict}
  \vspace{0.08cm}
\end{center}

\noindent\rule{\textwidth}{2pt}

\vspace{0.8cm}
\begin{center}
  \textit{\textbf{Soutenu le :}} 11 juin 2026\\[0.08cm]
  \textit{\textbf{devant le jury composé de :}}
\end{center}

\vspace{0.6cm}

%--- JURY ---
\begin{tabular}{ll}
  \textbf{Président}               & : Tarek M. Hamdani \\[1pt]
  \textbf{Rapporteur}              & : Emna Ben Ayed \\[1pt]
  \textbf{Encadrant Académique}    & : Ben Aichaa Takwa \\[1pt]
  \textbf{Encadrant Professionnel} & : Med Amine Gharbi \\
\end{tabular}

\vfill

\begin{center}
  Année universitaire 2025-2026
\end{center}
\end{titlepage}

% -------------------------------------------------------------------------
% DÉDICACE
% -------------------------------------------------------------------------
\chapter*{Dédicace}
\thispagestyle{empty}

\vspace*{1cm}

\begin{center}
\itshape
À celles et ceux qui nous ont soutenus sans réserve,\\
à nos familles pour leurs sacrifices,\\
à nos proches pour leur confiance et leur patience.

\vspace{0.4cm}

Pour leur encouragement constant, leur bienveillance et leur présence\\
dans les moments d’effort comme dans les instants de doute.

\vspace{0.6cm}

Ce travail leur est dédié.
\end{center}

\vspace{1cm}

\begin{flushright}
\large\itshape
\EtudiantUn \\
\EtudiantDeux
\end{flushright}

\clearpage

% -------------------------------------------------------------------------
% REMERCIEMENTS
% -------------------------------------------------------------------------
\chapter*{Remerciements} 
\thispagestyle{empty}

Nous tenons à exprimer notre profonde gratitude à toutes les personnes qui ont contribué, de près ou de loin, à la réalisation de ce projet.

Nos sincères remerciements s’adressent à notre encadrant pour son accompagnement rigoureux, ses conseils avisés et sa disponibilité tout au long de ce travail.

Nous remercions également les membres du jury pour l’intérêt qu’ils portent à ce mémoire et pour le temps consacré à son évaluation.

Notre reconnaissance va aussi à l’ensemble des enseignants et du personnel de l’Institut Supérieur d’Informatique de Mahdia (ISIMA) pour la qualité de la formation reçue.

Enfin, nous adressons une pensée particulière à nos familles et à nos proches pour leur soutien indéfectible et leurs encouragements constants.

\clearpage

% -------------------------------------------------------------------------
% RÉSUMÉ / ABSTRACT
% -------------------------------------------------------------------------
\chapter*{Résumé}


% -------------------------------------------------------------------------
% TABLE DES MATIÈRES ET LISTES
% -------------------------------------------------------------------------
\tableofcontents
\clearpage

\listoffigures
\clearpage

\listoftables
\clearpage

% -------------------------------------------------------------------------
% LISTE DES ABRÉVIATIONS
% -------------------------------------------------------------------------
\chapter*{Liste des Abréviations}
\addcontentsline{toc}{chapter}{Liste des Abréviations}

\begin{longtable}{|P{0.20\textwidth}|P{0.72\textwidth}|}
\hline
\rowcolor{TableHeaderBlue}
\color{white}\textbf{Acronyme} & \color{white}\textbf{Signification} \\
\hline
\endfirsthead

\hline
\rowcolor{TableHeaderBlue}
\color{white}\textbf{Acronyme} & \color{white}\textbf{Signification} \\
\hline
\endhead

\hline
\endfoot

\textbf{AI} & Intelligence Artificielle (\textit{Artificial Intelligence}) \\
\hline
\textbf{API} & Interface de Programmation Applicative \\
\hline
\textbf{ASGI} & Interface de Passerelle de Serveur Asynchrone (\textit{Asynchronous Server Gateway Interface}) \\
\hline
\textbf{BCrypt} & Algorithme de hachage de mots de passe adaptatif \\
\hline
\textbf{BI} & Intelligence d'Affaires (\textit{Business Intelligence}) \\
\hline
\textbf{CI/CD} & Intégration Continue / Déploiement Continu (\textit{Continuous Integration / Continuous Deployment}) \\
\hline
\textbf{CRF} & Champ Aléatoire Conditionnel (\textit{Conditional Random Field}) \\
\hline
\textbf{CRUD} & Créer, Lire, Mettre à jour, Supprimer (\textit{Create, Read, Update, Delete}) \\
\hline
\textbf{CSS} & Feuilles de Style en Cascade (\textit{Cascading Style Sheets}) \\
\hline
\textbf{CVE} & Vulnérabilités et Expositions Communes (\textit{Common Vulnerabilities and Exposures}) \\
\hline
\textbf{DISC} & Modèle psychométrique Dominance, Influence, Stabilité, Conformité \\
\hline
\textbf{DTO} & Objet de Transfert de Données (\textit{Data Transfer Object}) \\
\hline
\textbf{ERD} & Diagramme Entité-Relation (\textit{Entity-Relationship Diagram}) \\
\hline
\textbf{GPEC} & Gestion Prévisionnelle des Emplois et des Compétences \\
\hline
\textbf{HTML} & Langage de Balisage Hypertexte (\textit{HyperText Markup Language}) \\
\hline
\textbf{HTTP} & Protocole de Transfert Hypertexte (\textit{HyperText Transfer Protocol}) \\
\hline
\textbf{HTTPS} & Protocole de Transfert Hypertexte Sécurisé (\textit{HyperText Transfer Protocol Secure}) \\
\hline
\textbf{IaC} & Infrastructure en tant que Code (\textit{Infrastructure as Code}) \\
\hline
\textbf{IAM} & Gestion des Identités et des Accès (\textit{Identity and Access Management}) \\
\hline
\textbf{JPA} & Interface de Persistance Java (\textit{Java Persistence API}) \\
\hline
\textbf{JSON} & Notation Objet JavaScript \\
\hline
\textbf{JWT} & Jeton Web JSON \\
\hline
\textbf{KPI} & Indicateur Clé de Performance (\textit{Key Performance Indicator}) \\
\hline
\textbf{LLM} & Grand Modèle de Langage (\textit{Large Language Model}) \\
\hline
\textbf{LXP} & Plateforme d'Expérience d'Apprentissage (\textit{Learning Experience Platform}) \\
\hline
\textbf{MCQ / QCM} & Questionnaire à Choix Multiples (\textit{Multiple Choice Questions}) \\
\hline
\textbf{ML} & Apprentissage Automatique (\textit{Machine Learning}) \\
\hline
\textbf{MVC} & Modèle-Vue-Contrôleur (\textit{Model-View-Controller}) \\
\hline
\textbf{MVP} & Produit Minimum Viable (\textit{Minimum Viable Product}) \\
\hline
\textbf{NER} & Reconnaissance d'Entités Nommées (\textit{Named Entity Recognition}) \\
\hline
\textbf{NLP} & Traitement Automatique du Langage Naturel (\textit{Natural Language Processing}) \\
\hline
\textbf{ORM} & Mapping Objet-Relationnel (\textit{Object-Relational Mapping}) \\
\hline
\textbf{PCM} & Modèle de Communication de Processus (\textit{Process Communication Model}) \\
\hline
\textbf{PDF} & Format de Document Portable (\textit{Portable Document Format}) \\
\hline
\textbf{QA} & Assurance Qualité (\textit{Quality Assurance}) \\
\hline
\textbf{RBAC} & Contrôle d'Accès Basé sur les Rôles (\textit{Role-Based Access Control}) \\
\hline
\textbf{REST} & Transfert d'État Représentationnel (\textit{Representational State Transfer}) \\
\hline
\textbf{RGPD} & Règlement Général sur la Protection des Données \\
\hline
\textbf{RH} & Ressources Humaines \\
\hline
\textbf{SaaS} & Logiciel en tant que Service (\textit{Software as a Service}) \\
\hline
\textbf{SGBD} & Système de Gestion de Base de Données \\
\hline
\textbf{SIRH} & Système d'Information des Ressources Humaines \\
\hline
\textbf{SMTP} & Protocole Simple de Transfert de Courrier (\textit{Simple Mail Transfer Protocol}) \\
\hline
\textbf{SPA} & Application à Page Unique (\textit{Single Page Application}) \\
\hline
\textbf{SQL} & Langage de Requête Structuré (\textit{Structured Query Language}) \\
\hline
\textbf{SRP} & Principe de Responsabilité Unique (\textit{Single Responsibility Principle}) \\
\hline
\textbf{SSE} & Événements Envoyés par le Serveur (\textit{Server-Sent Events}) \\
\hline
\textbf{SSR} & Rendu Côté Serveur (\textit{Server-Side Rendering}) \\
\hline
\textbf{TLS} & Sécurité de la Couche de Transport (\textit{Transport Layer Security}) \\
\hline
\textbf{UML} & Langage de Modélisation Unifié (\textit{Unified Modeling Language}) \\
\hline
\textbf{UUID} & Identifiant Universel Unique (\textit{Universally Unique Identifier}) \\
\hline
\textbf{XHTML} & Langage de Balisage Hypertexte Extensible (\textit{Extensible HyperText Markup Language}) \\
\hline
\textbf{XSS} & Script Intersite (\textit{Cross-Site Scripting}) \\
\hline
\caption{Liste des abréviations utilisées dans le rapport}
\end{longtable}

\clearpage

% -------------------------------------------------------------------------
% CORPS DU DOCUMENT
% -------------------------------------------------------------------------
\pagenumbering{arabic}
\chapter*{Introduction Générale}
\addcontentsline{toc}{chapter}{Introduction Générale}

Dans un contexte professionnel marqué par la transformation numérique des processus métiers et par l’intégration croissante de l’intelligence artificielle dans les systèmes d’aide à la décision, la gestion des compétences constitue aujourd’hui un enjeu stratégique pour les organisations. Les entreprises doivent non seulement identifier les compétences disponibles en interne, mais aussi évaluer leur niveau de maîtrise, détecter les écarts par rapport aux besoins métier et proposer des parcours de développement adaptés. Or, les méthodes classiques d’évaluation des talents, souvent fondées sur des CV déclaratifs, des entretiens ponctuels ou des appréciations subjectives, montrent leurs limites lorsqu’il s’agit de produire une analyse objective, continue et exploitable à grande échelle.

C’est dans ce cadre que s’inscrit le présent projet de fin d’études, réalisé au sein de la société AgileX Management et intitulé TalentPredict AI. Ce projet vise à concevoir et développer une plateforme intelligente dédiée à l’évaluation automatisée des compétences humaines, en prenant en compte à la fois les compétences techniques et les compétences comportementales. La problématique centrale peut ainsi être formulée comme suit : comment mettre en place une solution logicielle capable d’analyser de manière structurée, fiable et sécurisée les profils des collaborateurs, tout en générant des recommandations de formation personnalisées et utiles à la prise de décision RH ?

La solution proposée repose sur une approche hybride combinant plusieurs sources d’information et plusieurs traitements complémentaires. L’évaluation des compétences techniques s’appuie sur l’analyse du CV, l’exploitation du profil GitHub, des tests QCM adaptatifs et des challenges de code. L’évaluation des compétences comportementales mobilise le modèle psychométrique PCM, l’analyse sémantique locale à l’aide d’Ollama, ainsi que des scénarios d’évaluation générés par l’intelligence artificielle. Ces traitements sont intégrés dans une architecture moderne articulée autour de n8n pour l’orchestration des workflows, FastAPI pour les services d’intelligence artificielle, Spring Boot pour le backend métier, Angular pour l’interface utilisateur et Docker pour la conteneurisation. Ce choix architectural vise à assurer la modularité de la plateforme, la portabilité de l’environnement, la souveraineté des données et une meilleure conformité aux exigences de confidentialité, notamment celles liées au RGPD.

Ce mémoire est structuré en six chapitres. Le premier chapitre présente le cadre général du projet, l’organisme d’accueil, la problématique métier, les objectifs, les besoins fonctionnels et non fonctionnels, ainsi que la méthodologie Agile Scrum adoptée. Le deuxième chapitre propose un état de l’art des concepts, technologies et solutions existantes liés à l’évaluation intelligente des compétences et au domaine de la HR Tech. Les chapitres suivants décrivent le déroulement du projet selon les sprints de développement : le Sprint 1 est consacré à la mise en place de l’architecture, de la sécurité et du système d’authentification ; le Sprint 2 porte sur le développement du moteur d’évaluation multi-sources des compétences ; le Sprint 3 traite du système de recommandation et d’upskilling piloté par l’intelligence artificielle ; enfin, le Sprint 4 présente les tableaux de bord analytiques, le pilotage stratégique et la mise en production de la plateforme.
\clearpage

% =========================================================
% Chapitre 1
% =========================================================

\chapter{Cadre général du projet}

\section{Introduction}
\phantomsection
\addcontentsline{toc}{section}{Introduction}
\phantomsection 

Ce chapitre présente le cadre général dans lequel s’inscrit le projet TalentPredict AI. Il introduit d’abord l’organisme d’accueil et le contexte métier ayant motivé la conception de la solution. Il précise ensuite la problématique, les objectifs, les besoins fonctionnels et non fonctionnels, ainsi que les acteurs impliqués. Enfin, il décrit la méthodologie Scrum adoptée, l’organisation de l’équipe projet et les principaux éléments de planification qui structurent le déroulement du travail.

\section{Présentation du contexte et de l’enjeu}

La compréhension du contexte organisationnel constitue une étape essentielle pour situer le projet dans son environnement réel. Cette section présente l’entreprise d’accueil, ses domaines d’expertise et les enjeux métiers auxquels le projet tente d’apporter une réponse. Elle permet également de montrer la cohérence entre les besoins de l’entreprise et la solution proposée.

\subsection{Présentation de l’organisme d’accueil}

L’organisme d’accueil joue un rôle important dans l’orientation fonctionnelle et technique du projet. Sa culture de travail, ses pratiques de développement et ses priorités stratégiques influencent directement les choix retenus durant la conception et la réalisation de la plateforme.

\subsubsection{Présentation générale d’AgileX}

AgilEX est une société spécialisée dans le développement de solutions logicielles sur mesure et dans l’accompagnement des entreprises dans leur transformation numérique. Elle intervient principalement dans la conception d’applications web et mobiles, l’intégration de systèmes, le conseil en architecture logicielle et le déploiement de solutions cloud.

Le nom AgileX reflète une orientation vers l’agilité, l’adaptation et l’excellence technique. L’entreprise privilégie des approches de développement itératives permettant de produire rapidement des incréments fonctionnels, de réduire les risques et d’intégrer progressivement les retours des parties prenantes.

La figure~\ref{fig:logo_orga} présente le logo de l’entreprise d’accueil. Elle permet d’identifier visuellement l’organisme dans lequel le projet a été réalisé.

\begin{figure}[H]
  \centering
  \includegraphics[width=0.3\textwidth]{images/logo_orga.png}
  \caption{Logo de l’organisme d’accueil AgileX }
  \label{fig:logo_orga}
\end{figure}
%\chapterfigure[0.22\textwidth]{images/agilx.png}{Logo de l’organisme d’accueil AgileX}

La culture de travail d’Agilx repose sur une organisation AgileX centrée sur la collaboration, la livraison continue de valeur et l’amélioration progressive des solutions développées. Cette culture est en adéquation avec le choix de la méthodologie Scrum retenue pour piloter le projet.

La figure~\ref{fig:struc_orga} illustre l’organisation structurelle de l’entreprise. Elle met en évidence les relations entre les pôles de travail et permet de mieux comprendre le contexte professionnel dans lequel s’inscrit le projet.
\begin{figure}[H]
  \centering
  \includegraphics[width=0.75\textwidth]{images/philosophie_orga.png}
  \caption{Organisation structurelle de l’entreprise AgileX}
  \label{fig:struc_orga}
\end{figure}


Le tableau~\ref{tab:fiche_signal_agilx} présente une fiche signalétique synthétique de l’entreprise. Il regroupe les informations principales permettant de situer l’organisme d’accueil du point de vue administratif, sectoriel et méthodologique.

\begin{longtable}{|P{0.32\textwidth}|P{0.58\textwidth}|}
\label{tab:fiche_signal_agilx}\\
\hline
\rowcolor{TableHeaderBlue}
\color{white}\textbf{Élément} & \color{white}\textbf{Information} \\
\hline
\endfirsthead

\hline
\rowcolor{TableHeaderBlue}
\color{white}\textbf{Élément} & \color{white}\textbf{Information} \\
\hline
\endhead

\hline
\endfoot

Raison sociale & AgileX \\
\hline
Secteur d’activité & Développement de logiciels et solutions digitales \\
\hline
Domaines d’intervention & Applications web et mobiles, architecture logicielle, intelligence artificielle, cloud et intégration de systèmes \\
\hline
Méthodologie de travail & Approche Agile, principalement inspirée de Scrum~\cite{scrum2020} \\
\hline
Encadrant professionnel & M. Mohamed Amine Gharbi, Manager SI \\
\hline
\caption{Fiche signalétique de l’entreprise AgileX}
\end{longtable}

\subsubsection{Portefeuille de services et domaines d’expertise}

L’offre de services d’AgileX s’articule autour de plusieurs domaines complémentaires. Cette diversité d’expertise a constitué un environnement favorable pour la réalisation de TalentPredict AI   , dans la mesure où le projet combine développement web, architecture logicielle, intelligence artificielle, orchestration de workflows et déploiement conteneurisé.

Le tableau~\ref{tab:portefeuille_agilx} synthétise les principaux pôles d’expertise de l’entreprise et précise leur contribution au projet. Cette présentation permet de relier les compétences de l’organisme d’accueil aux choix techniques et fonctionnels adoptés dans la solution.

\begin{longtable}{|P{0.24\textwidth}|P{0.34\textwidth}|P{0.32\textwidth}|}

\label{tab:portefeuille_agilx}\\
\hline
\rowcolor{TableHeaderBlue}
\color{white}\textbf{Pôle d’expertise} &
\color{white}\textbf{Description} &
\color{white}\textbf{Contribution à \textit{TalentPredict AI}} \\
\hline
\endfirsthead

\hline
\rowcolor{TableHeaderBlue}
\color{white}\textbf{Pôle d’expertise} &
\color{white}\textbf{Description} &
\color{white}\textbf{Contribution à \textit{TalentPredict AI}} \\
\hline
\endhead

\hline
\endfoot

Développement Full-Stack et mobilité &
Conception et développement d’applications web modernes, d’applications mobiles et de plateformes SaaS. &
Mise en œuvre du frontend Angular, des interfaces utilisateurs et des interactions avec le backend. \\
\hline

Ingénierie Cloud et DevOps &
Mise en place d’environnements cloud, conteneurisation, automatisation et amélioration de la portabilité des applications. &
Utilisation de Docker pour faciliter le déploiement, l’isolation des services et la reproductibilité de l’environnement. \\
\hline

Conseil en architecture logicielle &
Définition d’architectures robustes, évolutives et adaptées aux besoins métiers. &
Adoption d’une architecture modulaire séparant l’interface, la logique métier, les services IA et l’orchestration. \\
\hline

Intégration d’intelligence artificielle &
Intégration de capacités d’analyse sémantique, d’automatisation et d’aide à la décision dans les processus métiers. &
Utilisation d’Ollama, de FastAPI et de n8n pour l’analyse des compétences et la génération de recommandations. \\
\hline
\caption{Portefeuille de services et contribution au projet}
\end{longtable}

\subsection{Enjeu général du projet}

Dans les organisations modernes, la gestion des compétences ne se limite plus à la consultation des CV ou à l’organisation d’entretiens ponctuels. Elle nécessite une vision dynamique, objective et exploitable des compétences réellement disponibles au sein de l’entreprise. Cette exigence devient particulièrement importante dans les métiers du numérique, où les technologies évoluent rapidement et où les besoins en formation doivent être identifiés de manière continue.

Le projet \textit{TalentPredict AI} répond à cet enjeu en proposant une plateforme capable de croiser plusieurs sources d’information afin d’évaluer les compétences techniques et comportementales des collaborateurs. L’objectif n’est pas de remplacer l’expertise humaine des responsables RH, mais de fournir un outil d’aide à la décision capable de structurer les analyses, de réduire les biais déclaratifs et d’orienter les parcours de développement professionnel.

\section{Présentation du projet}

Après la présentation de l’organisme d’accueil, il est nécessaire de préciser la nature du projet, ses objectifs et les acteurs concernés. Cette section permet de délimiter le périmètre fonctionnel de la solution et d’expliciter la valeur attendue pour les utilisateurs finaux et les décideurs RH.

\subsection{Contexte et problématique métier}

Dans le secteur informatique, l’évaluation des compétences constitue un enjeu stratégique pour les entreprises qui souhaitent mieux connaître les capacités réelles de leurs collaborateurs . Cependant, les informations nécessaires à cette évaluation sont souvent réparties entre plusieurs sources hétérogènes, telles que les CV, les profils GitHub, les profils LinkedIn, les portfolios, les tests techniques, les feedbacks RH et les questionnaires comportementaux. Cette dispersion rend le processus d’évaluation plus long, plus complexe et moins facilement exploitable dans une démarche de formation continue.

Comme l’illustre la figure~\ref{fig:problematique-metier-talentpredict}, le problème principal réside dans la fragmentation des données et dans l’absence d’un mécanisme unifié permettant de croiser les informations disponibles. Les évaluations traditionnelles restent souvent dépendantes des déclarations du collaborateur, ce qui peut introduire des biais ou limiter la fiabilité des résultats. De plus, les compétences comportementales sont généralement plus difficiles à mesurer, car elles nécessitent une analyse plus fine du profil, des interactions, des expériences et des réponses à des situations professionnelles.

\begin{figure}[H]
    \centering
    \includegraphics[width=\textwidth]{images/contexte.png}
    \caption{Problématique métier liée à la fragmentation des données d’évaluation des compétences}
    \label{fig:problematique-metier-talentpredict}
\end{figure}

Ces limites peuvent conduire à des décisions RH moins précises, à des parcours de formation peu personnalisés et à une exploitation insuffisante du potentiel interne. Dans ce contexte, le besoin d’une plateforme intelligente devient essentiel afin de passer d’une évaluation dispersée et déclarative à une évaluation unifiée, multi-sources et exploitable.

La problématique centrale du projet peut donc être formulée comme suit : comment concevoir et développer une plateforme intelligente permettant d’évaluer de manière structurée les compétences techniques et comportementales des collaborateurs, tout en générant des recommandations de formation personnalisées, sécurisées et exploitables par l’entreprise ?

\subsection{Objectifs du projet}

L’objectif principal du projet est de concevoir et développer TalentPredict AI, une plateforme intelligente dédiée à l’évaluation et au développement des compétences. La solution vise à fournir une analyse plus complète du profil collaborateur en combinant des données déclaratives, des preuves techniques observables et des traitements d’intelligence artificielle. Elle permet ainsi de renforcer l’objectivité du processus d’évaluation et de produire des résultats plus utiles à la prise de décision RH.

Les objectifs spécifiques du projet sont les suivants :

\begin{itemize}[leftmargin=0.9cm]
    \item analyser les compétences techniques à partir du CV, du profil GitHub, de tests QCM et de challenges de code ;
    \item évaluer les compétences comportementales à travers le modèle PCM, l’analyse sémantique locale et des scénarios générés par l’intelligence artificielle ;
    \item proposer des recommandations de formation adaptées aux écarts de compétences identifiés ;
    \item offrir un tableau de bord permettant aux employés et aux responsables RH de suivre les résultats et l’évolution des compétences ;
    \item garantir un traitement sécurisé des données personnelles et professionnelles, notamment à travers une inférence locale et une architecture maîtrisée.
\end{itemize}

Ces objectifs permettent à la plateforme de ne pas se limiter à une simple notation des compétences. Ils orientent plutôt la solution vers une logique d’accompagnement, dans laquelle l’évaluation devient un point de départ pour construire des parcours de progression individualisés et alignés avec les besoins de l’entreprise.

\subsection{Acteurs concernés}

Le projet implique plusieurs catégories d’acteurs, ayant chacune des besoins, des responsabilités et des interactions spécifiques avec la plateforme. L’identification de ces acteurs est indispensable pour définir les fonctionnalités attendues, organiser les droits d’accès et assurer une communication claire entre les utilisateurs humains et les services automatisés.

Les principaux acteurs concernés sont l’employé, l’administrateur RH et les services automatisés intégrés à la plateforme. L’employé utilise la solution pour soumettre ses données, évaluer ses compétences, consulter ses résultats et suivre ses parcours de formation. L’administrateur RH supervise les profils, valide certaines informations, consulte les indicateurs globaux et exploite les rapports générés pour orienter les décisions de gestion des talents. Les services automatisés assurent l’analyse intelligente, l’orchestration des traitements, l’envoi des notifications et la communication en temps réel.

La figure~\ref{fig:acteurs-talentpredict} présente une vue synthétique des acteurs concernés et de leurs interactions avec \textit{TalentPredict AI}. Elle met en évidence le rôle central de la plateforme, qui coordonne les échanges entre les utilisateurs, les services d’intelligence artificielle, le service mail et le mécanisme de communication temps réel.

\begin{figure}[H]
    \centering
    \includegraphics[width=\textwidth]{images/acteurs.png}
    \caption{Acteurs concernés et interactions principales autour de la plateforme TalentPredict AI}
    \label{fig:acteurs-talentpredict}
\end{figure}

Cette organisation permet de mieux structurer les responsabilités de chaque acteur dans le système. Elle garantit également une séparation claire entre les actions humaines, comme la consultation des résultats ou la supervision RH, et les traitements automatisés, comme l’analyse des données, la génération des recommandations et l’envoi des notifications.

\section{Spécification des besoins}

La spécification des besoins permet de traduire la problématique métier en exigences concrètes. Elle distingue les besoins fonctionnels, qui décrivent les services attendus par les utilisateurs, et les besoins non fonctionnels, qui définissent les contraintes de qualité, de sécurité, de performance et d’évolutivité du système.

\subsection{Besoins fonctionnels}

Les besoins fonctionnels définissent les principales fonctionnalités attendues de la plateforme. Ils couvrent l’authentification, l’analyse des compétences, la recommandation de formation, le suivi des parcours et la restitution des résultats à travers des tableaux de bord.

\subsubsection{Module d’authentification et de gestion des utilisateurs}

Ce module constitue le point d’entrée sécurisé de la plateforme. Il permet aux utilisateurs de créer un compte, de se connecter, de gérer leur profil et d’accéder aux fonctionnalités correspondant à leur rôle. Il repose sur une authentification par jeton JWT et sur une gestion différenciée des rôles USER et ADMIN.

Les fonctionnalités attendues sont les suivantes :

\begin{itemize}[leftmargin=0.9cm]
    \item inscription avec adresse email et mot de passe ;
    \item connexion sécurisée avec génération d’un jeton JWT ;
    \item modification des informations du profil utilisateur ;
    \item gestion des rôles par l’administrateur ;
    \item récupération de mot de passe à travers un processus contrôlé.
\end{itemize}

\subsubsection{Module d’analyse des compétences comportementales}

Ce module vise à analyser les dimensions comportementales du collaborateur en combinant un modèle psychométrique, des données textuelles et une analyse sémantique locale. Le modèle PCM~\cite{kahler1996} est utilisé comme cadre de référence pour structurer les profils comportementaux, tandis qu’Ollama~\cite{ollama}, n8n~\cite{n8n} et FastAPI~\cite{fastapi} interviennent dans l’orchestration et l’analyse des données.

Les fonctionnalités attendues sont les suivantes :

\begin{itemize}[leftmargin=0.9cm]
    \item extraction sécurisée du texte des CV côté client ;
    \item analyse du profil GitHub afin d’identifier certains indicateurs comportementaux observables ;
    \item passation d’un test PCM ;
    \item orchestration d’un flux d’analyse à travers n8n ;
    \item calcul de scores liés à la communication, au leadership, à la collaboration, à la curiosité, à l’ownership et à la discipline ;
    \item identification du type de personnalité dominant ;
    \item génération de scénarios d’évaluation comportementale ;
    \item historisation des résultats afin de suivre l’évolution du profil.
\end{itemize}

\subsubsection{Module d’analyse des compétences techniques}

Ce module a pour objectif d’évaluer les compétences techniques du collaborateur à partir de plusieurs sources complémentaires. L’intérêt de cette approche multi-sources est de réduire la dépendance aux déclarations du CV et de confronter les compétences annoncées à des preuves techniques plus observables.

Les fonctionnalités attendues sont les suivantes :

\begin{itemize}[leftmargin=0.9cm]
    \item extraction et parsing du CV~\cite{celik2013,chien2020} ;
    \item analyse de l’activité GitHub~\cite{kalliamvakou2014} et des langages utilisés ;
    \item prise en compte éventuelle du profil LinkedIn ou d’un portfolio ;
    \item génération de QCM techniques adaptatifs à l’aide d’Ollama~\cite{ollama} ;
    \item génération et évaluation automatique de challenges de code avec Code Llama~\cite{roziere2023codellama} ;
    \item calcul d’un score de confiance par fusion multi-sources ;
    \item détection des incohérences entre compétences déclarées et compétences observées ;
    \item attribution d’un niveau de maîtrise par compétence.
\end{itemize}

\subsubsection{Module de recommandations et de développement des compétences}

Ce module transforme les résultats de l’évaluation en plans d’action exploitables. Il ne se limite pas à produire un score, mais propose des parcours de développement adaptés aux écarts de compétences identifiés et au profil du collaborateur.

Pour l’employé, le module permet de consulter des recommandations de formation, de suivre une roadmap d’apprentissage, de visualiser sa progression et de passer des mini-tests de validation. Pour l’administrateur RH, il permet de suivre les parcours, de valider ou refuser certaines demandes, de proposer des formations ciblées et d’observer les taux d’avancement.

\begin{itemize}[leftmargin=0.9cm]
    \item génération de recommandations personnalisées de formation~\cite{drachsler2009} ;
    \item construction d’une roadmap d’apprentissage ;
    \item suivi des statuts des formations ;
    \item intégration de mécanismes de progression et de gamification ;
    \item validation des acquis par mini-test ;
    \item génération d’un certificat interne lorsque les conditions sont satisfaites ;
    \item supervision des parcours par l’administrateur RH.
\end{itemize}

\subsubsection{Module de tableau de bord analytique}

Le tableau de bord constitue l’espace de restitution des résultats. Il permet de présenter les analyses sous une forme lisible, synthétique et exploitable. Les visualisations doivent être adaptées à la fois aux employés, qui consultent leurs propres résultats, et aux administrateurs RH, qui suivent les compétences à l’échelle de l’organisation.

Les fonctionnalités attendues sont les suivantes :

\begin{itemize}[leftmargin=0.9cm]
    \item visualisation des compétences à travers des radars, histogrammes et courbes d’évolution ;
    \item consultation du profil comportemental et des résultats techniques ;
    \item suivi de l’avancement des formations ;
    \item génération d’un Talent Passport individuel ;
    \item cartographie globale des compétences ;
    \item identification des écarts de compétences ;
    \item génération de rapports RH consolidés.
\end{itemize}

\subsection{Besoins non fonctionnels}

Les besoins non fonctionnels définissent les exigences de qualité qui encadrent le fonctionnement de la plateforme. Ils garantissent que la solution ne soit pas seulement fonctionnelle, mais également sécurisée, maintenable, performante et adaptée à une utilisation réelle en entreprise.

Les principales exigences non fonctionnelles sont résumées dans le tableau~\ref{tab:besoins_non_fonctionnels}. Elles couvrent les dimensions de performance, sécurité, disponibilité, maintenabilité, ergonomie, interopérabilité et confidentialité.

\begin{longtable}{|P{0.24\textwidth}|P{0.66\textwidth}|}
\label{tab:besoins_non_fonctionnels}\\
\hline
\rowcolor{TableHeaderBlue}
\color{white}\textbf{Critère} & \color{white}\textbf{Exigence} \\
\hline
\endfirsthead
\hline
\rowcolor{TableHeaderBlue}
\color{white}\textbf{Critère} & \color{white}\textbf{Exigence} \\
\hline
\endhead
\hline
\endfoot
Performance &
Les opérations courantes de l’API doivent répondre dans un délai raisonnable. Les traitements LLM, plus coûteux, doivent être accompagnés d’un retour visuel afin d’informer l’utilisateur de l’avancement. \\
\hline
Sécurité &
L’authentification doit reposer sur \texttt{JWT}, les mots de passe doivent être hachés avec \texttt{BCrypt}, et les accès doivent être contrôlés selon les rôles. Les protections contre les injections SQL, XSS et CSRF doivent être prises en compte. \\
\hline
Disponibilité et fiabilité &
La plateforme doit viser une disponibilité élevée, prévoir des mécanismes de reprise en cas d’erreur et limiter les interruptions lors des appels aux services externes ou aux modules IA. \\
\hline
Maintenabilité &
L’architecture doit être modulaire, documentée et organisée selon des conventions de codage claires pour Java, Python et Angular. \\
\hline
Évolutivité &
La solution doit pouvoir intégrer de nouveaux modèles IA, de nouvelles sources de données et de nouveaux modules fonctionnels sans remise en cause majeure de l’architecture. \\
\hline
Ergonomie et accessibilité &
L’interface doit être intuitive, responsive et adaptée à des utilisateurs non techniques. Les visualisations doivent être lisibles, cohérentes et accompagnées de légendes explicites. \\
\hline
Interopérabilité &
Les échanges entre services doivent reposer sur des API REST et des formats structurés, principalement JSON. Les appels aux modèles IA doivent être encapsulés pour faciliter un éventuel changement de fournisseur. \\
\hline
Confidentialité et conformité &
Le système doit limiter l’exposition des données sensibles, favoriser le traitement local lorsque cela est possible, recueillir le consentement explicite des utilisateurs et respecter les principes du RGPD~\cite{rgpd2016}. \\
\hline
\caption{Synthèse des besoins non fonctionnels du projet}
\end{longtable}
\subsection{Identification des utilisateurs et des rôles}

L’identification des utilisateurs permet de définir les responsabilités, les droits d’accès et les interactions attendues avec la plateforme. Dans le cadre de \textit{TalentPredict AI}, les acteurs ne sont pas uniquement humains : certains services automatisés jouent également un rôle central dans le fonctionnement global du système.

Le tableau~\ref{tab:acteurs_systeme} présente les principaux acteurs du système, leur rôle et leurs interactions clés.

\begin{longtable}{|P{0.18\textwidth}|P{0.20\textwidth}|P{0.27\textwidth}|P{0.25\textwidth}|}
\label{tab:acteurs_systeme}\\
\hline
\rowcolor{TableHeaderBlue}
\color{white}\textbf{Acteur} &
\color{white}\textbf{Rôle} &
\color{white}\textbf{Objectif principal} &
\color{white}\textbf{Interactions clés} \\
\hline
\endfirsthead

\hline
\rowcolor{TableHeaderBlue}
\color{white}\textbf{Acteur} &
\color{white}\textbf{Rôle} &
\color{white}\textbf{Objectif principal} &
\color{white}\textbf{Interactions clés} \\
\hline
\endhead

\hline
\endfoot

Employé &
Utilisateur principal &
Évaluer ses compétences, consulter ses résultats et suivre ses recommandations de formation. &
Soumission du CV, ajout du profil GitHub, passation des tests, consultation du tableau de bord. \\
\hline

Administrateur RH &
Administrateur métier &
Superviser les profils, analyser les résultats et suivre les parcours de développement. &
Validation des compétences, affectation des formations, consultation des rapports et indicateurs globaux. \\
\hline

Système IA &
Acteur automatisé d’analyse &
Analyser les données hétérogènes et produire des résultats exploitables par la plateforme. &
Inférence locale, analyse sémantique, génération de scénarios et recommandations. \\
\hline

Communication Hub &
Acteur automatisé temps réel &
Assurer la communication instantanée entre les traitements backend et l’interface utilisateur. &
Émission d’événements SSE et notification de l’état d’avancement des analyses. \\
\hline

Service Mail &
Acteur automatisé asynchrone &
Envoyer les notifications importantes et assurer la traçabilité de certaines actions. &
Envoi d’emails de confirmation, de récupération de mot de passe, de certificats et d’alertes. \\
\hline
\caption{Acteurs du système et interactions principales}
\end{longtable}

\section{Méthodologie de développement Scrum}

Le projet a été conduit selon une approche Agile afin de tenir compte de la complexité fonctionnelle et technique de la solution. La méthodologie Scrum a été retenue car elle permet une progression incrémentale, une adaptation continue aux retours des parties prenantes et une meilleure maîtrise des risques liés à l’intégration de technologies émergentes.

La figuree~\ref{fig:organisation-scrum} illustre de manière synthétique l'organisation globale de la méthode Scrum adoptée pour la réalisation de ce projet TalentPredict AI. Elle met en évidence la répartition des rôles, les artefacts et cérémonies mobilisés, ainsi que le découpage fonctionnel des quatre sprints.
\begin{figure}[H]
    \centering
    \includegraphics[width=\textwidth]{images/Organisation Scrum du projet.png}
    \caption{Organisation Scrum du projet}
    \label{fig:organisation-scrum}
\end{figure}
\subsection{Choix et justification de Scrum}

Le choix de Scrum se justifie par la nécessité de produire des incréments fonctionnels réguliers, de valider progressivement les modules développés et de maintenir une communication continue entre les membres de l’équipe. Dans le cadre de TalentPredict AI, cette méthodologie est particulièrement adaptée, car le projet combine des fonctionnalités classiques de développement logiciel avec des composants plus expérimentaux liés à l’intelligence artificielle.

Scrum a également permis d’organiser le travail autour de sprints successifs, chacun ayant un objectif clair et des livrables vérifiables. Cette organisation réduit le risque d’accumulation d’erreurs en fin de projet et favorise une amélioration progressive de la qualité technique.

\subsection{Principes fondamentaux appliqués}

Les principes Scrum appliqués durant le projet reposent sur l’itération, la collaboration, la priorisation et l’inspection continue. Ces principes ont guidé la planification, le développement, les revues et les ajustements réalisés au fil des sprints.

\begin{itemize}[leftmargin=0.9cm]
    \item organisation du travail en itérations courtes ;
    \item priorisation des fonctionnalités selon leur valeur métier ;
    \item validation régulière avec le Product Owner ;
    \item adaptation progressive du backlog selon les retours obtenus ;
    \item intégration continue des développements réalisés ;
    \item identification proactive des risques techniques.
\end{itemize}

\subsection{Rôles et artefacts Scrum}

La mise en œuvre de Scrum nécessite une clarification des rôles et des responsabilités. Dans ce projet, les rôles Scrum ont été adaptés au contexte académique et professionnel du PFE, tout en respectant la logique générale du cadre méthodologique.

Le tableau~\ref{tab:roles_scrum} présente les rôles identifiés au sein de l’équipe projet.

\begin{longtable}{|P{0.24\textwidth}|P{0.24\textwidth}|P{0.42\textwidth}|}
\hline
\rowcolor{TableHeaderBlue}
\color{white}\textbf{Rôle Scrum} &
\color{white}\textbf{Acteur} &
\color{white}\textbf{Responsabilités principales} \\
\hline
\endfirsthead

\hline
\rowcolor{TableHeaderBlue}
\color{white}\textbf{Rôle Scrum} &
\color{white}\textbf{Acteur} &
\color{white}\textbf{Responsabilités principales} \\
\hline
\endhead

\hline
\endfoot

Product Owner &
M. Mohamed Amine Gharbi &
Définition de la vision produit, priorisation du Product Backlog, clarification des besoins et validation des livrables. \\
\hline

Équipe de développement &
Rahma Barouni et Mohamed Abdeladhim &
Conception, développement, test et intégration des modules fonctionnels, techniques et IA de la plateforme. \\
\hline

Scrum Master  &
Mme. Takwa Ben Aicha &
Suivi méthodologique, validation de la cohérence scientifique et accompagnement dans la structuration du rapport. \\
\hline
\caption{Rôles Scrum identifiés dans le projet}
\label{tab:roles_scrum}\\
\end{longtable}

Les artefacts Scrum utilisés structurent le suivi du projet et assurent la traçabilité des fonctionnalités. Le tableau~\ref{tab:artefacts_scrum} résume les principaux artefacts mobilisés.

\begin{longtable}{|P{0.26\textwidth}|P{0.46\textwidth}|P{0.18\textwidth}|}
\hline
\rowcolor{TableHeaderBlue}
\color{white}\textbf{Artefact} &
\color{white}\textbf{Description} &
\color{white}\textbf{Fréquence d’utilisation} \\
\hline
\endfirsthead

\hline
\rowcolor{TableHeaderBlue}
\color{white}\textbf{Artefact} &
\color{white}\textbf{Description} &
\color{white}\textbf{Fréquence d’utilisation} \\
\hline
\endhead

\hline
\endfoot

Product Backlog &
Liste priorisée des besoins, fonctionnalités, user stories et améliorations à réaliser. &
Mise à jour continue \\
\hline

Sprint Backlog &
Sous-ensemble des user stories sélectionnées pour un sprint donné, avec les tâches associées. &
À chaque sprint \\
\hline

Incrément &
Version fonctionnelle, testée et potentiellement livrable de la plateforme. &
Fin de sprint \\
\hline

Definition of Ready &
Ensemble de critères indiquant qu’une user story est suffisamment claire pour être développée. &
Avant planification \\
\hline

Definition of Done &
Ensemble de critères permettant de considérer une fonctionnalité comme terminée. &
Fin de tâche ou de sprint \\
\hline
\caption{Artefacts Scrum utilisés}
\label{tab:artefacts_scrum}\\
\end{longtable}

\subsection{Organisation des sprints}

Le projet est organisé en quatre sprints principaux, chacun correspondant à une étape cohérente du développement. Cette structuration permet de construire progressivement la plateforme, depuis les fondations techniques jusqu’aux fonctionnalités avancées d’analyse, de recommandation et de pilotage.

Le tableau~\ref{tab:recap_sprints_ch1} présente la planification globale des sprints.

\begin{longtable}{|P{0.15\textwidth}|P{0.15\textwidth}|P{0.30\textwidth}|P{0.30\textwidth}|}
\hline
\rowcolor{TableHeaderBlue}
\color{white}\textbf{Sprint} &
\color{white}\textbf{Période} &
\color{white}\textbf{Objectif principal} &
\color{white}\textbf{Livrables clés} \\
\hline
\endfirsthead

\hline
\rowcolor{TableHeaderBlue}
\color{white}\textbf{Sprint} &
\color{white}\textbf{Période} &
\color{white}\textbf{Objectif principal} &
\color{white}\textbf{Livrables clés} \\
\hline
\endhead

\hline
\endfoot

Sprint 1 &
Semaines 1 à 4 &
Mise en place de l’architecture, de la sécurité et de l’authentification. &
Configuration Docker~\cite{docker}, backend sécurisé, authentification JWT, premières interfaces Angular. \\
\hline

Sprint 2 &
Semaines 5 à 8 &
Développement des moteurs d’analyse des compétences. &
Workflows n8n~\cite{n8n}, extraction de texte, analyse GitHub~\cite{kalliamvakou2014}, parsing CV et premiers scores. \\
\hline

Sprint 3 &
Semaines 9 à 12 &
Intégration de l’IA, des tests dynamiques et du moteur de recommandation. &
Ollama~\cite{ollama}, Code Llama~\cite{roziere2023codellama}, QCM adaptatifs, challenges de code et recommandations. \\
\hline

Sprint 4 &
Semaines 13 à 16 &
Finalisation des tableaux de bord, tests et déploiement. &
Dashboard Angular~\cite{angular}, rapports PDF, tests d’intégration, déploiement final et validation globale. \\
\hline
\caption{Récapitulatif de la planification des sprints}
\label{tab:recap_sprints_ch1}\\
\end{longtable}

\section{Constitution de l’équipe projet}

La constitution de l’équipe projet permet d’identifier clairement les responsabilités de chaque intervenant. Cette clarification facilite la coordination des tâches, la gestion des dépendances et la validation progressive des livrables.

Le tableau~\ref{tab:equipe_projet_ch1} présente les membres de l’équipe et leurs responsabilités principales dans la réalisation de \textit{TalentPredict AI}.

\begin{longtable}{|P{0.28\textwidth}|P{0.24\textwidth}|P{0.38\textwidth}|}
\hline
\rowcolor{TableHeaderBlue}
\color{white}\textbf{Membre} &
\color{white}\textbf{Rôle} &
\color{white}\textbf{Responsabilités principales} \\
\hline
\endfirsthead

\hline
\rowcolor{TableHeaderBlue}
\color{white}\textbf{Membre} &
\color{white}\textbf{Rôle} &
\color{white}\textbf{Responsabilités principales} \\
\hline
\endhead

\hline
\endfoot

M. Mohamed Amine Gharbi &
Product Owner et encadrant professionnel &
Définition des besoins, priorisation des fonctionnalités, validation des sprints et suivi professionnel du projet. \\
\hline

Takwa Ben Aicha &
Scrum Master &
Garantie du respect du cadre méthodologique Agile Scrum, animation des cérémonies (Sprint Planning, Daily Stand-up, Sprint Review, Retrospective), facilitation de la communication et élimination des obstacles entravant l'équipe. \\
\hline

Rahma Barouni &
Développeuse Full-Stack &
Développement du module des compétences comportementales, workflows n8n, analyse PCM, intégration d’Ollama et traitements via FastAPI~\cite{fastapi}. \\
\hline

Mohamed Abdeladhim &
Développeur Full-Stack &
Développement du module des compétences techniques, parsing de CV, intégration GitHub et LinkedIn, QCM adaptatifs, évaluation de code et persistance des résultats. \\
\hline
\caption{Constitution de l’équipe projet}
\label{tab:equipe_projet_ch1}\\
\end{longtable}

\section{Solution proposée : \textit{TalentPredict AI}}

La solution proposée vise à répondre aux limites identifiées dans les processus classiques d’évaluation des compétences. Elle repose sur une approche hybride combinant des sources de données hétérogènes, des traitements d’intelligence artificielle et des mécanismes de restitution adaptés aux besoins des employés et des responsables RH.

\subsection{Description générale de la solution}

TalentPredict AI est une plateforme d’évaluation et de développement des compétences destinée à accompagner les entreprises dans l’analyse de leurs talents. Elle permet de collecter des données issues du CV, du profil GitHub, des tests techniques, des questionnaires comportementaux et des résultats d’évaluation, puis de les transformer en indicateurs exploitables.

La valeur ajoutée de la solution repose sur trois axes principaux. Le premier est l’analyse duale des compétences, qui combine les dimensions techniques et comportementales. Le deuxième est la validation multi-sources, qui permet de réduire la dépendance aux informations déclaratives. Le troisième est la souveraineté des données, rendue possible par l’utilisation de traitements locaux et d’une architecture contrôlée.

\subsection{Technologies de la plateforme}

La plateforme repose sur un ensemble de technologies complémentaires choisies pour leur maturité, leur compatibilité et leur adéquation avec les exigences du projet. Le tableau~\ref{tab:technologies_talentpredict_ch1} présente les principales technologies utilisées, ainsi que leur rôle dans l’architecture.

\begin{longtable}{|P{0.18\textwidth}|P{0.18\textwidth}|P{0.27\textwidth}|P{0.27\textwidth}|}
\hline
\rowcolor{TableHeaderBlue}
\color{white}\textbf{Technologie} &
\color{white}\textbf{Version indicative} &
\color{white}\textbf{Rôle dans le projet} &
\color{white}\textbf{Justification technique} \\
\hline
\endfirsthead

\hline
\rowcolor{TableHeaderBlue}
\color{white}\textbf{Technologie} &
\color{white}\textbf{Version indicative} &
\color{white}\textbf{Rôle dans le projet} &
\color{white}\textbf{Justification technique} \\
\hline
\endhead

\hline
\endfoot

Angular~\cite{angular} &
19+ &
Développement du frontend, des formulaires, des tableaux de bord et des visualisations. &
Architecture par composants, typage TypeScript, routage intégré et bonne intégration avec les interfaces interactives. \\
\hline

Spring Boot~\cite{springboot} / Java &
3.4.2 / 21 &
Backend métier, API REST, sécurité, persistance et génération de documents. &
Cadre robuste pour développer des services REST sécurisés, avec intégration native de Spring Security et JPA. \\
\hline

PostgreSQL~\cite{postgresql} &
15+ &
Stockage relationnel des utilisateurs, rôles, profils, compétences et résultats. &
SGBD fiable, conforme aux propriétés ACID et adapté à la gestion de données structurées et semi-structurées. \\
\hline

Python / FastAPI~\cite{fastapi} &
3.11 / 0.110+ &
Microservice IA pour l’analyse des CV, des profils, des QCM et des scénarios. &
Framework léger, asynchrone et adapté à la création rapide d’API pour les traitements IA. \\
\hline

n8n~\cite{n8n} &
latest &
Orchestration des workflows d’analyse et coordination des appels entre services. &
Permet de modéliser visuellement les flux, de modifier les workflows sans recompilation et de faciliter l’évolution des traitements. \\
\hline

Ollama~\cite{ollama} &
latest &
Exécution locale de modèles de langage pour l’analyse sémantique et la génération de recommandations. &
Permet une inférence locale, limitant l’externalisation des données sensibles et renforçant la souveraineté. \\
\hline

Code Llama~\cite{roziere2023codellama} &
latest &
Génération et évaluation de challenges de code. &
Modèle spécialisé dans le traitement du code, adapté aux tâches d’évaluation algorithmique. \\
\hline

Docker / Docker Compose~\cite{docker} &
24+ &
Conteneurisation et orchestration locale des services. &
Garantit la reproductibilité de l’environnement, l’isolation des composants et la portabilité du déploiement. \\
\hline

pdfjs-dist~\cite{pdfjs} &
5.6.205 &
Extraction de texte à partir de fichiers PDF côté client. &
Réduit l’exposition des CV bruts côté serveur et contribue à une logique de Privacy by Design. \\
\hline

Chart.js~\cite{chartjs} &
4+ &
Visualisation des scores, courbes d’évolution et radars de compétences. &
Bibliothèque légère, adaptée aux tableaux de bord et facilement intégrable avec Angular. \\
\hline
\caption{Technologies principales de TalentPredict AI}
\label{tab:technologies_talentpredict_ch1}\\
\end{longtable}

\subsection{Architecture générale de la solution}

L’architecture de TalentPredict AI est organisée autour de plusieurs composants spécialisés. Cette séparation des responsabilités permet de mieux maîtriser la complexité, de faciliter la maintenance et de préparer l’évolution future de la plateforme.

La figure présente l’architecture globale de la solution. Elle montre les interactions entre l’interface utilisateur, le backend métier, les services IA, l’orchestrateur de workflows, la base de données et les services de communication.
\begin{figure}[H]
    \centering
    \includegraphics[width=\textwidth]{images/archi Glob.png}
    \caption{Architecture globale de la plateforme TalentPredict AI}
    \label{fig:architecture-glob}
\end{figure}
\begin{figure}[H]
    \centering
    \includegraphics[width=\textwidth]{images/Architecture en couches2.png}
    \caption{Architecture globale de la plateforme TalentPredict AI}
    \label{fig:architecture-glob2}
\end{figure}

\subsection{Vue d’ensemble en couches}

L’architecture retenue adopte un modèle en couches afin de séparer la présentation, la logique métier, l’analyse intelligente, la persistance et le déploiement. Cette organisation améliore la lisibilité du système et permet de faire évoluer chaque couche de manière relativement indépendante.

La figure suivante présente une vue synthétique de l’architecture en couches. Elle permet de comprendre le rôle de chaque niveau et la circulation générale des données au sein de la plateforme.

\chapterfigure[\textwidth]{images/archi_couche.png}{Vue d’ensemble de l’architecture en couches de TalentPredict AI}

\section{Product Backlog}

Le Product Backlog regroupe les principales fonctionnalités attendues dans la première version de la plateforme. Il constitue une base de priorisation pour l’équipe Scrum et permet de transformer les besoins métiers en user stories exploitables dans les sprints.

Le tableau~\ref{tab:product_backlog_ch1} présente un extrait du Product Backlog initial, centré sur les fonctionnalités prioritaires du MVP.

\begin{longtable}{|P{0.10\textwidth}|P{0.15\textwidth}|P{0.15\textwidth}|P{0.18\textwidth}|P{0.22\textwidth}|P{0.10\textwidth}|}
\hline
\rowcolor{TableHeaderBlue}
\color{white}\textbf{ID} &
\color{white}\textbf{Epic} &
\color{white}\textbf{En tant que} &
\color{white}\textbf{Je souhaite} &
\color{white}\textbf{Afin de} &
\color{white}\textbf{Priorité} \\
\hline
\endfirsthead

\hline
\rowcolor{TableHeaderBlue}
\color{white}\textbf{ID} &
\color{white}\textbf{Epic} &
\color{white}\textbf{En tant que} &
\color{white}\textbf{Je souhaite} &
\color{white}\textbf{Afin de} &
\color{white}\textbf{Priorité} \\
\hline
\endhead

\hline
\endfoot

US-001 &
Authentification et sécurité &
Utilisateur / RH &
M’authentifier via un jeton JWT &
Accéder à l’application et à mes données de manière sécurisée &
Haute \\
\hline

US-002 &
Compétences comportementales &
Utilisateur &
Passer le test de personnalité PCM~\cite{kahler1996} &
Générer mon profil comportemental et identifier mon type dominant &
Haute \\
\hline

US-003 &
Compétences comportementales et IA &
Système IA &
Analyser sémantiquement les profils avec un modèle local &
Déduire des indicateurs comportementaux tout en préservant la confidentialité des données &
Haute \\
\hline

US-004 &
Analyse technique &
Utilisateur &
Transmettre mon CV et mon pseudonyme GitHub &
Permettre l’extraction automatique des données techniques &
Haute \\
\hline

US-005 &
Analyse technique &
Utilisateur &
Fournir l’URL de mon portfolio ou profil LinkedIn &
Enrichir l’évaluation par des signaux professionnels supplémentaires &
Haute \\
\hline

US-006 &
Scoring multi-sources &
Système d’évaluation &
Fusionner les résultats issus des différentes sources &
Attribuer un niveau de maîtrise et un score de confiance &
Haute \\
\hline

US-007 &
Recommandations &
Système de recommandation &
Croiser les compétences détectées avec les besoins du poste &
Proposer des parcours de formation pertinents &
Haute \\
\hline

US-008 &
Orchestration &
Développeur &
Modéliser les workflows d’analyse dans n8n &
Coordonner les appels entre frontend, backend et services IA &
Haute \\
\hline

US-009 &
Communication &
Administrateur RH &
Notifier les collaborateurs lors des étapes importantes &
Assurer un suivi opérationnel centralisé &
Haute \\
\hline
\caption{Product Backlog initial du projet}
\label{tab:product_backlog_ch1}\\
\end{longtable}

\section{Critères d’acceptation des user stories}

Les critères d’acceptation permettent de vérifier qu’une user story est suffisamment claire avant son développement et qu’elle peut être considérée comme terminée après réalisation. Ils contribuent à réduire les ambiguïtés, à améliorer la qualité des livrables et à faciliter la validation par le Product Owner.

\subsection{Definition of Ready}

Une user story est considérée comme prête lorsqu’elle peut être comprise, estimée et développée par l’équipe sans ambiguïté majeure. Dans ce projet, les critères de préparation suivants ont été appliqués :

\begin{itemize}[leftmargin=0.9cm]
    \item la user story est formulée selon le format « En tant que, je souhaite, afin de » ;
    \item les critères d’acceptation sont clairement définis ;
    \item les dépendances techniques sont identifiées ;
    \item les maquettes ou descriptions fonctionnelles nécessaires sont disponibles ;
    \item la priorité est définie dans le Product Backlog ;
    \item la story est validée par le Product Owner.
\end{itemize}

\subsection{Definition of Done}

Une fonctionnalité est considérée comme terminée lorsqu’elle respecte les exigences fonctionnelles, techniques et qualitatives définies par l’équipe. Ces critères assurent que l’incrément livré est stable, testable et potentiellement déployable.
\begin{itemize}[leftmargin=0.9cm]
    \item le code est implémenté, relu et intégré dans la branche principale ;
    \item les tests unitaires ou fonctionnels associés sont validés ;
    \item les API concernées sont documentées ;
    \item l’interface ne présente pas d’erreurs bloquantes ;
    \item la fonctionnalité respecte les règles de sécurité définies ;
    \item le livrable est validé par le Product Owner lors de la revue de sprint.
\end{itemize}
\section{Planification générale}

La planification générale permet de visualiser l’organisation globale du projet et les interactions attendues entre les acteurs et le système. Elle regroupe les représentations fonctionnelles et temporelles nécessaires à la compréhension du déroulement du projet.
\subsection{Diagramme des cas d’utilisation global}

Le diagramme des cas d’utilisation~\ref{fig:D-use-glob} offre une vue fonctionnelle d’ensemble de la plateforme. Il permet d’identifier les acteurs principaux, les services externes et les fonctionnalités majeures auxquelles chacun participe.

La figure~\ref{fig:D-use-glob} présente le diagramme global des cas d’utilisation de TalentPredict AI. Elle illustre les interactions entre l’employé, l’administrateur RH, les services IA, le service mail, le service SSE et les principales fonctionnalités de la plateforme.
\begin{figure}[H]
    \centering
    \includegraphics[width=\textwidth]{images/D use case glob platfome.png}
    \caption{Diagramme global des cas d’utilisation de la plateforme TalentPredict AI}
    \label{fig:D-use-glob}
\end{figure}

À noter que ce diagramme~\ref{fig:D-use-glob} se concentre uniquement sur les fonctionnalités de haut niveau. Les détails techniques internes, tels que les workflows d'intelligence artificielle ou les appels d'API spécifiques, ne sont pas représentés ici et seront détaillés dans les diagrammes d'architecture et de séquence.

\color{red}
\begin{itemize}[noitemsep, topsep=2pt]
    \item \textit{Diagramme 1 :} Remplacer le titre de la frontière du système par \textit{TalentPredict AI} afin de garder le même nom que dans le rapport.
    \item \textit{Diagramme 1 :} Supprimer les flèches sur les associations entre les acteurs et les cas d’utilisation ; en UML, une association acteur--cas d’utilisation est une ligne simple sans flèche.
    \item \textit{Diagramme 1 :} Renommer l’acteur \textit{Administrateur / RH} en \textit{Administrateur RH} pour harmoniser le vocabulaire avec le reste du rapport.
    \item \textit{Diagramme 1 :} Remplacer le cas \textit{Recevoir des recommandations de formations} par \textit{Consulter les recommandations de formation}, car l’acteur ne reçoit pas directement une fonctionnalité mais consulte un résultat généré par la plateforme.
    \item \textit{Diagramme 1 :} Remplacer le cas \textit{Générer une feuille de route} par \textit{Consulter la feuille de route d’upskilling} si l’acteur est l’employé ; la génération doit être réalisée par le système ou le service IA.
    \item \textit{Diagramme 1 :} Ajouter le cas d’utilisation \textit{Créer un compte}, car il fait partie du parcours d’accès initial de l’employé.
    \item \textit{Diagramme 1 :} Ajouter le cas d’utilisation \textit{Gérer son profil} pour l’employé et l’administrateur RH.
    \item \textit{Diagramme 1 :} Ajouter une relation \texttt{<<include>>} entre \textit{Télécharger le passeport de compétences} et \textit{Générer le passeport de compétences}.
    \item \textit{Diagramme 1 :} Vérifier que la relation \texttt{<<include>>} part du cas principal \textit{Réaliser une évaluation multi-source} vers les cas inclus \textit{Effectuer les tests} et \textit{Soumettre données}.
    \item \textit{Diagramme 1 :} Reformuler \textit{Soumettre données (CV, GitHub, LinkedIn)} en plusieurs cas plus clairs : \textit{Soumettre son CV}, \textit{Renseigner son profil GitHub} et \textit{Renseigner son profil LinkedIn}.
    \item \textit{Diagramme 1 :} Relier l’acteur \textit{API GitHub} uniquement au cas \textit{Analyser le profil GitHub} ou \textit{Renseigner son profil GitHub}, et non à un cas général de soumission de données.
    \item \textit{Diagramme 1 :} Relier l’acteur \textit{Système IA} uniquement aux cas d’analyse ou de génération, par exemple \textit{Évaluer les compétences}, \textit{Générer les recommandations} et \textit{Générer la feuille de route}.
    \item \textit{Diagramme 1 :} Remplacer \textit{Recevoir les notifications asynchrones} par \textit{Consulter les notifications}, car le cas d’utilisation doit représenter une action observable de l’utilisateur.
    \item \textit{Diagramme 1 :} Relier \textit{Service Mail} au cas \textit{Récupérer son mot de passe} à travers un cas inclus \textit{Envoyer un lien de réinitialisation}.
    \item \textit{Diagramme 1 :} Relier \textit{Service SSE} uniquement au cas \textit{Diffuser les notifications en temps réel} ou \textit{Mettre à jour les notifications}.
    \item \textit{Diagramme 1 :} Séparer le cas \textit{Consulter les résultats et le Job Matching} en deux cas distincts : \textit{Consulter les résultats d’évaluation} et \textit{Consulter le Job Matching}.
    \item \textit{Diagramme 1 :} Ajouter une légende expliquant les acteurs, les cas d’utilisation, les relations \texttt{<<include>>} et les associations simples.
    \item \textit{Diagramme 1 :} Réorganiser les cas d’utilisation pour réduire les croisements de lignes, en plaçant les cas de l’employé à gauche, les cas de l’administrateur RH en bas ou à droite, et les services externes autour de la frontière du système.
    \item \textit{Diagramme 1 :} Ajouter une phrase avant la figure pour expliquer que le diagramme présente les interactions fonctionnelles globales entre les acteurs et la plateforme.
    \item \textit{Diagramme 1 :} Ajouter une phrase après la figure pour préciser que les détails techniques internes, comme les workflows IA ou les appels API, seront détaillés dans les diagrammes d’architecture ou de séquence.
\end{itemize}
\color{black}
\subsection{Diagramme de Gantt}

Le diagramme de Gantt permet de représenter la planification temporelle du projet. Il met en évidence l’enchaînement des sprints, les grandes phases de développement et la progression des livrables dans le temps.

La figure~\ref{fig:gantt-projet-ch1} l’emplacement prévu pour le diagramme de Gantt du projet. Elle devra représenter les quatre sprints principaux, leurs périodes respectives et les livrables associés.
\begin{figure}[H]
    \centering
    \includegraphics[width=\textwidth]{images/D gantt.png}
    \caption{Diagramme de Gantt du projet}
    \label{fig:gantt-projet-ch1}
\end{figure}
%\figureplaceholder{Diagramme de Gantt du projet}{fig:gantt-projet-ch1}

\section*{Conclusion}
\phantomsection
\addcontentsline{toc}{section}{Conclusion}

Ce chapitre a permis de poser le cadre général du projet \textit{TalentPredict AI}. Il a présenté l’organisme d’accueil, la problématique métier, les objectifs, les besoins, les acteurs, la méthodologie Scrum et la planification générale. L’analyse menée montre que le projet répond à un besoin réel d’évaluation structurée, sécurisée et évolutive des compétences en entreprise. Les choix méthodologiques et architecturaux introduits dans ce chapitre constituent ainsi une base cohérente pour les chapitres suivants, qui détailleront l’état de l’art puis la réalisation progressive de la plateforme à travers les différents sprints.


\chapter{État de l'Art}

\textbf{Introduction}

Ce chapitre présente une revue systématique des travaux existants et des technologies fondatrices sur lesquels repose la plateforme TalentPredict AI. Le domaine de l'évaluation intelligente des compétences en entreprise se situe à l'intersection de plusieurs champs disciplinaires en évolution rapide : la gestion des ressources humaines augmentée par l'intelligence artificielle (HR Tech), le traitement automatique du langage naturel (NLP), les modèles psychométriques computationnels et les architectures logicielles distribuées orientées IA. L'objectif de ce chapitre n'est pas de dresser un inventaire exhaustif, mais d'analyser avec esprit critique les approches et solutions existantes afin d'identifier leurs limites structurelles et de justifier les choix conceptuels et techniques de TalentPredict AI.

Le chapitre s'organise en cinq sections. La première délimite le domaine et ses enjeux fondamentaux. La deuxième présente les concepts et technologies clés. La troisième analyse les travaux et solutions similaires. La quatrième propose une synthèse comparative. La cinquième positionne TalentPredict par rapport aux lacunes identifiées.

\section{Présentation du domaine}
L'optimisation du capital humain requiert aujourd'hui des outils capables d'évaluer le profil global des collaborateurs de manière objective et continue. Pour bien comprendre le positionnement de notre plateforme, nous explorerons dans cette section le contexte stratégique de la gestion des compétences internes, les dimensions complémentaires qui les composent, ainsi que les principaux verrous qu'un système d'évaluation automatisé par IA se doit de lever.
\subsection{La gestion des compétences internes : un enjeu stratégique majeur}

Dans une économie de la connaissance caractérisée par l'obsolescence rapide des compétences techniques et la raréfaction des profils spécialisés, les organisations ne peuvent plus se contenter d'une gestion réactive des talents fondée sur l'acquisition de ressources externes. La tendance de fond, documentée depuis les travaux précurseurs de Ulrich (1997) sur le HR Business Partner~\cite{ch2-ulrich1997}, est celle d'une gestion proactive et continue des compétences internes, désignée sous les termes d'Upskilling (montée en compétences sur le poste actuel) et de Reskilling (reconversion vers un nouveau métier).

Cette transition impose aux départements RH de disposer d'outils capables d'évaluer objectivement et en continu le niveau de compétence réel de chaque collaborateur  et non plus seulement leur niveau déclaré  afin d'identifier les écarts (skill gaps) entre les capacités actuelles et les exigences stratégiques de l'organisation. C'est précisément cette problématique que TalentPredict AI cherche à résoudre à travers un pipeline d'évaluation automatisé, multi-sources et piloté par l'intelligence artificielle.

\subsection{Les deux dimensions complémentaires de la compétence professionnelle}

La littérature en psychologie du travail et des organisations distingue classiquement deux grandes catégories de compétences, dont l'évaluation requiert des approches méthodologiques fondamentalement différentes.

Les compétences techniques (hard skills) désignent les savoirs opérationnels vérifiables et mesurables : maîtrise d'un langage de programmation, d'un framework, d'une plateforme cloud ou d'une méthodologie de développement logiciel. Leur évaluation peut s'appuyer sur des preuves objectives  résultats de tests techniques, analyse de code produit, historique de contributions réduisant ainsi la part de subjectivité.

Les compétences comportementales (soft skills) désignent les aptitudes transverses qui déterminent la manière dont un individu mobilise ses compétences techniques en contexte professionnel : capacité de communication, leadership, collaboration, gestion du stress, adaptabilité. Leur évaluation est intrinsèquement plus complexe car elle requiert l'observation de comportements en situation, difficilement capturables par un simple questionnaire auto-déclaratif~\cite{ch2-salgado1997}.

L'originalité de TalentPredict réside précisément dans l'intégration de ces deux dimensions au sein d'un même pipeline d'évaluation, permettant de produire un profil collaborateur à 360\textdegree{} combinant preuves objectives et analyse comportementale.

\subsection{Les défis de l'évaluation automatisée des compétences}

Trois défis fondamentaux caractérisent l'automatisation de l'évaluation des compétences en entreprise et constituent le fil conducteur de cette revue de l'état de l'art.

Le premier défi est celui du biais déclaratif : la quasi-totalité des systèmes d'évaluation existants reposent, en tout ou partie, sur ce que le collaborateur affirme de lui-même. Or, la psychologie sociale a établi de longue date que les auto-évaluations sont systématiquement biaisées par des effets de désirabilité sociale~\cite{ch2-donaldson2002} et de surestimation des compétences propres (Dunning-Kruger effect~\cite{ch2-kruger1999}).

Le second défi est celui de la fragmentation technique,comportementale : les solutions actuelles traitent séparément l'évaluation technique et l'évaluation comportementale, sans mécanisme de croisement permettant, par exemple, de vérifier la cohérence entre un profil psychométrique déclaré et des comportements observables dans des données asynchrones (historique de contributions, dynamiques de collaboration en équipe).

Le troisième défi est celui de la souveraineté des données : l'exploitation des grands modèles de langage (LLM) pour l'analyse sémantique des profils RH implique le traitement de données à caractère sensible (données psychologiques, données professionnelles personnelles), ce qui soulève des problématiques de conformité réglementaire, en particulier au regard du Règlement Général sur la Protection des Données (RGPD)~\cite{ch2-rgpd2016}.

\section{Concepts et technologies clés}
Pour répondre aux défis de la gestion des talents, notre plateforme s'appuie sur une convergence entre psychométrie éprouvée et intelligence artificielle de pointe. Cette section dresse un panorama des concepts clés du projet : elle débute par une analyse des méthodes d'évaluation traditionnelles pour ensuite présenter les ruptures technologiques (analyse sémantique, Edge AI, microservices) permettant d'automatiser et d'objectiver ce processus.
\subsection{Méthodes traditionnelles d'évaluation des compétences}

\textbf{L'évaluation managériale structurée:} L'échange d'évaluation (ou bilan de compétences face-à-face) constitue, depuis des décennies, l'outil central de la gestion des compétences en entreprise. Dans sa forme structurée  où les questions sont standardisées et les réponses cotées selon une grille prédéfinie  il présente une validité prédictive de la performance professionnelle significativement supérieure aux évaluations non structurées. McDaniel et al. (1994), dans une méta-analyse portant sur 114 études, ont établi une validité prédictive (r corrigé) de 0,51 pour les mises en situation structurées, contre 0,38 pour les échanges non structurés~\cite{ch2-mcdaniel1994}. Toutefois, ce format d'évaluation reste coûteux en temps RH, difficile à standardiser à grande échelle et sujet au biais de l'évaluateur (halo effect, biais de similarité).

\textbf{Les tests psychométriques standardisés:} Les outils psychométriques constituent le second pilier de l'évaluation comportementale traditionnelle. Parmi les modèles les plus utilisés en contexte organisationnel, on distingue :
\begin{itemize}[leftmargin=1.5cm]
  \item Le modèle des Big Five (OCEAN : Ouverture, Conscienciosité, Extraversion, Agréabilité, Névrosisme), issu de la psychologie différentielle et validé par plusieurs décennies de recherche empirique~\cite{ch2-costa1992}. Sa solidité scientifique est reconnue, mais son opérationnalisation en recommandations de formation concrètes reste limitée.
  \item Le modèle DISC (Dominance, Influence, Stabilité, Conformité), largement utilisé en développement organisationnel mais dont la validité scientifique est contestée ; plusieurs études ont mis en évidence une faible fidélité test-retest et une validité prédictive modeste~\cite{ch2-pittenger1993}.
  \item Le Process Communication Model (PCM), développé par le Dr Taibi Kahler dans les années 1970 et validé cliniquement dans le cadre du programme de sélection des astronautes de la NASA~\cite{ch2-kahler1996}. Le PCM structure la personnalité en six types fondamentaux (Empathique, Travaillomane, Persévérant, Rêveur, Rebelle, Promoteur), chacun associé à des canaux de communication préférentiels, des besoins psychologiques spécifiques et des comportements prévisibles sous stress. Sa richesse opérationnelle --- notamment sa dimension prescriptive sur la modalité de communication optimale --- en fait un cadre particulièrement adapté à la personnalisation des recommandations de formation. C'est ce modèle qui a été retenu dans TalentPredict AI pour piloter le module d'évaluation comportementale.
\end{itemize}

\textbf{Les centres d'évaluation (Assessment Centers):} Les assessment centers sont des dispositifs multi-méthodes combinant simulations de situations professionnelles, mises en situation collectives et tests psychométriques, évalués par plusieurs assesseurs formés. Ils présentent la validité prédictive la plus élevée parmi les outils d'évaluation des compétences managériales (r = 0,36 à 0,45 selon les méta-analyses)~\cite{ch2-thornton2006}. Toutefois, leur coût élevé (plusieurs milliers d'euros par collaborateur), leur durée (une à deux journées) et la difficulté de les déployer à l'échelle d'une organisation entière en font un outil réservé aux populations managériales, inapplicable à une évaluation systématique des compétences de l'ensemble du personnel.

\subsection{Traitement automatique du langage naturel appliqué aux profils professionnels}

Le NLP constitue la brique technologique centrale des systèmes modernes d'évaluation automatisée des compétences. Son application à l'analyse des documents professionnels (CV internes, dossiers de compétences, profils professionnels) a connu une évolution majeure avec l'avènement des architectures Transformer.

\textbf{Approches par règles et expressions régulières:} Les premières générations de parseurs de CV reposaient sur des règles heuristiques et des expressions régulières (regex) pour identifier les entités nommées (technologies, formations, expériences). Ces approches présentent l'avantage d'être déterministes, rapides et transparentes, mais souffrent de limites structurelles importantes : elles sont sensibles aux variations de formulation, ne capturent pas le contexte sémantique et nécessitent une maintenance manuelle intensive à mesure que le vocabulaire technique évolue. Les études de performance publiées font état de scores F1 typiquement compris entre 0,65 et 0,75 pour l'extraction de technologies depuis des CV~\cite{ch2-celik2013}.

\textbf{Approches par modèles Transformer (BERT, LLaMA)} L'avènement des modèles de langage pré-entraînés a profondément transformé les performances des systèmes d'extraction d'information. Chien \& Chia-Hui (2020) ont démontré que les architectures de type Transformer surpassent significativement les approches basées sur les champs aléatoires conditionnels (CRF) pour les tâches de reconnaissance d'entités nommées (NER) appliquées aux CV, avec un gain de performance de 12 à 18 points de F1~\cite{ch2-chien2020}. Des benchmarks plus récents sur des modèles de la famille LLaMA confirment des scores F1 supérieurs à 0,85 pour l'extraction de compétences techniques depuis des documents textuels non structurés, issues de la littérature~\cite{ch2-shi2023}. Dans TalentPredict AI, cette approche est mise en œuvre via le module FastAPI/Ollama qui complète l'extraction par règles d'un second niveau d'analyse sémantique LLM, permettant de capturer les technologies implicitement mentionnées ou formulées de manière non standard.

\textbf{Extraction d'indicateurs comportementaux depuis les données textuelles:} Au-delà de l'extraction de compétences techniques, le NLP permet d'inférer des signaux comportementaux depuis le texte d'un CV ou d'un profil professionnel : détection de termes évoquant le travail en équipe, le leadership ou l'adaptabilité, analyse de la progression de carrière, cohérence du parcours professionnel. Ces approches restent toutefois limitées par le caractère auto-déclaratif des sources analysées, qui reproduit le biais documenté en section 2.1.3.

La figure~\ref{fig:NLP-profile} synthétise l'application concrète de ces approches d'extraction au sein de TalentPredict AI. Elle illustre le pipeline NLP complet, détaillant le cheminement d'un document brut (CV ou profil professionnel) à travers les phases d'extraction, de nettoyage sémantique, de détection par modèles d'intelligence artificielle et de normalisation, jusqu'à la production d'un score quantifié.
\begin{figure}[H]
    \centering
    \includegraphics[width=\textwidth]{images/Pipeline NLP profile.png}
    \caption{Architecture en 6 étapes du pipeline NLP appliqué à l'extraction, la normalisation et le scoring des compétences professionnelles.}
    \label{fig:NLP-profile}
\end{figure}

\subsection{Analyse des dépôts de code source comme proxy de compétences}

L'exploitation des plateformes de versioning (GitHub, GitLab) comme source de données pour l'évaluation des compétences techniques constitue un paradigme émergent dans la recherche en ingénierie logicielle. Kalliamvakou et al. (2014) ont posé les jalons de cette approche en démontrant que les métriques d'activité GitHub (nombre de dépôts, fréquence de contribution, diversité des langages utilisés, participation aux pull requests) constituent des proxies valides  mais imparfaits  de l'expertise technique réelle d'un développeur~\cite{ch2-kalliamvakou2014}.

Les travaux de Marlow et al. (2013) ont par ailleurs établi que certaines dynamiques de collaboration observables dans les données GitHub fréquence des revues de code, vitesse de réponse aux issues, diversité des collaborateurs  présentent une corrélation modérée mais statistiquement significative avec des indicateurs de performance professionnelle évalués par les managers~\cite{ch2-marlow2013}. Ces résultats suggèrent que les données GitHub peuvent servir de proxy partiel pour certaines dimensions comportementales (ownership, collaboration, consistance) en complément des outils psychométriques classiques, ce qui est précisément le principe mis en œuvre dans le workflow d'analyse GitHub de TalentPredict.

La figure~\ref{fig:exploi-github} modélise ce workflow d'analyse spécifique à notre plateforme. Elle illustre concrètement comment les différentes métriques extraites via l'API GitHub (diversité des langages, volume de commits, régularité via la heat-map, et dynamiques de collaboration) sont transformées en indicateurs pondérés pour générer un score technique global et objectif.

\begin{figure}[H]
    \centering
    \includegraphics[width=\textwidth]{images/exploi github data en preuve tech.png}
    \caption{Exploitation des données
GitHub comme preuve
techniqu}
    \label{fig:exploi-github}
\end{figure}

Néanmoins, plusieurs limites à cette approche ont été identifiées dans la littérature : le biais en faveur des développeurs contribuant à des projets open source publics (défavorable aux collaborateurs travaillant exclusivement sur des dépôts privés en entreprise), la difficulté à distinguer quantité et qualité des contributions, et la variabilité des pratiques de versioning entre organisations~\cite{ch2-kalliamvakou2014}.

\subsection{Grands modèles de langage : architectures et performances comparées}
Le traitement sémantique avancé requis par TalentPredict s'appuie sur l'utilisation de Grands Modèles de Langage (LLM). Cette section dresse un état des lieux des architectures existantes, propose une analyse comparative de plusieurs modèles candidats, et justifie les choix technologiques retenus pour notre moteur d'inférence.
\textbf{Panorama des architectures LLM} Depuis la publication de l'architecture Transformer par Vaswani et al. (2017)~\cite{ch2-vaswani2017}, le champ des modèles de langage de grande taille a connu une évolution exponentielle. Deux grandes familles se distinguent pour les applications en entreprise :
\begin{itemize}[leftmargin=1.5cm]
  \item Les modèles propriétaires hébergés dans le cloud (GPT-4 d'OpenAI, Claude d'Anthropic, Gemini de Google) offrent des performances de pointe sur les tâches de compréhension sémantique complexe et de génération structurée, mais impliquent une dépendance à un fournisseur tiers, des coûts d'utilisation à l'appel potentiellement élevés, et surtout une externalisation du traitement des données vers des infrastructures non maîtrisées par l'organisation.
  \item Les modèles open-source déployables localement (famille LLaMA de Meta, Mistral AI, Phi-3 de Microsoft) permettent un déploiement on-premise ou dans une infrastructure privée, garantissant ainsi la souveraineté des données. Leurs performances, initialement inférieures aux modèles propriétaires, ont considérablement progressé avec les récentes générations.
\end{itemize}

\textbf{Comparaison quantitative et qualitative des modèles open-source candidats} Plusieurs modèles ont été évalués comme candidats potentiels pour le moteur d'inférence de TalentPredict AI. Le tableau suivant synthétise leurs caractéristiques selon les critères pertinents pour ce projet :

\begin{table}[H]
\centering
\small
\renewcommand{\arraystretch}{1.35}
\begin{tabularx}{\textwidth}{|>{\RaggedRight\arraybackslash}p{0.22\textwidth}|Y|Y|Y|Y|}
\hline
\textbf{Critère} & \textbf{LLaMA 3.2 (3B/8B)} & \textbf{Mistral 7B} & \textbf{Phi-3 Mini (3.8B)} & \textbf{GPT-4 (cloud)} \\
\hline
Taille (paramètres) & 3B / 8B & 7B & 3.8B & \textasciitilde{}1T (estimé) \\
\hline
Déploiement local & Oui, via Ollama & Oui, via Ollama & Oui, via Ollama & Non, API cloud \\
\hline
Performance NLP générale (MMLU) & 63\% / 73\% & 64\% & 69\% & 86\% \\
\hline
Qualité génération JSON structuré & Bonne & Bonne & Moyenne & Excellente \\
\hline
Latence CPU (génération 200 tokens) & \textasciitilde{}8s / \textasciitilde{}18s & \textasciitilde{}15s & \textasciitilde{}6s & \textasciitilde{}3s (réseau) \\
\hline
Coût par requête & 0 (local) & 0 (local) & 0 (local) & \textasciitilde{}0,03€ \\
\hline
Conformité RGPD & Totale & Totale & Totale & Limitée : données cloud \\
\hline
Support CodeLlama (code) & Dédié & Partiel & Limité & Non disponible \\
\hline
Licence & Meta LLaMA & Apache 2.0 & MIT & Propriétaire \\
\hline
\end{tabularx}
\renewcommand{\arraystretch}{1.05}
\caption{Comparaison des modèles LLM candidats pour TalentPredict AI}
\end{table}

Sources : benchmarks MMLU publiés par les éditeurs ; mesures de latence estimées d'après la littérature sur du matériel standard (CPU 8 cœurs, 16 Go RAM)~\cite{ch2-meta2024,ch2-mistral2023,ch2-abdin2024}.

\textbf{Justification du choix de llama3.2 et CodeLlama} Au terme de cette analyse comparative, le modèle llama3.2 (variante 3B pour les analyses légères, 8B pour les synthèses comportementales complexes) a été retenu comme moteur d'inférence principal de TalentPredict AI, pour les raisons suivantes. Touvron et al. (2023) ont démontré que les architectures LLaMA offrent un rapport performance/taille remarquable pour les tâches de NLP standard, avec des performances compétitives par rapport aux modèles hébergés dans le cloud sur les tâches de classification et d'extraction sémantique~\cite{ch2-touvron2023}. Des évaluations ultérieures (Xu et al., 2024) ont confirmé la pertinence de llama3.2 quantifié en Q4 pour les tâches NLP appliquées aux ressources humaines, avec des latences acceptables sur des processeurs standards~\cite{ch2-xu2024}.

Le modèle CodeLlama, dérivé de LLaMA et spécialisé pour la compréhension et la génération de code, a été retenu pour le module de challenge de code en raison de ses performances supérieures aux modèles généralistes sur les tâches d'évaluation algorithmique. Mistral 7B a été écarté en raison d'une latence CPU plus élevée sans gain de performance significatif sur les tâches ciblées. Phi-3 Mini, bien que plus rapide, présentait une qualité de génération JSON structuré insuffisante pour les workflows d'orchestration.

\subsection{Edge AI et souveraineté des données}

Le concept d'Edge AI désigne le déploiement de modèles d'intelligence artificielle directement sur l'infrastructure locale de l'organisation, sans recours à des services cloud externes. Dans le contexte de l'évaluation des compétences RH, cette approche revêt une importance particulière au regard du cadre réglementaire européen.

L'Article 9 du RGPD classe les données psychologiques dans la catégorie des «~données à caractère sensible~», dont le traitement est soumis à des exigences renforcées~\cite{ch2-rgpd2016}. Les analyses de personnalité issues du modèle PCM, les scores comportementaux inférés depuis les profils professionnels et les résultats de tests psychométriques entrent directement dans cette catégorie. Leur transmission vers des infrastructures cloud tierces (serveurs hors UE, notamment pour les modèles américains) soulève des risques de non-conformité réglementaire que les DPO (Data Protection Officers) des grandes organisations ne peuvent pas accepter.

L'utilisation d'Ollama comme serveur d'inférence local dans TalentPredict AI résout structurellement ce problème : l'intégralité du traitement sémantique s'effectue sur l'infrastructure interne, les données ne quittent jamais le réseau privé Docker de la plateforme, et la conformité au principe de Privacy by Design (Article 25 du RGPD) est garantie par architecture.

La figure~\ref{fig:inference-local} illustre schématiquement cette rupture architecturale. Elle met en opposition l'approche cloud traditionnelle écartée en raison de ses risques inhérents de non conformité  avec l'approche d'inférence locale (via Ollama) retenue pour TalentPredict AI, garantissant ainsi une souveraineté et une confidentialité totales des données par conception.
\begin{figure}[H]
    \centering
    \includegraphics[width=\textwidth]{images/inference local ollama.png}
    \caption{Inférence locale avec Ol-
lama}
    \label{fig:inference-local}
\end{figure}

\subsection{Orchestration de workflows et architectures microservices}
L'intégration efficace de modèles d'intelligence artificielle au sein d'une application d'entreprise complexe exige une architecture logicielle robuste et modulaire. Cette section explore les principes d'orchestration et de conception en microservices qui sous-tendent l'infrastructure technique de TalentPredict, garantissant ainsi la flexibilité et la scalabilité de ses traitements.
\textbf{L'orchestration low-code:} Les plateformes d'orchestration visuelles de workflows (n8n, Zapier, Make, Apache Airflow) permettent de modéliser des séquences complexes d'appels à des services hétérogènes sans code impératif. van der Aalst (2013) a théorisé les bénéfices du découplage de la logique d'orchestration hors du noyau applicatif : meilleure maintenabilité, traçabilité des flux, et capacité à modifier les comportements (notamment les prompts LLM) sans recompilation du code source~\cite{ch2-vanderaalst2013}. Dans TalentPredict AI, n8n a été retenu pour orchestrer les workflows d'analyse comportementale (GitHub Analysis, CV Parser, PCM Test, Master Agent), permettant aux équipes RH ou IA de modifier les prompts directement depuis l'interface graphique sans intervention du développeur.

\textbf{Les architectures microservices pour l'IA:} Le pattern de découplage entre la logique métier (Spring Boot) et les traitements d'inférence IA (FastAPI/Python) s'inscrit dans les recommandations architecturales des systèmes MLOps modernes. Ce découplage permet une évolution indépendante des deux couches --- notamment le remplacement ou la mise à jour du modèle LLM sans impact sur le backend Java et facilite la gestion des ressources computationnelles spécifiques à l'inférence (GPU, RAM étendue). Les stratégies de timeout et de fallback heuristique, mises en œuvre dans TalentPredict pour gérer la latence variable des appels Ollama, constituent une bonne pratique documentée dans la littérature sur la résilience des systèmes distribués~\cite{ch2-newman2015}.

\textbf{Les systèmes de recommandation pédagogique:} Drachsler et al. (2009) ont établi, dans le contexte des environnements d'apprentissage enrichis par la technologie (Technology Enhanced Learning), que la contextualisation des recommandations de formation en fonction du profil précis de l'apprenant  plutôt qu'une approche générique par filtrage collaboratif  augmente le taux de réussite de 23 à 35\%~\cite{ch2-drachsler2009}. Ce résultat fonde scientifiquement le choix de TalentPredict de générer des recommandations prescriptives croisant le profil PCM du collaborateur avec ses lacunes techniques identifiées, plutôt que de se limiter à des suggestions génériques de cours.

\section{Travaux existants et solutions similaires}
Afin de positionner l'apport de notre plateforme par rapport à l'existant, cette section propose une analyse critique des solutions actuellement disponibles sur le marché. Celles-ci ont été classées en trois grandes catégories : les outils d'évaluation psychométrique, les plateformes de tests techniques dédiées aux développeurs, et enfin les systèmes d'apprentissage pilotés par l'IA (LXP).
\subsection{Solutions d'évaluation psychométrique}
La première catégorie d'outils étudiée concerne l'évaluation des compétences comportementales (soft skills). Cette sous-section analyse les plateformes psychométriques les plus représentatives du marché, qu'elles s'appuient sur des questionnaires auto-déclaratifs classiques ou sur l'inférence de traits de personnalité par l'analyse algorithmique
\textbf{PerformanSe} est une plateforme psychométrique française proposant des questionnaires basés sur le Big Five et d'autres modèles validés scientifiquement, destinés à l'évaluation des compétences comportementales des collaborateurs. Son principe repose sur des questionnaires auto-déclaratifs en ligne, avec restitution d'un rapport de personnalité et de recommandations de développement.

\textbf{Apports :} validation psychométrique rigoureuse des outils, restitution structurée du profil comportemental, utilisation en contexte RH français bien établie.

\textbf{Limites :} le biais déclaratif n'est pas atténué par des mécanismes de validation externe ; les recommandations de formation restent génériques et non corrélées à une analyse des compétences techniques du collaborateur ; aucune automatisation du cycle d'évaluation n'est proposée.

\textbf{Crystal Knows} prédit le type DISC d'un individu à partir de l'analyse NLP de son profil LinkedIn public. Son approche est intéressante car elle tente d'inférer des traits comportementaux depuis des données numériques observables plutôt que déclaratives.

\textbf{Apports :} réduction partielle du biais déclaratif en s'appuyant sur des données textuelles publiques.

\textbf{Limites :} la validité prédictive du modèle DISC est contestée dans la littérature scientifique~\cite{ch2-pittenger1993} ; l'accuracy du classificateur n'est pas publiée par l'éditeur ; la solution est totalement déconnectée de toute évaluation technique ; la dépendance aux données LinkedIn introduit un biais lié à la complétude et au style rédactionnel des profils.

\subsection{Systèmes d'évaluation technique par tests}
Les systèmes d’évaluation technique traditionnels se concentrent sur la mesure de l'expertise opérationnelle par le biais de challenges et de QCM. Nous examinons ici deux solutions majeures de ce domaine afin d'en dégager les apports et les limites.
\textbf{HackerRank} est l'une des plateformes de référence pour l'évaluation des compétences en programmation, proposant des challenges algorithmiques dans une interface de développement intégrée, avec correction automatique et scoring.

\textbf{Apports :} objectivité des scores (la solution produit ou non le résultat attendu), large bibliothèque de problèmes couvrant de nombreux langages, fiabilité du scoring technique.

\textbf{Limites :} les tests sont statiques et ne s'adaptent pas dynamiquement au profil du collaborateur ; aucune donnée comportementale n'est collectée ni analysée ; l'absence de croisement avec le profil GitHub ou le CV ne permet pas de détecter les incohérences entre compétences déclarées et compétences démontrées.

\textbf{TestGorilla} propose une approche combinant tests techniques standardisés et questionnaires comportementaux, mais sans lien organique entre les deux dimensions d'évaluation.

\textbf{Apports :} facilité de déploiement, variété des tests disponibles.

\textbf{Limites :} absence d'adaptation dynamique de la difficulté ; les tests comportementaux restent auto-déclaratifs sans validation par des données comportementales asynchrones ; aucun mécanisme de suivi de la progression dans le temps n'est intégré.

\subsection{Plateformes d'expérience d'apprentissage pilotées par IA}
Le troisième volet de notre étude de l'existant se concentre sur l'étape de prescription de formations. Nous analysons ici les solutions d'apprentissage modernes, dites LXP, afin de comprendre comment elles exploitent les algorithmes de recommandation, mais aussi pourquoi leurs suggestions manquent souvent de précision objective.
\textbf{Degreed} est une LXP (Learning Experience Platform) qui agrège des contenus de formation issus de multiples sources (MOOCs, articles, vidéos) et utilise des algorithmes de recommandation pour suggérer des parcours adaptés au profil déclaré du collaborateur.

\textbf{Apports :} large catalogue de ressources, interface utilisateur soignée, intégration avec les outils RH existants.

\textbf{Limites :} les recommandations sont basées sur des compétences déclarées et non validées par des preuves objectives ; l'utilisation de LLM cloud pour la personnalisation expose les données RH à des risques de conformité RGPD ; aucune évaluation des compétences techniques n'est intégrée en amont de la recommandation.

\textbf{360Learning} se distingue par son approche d'apprentissage collaboratif où les experts internes créent les contenus de formation.

\textbf{Apports :} valorisation de l'expertise interne, forte dimension sociale de l'apprentissage.

\textbf{Limites :} la recommandation de formations est faiblement automatisée et repose largement sur des décisions manuelles des équipes L\&D ; l'absence d'analyse automatique des écarts de compétences (skill gap analysis) conduit à des propositions de formation non nécessairement alignées sur les besoins réels du collaborateur.

\section{Analyse comparative}
Afin de situer clairement TalentPredict AI dans son écosystème fonctionnel et technologique, cette section propose de confronter notre solution aux principales plateformes d'évaluation et de formation actuellement disponibles sur le marché RH.

\subsection{Tableau comparatif multi-critères}

Le tableau~\ref{tab:comparatif-multi-carecteres} synthétise ce positionnement face à six solutions concurrentes reconnues, en les évaluant selon des critères déterminants pour notre problématique : la capacité d'évaluation multidimensionnelle, les mécanismes de validation croisée, la souveraineté des données (conformité RGPD) et le modèle économique.
\begin{table}[H]
\centering
\tiny
\renewcommand{\arraystretch}{1.32}
\setlength{\tabcolsep}{2pt}
\begin{tabularx}{\textwidth}{|>{\RaggedRight\arraybackslash}p{0.10\textwidth}|*{7}{>{\RaggedRight\arraybackslash}X|}}
\hline
\textbf{Critère} & \textbf{Performan\-Se} & \textbf{Crystal Knows} & \textbf{Hacker\-Rank} & \textbf{Test\-Gorilla} & \textbf{Degreed} & \textbf{360\-Learning} & \textbf{Talent\-Predict AI} \\
\hline
Évaluation technique & Non & Non & Oui, tests statiques & Oui, partielle & Non & Non & Oui, multi-sources \\
\hline
Évaluation comportementale & Oui, Big Five & Oui, DISC & Non & Oui, déclarative & Non & Non & Oui, PCM + IA locale \\
\hline
Validation croisée & Non & Non & Non & Non & Non & Non & Oui, GitHub -- CV -- Tests \\
\hline
Analyse GitHub & Non & Non & Non & Non & Non & Non & Oui \\
\hline
Recom\-mandation prescriptive & Partielle & Non & Non & Non & Oui, générique & Oui, générique & Oui, adaptée PCM \\
\hline
Souveraineté des données & Non, cloud & Non, cloud & Non, cloud & Non, cloud & Non, LLM cloud & Non, cloud & Oui, Edge AI / on-premise \\
\hline
Scoring multi-sources & Non & Non & Non & Non & Non & Non & Oui, 500+ alias \\
\hline
Suivi de progression & Non & Non & Non & Non & Oui & Oui & Oui, Kanban + gamification \\
\hline
Conformité RGPD données sensibles & Partielle & Non & Non & Non & Non & Non & Oui, totale \\
\hline
Coût d'exploitation & Licence élevée & Abonne\-ment & Licence & Licence & Licence élevée & Licence élevée & Open-source \\
\hline
\end{tabularx}
\renewcommand{\arraystretch}{1.05}
\caption{Analyse comparative multi-critères des solutions existantes et de TalentPredict AI}
\label{tab:comparatif-multi-carecteres}
\end{table}

Afin de synthétiser cette étude comparative, la figure~\ref{fig:positionnement-existant} propose une lecture graphique du marché sous la forme d'une matrice de positionnement. Cette cartographie permet de visualiser l'écart entre les solutions actuelles et TalentPredict AI, en croisant le degré d'objectivité de l'évaluation avec le niveau de souveraineté des données garanti par l'architecture.
\begin{figure}[H]
    \centering
    \includegraphics[width=\textwidth]{images/position existant.png}
    \caption{Matrice de positionnement concurrentiel de TalentPredict AI selon les axes de souveraineté et d'objectivité.}
    \label{fig:positionnement-existant}
\end{figure}
\subsection{Synthèse analytique des différenciateurs}

L'analyse comparative révèle trois fractures structurelles dans l'offre existante.

La première fracture est l'opposition irrésolue entre objectivité et complétude. Les solutions techniques (HackerRank, TestGorilla) offrent une bonne objectivité mais ignorent la dimension comportementale. Les solutions psychométriques (PerformanSe, Crystal Knows) couvrent la dimension comportementale mais de manière déclarative. Aucune solution n'intègre les deux dimensions dans un pipeline unifié avec validation croisée.

La deuxième fracture est le manque de validation factuelle des compétences déclarées. Toutes les solutions analysées acceptent sans vérification les compétences que le collaborateur ou son manager affirme maîtriser. TalentPredict introduit une rupture en confrontant systématiquement les déclarations à des preuves asynchrones (historique GitHub, résultats de tests génératifs).

La troisième fracture est l'incompatibilité entre puissance analytique LLM et souveraineté des données. Les solutions LXP modernes commencent à intégrer des LLM pour la personnalisation, mais via des APIs cloud qui exposent des données RH sensibles à des tiers. TalentPredict résout cette tension par l'inférence locale via Ollama.

\section{Positionnement du projet TalentPredict AI}
À la lumière de l'état de l'art et de l'analyse comparative présentés précédemment, cette section a pour objectif de définir le positionnement exact de TalentPredict AI. Pour ce faire, nous commencerons par formaliser les limites majeures des approches actuelles, justifiant ainsi les choix architecturaux et fonctionnels de notre solution.
\subsection{Lacunes identifiées dans la littérature et sur le marché}

L'analyse des travaux existants et des solutions disponibles permet d'identifier trois lacunes principales qui restent insuffisamment traitées.

\textbf{Lacune 1 : Absence de pipeline open-source unifié combinant évaluation technique et comportementale:} Aucune solution open-source ne propose, à ce jour, un système intégré permettant d'évaluer simultanément les compétences techniques d'un collaborateur (via analyse de code, parsing de CV, tests adaptatifs) et ses compétences comportementales (via modèle psychométrique et analyse sémantique), en croisant les résultats des deux dimensions pour produire un profil unifié et des recommandations prescriptives.

\textbf{Lacune 2 : Absence de validation croisée multi-sources avec déduplication sémantique:} Les systèmes d'évaluation existants traitent chaque source de données (CV, tests, LinkedIn) de manière indépendante. Aucun mécanisme de déduplication et de normalisation sémantique entre sources hétérogènes n'est proposé, conduisant à des évaluations fragmentées et potentiellement contradictoires. TalentPredict répond à cette lacune avec son algorithme de normalisation par dictionnaire d'alias (500+ variantes canoniques) et son moteur de scoring par triangulation multi-sources.

\textbf{Lacune 3 : Dépendance au cloud pour les solutions LLM appliquées aux RH:} Les solutions utilisant des LLM pour l'analyse des compétences comportementales externalisent toutes leur traitement vers des infrastructures cloud, compromettant la conformité RGPD pour les données psychologiques (Article 9). TalentPredict répond à cette lacune par un déploiement 100\% local via Ollama, garantissant la souveraineté totale des données.

\subsection{Contributions spécifiques de TalentPredict AI}

TalentPredict AI apporte trois contributions originales répondant directement aux lacunes identifiées.

La première contribution est une architecture hybride n8n/FastAPI pour l'Edge AI RH. L'orchestration des workflows comportementaux via n8n permet de séparer la logique de prompt engineering de la logique applicative, offrant une maintenabilité supérieure aux approches monolithiques tout en garantissant l'inférence locale. Cette architecture constitue, à notre connaissance, une approche originale dans le domaine de la HR Tech open-source.

La deuxième contribution est un algorithme de scoring multi-sources avec normalisation sémantique. Le pipeline d'évaluation technique de TalentPredict implémente un mécanisme de déduplication basé sur un dictionnaire de 500+ alias canoniques (par exemple : reactjs $\rightarrow$ React, k8s $\rightarrow$ Kubernetes, sklearn $\rightarrow$ Scikit-learn), couplé à un algorithme de nivellement par triangulation (Expert, Advanced, Intermediate, Beginner) selon le nombre de sources confirmant la compétence et la force du signal dans chacune. Ce mécanisme atténue structurellement le biais déclaratif en pondérant les sources objectives (GitHub, tests) davantage que les sources déclaratives (CV).

La troisième contribution est un pipeline PCM heuristique à reproductibilité totale couplé à une validation LLM sémantique. Le calcul des scores PCM est implémenté de manière entièrement algorithmique en JavaScript, garantissant une reproductibilité de 100\% (même entrée $\rightarrow$ même sortie, quelle que soit l'exécution), ce qu'un calcul LLM ne peut pas assurer en raison de la stochasticité inhérente à la génération de texte (accord inter-exécutions typiquement de 70 à 80\% pour les tâches de classification psychométrique~\cite{ch2-xu2024}). La dimension LLM intervient uniquement pour l'analyse sémantique des données textuelles et la synthèse narrative, où la variabilité est acceptable et même souhaitable.

Le positionnement scientifique et technique de TalentPredict AI repose sur trois piliers fondamentaux, ou contributions, qui structurent l'innovation du projet. Le schéma~\ref{fig:position-platfome} détaille ces contributions (C1, C2 et C3) en mettant en relation les objectifs méthodologiques avec la stack technologique et les mécanismes d'orchestration mis en œuvre.
\begin{figure}[H]
    \centering
    \includegraphics[width=\textwidth]{images/position platforme.png}
    \caption{Positionnement scienti-
fique et technique de Ta-
lentPredict AI}
    \label{fig:position-platfome}
\end{figure}
\subsection{Valeur ajoutée synthétique}
Afin de résumer la proposition de valeur de notre plateforme, le tableau~\ref{tab:valeur-ajoutee-synthetique} illustre comment TalentPredict AI apporte une réponse technologique et méthodologique directe à chacune des lacunes identifiées dans l'état de l'art.
\begingroup
\small
\renewcommand{\arraystretch}{1.35}
\setlength{\tabcolsep}{5pt}
\begin{longtable}{|>{\RaggedRight\arraybackslash}p{0.30\textwidth}|>{\RaggedRight\arraybackslash}p{0.62\textwidth}|}
\label{tab:valeur-ajoutee-synthetique}\\
\hline
\textbf{Lacune identifiée} & \textbf{Réponse de TalentPredict AI} \\
\hline
\endfirsthead
\hline
\textbf{Lacune identifiée} & \textbf{Réponse de TalentPredict AI} \\
\hline
\endhead
Fragmentation technique/comportementale & Pipeline unifié combinant l'évaluation technique (GitHub, CV, QCM et code) avec l'évaluation comportementale (PCM, Ollama et scénarios IA). \\
\hline
Biais déclaratif non contrôlé & Validation croisée multi-sources : les compétences déclarées sont confrontées à des preuves asynchrones et objectives. \\
\hline
Dépendance cloud / RGPD & Inférence locale via Ollama, sans externalisation des données RH sensibles. \\
\hline
Recommandations génériques & Recommandations prescriptives avec une modalité pédagogique adaptée au profil PCM du collaborateur. \\
\hline
Coûts de licence élevés & Stack open-source fondée sur n8n, Ollama, FastAPI, Spring Boot, PostgreSQL et Angular. \\
\hline
\caption{Synthèse des contributions de TalentPredict AI face aux lacunes identifiées}
\end{longtable}
\renewcommand{\arraystretch}{1.05}
\endgroup

\section*{Conclusion}
\phantomsection
\addcontentsline{toc}{section}{Conclusion}

Ce chapitre établit le socle théorique de TalentPredict AI en démontrant que face aux limites des méthodes d'évaluation traditionnelles (biais déclaratif, coûts prohibitifs, fragmentation) et à l'incapacité des solutions actuelles du marché à garantir simultanément l'objectivité factuelle, l'analyse technique et comportementale, ainsi que la souveraineté des données, la conception d'une plateforme innovante s'impose. Pour combler ces lacunes, notre étude a directement dicté les choix architecturaux qui seront détaillés lors des quatre prochains sprints : l'exploitation des données GitHub comme indicateur comportemental, combinée à la puissance des LLMs via le déploiement local de llama3.2 sous Ollama pour assurer la conformité RGPD. Ces technologies s'articulent autour du modèle scientifique PCM, d'une orchestration des workflows par n8n, et d'un scoring multi-sources indispensable pour neutraliser efficacement les biais d'auto-évaluation.




\chapter{ Sprint 1 : Mise en place de l’Architecture, de la Sécurité et du Système d’Authentification}
\section*{Introduction}
\phantomsection
\addcontentsline{toc}{section}{Introduction}
\phantomsection 
\noindent
Le premier sprint de TalentPredict a été conçu comme une phase de fondation stratégique, essentielle à la pérennité de notre solution. Avant d'intégrer l'intelligence artificielle, nous avons jugé primordial de stabiliser l'écosystème technique pour offrir un environnement à la fois sécurisé et performant. Cette étape s'est concrétisée par une architecture conteneurisée garantissant une agilité de déploiement, doublée d'un protocole de sécurité strict reposant sur les standards JWT et BCrypt. En érigeant ce socle robuste, nous avons ainsi préparé le terrain pour une montée en charge fluide des futurs modules prédictifs.
 \section{Objectifs du Sprint 1}
L'objectif majeur de ce sprint réside dans la livraison d'un incrément logiciel opérationnel, constituant le cœur névralgique de TalentPredict. Loin de se limiter à un simple prototype, cette phase vise à instaurer une infrastructure mature et structurée, conçue pour intégrer harmonieusement les modules métiers complexes des étapes ultérieures. Notre ambition est de garantir, dès l'achèvement de ce sprint, une expérience utilisateur fluide au sein d'un environnement sécurisé, normé et optimisé pour répondre aux exigences de performance du projet.
\subsubsection{1. Objectifs Fonctionnels et Interactifs}
\begin{itemize}
    \item \textbf{Système d'Authentification Complet (IAM) :}
    \begin{itemize}
        \item Mise en œuvre du cycle d’inscription (Register), consolidé par un système de validation d’adresse email afin de garantir l’authenticité des comptes créés.
        \item Déploiement d’une authentification sécurisée (Login), enrichie d’une gestion rigoureuse des erreurs de saisie pour optimiser l’expérience utilisateur et la sécurité des accès.
        \item Implémentation de la procédure de récupération de mot de passe (Forgot Password), s'appuyant sur l'émission de tokens temporaires pour assurer la confidentialité et l'intégrité des données.
    \end{itemize}
    
    \item \textbf{Gestion Dynamique des Profils :}
    \begin{itemize}
        \item Interface de consultation du profil utilisateur permettant de visualiser de manière structurée les données d’identité ainsi que les rôles associés.
        \item Formulaires d'édition optimisés pour la mise à jour des informations professionnelles, telles que l'intitulé du poste, la biographie et les réseaux sociaux.
    \end{itemize}
    
    \item \textbf{Interface Utilisateur (UX/UI) de Base :}
    \begin{itemize}
        \item Déploiement d’une page d'accueil responsive, intégrant les sections stratégiques : Hero, présentation des fonctionnalités et appels à l’action (CTA).
        \item Mise en place du Dashboard Shell, incluant une navigation latérale (Sidenav) ergonomique et une barre supérieure (Topbar) dotée du menu utilisateur.
        
    \end{itemize}
    
    \item \textbf{Autorisation Granulaire :}
    \begin{itemize}
        \item Développement de gardes de navigation (Route Guards) afin de sécuriser les accès et d'interdire systématiquement l'affichage des pages protégées aux utilisateurs non authentifiés.
        \item Adaptation dynamique de l’interface en fonction des privilèges de l'utilisateur, permettant notamment de restreindre la visibilité de certaines sections, telle que la gestion des administrateurs, aux seuls rôles autorisés.
    \end{itemize}
\end{itemize}

\subsubsection{2. Objectifs Techniques et Opérationnels}

\begin{itemize}
    \item \textbf{Standardisation de l'Architecture :}
    \begin{itemize}
        \item Initialisation du projet en \textit{Monolithe Modulaire} sous \texttt{Spring Boot 3.4.2}.
        \item Organisation rigoureuse des \textit{packages} (\texttt{Controller}, \texttt{Service}, \texttt{Repository}, \texttt{DTO}).
    \end{itemize}
    
    \item \textbf{Sécurité et Cryptographie :}
    \begin{itemize}
        \item Mise en place de \texttt{Spring Security} avec une politique de session \textit{stateless}.
        \item Implémentation du standard \texttt{JWT} (Génération, Signature HS256, Extraction et Validation).
        \item Hachage des secrets via l'algorithme \texttt{BCrypt}.
    \end{itemize}
    
    \item \textbf{Infrastructure et DevOps :}
    \begin{itemize}
        \item Orchestration complète via \texttt{Docker Compose} pour assurer la parité des environnements.
        \item Configuration sécurisée des variables d'environnement (\texttt{.env}).
    \end{itemize}
    
    \item \textbf{Robustesse et Qualité de Code :}
    \begin{itemize}
        \item Centralisation de la gestion des exceptions via un \texttt{GlobalExceptionHandler}.
        \item Intercepteur HTTP côté \texttt{Angular} pour l'injection automatique du \textit{token} dans les \textit{headers}.
        \item Persistance via \texttt{PostgreSQL}; initialisation automatique du schéma avec \texttt{JPA} et \texttt{Hibernate}.
    \end{itemize}
    
    \item \textbf{Exposition et Test de l'API :}
    \begin{itemize}
        \item Configuration de \texttt{Swagger} pour la documentation et les tests en temps réel.
        \item Mise en œuvre des politiques de \texttt{CORS} restrictives.
    \end{itemize}
\end{itemize}

\paragraph{Critères d’acceptation (Definition of Done)}

Pour valider chaque fonctionnalité de ce sprint, les critères suivants ont été appliqués :

\begin{itemize}
    \item \textbf{Fonctionnalité} : Conformité totale avec la \textit{User Story} et validation des tests unitaires.
    \item \textbf{Sécurité} : Endpoints protégés par JWT et mots de passe hachés via \textit{BCrypt}.
    \item \textbf{Qualité UI} : Interface responsive (mobile/desktop) et absence d'erreurs console.
    \item \textbf{Documentation} : Exposition des APIs via \textit{Swagger} et code documenté.
    \item \textbf{Déploiement} : Construction réussie de l'image \textit{Docker} et intégration au \textit{Compose}.
\end{itemize}



\section{Séance de Sprint Planning}

La séance de Sprint Planning a réuni l’équipe de développement, composée de Mohamed Ben Abdeladhim et Rahma Barouni, ainsi que le Product Owner, Mr Med Amine Gharbi, et notre encadrante pédagogique, Mme Ben Aïcha Takwa. Cette réunion a permis de définir le périmètre fonctionnel et technique de ce premier sprint, crucial pour la solidité de la plateforme TalentPredict.

Les points clés suivants ont été déterminés lors de cette séance :
\begin{itemize}
    \item \textbf{Durée :} La durée de cette itération a été fixée à quatre semaines. Ce délai a été jugé nécessaire pour mettre en place une infrastructure conteneurisée robuste et implémenter un système de sécurité complexe (\texttt{JWT}) garantissant l'intégrité des futures données d'évaluation.
    
    \item \textbf{Priorités Absolues :}
    \begin{itemize}
        \item Mise en place de l'architecture \textit{Monolithe Modulaire} sous \texttt{Spring Boot} et du schéma relationnel \texttt{PostgreSQL}.
        \item Développement du module \textit{Authentication \& Security} (Gestion des \textit{tokens} \texttt{JWT}, hachage \texttt{BCrypt} et contrôle d'accès \texttt{RBAC}).
        \item Implémentation du système d'Inscription et de Connexion pour les deux rôles principaux : Employé et Administrateur RH.
        \item Création de la page d’accueil et du \textit{Dashboard Shell} (\textit{Sidenav}, \textit{Topbar}) pour stabiliser l'interface utilisateur.
    \end{itemize}
    
    \item \textbf{Priorités Secondaires :}
    \begin{itemize}
        \item Configuration de la documentation interactive de l'API via \texttt{Swagger}.
        \item Centralisation de la gestion des erreurs (\textit{Global Exception Handling}) et configuration des politiques \texttt{CORS}.
        \item Mise en place du service de récupération de mot de passe et de gestion des profils utilisateurs.
    \end{itemize}
    
    \item \textbf{Prérequis :} La disponibilité de l'environnement de virtualisation \texttt{Docker} et l'initialisation du \textit{repository} \texttt{Git} sont essentielles pour permettre une collaboration fluide et un déploiement continu dès les premières phases du projet.
\end{itemize}

À la fin de cette réunion, les besoins ont été décomposés en User Stories (US) et estimés en jours-homme. La priorité maximale a été donnée aux briques de sécurité et à l'infrastructure, car elles constituent la condition \textit{sine qua non} pour le déploiement des modules d'intelligence artificielle prévus dans les sprints suivants.
\section{Backlog du Sprint }

Le \textit{Sprint Backlog} constitue le plan opérationnel détaillé de cette itération. Il résulte de la décomposition des \textit{User Stories} (besoins utilisateurs) en tâches techniques élémentaires, permettant ainsi une exécution structurée et un suivi précis de l'avancement.

Pour chaque tâche, une estimation en Jours-Homme (JH) a été réalisée afin de quantifier la charge de travail et d'assurer une répartition équitable au sein du binôme de développement. Le tableau \ref{tab:backlog_sprint1} ci-dessous présente la feuille de route exhaustive de ce premier sprint :
\begin{longtable}{|c|>{\raggedright\arraybackslash}p{3cm}|>{\raggedright\arraybackslash}p{7cm}|c|>{\centering\arraybackslash}p{1.8cm}|}
    % --- First Page Header ---
    \hline
    \rowcolor[HTML]{F2F2F2}
    \textbf{ID} & \textbf{Module / US} & \textbf{Tâches de Développement} & \textbf{JH} & \textbf{Priorité} \\
    \hline
    \endfirsthead

    % --- Subsequent Pages Header ---
    \hline
    \rowcolor[HTML]{F2F2F2}
    \textbf{ID} & \textbf{Module / US} & \textbf{Tâches de Développement} & \textbf{JH} & \textbf{Priorité} \\
    \hline
    \endhead

    % --- Footer for pages before the last ---
    \hline
    \multicolumn{5}{|r|}{\textit{Suite à la page suivante...}} \\
    \hline
    \endfoot

    % --- Final Footer ---
    \hline
    \endlastfoot

    % --- Table Data ---
    1.1 & Auth / Enregistrement & 
    \begin{itemize}[leftmargin=*, nosep]
        \item Formulaire de \textit{Registration} (\texttt{Angular}).
        \item Service \textit{Backend} et hachage \texttt{BCrypt}.
        \item Validation et gestion des doublons.
    \end{itemize} & 4 & Haute \\
    \hline
    
    1.2 & Auth / Connexion JWT & 
    \begin{itemize}[leftmargin=*, nosep]
        \item Configuration \texttt{Spring Security} (\textit{Stateless}).
        \item Développement du \texttt{TokenProvider} (JWT).
        \item Composant \textit{Login} et stockage du \textit{token}.
    \end{itemize} & 5 & Haute \\
    \hline
    
    1.3 & Auth / Forget Password & Interface "Mot de passe oublié" et logique de génération de \textit{token} de récupération. & 3 & Haute \\
    \hline
    
    1.4 & Services / Notifications & 
    \begin{itemize}[leftmargin=*, nosep]
        \item Intégration \texttt{Spring Mail} (SMTP).
        \item \textit{Templates} d'emails et envoi asynchrone.
    \end{itemize} & 3 & Moyenne \\
    \hline
    
    1.5 & Profil / Différencié & 
    \begin{itemize}[leftmargin=*, nosep]
        \item \textbf{User :} Champs techniques (GitHub, CV).
        \item \textbf{Admin :} Champs RH.
        \item \textit{Endpoints} CRUD et affichage conditionnel.
    \end{itemize} & 5 & Haute \\
    \hline
    
    1.6 & Profil / Sécurité & Changement de MDP, option de désactivation et audit des connexions. & 3 & Moyenne \\
    \hline
    
    1.7 & Interface / Navigation & 
    \begin{itemize}[leftmargin=*, nosep]
        \item \textit{Design} Landing Page et \textit{Layout}.
        \item \textit{Route Guards} et Intercepteurs JWT.
        \item Système de Toasts/\textit{Snackbars}.
    \end{itemize} & 5 & Haute \\
    \hline
    
    1.8 & Infra / Environnement & 
    \begin{itemize}[leftmargin=*, nosep]
        \item Configuration \texttt{DOCKER Compose}.
        \item Setup \texttt{PostgreSQL} (JPA/Hibernate).
        \item Gestion des secrets via \texttt{.env}.
    \end{itemize} & 3 & Haute \\
\end{longtable}
\addtocounter{table}{-1} % Adjust counter for the manual caption below if needed
\begin{center}
    \small Table \refstepcounter{table}\thetable: Backlog du Sprint 1 \label{tab:backlog_sprint1}
\end{center}
\section{Démarrage du Sprint 1}

Une fois la phase de planification finalisée et les priorités clairement établies, l’équipe a entamé la phase de réalisation technique. Cette étape cruciale du cycle de vie Agile consiste à transformer les exigences du backlog en un incrément logiciel fonctionnel à travers un processus itératif de codage, d'intégration et de test. L’exécution de ce premier sprint a été guidée par une volonté de rigueur architecturale, visant à stabiliser le socle de TalentPredict avant d'entamer les développements liés à l'intelligence artificielle. Cette section retrace le parcours technique de ce sprint, en détaillant les activités menées, les résultats concrets produits, ainsi que la gestion des imprévus technologiques qui ont jalonné le développement.
\clearpage
\subsubsection{1. Activités réalisées}

L'équipe a procédé à la mise en place des fondations logicielles et infrastructurelles suivantes :

\begin{itemize}
    \item \textbf{Initialisation du projet :} Création de la structure modulaire du \textit{backend} (\texttt{Spring Boot 3.4.2}) et du \textit{frontend} (\texttt{Angular 17}) avec mise en place des standards de nommage et de codage.
    \item \textbf{Développement de la couche de sécurité :} Implémentation du système d'authentification \texttt{JWT}, incluant la signature des \textit{tokens} et la configuration de \texttt{Spring Security} pour le contrôle d'accès par rôle (\texttt{RBAC}).
    \item \textbf{Conception et déploiement de la base de données :} Modélisation des entités \texttt{User} et \texttt{Profile} différenciés, et exécution des scripts de création via \texttt{JPA}/\texttt{Hibernate} sous \texttt{PostgreSQL}.
    \item \textbf{Intégration des services de notification :} Configuration du serveur \texttt{SMTP} et développement du service asynchrone d'envoi d'emails pour l'enregistrement et la récupération de mot de passe.
    \item \textbf{Conception de l'interface utilisateur (UI) :} Développement de la \textit{Landing Page}, des formulaires d'authentification et de la structure globale du \textit{Dashboard} (\textit{Sidenav}, \textit{Topbar}).
    \item \textbf{Mise en place de l'infrastructure DevOps :} Rédaction et test des configurations \texttt{Docker} pour garantir un environnement de développement reproductible.
\end{itemize}

\subsubsection{2. Livrables produits}

À l'issue de cette itération, les livrables suivants ont été intégrés :

\begin{itemize}
    \item \textbf{Module d'Authentification Fonctionnel :} Un système complet de \textit{Login/Register} sécurisé par \textit{token}.
    \item \textbf{Profils Différenciés :} Des interfaces et des \textit{endpoints} spécifiques pour les données "User" (Employé) et "Admin" (RH).
    \item \textbf{Socle d'Infrastructure Conteneurisé :} Un fichier \texttt{docker-compose.yml} orchestrant l'application, le \textit{frontend} et la base de données.
    \item \textbf{Documentation de l'API :} Une interface \texttt{Swagger UI} exposant tous les \textit{endpoints} réalisés pour faciliter les tests.
    \item \textbf{Composants UI Réutilisables :} Une bibliothèque de composants \texttt{Angular} (Boutons, \textit{Cards}, \textit{Toasts}) basée sur \texttt{Angular Material}.
\end{itemize}

\subsubsection{3. Difficultés rencontrées et traitement}

L'exécution de ce sprint a été marquée par trois défis techniques majeurs :

\begin{itemize}
    \item \textbf{Difficulté 1 : Conflits de CORS}
    \begin{itemize}
        \item \textit{Description :} Le navigateur bloquait les requêtes entre le \textit{frontend} (port 4200) et le \textit{backend} (port 8001).
        \item \textit{Traitement :} Implémentation d'un \textit{bean} de configuration globale \texttt{CorsConfiguration} autorisant explicitement les \textit{headers} et méthodes requis.
    \end{itemize}
    \vspace{2mm}
    \item \textbf{Difficulté 2 : Gestion de l'authentification \textit{Stateless}}
    \begin{itemize}
        \item \textit{Description :} Difficulté à injecter le \textit{token} \texttt{JWT} de manière transparente dans toutes les requêtes \textit{frontend}.
        \item \textit{Traitement :} Création d'un \textit{Intercepteur HTTP} côté \texttt{Angular} pour automatiser l'ajout du \textit{header} \texttt{Authorization} à chaque appel API.
    \end{itemize}
    \vspace{2mm}
    \item \textbf{Difficulté 3 : Parité des environnements avec \texttt{Docker}}
    \begin{itemize}
        \item \textit{Description :} Des comportements différents observés entre le déploiement \texttt{Docker} et l'exécution locale.
        \item \textit{Traitement :} Standardisation des variables d'environnement via des fichiers \texttt{.env} et utilisation de \texttt{Docker Multi-stage builds} pour isoler les étapes de construction.
    \end{itemize}
\end{itemize}
\subsection{Analyse}

Cette phase d'analyse identifie les interactions spécifiques des deux types d'acteurs de TalentPredict et définit les garde-fous techniques nécessaires à la sécurisation de leurs espaces respectifs.

\subsubsection*{1. Identification des Acteurs}
\begin{itemize}
    \item \textbf{L'Employé (Utilisateur) :} C'est l'acteur qui utilise la plateforme pour ses auto-évaluations. Il doit pouvoir s'enregistrer, se connecter et enrichir son profil (liens \texttt{GitHub}/\texttt{LinkedIn}) qui servira de base à l'analyse de ses compétences techniques et comportementales.
    \item \textbf{Le RH (Administrateur) :} C'est l'acteur responsable du suivi et de la gouvernance. Il possède un accès privilégié lui permettant de visualiser l'ensemble des talents de l'entreprise. Dès ce sprint, il doit disposer d'un accès sécurisé à son interface de gestion pour préparer le suivi des futurs résultats.
    \item \textbf{Service Mail}:C'est un système externe automatisé qui interagit avec le backend. Son rôle est critique dans le flux de sécurité : il est chargé d'envoyer les emails transactionnels de manière asynchrone, notamment lors du processus de récupération de compte ("Mot de passe oublié") ou de la confirmation de sécurité, garantissant ainsi l'autonomie et la protection de l'utilisateur.
\end{itemize}

\subsubsection*{2. Clarification du besoin}
Le besoin pour ce premier sprint est de construire le socle de confiance entre l'Employé et le Admin RH :
\begin{itemize}
    \item \textbf{Séparation des espaces :} Garantir que l'Employé n'accède qu'à ses propres données, tandis que le RH accède à une vue globale consolidée.
    \item \textbf{Collecte de données structurées :} Permettre à l'Employé de fournir les vecteurs d'analyse (liens externes) nécessaires aux sprints suivants.
    \item \textbf{Système de confiance :} Mise en place d'une authentification \textit{Stateless} (\texttt{JWT}) pour assurer que chaque action sur la plateforme est tracée et autorisée selon le rôle (\texttt{Admin RH} ou \texttt{Employé}).
\end{itemize}

\subsubsection*{3. Scénarios de validation}
\begin{itemize}
    \item \textbf{Enregistrement de l'Employé :} Création d'un compte, hachage du mot de passe et redirection vers son \textit{dashboard} personnel.
    \item \textbf{Connexion sécurisée :} Saisie des identifiants (Email/Mot de passe) et réception d'un \textit{token} \texttt{JWT} valide pour l'accès aux services.
    \item \textbf{Détention de Rôle :} Tentative d'accès d'un Employé à la zone RH ; le système doit rejeter l'accès (Erreur \texttt{403 Forbidden}).
    \item \textbf{Récupération d'accès :} Envoi d'un email de réinitialisation de mot de passe en cas d'oubli, utilisable par les deux acteurs.
\end{itemize}

\subsubsection*{4. Contraintes techniques}
\begin{itemize}
    \item \textbf{Modèle de données évolutif :} Utilisation de \texttt{PostgreSQL} pour stocker les profils différenciés de l'Employé et du RH.
    \item \textbf{Standard Sécurité \texttt{JWT} :} Utilisation de \textit{tokens} auto-porteurs pour une authentification performante et décentralisée.
    \item \textbf{Interfaces \texttt{Angular Material} :} Développement de deux vues distinctes au sein du même \textit{Shell} applicatif, s'adaptant dynamiquement au rôle détecté dans le \textit{token}.
    \item \textbf{Infrastructure \texttt{Docker} :} Déploiement unifié du \textit{Backend} \texttt{Spring Boot} et du \textit{Frontend} \texttt{Angular} dans des conteneurs isolés et sécurisés.
\end{itemize}
\subsubsection*{5. Diagramme de cas d'utilisation}
Le premier sprint, dédié à la Gestion des Utilisateurs et de l’Authentification, pose les fon-
dations de la plateforme TalentPredict. Ce diagramme de cas d’utilisation \ref{fig:usecase_du_sprint1} synthétise les
fonctionnalités essentielles permettant de sécuriser et de personnaliser l’accès à l’application.

\begin{figure}[H]
    \centering
    \includegraphics[width=\textwidth, keepaspectratio]{images/usecasSprint1.png}
    \caption{Diagramme de cas d'utilisation - Sprint 1 }
    \label{fig:usecase_du_sprint1}
\end{figure}

Le tableau \ref{tab:legende_sprint1} suivant résume la sémantique visuelle et l'application concrète des différentes relations modélisées dans notre diagramme pour ce premier sprint.

\begin{table}[H]
    \centering
    \renewcommand{\arraystretch}{1.5} % Aère les lignes du tableau
    \begin{tabularx}{\textwidth}{@{} >{\raggedright\arraybackslash}p{3cm} >{\raggedright\arraybackslash}p{4.5cm} X @{}}
        \toprule
        \textbf{Type de Relation} & \textbf{Représentation Graphique} & \textbf{Application Concrète } \\
        \midrule
        \textbf{Généralisation} \newline \textit{\small{(Héritage)}} & 
        Trait plein continu + flèche triangulaire vide pointant vers l'acteur parent. & 
        Les acteurs Employé et Admin RH héritent de l'acteur de base Utilisateur et de tous ses accès par défaut (ex: \textit{Se connecter}). \\
        
        \textbf{Association} \newline \textit{\small{(Communication)}} & 
        Ligne continue simple. & 
        Lien direct entre l'Utilisateur et les fonctions de base, ou le Service Mail pour l'envoi de liens. \\
        
        \textbf{Inclusion} \newline \textit{\small{(«include»)}} & 
        Flèche en pointillés avec le mot-clé \texttt{«include»}, orientée vers le cas inclus. & 
        L'action \textit{Mot de passe oublié} déclenche obligatoirement \textit{Envoyer un lien de récupération}. \\
        
        \textbf{Extension} \newline \textit{\small{(«extend»)}} & 
        Flèche en pointillés avec le mot-clé \texttt{«extend»}, orientée vers le cas de base. & 
        L'action \textit{Modifier le profil} offre des sous-options supplémentaires comme \textit{Gérer les préférences} ou \textit{Changer le mot de passe}. \\
        \bottomrule
    \end{tabularx}
    \caption{Légende des associations du diagramme  cas d’utilisation - Sprint 1}
    \label{tab:legende_sprint1}
\end{table}
\clearpage
\subsection{Conception}
Cette phase de conception, initiée suite à l'identification des acteurs et à la définition des besoins fonctionnels, constitue le pivot essentiel entre l'analyse théorique et la réalisation technique. Elle a pour objectif de structurer l'architecture globale du système, de modéliser avec précision le schéma des données et de spécifier les interfaces de communication (API). Cette étape s'avère déterminante pour doter TalentPredict d'un socle technique robuste et sécurisé, capable de soutenir l'intégration évolutive des modules d'intelligence artificielle programmés lors des cycles de développement ultérieurs.
\subsubsection*{1. Objectifs de conception}
\begin{itemize}
    \item \textbf{\textit{Stratégie de Monolithe Modulaire } :} Afin d'éviter la complexité prématurée des microservices tout en préparant le futur, nous avons opté pour un monolithe modulaire. Chaque fonctionnalité (Auth, Évaluation, IA) est encapsulée dans son propre \textit{package}, communiquant via des services définis, respectant ainsi le principe de responsabilité unique (SRP).
    \item \textbf{\textit{Authentification sans état (Stateless) et Montée en charge (Scalability)} :} En choisissant \texttt{JWT}, nous éliminons le besoin de \texttt{HttpSession} côté serveur. Cela libère de la mémoire RAM et permet au \textit{backend} de monter en charge sans se soucier de la réplication des sessions.
    \item \textbf{Intégrité des données (\textit{Data Integrity}) :} Assurer que chaque utilisateur possède un profil unique dès sa création via des contraintes de clés étrangères au niveau de la base de données.
\end{itemize}

\subsubsection*{2. Architecture Technique }
L’architecture de TalentPredict repose sur une structure N-Tier (multicouche) moderne, conçue pour être modulaire, sécurisée et facilement déployable. \ref{fig:archi_globale_s1} expose l’organisation globale de l’infrastructure technique ainsi que l’articulation des flux de données entre les différents composants :
\begin{figure}[H]
    \centering
    \includegraphics[width=\textwidth]{images/Diagramme d'Architecture.png}
    \vspace{2mm}
    \caption{Architecture globale et flux de données du Sprint 1}
    \label{fig:archi_globale_s1}
\end{figure}

\paragraph*{Infrastructure et environnement \texttt{Docker}}
\begin{itemize}
    \item L’ensemble du système est orchestré via \texttt{Docker Compose}, assurant une portabilité complète entre les environnements de développement et de production, ainsi qu’une architecture multi-conteneurs.
    \item Le réseau \texttt{Docker} privé \texttt{talentpredict-net} permet une communication isolée et contrôlée entre les services, garantissant qu’aucun accès direct à la base de données n’est possible depuis l’extérieur. L’accès aux données est exclusivement réalisé via le \textit{backend}.
    \item La persistance des données est assurée par un \textit{Docker Volume} attaché au conteneur \texttt{PostgreSQL}, permettant de conserver les informations même après arrêt ou redémarrage des conteneurs.
    \item La gestion des configurations sensibles est externalisée via un fichier \texttt{.env}, contenant les secrets (\texttt{JWT\_SECRET}, mots de passe de base de données, identifiants \texttt{SMTP}). Cette approche respecte les bonnes pratiques du modèle \textit{Twelve-Factor App} et améliore la sécurité globale du système.
\end{itemize}

\paragraph*{Couche Client (\textit{Frontend Tier})}
\begin{itemize}
    \item L’interface utilisateur est développée en \texttt{Angular 17} et déployée via un serveur \texttt{Nginx} sous forme d’application SPA (\textit{Single Page Application}).
    \item Un intercepteur HTTP \texttt{Angular} joue un rôle central dans le mécanisme de sécurité en injectant automatiquement le jeton \texttt{JWT} dans l’en-tête \texttt{Authorization} des requêtes HTTP, permettant ainsi l’accès aux ressources protégées.
\end{itemize}

\paragraph*{Couche Application (\textit{Backend Tier})}
\begin{itemize}
    \item Le cœur métier est implémenté avec \texttt{Spring Boot 3} selon une architecture MVC enrichie.
    \item Le module \texttt{Spring Security} fonctionne en mode \textit{stateless} et repose sur \texttt{JWT}. Chaque requête traverse la \textit{Security Filter Chain}, qui valide le \textit{token} et applique les règles \textbf{RBAC} (\textit{Role-Based Access Control}).
    \item Les contrôleurs REST exposent les \textit{endpoints} de l’API et manipulent des DTOs afin d’assurer une séparation claire entre les couches. La logique métier est encapsulée dans des services tels que \texttt{AuthService} et \texttt{UserService}.
    \item La sécurité des données est renforcée par l’utilisation de \texttt{BCrypt} pour le hachage des mots de passe. L’accès à la base de données est géré via \texttt{Spring Data JPA}, réduisant le risque d’injection SQL grâce à l’utilisation de requêtes paramétrées.
\end{itemize}

\paragraph*{Couche Données (\textit{Data Tier})}
La persistance des données est assurée par \texttt{PostgreSQL 16}, choisi pour sa robustesse, sa fiabilité et sa gestion avancée des transactions. Le schéma du Sprint 1 se concentre principalement sur les entités liées à la gestion des utilisateurs, des rôles et des mécanismes d’authentification.
\paragraph*{Services externes (SMTP)}
Le \textit{backend} communique avec un service \texttt{SMTP} externe (par exemple \texttt{Gmail SMTP}) afin de gérer l’envoi d’e-mails transactionnels, notamment pour la réinitialisation des mots de passe. Ce mécanisme améliore l’autonomie des utilisateurs et renforce la sécurité du processus de récupération de compte.

\subsubsection*{3. Modèle de données }
Le diagramme \ref{fig:schema_donnees_s1} présente le modèle de données du Sprint 1, couvrant les entités nécessaires à la gestion des utilisateurs, des profils, des rôles et de la sécurité d'authentification.

\begin{figure}[H]
    \centering
    % Remplacez 'chemin_schema_donnees.png' par le nom exact de votre image
    \includegraphics[width=1\textwidth]{images/ERDsprint1.png}
    \vspace{2mm}
    \caption{Diagramme entité-association et schéma de données du Sprint 1}
    \label{fig:schema_donnees_s1}
\end{figure}

Le modèle physique de données (ERD) est composé des entités suivantes :
\begin{itemize}
    \item \textbf{Table \texttt{ROLE}} : Stocke les rôles de sécurité applicatifs pour le contrôle d'accès (RBAC).
    \\ \textit{Points clés} : Unicité stricte du libellé via le champ \texttt{name} (UK) pour interdire les doublons sur \texttt{ROLE\_USER} et \texttt{ROLE\_ADMIN}.
    
    \item \textbf{Table \texttt{USER}} : Gère l'authentification globale et l'activation logique des comptes.
    \\ \textit{Points clés} : Identifiant de connexion unique (\texttt{email} UK), mot de passe sécurisé (\texttt{password\_hash} haché par BCrypt), contact (\texttt{phone}) et statut du compte (\texttt{is\_active} pour le blocage et la conformité RGPD).
    
    \item \textbf{Table \texttt{PROFILE}} : Contient les métadonnées professionnelles et les liens d'intégration (LinkedIn, GitHub, CV).
    \\ \textit{Points clés} : Clé étrangère unique \texttt{user\_id} (FK, UK) garantissant une relation de composition exclusive 1-to-1 (un profil par utilisateur).
    
    \item \textbf{Table \texttt{PASSWORD\_RESET\_TOKEN}} : Sécurise le processus de réinitialisation des mots de passe oubliés.
    \\ \textit{Points clés} : Jeton d'authentification unique (\texttt{token} UK) sous format UUID et sécurité temporelle via la date de fin de validité (\texttt{expiry\_date}).
\end{itemize}


Ce modèle garantit une structure normalisée, évolutive et adaptée aux exigences de sécurité ainsi qu’aux futurs modules d’analyse intelligente.
\subsubsection{4. Hachage des Mots de Passe}

Pour le stockage des mots de passe en base de données, nous avons retenu l'algorithme BCrypt avec un facteur de coût de 10, valeur par défaut recommandée par Spring Security. Ce choix s'appuie sur une analyse comparative de plusieurs standards de hachage, détaillée dans le Tableau \ref{tab:comparaison_hachage} ci-dessous:

\begin{table}[H]
    \centering
    \renewcommand{\arraystretch}{1.3} % Aère les lignes du tableau pour une meilleure lisibilité
    \begin{tabular}{|l|c|c|c|}
        \hline
        \textbf{Critère} & \textbf{BCrypt (coût 10)} & \textbf{MD5} & \textbf{SHA-256} \\
        \hline
        Résistance force brute & Élevée & Faible & Faible \\
        \hline
        Sel automatique (\textit{Salt}) & Oui & Non & Non \\
        \hline
        Adaptatif & Oui & Non & Non \\
        \hline
        Temps de hachage & $\sim$100ms & $<$1ms & $<$1ms \\
        \hline
    \end{tabular}
    \caption{Comparaison des algorithmes de hachage de mots de passe}
    \label{tab:comparaison_hachage}
\end{table}

Cette implémentation garantit un haut niveau de sécurité pour trois raisons fondamentales :
\begin{itemize}
    \item \textbf{Protection contre la force brute :} Le facteur de coût 10 génère environ $2^{10}$ (soit 1024) itérations de hachage. Cela rend les attaques par dictionnaire et par force brute économiquement non viables en termes de temps de calcul pour un attaquant.
    \item \textbf{Résistance aux Rainbow Tables :} Un sel cryptographique (\textit{salt}) est généré automatiquement et de manière unique pour chaque mot de passe lors de sa création, éliminant ainsi les vulnérabilités liées aux tables d'empreintes précalculées.
    \item \textbf{Pérennité (Adaptabilité) :} L'aspect adaptatif de l'algorithme permet d'augmenter le facteur de coût dans les versions futures de la plateforme pour suivre l'évolution de la puissance de calcul matérielle, et ce, sans nécessiter de refonte de l'architecture de sécurité.
\end{itemize}

\subsubsection*{5. Contrats API }

Pour assurer l'interopérabilité entre le \textit{Frontend} \texttt{Angular} et le \textit{Backend} \texttt{Spring Boot}, nous avons défini des contrats d'interface clairs et rigoureux. Voici l'exemple ci-dessous de type du service d'authentification, qui constitue la porte d'entrée sécurisée de l'application :

\vspace{2mm}
\textbf{Point d'accès :} \texttt{POST /api/auth/login}

\textbf{Entrée (\textit{Request Body}) :}
\begin{lstlisting}[language=json, caption={Structure de la requete de connexion}]
{
  "email": "user@talentpredict.exemple",
  "password": "securePassword123"
}
\end{lstlisting}

\textbf{Sortie (\textit{Response} - Success \texttt{200 OK}) :}
\begin{lstlisting}[language=json, caption={Structure de la reponse d'authentification (JWT)}]
{
  {
  "token": "eyJhbGciOiJIUzI1NiIsInR5cCI6IkpXVCJ9...",
  "type": "Bearer",
  "id": "550e8400-e29b-41d4-a716-446655440000",
  "email": "user@talentpredict.exemple",
  "role": "ROLE_USER",
  "nom": "Ben Ali",
  "prenom": "Wajih",
  "redirectUrl": "/dashboard",
  "emailVerified": true
}

}
\end{lstlisting}

\vspace{2mm}
\textbf{Précisions techniques :}
\begin{itemize}
    \item \textbf{\texttt{token}} : Jeton \texttt{JWT} signé par le serveur (\texttt{HS256}) permettant l'accès sécurisé aux ressources protégées.
    \item \textbf{\texttt{type}} : Identifie le schéma d'authentification utilisé (\textit{Bearer Authentication}).
    \item \textbf{\texttt{role}} : Information cruciale permettant au \textit{Frontend} \texttt{Angular} d'adapter dynamiquement l'affichage et de configurer les gardes de navigation (\textit{Route Guards}).
    \item \textbf{\texttt{nom} / \texttt{prenom}} : Données d'identité permettant une personnalisation immédiate de l'interface utilisateur dès la connexion.
    \item \textbf{\texttt{redirectUrl}} : \texttt{URL} de redirection suggérée par le serveur en fonction du profil ou du statut de l'utilisateur.
    \item \textbf{\texttt{emailVerified}} : Indicateur de sécurité permettant au \textit{Frontend} d'inviter l'utilisateur à vérifier son compte si nécessaire.
\end{itemize}


\subsubsection*{6. Gestion des Sessions et Sécurité JWT}
\subsubsection{Architecture d'Authentification \textit{Stateless} et \texttt{JWT}}

L’architecture ci-dessous \ref{fig:auth_flow} de TalentPredict adopte une authentification \textit{Stateless} basée sur les jetons \texttt{JWT} (\textit{JSON Web Token}), complétée par une stratégie de Rotation de Jeton (\textit{Token Rotation}) pour une sécurité maximale.
\begin{figure}[H]
    \centering
    \includegraphics[width=0.88\textwidth,height=0.48\textheight,keepaspectratio]{images/archdesecSprint1.png}
    \vspace{-2mm}
    \caption{Flux de la requête d'authentification et de sécurité via JWT}
    \label{fig:auth_flow}
\end{figure}
\vspace{-4mm}
\paragraph*{Structure et Double Jeton (\textit{Dual Token Strategy})}
Le système ne repose pas sur un seul jeton, mais sur un duo permettant d'équilibrer sécurité et expérience utilisateur :
\begin{itemize}[noitemsep, topsep=2pt]
    \item \textbf{\textit{Access Token} (12 heures) :} Un jeton à durée de vie très courte utilisé pour signer chaque requête \texttt{API}. Il contient les rôles RBAC (\textit{User}/\textit{Admin}).
    \item \textbf{\textit{Refresh Token} (7 jours) :} Un jeton persistant stocké de manière sécurisée, permettant de renouveler l'\textit{Access Token} sans obliger l'utilisateur à se reconnecter manuellement.
\end{itemize}
\paragraph*{Mécanisme \textit{Stateless} et Rotation}
Côté \textit{backend}, aucune session n’est conservée en mémoire vive. Pour chaque requête, le filtre \texttt{JwtAuthFilter} (illustré dans la Figure \ref{fig:auth_flow}) valide la signature (\texttt{HMAC-SHA256}). Nous avons implémenté la \textit{Token Rotation} : chaque demande de nouvel \textit{Access Token} invalide l'ancien \textit{Refresh Token} et en génère un nouveau, empêchant ainsi les attaques par rejeu (\textit{Replay Attacks}).
\paragraph*{Sécurité et RBAC}
Les rôles encodés dans le jeton contrôlent l’accès aux ressources. Toute tentative d’accès non autorisé est rejetée avec une erreur HTTP \texttt{403 Forbidden}. Le stockage côté client est géré de manière à minimiser les risques d'attaques \texttt{XSS}, tandis que les échanges sont protégés par le protocole \texttt{HTTPS}.
\subsubsection*{7. Diagramme de seqence }
\subsubsection{Cas 1 : Cycle de vie de l'Employé}
La figure~\ref{fig:seq_employe_public} représente le diagramme de séquence relatif à l'inscription, la session et la récupération complète du mot de passe pour l'employé.

\begin{figure}[H]
    \centering
    % Remplace le chemin ci-dessous par le nom réel de ton image
    \includegraphics[width=\textwidth]{images/diagramme_3A.png} 
    \caption{Diagramme de séquence : Inscription, Session et Récupération complète (Employé)}
    \label{fig:seq_employe_public}
\end{figure}

Ce diagramme illustre les processus publics accessibles à l'employé avant son authentification. Il se divise en trois phases principales :
\begin{itemize}
    \item \textbf{Phase d'inscription :} L'utilisateur soumet ses informations via l'interface Angular. Le contrôleur délègue la vérification à la couche de service pour s'assurer que l'email ou le numéro de téléphone n'existe pas déjà. Si les données sont valides, le mot de passe est haché à l'aide de BCrypt, le profil est inséré dans la base de données PostgreSQL, et un email de bienvenue est expédié.
    \item \textbf{Phase de récupération de mot de passe :} En cas d'oubli, l'employé demande une réinitialisation. Le système vérifie l'existence du compte, génère un jeton unique (token), l'enregistre en base, et envoie un lien sécurisé par email. Une fois le lien cliqué et le nouveau mot de passe soumis, le système valide le jeton et met à jour le mot de passe haché.
    \item \textbf{Phase de connexion et déconnexion :} Lors de l'authentification, les identifiants sont vérifiés (comparaison BCrypt). En cas de succès, un jeton JWT est généré et stocké localement côté frontend. La déconnexion consiste en la suppression de ce jeton du stockage local.
\end{itemize}

\vspace{0.5cm}

La figure~\ref{fig:seq_employe_jwt} représente le diagramme de séquence détaillant la gestion du profil, des préférences et du mot de passe de l'employé sous authentification sécurisée.

\begin{figure}[H]
    \centering
    % Remplace le chemin ci-dessous par le nom réel de ton image
    \includegraphics[width=\textwidth]{images/diagramme_3B.png} 
    \caption{Diagramme de séquence : Profil, Préférences et Mot de passe (Employé - JWT)}
    \label{fig:seq_employe_jwt}
\end{figure}

Ce diagramme détaille les opérations sécurisées nécessitant une authentification préalable. Toutes les requêtes de ce flux incluent le jeton JWT dans l'en-tête HTTP.
\begin{itemize}
    \item \textbf{Interception de sécurité :} Avant d'atteindre le contrôleur, chaque requête traverse le \texttt{SecurityFilter} qui valide la signature et l'expiration du jeton JWT.
    \item \textbf{Gestion du Profil et Préférences :} Une fois le jeton validé, l'employé peut mettre à jour ses informations personnelles ou ses préférences de notification. Le contrôleur transmet les données au \texttt{ProfileService} qui exécute les requêtes \texttt{UPDATE} correspondantes en base de données.
    \item \textbf{Changement de mot de passe interne :} L'employé saisit son ancien et son nouveau mot de passe. Le service récupère le mot de passe actuel en base de données, le compare avec l'ancien mot de passe saisi via BCrypt. Si la correspondance est validée, le nouveau mot de passe est haché et enregistré.
\end{itemize}

\subsubsection{Cas 2 : Gouvernance et Accès de l'Administrateur RH}

La figure~\ref{fig:seq_admin_public} représente le diagramme de séquence relatif à l'authentification, la gouvernance globale et la récupération du mot de passe pour l'administrateur RH.

\begin{figure}[H]
    \centering
    % Remplace le chemin ci-dessous par le nom réel de ton image
    \includegraphics[width=\textwidth]{images/diagramme_4A.png} 
    \caption{Diagramme de séquence : Authentification, Gouvernance et Récupération complète (Admin RH)}
    \label{fig:seq_admin_public}
\end{figure}

Ce diagramme modélise l'accès et les actions globales de l'Administrateur RH, en mettant l'accent sur les restrictions de rôles.
\begin{itemize}
    \item \textbf{Récupération et Connexion (Phases publiques) :} Le processus de mot de passe oublié est similaire à celui de l'employé. Cependant, lors de la phase de connexion, la requête SQL vérifie non seulement la validité des identifiants (via BCrypt), mais s'assure également que l'utilisateur possède explicitement le rôle \texttt{ROLE\_ADMIN}. Si ce rôle est absent, une exception \texttt{BadCredentialsException} est levée.
    \item \textbf{Gouvernance et Tableau de bord (Phase sécurisée) :} Pour accéder aux statistiques globales, la requête GET est interceptée par le \texttt{SecurityFilter} qui extrait le JWT et valide la présence du rôle Administrateur. En cas d'autorisation, l'\texttt{AdminService} exécute les requêtes d'agrégation (comme le comptage des utilisateurs) sur la base PostgreSQL et retourne un objet DTO formaté en JSON pour peupler le tableau de bord côté Angular.
\end{itemize}

\vspace{0.5cm}

La figure~\ref{fig:seq_admin_jwt} représente le diagramme de séquence illustrant la mise à jour des informations organisationnelles et des paramètres de sécurité de l'administrateur RH.

\begin{figure}[H]
    \centering
    % Remplace le chemin ci-dessous par le nom réel de ton image
    \includegraphics[width=\textwidth]{images/diagramme_4B.png} 
    \caption{Diagramme de séquence : Profil, Préférences et Mot de passe (Admin RH - JWT)}
    \label{fig:seq_admin_jwt}
\end{figure}

Ce diagramme présente la gestion du compte personnel de l'Administrateur RH.
\begin{itemize}
    \item \textbf{Mécanisme :} Le flux est techniquement identique à celui de l'employé (Diagramme 3.B), reposant sur l'en-tête d'autorisation Bearer et la validation du JWT par le \texttt{SecurityFilter}.
    \item \textbf{Opérations :} Il couvre la mise à jour des informations organisationnelles de l'administrateur, la gestion de ses propres préférences de sécurité, ainsi que la procédure stricte de modification de son mot de passe nécessitant la vérification préalable de l'ancien mot de passe haché.
\end{itemize}

\subsubsection*{8. Diagramme de Classes }
Le diagramme de classes \ref{fig:class_diagram_sprint1} illustre l'architecture logicielle du module d'authentification et de gestion de profil, mis en place lors du premier sprint du projet TalentPredict. Il s'appuie sur une conception robuste en couches (N-tiers), caractéristique des applications modernes basées sur Spring Boot. Cette approche architecturale garantit une séparation stricte des responsabilités (\textit{Separation of Concerns}), facilitant ainsi la maintenance du code, l'évolutivité du système et l'application rigoureuse des principes de \textit{Security-by-Design}.
\begin{figure}[H]
    \centering
    % Remplacez 'chemin_diagramme_classes_sprint1.png' par le nom exact de votre image
    \includegraphics[width=\textwidth]{images/classSprint1.png.png}
    \vspace{2mm}
    \caption{Diagramme de classes : Architecture en couches du Sprint 1}
    \label{fig:class_diagram_sprint1}
\end{figure}

'architecture est structurée autour de six composants ou couches principales, interagissant de manière unidirectionnelle pour traiter les requêtes des utilisateurs :

\begin{itemize}
    \item \textbf{Couche Modèle (Entités) :} Constitue le cœur du domaine de l'application. Elle regroupe les objets persistants mappés directement aux tables de la base de données (via JPA/Hibernate), tels que \texttt{User}, \texttt{Profile}, \texttt{PasswordResetToken} et \texttt{UserPrivacySettings}. Les énumérations, comme \texttt{Role} (USER, ADMIN), y sont également définies pour normaliser les données.
    
    \item \textbf{Couche Persistance (Repositories) :} Agit comme l'interface de communication avec la base de données. En étendant les interfaces de Spring Data JPA (comme \texttt{UserRepository} ou \texttt{ProfileRepository}), cette couche encapsule la logique d'accès aux données (opérations CRUD) et permet de récupérer ou manipuler les entités sans exposer les requêtes SQL complexes aux couches supérieures.
    
    \item \textbf{Couche DTOs (Data Transfer Objects) :} Joue un rôle de bouclier de protection et de formatage. Les objets tels que \texttt{LoginRequest}, \texttt{RegisterRequest} ou \texttt{AuthResponse} sont utilisés pour transporter strictement les données nécessaires entre le client (frontend Angular) et le serveur, évitant ainsi d'exposer directement la structure interne des entités de la base de données.
    
    \item \textbf{Couche Sécurité (Spring Security) :} Centralise les mécanismes de protection des ressources. Elle intègre les composants essentiels pour une authentification \textit{stateless}, notamment \texttt{JwtService} pour la génération et la validation des tokens JWT, et \texttt{JwtAuthenticationFilter} qui intercepte chaque requête entrante pour vérifier les habilitations de l'utilisateur avant d'autoriser l'accès aux API.
    
    \item \textbf{Couche Métier (Services) :} Contient la logique fonctionnelle et les règles métier de l'application. Les classes comme \texttt{AuthServiceImpl} et \texttt{ProfileService} orchestrent les opérations complexes (inscription, hachage des mots de passe, envoi d'e-mails de vérification, gestion des tokens de réinitialisation) en faisant le lien entre les contrôleurs et les repositories.
    
    \item \textbf{Couche API (Controllers) :} Représente le point d'entrée du backend (API REST). Les contrôleurs (\texttt{AuthController}, \texttt{ProfileController}) interceptent les requêtes HTTP (GET, POST, PUT), délèguent le traitement à la couche métier appropriée, puis formatent la réponse finale sous forme de \texttt{ResponseEntity} pour la renvoyer au client.
\end{itemize}




\subsection{Codage}

Cette sous-section détaille le passage de la conception à l'implémentation logicielle. Pour ce premier sprint, l'objectif principal était de bâtir un socle technologique résilient, sécurisé et prêt à monter en charge (\textit{Scalable}) pour accueillir les futurs modules d'intelligence artificielle.

\paragraph*{1. Écosystème Technologique (\textit{Stack} Technique)}
La sélection des technologies s'est portée sur des standards de l'industrie offrant un haut niveau de sécurité et de maintenabilité :

\begin{itemize}
    \item \textbf{Couche Métier et Sécurité (\textit{Backend}) :} Propulsée par \texttt{Java 17} et le \textit{framework} \texttt{Spring Boot 3.4.2}. La persistance des entités est orchestrée par \texttt{Spring Data JPA} (\texttt{Hibernate}). La sécurisation des \textit{endpoints} repose sur \texttt{Spring Security} couplé à la librairie \texttt{JJWT} pour la cryptographie des jetons (\texttt{HMAC-SHA256}).
    \item \textbf{Couche Présentation (\textit{Frontend}) :} Implémentée avec \texttt{Angular 17} (\texttt{TypeScript}). Nous avons privilégié une architecture réactive (\texttt{RxJS}) avec une interface utilisateur adaptative basée sur des composants UI modernes.
    \item \textbf{Persistance Relationnelle :} \texttt{PostgreSQL 16}, sélectionné pour sa robustesse transactionnelle (ACID) et son extensibilité (capacités futures de recherche vectorielle ou \texttt{JSONB} pour l'IA).
\end{itemize}

\subsubsection{2. Configuration du Service de Messagerie (Gmail SMTP)}

Pour assurer la fiabilité des communications transactionnelles (activation de compte, réinitialisation de mot de passe, alertes système), TalentPredict intègre le service \texttt{Gmail SMTP}. La configuration s'appuie sur le serveur \texttt{smtp.gmail.com} tilisant le protocole sécurisé \texttt{STARTTLS} sur le port \texttt{587}.

\begin{lstlisting}[language=properties, caption=Configuration SMTP dans le fichier application.properties]
spring.mail.host=smtp.gmail.com
spring.mail.port=587
spring.mail.username=${MAIL_USERNAME}
spring.mail.password=${MAIL_PASSWORD}
spring.mail.properties.mail.smtp.auth=true
spring.mail.properties.mail.smtp.starttls.enable=true
\end{lstlisting}

Les identifiants sont externalisés dans le fichier \texttt{.env} pour respecter la sécurité des secrets et éviter toute compromission lors du versionnage. L'envoi des e-mails est configuré en mode asynchrone (via l'annotation \texttt{@Async}) afin de ne pas bloquer le \textit{thread} principal lors des requêtes utilisateur, garantissant ainsi un temps de réponse optimal de l'API. Enfin, un \textit{App Password} (mot de passe d'application) Google est utilisé à la place du mot de passe principal du compte, conformément aux recommandations strictes de sécurité.

\subsubsection*{3. Organisation Modulaire et \textit{Clean Architecture}}
Le projet rejette l'approche monolithique classique au profit d'un \textit{Monolithe Modulaire}. Le code \textit{backend} est scindé en domaines fonctionnels stricts, par exemple \texttt{modules.auth}. Chaque module respecte le principe de ségrégation des responsabilités (SRP) :

\begin{itemize}
    \item \textbf{\texttt{Controllers} :} Exclusivement dédiés au routage HTTP et à l'interception des requêtes.
    \item \textbf{\texttt{Services} :} Encapsulent la logique métier complexe (algorithmes, orchestration).
    \item \textbf{\texttt{Repositories} :} Couche d'abstraction (Interfaces) pour l'accès aux données.
    \item \textbf{\texttt{DTOs} (\textit{Data Transfer Objects}) :} Objets immuables servant de contrat entre le \textit{Frontend} et le \textit{Backend}, empêchant les failles de sécurité de type \textit{Mass Assignment}.
\end{itemize}

Côté \texttt{Angular}, l'arborescence suit le \textit{pattern Core/Shared/Feature} :
\begin{itemize}
    \item Le dossier \texttt{core/} centralise les singletons (Intercepteurs HTTP pour injecter le \texttt{JWT}, \textit{Guards} pour sécuriser le routage).
    \item Le dossier \texttt{modules/} isole les composants (ex. : \textit{Login}, \textit{Dashboard}).
\end{itemize}

\subsubsection*{5. Fonctionnalités Implémentées}
\begin{itemize}
    \item \textbf{Authentification \textit{Stateless} (\texttt{JWT}) :} Mise en place d'un flux de connexion où l'état de session n'est pas conservé sur le serveur. Les mots de passe sont hachés via \texttt{BCrypt} (avec ajout d'un "sel" cryptographique) pour contrer les attaques par dictionnaire.
    \item \textbf{Contrôle d'Accès Basé sur les Rôles (RBAC) :} Configuration de la sécurité bloquant l'accès aux routes d'administration (\texttt{/api/admin/**}) pour tout utilisateur ne possédant pas le \texttt{ROLE\_ADMIN} (renvoi d'une erreur HTTP \texttt{403}).
    \item \textbf{Collecte des Vecteurs IA (Profil) :} Développement des interfaces permettant aux employés de lier dynamiquement leurs URLs professionnelles (\texttt{GitHub}, \texttt{LinkedIn}, \texttt{Portfolio}) et de téléverser leur CV, données indispensables pour l'extraction de compétences prévue au Sprint 2.
\end{itemize}

\subsubsection*{6. Ingénierie DevOps et Conteneurisation}
Afin de garantir une parité absolue entre les environnements de développement et de production, une stratégie d’\textit{Infrastructure as Code} (IaC) a été adoptée dès l'initialisation du projet.

\subsubsection*{A. Orchestration de la Base de Données avec Docker}
Le SGBD PostgreSQL est instancié via un fichier \texttt{docker-compose.yml}. Ce dernier définit les limites de ressources, les réseaux isolés et surtout le montage des volumes locaux (\texttt{volumes}) pour assurer la persistance des données. Cette approche garantit qu'en cas de destruction ou de redémarrage du conteneur, aucune donnée utilisateur n'est perdue.

\subsubsection{B. Sécurisation des Configurations via le fichier \texttt{.env}}
Conformément aux meilleures pratiques de sécurité, aucun secret n'est stocké en dur dans le code source ou versionné sur Git. L'application consomme ses paramètres sensibles (identifiants de base de données, clés cryptographiques, configuration SMTP) via un fichier d'environnement local \texttt{.env}.

Voici un aperçu de la structure du fichier \texttt{.env} utilisé pour notre backend :

\begin{lstlisting}[language=bash, caption={Structure du fichier de configuration sécurisé (.env)}, label={lst:env_file}]
# --- DATABASE CONFIGURATION ---
DB_URL=jdbc:postgresql://localhost:5432/talentpredict_db
DB_USERNAME=tp_admin
DB_PASSWORD=******** # Secret masque

# --- SECURITY (JWT) ---
JWT_SECRET=4a2b6c8d9e0f1g2h3i4j5k6l7m8n9o0p1q2r3s4t5u6v7w8x9y0z
JWT_EXPIRATION=86400000  # 24 Heures

# --- MAIL SERVER (SMTP) ---
SMTP_HOST=smtp.gmail.com
SMTP_PORT=587
SMTP_USERNAME=noreply.talentpredict@gmail.com
SMTP_PASSWORD=app_password_security
\end{lstlisting}

Cette approche repose sur deux piliers fondamentaux de l'ingénierie logicielle :
\begin{itemize}
    \item \textbf{Séparation stricte du Code et de la Configuration :} Le code source reste identique quel que soit l'environnement (Développement, Staging, Production). Seul le fichier \texttt{.env} diffère, ce qui facilite grandement le déploiement continu.
    \item \textbf{Sécurité et Confidentialité :} L'inclusion du fichier \texttt{.env} dans le fichier \texttt{.gitignore} prévient toute fuite accidentelle d'identifiants sur des dépôts distants (GitHub/GitLab), protégeant ainsi l'infrastructure contre les accès non autorisés.
\end{itemize}

\subsubsection*{C. Sécurisation des Échanges via CORS}
Enfin, pour protéger l'API REST contre les attaques de type \textit{Cross-Site Request Forgery} (CSRF), un paramétrage strict du filtre CORS (\textit{Cross-Origin Resource Sharing}) a été appliqué au niveau de Spring Boot. Ce filtre n'autorise que les requêtes provenant de l'origine de confiance de notre application front-end (\texttt{http://localhost:4200}), bloquant systématiquement toute tentative d'interaction provenant de domaines non reconnus.

\subsubsection*{7. Assurance Qualité et Validation des Contrats API}
Avant toute connexion au client \texttt{Angular}, les APIs ont été éprouvées selon des méthodes d'ingénierie rigoureuses :

\begin{itemize}
    \item \textbf{Documentation OpenAPI (\texttt{Swagger 3}) :} Intégration de \texttt{springdoc-openapi} pour exposer dynamiquement un portail interactif listant les schémas de requêtes/réponses et les codes HTTP attendus.

\begin{figure}[H]
    \centering
    % Remplacez 'chemin_swagger_ui.png' par le nom exact de votre image
    \includegraphics[width=0.8\textwidth]{images/swager.png}
    \vspace{2mm}
    \caption{Documentation interactive des APIs REST via Swagger UI}
    \label{fig:swagger_ui_s1}
\end{figure}

La documentation interactive des APIs REST a été générée automatiquement grâce à l’intégration de \texttt{Springdoc OpenAPI} (\texttt{Swagger 3}). Cette interface permet de visualiser dynamiquement l’ensemble des \textit{endpoints} exposés par le \textit{backend} \texttt{Spring Boot}, leurs paramètres d’entrée, les schémas \texttt{JSON} attendus ainsi que les codes de réponses HTTP possibles.

\texttt{Swagger UI} a également permis de tester directement les routes sécurisées nécessitant un jeton \texttt{JWT} en injectant le \textit{token} dans le \textit{header} \texttt{Authorization}. Cette approche facilite considérablement les phases de validation, de débogage et d’intégration entre le \textit{frontend} \texttt{Angular} et le \textit{backend} REST.

L’utilisation d’\texttt{OpenAPI} améliore également la maintenabilité du projet en fournissant une documentation technique centralisée et synchronisée automatiquement avec le code source.


    \item \textbf{Tests Automatisés (\texttt{Postman}) :} Création de collections \texttt{Postman} intégrant des scripts de pré-requête (\textit{Pre-request scripts}). Ces scripts capturent automatiquement le jeton \texttt{JWT} renvoyé par la route \texttt{/login} et le placent dynamiquement dans le \textit{header} \texttt{Authorization: Bearer} pour valider les routes protégées, simulant un comportement utilisateur réel.
\end{itemize}
\subsection{Realisation}
ette section présente les interfaces majeures développées pour assurer l'interaction avec le socle de sécurité et de gestion des profils.

\subsubsection*{1. La Page d'Accueil (Landing Page)}
L’interface de \textit{Page d'Accueil} \ref{fig:ui_landing_s1} constitue la vitrine technologique de TalentPredict. Son rôle est de présenter la proposition de valeur de la plateforme : l'utilisation de l'intelligence artificielle pour l'évaluation prédictive des talents. Conçue avec une approche \textit{Mobile-First}, elle offre une navigation fluide vers les modules d'authentification. Elle utilise des composants visuels modernes (dégradés, animations légères) pour instaurer une image de marque innovante et professionnelle dès le premier contact avec l'utilisateur.

\begin{figure}[H]
    \centering
    \includegraphics[width=\textwidth]{images/landingpage.png}
    \vspace{2mm}
    \caption{Interface de la Landing Page de TalentPredict}
    \label{fig:ui_landing_s1}
\end{figure}

\subsubsection*{2. Interface d'Inscription (Register)}

L'interface d'inscription \ref{fig:ui_register_s1} est le point d'entrée pour les nouveaux employés. Elle repose sur un formulaire intelligent utilisant les \texttt{Angular Reactive Forms}, permettant une validation instantanée des données saisies (unicité de l'email, robustesse du mot de passe). Une fois le formulaire soumis, l'application communique avec l'\textit{endpoint} \texttt{/api/auth/register} du \textit{backend} \texttt{Spring Boot}, qui procède à la création sécurisée du compte et à l'initialisation automatique d'un profil vide, prêt à être enrichi par le candidat.

\begin{figure}[H]
    \centering
    \includegraphics[width=\textwidth]{images/signin.png}
    \vspace{2mm}
    \caption{Interface d'inscription des nouveaux employés}
    \label{fig:ui_register_s1}
\end{figure}

\subsubsection*{3. Interface de Connexion (Login)}

Cette interface \ref{fig:ui_login_s1} centralise l'accès sécurisé pour l'ensemble des utilisateurs (Employés et RH). Elle a été conçue pour être minimaliste et efficace, incluant des mécanismes de gestion d'erreurs clairs en cas d'identifiants incorrects. C'est ici que s'opère le flux critique d'authentification : en échange des identifiants valides, le serveur délivre un \textit{Token} \texttt{JWT}. Ce jeton est alors intercepté et stocké par le \textit{frontend} pour autoriser l'accès aux zones protégées de l'application selon le rôle de l'utilisateur (RBAC). Elle intègre également le lien vers le flux de récupération de compte en cas d'oubli du mot de passe.

\begin{figure}[H]
    \centering
    \includegraphics[width=\textwidth]{images/login.png}
    \vspace{2mm}
    \caption{Interface de connexion et gestion d'accès}
    \label{fig:ui_login_s1}
\end{figure}
\subsubsection*{4. Le Flux de Récupération de Compte (Self-Service Password Reset)}

Afin de garantir une autonomie totale des utilisateurs et une sécurité optimale, nous avons implémenté un flux de récupération de mot de passe complet et sécurisé. Ce mécanisme repose sur trois étapes clés :

\begin{itemize}
    \item \textbf{A. Demande de Récupération :} L'utilisateur saisit son adresse email sur une interface dédiée. Le système vérifie l'existence du compte et génère un jeton de réinitialisation unique (\textit{Token} \texttt{UUID}) à durée de vie limitée (7 heures). Cette étape prévient les tentatives d'usurpation grâce à un canal de communication chiffré.
    
    \begin{figure}[H]
        \centering
        \includegraphics[width=\textwidth]{images/forgetpassword.png}
        \caption{Interface de demande de récupération de compte}
        \label{fig:ui_forget_password}
    \end{figure}

    \item  \textbf{B. Email Transactionnel :} Un email professionnel est instantanément envoyé à l'utilisateur via un service \texttt{SMTP}. Cet email contient un lien sécurisé pointant vers l'application \textit{frontend}. Comme l'illustre la capture de la boîte de réception , le message fournit les instructions nécessaires pour poursuivre la démarche en toute confiance.

    \begin{figure}[H]
        \centering
        \includegraphics[width=\textwidth]{images/mail.png}
        \caption{Exemple d'email transactionnel reçu par l'utilisateur}
        \label{fig:ui_mail_reset}
    \end{figure}

    \item \textbf{C. Réinitialisation du Mot de Passe :} En cliquant sur le lien, l'utilisateur est redirigé vers une interface "Etape Finale" . Cette page valide le jeton présent dans l'URL et permet de définir un nouveau mot de passe robuste. L'interface bénéficie d'un design moderne et cohérent avec l'identité visuelle de TalentPredict.

    \begin{figure}[H]
        \centering
        \includegraphics[width=\textwidth]{images/modifierpassword.png}
        \caption{Interface de définition du nouveau mot de passe}
        \label{fig:ui_modify_password}
    \end{figure}
\end{itemize}
\subsubsection*{5. Contrôle d'Accès Dynamique et Navigation Différenciée}

L'une des réussites techniques majeures de ce sprint réside dans l'implémentation d'une navigation adaptative pilotée par le contexte d'authentification. Grâce à l'extraction des \textit{claims} (revendications) du jeton JWT (JSON Web Token), l'interface utilisateur (UI) se reconfigure dynamiquement en temps réel pour ne présenter que les fonctionnalités autorisées selon le rôle de l'utilisateur.
\begin{figure}[H]
    \centering
    \begin{minipage}{0.2\textwidth}
        \centering
        \includegraphics[width=\textwidth]{images/navbarUser.png} % Remplacez par le nom de l'Image 6
    \end{minipage}\hfill
    \begin{minipage}{0.2\textwidth}
        \centering
        \includegraphics[width=\textwidth]{images/navbarAdmin.png} % Remplacez par le nom de l'Image 7
    \end{minipage}
    \vspace{1mm}
    \caption{Interface de navigation côté administrateur et côté utilisateur}
    \label{fig:avigation côté administrateur et côté utilisateur}
\end{figure}

\subsubsection*{A. Menu Employé (Vue Employé)}
La barre de navigation destinée à l'employé est centrée sur son parcours d'évaluation et de montée en compétences :
\begin{itemize}
    \item \textbf{Tests :} Accès aux évaluations des \textit{Compétences comportementales} et \textit{Compétences techniques}.
    \item \textbf{Résultats :} Consultation des scores et suivi de la progression individuelle.
    \item \textbf{Upskilling :} Accès direct aux recommandations de formation personnalisées.
\end{itemize}

\subsubsection*{B. Menu Administrateur (Vue RH / Manager)}
À l'opposé, l'interface administrateur propose des outils de pilotage stratégique et de gouvernance globale :
\begin{itemize}
    \item \textbf{Effectifs :} Gestion de la base de données des collaborateurs.
    \item \textbf{Plan de Formation :} Supervision du déploiement des modules d'apprentissage.
    \item \textbf{Campagnes :} Gestion des campagnes de communication et de recrutement interne.
\end{itemize}

\subsubsection*{6.Gestion du Profil et Préférences de Sécurité}

Ces deux interfaces représentent le centre de contrôle de l'utilisateur au sein de la plateforme TalentPredict. Elles allient la collecte de données métier à la protection de la vie privée.

\subsubsection*{A. Interface "Mon Profil"}
Cette interface \ref{fig:ui_profile_user} est conçue comme un tableau de bord personnel permettant à l'employé de valoriser ses compétences. Elle centralise les vecteurs de données critiques qui seront analysés par le moteur d'IA lors du Sprint 2 :
\begin{figure}[htbp]
    \centering
    \includegraphics[width=0.85\textwidth]{images/profileUser.png}
    \caption{Interface de gestion du profil employé }
    \label{fig:ui_profile_user}
\end{figure}

\begin{itemize}
    \item \textbf{Données professionnelles :} Nombre de dépôts \texttt{GitHub}, lien vers le profil \texttt{LinkedIn} ,ien vers le profil \texttt{Portfolio} et téléversement du \texttt{CV} numérique.
    \item \textbf{Indicateurs d'engagement :} L'intégration d'une jauge de complétion et de conseils personnalisés sur les champs manquants garantit la qualité des données collectées pour le système.
\end{itemize}
\subsubsection*{B. Interface "Sécurité \& Confidentialité"}
Cette interface \ref{fig:ui_security_user} assure la confiance de l'utilisateur envers la plateforme en respectant les standards de \textit{Security by Design}. Elle offre une transparence totale sur l'utilisation du compte :
\begin{figure}[H]
    \centering
    \includegraphics[width=0.85\textwidth]{images/sécuritéUser.png}
    \caption{Interface de sécurité, historique de connexion et paramètres RGPD}
    \label{fig:ui_security_user}
\end{figure}

\begin{itemize}
    \item \textbf{Gestion des accès :} Fonctionnalités de modification sécurisée du mot de passe.
    \item \textbf{Traçabilité :} Historique de connexion détaillé incluant la date, l'adresse IP et le type d'appareil utilisé.
    \item \textbf{Conformité :} Un volet dédié au \textbf{RGPD} permet à l'utilisateur de gérer son consentement et ses données personnelles.
\end{itemize}


\subsubsection{6. Gestion du Profil et Sécurité Administrateur}

Dans l'optique de garantir une gestion rigoureuse des comptes à hauts privilèges, la plateforme \textit{TalentPredict} intègre des interfaces de gestion de profil spécifiquement configurées pour les utilisateurs possédant le rôle \texttt{ROLE\_ADMIN}.

\subsubsection{A. Profil Professionnel de l'Administrateur}
Cette interface \ref{fig:profil_admin} permet au gestionnaire RH de maintenir ses informations professionnelles et son expertise. Bien que la cohérence visuelle soit maintenue avec l'interface employé, cet espace est strictement isolé pour permettre une identification claire des responsables au sein de l'organisation :
\begin{figure}[H]
    \centering
    \includegraphics[width=0.7\textwidth]{images/profileAdmin.png}
    \caption{Interface de gestion du profil administrateur}
    \label{fig:profil_admin}
\end{figure}
\begin{itemize}
    \item \textbf{Informations Générales :} Mise à jour du titre professionnel, des années d'expérience et du niveau d'études.
    \item \textbf{Identité Numérique :} Intégration de liens vers les réseaux professionnels (LinkedIn, GitHub) et téléchargement de documents (CV, Photo de profil).
    \item \textbf{Bio \& Parcours :} Espace libre pour décrire la vision du recrutement et le parcours de l'administrateur.
\end{itemize}
\Needspace{18\baselineskip}

\subsubsection{B. Sécurité du Compte Privilégié}
Compte tenu des droits étendus de l'administrateur sur les données sensibles du personnel, l'interface de sécurité \ref{fig:security_admin} impose des contrôles rigoureux :
\begin{figure}[H]
    \centering
    \includegraphics[width=0.8\textwidth]{images/modifermdpAdmin.png}
    \caption{Interface de modification de mot de passe administrateur}
    \label{fig:security_admin}
\end{figure}
\begin{itemize}
    \item \textbf{Renouvellement des Identifiants :} Formulaire dédié imposant la saisie du mot de passe actuel pour toute modification.
    \item \textbf{Standards de Sécurité :} Le système exige un nouveau mot de passe robuste (minimum 8 caractères) pour prévenir les tentatives d'usurpation d'identité.
    \item \textbf{Confidentialité :} Accès rapide aux paramètres de visibilité des informations du compte.
\end{itemize}

\Needspace{18\baselineskip}

\subsubsection*{7. Processus de Déconnexion Sécurisée}

La clôture d'une session sur la plateforme \textit{TalentPredict} a été conçue pour allier confort d'utilisation et sécurité des données. Ce processus marque la fin du cycle d'activité de l'utilisateur et assure l'invalidité des accès après son départ.

\begin{figure}[H]
    \centering
    \includegraphics[width=0.8\textwidth]{images/logout.png}
    \caption{Interface de confirmation de déconnexion }
    \label{fig:logout_modal}
\end{figure}

\subsection{Phase de Test et Déploiement}

Cette dernière phase du Sprint 1 a pour objectif de s'assurer que les fonctionnalités développées répondent strictement aux exigences fixées, qu'elles sont sécurisées, et que le système est stable avant d'être déployé.

\subsubsection*{1. Stratégie d’Assurance Qualité (QA)}
Pour garantir la fiabilité de la plateforme TalentPredict, nous avons adopté une stratégie de test multi-niveaux :

\begin{itemize}
    \item \textbf{Validation des Règles Métier (\textit{Backend}) :} Utilisation de \texttt{JUnit 5} et \texttt{Mockito} pour tester unitairement les services critiques (notamment le hachage des mots de passe et la génération du jeton \texttt{JWT}).
    \item \textbf{Qualité du Code (Analyse Statique) :} Mise en place du \textit{plugin} Maven \texttt{Checkstyle} (via \texttt{checkstyle.xml}) pour imposer le respect des conventions de codage standard de Java.
    \item \textbf{Tests d'Intégration API :} Utilisation de \texttt{Postman} pour exécuter des requêtes sur tous les \textit{endpoints} REST, validant ainsi les codes de statut HTTP (\texttt{200 OK}, \texttt{201 Created}, \texttt{401 Unauthorized}, \texttt{403 Forbidden}, \texttt{409 Conflict}).
\end{itemize}
\subsubsection*{2. Synthèse des Tests d'Authentification Réalisés}

Pour valider le socle de sécurité du Sprint 1, une série de tests fonctionnels a été exécutée de bout en bout :

\begin{itemize}
    \item \textbf{Inscription (\texttt{/api/auth/register}) :} Validation de la création d'un utilisateur et de son profil associé (\texttt{201 Created}), du hachage sécurisé du mot de passe (\textbf{BCrypt}) et du rejet automatique des doublons d'adresse e-mail (\texttt{409 Conflict}).
    
    \item \textbf{Connexion (\texttt{/api/auth/login}) :} Validation de la vérification des identifiants, de la génération du jeton de session \textbf{JWT} signé, et du blocage des attaques par force brute via un limiteur de débit (\textbf{Bucket4j}) renvoyant un code d'erreur \texttt{429 Too Many Requests}.
    
    \item \textbf{Récupération de mot de passe :} Validation du cycle de réinitialisation sécurisée via la génération d'un jeton temporaire et à usage unique (\texttt{PasswordResetToken}), l'envoi du lien par courrier électronique, et la mise à jour effective du nouveau mot de passe haché.
    
    \item \textbf{Routes protégées et Rôles (RBAC) :} Validation de l'interception des requêtes HTTP par le filtre \texttt{JwtAuthFilter} et les Guards Angular. L'accès sans jeton valide renvoie systématiquement un code \texttt{401 Unauthorized}, tandis qu'une tentative d'accès d'un compte de rôle standard (\texttt{USER}) aux ressources réservées aux administrateurs (\texttt{/admin/**}) est bloquée avec un code de retour \texttt{403 Forbidden}.
\end{itemize}

\subsubsection*{3. Couverture de Code et Résultats des Tests}:

Afin de garantir la fiabilité des modules critiques, nous avons mesuré la couverture de code via le \textit{plugin} \texttt{JaCoCo} intégré à \texttt{Maven}. Les tests unitaires ont été rédigés avec \texttt{JUnit 5} et \texttt{Mockito}, en isolant la logique métier de la couche de persistance. Les résultats obtenus sur les modules de sécurité et d'authentification sont les suivants \ref{tab:coverage_s1} :

\begin{table}[H]
    \centering
    \renewcommand{\arraystretch}{1.5}
    \small
    \begin{tabular}{|l|c|}
        \hline
        \rowcolor[HTML]{F2F2F2}
        \textbf{Module} & \textbf{Couverture} \\
        \hline
        \texttt{AuthService} & $92\%$ \\
        \hline
        \texttt{JwtUtil} & $95\%$ \\
        \hline
        \texttt{ProfileService} & $82\%$ \\
        \hline
        \texttt{GlobalExceptionHandler} & $88\%$ \\
        \hline
        \rowcolor[HTML]{EAEAEA}
        \textbf{Total modules critiques} & \textbf{85\%} \\
        \hline
    \end{tabular}
    \vspace{2mm}
    \caption{Résultats de la couverture de code par module (Sprint 1)}
    \label{tab:coverage_s1}
\end{table}

Ces résultats confirment la robustesse du socle de sécurité livré lors de ce premier sprint.

\subsubsection*{4. Plan de Validation Fonctionnelle (Tests \textit{End-to-End})}
Le tableau \ref{tab:validation_sprint1} suivant résume les principaux scénarios de tests exécutés pour valider les flux de bout en bout de l'application :

\begin{table}[H]
    \centering
    \renewcommand{\arraystretch}{1.5}
    \small
    \begin{tabularx}{\textwidth}{|c|>{\raggedright\arraybackslash}p{3cm}|X|X|}
        \hline
        \rowcolor[HTML]{F2F2F2}
        \textbf{ID} & \textbf{Scénario Testé} & \textbf{Pré-conditions \& Étapes} & \textbf{Résultat Attendu} \\
        \hline
        TC-01 & Inscription d'un nouvel Employé & \textbf{Pré-cond. :} Utilisateur sur la \textit{Landing Page}. \newline \textbf{Étapes :} Saisir nom, email, mot de passe $\rightarrow$ Soumettre. & Redirection vers la page de \textit{login}. Code HTTP \texttt{201}. Création d'un \texttt{User} et d'un \texttt{Profile} vide en base. \\
        \hline
        TC-02 & Gestion des doublons (Inscription) & \textbf{Pré-cond. :} L'email "\texttt{test@tp.com}" existe déjà. \newline \textbf{Étapes :} Tenter une inscription avec cet email. & Le système bloque l'inscription. Affichage d'un message d'erreur. Code HTTP \texttt{409 Conflict}. \\
        \hline
        TC-03 & Connexion \& Contrôle d'Accès (RBAC) & \textbf{Pré-cond. :} Posséder un compte Employé valide. \newline \textbf{Étapes :} Se connecter $\rightarrow$ Tenter d'accéder à l'URL \texttt{/admin/dashboard}. & La connexion réussit (\textit{Token} \texttt{JWT} généré), mais l'accès au Panel RH est bloqué (Erreur \texttt{403 Forbidden}). \\
        \hline
        TC-04 & Enrichissement du profil (Vecteurs IA) & \textbf{Pré-cond. :} L'employé est connecté (\textit{Token} actif). \newline \textbf{Étapes :} Accéder au profil $\rightarrow$ Saisir lien \texttt{GitHub} $\rightarrow$ Enregistrer. & Mise à jour réussie en base de données. Message de succès côté interface (Code HTTP \texttt{200}). \\
        \hline
    \end{tabularx}
    \caption{Plan de validation fonctionnelle du Sprint 1}
    \label{tab:validation_sprint1}
\end{table}

\subsubsection*{5. Résultats obtenus et Corrections apportées (Débogage)} Au cours de l'exécution des tests, plusieurs défis techniques ont été identifiés et résolus, renforçant ainsi la robustesse du système :

\begin{itemize}
    \item \textbf{Erreur de Politique CORS (\textit{Cross-Origin Resource Sharing}) :}
    \begin{itemize}
        \item \textit{Problème :} Le navigateur bloquait les requêtes d'\texttt{Angular} (port 4200) vers \texttt{Spring Boot} (port 8081).
        \item \textit{Correction :} Implémentation d'une classe de configuration \texttt{WebMvcConfigurer} dans \texttt{Spring Boot} pour autoriser explicitement les requêtes, les en-têtes (notamment \texttt{Authorization}), et les méthodes HTTP provenant de l'interface \textit{front-end}.
    \end{itemize}
    \vspace{2mm}
    \item \textbf{Perte de contexte de sécurité (Filtre \texttt{JWT}) :}
    \begin{itemize}
        \item \textit{Problème :} Les rôles de l'utilisateur n'étaient pas reconnus après la connexion.
        \item \textit{Correction :} Ajout de l'extraction des \textit{Claims} (Rôles) directement depuis le \textit{Payload} du \textit{token} \texttt{JWT} dans la classe \texttt{JwtAuthFilter} et injection dans le \texttt{SecurityContextHolder} de \texttt{Spring Security}.
    \end{itemize}
    \vspace{2mm}
    \item \textbf{Problème de communication inter-conteneurs :}
    \begin{itemize}
        \item \textit{Problème :} Le \textit{backend} n'arrivait pas à se connecter à \texttt{PostgreSQL} lors du lancement \texttt{Docker}.
        \item \textit{Correction :} Ajustement de la configuration Docker pour pointer vers le nom du service (\texttt{postgres:5432}) au lieu de \texttt{localhost}.
    \end{itemize}
\end{itemize}

\subsubsection*{6. Renforcement de la Sécurité Applicative}

Afin de sécuriser la plateforme \textbf{TalentPredict} contre les vulnérabilités web courantes, plusieurs mécanismes complémentaires ont été intégrés dès le Sprint 1 selon une approche \textit{Security by Design}.

\begin{itemize}
    \item \textbf{Protection contre les injections SQL :} L’accès à la base de données repose sur \texttt{Spring Data JPA} et \texttt{Hibernate}, utilisant des requêtes paramétrées et des mécanismes ORM empêchant l’exécution directe de requêtes SQL dynamiques non sécurisées.

    \item \textbf{Validation stricte des entrées utilisateur :} Les données saisies sont validées à plusieurs niveaux. Côté \textit{frontend} \texttt{Angular}, les \texttt{Reactive Forms} assurent une validation immédiate des champs (format email, longueur minimale, champs obligatoires). Côté \textit{backend} \texttt{Spring Boot}, les annotations \textit{Bean Validation} (\texttt{@Valid}, \texttt{@NotNull}, \texttt{@Email}, \texttt{@Size}) garantissent l’intégrité des données reçues par les \textit{endpoints} REST.

    \item \textbf{Protection contre les accès non autorisés :} L’ensemble des routes sensibles est sécurisé via \texttt{Spring Security} et des filtres \texttt{JWT} assurant l’authentification et la vérification des rôles utilisateurs avant l’accès aux ressources protégées.

    \item \textbf{Restriction des communications \textit{Cross-Origin} (CORS) :} Une politique CORS stricte a été configurée afin d’autoriser uniquement les requêtes provenant du \textit{frontend} \texttt{Angular} officiel, limitant ainsi les risques d’accès non autorisés depuis des domaines externes.

    \item \textbf{Hachage sécurisé des mots de passe :} Les mots de passe utilisateurs sont stockés sous forme hachée à l’aide de \texttt{BCrypt} avec salage cryptographique automatique afin de prévenir les attaques par dictionnaire et \textit{Rainbow Tables}.

    \item \textbf{Sécurisation des communications réseau :} Les échanges entre le \textit{frontend} \texttt{Angular}, le \textit{backend} \texttt{Spring Boot} et les services externes sont conçus pour fonctionner via le protocole HTTPS/TLS afin de garantir la confidentialité et l’intégrité des données transmises.

    \item \textbf{Journalisation et traçabilité :} Les opérations sensibles telles que les connexions utilisateurs, les tentatives d’accès refusées et les opérations critiques sont journalisées afin d’assurer une traçabilité minimale des activités de sécurité.
    \item \textbf{Protection contre les attaques par force brute (Rate Limiting) :} La bibliothèque Bucket4j a été intégrée pour limiter le nombre de tentatives d'authentification à 5 requêtes par minute par adresse IP sur les endpoints /api/auth/**. Au-delà de ce seuil, le système retourne une erreur HTTP 429 Too Many Requests, protégeant ainsi la plateforme contre les attaques automatisées par force brute.
\end{itemize}

Cette approche multicouche permet de renforcer significativement la robustesse de la plateforme tout en préparant l’intégration des futurs services d’intelligence artificielle dans un environnement sécurisé et \textit{scalable}.

\subsubsection*{7. Processus de Déploiement et Exécution}


Pour valider l’environnement cible et garantir la répétabilité de l'installation, l’application TalentPredict a été déployée en utilisant une approche d’\textit{Infrastructure as Code} (IaC). Cette méthode permet de décrire l'intégralité de l'infrastructure logicielle à travers des fichiers de configuration versionnés. Le déploiement s’effectue de manière automatisée via l'orchestrateur \texttt{Docker Compose}, selon le cycle suivant :

\begin{itemize}
    \item \textbf{Phase de Build (\textit{Backend} \& \textit{Frontend}) :} Compilation du code source Java via \texttt{Maven} (\texttt{mvn clean package}) et construction des \textit{bundles} \texttt{Angular}. Ces artefacts sont ensuite encapsulés dans des images \texttt{Docker} optimisées (utilisation de \textit{Multi-stage builds} pour réduire la taille des images).
    
    \item \textbf{Conteneurisation :} L’exécution de la commande \texttt{docker-compose up -d --build} automatise la création des conteneurs, l'allocation des ports réseau et le montage des volumes de persistance.
    
    \item \textbf{Orchestration et Dépendances :} \texttt{Docker} gère l'ordonnancement du démarrage. Le service \texttt{talentpredict-db} (\texttt{PostgreSQL 16}) est lancé en priorité. Grâce à une directive de \textit{Healthcheck}, le service \texttt{talentpredict-backend} (\texttt{Spring Boot 3.4.2}) attend que la base de données soit opérationnelle avant d'initier sa propre connexion. Enfin, le serveur \texttt{Nginx} servant l'interface \texttt{Angular} est déployé pour exposer l'application à l'utilisateur final.
\end{itemize}

\begin{figure}[H]
    \centering
    % Remplacez 'chemin_architecture_deploiement.png' par le nom exact de votre image
    \includegraphics[width=0.9\textwidth]{images/Deployment Architecture.png}
    \vspace{2mm}
    \caption{Architecture de déploiement Docker du Sprint 1}
    \label{fig:deploiement_docker_s1}
\end{figure}


Le schéma ci-dessus \ref{fig:deploiement_docker_s1} met en évidence la robustesse de l'infrastructure mise en place :

\begin{itemize}
    \item \textbf{Isolation Réseau :} L'utilisation d'un réseau privé \texttt{Docker} (\texttt{talentpredict-net}) assure qu'aucune brique sensible (\textit{Backend}, Base de données) n'est exposée directement sur Internet. Seul le serveur \texttt{Nginx} (Port 4200/443) est ouvert aux requêtes des utilisateurs.
    
    \item \textbf{Gestion de la Persistance :} Le volume \texttt{pg\_data} garantit la survie des données. En cas de mise à jour du conteneur \texttt{PostgreSQL}, les données restent intactes sur l'hôte, assurant une continuité de service.
    
    \item \textbf{Flux de Communication :} Le diagramme illustre le flux critique allant de la requête utilisateur (\textit{Browser}) vers les services métiers (\texttt{Spring Boot}), puis vers la persistance (\texttt{JDBC}). L'intégration d'un service \texttt{SMTP} externe complète ce dispositif pour les notifications critiques.
\end{itemize}


\section{Suivi et Pilotage du Sprint 1}
Le suivi et le pilotage du Sprint 1 s'inscrivent dans une démarche de gestion agile rigoureuse, visant à assurer la transparence de l'avancement et le respect des objectifs fixés. À travers l'analyse des indicateurs de vélocité et l'actualisation du backlog, ce contrôle continu permet de garantir l'alignement parfait entre les développements techniques et les exigences de la plateforme.
\subsection{Scrum Board}

Afin de garantir la transparence et l’inspection continue, essentielles dans la méthodologie \textit{Scrum}, nous avons utilisé des outils visuels pour gérer cette itération. Le sprint a duré trois semaines (soit 15 jours ouvrés), avec une charge de travail estimée à 38 Story Points.

Le suivi quotidien des tâches s’est matérialisé par un tableau de bord (\textit{Scrum Board}) reflétant l’état d’avancement des \textit{User Stories}. Ce management visuel a permis à l’équipe de développement d’identifier rapidement les goulots d’étranglement et de fluidifier le passage des tâches de la phase d'implémentation \textit{backend} vers la phase d'intégration et de test \textit{frontend}.

\begin{figure}[H]
    \centering
    % Remplacez 'scrum_board_sprint1' par le nom de votre fichier image
    \includegraphics[width=1\textwidth]{images/scrumboardSprint1.png}
    \vspace{2mm}
    \caption{État final du Scrum Board - Sprint 1}
    \label{fig:scrum_board_s1}
\end{figure}

\subsection{Le Burndown Chart}

Pour évaluer notre vélocité et s’assurer de l’atteinte de l’objectif dans le temps imparti, nous avons maintenu à jour le \textit{Burndown Chart} lors de nos \textit{Daily Stand-ups}.

\begin{figure}[H]
    \centering
    \begin{tikzpicture}
        \begin{axis}[
            width=16cm,
            height=8cm,
            title={\textbf{Burndown Chart - Sprint 1 }},
            xlabel={Jours ouvrés du Sprint (3 Semaines)},
            ylabel={Story Points restants},
            xmin=0, xmax=15,
            ymin=0, ymax=40,
            xtick={0,1,2,3,4,5,6,7,8,9,10,11,12,13,14,15},
            ytick={0,10,20,30,38,40},
            legend pos=north east,
            ymajorgrids=true,
            grid style=dashed,
            smooth
        ]
        
        % Ligne idéale (Tendance Idéale)
        \addplot[
            color=blue,
            dashed,
            mark=none,
            very thick
            ]
            coordinates {
            (0,38)(15,0)
            };
        \addlegendentry{Tendance Idéale}
        
        % Ligne réelle (Avancement Réel avec stagnation J7-J9)
        \addplot[
            color=red,
            mark=square*,
            very thick
            ]
            coordinates {
            (0,38)(1,36)(2,33)(3,30)(4,28)(5,25)(6,22)(7,22)(8,22)(9,20)(10,15)(11,10)(12,5)(13,2)(14,0)(15,0)
            };
        \addlegendentry{Avancement Réel}
        
        \end{axis}
    \end{tikzpicture}
    \caption{Burndown Chart du Sprint 1}
    \label{fig:burndown_sprint1}
\end{figure}

Le graphique montre la courbe de l’effort restant en \textit{Story Points} au fil des 15 jours du sprint. Nous avons rencontré quelques difficultés techniques notables :
\begin{itemize}
    \item \textbf{Stagnation (J7 à J9) :} Observable sur la courbe, cette période s'explique par les défis techniques liés à la configuration du filtre \texttt{JWT} dans \texttt{Spring Security} et la résolution des conflits de politique \texttt{CORS} entre l'interface \texttt{Angular} et le \textit{backend}.
    \item \textbf{Accélération finale :} La résolution de ces obstacles architecturaux a permis une forte accélération de la validation finale, clôturant ainsi le sprint dans les délais impartis.
\end{itemize}
\subsection{Indicateurs de Performance (KPIs) du Sprint 1}

Pour évaluer l'efficacité de cette première itération, nous avons défini un ensemble d'indicateurs de performance répartis selon deux axes : la gestion de projet Agile et la qualité technique.

\paragraph*{1. Indicateurs de Pilotage (Agilité)}
\begin{itemize}
    \item \textbf{Vélocité de l'Équipe :} L'objectif était de réaliser 38 \textit{Story Points}. Au terme du sprint, l'intégralité des points a été validée, soit un taux de complétion de 100\%.
    \item \textbf{Respect des Délais :} Le sprint a été clôturé en 15 jours, conformément au planning initial, malgré les défis techniques liés à l'intégration du module \texttt{Spring Security}.
    \item \textbf{Qualité du \textit{Backlog} :} Aucun \textit{Impediment} (blocage) majeur n'a nécessité le report d'une \textit{User Story} au sprint suivant.
\end{itemize}

\paragraph*{2. Indicateurs Techniques et Qualité}
Le tableau suivant résume les métriques de qualité logicielle obtenues :

\begin{table}[H]
    \centering
    \renewcommand{\arraystretch}{1.5}
    \small
    \begin{tabularx}{0.9\textwidth}{|X|c|c|c|}
        \hline
        \rowcolor[HTML]{F2F2F2}
        \textbf{Indicateur} & \textbf{Cible} & \textbf{Résultat Obtenu} & \textbf{État} \\
        \hline
        Couverture de code (Tests) & $> 80\%$ & $85\%$ & $\checkmark$ Validé \\
        \hline
        Temps de réponse moyen API & $< 200$ ms & $120$ ms & $\checkmark$ Validé \\
        \hline
        Failles de sécurité (\texttt{JWT}) & $0$ & $0$ & $\checkmark$ Validé \\
        \hline
        Disponibilité (\textit{Uptime} \texttt{Docker}) & $100\%$ & $100\%$ & $\checkmark$ Validé \\
        \hline
        Poids des images \texttt{Docker} & $< 500$ Mo & $320$ Mo & $\checkmark$ Validé \\
        \hline
    \end{tabularx}
    \caption{Tableau de bord des métriques techniques du Sprint 1}
    \label{tab:kpi_tech_s1}
\end{table}

\paragraph*{3. Analyse du Bilan}
Le bilan du Sprint 1 est extrêmement positif. La réussite des tests d'intégration et la stabilité de l'infrastructure \texttt{Docker} confirment que le socle de TalentPredict est sain. L'indicateur le plus significatif est la vélocité constante, qui démontre une bonne maîtrise de la \textit{stack} technologique (\texttt{Angular} / \texttt{Spring Boot}) par l'équipe de développement.
\section*{Conclusion}
\addcontentsline{toc}{section}{Conclusion du Sprint 1}
Le premier \textit{sprint} du projet TalentPredict s'achève sur un bilan très positif. Cette itération, consacrée aux fondations techniques et à la gestion de l'identité, a permis de mettre en place les bases essentielles au bon fonctionnement de la plateforme. 

En appliquant rigoureusement la méthodologie \textit{Scrum}, nous sommes passés d'un concept théorique à une application fonctionnelle, sécurisée et conteneurisée. Les principaux objectifs ont été atteints avec succès : l'architecture a été stabilisée avec la mise en place d'une structure N-tiers robuste basée sur \texttt{Spring Boot} et \texttt{Angular} ; la sécurité a été assurée grâce à l'implémentation de l'authentification et de l'autorisation via des jetons \texttt{JWT} ; et la modélisation des données a permis de représenter les principaux éléments du profil professionnel, incluant le CV, \texttt{GitHub}, \texttt{LinkedIn} et le \textit{portfolio}.
Ce \textit{sprint} a non seulement validé la solidité du socle technique, mais a également démontré notre capacité à surmonter les défis d'implémentation, notamment à travers des tests d'\texttt{API} systématiques avec \texttt{Postman} et l'utilisation d'un environnement \texttt{DevOps} basé sur \texttt{Docker}. 
Cependant, un profil utilisateur reste une donnée statique tant qu'il n'est pas interprété par une couche d'intelligence. C'est pourquoi le Sprint 2 constitue une étape clé du projet, avec l'introduction du Moteur d'Évaluation Multidimensionnelle, visant à passer de la simple collecte de données à la génération de valeur, grâce à l'intégration d'une analyse automatique des compétences basée sur l'IA et d'une évaluation psychométrique (PCM), ouvrant ainsi la voie vers le cœur prédictif de TalentPredict.



\chapter{Sprint 2  Évaluation Multi-Source des Compétences}

\section*{Introduction}
\phantomsection
\addcontentsline{toc}{section}{Introduction}

Ce troisième chapitre documente la réalisation du Sprint 2, dont l'objectif est de livrer le module complet d'évaluation des compétences techniques et comportementales de TalentPredict. Contrairement aux fondations posées lors du Sprint 1 (sécurité, authentification, architecture de base), ce sprint se concentre sur le cœur métier de la plateforme : l'analyse intelligente multi-source des talents.
Le module se décompose en deux sous-systèmes complémentaires développés en parallèle par les deux membres du binôme : le module Compétences comportementales  et le module Compétences techniques . Ces deux modules partagent une architecture commune mais reposent sur des pipelines technologiques distincts, reflétant la nature différente des données analysées.
La structure de ce chapitre suit la méthodologie Agile Scrum adoptée par l'équipe. Nous présentons successivement : l'objectif du sprint et ses livrables attendus , la séance de planification et les décisions d'architecture , le backlog estimé , puis le déroulement technique du sprint organisé selon les phases d'analyse, conception, codage et validation . 



\section{Objectifs du Sprint 2}

Avant de décrire le déroulement du sprint, il est essentiel de définir précisément le périmètre fonctionnel visé et les critères de validation qui guideront l’équipe tout au long des trois semaines de développement.

L’objectif majeur de ce sprint est de livrer un moteur d’évaluation multi-sources complet et autonome. L'enjeu est de faire évoluer la plateforme TalentPredict : d'une simple base de données déclarative, elle devient un véritable système de diagnostic intelligent capable de profiler un collaborateur sous toutes ses coutures en agrégeant de multiples sources d’analyse complémentaires pour produire une cartographie unifiée (compétences techniques et comportementales).

\paragraph*{1. Objectifs fonctionnels et interactifs :}
\begin{itemize}[leftmargin=1cm]
    \item \textbf{Évaluation comportementale :} Déploiement d'un workflow automatisé intégrant le test heuristique PCM, l'analyse sémantique du CV et de GitHub, ainsi que la résolution de scénarios immersifs générés par l'IA.
    \item \textbf{Évaluation technique :} Mise en œuvre d'un parcours de validation factuel combinant l'extraction de la \textit{stack} technique (CV, LinkedIn), l'analyse de l'historique GitHub, et des tests interactifs (QCM adaptatif et challenge de code).
    \item \textbf{Job Matching :} Intégration d'un algorithme de correspondance permettant de calculer un score de compatibilité entre le profil consolidé du collaborateur et plus de 20 profils métiers standards.
\end{itemize}

\paragraph*{2. Objectifs techniques et opérationnels :}
\begin{itemize}[leftmargin=1cm]
    \item \textbf{Orchestration asynchrone :} Configuration et déploiement de workflows n8n (\textit{Master Agent}, \textit{CV Parser}, \textit{GitHub Analysis}) pour piloter les flux de données lourds sans surcharger le backend central.
    \item \textbf{Inférence IA locale :} Implémentation du microservice Python (FastAPI) communiquant avec l'écosystème local Ollama (Llama 3 / CodeLlama) pour la génération, la correction et l'évaluation des scénarios et des défis algorithmiques.
    \item \textbf{Persistance unifiée :} Conception et migration du schéma de base de données PostgreSQL via Spring Boot pour stocker les entités \texttt{Skill} (en distinguant les énumérations \texttt{TypeSkill.TECH} et \texttt{TypeSkill.SOFT}).
\end{itemize}

\paragraph*{Critères d’acceptation (Definition of Done) :}
La validation du pipeline d'évaluation repose sur les critères stricts suivants :
\begin{itemize}[leftmargin=1cm]
    \item \textbf{Cohérence multi-sources :} Le  Master Agent doit agréger avec succès les scores comportementaux (0--100) issus de GitHub, du CV, du PCM et du scénario immersif pour générer un rapport unifié cohérent.
    \item \textbf{Exactitude de l'extraction technique :} Le \textit{parser} doit être capable d'identifier avec précision un catalogue de plus de 150 technologies avec leurs années d'expérience associées à partir d'un fichier brut.
    \item \textbf{Évaluation algorithmique :} Le challenge de code généré par CodeLlama doit être évalué de manière autonome selon 4 critères précis (syntaxe, logique, complexité, sécurité) pour en déduire un niveau (\textit{Beginner} à \textit{Expert}).
    \item \textbf{Intégrité des données :} Chaque nouvelle compétence détectée et validée par le pipeline doit être persistée instantanément et correctement catégorisée en base de données.
\end{itemize}

\vspace{0.5cm}
Le Sprint 2 vise à livrer treize composants techniques répartis sur les deux modules parallèles d'évaluation. Le tableau suivant formalise ces livrables avec leurs critères de validation respectifs :

\begin{table}[H]
\centering
\small
\renewcommand{\arraystretch}{1.3}
\begin{tabularx}{\textwidth}{|c|>{\bfseries}P{0.35\textwidth}|Y|}
\hline
\textbf{\#} & \textbf{Livrable technique} & \textbf{Critère de validation détaillé} \\
\hline
1 & Analyse GitHub comportementale & Génération de 4 scores comportementaux (0--100). \\
\hline
2 & Analyse CV comportementale & Calcul du score CV et extraction des indices sémantiques. \\
\hline
3 & Test heuristique PCM & Identification des 6 dimensions cognitives et du type dominant via 18 questions. \\
\hline
4 & Synthèse Master Agent & Agrégation réussie produisant un rapport pondéré unifié. \\
\hline
5 & Scénario IA comportemental & Évaluation sémantique produisant 4 scores psychologiques. \\
\hline
6 & Analyse GitHub technique & Extraction avérée des langages, frameworks et du niveau d'engagement. \\
\hline
7 & Parsing CV technique & Identification précise sur une base de 150+ technologies et années d’expérience. \\
\hline
8 & Analyse LinkedIn & Établissement du positionnement professionnel et score /10. \\
\hline
9 & QCM technique adaptatif & Génération dynamique des questions et \textit{scoring} pondéré validé. \\
\hline
10 & Challenge de code (CodeLlama) & Génération et évaluation algorithmique fonctionnelles (basées sur 4 critères). \\
\hline
11 & Scoring multi-sources & Classification de la maîtrise technique de \textit{Beginner} à \textit{Expert}. \\
\hline
12 & Moteur de Job matching & Calcul du score de compatibilité (/100) face à 20 profils métiers cibles. \\
\hline
13 & Persistance des compétences & Stockage relationnel réussi de l'entité \texttt{Skill} en base de données PostgreSQL. \\
\hline
\end{tabularx}
\renewcommand{\arraystretch}{1.0}
\caption{Livrables attendus et validation du Sprint 2}
\end{table}
La figure~\ref{fig:pip-glob-sprint2} illustre le pipeline global d'analyse implémenté durant cette phase. Elle retrace le flux complet des données, depuis l'ingestion des sources hétérogènes (CV, API GitHub, résultats de QCM, challenges de code et analyse comportementale PCM) jusqu'à leur agrégation en un score d'évaluation unifié.
\begin{figure}[H]
  \centering
  \includegraphics[width=0.9\textwidth]{images/pipline glob sprint2.png}
  \caption{Pipeline global d’analyse multi-sources}
  \label{fig:pip-glob-sprint2}
\end{figure}

\section{Séance du Sprint Planning}
La séance de Sprint Planning s'est tenue  en présence de l'ensemble de l'équipe projet. Cette réunion structurante a permis de traduire les exigences fonctionnelles en choix techniques concrets et de définir les modalités de collaboration entre les deux développeurs. Les décisions prises durant cette séance conditionnent l'architecture finale du module
Participants : Barouni Rahma, Abdeladhim Mohamed  , M. Mohamed Amine Gharbi (Product Owner/ encadrant professionnel et Takwa ben Aicha ( encadrente Academique )).

\textbf{Durée du sprint :} quatre semaines, correspondant aux semaines 5 à 8 du diagramme de Gantt du projet.

 

\textbf{Éléments discutés et décisions clés}

\textbf{Module compétences comportementales  :}
\begin{itemize}[leftmargin=1cm]
  \item n8n est retenu comme orchestrateur des workflows d'analyse GitHub, CV et PCM afin de dissocier la logique d'orchestration du backend applicatif et de permettre la modification des prompts sans recompilation.
  \item Ollama \texttt{llama3.2:latest} est choisi comme moteur LLM local pour l'ensemble du module compétences comportementales, ce qui favorise une conformité RGPD par conception.
  \item Le scoring PCM est implémenté en JavaScript pur, sans appel LLM, afin de garantir la reproductibilité et de réduire la latence d'exécution.
  \item Une séparation claire des responsabilités est établie : n8n gère l'analyse de profil statique, tandis que FastAPI prend en charge l'évaluation comportementale en temps réel via les scénarios IA.
\end{itemize}

\textbf{Module compétences techniques  :}
\begin{itemize}[leftmargin=1cm]
  \item Une architecture  :  Ollama \texttt{llama3.2} en mode local avec pipeline direct.
  \item \texttt{pdfplumber} est retenu pour l'extraction PDF côté serveur, car il offre une meilleure fiabilité que \texttt{PyPDF2} sur les documents PDF vectoriels.
  \item \texttt{CodeLlama} est utilisé spécifiquement pour les challenges de code, grâce à de meilleures performances sur la génération et l'évaluation de code que \texttt{llama3.2}.
  \item Le service \texttt{TalentPredictAiProxyService} de Spring Boot assure le pont entre l'API Java et le microservice Python FastAPI avec un délai d'attente configurable de 120 secondes.
\end{itemize}

\textbf{Priorités retenues :}
\begin{enumerate}[leftmargin=1cm]
  \item Workflows n8n : GitHub Analysis, CV Parser, PCM Test et Master Agent.
  \item Scénario comportemental IA dans FastAPI.
  \item CV Parser et GitHub Analyzer pour les compétences techniques.
  \item QCM adaptatif et challenge de code avec FastAPI et CodeLlama.
  \item Persistance des compétences via Spring Boot et PostgreSQL.

\end{enumerate}

\textbf{Hypothèses et contraintes :}
\begin{itemize}[leftmargin=1cm]
  \item \textbf{Latence Ollama :} les délais d'attente sont configurés avec des fallbacks heuristiques afin de garantir une réponse exploitable même en cas de lenteur du modèle local.
  \item \textbf{Format JSON instable :} les réponses produites par Ollama font l'objet d'un nettoyage systématique avant leur exploitation par les services applicatifs.
  \item \textbf{GitHub rate limit :} un \texttt{GITHUB\_TOKEN} est requis en production, car l'API GitHub est limitée à 60 requêtes par heure sans authentification.
  \item \textbf{PDF scannés :} un avertissement est retourné et une analyse partielle est réalisée lorsque le texte du document n'est pas extractible.
\end{itemize}
Ces décisions d'architecture ont été formalisées dans le Sprint Backlog présenté en section suivante.

\section{Sprint Backlog Estimé}
Sur la base des décisions prises lors du Sprint Planning, l'équipe a décomposé l'objectif global en User Stories estimées. Le backlog ci-dessous reflète la priorisation métier validée par le Product Owner et l'effort estimé en jours-homme selon la méthode du Planning Poker. La charge totale de 37 jours-homme (20j pour Compétences techniques + 17j pour Compétences comportementales) correspond à un sprint de 3 semaines avec deux développeurs à temps plein.

Le tableau \ref{tab:backlog-hard-skills} présente le backlog du module Compétences techniques.

\begin{table}[H]
\centering
\small
\setlength{\tabcolsep}{3pt}
\begin{tabular}{|P{1.0cm}|P{2.4cm}|P{5.8cm}|C{0.7cm}|P{1.5cm}|}
\hline
\textbf{ID} & \textbf{Module / US} & \textbf{Tâches de développement} & \textbf{JH} & \textbf{Priorité} \\
\hline
HS-01 & Analyse GitHub technique & Récupération des repos via l'API REST ; détection des langages, frameworks et niveau d'activité. & 3 & Haute \\
\hline
HS-02 & Parsing CV PDF/DOCX & Extraction de 150+ technologies par regex via \texttt{pdfplumber} et \texttt{python-docx} ; estimation des années d'expérience. & 2 & Haute \\
\hline
HS-03 & Analyse LinkedIn & Analyse via Ollama du positionnement, arc de carrière et cohérence CV/GitHub. & 2 & Haute \\
\hline
HS-04 & Génération QCM adaptatif & Moteur de génération de 6 à 12 questions par compétence et niveau via Ollama. & 3 & Haute \\
\hline
HS-05 & Scoring QCM pondéré & Évaluation selon trois dimensions : exactitude $\times$ confiance $\times$ temps de réponse. & 2 & Haute \\
\hline
HS-06 & Challenge de code CodeLlama & Génération \texttt{fix\_bugs} / \texttt{complete} / \texttt{refactor} ; évaluation sur 4 critères avec pénalité par indice. & 3 & Moyenne \\
\hline
HS-07 & Scoring multi-sources & Déduplication de 500+ alias ; nivellement Beginner $\rightarrow$ Expert ; job match sur 20+ profils métiers. & 3 & Haute \\
\hline
HS-08 & Persistance compétences TECH & Entité \texttt{Skill} (\texttt{TypeSkill.TECH}) avec endpoints Spring Boot pour sauvegarde et validation ADMIN. & 2 & Haute \\
\hline
\multicolumn{3}{|r|}{\textbf{Total}} & \textbf{20} & \\
\hline
\end{tabular}
\caption{Sprint Backlog estimé du module Compétences techniques}\label{tab:backlog-hard-skills}
\end{table}

Le tableau \ref{tab:backlog-soft-skills} détaille le backlog du module Compétences comportementales.

\begin{table}[H]
\centering
\small
\setlength{\tabcolsep}{3pt}
\begin{tabular}{|P{1.0cm}|P{2.5cm}|P{5.7cm}|C{0.7cm}|P{1.5cm}|}
\hline
\textbf{ID} & \textbf{Module / US} & \textbf{Tâches de développement} & \textbf{JH} & \textbf{Priorité} \\
\hline
SS-01 & Analyse GitHub comportementale & Workflow n8n : scoring heuristique JS (repos, stars, forks) ; analyse Ollama des 4 dimensions comportementales. & 3 & Haute \\
\hline
SS-02 & Extraction compétences comportementales CV & Workflow n8n CV Parser : envoi du texte brut à Ollama ; extraction des indices comportementaux et score CV global. & 2 & Haute \\
\hline
SS-03 & Test PCM (18 questions) & Workflow n8n : calcul 100\% heuristique JS des 6 dimensions PCM, sans appel LLM. & 2 & Haute \\
\hline
SS-04 & Master Agent --- Synthèse pondérée & Orchestrateur n8n : appels parallèles aux 3 sous-workflows ; synthèse Ollama CV 40\% / PCM 30\% / GitHub 30\%. & 2 & Haute \\
\hline
SS-05 & Génération de scénario IA & FastAPI + Ollama : génération d'un dilemme professionnel contextualisé par rôle et niveau. & 2 & Moyenne \\
\hline
SS-06 & Évaluation de la réponse scénario & FastAPI + Ollama : évaluation sur 4 dimensions psychologiques (empathie, assertivité, pragmatisme, clarté). & 2 & Moyenne \\
\hline
SS-07 & Entretien vocal IA multi-tours & FastAPI + Ollama : entretien de 5 tours avec pivot automatique vers la dimension la plus faible (score < 55/100). & 3 & Moyenne \\
\hline
SS-08 & Persistance compétences SOFT & Sauvegarde de l'entité \texttt{Skill} (\texttt{TypeSkill.SOFT}) depuis les résultats du Master Agent et scores PCM via Spring Boot. & 1 & Haute \\
\hline
\multicolumn{3}{|r|}{\textbf{Total}} & \textbf{17} & \\
\hline
\end{tabular}
\caption{Sprint Backlog estimé du module Compétences comportementales}\label{tab:backlog-soft-skills}
\end{table}

\section{Démarrage du Sprint 2}

Le début du Sprint 2 a été caractérisé par une phase d’initialisation technique approfondie en préparation des développements relatifs au moteur d'évaluation multi-sources. L’équipe s’est aussitôt focalisée sur l'orchestration des flux de données via \texttt{n8n} et la configuration de l'intelligence artificielle locale (Ollama), tout en établissant les ponts de communication avec le backend Spring Boot. Cette phase essentielle a posé les fondations nécessaires pour l’analyse croisée des compétences techniques et comportementales, garantissant ainsi un passage en douceur de la modélisation théorique à la mise en œuvre effective des pipelines d'évaluation.

\subsubsection*{1. Activités réalisées}

Durant cette itération, l’équipe a mené de front des développements backend complexes et des intégrations poussées d’intelligence artificielle :
\begin{itemize}[leftmargin=1cm]
    \item \textbf{Ingénierie des flux (n8n) :} Conception et automatisation des workflows d'extraction et d'analyse pour décharger le backend des traitements asynchrones lourds.
    \item \textbf{Développement IA (FastAPI \& Ollama) :} Implémentation du microservice Python pour le \textit{parsing} sémantique des CV, la génération de scénarios immersifs et l'évaluation du code source.
    \item \textbf{Développement Backend (Spring Boot) :} Création des \textit{endpoints} REST pour la collecte des données brutes (CV, GitHub) et implémentation du stockage relationnel des entités de compétences.
    \item \textbf{Développement Frontend (Angular) :} Création des interfaces de passation de tests (QCM, questionnaire PCM, IDE de code) et des formulaires de soumission de profil.
    \item \textbf{Intégration d'API tierces :} Connexion sécurisée à l'API GitHub pour récupérer l'historique d'activité et les statistiques de collaboration des employés.
\end{itemize}

\subsubsection*{2. Livrables produits}

Les livrables suivants ont été validés et intégrés à la branche principale du projet :
\begin{itemize}[leftmargin=1cm]
    \item \textbf{Module d'évaluation comportementale :} Workflow complet de détermination des 6 dimensions cognitives.
    \item \textbf{Module d'évaluation technique :} Pipeline d'analyse factuelle du niveau de maîtrise (\textit{Beginner} à \textit{Expert}).
    \item \textbf{Service d'Intelligence Artificielle :} API locale propulsée de manière autonome par Ollama et CodeLlama.
    \item \textbf{Moteur de Job Matching :} Algorithme de calcul de compatibilité basé sur plus de 20 profils métiers.
    \item \textbf{Interfaces Utilisateur :} Écrans interactifs de passation de tests dynamiques et adaptatifs.
\end{itemize}

\subsubsection*{3. Difficultés rencontrées et traitement}

L’exécution de ce sprint a été marquée par trois défis techniques majeurs liés à l'intelligence artificielle et au traitement des données :
\begin{itemize}[leftmargin=1cm]
    \item \textbf{Difficulté 1 : Hétérogénéité des formats de CV}
    \begin{itemize}[label=--]
        \item \textit{Problème :} L'extraction de texte brut à partir de fichiers PDF générés par différents outils entraînait des pertes de données sémantiques ou des erreurs d'encodage.
        \item \textit{Traitement :} Implémentation de la bibliothèque \texttt{pdfplumber} dans le microservice FastAPI, couplée à des expressions régulières avancées (Regex) pour nettoyer et normaliser le texte avant son ingestion par le LLM.
    \end{itemize}
    \item \textbf{Difficulté 2 : Limites de requêtes de l'API GitHub (\textit{Rate Limits})}
    \begin{itemize}[label=--]
        \item \textit{Problème :} Les requêtes multiples pour extraire l'historique détaillé des commits provoquaient des blocages temporaires imposés par GitHub.
        \item \textit{Traitement :} Optimisation des appels via \texttt{n8n} en ciblant uniquement les métadonnées essentielles et mise en place d'un système de mise en cache temporaire.
    \end{itemize}
    \item \textbf{Difficulté 3 : Temps d'inférence de l'IA locale}
    \begin{itemize}[label=--]
        \item \textit{Problème :} La génération asynchrone des scénarios professionnels complexes par Ollama causait des \textit{timeouts} côté frontend (Erreur 503).
        \item \textit{Traitement :} Transition vers une architecture événementielle en s'appuyant sur les \textit{Server-Sent Events} (SSE) et le Communication Hub pour libérer la connexion HTTP et notifier l'utilisateur dès la fin du traitement de l'IA.
    \end{itemize}
\end{itemize}

\subsection{Analyse}

La phase d’analyse vise à définir le périmètre fonctionnel du module d'évaluation et à identifier avec précision les interactions entre le système et les utilisateurs. Ce sprint introduit une forte automatisation grâce à l'Intelligence Artificielle et à l'orchestration des flux.

\subsubsection*{1. Identification des Acteurs}

Le système interagit avec plusieurs acteurs principaux, parmi lesquels se trouvent des entités humaines et des acteurs automatisés :
\begin{itemize}[leftmargin=1cm]
    \item \textbf{Utilisateur (Employé) :} Constitue le cœur du processus de collecte de données. Il soumet ses preuves documentaires (CV, identifiant GitHub) et s'engage activement dans l'évaluation en répondant au test PCM, au QCM adaptatif, en résolvant le challenge de code, puis en interagissant avec le scénario immersif IA.
    \item \textbf{Administrateur (RH) :} Supervise le bon déroulement des évaluations de son équipe. Il consulte les résultats consolidés, valide formellement les compétences détectées par la machine et s'appuie sur le système pour générer les rapports PDF d'aide à la décision.
    \item \textbf{Moteurs IA et Orchestrateurs (« Systèmes ») :} Agissent de manière autonome en arrière-plan pour traiter la complexité algorithmique :
    \begin{itemize}[label=--]
        \item \textit{Ollama \& CodeLlama} : Réalisent l’analyse sémantique textuelle, la synthèse comportementale, ainsi que la génération et la correction intelligente des scénarios et du code.
        \item \textit{Microservice FastAPI} : Exécute le \textit{parsing} des CV, pré-traite les données brutes, et assure l'interface directe avec les modèles LLM.
        \item \textit{n8n} : Ordonnance les appels asynchrones entre le backend, les API externes et les services d'IA, garantissant la bonne chronologie de l'évaluation.
    \end{itemize}
    \item \textbf{Sources externes (GitHub API) :} Agit comme un acteur de confiance tiers, fournissant des données factuelles (dépôts, commits, langages, \textit{pull requests}) pour objectiver les déclarations initiales du collaborateur.
\end{itemize}

\subsubsection{Spécification des besoins fonctionnels}

\textbf{BF-SS-01 : Analyse de l'activité GitHub (compétences comportementales)}\par
Le système calcule un score heuristique JavaScript (repos 40\%, stars 40\%, forks 20\%), puis soumet les données à Ollama pour évaluer quatre dimensions comportementales : curiosité, collaboration, ownership et consistance. Un fallback heuristique est retourné en cas d'échec, avec des timeouts de 10 s (GitHub API) et 20 s (Ollama).

\textbf{BF-SS-02 : Extraction compétences comportementales depuis CV}\par
Le texte brut extrait via \texttt{pdfjs-dist} est analysé par Ollama pour détecter les indices comportementaux (travail en équipe, leadership, adaptabilité), le niveau d'éducation et les années d'expérience, produisant un score global de 0 à 100. Le prompt exclut explicitement les technologies pour éviter toute contamination avec le module compétences techniques (timeout : 15 s).

\textbf{BF-SS-03 : Test de personnalité PCM}\par
Les scores des six dimensions PCM (communication, discipline, curiosité, collaboration, ownership, leadership) sont calculés algorithmiquement en JavaScript pur, chaque dimension correspondant à la moyenne de trois questions sur une échelle de 0 à 5, sans appel LLM, garantissant une reproductibilité totale en moins de 100 ms.
\begin{figure}[H]
  \centering
  \includegraphics[width=0.9\textwidth]{images/pipline pcm sprint2.png}
  \caption{Pipeline d’évaluation PCM}
  \label{fig:pip-PCM-sprint2}
\end{figure}
\textbf{BF-SS-04 : Synthèse Master Agent}\par
Le Master Agent orchestre en parallèle via \texttt{Promise.all} les trois sous-workflows (GitHub 30\%, CV 40\%, PCM 30\%), puis soumet les résultats agrégés à Ollama pour produire un profil unifié (score global, type de personnalité, forces, recommandations). En cas d'indisponibilité GitHub, la pondération bascule dynamiquement vers CV 60\% / PCM 40\% (timeout global : 30 s).

\textbf{BF-SS-05 : Génération et évaluation de scénarios comportementaux IA}\par
Ollama génère un dilemme professionnel parmi cinq types (deadline, conflit, feedback client, décision, intégration) contextualisé par rôle et niveau (température 0.7). La réponse libre du employé est ensuite évaluée sur quatre dimensions psychologiques : empathie, assertivité, pragmatisme et clarté, avec une température de 0.2 pour garantir la reproductibilité, produisant scores, forces, axes d'amélioration et profil culture-add.

\textbf{BF-HS-01 : Analyse GitHub compétences techniques}par
Le système récupère les dépôts publics via l'API REST GitHub, calcule les statistiques de langages et détecte les frameworks depuis les topics et descriptions, avant enrichissement par Ollama. Un  GITHUB\_TOKEN  est requis en production (limite : 60 req/h sans token ; timeout : 30 s). Les champs  language\_stats et detected\_frameworks alimentent ensuite le moteur de scoring.

\textbf{BF-HS-02 : Extraction de technologies depuis CV}\par
\texttt{pdfplumber} (PDF vectoriel) ou \texttt{python-docx} (DOCX) extrait le texte, puis 150+ technologies sont détectées par expressions régulières organisées par catégorie (langages, frameworks, cloud, DevOps, IA/ML, bases de données). Les années d'expérience sont estimées via patterns linguistiques. En cas de PDF scanné, une analyse partielle est poursuivie avec avertissement.
\begin{figure}[H]
  \centering
  \includegraphics[width=0.9\textwidth]{images/pipline cv sprint2.png}
  \caption{Pipeline d’analyse du CV}
  \label{fig:pip-cv-tech-sprint2}
\end{figure}
\textbf{BF-HS-03 : Analyse LinkedIn}\par
Ollama analyse le profil LinkedIn (URL ou texte) pour extraire le rôle actuel, l'arc de carrière, les certifications et la cohérence avec le CV et GitHub. Un score de complétude de 0 à 10 est attribué, permettant de détecter les incohérences entre les déclarations du employé et son activité réelle.

\textbf{BF-HS-04 : Test QCM technique adaptatif}\par
Ollama génère dynamiquement 6 à 12 questions par compétence et niveau (Junior / Senior / Lead) à température 0.4, avec fallback sur une banque prédéfinie balisée \texttt{source: fallback}. Le scoring combine trois dimensions : exactitude $\times$ confiance $\times$ temps, avec un multiplicateur 1.2 pour haute confiance correcte, une pénalité de surconfiance pour haute confiance incorrecte, et un bonus de 0.05 pour bonne réponse difficile en moins de 15 secondes.
\begin{figure}[H]
  \centering
  \includegraphics[width=0.9\textwidth]{images/QCM.png}
  \caption{Génération de QCM adaptatifs par IA}
  \label{fig:pip-QCM-sprint2}
\end{figure}
\textbf{BF-HS-05 : Challenge de code avec CodeLlama}\par
CodeLlama génère un défi de type \texttt{fix\_bugs}, \texttt{complete} ou \texttt{refactor} et évalue la soumission selon quatre critères pondérés : correctness (40 pts), code quality (30 pts), efficiency (20 pts) et readability (10 pts). Une pénalité de 10 points par indice utilisé est appliquée, soit : score final = total brut $-$ (\texttt{hints\_used} $\times$ 10).

\textbf{BF-HS-06 : Scoring multi-sources et job matching}\par
Les compétences issues des trois sources~\ref{fig:Fusion-multi-scores-sprint2} sont consolidées via 500+ alias normalisés dans \texttt{skill\_extractor.py}, puis nivelées par triangulation : Expert ($\geq$ 3 sources, score $\geq$ 90), Advanced ($\geq$ 2 sources, 70--89), Intermediate (1 source forte, 40--69), Beginner ($\sim$25). Le job matching compare le profil avec 20+ profils métiers selon la formule :
\[
  \texttt{job\_match\_score} = \left(\frac{\texttt{earned\_weight}}{\texttt{total\_weight}}\right) \times 85 + \texttt{bonus\_extra\_skills},
\]
ajustée par l'\texttt{experience\_score}.
\begin{figure}[H]
  \centering
  \includegraphics[width=0.9\textwidth]{images/pipline cv sprint2.png}
  \caption{Fusion multi-sources des
scores}
  \label{fig:Fusion-multi-scores-sprint2}
\end{figure}

\subsubsection{Objectifs non fonctionnels}

\begin{table}[H]
\centering
\renewcommand{\arraystretch}{1.3}
\begin{tabularx}{\TableWidth}{|P{2.6cm}|Y|}
\hline
\textbf{Critère} & \textbf{Exigence} \\
\hline
Performance & Analyse complète compétences comportementales en moins de 15 secondes et APIs REST courantes en moins de 2 secondes \\
\hline
Confidentialité & Traitement 100\% local via Ollama pour les données RH sensibles \\
\hline
Fiabilité & Fallbacks systématiques pour chaque appel LLM avec timeouts configurés et réponses heuristiques par défaut \\
\hline
Scalabilité & Architecture conteneurisée Docker autorisant le déploiement multi-instances \\
\hline
Sécurité & Authentification JWT Bearer Token sur les endpoints Spring Boot et RBAC USER/ADMIN \\
\hline
\end{tabularx}
\caption{Objectifs non fonctionnels du Sprint 2}
\label{tab:objectifs-non-fonctionnels}
\end{table}

\subsubsection{Contraintes techniques}
L'analyse préliminaire des technologies retenues a révélé quatre contraintes majeures susceptibles d'impacter la fiabilité et les performances du pipeline d'évaluation. Ces contraintes ont été anticipées dès la phase de planification afin de concevoir des mécanismes de mitigation adaptés, notamment des fallbacks heuristiques et des timeouts explicites, avant toute décision d'architecture.
\begin{itemize}[leftmargin=1cm]
  \item \textbf{Latence Ollama :} le modèle \texttt{llama3.2} peut dépasser 20 secondes, d'où des timeouts explicites et des fallbacks heuristiques immédiats.
  \item \textbf{Format JSON des modèles locaux :} les réponses Ollama contiennent fréquemment des artefacts Markdown ; un nettoyage systématique est donc appliqué dans les workflows n8n et les services Python.
  \item \textbf{PCM sans LLM :} le scoring des 18 questions est entièrement algorithmique afin d'assurer une forte reproductibilité et une exécution inférieure à 100 ms.
  \item \textbf{CodeLlama :} ce modèle étant plus lourd, le timeout des challenges de code est porté à 600 secondes.

  \item \textbf{GitHub rate limit :} en production, \texttt{GITHUB\_TOKEN} est requis pour éviter la limite de 60 requêtes par heure.
  \item \textbf{PDF scannés :} \texttt{pdfplumber} ne peut extraire le texte d'un PDF image, ce qui conduit à une analyse partielle accompagnée d'un avertissement.
\end{itemize}

%\figureplaceholder{Diagramme de contexte unifié compétences comportementales + compétences techniques}{fig:contexte-unifie}



\subsubsection{Cas d'utilisation}

\textbf{1. Diagramme de cas d'utilisation global (Sprint 2)}

Le diagramme global~\ref{fig:usecase-global-sprint2} illustre l'ensemble des interactions fonctionnelles entre les acteurs principaux  l'Employé et l'Administrateur RH  et le système TalentPredict AI dans le cadre du Sprint 2. Il regroupe les fonctionnalités liées à l'évaluation des compétences techniques et comportementales. L'employé soumet ses données (CV, profil GitHub), passe ses évaluations (test PCM, QCM adaptatif, scénario immersif, challenge de code) et consulte les recommandations générées. L'Administrateur RH supervise les évaluations, valide les résultats et génère des rapports. Les services externes   Service IA, Service GitHub et Service de communication temps réel   interviennent respectivement pour l'analyse intelligente, la récupération des données de dépôts, et la diffusion des mises à jour en temps réel.

\begin{figure}[H]
  \centering
  \includegraphics[width=0.9\textwidth]{images/D use case sprint2.png}
  \caption{Diagramme de cas d'utilisation : Passation et consultation de l'évaluation globale}
  \label{fig:usecase-global-sprint2}
\end{figure}

\begin{table}[H]
\centering
\small
\renewcommand{\arraystretch}{1.25}
\begin{tabularx}{\textwidth}{|P{0.26\textwidth}|Y|}
\hline
\textbf{Élément} & \textbf{Description} \\
\hline
Nom du cas & Passation et consultation de l'évaluation globale \\
\hline
Acteur principal & Employé \\
\hline
Acteurs secondaires & Administrateur RH, Service IA, Service GitHub, Service de communication temps réel \\
\hline
Pré-conditions & L'employé possède un compte actif et s'est authentifié sur TalentPredict. \\
\hline
Scénario nominal & \begin{minipage}[t]{\linewidth}
\begin{enumerate}[leftmargin=*,nosep]
  \item L'employé s'authentifie puis accède au module d'évaluation.
  \item Il soumet son CV et renseigne son nom d'utilisateur GitHub pour déclencher les analyses comportementale et technique.
  \item Le système évalue les compétences comportementales (CV sémantique, collaboration GitHub, test PCM, scénario immersif) et les compétences techniques (CV technique, activité GitHub, QCM adaptatif, challenge de code).
  \item Les résultats d'évaluation sont produits et des recommandations de formation sont générées par le Service IA.
  \item L'employé consulte les recommandations et le score de correspondance au poste.
  \item Le Service de communication temps réel diffuse les mises à jour ; l'Administrateur RH supervise, valide les résultats et génère les rapports.
\end{enumerate}
\end{minipage} \\
\hline
Scénario alternatif & Si un module est indisponible (ex.\ : timeout du Service IA ou profil GitHub introuvable), l'évaluation partielle est sauvegardée et le Service de communication temps réel prévient l'employé de la reprise future. \\
\hline
Post-conditions & Le profil complet du collaborateur est généré (compétences techniques et comportementales), validé par l'Administrateur RH, et prêt pour l'intégration au plan de formation. \\
\hline
\end{tabularx}
\renewcommand{\arraystretch}{1.05}
\caption{Fiche descriptive du cas d'utilisation : Passation et consultation de l'évaluation globale}
\end{table}

\textbf{2. Diagramme de cas d'utilisation : Module Compétences comportementales}

Ce diagramme~\ref{fig:usecase-soft-skills-sprint2} présente les fonctionnalités du module d'évaluation comportementale. L'employé soumet son CV et renseigne son nom d'utilisateur GitHub, ce qui déclenche respectivement l'extraction sémantique du CV et l'analyse de sa dynamique de collaboration. Il passe également le test PCM et répond à un scénario métier immersif. L'ensemble de ces sous-évaluations constitue le cas central \og Évaluer les compétences comportementales \fg{}. L'employé consulte ensuite son profil et son rapport comportemental. L'Administrateur RH valide le profil et génère le rapport Soft Skills. Le Service de communication temps réel diffuse les notifications de manière optionnelle via une relation \og extend \fg{}.

\begin{figure}[H]
  \centering
  \includegraphics[width=0.9\textwidth]{images/D use case soft.png}
  \caption{Diagramme de cas d'utilisation : Évaluation et restitution des Compétences comportementales}
  \label{fig:usecase-soft-skills-sprint2}
\end{figure}

\begin{table}[H]
\centering
\small
\renewcommand{\arraystretch}{1.25}
\begin{tabularx}{\textwidth}{|P{0.26\textwidth}|Y|}
\hline
\textbf{Élément} & \textbf{Description} \\
\hline
Nom du cas & Évaluation et restitution du profil comportemental \\
\hline
Acteur principal & Employé \\
\hline
Acteurs secondaires & Administrateur RH, Service IA, Service GitHub, Service de communication temps réel \\
\hline
Pré-conditions & L'employé s'est authentifié et a accédé au module \og Compétences comportementales \fg{} de son espace d'évaluation. \\
\hline
Scénario nominal & \begin{minipage}[t]{\linewidth}
\begin{enumerate}[leftmargin=*,nosep]
  \item L'employé soumet son CV ; le système en extrait le ton et la sémantique via le Service IA.
  \item L'employé renseigne son nom d'utilisateur GitHub ; le Service GitHub et le Service IA analysent sa dynamique de collaboration.
  \item L'employé passe le test PCM ; les scores des six dimensions sont calculés par le Service IA.
  \item L'employé répond à un scénario métier immersif généré et évalué par le Service IA.
  \item Le système consolide les quatre sources et produit le profil comportemental avec le type de personnalité dominant.
  \item L'employé consulte son profil et son rapport comportemental. L'Administrateur RH valide le profil et génère le rapport Soft Skills.
\end{enumerate}
\end{minipage} \\
\hline
Scénario alternatif & Si l'employé ne termine pas le scénario professionnel, le système sauvegarde l'état actuel et le Service de communication temps réel programme une relance automatique. \\
\hline
Post-conditions & La cartographie de personnalité de l'employé est finalisée, permettant de prescrire les meilleurs leviers de communication et d'apprentissage. \\
\hline
\end{tabularx}
\renewcommand{\arraystretch}{1.05}
\caption{Fiche descriptive du cas d'utilisation : Évaluation et restitution du profil comportemental}
\end{table}

\textbf{3. Diagramme de cas d'utilisation : Module Compétences techniques}

Ce diagramme~\ref{fig:usecase-hard-skills-sprint2} présente les fonctionnalités du module d'évaluation technique. L'employé soumet son CV technique et renseigne son nom d'utilisateur GitHub, déclenchant respectivement l'extraction des compétences techniques du CV et l'analyse de l'activité GitHub. Il passe également un QCM technique adaptatif et réalise un challenge de code, tous deux supervisés par le Service IA. Le cas central \og Évaluer les compétences techniques \fg{} inclut l'ensemble de ces analyses ainsi que le calcul automatique du niveau technique. L'employé consulte ensuite son niveau, son score de correspondance au poste et son rapport technique. L'Administrateur RH valide le niveau certifié et génère le rapport Hard Skills.

\begin{figure}[H]
  \centering
  \includegraphics[width=0.9\textwidth]{images/D uses case hard.png}
  \caption{Diagramme de cas d'utilisation : Évaluation multi-sources des Compétences techniques}
  \label{fig:usecase-hard-skills-sprint2}
\end{figure}

\begin{table}[H]
\centering
\small
\renewcommand{\arraystretch}{1.25}
\begin{tabularx}{\textwidth}{|P{0.26\textwidth}|Y|}
\hline
\textbf{Élément} & \textbf{Description} \\
\hline
Nom du cas & Évaluation multi-sources et certification des Compétences techniques \\
\hline
Acteur principal & Employé \\
\hline
Acteurs secondaires & Administrateur RH, Service IA, Service GitHub, Service de communication temps réel \\
\hline
Pré-conditions & L'employé s'est authentifié et a accédé au module \og Compétences techniques \fg{}. \\
\hline
Scénario nominal & \begin{minipage}[t]{\linewidth}
\begin{enumerate}[leftmargin=*,nosep]
  \item L'employé soumet son CV technique ; le Service IA en extrait les compétences techniques.
  \item L'employé renseigne son nom d'utilisateur GitHub ; le Service GitHub et le Service IA analysent son activité et sa stack technique.
  \item L'employé passe un QCM technique adaptatif généré par le Service IA selon ses compétences détectées.
  \item L'employé réalise un challenge de code évalué par le Service IA sur quatre critères.
  \item Le Service IA calcule le niveau technique par triangulation multi-sources (\textit{Beginner} à \textit{Expert}).
  \item L'employé consulte son niveau technique, son score de correspondance au poste et son rapport technique.
  \item L'Administrateur RH valide le niveau certifié et génère le rapport Hard Skills.
\end{enumerate}
\end{minipage} \\
\hline
Scénario alternatif & Si le nom d'utilisateur GitHub est erroné ou introuvable, le système en avertit l'employé via le Service de communication temps réel et base le calcul du niveau uniquement sur le CV, le QCM et le challenge de code. \\
\hline
Post-conditions & Les compétences techniques sont mesurées, classées et le système est prêt à recommander des formations ciblées pour combler les lacunes identifiées (\textit{Upskilling}). \\
\hline
\end{tabularx}
\renewcommand{\arraystretch}{1.05}
\caption{Fiche descriptive du cas d'utilisation : Évaluation multi-sources et certification des Compétences techniques}
\end{table}

\subsubsection{Sources d'évaluation et dimensions mesurées}
L'approche multi-sources retenue croise données objectives, évaluations auto-déclarées et tests interactifs afin de produire un profil employé complet. Chaque source apporte un signal complémentaire irréductible aux autres, comme détaillé dans les tableaux des deux modules ci-dessous.

\textbf{Module compétences comportementales}

\begin{table}[H]
\centering
\renewcommand{\arraystretch}{1.3}
\begin{tabularx}{\TableWidth}{|P{2.4cm}|P{3.2cm}|Y|}
\hline
\textbf{Source} & \textbf{Technologie} & \textbf{Compétences comportementales évaluées} \\
\hline
GitHub & n8n + GitHub API + Ollama & Curiosité, collaboration, ownership et consistance \\
\hline
CV (texte brut) & n8n + Ollama & Indices de compétences comportementales, années d'expérience et niveau d'éducation \\
\hline
Test PCM (18 questions) & n8n (JavaScript heuristique, sans LLM) & Communication, discipline, curiosité, collaboration, ownership et leadership \\
\hline
Scénario IA & FastAPI + Ollama & Empathie, assertivité, pragmatisme et clarté de communication \\
\hline
\end{tabularx}
\caption{Sources d'évaluation du module compétences comportementales}
\label{tab:sources-soft-skills}
\end{table}

\textbf{Module compétences techniques}

\begin{table}[H]
\centering
\small
\renewcommand{\arraystretch}{1.3}
\begin{tabularx}{\textwidth}{|P{2.8cm}|P{4.0cm}|Y|}
\hline
\textbf{Source} & \textbf{Technologie} & \textbf{Ce qui est évalué} \\
\hline
GitHub & GitHub API REST + Ollama & Langages, frameworks détectés et niveau d'activité \\
\hline
CV (PDF/DOCX) & pdfplumber + regex sur 150+ mots-clés & Technologies mentionnées et années d'expérience \\
\hline
LinkedIn & Agent Ollama & Rôle actuel, arc de carrière, certifications et cohérence CV/GitHub \\
\hline
Test QCM & Ollama (llama3.2) & Maîtrise technique par compétence selon exactitude, confiance et temps \\
\hline
Challenge de code & CodeLlama & Correctness (40), code quality (30), efficiency (20) et readability (10) \\
\hline
\end{tabularx}
\caption{Sources d'évaluation du module compétences techniques}
\label{tab:sources-hard-skills}
\end{table}



\subsubsection{4.5.2 Conception}

Cette phase a pour but de traduire les exigences fonctionnelles et les cas d’utilisation identifiés précédemment en solutions techniques concrètes. L’enjeu majeur de ce deuxième sprint réside dans la conception d’une architecture capable d’ingérer et de croiser des données hétérogènes (fichiers PDF, données d'API tierces, réponses dynamiques) via une intelligence artificielle locale, tout en garantissant la fiabilité et la vélocité de l'évaluation globale.

\paragraph*{1. Objectifs de conception}
\begin{itemize}[leftmargin=1cm]
    \item \textbf{Découplage et Modularité :} Isoler les traitements cognitifs lourds (\textit{parsing} sémantique, inférence LLM) dans un microservice dédié (FastAPI) et confier l'orchestration des flux de données à \texttt{n8n} afin d'éviter la saturation du cœur transactionnel (Spring Boot).
    \item \textbf{Fiabilité Multi-Sources :} Concevoir des algorithmes capables de corroborer de manière robuste les données déclaratives (CV, profil LinkedIn) avec des preuves factuelles (historique de \textit{commits} GitHub, résultats aux QCM et aux défis de code).
    \item \textbf{Asynchronisme et Réactivité :} Mettre en œuvre des flux d'exécution hautement asynchrones, couplés au Communication Hub (SSE), pour garantir que les temps d'inférence de l'IA (Ollama) n'impactent pas la fluidité et l'expérience utilisateur sur le \textit{frontend}.
\end{itemize}

\paragraph*{2. Architecture globale du module}
L'architecture repose sur quatre  couches aux responsabilités strictement séparées : Angular collecte les données et restitue les résultats ; n8n orchestre les workflows compétences comportementales via GitHub API et Ollama ; FastAPI gère le module compétences techniques et les scénarios comportementaux avec llama3.2 et CodeLlama ; Spring Boot assure la persistance PostgreSQL, la sécurité JWT et le proxy vers FastAPI ; l'API GitHub constitue la source externe commune aux deux modules.
\begin{figure}[H]
  \centering
  \includegraphics[width=0.9\textwidth]{images/archiglob.png}
  \caption{Architecture globale du module d'évaluation des compétences }
  \label{fig:architecture-globale}
\end{figure}

%\figureplaceholder{Architecture globale du module d'évaluation des compétences (soft + compétences techniques)}{fig:architecture-globale}

L'architecture repose sur un modèle hybride orchestré autour de quatre couches technologiques :

\begin{table}[H]
\centering
\small
\renewcommand{\arraystretch}{1.25}
\begin{tabularx}{\TableWidth}{|P{2.8cm}|P{3.4cm}|Y|}
\hline
\textbf{Couche} & \textbf{Technologie} & \textbf{Rôle} \\
\hline
Orchestration IA & n8n (workflows visuels) & Analyse GitHub soft, extraction CV, calcul PCM et synthèse finale \\
\hline
API IA comportementale & FastAPI + Ollama & Scénarios IA, analyse GitHub/CV/LinkedIn compétences techniques, QCM et code challenge \\
\hline
Backend métier & Spring Boot (Java 21) & Persistance PostgreSQL, sécurité JWT, API REST et rapports PDF \\
\hline
Frontend & Angular 18~\cite{angular}+ & Interface employé/admin/RH, visualisations et restitution des résultats \\
\hline
\end{tabularx}
\caption{Couches technologiques de l'architecture hybride}
\label{tab:couches-architecture-hybride}
\end{table}
L'architecture retenue repose sur un modèle hybride structuré autour de deux pipelines parallèles.





%\figureplaceholder{Diagramme architecture globale --- Vue complète des couches Frontend, n8n, FastAPI, Spring Boot, PostgreSQL, Ollama/CodeLlama avec flux de données}{fig:architecture-globale-couches-flux}

\subsubsection{Architecture n8n  Quatre workflows}



Le pipeline compétences comportementales s'articule autour du Master Agent n8n, qui déclenche
en parallèle via Promise.all() trois sous-workflows : CV Parser (analyse sémantique Ollama),
GitHub Analysis (score heuristique JavaScript + évaluation Ollama) et PCM Test (calcul algorithmique
pur sans LLM). Les résultats sont agrégés selon la pondération CV~40\%, PCM~30\%, GitHub~30\%,
avec basculement automatique vers CV~60\%/PCM~40\% si GitHub est indisponible. Ollama produit
ensuite le profil final (score global, type de personnalité, forces, recommandations). FastAPI
complète l'évaluation via un scénario IA sur quatre dimensions psychologiques en temps réel.
\begin{figure}[H]
  \centering
  \includegraphics[width=0.9\textwidth]{images/piplinesoft1.png}
   \caption{Pipeline d'évaluation des compétences comportementales}
  \label{fig:pipeline-soft-skills}
\end{figure}
Ce pipeline~\ref{fig:pipeline-soft-skills} combine les résultats du CV, du test PCM, de l'analyse GitHub et du scénario
comportemental IA afin de produire un profil soft skills unifié. Une étape de validation des
données en amont garantit la présence du CV, de l'identifiant GitHub et des réponses PCM avant
le lancement des workflows. Les scores agrégés sont ensuite sauvegardés en base de données et
restitués à l'interface utilisateur.


\color{red}
\begin{itemize}
    \item \textit{Diagramme 9 :} Le diagramme présenté n’est pas un diagramme de séquence UML ; il ressemble plutôt à un diagramme d’activité ou à un pipeline de traitement.
    \item \textit{Diagramme 9 :} Renommer la figure en \textit{Pipeline d’évaluation des compétences comportementales} si vous gardez cette représentation sous forme de workflow.
    \item \textit{Diagramme 9 :} Si vous souhaitez conserver le titre \textit{Diagramme de séquence}, il faut reconstruire le diagramme avec des lignes de vie verticales, des messages horizontaux et un ordre chronologique clair.
    \item \textit{Diagramme 9 :} Ajouter les lignes de vie nécessaires pour un vrai diagramme de séquence : \textit{Employé}, \textit{Frontend}, \textit{Backend}, \textit{n8n Master Agent}, \textit{Workflow CV Parser}, \textit{Workflow GitHub Analysis}, \textit{Workflow PCM Test}, \textit{FastAPI}, \textit{Ollama}, \textit{GitHub API} et \textit{Base de données}.
    \item \textit{Diagramme 9 :} Remplacer les blocs de traitement par des messages UML explicites, par exemple \textit{soumettreProfil()}, \textit{lancerAnalyseComportementale()}, \textit{extraireSoftSkillsCV()}, \textit{analyserGitHub()}, \textit{calculerScorePCM()} et \textit{agrégerRésultats()}.
    \item \textit{Diagramme 9 :} Représenter l’exécution parallèle des workflows CV, GitHub et PCM avec un fragment combiné \texttt{par}, au lieu d’utiliser uniquement des flèches en pointillés.
    \item \textit{Diagramme 9 :} Représenter le cas \textit{GitHub disponible ?} avec un fragment \texttt{alt} contenant deux branches claires : \textit{GitHub disponible} et \textit{GitHub non disponible}.
    \item \textit{Diagramme 9 :} Clarifier la pondération des scores : le bloc indique \textit{CV 60\% + PCM 40\%}, alors que l’agrégation indique \textit{CV 40\% + PCM 30\% + GitHub 30\%}. Il faut choisir une seule règle cohérente.
    \item \textit{Diagramme 9 :} Éviter les détails trop techniques dans les blocs, comme \textit{temp: 0.7} ou \textit{temp: 0.2} ; ces paramètres doivent être expliqués dans le texte technique, pas dans le diagramme principal.
    \item \textit{Diagramme 9 :} Remplacer les libellés mixtes français/anglais par des intitulés uniformes en français, par exemple \textit{Nœud maître d’analyse IA}, \textit{Score global}, \textit{Type de personnalité}, \textit{Forces principales} et \textit{Recommandations}.
    \item \textit{Diagramme 9 :} Ajouter une étape de validation des données avant le lancement des workflows, notamment pour vérifier la présence du CV, du nom d’utilisateur GitHub et des réponses PCM.
    \item \textit{Diagramme 9 :} Ajouter une étape de stockage des résultats dans la base de données après l’agrégation des scores et la génération du profil soft skills.
    \item \textit{Diagramme 9 :} Ajouter une étape de retour vers l’interface utilisateur afin de montrer comment l’employé consulte le profil comportemental unifié.
    \item \textit{Diagramme 9 :} Clarifier le rôle de \textit{FastAPI} : il doit apparaître comme une ligne de vie ou un composant appelé par n8n, et non seulement comme un texte dans un bloc.
    \item \textit{Diagramme 9 :} Clarifier le rôle d’\textit{Ollama} : il doit être représenté comme un service IA appelé pour l’analyse sémantique, la synthèse et l’évaluation du scénario.
    \item \textit{Diagramme 9 :} Corriger les flèches de retour afin de distinguer les appels, les réponses et les résultats intermédiaires.
    \item \textit{Diagramme 9 :} Réduire les longues boucles en pointillés, car elles rendent le diagramme difficile à lire et ne correspondent pas à une notation UML standard.
    \item \textit{Diagramme 9 :} Ajouter un fragment \texttt{opt} pour le scénario IA comportemental si cette évaluation est optionnelle ou déclenchée seulement dans certains cas.
    \item \textit{Diagramme 9 :} Ajouter une phrase avant la figure pour préciser si elle représente un pipeline de traitement ou un diagramme de séquence.
    \item \textit{Diagramme 9 :} Ajouter une phrase après la figure pour expliquer que l’évaluation comportementale combine les résultats du CV, du test PCM, de GitHub et du scénario IA afin de produire un profil soft skills unifié.
\end{itemize}
\color{black}
\subsubsection{Architecture FastAPI Structure du code}

talent\_agent.run\_agent() orchestre en parallèle trois sources :
l'analyseur GitHub (API GitHub REST), l'analyseur CV (extraction de plus de 150~technologies
via pdfplumber et expressions régulières) et un agent Ollama pour l'analyse LinkedIn.
Les PDF scannés déclenchent une analyse partielle avec avertissement transmis au frontend.
Le module de fusion et déduplication des compétences normalise les résultats via plus de
500~alias canoniques\,; le moteur de calcul du niveau de maîtrise attribue un niveau par
triangulation multi-sources (Débutant, Intermédiaire, Avancé, Expert)\,; le moteur de
correspondance au poste compare le profil avec plus de 20~profils métiers. Des tests interactifs
optionnels (QCM adaptatif Ollama, challenge de code CodeLlama) enrichissent le profil final
sans bloquer la génération des résultats. L'ensemble est persisté en PostgreSQL via Spring Boot
sous l'entité Skill de type TECH.

\begin{figure}[H]
  \centering
  \includegraphics[width=0.9\textwidth]{images/piplinehard1.png}
  \caption{Pipeline d'évaluation des compétences techniques}
  \label{fig:pipeline-hard-skills}
\end{figure}
Ce pipeline~\ref{fig:pipeline-hard-skills} plusieurs sources d'information (GitHub, CV, LinkedIn), applique une fusion
et une déduplication des compétences, calcule un niveau de maîtrise par triangulation et produit
un profil final enrichi par des tests interactifs optionnels. Les sorties principales sont les
compétences détectées avec leur niveau estimé, le score d'expérience, le score de correspondance
au poste cible et les recommandations associées.
\color{red}
\begin{itemize}
    \item \textit{Diagramme 10 :} Le schéma présenté n’est pas un diagramme de séquence UML ; il s’agit plutôt d’un pipeline de traitement ou d’un diagramme d’activité.
    \item \textit{Diagramme 10 :} Renommer la figure en \textit{Pipeline d’évaluation des compétences techniques} ou \textit{Workflow du module de compétences techniques}.
    \item \textit{Diagramme 10 :} Uniformiser la langue du diagramme, car les libellés mélangent le français et l’anglais.
    \item \textit{Diagramme 10 :} Remplacer les noms de scripts internes comme \texttt{github\_analyzer.py}, \texttt{cv\_parser.py}, \texttt{skill\_extractor.py}, \texttt{scoring\_engine.py} et \texttt{job\_matching.py} par des intitulés fonctionnels plus lisibles.
    \item \textit{Diagramme 10 :} Réduire les détails trop techniques dans les blocs, comme \textit{temp: 0.4}, \textit{6--12 questions}, \textit{500+ alias} ou les formules de scoring détaillées, et déplacer ces précisions dans le texte du rapport.
    \item \textit{Diagramme 10 :} Clarifier le rôle de LinkedIn, car il est cité dans l’entrée mais aucun bloc d’analyse dédié n’apparaît clairement dans le pipeline.
    \item \textit{Diagramme 10 :} Si LinkedIn est analysé par le bloc \textit{Agent Ollama}, le préciser explicitement dans le nom du bloc ou ajouter un bloc séparé pour éviter l’ambiguïté.
    \item \textit{Diagramme 10 :} Clarifier la signification des flèches en pointillés, car elles ne permettent pas de distinguer clairement les flux de données, les dépendances et les branches conditionnelles.
    \item \textit{Diagramme 10 :} Ajouter une légende expliquant la signification des flèches, des décisions et des traitements parallèles.
    \item \textit{Diagramme 10 :} Représenter le parallélisme avec une notation plus claire, par exemple en indiquant explicitement les trois branches parallèles : GitHub, CV et LinkedIn/Ollama.
    \item \textit{Diagramme 10 :} Corriger la décision \textit{PDF extractible ?} en ajoutant clairement les deux branches \textit{Oui} et \textit{Non}, avec un chemin de traitement explicite pour chaque cas.
    \item \textit{Diagramme 10 :} Mieux relier le bloc \textit{Warning frontend + analyse partielle} à la branche d’erreur correspondante afin de montrer quand cette alerte est déclenchée.
    \item \textit{Diagramme 10 :} Vérifier la cohérence des règles de nivellement dans \textit{scoring\_engine.py}, car les conditions affichées sont hétérogènes et difficiles à interpréter.
    \item \textit{Diagramme 10 :} Reformuler les niveaux techniques avec une logique homogène, par exemple en utilisant uniquement des intervalles de scores ou uniquement des critères de sources/occurrences.
    \item \textit{Diagramme 10 :} Clarifier l’enchaînement entre \textit{scoring\_engine} et \textit{job\_matching}, en précisant quelles données sont transmises d’un bloc à l’autre.
    \item \textit{Diagramme 10 :} Vérifier la cohérence entre le module de compétences techniques et l’entité stockée dans Spring Boot, car la mention \textit{SkillTech\{SOFT\}} semble contradictoire.
    \item \textit{Diagramme 10 :} Remplacer \textit{SkillTech\{SOFT\}} par une entité cohérente avec le module, par exemple \textit{SkillTech} ou \textit{HardSkillsProfile}.
    \item \textit{Diagramme 10 :} Clarifier si les tests interactifs sont réellement optionnels et, si oui, montrer qu’ils enrichissent le profil sans bloquer la génération du résultat final.
    \item \textit{Diagramme 10 :} Mieux montrer que les résultats du QCM adaptatif et du challenge de code alimentent le profil final ou le service Spring Boot.
    \item \textit{Diagramme 10 :} Simplifier le bloc final \textit{Profil Hard Skills Complet} en gardant uniquement les sorties métier principales : compétences détectées, niveau estimé, score d’expérience, score de job matching et recommandations.
    \item \textit{Diagramme 10 :} Ajouter une phrase avant la figure pour préciser qu’elle présente le pipeline global du module d’évaluation des compétences techniques.
    \item \textit{Diagramme 10 :} Ajouter une phrase après la figure pour expliquer que le pipeline combine plusieurs sources, applique une fusion des compétences, calcule un niveau technique et produit un profil final enrichi par des tests interactifs.
\end{itemize}
\color{black}
\subsubsection{Modèles de données}

\textbf{Modèle entité-relation Spring Boot / PostgreSQL unifié}

Le modèle de persistance PostgreSQL s'articule autour d'une entité centrale Utilisateur reliée à huit tables satellites par des relations many-to-one. Competence enregistre chaque compétence détectée (TECHNIQUE ou COMPORTEMENTALE) avec son niveau (DEBUTANT, INTERMEDIAIRE, AVANCE, EXPERT), sa source (CV, GITHUB, PCM, IA\_PYTHON), un score de confiance (0-100) et son statut de validation administrative via le booléen validee.

\texttt{RésultatTest} enregistre les résultats détaillés des différents types de tests (QCM, CODE, SCENARIO, PCM) avec leur score et leur statut (EN\_ATTENTE, TERMINE, ECHOUE). Prédiction conserve la synthèse globale du profil candidat avec les scores décomposés en trois axes : score\_global, score\_technique etscore\_comportemental, ainsi que le type de personnalité et le score d'adéquation emploi.

\texttt{ProfilComportemental} dédie une table séparée aux données de personnalité (MBTI, leadership, collaboration, adaptabilité, communication) pour une meilleure traçabilité. Recommandationstocke les suggestions de formation personnalisées avec leur priorité (BASSE, MOYENNE, HAUTE) et leur type (MONTEE\_EN\_COMPETENCES, RECONVERSION). 

\texttt{ParcoursFormation} suit la progression des parcours d'upskilling avec un statut (ACTIF, TERMINE, EN\_PAUSE) et un pourcentage de progression. \texttt{PasseportCompétence} historise les générations du passeport de compétences (PDF, JSON, XML) avec un contrôle de version.

Les entités \texttt{Invitation} et \texttt{Équipe} gèrent les interactions collaboratives entre utilisateurs. Cette architecture complète permet de reconstruire intégralement le profil employé depuis son identifiant unique, avec l'historique complet de ses apprentissages, ses évaluations, ses recommandations et ses parcours de formation.

\begin{figure}[H]
  \centering
  \includegraphics[width=0.85\textwidth]{images/modele entitee relation.png}
  \caption{Diagramme entité-relation du modèle de persistance unifié de TalentPredict AI}
  \label{fig:entite-relation-skill}
\end{figure}

\textbf{Traçabilité et intégrité des données}

Toutes les tables incluent les champs date\_creation et date\_modification pour assurer une traçabilité complète. Les clés étrangères sont explicitement nommées (id\_utilisateur, id\_competence, etc.) et les contraintes d'unicité sont appliquées sur nom\_utilisateur et email. Le mot de passe est stocké sous forme de hash\_mot\_de\_passe pour garantir la sécurité des données sensibles. Les relations many-to-one garantissent l'intégrité référentielle entre l'entité centrale \texttt{Utilisateur} et ses données associées.


\color{red}
\begin{itemize}
    \item \textit{Diagramme 11 :} Renommer la figure en \textit{Diagramme entité-relation du modèle de persistance} afin de garder un titre clair et académique.
    \item \textit{Diagramme 11 :} Ajouter un titre visible au-dessus du schéma, par exemple \textit{Modèle de persistance de TalentPredict AI}.
    \item \textit{Diagramme 11 :} Corriger les multiplicités entre \textit{User} et \textit{Skill} : un utilisateur peut posséder plusieurs compétences, donc la relation doit être \textit{User 1} vers \textit{Skill 0..*}.
    \item \textit{Diagramme 11 :} Corriger les multiplicités entre \textit{User} et \textit{TestResult} : un utilisateur peut avoir plusieurs résultats de tests, donc la relation doit être \textit{User 1} vers \textit{TestResult 0..*}.
    \item \textit{Diagramme 11 :} Corriger les multiplicités entre \textit{User} et \textit{Prediction} : un utilisateur peut avoir plusieurs prédictions dans le temps, donc la relation doit être \textit{User 1} vers \textit{Prediction 0..*}, sauf si le système conserve uniquement la dernière prédiction.
    \item \textit{Diagramme 11 :} Ajouter explicitement les clés étrangères \textit{user\_id} dans les relations, afin de montrer que \textit{Skill}, \textit{TestResult} et \textit{Prediction} dépendent de \textit{User}.
    \item \textit{Diagramme 11 :} Ajouter les cardinalités de manière lisible près de chaque extrémité des relations, car elles sont actuellement peu visibles.
    \item \textit{Diagramme 11 :} Remplacer l’attribut \textit{password} s’il existe dans le modèle réel par \textit{password\_hash}, afin d’indiquer que le mot de passe n’est pas stocké en clair.
    \item \textit{Diagramme 11 :} Ajouter les attributs importants de la table \textit{User}, comme \textit{first\_name}, \textit{last\_name}, \textit{phone}, \textit{is\_active}, \textit{created\_at} et \textit{updated\_at}, si ces champs sont utilisés dans l’application.
    \item \textit{Diagramme 11 :} Vérifier si l’attribut \textit{username} est unique ; si oui, indiquer clairement la contrainte \textit{unique}.
    \item \textit{Diagramme 11 :} Vérifier si l’attribut \textit{email} est unique ; si oui, indiquer clairement la contrainte \textit{unique}.
    \item \textit{Diagramme 11 :} Dans la table \textit{Skill}, remplacer \textit{nom} par \textit{name} ou uniformiser toute la base en français ; il faut éviter le mélange entre français et anglais.
    \item \textit{Diagramme 11 :} Dans la table \textit{Skill}, vérifier si l’attribut \textit{niveau} doit être un entier simple ou une énumération, par exemple \textit{Beginner}, \textit{Intermediate}, \textit{Advanced}, \textit{Expert}.
    \item \textit{Diagramme 11 :} Dans la table \textit{Skill}, remplacer \textit{validee} par \textit{validated} ou garder \textit{validée}, selon la langue choisie pour tout le modèle.
    \item \textit{Diagramme 11 :} Ajouter un attribut \textit{confidence\_score} dans \textit{Skill} si le système calcule un niveau de confiance à partir des sources CV, GitHub, PCM ou IA.
    \item \textit{Diagramme 11 :} Dans la table \textit{TestResult}, remplacer \textit{skill\_name} par \textit{skill\_id} si le résultat doit être relié directement à une compétence existante.
    \item \textit{Diagramme 11 :} Ajouter une relation entre \textit{TestResult} et \textit{Skill} si chaque résultat de test concerne une compétence précise.
    \item \textit{Diagramme 11 :} Dans la table \textit{TestResult}, vérifier si \textit{test\_type} doit contenir aussi \textit{PCM}, car le projet évalue à la fois les compétences techniques et comportementales.
    \item \textit{Diagramme 11 :} Dans la table \textit{Prediction}, clarifier le rôle de \textit{overall\_score} : préciser s’il représente le score global, le score technique, le score comportemental ou un score combiné.
    \item \textit{Diagramme 11 :} Vérifier si \textit{personality\_type} doit être placé dans \textit{Prediction} ou dans une table séparée dédiée au profil comportemental.
    \item \textit{Diagramme 11 :} Ajouter une table \textit{Recommendation} si les recommandations de formation sont stockées séparément et consultées par l’utilisateur.
    \item \textit{Diagramme 11 :} Ajouter une table \textit{TrainingPath} ou \textit{Formation} si le système gère des parcours de formation ou des recommandations d’upskilling.
    \item \textit{Diagramme 11 :} Ajouter une table \textit{TalentPassport} si le passeport de compétences est généré, historisé ou téléchargé comme document.
    \item \textit{Diagramme 11 :} Ajouter les champs \textit{created\_at} et \textit{updated\_at} dans toutes les tables principales pour assurer la traçabilité des données.
    \item \textit{Diagramme 11 :} Ajouter les contraintes principales, notamment \textit{pk}, \textit{fk}, \textit{unique}, \textit{not null} et les valeurs possibles des énumérations.
    \item \textit{Diagramme 11 :} Harmoniser le style des noms des tables et attributs : utiliser soit le style anglais \textit{snake\_case}, soit le style français, mais ne pas mélanger les deux.
    \item \textit{Diagramme 11 :} Ajouter une phrase avant la figure pour expliquer que le diagramme présente les principales entités persistées par la plateforme.
    \item \textit{Diagramme 11 :} Ajouter une phrase après la figure pour expliquer que le modèle relie chaque utilisateur à ses compétences, ses résultats de tests et ses prédictions globales.
\end{itemize}
\color{black}

\subsubsection{Diagrammes de séquence}

\textbf{1. Diagramme de Séquence : Évaluation des Compétences Comportementales}
Ce diagramme de séquence modélise le flux d'exécution asynchrone régissant l'évaluation comportementale complète du collaborateur. L'architecture met en exergue l'orchestration des microservices via n8n. Le processus débute par la collecte d'informations non structurées : le texte brut du CV est extrait sémantiquement via FastAPI, et l'historique de collaboration (messages de commit, dynamiques d'équipe sur les Pull Requests) est récupéré via l'API GitHub.
Parallèlement, le collaborateur répond au questionnaire heuristique PCM. L'orchestrateur n8n agrège ces trois vecteurs de données textuelles et les soumet à l'IA locale (Ollama) pour générer un scénario immersif hyper-personnalisé. Les réactions de l'employé face à cette mise en situation sont ensuite analysées pour déduire les six dimensions cognitives de son profil. La persistance de ces métriques déclenche un événement de propagation (Server-Sent Events) par l'intermédiaire du Communication Hub, assurant ainsi une notification fluide sur l'interface utilisateur.

\begin{figure}[H]
  \centering
  \includegraphics[width=0.9\textwidth]{images/sequencesoft2.png}
  \caption{ Diagramme de sequence d'Évaluation des Compétences Comportementales}
  \label{fig:seq_soft}
\end{figure}

\textbf{2. Diagramme de Séquence : Évaluation des compétences technique} 
Ce second diagramme de séquence formalise le paradigme d'évaluation multi-sources implémenté pour objectiver l'expertise technique. Le flux d'exécution débute par la délégation de l'extraction documentaire (CV) au microservice Python (FastAPI). En parallèle, le noyau Spring Boot interroge l'API GitHub pour corroborer les compétences déclarées avec des traces d'activité réelles.

L'agrégation de ces vecteurs d'information permet au système de solliciter l'IA (CodeLlama/Ollama) pour instancier un environnement d'évaluation adaptatif en deux phases : un QCM sur-mesure pour valider les connaissances théoriques, suivi d'un challenge de code pour éprouver la logique pratique. La correction automatisée de cette double soumission par le LLM permet le calcul d'un Score de Confiance global. Ce triptyque de validation (Déclaratif, Factuel, Évaluatif) est persistant en base de données, et l'utilisateur est instantanément notifié de son nouveau niveau d'expertise via le Communication Hub.
\begin{figure}[H]
  \centering
  \includegraphics[width=0.9\textwidth]{images/sequencehard2.png}
  \caption{Diagramme de sequence d'evaluation des competences technique}
  \label{fig:seq_hard}
\end{figure}

\subsubsection{Diagramme de classes  Module scénarios IA }

Le microservice FastAPI s'articule autour de cinq classes Pydantic principales : ScenarioRequest (entrées rôle/niveau), ScenarioResponse (titre, description, compétences testées), EvaluationResult (quatre scores psychologiques, forces, axes d'amélioration, profil culture-add), EmployéeAnalysis (données consolidées GitHub/CV/LinkedIn) et Skill (attributs name, level, score, sources), partagée entre les deux modules. Toutes les entrées sont validées automatiquement par Pydantic avant traitement.
\begin{figure}[H]
  \centering
  \includegraphics[width=0.7\textwidth]{images/classes_fastapi.png}
  \caption{Diagramme de classes du module scénarios IA et compétences techniques (FastAPI) }
  \label{fig:classes-fastapi}
\end{figure}
%\figureplaceholder{Diagramme de classes du module scénarios IA et compétences techniques (FastAPI)}{fig:classes-fastapi}


\subsection{Codage}
Cette section documente la phase d'implémentation effective du sprint. Elle présente les technologies déployées, l'organisation du code source, et les bonnes pratiques appliquées pour garantir la maintenabilité et la lisibilité du système. Chaque composant clé est détaillé avec des extraits de code commentés illustrant les choi-x techniques retenus.
\subsubsection{Stack technologique }

{\footnotesize
\setlength{\tabcolsep}{3pt}
\begin{table}[H]
\centering
\footnotesize
\setlength{\tabcolsep}{3pt}
\begin{tabular}{|P{2.3cm}|P{2.2cm}|P{1.2cm}|P{4.3cm}|}
\hline
\textbf{Composant} & \textbf{Technologie} & \textbf{Version} & \textbf{Justification} \\
\hline
Orchestration workflows & n8n & latest & Interface visuelle, modification des prompts sans recompilation \\
\hline
LLM local généraliste & Ollama (llama3.2) & latest & Déploiement local et conformité RGPD \\
\hline
LLM spécialisé code & CodeLlama & latest & Meilleures performances pour la génération et l'évaluation de code \\
\hline
API IA & FastAPI & 0.110+ & Support asynchrone et auto-documentation OpenAPI \\
\hline
Extraction PDF client & pdfjs-dist & 5.6.205 & Extraction navigateur pour le module compétences comportementales \\
\hline
Extraction PDF serveur & pdfplumber & 0.10+ & Parsing précis pour le module compétences techniques \\
\hline
Backend métier & Spring Boot + Java 21 & 3.x & API REST, persistance JPA et sécurité JWT \\
\hline
Base de données & PostgreSQL & 15+ & Stockage \texttt{Skill}, \texttt{User}, \texttt{Prediction} et \texttt{TestResult} \\
\hline
Conteneurisation & Docker Compose & 24+ & Isolation des services et reproductibilité \\
\hline
\end{tabular}
\caption{Stack technologique du Sprint 2}\label{tab:stack-technologique}
\end{table}
}

\subsubsection{Implémentation des workflows n8n}

\paragraph{Workflow 1 : GitHub Analysis}
L'objectif de ce workflow est d'analyser l'activité open source du employé afin de déduire quatre compétences comportementales objectivables : curiosité, collaboration, ownership et consistance.
\begin{figure}[H]
  \centering
  \includegraphics[width=0.9\textwidth]{images/workflow-github.png}
  \caption{Capture du workflow GitHub Analysis dans n8n}
  \label{fig:workflow-github-analysis}
\end{figure}
%\figureplaceholder{Capture du workflow GitHub Analysis dans n8n}{fig:workflow-github-analysis}

\textbf{Nœud \og Compute GitHub Score \fg{}}

Le score GitHub est calculé de manière purement heuristique en JavaScript, sans appel LLM, afin de garantir rapidité et reproductibilité. Le calcul repose sur trois indicateurs pondérés selon le tableau ci-dessous.

\begin{table}[H]
\centering
\small
\renewcommand{\arraystretch}{1.25}
\begin{tabularx}{\TableWidth}{|P{3.0cm}|Y|P{2.1cm}|}
\hline
\textbf{Indicateur} & \textbf{Formule} & \textbf{Poids maximal} \\
\hline
Nombre de dépôts & \texttt{Math.min((repoCount / 20) $\times$ 40, 40)} & 40 pts \\
\hline
Nombre d'étoiles & \texttt{Math.min((totalStars / 50) $\times$ 40, 40)} & 40 pts \\
\hline
Nombre de forks & \texttt{Math.min((totalForks / 20) $\times$ 20, 20)} & 20 pts \\
\hline
\textbf{Score final} & \texttt{Math.round(repoScore + starScore + forkScore)} & \textbf{100 pts} \\
\hline
\end{tabularx}
\caption{Calcul heuristique du score GitHub}
\label{tab:compute-github-score}
\end{table}

%\figureplaceholder{Capture du code JavaScript du nœud Compute GitHub Score}{fig:compute-github-score}

Le nœud \og Analyze Compétences comportementales \fg{} soumet ensuite les données à Ollama à travers un prompt structuré centré sur les variables de contribution et de diversité technique.


\begin{lstlisting}
 Tu es un expert RH chargé d'analyser le profil GitHub d'un développeur. À partir des données suivantes : \{repo\_count\} dépôts, \{total\_stars\} étoiles, \{total\_forks\} forks et langages utilisés : \{languages\}, évalue uniquement les quatre compétences comportementales objectivement mesurables via GitHub. Retourne exclusivement un JSON valide contenant : curiosity (0--100), collaboration (0--100), ownership (0--100) et consistency (0--100).
\end{lstlisting}


%\figureplaceholder{Capture du nœud Ollama --- prompt Analyze Compétences comportementales}{fig:prompt-github-soft}
%\figureplaceholder{Exemple de sortie JSON du workflow GitHub Analysis}{fig:output-github-analysis}

\paragraph{Workflow 2 : CV Parser}
Ce workflow a pour objectif d'extraire les indices de compétences comportementales depuis le texte brut du CV via une analyse sémantique Ollama.

%\figureplaceholder{Capture du workflow CV Parser dans n8n}{fig:workflow-cv-parser}
\begin{figure}[H]
  \centering
  \includegraphics[width=0.85\textwidth]{images/workflow_cv.png}
  \caption{Capture du workflow CV Parser dans n8n}
  \label{fig:workflow-cv-parser}
\end{figure}

\begin{lstlisting}
Tu es un expert RH analysant un CV. Extrais uniquement les signaux de compétences comportementales depuis ce texte. Retourne uniquement un JSON valide avec : behavioral_skills_hints, experience_years, education_level et cv_score.
\end{lstlisting}


%\figureplaceholder{Capture du nœud Ollama --- prompt CV Parser}{fig:prompt-cv-parser}
%\figureplaceholder{Exemple de sortie JSON du workflow CV Parser}{fig:output-cv-parser}

\paragraph{Workflow 3 : PCM Test}
L'objectif est de calculer les scores PCM sur six dimensions à partir de 18 réponses auto-évaluées, sans recours à un LLM.

%\figureplaceholder{Capture du workflow PCM Test dans n8n}{fig:workflow-pcm}
\begin{figure}[H]
  \centering
  \includegraphics[width=0.8\textwidth]{images/workflow_pcm.png}
  \caption{Capture du workflow PCM Test dans n8n}
  \label{fig:workflow-pcm}
\end{figure}
\textbf{Nœud \og Score PCM Answers \fg{}}

Le scoring PCM est entièrement algorithmique. Chaque dimension est calculée comme la moyenne arrondie à une décimale de ses trois questions associées, puis le score global correspond à la moyenne des six dimensions.

\begin{table}[H]
\centering
\small
\renewcommand{\arraystretch}{1.25}
\begin{tabularx}{\TableWidth}{|P{2.6cm}|P{2.2cm}|Y|}
\hline
\textbf{Dimension} & \textbf{Questions} & \textbf{Formule} \\
\hline
Communication & q1, q2, q3 & \texttt{Math.round((q1+q2+q3) / 3 $\times$ 10) / 10} \\
\hline
Discipline & q4, q5, q6 & \texttt{Math.round((q4+q5+q6) / 3 $\times$ 10) / 10} \\
\hline
Curiosité & q7, q8, q9 & \texttt{Math.round((q7+q8+q9) / 3 $\times$ 10) / 10} \\
\hline
Collaboration & q10, q11, q12 & \texttt{Math.round((q10+q11+q12) / 3 $\times$ 10) / 10} \\
\hline
Ownership & q13, q14, q15 & \texttt{Math.round((q13+q14+q15) / 3 $\times$ 10) / 10} \\
\hline
Leadership & q16, q17, q18 & \texttt{Math.round((q16+q17+q18) / 3 $\times$ 10) / 10} \\
\hline
\textbf{Score global} & Toutes & Moyenne des six dimensions, arrondie à une décimale \\
\hline
\end{tabularx}
\caption{Calcul algorithmique des scores PCM}
\label{tab:score-pcm-answers}
\end{table}

%\figureplaceholder{Capture du code JavaScript du nœud Compute PCM Scores}{fig:compute-pcm-scores}
%\figureplaceholder{Exemple de sortie JSON du workflow PCM Test}{fig:output-pcm}

\paragraph{Workflow 4 : Master Compétences comportementales Agent}
Ce workflow orchestre les trois sous-workflows en parallèle puis produit une synthèse pondérée via Ollama.

%\figureplaceholder{Capture du workflow Master Compétences comportementales Agent dans n8n}{fig:workflow-master-agent}
\begin{figure}[H]
  \centering
  \includegraphics[width=0.9\textwidth]{images/workflow_master.png}
  \caption{Capture du workflow Master Compétences comportementales Agent dans n8n}
  \label{fig:workflow-master-agent}
\end{figure}

\textbf{Nœud \og Prepare \& Call Sub-Workflows \fg{}}

L'orchestration des trois sous-workflows est réalisée en parallèle via \texttt{Promise.all}, ce qui réduit la latence globale du Master Agent en évitant les appels séquentiels. Le tableau ci-dessous détaille les paramètres transmis à chaque sous-workflow.

\begin{table}[H]
\centering
\small
\renewcommand{\arraystretch}{1.25}
\begin{tabularx}{\TableWidth}{|P{3.0cm}|P{3.2cm}|Y|}
\hline
\textbf{Sous-workflow} & \textbf{Endpoint appelé} & \textbf{Paramètres transmis} \\
\hline
CV Parser & \texttt{behavioral-}\newline\texttt{skills-}\newline\texttt{extract} & \texttt{extracted\_cv\_text}, \texttt{full\_name}, \texttt{email} \\
\hline
PCM Test & \texttt{pcm-test} & \texttt{full\_name}, \texttt{q1} à \texttt{q18} \\
\hline
GitHub Analysis & \texttt{github-analysis} & \texttt{github\_username} \\
\hline
\end{tabularx}
\caption{Paramètres transmis aux sous-workflows du Master Agent}
\label{tab:prepare-call-subworkflows}
\end{table}

Les trois appels sont déclenchés simultanément via \texttt{await Promise.all([...])} et leurs résultats sont ensuite agrégés avant d'être soumis au nœud d'analyse Ollama pour la synthèse pondérée finale.

%\figureplaceholder{Capture du code JavaScript du nœud Prepare \& Call Sub-Workflows}{fig:prepare-call-subworkflows}

Le nœud \og Master AI Analysis \fg{} envoie les résultats agrégés à Ollama afin de produire un profil comportemental unifié avec score global, type de personnalité, compétences consolidées, forces et recommandations.

Prompt Ollama clé Master Agent synthèse :

\begin{lstlisting}
You are an expert in organizational psychology. Synthesize this multi-source behavioral data.
Data sources and weights:
- GitHub (30\%): behavioral signals from open-source activity
- CV (40\%): behavioral signals from professional background
- PCM (30\%): self-assessed personality dimensions
Return ONLY valid JSON:
{"overall_score": <0-100>, "personality_type": "<French>", 
 "merged_behavioral_skills": {...}, "key_strengths": [...], 
 "training_recommendations": "..."}
\end{lstlisting}
Output exemple :
\begin{lstlisting}
{
    "overall_score": 72,
    "personality_type": "Innovateur",
    "skills": [
        {
            "name": "communication",
            "score": 80,
            "source": "pcm"
        },
        {
            "name": "problem_solving",
            "score": 90,
            "source": "github"
        },
        {
            "name": "adaptability",
            "score": 70,
            "source": "cv"
        },
        {
            "name": "leadership",
            "score": 60,
            "source": "pcm"
        },
        {
            "name": "teamwork",
            "score": 50,
            "source": "github"
        }
    ],
    "merged_behavioral_skills": {
        "communication": 80,
        "leadership": 60,
        "teamwork": 50,
        "adaptability": 70,
        "problem_solving": 90
    },
    "source_data": {
        "cv_score": 75,
        "pcm_score": 85,
        "github_score": 92
    },
    "summary": "",
    "key_strengths": [],
    "training_recommendations": ""
}
\end{lstlisting}

%\figureplaceholder{Capture du nœud Ollama --- prompt Master AI Analysis}{fig:prompt-master-agent}

%\figureplaceholder{Exemple de sortie JSON du Master Agent}{fig:output-master-agent}


\subsubsection{Implémentation FastAPI}

Le microservice FastAPI centralise l'ensemble de la logique d'évaluation comportementale et technique en deux catégories de composants.

Les services à appel LLM sollicitent Ollama avec des températures calibrées selon le degré de créativité requis : \texttt{generate\_soft\_skills\_scenario} (température 0.7) produit un dilemme professionnel contextualisé par rôle et niveau, tandis que \texttt{evaluate\_scenario\_response} (température 0.2) évalue la réponse libre du employé sur quatre dimensions psychologiques. \texttt{analyze\_github\_profile} (température 0.3) enrichit les statistiques de dépôts publics par une synthèse sémantique Ollama.

Les services à logique purement algorithmique n'effectuent aucun appel LLM, garantissant des temps de réponse inférieurs à 200 ms et une reproductibilité totale : \texttt{analyze\_cv} extrait 150+ technologies par expressions régulières depuis un PDF ou DOCX, \texttt{extract\_skills} normalise les compétences issues des trois sources via un dictionnaire de 500+ alias, \texttt{\_per\_question\_weight} calcule le poids par réponse selon la formule exactitude $\times$ confiance $\times$ temps, \texttt{\_determine\_level} attribue un niveau Beginner $\rightarrow$ Expert par triangulation multi-sources, et \texttt{match\_job} produit un score de correspondance pour 20+ profils métiers.


\subsubsection{Bonnes pratiques appliquées}

Quatre bonnes pratiques transversales ont guidé l'implémentation des deux modules afin de garantir performance, fiabilité et maintenabilité du pipeline.

\textbf{1. Séparation calcul heuristique / analyse LLM}

Le LLM n'est sollicité que lorsque l'analyse sémantique est strictement nécessaire. Toute métrique déterministe est traitée sans appel Ollama : le score GitHub et les dimensions PCM sont calculés en JavaScript pur ($< 100$ ms), la détection de 150+ technologies par regex et la déduplication de 500+ alias en Python ($< 200$ ms), et le scoring QCM par algorithme pondéré ($< 100$ ms). Seules l'analyse sémantique GitHub, la synthèse Master Agent et l'évaluation de scénario sollicitent Ollama (5--25 s). Cette séparation a réduit de 60\% le nombre d'appels LLM par rapport à une approche tout-Ollama.

\textbf{2. Nettoyage systématique des réponses JSON Ollama}

Le modèle \texttt{llama3.2} produit fréquemment des réponses enveloppées dans des blocs Markdown ou contenant des artefacts textuels. Un mécanisme de nettoyage en quatre étapes est appliqué systématiquement dans tous les workflows n8n et services Python : suppression des balises Markdown, tentative de parsing direct via \texttt{JSON.parse()}, extraction du premier objet JSON valide par regex \texttt{\textbackslash\{[\textbackslash s\textbackslash S]*\textbackslash\}}, puis retour d'un fallback sécurisé avec scores à 0 si le parsing demeure impossible. Ce mécanisme garantit qu'aucun appel Ollama ne bloque le pipeline.

\textbf{3. Gestion des timeouts et fallbacks}

Chaque appel externe est encapsulé dans un mécanisme de protection à deux niveaux : un timeout explicite adapté à la complexité de la tâche et un fallback heuristique immédiat en cas d'échec. Les valeurs retenues vont de 10 s (GitHub API) à 600 s (CodeLlama), chaque échec déclenchant automatiquement une réponse de substitution exploitable --- scénario prédéfini, banque de questions fallback ou \texttt{starter\_code} générique. La fonction \texttt{call\_ollama\_json} centralise cette logique dans FastAPI en interceptant toute exception \texttt{httpx.TimeoutException} ou erreur réseau.

\textbf{4. Déduplication des compétences --- Dictionnaire d'alias}

La fusion de sources hétérogènes (GitHub, CV, LinkedIn) génère des doublons lorsqu'une même technologie est désignée par des noms différents. Le dictionnaire \texttt{\_ALIASES} dans \texttt{skill\_extractor.py} normalise 500+ variantes vers leur forme canonique --- par exemple \texttt{reactjs} $\rightarrow$ \texttt{React}, \texttt{k8s} $\rightarrow$ \texttt{Kubernetes}, \texttt{tf} $\rightarrow$ \texttt{TensorFlow}, \texttt{postgres} $\rightarrow$ \texttt{PostgreSQL}. La fonction \texttt{\_normalize\_skill\_name} convertit le nom brut en minuscules, supprime tirets et underscores, puis recherche la correspondance dans \texttt{\_ALIASES}, retournant à défaut la forme \texttt{title()} du nom original. Cette normalisation est critique pour éviter qu'une compétence présente dans plusieurs sources ne soit comptabilisée avec des scores fragmentés.


\subsection{Réalisation}

Cette section présente les interfaces développées lors du Sprint 2, illustrant la concrétisation des pipelines d'évaluation multi-source (\textit{compétences techniques et comportementales}).

\subsubsection{Module d'Évaluation Comportementale (\textit{Compétences comportementales})}

\textbf{1. Extraction de données (GitHub \& CV)}

L'interface permet au employé de soumettre son identifiant GitHub et son CV. Le système traite ces entrées pour extraire instantanément de premières tendances comportementales basées sur l'historique d'engagement et la sémantique du texte.
\begin{figure}[H]
  \centering
  \includegraphics[width=0.6\textwidth]{images/soft1.png}
  \caption{Interface de soumission des profils et extraction CV/GitHub}
  \label{fig:soft-submission-cv-github}
\end{figure}
%\figureplaceholder{Interface de soumission des profils et extraction CV/GitHub}{fig:soft-submission-cv-github}

\textbf{2. Test Process Communication Model (PCM)}

L'employe complète un questionnaire de personnalité de 18 questions. L'interface restitue immédiatement le profil psychologique dominant grâce à un calcul purement algorithmique.
\begin{figure}[H]
  \centering
  \includegraphics[width=0.5\textwidth]{images/softPCM.png}
  \caption{Interface du test PCM et restitution du profil dominant}
  \label{fig:soft-pcm-test-profile}
\end{figure}
%\figureplaceholder{Interface du test PCM et restitution du profil dominant}{fig:soft-pcm-test-profile}

\textbf{3. Paramétrage du Scénario (Rôle et Niveau)}

Pour garantir la pertinence de la mise en situation, le employé définit son rôle métier cible et son niveau de séniorité. Ces paramètres permettent au modèle de calibrer dynamiquement le contexte du test.
\begin{figure}[H]
  \centering
  \includegraphics[width=0.65\textwidth]{images/SoftScenarioGenerate.png}
  \caption{Paramétrage contextuel de l'évaluation IA}
  \label{fig:soft-scenario-parameters}
\end{figure}
%\figureplaceholder{Paramétrage contextuel de l'évaluation IA}{fig:soft-scenario-parameters}

\textbf{4. Évaluation par Scénario IA}

L'agent conversationnel soumet un dilemme professionnel sur mesure. Le système analyse la réponse en texte libre du employé pour évaluer quatre dimensions clés : empathie, assertivité, pragmatisme et clarté.
\begin{figure}[H]
  \centering
  \includegraphics[width=0.7\textwidth]{images/softScenario.png}
  \caption{Interaction avec le scénario comportemental généré par IA}
  \label{fig:soft-ai-scenario-interaction}
\end{figure}
%\figureplaceholder{Interaction avec le scénario comportemental généré par IA}{fig:soft-ai-scenario-interaction}

\textbf{5. Résultats Comportementaux}

Le module consolide l'ensemble des scores \textit{Compétences comportementales} et génère un tableau de bord synthétique, offrant un diagnostic sous forme de graphique radar et des recommandations détaillées.
\begin{figure}[H]
  \centering
  \includegraphics[width=0.4\textwidth]{images/softResult1.png}\hfill
  \includegraphics[width=0.45\textwidth]{images/softResult2.png}
  \caption{Restitution globale des résultats Compétences comportementales}
  \label{fig:soft-skills-results-global}
\end{figure}
%\figureplaceholder{Restitution globale des résultats Compétences comportementales}{fig:soft-skills-results-global}

\subsubsection{Module d'Évaluation Technique (\textit{Compétences techniques})}

\textbf{1. Collecte Technique (GitHub, LinkedIn, CV)}

Le employé renseigne ses liens professionnels. Un pipeline d'extraction robuste croise ces trois sources (profils publics et document parsé) pour identifier exhaustivement les technologies manipulées.
\begin{figure}[H]
  \centering
  \includegraphics[width=0.65\textwidth]{images/hard1.png}
  \caption{Formulaire de collecte des sources techniques}
  \label{fig:hard-technical-sources-collection}
\end{figure}
%\figureplaceholder{Formulaire de collecte des sources techniques}{fig:hard-technical-sources-collection}

\textbf{2. Cartographie et \textit{Job Matching}}

Le système restitue le dictionnaire des compétences extraites, classées par catégorie. Un algorithme calcule instantanément le score d'adéquation (\textit{Job Matching}) du employé face à des profils métiers standards.
\begin{figure}[H]
  \centering
  \includegraphics[width=0.75\textwidth]{images/hard2.png}
  \caption{Cartographie des compétences et taux de Job Matching}
  \label{fig:hard-skills-job-matching}
\end{figure}
%\figureplaceholder{Cartographie des compétences et taux de Job Matching}{fig:hard-skills-job-matching}

\textbf{3. QCM Technique Adaptatif}

Afin de valider le niveau de maîtrise déclaré, la plateforme génère dynamiquement un questionnaire à choix multiples. La difficulté s'adapte aux technologies extraites précédemment.
\begin{figure}[H]
  \centering
  \includegraphics[width=0.75\textwidth]{images/hard3.png}
  \caption{Interface du QCM technique adaptatif}
  \label{fig:hard-adaptive-qcm}
\end{figure}
%\figureplaceholder{Interface du QCM technique adaptatif}{fig:hard-adaptive-qcm}
\textbf{4. Résultats de l'Évaluation Technique}
À l'issue du QCM adaptatif et des épreuves techniques, le système calcule et affiche un score de performance global. L'interface détaille la justesse des réponses par technologie et attribue un niveau de maîtrise consolidé, venant ainsi valider les compétences initialement extraites.

\begin{figure}[H]
  \centering
  \includegraphics[width=0.7\textwidth]{images/hardres1.png}
  \caption{Résultats détaillés de l'évaluation technique et score final}
  \label{fig:hard-skills-results-global}
\end{figure}
\subsubsection{Bilan Consolidé}
\textbf{1. Profil de Compétences et Carrière (Global)}

L'aboutissement de ce sprint est le tableau de bord global. Il fusionne les évaluations techniques et comportementales, fournissant aux admin/RH une vision unifiée du potentiel du employé et des opportunités d'évolution (\textit{Upskilling}).
\begin{figure}[H]
  \centering
  \includegraphics[width=0.9\textwidth]{images/res_tech_soft.png}
  \caption{Tableau de bord global : profil de compétences unifié compétences techniques et comportementales}
  \label{fig:global-hard-soft-dashboard}
\end{figure}
%\figureplaceholder{Tableau de bord global : profil de compétences unifié compétences techniques et comportementales}{fig:global-hard-soft-dashboard}

\subsection{Tests et Déploiement}
La validation du sprint repose sur une stratégie de tests multi-niveaux : tests unitaires des workflows n8n, tests d'intégration des endpoints FastAPI, tests de robustesse des fallbacks en cas de défaillance LLM, et validation fonctionnelle end-to-end. Cette section présente les résultats de ces campagnes de tests ainsi que le processus de déploiement containerisé
\subsubsection{Tests réalisés}

\textbf{Tests des workflows n8n (module compétences comportementales)}
Chaque workflow a été validé directement depuis l'interface n8n en mode « Test Workflow », permettant de vérifier en temps réel la validité du JSON retourné, le bon enchaînement des nœuds et le déclenchement correct des mécanismes de fallback. Le tableau ci-dessous synthétise les cas testés et les résultats obtenus pour les quatre workflows du module compétences comportementales.
\begin{table}[H]
\centering
\small
\renewcommand{\arraystretch}{1.25}
\begin{tabularx}{\TableWidth}{|P{2.3cm}|P{2.5cm}|X|X|}
\hline
\textbf{Workflow} & \textbf{Webhook} & \textbf{Cas testés} & \textbf{Résultat} \\
\hline
GitHub Analysis & POST /webhook/\newline github-analysis & Nom d'utilisateur valide, vide et inexistant & Succès avec fallback si vide \\
\hline
CV Parser & POST /webhook/\newline behavioral-\newline skills-extract & CV riche et CV vide inférieur à 50 caractères & JSON structuré avec cv\_score = 0 si vide \\
\hline
PCM Test & POST /webhook/\newline pcm-test & 18 réponses numériques & Six scores, top strength et faiblesse principale \\
\hline
Master Agent & POST /webhook/\newline master-agent & Payload complet et GitHub indisponible & Synthèse produite avec pondération dynamique \\
\hline
\end{tabularx}
\caption{Tests des workflows n8n}
\label{tab:tests-n8n}
\end{table}

\textbf{Tests FastAPI   routes compétences comportementales}
Les endpoints comportementaux du module compétences comportementales ont été testés via la documentation interactive Swagger UI . Les deux routes principales ont été soumises à des cas représentatifs couvrant la génération de scénarios et l'évaluation de réponses libres. Le tableau ci-dessous présente les résultats obtenus.
\begin{table}[H]
\centering
\small
\renewcommand{\arraystretch}{1.25}
\begin{tabularx}{\textwidth}{|P{4.8cm}|Y|Y|}
\hline
\textbf{Endpoint} & \textbf{Cas testé} & \textbf{Résultat} \\
\hline
POST \texttt{/api/test/}\newline\texttt{scenario/generate} & Role = SWE, Level = Junior & JSON avec \texttt{scenario\_title} et \texttt{skills\_tested} \\
\hline
POST \texttt{/api/test/}\newline\texttt{scenario/evaluate} & Réponse employée libre & Quatre scores psychologiques sur 100 \\
\hline
\end{tabularx}
\caption{Tests FastAPI pour les routes compétences comportementales}
\label{tab:tests-fastapi-soft}
\end{table}

\textbf{Tests FastAPI   routes compétences techniques}
L'ensemble des routes du module compétences techniques a été validé via Swagger UI en soumettant des payloads réels incluant un profil GitHub actif, un CV PDF vectoriel et une URL LinkedIn. Les tests couvrent le pipeline complet d'analyse employé ainsi que les deux types d'évaluation interactive : QCM adaptatif et challenge de code. Le tableau ci-dessous détaille les cas testés et les résultats observés.
\begin{table}[H]
\centering
\small
\renewcommand{\arraystretch}{1.25}
\begin{tabularx}{\TableWidth}{|P{3.2cm}|P{1.7cm}|X|X|}
\hline
\textbf{Endpoint} & \textbf{Méthode} & \textbf{Cas testé} & \textbf{Résultat} \\
\hline
/analyze-employée & POST (multipart) & GitHub + CV PDF + LinkedIn & JSON complet avec \texttt{skills[]} et \texttt{job\_match} \\
\hline
/api/test/generate & POST & Skill = Python, Level = Intermediate, n = 8 & Huit questions MCQ valides \\
\hline
/api/test/evaluate & POST & Huit réponses avec confiance déclarée & score réel, scores par compétence et précision de confiance \\
\hline
/api/test/code-challenge & POST & Skill = Python, Level = Junior & Type, starter\_code, hints[] et language \\
\hline
/api/test/code-submit & POST & Code soumis avec un indice utilisé & Score sur 100 avec détail des quatre critères \\
\hline
\end{tabularx}
\caption{Tests FastAPI pour les routes compétences techniques}
\label{tab:tests-fastapi-hard}
\end{table}

\textbf{Tests Spring Boot  API Skills}
La couche de persistance exposée par Spring Boot a été testée via la documentation Swagger UI du backend Java accessible à http://localhost:8080/swagger-ui.html, en s'authentifiant avec un token JWT valide. Les quatre opérations CRUD de l'entité Skill ont été vérifiées pour les deux rôles du système : USER pour la création et la consultation, et ADMIN pour la validation et la suppression. Le tableau ci-dessous résume les résultats obtenus.
\begin{table}[H]
\centering
\small
\renewcommand{\arraystretch}{1.25}
\begin{tabularx}{\textwidth}{|P{4.8cm}|Y|P{2.8cm}|}
\hline
\textbf{Endpoint} & \textbf{Test} & \textbf{Résultat} \\
\hline
POST \texttt{/api/skills/}\newline\texttt{accounts/\{id\}} & Création d'une skill TECH Python niveau 3 & UUID retourné \\
\hline
GET /api/skills/\newline accounts/\{id\}/\newline type/TECH & Liste des skills techniques & Liste filtrée par type \\
\hline
PUT \texttt{/api/skills/}\newline\texttt{\{id\}/valider} & Validation d'une skill par ADMIN & validee = true \\
\hline
DELETE /api/skills/\{id\} & Suppression d'une skill & Code HTTP 204 \\
\hline
\end{tabularx}
\caption{Tests Spring Boot pour l'API Skills}
\label{tab:tests-spring-skills}
\end{table}

\textbf{Tests de robustesse   fallbacks}
La robustesse du pipeline face aux défaillances des services externes constitue une exigence critique du module. Des tests de panne simulée ont été conduits en arrêtant successivement Ollama, CodeLlama et en soumettant des PDF scannés sans couche texte extractible. Le tableau ci-dessous confirme que chaque scénario de défaillance déclenche bien le mécanisme de fallback prévu sans interruption du service.
\begin{table}[H]
\centering
\small
\renewcommand{\arraystretch}{1.25}
\begin{tabularx}{\TableWidth}{|X|P{2.3cm}|Y|P{1.2cm}|}
\hline
\textbf{Cas} & \textbf{Module} & \textbf{Comportement} & \textbf{Résultat} \\
\hline
Ollama hors ligne & GitHub Analysis n8n & Le nœud JavaScript retourne des scores nuls avec parse\_error = true & OK \\
\hline
Ollama hors ligne & Scénario FastAPI & Retour d'un scénario de secours prédéfini & OK \\
\hline
Ollama hors ligne & QCM FastAPI & Questions de fallback prédéfinies par catégorie & OK \\
\hline
CodeLlama hors ligne & Code Challenge & \texttt{starter\_code} générique de secours & OK \\

\hline
PDF scanné & CV Parser hard & Warning frontend et analyse partielle sans CV & OK \\
\hline
JSON malformé Ollama & Tous modules & Extraction du premier objet JSON valide et fallback sécurisé & OK \\
\hline
\end{tabularx}
\caption{Tests de robustesse et mécanismes de fallback}
\label{tab:tests-fallbacks}
\end{table}

\textbf{Tests de scoring multi-sources}
La cohérence de l'algorithme de nivellement a été vérifiée en construisant des jeux de données synthétiques couvrant les quatre niveaux de maîtrise possibles, de Beginner à Expert, avec des combinaisons variables de sources. Ces tests permettent de s'assurer que la logique de triangulation multi-sources produit des scores stables et conformes aux seuils définis dans \texttt{scoring\_engine.py}. Le tableau ci-dessous présente les cas de validation retenus et les résultats obtenus.
\begin{table}[H]
\centering
\small
\renewcommand{\arraystretch}{1.25}
\begin{tabularx}{\TableWidth}{|Y|Y|P{1.8cm}|}
\hline
\textbf{Cas} & \textbf{Résultat attendu} & \textbf{Résultat obtenu} \\
\hline
Skill présente dans GitHub, CV et portfolio & Expert avec score supérieur ou égal à 90 & OK \\
\hline
Skill présente dans GitHub et CV uniquement & Advanced avec score entre 70 et 89 & OK \\
\hline
Skill présente uniquement dans le CV avec une mention & Beginner avec score proche de 25 & OK \\
\hline
Job match Python Developer avec 8 compétences sur 14 & Score proche de 68/100 & OK \\
\hline
\end{tabularx}
\caption{Validation du scoring multi-sources}
\label{tab:tests-scoring-multisources}
\end{table}


\subsubsection{Déploiement}

La plateforme est déployée via \textit{Docker Compose} en six services conteneurisés et interconnectés. Le tableau ci-dessous détaille la configuration de chaque service.

\begin{table}[H]
\centering
\scriptsize
\renewcommand{\arraystretch}{1.2}
\begin{tabularx}{\TableWidth}{|P{2.4cm}|P{2.5cm}|P{1.2cm}|Y|}
\hline
\textbf{Service} & \textbf{Image / Build} & \textbf{Port exposé} & \textbf{Variables d'environnement clés} \\
\hline
n8n & \texttt{n8nio/n8n} & 5678 & Volume persistant \texttt{n8n\_data} \\
\hline
Ollama & \texttt{ollama/ollama} & 11434 & Volume persistant \texttt{ollama\_data} \\
\hline
\texttt{talentpredict-}\newline\texttt{ai} (FastAPI) & \texttt{./talentpredict-}\newline\texttt{ai} & 8000 & \texttt{ollama}, \texttt{GITHUB\_TOKEN}, \texttt{llama3.2}, \texttt{CODE\_CHALLENGE\_MODEL=codellama} \\
\hline
\texttt{backend} (Spring Boot) & \texttt{./BackEnd} & 8080 & \texttt{SPRING\_DATASOURCE\_URL}, \texttt{TALENTPREDICT\_AI\_BASE\_URL}, timeout 120 s \\
\hline
\texttt{postgres} & \texttt{postgres:15} & 5432 & \texttt{POSTGRES\_DB=talentpredict}, volume persistant \\
\hline
\texttt{frontend}\newline (Angular) & \texttt{./FrontEnd} & 4200 & --- \\
\hline
\end{tabularx}
\caption{Configuration des services Docker Compose}
\label{tab:docker-compose-services}
\end{table}

Les services communiquent via le réseau Docker interne, ce qui évite toute exposition directe des composants sensibles (PostgreSQL, Ollama) sur l'interface publique. Le timeout configurable de 120 secondes entre Spring Boot et FastAPI absorbe la latence inhérente aux appels CodeLlama les plus lourds.


%\begin{table}[H]
%\centering
%%\label{tab:urls-acces}
%\small
%\renewcommand{\arraystretch}{1.25}
%\begin{tabularx}{\TableWidth}{|P{3.4cm}|Y|}
%\hline
%\textbf{Service} & \textbf{URL} \\
%\hline
%n8n (workflows visuels) & \url{http://localhost:5678} %\\
%\hline
%FastAPI (Swagger UI) & \url{http://localhost:8000/docs} \\
%\hline
%Spring Boot (Swagger UI) & \url{http://localhost:8080/swagger-ui.html} \\
%\hline
%Frontend Angular & \url{http://localhost:4200} \\
%\hline
%Ollama & \url{http://localhost:11434} \\
%\hline
%PostgreSQL & \texttt{localhost:5432} \\
%\hline
%\end{tabularx}
%
%\caption{URLs d'accès après déploiement}
%\end{table}


\section{Suivi et Pilotage du Sprint 2}

\subsection{Scrum Board}

Pour garantir une gestion transparente, ce sprint de trois semaines, estimé à 37 \textit{Story Points}, a été piloté via un Scrum Board. Ce management visuel a permis à l'équipe d'identifier immédiatement les goulots d'étranglement et de fluidifier le passage des tâches liées à l'IA vers la phase d'intégration finale.
\begin{figure}[H]
  \centering
  \includegraphics[width=0.9\textwidth]{images/Scrum Board.png}
  \caption{État d'avancement du Scrum Board  Sprint 2}
  \label{fig:scrum-board-sprint2}
\end{figure}
%\figureplaceholder{État d'avancement du Scrum Board --- Sprint 2}{fig:scrum-board-sprint2}

\subsection{Le Burndown Chart}

Afin d'évaluer notre vélocité, le Burndown Chart a été mis à jour quotidiennement. Le graphique révèle un léger ralentissement entre le Jour 7 et le Jour 10, dû aux défis d'intégration du microservice IA, notamment l'instabilité des réponses JSON. La résolution de ces obstacles a provoqué une franche accélération, permettant d'achever le sprint dans les délais impartis.
\begin{figure}[H]
  \centering
  \includegraphics[width=0.9\textwidth]{images/Burndown Chart.png}
  \caption{Burndown Chart : vélocité de l'équipe face aux défis techniques  Sprint 2}
  \label{fig:burndown-chart-sprint2}
\end{figure}
%\figureplaceholder{Burndown Chart : vélocité de l'équipe face aux défis techniques --- Sprint 2}{fig:burndown-chart-sprint2}
\subsection{Indicateurs de Performance (KPIs) du Sprint 2}

Afin de valider cette itération dédiée à l'évaluation multi-source (\textit{compétences techniques et comportementales}), nous avons mesuré l'efficacité de nos pipelines de collecte et de nos modèles d'IA.
\paragraph*{1. Pilotage et Agilité}
Malgré la complexité de l'intégration entre l'orchestrateur \texttt{n8n} et le microservice \texttt{FastAPI}, l'équipe a maintenu un rythme constant. La totalité des 37 \textit{Story Points} planifiés a été livrée, sans aucun report (\textit{Carry-over}), validant ainsi la précision de nos estimations initiales.
\paragraph*{2. Métriques Techniques}
Le tableau ci-dessous résume les performances clés atteintes à la fin de ce sprint, confirmant la viabilité et la rapidité de l'architecture d'évaluation.

\begin{table}[H]
    \centering
    \renewcommand{\arraystretch}{1.5}
    \small
    \begin{tabularx}{0.95\textwidth}{|X|c|c|c|}
        \hline
        \textbf{Indicateur} & \textbf{Cible} & \textbf{Résultat Obtenu} & \textbf{État} \\
        \hline
        \textit{Story Points} complétés & 37 & 37 & Validé \\
        \hline
        Temps d'exécution workflow \texttt{n8n} & $< 10$ s & $\sim 6.2$ s & Validé \\
        \hline
        Précision de l'extraction CV (150+ technos) & $> 85\%$ & $91\%$ & Validé \\
        \hline
        Temps de réponse IA (Évaluation QCM/Code) & $< 8$ s & $\sim 5.5$ s & Validé \\
        \hline
        Taux de succès des tests d'intégration & $> 90\%$ & $100\%$ & Validé \\
        \hline
    \end{tabularx}
    \caption{Indicateurs de Performance du Sprint 2}
    \label{tab:kpi-tech-s2}
\end{table}

\paragraph*{3. Bilan de l'itération}
Le bilan du Sprint 2 est hautement satisfaisant. La triangulation des données (CV, GitHub, LinkedIn) via les workflows \texttt{n8n} couplée à l'analyse sémantique du modèle LLM s'est avérée robuste. Les temps de réponse obtenus respectent largement les contraintes métiers, garantissant une expérience employé fluide lors des évaluations.

\section{ Conclusion}

Dans ce chapitre, nous avons présenté la conception et la réalisation du moteur d'évaluation multi-sources, pilier diagnostique de TalentPredict. Grâce à une modélisation UML rigoureuse, nous avons structuré l'orchestration complexe des données (CV, GitHub, PCM) vers l'IA locale (Ollama). Le développement d'interfaces Angular interactives, couplées à une architecture \textit{backend} hybride (Spring Boot, FastAPI, n8n), a permis de concrétiser l'enjeu du sprint : objectiver avec précision les compétences techniques et comportementales des collaborateurs.

Par ailleurs, l'intégration du Communication Hub (SSE) a garanti une expérience utilisateur fluide malgré les délais d'inférence du modèle LLM. L'application stricte de la méthode Agile Scrum a prouvé la fiabilité de ce pipeline d'évaluation, aboutissant à un calcul pertinent du score de \textit{Job Matching}.

Avec la validation de ce deuxième sprint, la plateforme dispose désormais d'une cartographie exhaustive et certifiée des talents. Le chapitre suivant s’attachera à présenter le troisième sprint, dédié à l'exploitation de ces résultats à travers le module d'Upskilling et de gouvernance RH.

\chapter{Sprint 3 : Moteur d'IA et Gestion de l'Upskilling}

\section*{Introduction}
\phantomsection
\addcontentsline{toc}{section}{Introduction}
\phantomsection 
Suite à la mise en place des mécanismes d'évaluation technique et comportementale lors du sprint précédent, le défi principal de cette troisième itération est de transformer ces analyses en un véritable outil de développement des compétences. En effet, identifier une lacune de compétences ne sert à rien si le système ne propose pas de solutions de remédiation ciblées et automatisées.

Ce chapitre est consacré à la conception et à la mise en œuvre du Moteur d'Upskilling Intelligent lié au Module de Gouvernance RH. Ce module est le « cœur apprenant » de la plateforme \textbf{TalentPredict} et a pour but d'organiser un cycle d'apprentissage complet. Il assure une transition fluide depuis la recommandation prédictive de parcours de formation, basée sur l'Intelligence Artificielle, jusqu'à la certification des acquis, tout en offrant une supervision administrative rigoureuse et centralisée.
\section{Objectifs du Sprint 3}

L'objectif majeur de ce sprint est de livrer un module opérationnel dédié à la gestion de \textit{l'Upskilling} et à \textit{la Gouvernance RH} . L'enjeu est de faire évoluer la plateforme \textbf{TalentPredict} : d'un simple outil d'évaluation, elle devient un véritable moteur de développement professionnel, capable de recommander, de suivre et de certifier la montée en compétences des collaborateurs de manière intelligente et automatisée


\textbf{1. Objectifs fonctionnels et interactifs:}
\begin{itemize}
    \item \textbf{Upskilling Personnalisé :} Déploiement d'un moteur de recommandation de formations ciblé, s'appuyant sur l'analyse précise des écarts de compétences (\textit{Skill Gaps}).
    \item \textbf{Tracking \& Progression :} Mise en œuvre d’un \textit{Dashboard} global incluant une \textit{Roadmap} dynamique et un tableau \textit{Kanban} interactif pour le suivi des apprentissages des utilisateurs.
    \item \textbf{Validation \& Gamification :} Intégration d’un module de mini-tests générés par l’IA et d’un \textit{Leaderboard} pour stimuler l'engagement.
    \clearpage
    \item \textbf{Espace Gouvernance RH :} Conception d'une interface d'administration complète permettant aux responsables de gérer (\texttt{approuver, suggérer ou refuser}) les parcours de formation.
\end{itemize}
 \textbf{2. Objectifs techniques et opérationnels  :}
\begin{itemize}
    \item \textbf{Intégration IA :} Optimisation des flux d'inférence via le microservice \textbf{Python} pour assurer une génération fluide et structurée des contenus pédagogiques.
    \item \textbf{Communication Synchrone :} Implémentation d'un système de notifications en temps réel pour garantir un suivi réactif des flux de validation \textit{workflows} .
    \item \textbf{Reporting \& Certification :} Développement d'un moteur permettant la génération dynamique et automatisée de certificats au format PDF à la volée.
    \item \textbf{Sécurité \& Workflow :} Renforcement de la sécurité via un modèle de contrôle d'accès basé sur les rôles (\textbf{RBAC}) , sécurisant ainsi les processus de décision administrative.
    \end{itemize}

\paragraph{Critères d’acceptation (Definition of Done)}

La validation des modules d'Upskilling et de Gouvernance repose sur les critères suivants :

\begin{itemize}
    \item \textbf{Cohérence IA} : Les recommandations de formation et les mini-tests générés par le microservice Python doivent être cohérents avec le profil de l'utilisateur.
    \item \textbf{Réactivité (SSE)} : Les notifications de validation RH doivent être reçues instantanément par l'utilisateur sans latence perceptible.
    \item \textbf{Conformité PDF} : Les certificats générés doivent être conformes au modèle visuel et inclure les métadonnées de l'utilisateur sans erreur de rendu.
    \item \textbf{Sécurité RBAC} : Les routes d'administration (approbation/refus) doivent être inaccessibles aux utilisateurs ne possédant pas le rôle \textit{ADMIN}.
    \item \textbf{Intégrité Workflow} : Le passage d'une formation dans le tableau Kanban doit mettre à jour l'état de la base de données en temps réel.
\end{itemize}



\section{Séance du Sprint planning }

La séance de \textit{Sprint Planning} a réuni l’équipe de développement, composée de \textbf{Mohamed Ben Abdeladhim} et \textbf{Rahma Barouni}, ainsi que le \textit{Product Owner} \textbf{Mr Med Amine Gharbi}, et notre encadrante pédagogique  \textbf{Mme Ben Aïcha Takwa}. Cette réunion a permis de définir le périmètre fonctionnel et technique de ce troisième sprint.

Les points clés suivants ont été déterminés lors de cette séance :

\begin{itemize}
    \item \textbf{Durée :} La durée de cette itération a été fixée à \textbf{un mois}, un délai nécessaire pour stabiliser les algorithmes de recommandation et assurer l'intégration fluide entre le \textit{backend} \texttt{Spring Boot} et le microservice d'IA.
    
    \item \textbf{Priorités Absolues :}
    \begin{itemize}
        \item Mise en place du \textbf{Tableau de Bord d'Upskilling} pour les utilisateurs et l'espace de \textbf{Gouvernance RH}.
        \item Développement du moteur de \textbf{Roadmap IA} et du catalogue de formations.
        \item Implémentation du module de \textbf{Mini-Tests} et du système de \textbf{Génération Automatisée des Certificats}.
        \item Gestion du cycle de vie des formations (Approbation, Suggestion, Refus par les RH).
    \end{itemize}
    
    \item \textbf{Priorités Secondaires :}
    \begin{itemize}
        \item Intégration des mécaniques de \textbf{Gamification} (\textit{Leaderboard} et Badges).
        \item Optimisation des performances du moteur d'IA pour la génération de questions.
        \item Mise à jour du pipeline \texttt{CI/CD} via \texttt{Docker} et \texttt{GitLab} pour inclure les nouveaux \textit{endpoints}.
    \end{itemize}
    
    \item \textbf{Prérequis :} La stabilité des API livrées lors des sprints précédents est essentielle, notamment le module de diagnostic technique et l'évaluation comportementale , qui servent de données d'entrée au moteur de recommandation.
\end{itemize}

À la fin de cette réunion, les besoins ont été décomposés en \textbf{User Stories (US)} et estimés en \textbf{jours-homme}, en donnant la priorité maximale aux livrables critiques pour assurer la continuité entre l'évaluation des talents et leur montée en compétences.
\section{Backlog du Sprint 3}
Le Sprint Backlog constitue le plan d'action opérationnel de cette troisième itération. Il rassemble l'ensemble des User Stories (US) extraites du Product Backlog afin de concrétiser les objectifs liés à l'Upskilling définis précédemment. Pour garantir une planification rigoureuse, chaque tâche est catégorisée par Epic (macro-fonctionnalités) et estimée en Story Points en s'appuyant sur la suite de Fibonacci. Cette méthode permet d'évaluer l'effort de développement nécessaire tout en assurant une vélocité réaliste et maîtrisée pour l'équipe.
\begin{table}[H]
    \centering
    \renewcommand{\arraystretch}{1.5}
    \small
    \begin{tabularx}{\textwidth}{|c|>{\raggedright\arraybackslash}p{2.5cm}|>{\raggedright\arraybackslash}X|c|>{\centering\arraybackslash}p{2cm}|}
        \hline
        \rowcolor[HTML]{EFEFEF} 
        \textbf{ID} & \textbf{Module / US} & \textbf{Tâches de Développement} & \textbf{JH} & \textbf{Priorité} \\
        \hline
        
        US 3.1 & Interface \textit{Upskilling} & Création du \textit{Dashboard} Global et du tableau \textit{Kanban} interactif. & 3 & Haute \\
        \hline
        
        US 3.2 & Moteur de \textit{Roadmap} & Intégration de l'API IA (\texttt{Python}) pour la génération de parcours basés sur le \textit{Skill Gap}. & 4 & \textbf{Critique} \\
        \hline
        
        US 3.3 & Catalogue Formations & Développement du service de recommandation et de la base de données des cours. & 2 & Haute \\
        \hline
        
        US 3.4 & Mini-Tests IA & Création du moteur de génération de QCM et du service de correction automatique. & 4 & \textbf{Critique} \\
        \hline
        
        US 3.5 & \textit{Gamification} & Implémentation du système de points, badges et du \textit{Leaderboard} temps réel. & 2 & Moyenne \\
        \hline
        
        US 3.6 & Gouvernance RH & Développement du \textit{workflow} d'approbation (Approuver/Suggérer/Refuser) et espace Admin. & 3 & Haute \\
        \hline
        
        US 3.7 & Certification & Mise en place du moteur \texttt{PDFBox} pour la génération dynamique des certificats. & 2 & Haute \\
        \hline
        
        TS 3.8 & Assurance Qualité & Tests unitaires, tests d'intégration (\texttt{Spring}/\texttt{Python}) et tests de montée en charge. & 2 & \textbf{Critique} \\
        \hline
        
    \end{tabularx}
    \vspace{2mm}
    \caption{Backlog du Sprint 3}
    \label{tab:planification_sprint3}
\end{table}

\section{Démarrage du Sprint 3}
Le début du Sprint 3 a été caractérisé par une phase d'initialisation technique approfondie en préparation des développements relatifs à l'Upskilling. L'équipe s'est aussitôt focalisée sur l'établissement des canaux de communication entre le backend Spring Boot et le microservice Python, tout en mettant en place les services de notifications en temps réel. Cette phase essentielle a établi les fondations nécessaires pour l'application des algorithmes de recommandation et du système de gouvernance RH, garantissant ainsi un passage en douceur de la planification à la mise en œuvre effective des modules.
\subsubsection*{1. Activités réalisées}

Durant cette itération, l'équipe a mené de front des développements \textit{backend} complexes et des intégrations d'intelligence artificielle :

\begin{itemize}
    \item \textbf{Ingénierie des Prompts :} Conception de \textit{prompts} système pour forcer l'IA à générer des \textit{roadmaps} et des tests au format JSON strict.
    \item \textbf{Développement Backend (\texttt{Spring Boot}) :} Implémentation des services métier pour la gestion du cycle de vie des formations et le calcul du \textit{Readiness Score}.
    \item \textbf{Développement Frontend (\texttt{Angular}) :} Création de composants dynamiques pour le tableau \textit{Kanban} et le module interactif de mini-tests.
    \item \textbf{Mise en place de la Gouvernance :} Développement de l'espace HR avec des \textit{dashboards} de suivi et des actions de validation (Approbation/Refus).
    \item \textbf{Intégration du moteur PDF :} Configuration du service de génération de certificats à la volée.
\end{itemize}

\subsubsection*{2. Livrables produits}

Les livrables suivants ont été validés et intégrés à la branche principale du projet :

\begin{itemize}
    \item \textbf{Module "\textit{Upskilling Engine}" :} Moteur fonctionnel de recommandation et de génération de \textit{roadmaps}.
    \item \textbf{Module "\textit{Evaluation Validation}" :} Interface de passage de tests et correction automatique.
    \item \textbf{Module "\textit{HR Governance}" :} Console d'administration pour la supervision des talents.
    \item \textbf{Service de Certification :} Générateur de documents PDF officiels.
    \item \textbf{Documentation API :} Mise à jour de \texttt{Swagger} pour inclure les \textit{endpoints} de formation.
\end{itemize}

\subsubsection*{3. Difficultés rencontrées et traitement}

L'exécution de ce sprint a été marquée par trois défis techniques majeurs :

\begin{itemize}
    \item \textbf{Difficulté 1 : Instabilité des formats JSON de l’IA}
    \begin{itemize}
        \item \textbf{Problème :} Les réponses des modèles LLM présentaient parfois des structures incohérentes, compromettant le \textit{parsing} des \textit{roadmaps}.
        \item \textbf{Traitement :} Renforcement du \textit{system prompting} et intégration de validateurs de schémas rigoureux côté \textit{backend} pour garantir l'intégrité des données.
    \end{itemize}
    \vspace{2mm}
    
    \item \textbf{Difficulté 2 : Conflits de dépendances PDF}
    \begin{itemize}
        \item \textbf{Problème :} L'intégration de la génération de certificats a révélé des incompatibilités entre les bibliothèques de rendu.
        \item \textbf{Traitement :} Isolation des dépendances critiques et migration vers \texttt{OpenHTMLtoPDF}, assurant ainsi une stabilité du moteur de génération et un rendu style fidèle.
    \end{itemize}
    \vspace{2mm}
    
    \item \textbf{Difficulté 3 : Latence des appels API externes}
    \begin{itemize}
        \item \textbf{Problème :} Les délais de génération des tests par l’IA impactaient la fluidité de l'interface.
        \item \textbf{Traitement :} Implémentation d'une communication asynchrone via les \textit{Server-Sent Events} (\texttt{SSE}), permettant d'informer l'utilisateur en temps réel de l'état d'avancement des processus en tâche de fond.
    \end{itemize}
\end{itemize}

\subsection{Analyse}

La phase d’analyse vise à définir le périmètre fonctionnel du module de Formation.
et identifier avec précision les interactions entre le système et les utilisateurs. Contrairement au
Le premier sprint du module introduit une grande automatisation grâce à \textbf{Intelligence Artificielle}.

\subsubsection*{1. Identification des Acteurs}

Le système interagit avec trois acteurs principaux, parmi lesquels se trouve un acteur automatisé :

\begin{itemize}
    \item \textbf{Utilisateur (Employé) :} Constitue le cœur du module. Il réalise les évaluations, examine ses compétences et suit sa progression d'apprentissage via la roadmap interactive. Il soumet ses preuves de réussite (mini-tests ou certificats) pour valider sa montée en compétences.
    
    \item \textbf{Administrateur (RH) :} garantit la politique de formation de l’entreprise.
Le superviseur de l’équipe dirige les membres. Le superviseur a le droit d’approuver ou de rejeter chaque demande. L’inscription aux cours est obligatoire.
    
    \item \textbf{Moteur IA (\enquote{System}) :} Agit de manière autonome en arrière-plan. Il est sollicité pour générer des recommandations ciblées suite à un \textit{Assessment} et pour créer dynamiquement les questions des mini-tests de validation.
\end{itemize}


\subsubsection*{2. Clarification du besoin}

La finalité de ce sprint est de répondre à la problématique de la \textbf{Remédiation Intelligente}. Le système doit automatiser le cycle de vie de l'\textit{Upskilling} en comblant le fossé entre les compétences détectées et les exigences du marché. L'enjeu est de fournir un environnement de \textbf{GPEC} (Gestion Prévisionnelle des Emplois et des Compétences) dynamique, où chaque recommandation est justifiée par une analyse de données et chaque réussite est sanctionnée par une preuve tangible (Certificat).

\subsubsection*{3. Scénarios de validation}

Pour garantir la robustesse fonctionnelle, trois parcours types ont été modélisés :

\begin{itemize}
    \item \textbf{Parcours Nominal (Succès) :} \\
    Identification du besoin $\rightarrow$ Requête de formation $\rightarrow$ Approbation RH $\rightarrow$ Suivi du module $\rightarrow$ Réussite au test $\rightarrow$ Certification.
    
    \item \textbf{Parcours Alternatif (Ajustement RH) :} \\
    Requête utilisateur $\rightarrow$ Évaluation de la pertinence par le RH $\rightarrow$ Recommandation d'une ressource alternative $\rightarrow$ Mise à jour de la \textit{Roadmap}.
    
    \item \textbf{Parcours d'Échec (Réévaluation) :} \\
    Tentative de validation infructueuse $\rightarrow$ Analyse des erreurs par le système $\rightarrow$ Prescription de révision $\rightarrow$ Réactivation du test après un délai de carence.
\end{itemize}

\subsubsection*{4. Contraintes techniques}

L'implémentation de ces processus repose sur le respect de contraintes strictes :

\begin{itemize}
    \item \textbf{Persistance et Intégrité :} Garantie de la cohérence des états de formation dans la base de données \texttt{PostgreSQL} lors des transitions du \textit{Kanban}.
    \item \textbf{Interfaçage Asynchrone :} Gestion des temps de réponse des modèles LLM via une architecture événementielle pour maintenir une interface fluide.
    \item \textbf{Sécurité et RBAC :} Implémentation de contrôles d'accès granulaires pour les \textit{endpoints} de gouvernance, limitant les actions sensibles aux seuls profils RH.
    \item \textbf{Standardisation de la Certification :} Respect des normes de rendu documentaire pour la génération de certificats PDF à haute fidélité visuelle.
\end{itemize}


\subsubsection*{5. Diagramme de Cas d’Utilisation du Sprint 3}
Afin d'illustrer ces interactions, le diagramme de cas d'utilisation ci-dessous met en évidence la séparation des responsabilités entre l'autonomie du Candidat, le contrôle du RH, et le travail de fond réalisé par l'IA.
\begin{figure}[H]
    \centering
    \makebox[\textwidth][c]{%%
        \includegraphics[width=1\textwidth]{images/usecaseSprint3.png}%%
    }
    \caption{Diagramme de Cas d'Utilisation globale du Sprint 3}
    \label{fig:uc_formation}
\end{figure}

\color{red}
\begin{itemize}
    \item \textit{Diagramme 12 :} Renommer le titre en \textit{Diagramme de cas d’utilisation global du Sprint 3} afin de préciser clairement qu’il s’agit du périmètre fonctionnel du Sprint 3.
    \item \textit{Diagramme 12 :} Remplacer le nom de la frontière \textit{TalentPredict - Vue Globale} par \textit{TalentPredict AI -- Sprint 3}, car le diagramme ne représente pas toute la plateforme mais seulement les fonctionnalités du sprint.
    \item \textit{Diagramme 12 :} Supprimer les flèches sur les associations entre les acteurs et les cas d’utilisation ; en UML, une association acteur--cas d’utilisation est une ligne simple sans flèche.
    \item \textit{Diagramme 12 :} Renommer l’acteur \textit{RH / Admin} en \textit{Administrateur RH} pour garder une terminologie homogène dans le rapport.
    \item \textit{Diagramme 12 :} Remplacer l’acteur \textit{GitHub API} par \textit{Service GitHub}, car il représente un service externe utilisé par la plateforme.
    \item \textit{Diagramme 12 :} Remplacer l’acteur \textit{Ollama} par \textit{Service IA}, puis préciser dans le texte que le service IA s’appuie sur Ollama.
    \item \textit{Diagramme 12 :} Remplacer l’acteur \textit{Communication Hub} par \textit{Service de communication temps réel} pour garder le même vocabulaire que dans les autres figures.
    \item \textit{Diagramme 12 :} Remplacer \textit{Évaluer Compétences Techniques} par \textit{Évaluer les compétences techniques}.
    \item \textit{Diagramme 12 :} Remplacer \textit{Évaluer Compétences Comportementales} par \textit{Évaluer les compétences comportementales}.
    \item \textit{Diagramme 12 :} Remplacer \textit{Analyser historique de code GitHub} par \textit{Analyser le profil GitHub}, formulation plus simple et plus cohérente avec le rapport.
    \item \textit{Diagramme 12 :} Remplacer \textit{Extraire stack technique (CV)} par \textit{Analyser le CV} ou \textit{Extraire les compétences techniques du CV}.
    \item \textit{Diagramme 12 :} Remplacer \textit{Passer QCM adaptatif} par \textit{Passer un QCM adaptatif}.
    \item \textit{Diagramme 12 :} Remplacer \textit{Résoudre challenge de code} par \textit{Réaliser un challenge de code}, car le cas d’utilisation doit décrire l’action de l’employé.
    \item \textit{Diagramme 12 :} Remplacer \textit{Analyser collaboration GitHub} par \textit{Analyser la collaboration GitHub} ou \textit{Analyser l’activité collaborative GitHub}.
    \item \textit{Diagramme 12 :} Remplacer \textit{Extraire sémantique CV} par \textit{Analyser la sémantique du CV}.
    \item \textit{Diagramme 12 :} Remplacer \textit{Résoudre scénario immersif} par \textit{Répondre à un scénario immersif}.
    \item \textit{Diagramme 12 :} Remplacer \textit{Passer test heuristique PCM} par \textit{Passer le test PCM}, car le détail heuristique doit être expliqué dans le texte et non dans le libellé du cas.
    \item \textit{Diagramme 12 :} Remplacer \textit{Recevoir notifications SSE et suivi} par \textit{Consulter les notifications et le suivi}, car le cas d’utilisation doit représenter une action observable de l’utilisateur.
    \item \textit{Diagramme 12 :} Remplacer \textit{S’authentifier (JWT)} par \textit{S’authentifier}, puis expliquer l’usage du JWT dans la partie technique.
    \item \textit{Diagramme 12 :} Séparer le cas \textit{Superviser, valider et générer rapports} en trois cas distincts : \textit{Superviser les évaluations}, \textit{Valider les résultats} et \textit{Générer des rapports}.
    \item \textit{Diagramme 12 :} Ajouter un cas central \textit{Produire les résultats d’évaluation} afin de relier les évaluations techniques et comportementales aux résultats consultables.
    \item \textit{Diagramme 12 :} Ajouter un cas \textit{Générer des recommandations de formation}, car le Sprint 3 semble concerner l’évaluation IA, le suivi et l’upskilling.
    \item \textit{Diagramme 12 :} Vérifier les relations \texttt{<<include>>} depuis \textit{Évaluer les compétences techniques} vers les sous-cas techniques ; elles sont correctes seulement si ces étapes sont obligatoires.
    \item \textit{Diagramme 12 :} Si le QCM adaptatif ou le challenge de code sont optionnels, remplacer les relations \texttt{<<include>>} correspondantes par des relations \texttt{<<extend>>}.
    \item \textit{Diagramme 12 :} Vérifier la relation \texttt{<<include>>} entre \textit{Évaluer les compétences comportementales} et \textit{Analyser la collaboration GitHub}, car ce lien doit être justifié dans le texte.
    \item \textit{Diagramme 12 :} Relier \textit{Service GitHub} uniquement aux cas qui exploitent GitHub, comme \textit{Analyser le profil GitHub} et éventuellement \textit{Analyser la collaboration GitHub}.
    \item \textit{Diagramme 12 :} Relier \textit{Service IA} uniquement aux cas d’analyse intelligente et de génération, et non à toutes les fonctionnalités du diagramme.
    \item \textit{Diagramme 12 :} Relier \textit{Service de communication temps réel} uniquement au cas lié aux notifications ou au suivi en temps réel.
    \item \textit{Diagramme 12 :} Ajouter une légende expliquant les associations simples et les relations \texttt{<<include>>}.
    \item \textit{Diagramme 12 :} Réorganiser les éléments pour réduire les croisements de lignes, en plaçant les actions de l’employé à gauche, les traitements IA au centre, les services externes à droite et les actions RH à droite ou en bas.
    \item \textit{Diagramme 12 :} Ajouter une phrase avant la figure pour préciser qu’elle présente les fonctionnalités principales du Sprint 3 liées à l’évaluation IA et au suivi.
    \item \textit{Diagramme 12 :} Ajouter une phrase après la figure pour expliquer que les détails techniques liés à Ollama, GitHub API, JWT et SSE seront détaillés dans les diagrammes d’architecture ou de séquence.
\end{itemize}
\color{black}


\subsubsection* {Cas 1 : Processus de Planification de l'Apprentissage :}
Ce cas décrit comment un employé prend en main sa montée en compétences en interagissant avec l'IA pour combler ses lacunes.

\begin{figure}[H]
    \centering
    \includegraphics[width=\textwidth]{images/case1Sprint3.png}
    \caption{Diagramme de cas d’utilisation : Processus de planification de l'apprentissage}
    \label{fig:case_cas1_sprint3}
\end{figure}

\color{red}
\begin{itemize}
    \item \textit{Diagramme 13 :} Renommer le titre en \textit{Diagramme de cas d’utilisation du processus de planification de l’apprentissage} pour alléger l’intitulé.
    \item \textit{Diagramme 13 :} Remplacer le nom de la frontière par \textit{Module de planification et de suivi de l’apprentissage} pour garder une formulation simple et académique.
    \item \textit{Diagramme 13 :} Supprimer les flèches sur les associations entre les acteurs et les cas d’utilisation ; en UML, une association acteur--cas d’utilisation est une ligne simple sans flèche.
    \item \textit{Diagramme 13 :} Renommer l’acteur \textit{Utilisateur (Employé)} en \textit{Employé} pour harmoniser le vocabulaire avec le reste du rapport.
    \item \textit{Diagramme 13 :} Remplacer l’acteur \textit{Moteur d’IA} par \textit{Service IA} pour garder une terminologie homogène avec les autres diagrammes.
    \item \textit{Diagramme 13 :} Supprimer les libellés sur les associations avec le service IA, comme \textit{Construit les étapes JSON}, \textit{Filtre les formations} et \textit{Fournit l’analyse initiale}, car ces détails relèvent du texte explicatif et non du diagramme de cas d’utilisation.
    \item \textit{Diagramme 13 :} Vérifier la relation \texttt{<<include>>} entre \textit{Gérer son parcours d’Upskilling} et \textit{S’authentifier} ; elle est correcte seulement si l’authentification est obligatoire pour ce cas.
    \item \textit{Diagramme 13 :} Remplacer \textit{Gérer son parcours d’Upskilling} par \textit{Gérer son parcours d’apprentissage} ou garder un seul terme de manière uniforme dans tout le rapport.
    \item \textit{Diagramme 13 :} Remplacer \textit{Consulter le Dashboard d’Upskilling} par \textit{Consulter le tableau de bord d’apprentissage} pour éviter le mélange français--anglais.
    \item \textit{Diagramme 13 :} Remplacer \textit{Afficher le Skill Gap actuel} par \textit{Afficher l’écart de compétences actuel}.
    \item \textit{Diagramme 13 :} Remplacer \textit{Consulter le Leaderboard} par \textit{Consulter le classement} si vous souhaitez uniformiser les libellés en français.
    \item \textit{Diagramme 13 :} Vérifier que \textit{Consulter le tableau de bord d’apprentissage} inclut bien les cas \textit{Consulter les cours recommandés}, \textit{Afficher l’écart de compétences actuel} et, si nécessaire, \textit{Consulter la roadmap personnalisée}.
    \item \textit{Diagramme 13 :} Ajouter un cas \textit{Consulter la roadmap personnalisée}, car la roadmap est générée dans le diagramme mais sa consultation par l’employé n’est pas représentée clairement.
    \item \textit{Diagramme 13 :} Vérifier la relation \texttt{<<include>>} entre \textit{Générer une roadmap personnalisée} et \textit{Spécifier le profil métier cible} ; elle est correcte si le profil cible est toujours nécessaire.
    \item \textit{Diagramme 13 :} Corriger la relation entre \textit{Piloter l’apprentissage via Kanban} et \textit{Ajouter un cours} ; si l’ajout d’un cours est une action optionnelle du pilotage, utiliser \texttt{<<extend>>} au lieu de \texttt{<<include>>}.
    \item \textit{Diagramme 13 :} Vérifier la relation \texttt{<<extend>>} entre \textit{Rechercher manuellement dans le catalogue} et \textit{Ajouter un cours} ; elle est cohérente seulement si la recherche manuelle constitue un scénario optionnel d’ajout.
    \item \textit{Diagramme 13 :} Corriger la relation \texttt{<<extend>>} entre \textit{Débloquer un badge de progression} et \textit{Consulter le classement}, car le déblocage du badge n’est pas déclenché par la consultation du classement.
    \item \textit{Diagramme 13 :} Relier plutôt \textit{Débloquer un badge de progression} à \textit{Ajouter un cours} ou à un cas du type \textit{Mettre à jour la progression}, puisque la note précise qu’il est déclenché lors de l’ajout du premier cours au Kanban.
    \item \textit{Diagramme 13 :} Vérifier si \textit{Débloquer un badge de progression} est réellement un cas d’utilisation visible pour l’employé ; sinon, remplacer ce cas par \textit{Consulter ses badges} et expliquer l’automatisme dans le texte.
    \item \textit{Diagramme 13 :} Ajouter un cas \textit{Recevoir des recommandations de formation} si le module met en avant la recommandation comme résultat du service IA.
    \item \textit{Diagramme 13 :} Vérifier si le service IA doit être relié uniquement à \textit{Générer une roadmap personnalisée}, \textit{Consulter les cours recommandés} et \textit{Afficher l’écart de compétences actuel}, afin d’éviter des liaisons trop générales.
    \item \textit{Diagramme 13 :} Réorganiser les cas d’utilisation pour réduire les croisements de lignes, en plaçant les cas du tableau de bord à gauche, les cas de génération au centre et les cas Kanban/classement à droite.
    \item \textit{Diagramme 13 :} Ajouter une légende expliquant les associations simples, les relations \texttt{<<include>>} et les relations \texttt{<<extend>>}.
\end{itemize}
\color{black}

\begin{table}[H]
    \centering
    \renewcommand{\arraystretch}{1.3}
    \small
    \begin{tabularx}{\textwidth}{|>{\raggedright\arraybackslash\bfseries}p{3.5cm}|X|}
        \hline
        \rowcolor[HTML]{F2F2F2}
        Paramètre & \textbf{Description} \\
        \hline
        Titre du Cas & Planifier l'apprentissage (\textit{Roadmap}, \textit{Kanban}, Gamification) \\
        \hline
        Acteur Principal & Utilisateur (Employé) \\
        \hline
        Acteur Secondaire & Moteur d'IA  \\
        \hline
        Préconditions & L'utilisateur est authentifié. Le \textit{Skill Gap} (calculé au Sprint 2) est disponible en base de données. \\
        \hline
        Scénario Nominal & 
        \vspace{-2mm}
        \begin{enumerate}[leftmargin=*, nosep]
            \item L'utilisateur accède à son \textit{Dashboard} d'\textit{Upskilling}.
            \item Le système affiche les recommandations de cours ciblant ses points faibles.
            \item L'utilisateur spécifie son poste cible et clique sur "Générer ma \textit{Roadmap}".
            \item Le Moteur d'IA analyse le profil et génère les étapes structurées.
            \item L'utilisateur valide les étapes et les ajoute à son tableau \textit{Kanban} ("À faire").
            \item Le système attribue des points d'initiative et met à jour le \textit{Leaderboard}.
        \end{enumerate} \\
        \hline
        Scénarios Alternatifs & \textbf{Service IA indisponible :} Si le microservice \texttt{Python} ne répond pas, le système désactive le bouton IA et redirige l'utilisateur vers le catalogue de cours pour une sélection manuelle. \\
        \hline
        Post-conditions & La \textit{roadmap} est sauvegardée en base, le \textit{Kanban} est initialisé, et le profil est mis à jour avec les points de gamification. \\
        \hline
    \end{tabularx}
    \caption{Fiche descriptive du Cas d'Utilisation : Planification de l'apprentissage}
    \label{tab:uc_cas1}
\end{table}



\vspace{5mm}

\subsubsection*{Cas 2 : Processus de Gouvernance et Supervision RH :}
Ce cas décrit comment le département RH contrôle le budget et la pertinence des formations demandées par les collaborateurs.

\begin{figure}[H]
    \centering
    \includegraphics[width=\textwidth]{images/case2Sprint3.png}
    \caption{Diagramme  cas d’utilisation :Gouvernance et Supervision RH}
    \label{fig:case_cas2_sprint3}
\end{figure}

\color{red}
\begin{itemize}
    \item \textit{Diagramme 14 :} Renommer le titre en \textit{Diagramme de cas d’utilisation du module de gouvernance et de supervision RH} pour alléger l’intitulé.
    \item \textit{Diagramme 14 :} Remplacer le nom de la frontière par \textit{Module de gouvernance et de supervision RH} avec une formulation homogène et sans majuscules inutiles.
    \item \textit{Diagramme 14 :} Renommer l’acteur \textit{Administrateur (RH)} en \textit{Administrateur RH} afin de garder la même terminologie dans tout le rapport.
    \item \textit{Diagramme 14 :} Remplacer l’acteur \textit{Service de Notifications} par \textit{Service de notification} ou \textit{Service de communication temps réel}, selon le vocabulaire adopté dans les autres diagrammes.
    \item \textit{Diagramme 14 :} Supprimer les flèches sur les associations entre les acteurs et les cas d’utilisation ; en UML, une association acteur--cas d’utilisation est une ligne simple sans flèche.
    \item \textit{Diagramme 14 :} Remplacer \textit{S’authentifier (Rôle Admin)} par \textit{S’authentifier}, puis expliquer la vérification du rôle administrateur dans le texte ou dans un diagramme de séquence.
    \item \textit{Diagramme 14 :} Vérifier la relation \texttt{<<include>>} entre \textit{Superviser la montée en compétences} et \textit{S’authentifier} ; elle est acceptable si l’authentification est obligatoire, mais elle peut être omise pour ne pas surcharger le diagramme.
    \item \textit{Diagramme 14 :} Remplacer \textit{Consulter le Dashboard RH} par \textit{Consulter le tableau de bord RH} afin d’éviter le mélange français--anglais.
    \item \textit{Diagramme 14 :} Remplacer \textit{Suivre les KPIs d’apprentissage de l’équipe} par \textit{Suivre les indicateurs d’apprentissage de l’équipe}.
    \item \textit{Diagramme 14 :} Remplacer \textit{Traiter une demande de formation} par \textit{Traiter les demandes de formation} si ce cas représente une fonctionnalité générale du module.
    \item \textit{Diagramme 14 :} Corriger les relations \texttt{<<extend>>} entre \textit{Approuver la demande}, \textit{Refuser la demande} et \textit{Traiter une demande de formation} ; ces deux actions sont bien des scénarios alternatifs ou conditionnels du traitement de la demande.
    \item \textit{Diagramme 14 :} Vérifier le sens des relations \texttt{<<extend>>} : la flèche doit partir de \textit{Approuver la demande} ou \textit{Refuser la demande} vers \textit{Traiter une demande de formation}.
    \item \textit{Diagramme 14 :} La relation \texttt{<<include>>} entre \textit{Refuser la demande} et \textit{Saisir un motif de refus} est correcte si le motif est obligatoire lors du refus.
    \item \textit{Diagramme 14 :} La relation \texttt{<<include>>} entre \textit{Refuser la demande} et \textit{Informer l’employé de la décision} est correcte si l’information est envoyée systématiquement après un refus.
    \item \textit{Diagramme 14 :} La relation \texttt{<<include>>} entre \textit{Approuver la demande} et \textit{Informer l’employé de la décision} est correcte si l’employé est toujours notifié après approbation.
    \item \textit{Diagramme 14 :} Corriger la relation entre \textit{Suggérer une formation alternative} et \textit{Traiter une demande de formation} ; si la suggestion est optionnelle, utiliser \texttt{<<extend>>} de \textit{Suggérer une formation alternative} vers \textit{Traiter une demande de formation}.
    \item \textit{Diagramme 14 :} Vérifier la relation \texttt{<<include>>} entre \textit{Suggérer une formation alternative} et \textit{Sélectionner un cours dans le catalogue} ; elle est correcte si la sélection d’un cours est toujours nécessaire pour proposer une alternative.
    \item \textit{Diagramme 14 :} Ajouter un cas \textit{Consulter les demandes de formation} avant \textit{Traiter les demandes de formation}, car l’administrateur RH doit d’abord visualiser les demandes à traiter.
    \item \textit{Diagramme 14 :} Ajouter un cas \textit{Consulter le profil employé} ou \textit{Consulter le profil de compétences} si l’analyse d’adéquation repose sur les données de l’employé.
    \item \textit{Diagramme 14 :} Remplacer \textit{Analyser l’adéquation Profil / Poste cible} par \textit{Analyser l’adéquation entre le profil et le poste cible} pour améliorer la formulation.
    \item \textit{Diagramme 14 :} Vérifier si \textit{Analyser l’adéquation entre le profil et le poste cible} est réalisé par l’administrateur RH ou par le service IA ; si c’est un traitement automatique, le relier au \textit{Service IA} plutôt qu’à l’administrateur.
    \item \textit{Diagramme 14 :} Remplacer les notes longues par des phrases plus courtes ou les déplacer dans le texte explicatif du rapport.
    \item \textit{Diagramme 14 :} Supprimer la note \textit{L’Admin remplace la demande par un autre cours du catalogue} si elle répète déjà le cas \textit{Suggérer une formation alternative}.
    \item \textit{Diagramme 14 :} Relier le \textit{Service de notification} uniquement au cas \textit{Informer l’employé de la décision}, car c’est le seul cas explicitement lié à l’envoi d’une notification.
    \item \textit{Diagramme 14 :} Éviter les longues lignes horizontales entre le service externe et les cas d’utilisation ; rapprocher le service de notification du cas \textit{Informer l’employé de la décision}.
    \item \textit{Diagramme 14 :} Réorganiser le diagramme pour réduire les croisements de lignes, en plaçant les cas de consultation à gauche, les décisions d’approbation/refus au centre et les notifications à droite.
    \item \textit{Diagramme 14 :} Ajouter une légende expliquant les associations simples, les relations \texttt{<<include>>} et les relations \texttt{<<extend>>}.

\end{itemize}
\color{black}
\begin{table}[H]
    \centering
    \renewcommand{\arraystretch}{1.5}
    \small
    \begin{tabularx}{\textwidth}{|>{\raggedright\arraybackslash\bfseries}p{3.5cm}|X|}
        \hline
        \rowcolor[HTML]{F2F2F2}
        Paramètre & \textbf{Description} \\
        \hline
        Titre du Cas & Traiter les demandes de formation et superviser les équipes \\
        \hline
        Acteur Principal & Administrateur (Responsable RH) \\
        \hline
        Acteur Secondaire & Service de Notification (\texttt{SSE}) \\
        \hline
        Préconditions & L'Administrateur est authentifié et possède le rôle \texttt{ROLE\_ADMIN}. Une demande de formation a été soumise par un candidat. \\
        \hline
        Scénario Nominal & 
        \vspace{-3mm}
        \begin{enumerate}[leftmargin=*, nosep]
            \item L'Admin accède à l'espace de Gouvernance (\textit{Dashboard} RH).
            \item Il sélectionne une demande de formation en attente.
            \item Le système affiche l'adéquation entre le cours demandé et le \textit{Skill Gap} du candidat.
            \item L'Admin clique sur "Approuver".
            \item Le \textit{Backend} débloque l'accès au cours pour le candidat.
            \item Le Service de Notification informe le candidat en temps réel de l'approbation.
        \end{enumerate} \\
        \hline
        Scénarios Alternatifs & 
        \vspace{-3mm}
        \begin{itemize}[leftmargin=*, nosep]
            \item \textbf{Suggestion (Alternative) :} L'Admin juge la formation trop complexe. Il clique sur "Suggérer", choisit un cours préparatoire, et le système notifie le candidat du changement de programme.
            \item \textbf{Refus :} L'Admin refuse (manque de budget/temps). Il saisit obligatoirement un motif qui est envoyé au candidat.
        \end{itemize} \\
        \hline
        Post-conditions & Le statut de la demande passe à \texttt{APPROVED}, \texttt{REJECTED} ou \texttt{SUGGESTED}. L'utilisateur est notifié et son accès est actualisé. \\
        \hline
    \end{tabularx}
    \caption{Fiche descriptive du Cas d'Utilisation : Gouvernance et Supervision RH}
    \label{tab:uc_cas2}
\end{table}


\vspace{5mm}

\subsubsection*{Cas 3 : Processus d'Évaluation et de Certification}
Ce cas est le plus abouti du sprint, faisant interagir l'humain et la machine pour garantir l'intégrité de la certification.

\begin{figure}[H]
    \centering
    \includegraphics[width=1\textwidth]{images/case3Sprint3.png}
    \caption{Diagramme  Cas d’Utilisation : Processus d’Évaluation et de Certification}
    \label{fig:case_cas3_sprint3}
\end{figure}


\color{red}
\begin{itemize}
    \item Renommer le titre en \textit{Diagramme de cas d’utilisation : processus d’évaluation et de certification} pour garder une formulation académique plus simple et homogène.
    \item Renommer l’acteur \textit{Utilisateur (Candidat)} en \textit{Candidat} afin d’utiliser un intitulé plus direct et cohérent avec le reste du rapport.
    \item Supprimer les flèches sur les associations entre les acteurs et les cas d’utilisation ; en UML, une association acteur--cas d’utilisation est une ligne simple sans flèche.
    \item Vérifier la relation \texttt{<<include>>} entre \textit{Passer l’évaluation de validation} et \textit{S’authentifier} ; elle est correcte si l’authentification est toujours obligatoire avant l’évaluation.
    \item Vérifier si \textit{Vérifier l’éligibilité (Cours terminé)} doit vraiment apparaître comme un cas d’utilisation ; si c’est une règle métier automatique, il vaut mieux l’intégrer implicitement au cas \textit{Passer l’évaluation de validation} ou la décrire dans le texte.
    \item Corriger la relation entre \textit{Auditer les résultats du test}, \textit{Approuver la certification} et \textit{Refuser et exiger une réévaluation} ; l’audit constitue le cas de base, tandis que l’approbation et le refus doivent être modélisés comme des alternatives.
    \item Vérifier le sens des relations \texttt{<<extend>>} : la flèche doit partir de \textit{Approuver la certification} et de \textit{Refuser et exiger une réévaluation} vers \textit{Auditer les résultats du test}.
    \item Ajouter une relation \texttt{<<include>>} entre \textit{Approuver la certification} et \textit{Notifier l’utilisateur} si la notification est systématique après approbation.
    \item Conserver la relation \texttt{<<include>>} entre \textit{Refuser et exiger une réévaluation} et \textit{Notifier l’utilisateur} uniquement si la notification est obligatoire dans tous les cas de refus.
    \item Corriger la relation entre \textit{Télécharger le certificat} et \textit{Générer le certificat PDF} ; le téléchargement ne doit pas nécessairement inclure la génération du PDF. Il est plus logique que \textit{Approuver la certification} inclue \textit{Générer le certificat PDF}, puis que le candidat puisse le télécharger.
    \item Vérifier si \textit{Corriger automatiquement} et \textit{Mettre à jour le Readiness Score} doivent figurer dans ce diagramme ; ce sont des traitements internes qui peuvent être déplacés vers un diagramme de séquence ou d’activité pour alléger le diagramme de cas d’utilisation.
    \item Vérifier si \textit{Générer dynamiquement le QCM} est formulé du point de vue de l’utilisateur ; si nécessaire, remplacer par un intitulé plus fonctionnel comme \textit{Préparer l’évaluation de validation}.
    \item Relier le \textit{Moteur d’IA} uniquement aux cas qu’il exécute réellement, par exemple \textit{Générer dynamiquement le QCM} et éventuellement \textit{Corriger automatiquement}, afin d’éviter les associations trop larges ou ambiguës.
    \item Ajouter un cas d’utilisation \textit{Consulter le résultat de l’évaluation} ou \textit{Consulter le statut de certification} si cette fonctionnalité existe, car elle n’apparaît pas clairement dans le diagramme.
    \item Réorganiser la disposition des cas d’utilisation pour réduire les croisements, notamment autour de \textit{Télécharger le certificat}, \textit{Notifier l’utilisateur} et \textit{Générer le certificat PDF}.
    \item Ajouter une phrase avant la figure pour préciser qu’elle présente le processus d’évaluation, d’audit, de décision RH et de délivrance du certificat.
    \item Ajouter une phrase après la figure pour expliquer que le candidat passe une évaluation de validation, que l’administrateur RH audite le résultat puis décide d’approuver ou de refuser la certification, avec notification et génération du certificat en cas d’approbation.
\end{itemize}
\color{black}
\begin{table}[H]
    \centering
    \renewcommand{\arraystretch}{1.5}
    \small
    \begin{tabularx}{\textwidth}{|>{\raggedright\arraybackslash\bfseries}p{3.5cm}|X|}
        \hline
        \rowcolor[HTML]{F2F2F2}
        Paramètre & \textbf{Description} \\
        \hline
        Titre du Cas & Évaluer les acquis et certifier la compétence \\
        \hline
        Acteur Principal & Utilisateur (Candidat) \\
        \hline
        Acteurs Secondaires & Administrateur RH, Moteur d'IA, Générateur PDF \\
        \hline
        Préconditions & Le cours est marqué comme "Terminé" dans le \textit{Kanban} de l'utilisateur. \\
        \hline
        Scénario Nominal & 
        \vspace{-3mm}
        \begin{enumerate}[leftmargin=*, nosep]
            \item L'utilisateur lance l'évaluation de fin de cours.
            \item Le système vérifie l'éligibilité, puis le Moteur d'IA génère un QCM dynamique.
            \item L'utilisateur saisit et soumet ses réponses dans le temps imparti.
            \item Le \textit{Backend} corrige, valide le seuil de réussite, et recalcule le \textit{Readiness Score}.
            \item Une demande de certification est envoyée au \textit{Dashboard} de l'Admin RH.
            \item L'Admin audite les résultats et approuve la certification.
            \item Le Moteur PDF génère le certificat officiel et l'utilisateur le télécharge.
        \end{enumerate} \\
        \hline
        Scénarios Alternatifs & 
        \vspace{-3mm}
        \begin{itemize}[leftmargin=*, nosep]
            \item \textbf{Échec au test :} Le score est insuffisant. Le processus de certification est bloqué. Le candidat reçoit un plan de révision et doit repasser le test ultérieurement.
            \item \textbf{Refus RH :} L'Admin détecte une fraude potentielle ou une anomalie de temps de réponse, et exige une réévaluation.
        \end{itemize} \\
        \hline
        Post-conditions & Le \textit{Readiness Score} global du candidat est mis à jour, des points de réussite sont attribués, et le diplôme PDF est stocké. \\
        \hline
    \end{tabularx}
    \caption{Fiche descriptive du Cas d'Utilisation : Évaluation et Certification}
    \label{tab:uc_cas3}
\end{table}

\subsection{Conception }

Cette phase a pour but de traduire les exigences fonctionnelles et les cas d'utilisation identifiés précédemment en solutions techniques concrètes. L'enjeu majeur de ce troisième sprint réside dans la conception d'un système capable d'orchestrer une intelligence artificielle générative tout en préservant la robustesse transactionnelle requise pour la certification et la gouvernance RH.

\paragraph*{1. Objectifs de conception}
\begin{itemize}
    \item \textbf{Interopérabilité :} Concevoir une communication fluide et tolérante aux pannes entre le cœur transactionnel (Java) et le service cognitif (Python).
    \item \textbf{Traçabilité :} Assurer que chaque décision de l'IA (\textit{Roadmap} générée) et chaque décision RH (Approbation/Refus) soient historisées en base de données.
    \item \textbf{Réactivité :} Mettre en place des mécanismes asynchrones (Notifications) pour que les décisions de gouvernance ne bloquent pas l'expérience utilisateur.
\end{itemize}

\subsubsection*{2. Architecture Technique}

\begin{itemize}
    \item \textbf{Client (\texttt{Angular}) :} Responsable du rendu interactif du tableau \textit{Kanban} (\textit{Drag \& Drop}), du chronomètre pour les mini-tests, et de la visualisation des recommandations.
    \item \textbf{\textit{Backend Core} (\texttt{Spring Boot}) :} Joue le rôle de chef d'orchestre. Il sécurise les accès via \texttt{Spring Security} (JWT + RBAC), calcule les \textit{Skill Gaps}, sauvegarde les scores, et utilise \texttt{PDFBox} pour la génération de certificats côté serveur.
    \item \textbf{Microservice IA (\texttt{Python}/\texttt{FastAPI}) :} Isolé du \textit{backend} principal, il prend en charge la construction des \textit{prompts} et la communication avec l'API du LLM en forçant des sorties structurées (JSON).
    \item \textbf{Couche de Communication (SSE) :} Utilisation des \textit{Server-Sent Events} pour "pousser" les notifications de l'Admin vers l'Utilisateur en temps réel, évitant ainsi de surcharger le réseau avec du \textit{Long Polling}.
\end{itemize}

\begin{figure}[H]
    \centering
    % Remplacez 'chemin_architecture_ia_fastapi.png' par le nom exact de votre image
     \includegraphics[width=0.75\textwidth]{images/ArchitectureFastAPI.png}
    \vspace{2mm}
    \caption{Architecture interne du microservice IA FastAPI}
    \label{fig:archi_ia_fastapi}
\end{figure}

Afin d’assurer une séparation claire entre la logique métier \texttt{Java} et les traitements d’intelligence artificielle, une architecture microservices basée sur \texttt{FastAPI} et \texttt{Ollama} a été mise en place. 

Le diagramme de la Figure présente l’architecture détaillée du \textit{pipeline} d’IA, de la sollicitation par l’utilisateur jusqu’à la restitution des prédictions. Ce flux a été conçu pour maximiser la réactivité et la fiabilité du système.

\paragraph*{A. Orchestration et Flux Temps Réel (SSE)}
Contrairement aux architectures classiques où l’utilisateur doit attendre la fin d’un calcul long, \textbf{TalentPredict} implémente un flux de \textit{Streaming} via \textit{Server-Sent Events} (\texttt{SSE}).
\begin{itemize}
    \item Lorsqu'une analyse est lancée, le \textit{Backend} \texttt{Spring Boot} établit une connexion persistante avec le \textit{Frontend} \texttt{Angular}.
    \item Dès que les premiers \textit{tokens} sont générés par le modèle \texttt{Llama 3.2 8B}, ils sont transmis progressivement à l'interface, offrant une expérience utilisateur fluide et dynamique.
\end{itemize}

\paragraph*{B. Microservice FastAPI : Performance et Validation}
Le microservice \texttt{Python} agit comme un filtre d’intelligence et de validation :
\begin{itemize}
    \item \textbf{Validation Stricte (\texttt{Pydantic}) :} Chaque donnée entrante est validée par des schémas \texttt{Pydantic} pour éviter les erreurs de traitement. De même, la réponse de l'IA est rigoureusement vérifiée pour s'assurer qu'elle respecte le format \texttt{JSON} attendu par le \textit{backend} \texttt{Java}.
    \item \textbf{Traitement Asynchrone :} Grâce au moteur \texttt{AsyncIO} de \texttt{FastAPI}, les appels vers \texttt{Ollama} sont non-bloquants, permettant de gérer des files d'attente d'analyses sans dégrader les performances globales.
\end{itemize}

\paragraph*{C. Ingénierie de Prompt et Robustesse}
La "logique métier" de l'IA est encapsulée dans le \textit{Prompt Engine} :
\begin{itemize}
    \item \textbf{Contraintes de Système :} Le modèle est contraint par des \textit{System Prompts} qui définissent son rôle d'expert RH et lui imposent des règles de sortie déterministes.
    \item \textbf{Gestion des Erreurs (\textit{Retry Logic}) :} Le \textit{pipeline} intègre un module de gestion des exceptions. En cas de latence excessive du moteur \texttt{Ollama} ou de réponse malformée, le système déclenche automatiquement une logique de ré-essai (\textit{Retry}), garantissant la continuité du service.
\end{itemize}

\paragraph*{D. Sécurité et Inférence Locale}
La sécurité est assurée à deux niveaux :
\begin{itemize}
    \item \textbf{Authentification Inter-Services :} Les échanges entre \texttt{Spring Boot} et \texttt{FastAPI} sont sécurisés par des jetons \texttt{JWT} internes, empêchant toute sollicitation non autorisée du moteur d'IA.
    \item \textbf{Souveraineté des Données :} L'utilisation d'\texttt{Ollama} permet une inférence locale optimisée. Les données sensibles des collaborateurs restent confinées dans le réseau privé \texttt{Docker}, garantissant une conformité totale avec le \textbf{RGPD}.
\end{itemize}
\subsubsection*{3. Spécification des API du Microservice IA}

Le microservice IA expose plusieurs \textit{endpoints} REST permettant l’interaction avec le moteur de génération :

\begin{itemize}
    \item \textbf{\texttt{GET /api/formations/recommendations}} : Liste les formations suggérées en fonction des lacunes identifiées.
    \item \textbf{\texttt{GET /api/predictions/roadmap}} : Récupère la trajectoire de carrière générée par l'IA.
    \item \textbf{\texttt{POST /api/predictions/generate}} : Déclenche une nouvelle analyse du \textit{Gap} par le moteur \texttt{Llama 3}.
\end{itemize}

Chaque \textit{endpoint} retourne une réponse strictement formatée en \texttt{JSON} afin d’assurer la compatibilité avec le \textit{backend} \texttt{Spring Boot} et le \textit{frontend} \texttt{Angular}.
\subsubsection*{4. Validation et Robustesse des Réponses IA}

L'un des défis majeurs lors de l'intégration de modèles de langage (LLM) comme \texttt{Llama 3} est de transformer une réponse textuelle probabiliste en une donnée structurée et fiable pour une application d'entreprise. Pour garantir cette robustesse, \textbf{TalentPredict} implémente une stratégie de validation à plusieurs niveaux.

\paragraph*{A. Le Forçage de Structure (\textit{JSON Mode} \& Pydantic)}
Pour éviter les "hallucinations" de format, nous utilisons des schémas \texttt{Pydantic} qui définissent contractuellement la structure attendue (champs obligatoires, types de données, énumérations).
\begin{itemize}
    \item Le microservice \texttt{FastAPI} impose au modèle \texttt{Llama 3} un \texttt{JSON Schema} strict.
    \item Si la réponse de l'IA ne respecte pas ce schéma (ex: champ manquant ou type erroné), la donnée est rejetée avant même d'atteindre le \textit{backend} \texttt{Java}, évitant ainsi des erreurs de désérialisation en cascade.
\end{itemize}

\paragraph*{B. Déterminisme et \textit{Prompt Engineering}}
La robustesse passe également par la réduction de la "créativité" non désirée du modèle :
\begin{itemize}
    \item \textbf{Température basse :} Nous avons configuré la température du modèle à une valeur proche de 0 pour favoriser des réponses constantes et déterministes.
    \item \textbf{\textit{System Prompting} :} Des instructions système explicites interdisent au modèle d'ajouter des commentaires textuels (ex: "Voici votre analyse :") en dehors du bloc \texttt{JSON}, garantissant une sortie "pure" et directement \textit{parsable}.
\end{itemize}

\paragraph*{C. Mécanismes de Résilience (\textit{Retry \& Fallback})}
Malgré ces précautions, les modèles d'IA peuvent parfois échouer ou subir des latences. Nous avons donc mis en place :
\begin{itemize}
    \item \textbf{Stratégie de Re-essai (\textit{Retry Logic}) :} En cas de réponse invalide, le système effectue automatiquement jusqu'à trois tentatives avec un ajustement léger du \textit{prompt} pour corriger l'erreur.
    \item \textbf{Gestion des \textit{Timeouts} :} Pour garantir la disponibilité de l'application, un délai d'attente maximum est défini. Si l'IA ne répond pas dans le temps imparti, le système bascule sur un mode dégradé (affichage des dernières données connues) et alerte l'administrateur.
\end{itemize}

\paragraph*{D. Validation Métier (\textit{Sanity Checks})}
Une fois le \texttt{JSON} extrait, le \textit{backend} \texttt{Java} effectue des vérifications métier finales :
\begin{itemize}
    \item Validation de la cohérence des scores (de 0 à 100\%).
    \item Vérification de l'existence des compétences référencées. Cela assure que l'intelligence artificielle reste un outil d'aide à la décision fiable et non une source d'erreurs logiques.
\end{itemize}

\subsubsection*{5. Gestion Globale des Erreurs du Système IA}

Le système implémente une gestion centralisée des erreurs afin de garantir la stabilité globale de l’application. Les principaux scénarios couverts sont :

\begin{itemize}
    \item \textbf{\textit{Timeout} des requêtes IA} : Dépassement du temps d'attente autorisé lors de la génération.
    \item \textbf{Réponses \texttt{JSON} invalides} : Échec de la validation structurelle par les schémas \texttt{Pydantic}.
    \item \textbf{Indisponibilité du modèle \texttt{LLM}} : Déconnexion ou surcharge du moteur \texttt{Ollama}.
    \item \textbf{Erreurs réseau} : Problèmes de communication entre les conteneurs \texttt{Spring Boot} et \texttt{FastAPI}.
\end{itemize}

Dans ces cas, le système applique une stratégie de \textit{fallback} permettant de retourner une réponse contrôlée ou les dernières données valides disponibles, évitant ainsi un blocage de l'interface utilisateur.

\subsubsection*{6. Modèle de données : }


\begin{figure}[H]
    \centering
    % Remplacez 'chemin_diagramme_erd.png' par le nom exact de votre image
    \includegraphics[width=1\textwidth]{images/ERDsprint3.png}
    \vspace{2mm}
    \caption{Diagramme Entité-Relation (ERD) de l'architecture TalentPredict}
    \label{fig:erd_talentpredict_s3}
\end{figure}


\color{red}
\begin{itemize}
    \item \textit{Diagramme 15 :} Renommer le titre en \textit{Diagramme entité-relation de l’architecture de persistance de TalentPredict AI} pour une formulation plus académique et cohérente.
    \item \textit{Diagramme 15 :} Remplacer le champ \textit{password} dans la table \textit{USERS} par \textit{password\_hash}, afin d’indiquer que le mot de passe ne doit jamais être stocké en clair.
    \item \textit{Diagramme 15 :} Vérifier si l’attribut \textit{email} est bien unique ; la contrainte \textit{UK} est correcte, mais elle doit aussi être mentionnée dans le texte explicatif.
    \item \textit{Diagramme 15 :} Ajouter les champs \textit{created\_at} et \textit{updated\_at} dans les tables principales pour assurer la traçabilité des données.
    \item \textit{Diagramme 15 :} Ajouter un champ \textit{phone} dans la table \textit{USERS} si le numéro de téléphone est utilisé dans l’inscription ou la récupération de compte.
    \item \textit{Diagramme 15 :} Vérifier la relation entre \textit{USERS} et \textit{PROFILES} ; si chaque utilisateur possède un seul profil, la cardinalité doit être clairement indiquée en \textit{1:1}.
    \item \textit{Diagramme 15 :} Corriger la relation entre \textit{PROFILES} et \textit{PCM\_RESULTS} ; si un profil peut avoir plusieurs résultats PCM dans le temps, la cardinalité \textit{1:N} est correcte, sinon elle doit être \textit{1:1}.
    \item \textit{Diagramme 15 :} Vérifier si \textit{PCM\_RESULTS} doit être reliée directement à \textit{USERS} plutôt qu’à \textit{PROFILES}, selon la logique réelle du backend.
    \item \textit{Diagramme 15 :} Dans la table \textit{SKILLS}, remplacer \textit{nom} par \textit{name} ou uniformiser tous les noms d’attributs en français ; il faut éviter le mélange français--anglais.
    \item \textit{Diagramme 15 :} Dans la table \textit{SKILLS}, préciser clairement les valeurs possibles de \textit{type}, par exemple \textit{TECH} et \textit{SOFT}.
    \item \textit{Diagramme 15 :} Dans la table \textit{SKILLS}, ajouter un champ \textit{confidence\_score} si le système calcule un niveau de confiance à partir des sources CV, GitHub ou IA.
    \item \textit{Diagramme 15 :} Dans la table \textit{SKILLS}, vérifier si l’attribut \textit{niveau} doit rester un entier de 1 à 5 ou être remplacé par une énumération plus lisible.
    \item \textit{Diagramme 15 :} Dans la table \textit{FORMATIONS}, ajouter les champs utiles comme \textit{description}, \textit{level}, \textit{duration}, \textit{url}, \textit{status} et \textit{progress}, si le suivi de formation est réellement géré.
    \item \textit{Diagramme 15 :} Vérifier la relation entre \textit{USERS} et \textit{FORMATIONS} ; si une formation peut être suivie par plusieurs utilisateurs, il faut ajouter une table d’association comme \textit{USER\_FORMATIONS}.
    \item \textit{Diagramme 15 :} Dans \textit{USER\_NOTIFICATIONS}, ajouter les champs \textit{type}, \textit{created\_at} et éventuellement \textit{read\_at} pour mieux gérer l’historique des notifications.
    \item \textit{Diagramme 15 :} Dans \textit{PREDICTIONS}, remplacer le champ \textit{insights} de type \textit{text} par un type \textit{json} si les résultats d’analyse IA sont structurés.
    \item \textit{Diagramme 15 :} Dans \textit{PREDICTIONS}, ajouter un champ \textit{score} ou \textit{overall\_score} si la prédiction produit une mesure globale.
    \item \textit{Diagramme 15 :} Vérifier si \textit{CANDIDATE\_TEST\_RESULTS} doit être reliée à \textit{SKILLS} par un champ \textit{skill\_id}, car un résultat de test concerne généralement une compétence précise.
    \item \textit{Diagramme 15 :} Ajouter un champ \textit{test\_type} dans \textit{CANDIDATE\_TEST\_RESULTS}, par exemple \textit{QCM}, \textit{CODE}, \textit{PCM} ou \textit{SCENARIO}.
    \item \textit{Diagramme 15 :} Ajouter un champ \textit{created\_at} dans \textit{CANDIDATE\_TEST\_RESULTS} pour historiser les passages de tests.
    \item \textit{Diagramme 15 :} Vérifier si \textit{REFRESH\_TOKENS} doit contenir au moins les champs \textit{id}, \textit{token}, \textit{expiry\_date}, \textit{revoked} et \textit{user\_id}.
    \item \textit{Diagramme 15 :} Vérifier si \textit{PASSWORD\_RESET\_TOKENS} doit contenir au moins les champs \textit{id}, \textit{token}, \textit{expiry\_date}, \textit{used} et \textit{user\_id}.
    \item \textit{Diagramme 15 :} Ajouter une table \textit{RECOMMENDATIONS} si les recommandations de formation sont stockées et consultées ultérieurement.
    \item \textit{Diagramme 15 :} Ajouter une table \textit{TALENT\_PASSPORTS} si le passeport de compétences est généré, téléchargé ou historisé.
    \item \textit{Diagramme 15 :} Ajouter une table \textit{ROADMAPS} ou \textit{LEARNING\_PATHS} si le système gère des parcours d’upskilling personnalisés.
    \item \textit{Diagramme 15 :} Vérifier si les cardinalités sont toutes lisibles ; certaines relations sont longues et difficiles à interpréter dans la version actuelle.
    \item \textit{Diagramme 15 :} Réorganiser le diagramme pour rapprocher les tables fortement liées, par exemple \textit{PROFILES} avec \textit{PCM\_RESULTS}, et \textit{SKILLS} avec \textit{CANDIDATE\_TEST\_RESULTS}.
    \item \textit{Diagramme 15 :} Réduire les longues relations partant de \textit{USERS} vers toutes les tables, car elles rendent le diagramme difficile à lire.
    \item \textit{Diagramme 15 :} Ajouter une légende expliquant les notations \textit{PK}, \textit{FK}, \textit{UK}, ainsi que les cardinalités \textit{1:1} et \textit{1:N}.

\end{itemize}
\color{black}
Le modèle de données de \textbf{TalentPredict} a été conçu pour supporter une architecture orientée "Intelligence de Données". Le diagramme ERD illustre la manière dont les informations sont centralisées autour de l'utilisateur pour permettre une analyse prédictive à 360°.

\paragraph*{A. Le Pivot Central : Identité et Gamification (Table \texttt{USERS})}
L'entité \texttt{USERS} est le cœur du système. Outre les informations d'authentification classique (Email, Mot de passe, Rôles), elle intègre des mécanismes de \textit{Gamification} via les attributs \texttt{xp} et \texttt{level}. Cette approche permet de suivre l'engagement du collaborateur au fur et à mesure qu'il complète des évaluations ou des formations.

\paragraph*{B. Le Profil Professionnel et l'Analyse IA (Table \texttt{PROFILES})}
Le \texttt{PROFILE} est une extension de l'utilisateur (relation 1:1) qui centralise les vecteurs d'entrée pour l'intelligence artificielle :
\begin{itemize}
    \item \textbf{Vecteurs Sources :} URLs \texttt{LinkedIn}, \texttt{GitHub} et chemin vers le CV numérique.
    \item \textbf{Intelligence de Stockage :} Le champ \texttt{ai\_summary} contient la synthèse textuelle générée par le modèle \texttt{Llama 3}, tandis que \texttt{skill\_real\_scores\_json} stocke les scores de \textit{matching} sous forme de données structurées, permettant une flexibilité totale dans l'affichage des résultats.
\end{itemize}

\paragraph*{C. Évaluation Multidimensionnelle (Tables \texttt{SKILLS}, \texttt{PCM\_RESULTS} \& \texttt{TESTS})}
La force de \textbf{TalentPredict} réside dans la séparation des types de compétences :
\begin{itemize}
    \item \texttt{SKILLS} : Stocke à la fois les \textit{Hard} et \textit{Compétences comportementales}, en précisant leur Source (CV, GitHub ou analyse IA). Cela permet de tracer l'origine de chaque compétence détectée.
    \item \texttt{PCM\_RESULTS} : Table dédiée aux résultats psychométriques (\textit{Process Communication Model}). Elle stocke le profil de personnalité, essentiel pour l'analyse des écarts comportementaux.
    \item \texttt{CANDIDATE\_TEST\_RESULTS} : Enregistre les scores bruts des tests techniques passés sur la plateforme.
\end{itemize}

\paragraph*{D. \textit{Upskilling} et Plan de Carrière (Tables \texttt{PREDICTIONS} \& \texttt{FORMATIONS})}
Ces deux entités représentent l'aboutissement de l'intelligence métier :
\begin{itemize}
    \item \texttt{PREDICTIONS} : Stocke l'analyse du \textit{Gap} (écart) produite par l'IA. Elle définit le chemin de carrière et les recommandations stratégiques.
    \item \texttt{FORMATIONS} : Agit comme un catalogue personnalisé. Une recommandation IA pointe directement vers un module de formation pour combler un écart de compétence spécifique.
\end{itemize}

\paragraph*{E. Hub de Communication (Table \texttt{USER\_NOTIFICATIONS})}
Cette table assure la liaison en temps réel. Elle archive chaque événement important (recommandation prête, test assigné), permettant une synchronisation parfaite entre le \textit{backend} et l'interface utilisateur via les flux asynchrones (\texttt{WebSockets}/\texttt{SSE}).


\subsubsection*{7. Algorithme de Calcul du Readiness Score: }
Dans notre architecture, le \textbf{Readiness Score} (Taux de Préparation) est 
l'indicateur central qui permet de quantifier l'adéquation technique d'un candidat 
par rapport aux exigences d'un poste cible. Afin de garantir un calcul robuste, 
transparent et vérifiable, nous avons implémenté l'algorithme suivant dans le 
backend Java (\texttt{CareerController.java}) :

\begin{equation}
R_{brut} = \left( \frac{\sum_{i=1}^{n} Niveau_i}{n \times 10} \right) \times 100
\end{equation}

\noindent Définition des variables :
\begin{itemize}
    \item $R_{brut}$ : le Readiness Score théorique exprimé en pourcentage.
    \item $Niveau_i$ : le niveau de maîtrise actuel de l'utilisateur pour la 
    compétence $i$, évalué sur une échelle de 1 à 10.
    \item $n$ : le nombre total de compétences évaluées dans la Skill Gap Analysis.
\end{itemize}

\noindent Pour éviter les valeurs aberrantes, le résultat brut est soumis à une 
fonction de bornage (\textit{clamping}) avant sa persistance en base de données :

\begin{equation}
Readiness\_Score = \max\!\left(10,\ \min\!\left(95,\ \text{arrondi}(R_{brut})\right)\right)
\end{equation}

\noindent Cette borne garantit qu'aucun collaborateur ne se voit attribuer un score 
décourageant de 0\% ni une perfection irréaliste de 100\%.


\subsubsection*{8. Contrats API (Exemple de Conception): }
Afin de garantir la séparation entre le \textit{Frontend} et le \textit{Backend}, des contrats d'interface stricts ont été définis. Voici un exemple illustrant la requête et la réponse pour la génération asynchrone d'un mini-test :

\textbf{Point d'accès :} \texttt{POST /api/assessment/mini-test/generate}

\textbf{Requête (\textit{Payload}) :}
\begin{lstlisting}[language=json, caption={Payload de la requête de génération de test}]
{
  "candidateId": "123e4567-e89b-12d3-a456-426614174000",
  "targetSkill": "Spring Boot",
  "currentLevel": 4,
  "difficulty": "INTERMEDIATE"
}
\end{lstlisting}

\textbf{Réponse (200 OK - \texttt{TestDto}) :}
\begin{lstlisting}[language=json, caption={Structure de réponse d'un mini-test généré}]
{
  "sessionId": "eval-9048-abcd",
  "timeLimitSeconds": 600,
  "totalQuestions": 5,
  "questions": [
    {
      "questionId": "q-101",
      "text": "Quelle annotation Spring Boot expose des endpoints REST ?",
      "options": [
        { "key": "A", "value": "@Controller" },
        { "key": "B", "value": "@RestController" },
        { "key": "C", "value": "@ApiComponent" },
        { "key": "D", "value": "@RestApi" }
      ]
    }
  ]
}

\end{lstlisting}


\subsubsection*{9. Diagramme de Séquence}

\subsubsection*{ Cas 1: Planification de l'Apprentissage}

Ce diagramme illustre le parcours de bout en bout d'un employé souhaitant planifier sa montée en compétences. Sur le plan architectural, la sécurité est garantie dès l'entrée par le filtre \texttt{Spring Security} qui valide le \textit{token} \texttt{JWT}. Le \textit{Backend} (\texttt{Java}) agit comme un orchestrateur : avant de solliciter l'intelligence artificielle, il récupère le profil de l'employé en base de données pour cibler ses lacunes (\textit{Skill Gap}). Cette donnée est envoyée au Microservice IA (\texttt{Python}), qui se charge de l'ingénierie de \textit{prompt} et interroge l'API du modèle LLM. La réponse structurée renvoyée par l'IA est ensuite convertie en entités relationnelles sauvegardées dans \texttt{PostgreSQL}. Le processus se termine par l'initialisation du tableau \textit{Kanban} et l'attribution de points d'initiative (\textit{Gamification}), démontrant une intégration complète entre l'IA et la logique métier de la plateforme.

\begin{figure}[H]
    \centering
    % Remplacez 'chemin_image_seq_cas1' par le nom exact de votre image
    \includegraphics[width=0.95\textwidth]{images/seq1Sprint3.png}
    \vspace{2mm}
    \caption{Diagramme de séquence : lanification de l'Apprentissage}
    \label{fig:seq_tech_cas1}
\end{figure}

\color{red}
\begin{itemize}
    \item \textit{Diagramme 16 :} Corriger le titre en \textit{Diagramme de séquence : planification de l’apprentissage}, car le mot \textit{lanification} contient une faute.
    \item \textit{Diagramme 16 :} Alléger le titre en supprimant \textit{Très détaillé -- Cas 1}, qui peut être conservé dans le texte du rapport mais pas nécessairement dans la figure.
    \item \textit{Diagramme 16 :} Renommer l’acteur \textit{Utilisateur (Employé)} en \textit{Employé} afin de garder une terminologie homogène avec les autres diagrammes.
    \item \textit{Diagramme 16 :} Séparer la ligne de vie \textit{Controllers \& Services (Java)} en deux lignes distinctes : \textit{Controller} et \textit{Service métier}, afin de mieux respecter la logique d’un diagramme de séquence.
    \item \textit{Diagramme 16 :} Renommer \textit{Security Filter} en \textit{Filtre de sécurité} ou garder tous les noms techniques en anglais ; il faut éviter le mélange non maîtrisé entre français et anglais.
    \item \textit{Diagramme 16 :} Remplacer \textit{Base de Données (PostgreSQL)} par \textit{Base de données PostgreSQL}, avec une formulation plus standard.
    \item \textit{Diagramme 16 :} Vérifier si le \textit{Microservice IA} appelle directement le \textit{Modèle LLM} ; si oui, cette interaction est correctement représentée, mais elle doit être décrite dans le texte.
    \item \textit{Diagramme 16 :} Ajouter un fragment \texttt{alt} après la validation du jeton JWT pour distinguer le cas \textit{jeton valide} du cas \textit{jeton invalide}.
    \item \textit{Diagramme 16 :} Ajouter un retour d’erreur possible \textit{401 Unauthorized} lorsque le jeton JWT est invalide ou expiré.
    \item \textit{Diagramme 16 :} Ajouter un fragment \texttt{alt} lors de la génération de la roadmap pour distinguer le cas \textit{réponse LLM valide} du cas \textit{réponse LLM invalide ou non exploitable}.
    \item \textit{Diagramme 16 :} Ajouter une étape de validation de la réponse JSON retournée par le modèle LLM avant l’enregistrement de la roadmap.
    \item \textit{Diagramme 16 :} Ajouter une étape de gestion d’erreur si la réponse JSON du modèle LLM n’est pas conforme au format attendu.
    \item \textit{Diagramme 16 :} Vérifier si le message \textit{POST /v1/chat/completions} correspond réellement à l’API utilisée par le modèle local ; sinon, le remplacer par un appel plus générique vers le service LLM.
    \item \textit{Diagramme 16 :} Remplacer \textit{Traitement IA basé sur le contexte métier} par une note plus courte, par exemple \textit{Génération IA de la roadmap selon le profil cible}.
    \item \textit{Diagramme 16 :} Ajouter une étape de récupération du profil utilisateur ou des compétences existantes avant \textit{fetchSkillGapAndRecommendations}, si ces données sont nécessaires au calcul du skill gap.
    \item \textit{Diagramme 16 :} Vérifier que \textit{fetchSkillGapAndRecommendations(userId)} retourne bien à la fois les écarts de compétences et les cours suggérés ; sinon, séparer ces deux récupérations.
    \item \textit{Diagramme 16 :} Remplacer \textit{Data (Skill Gap, Cours suggérés)} par \textit{Écarts de compétences et cours suggérés} pour améliorer la lisibilité.
    \item \textit{Diagramme 16 :} Ajouter une étape explicite de sélection ou de saisie du \textit{poste cible} avant l’appel de génération de roadmap.
    \item \textit{Diagramme 16 :} Remplacer \textit{targetRole} par \textit{posteCible} ou harmoniser tous les paramètres en anglais, selon la convention adoptée dans le rapport.
    \item \textit{Diagramme 16 :} Vérifier si l’initialisation du Kanban doit toujours suivre la génération de la roadmap ; si ce n’est pas automatique, représenter cette partie dans un fragment \texttt{opt}.
    \item \textit{Diagramme 16 :} Ajouter un fragment \texttt{opt} autour de la notification \textit{Badge débloqué}, car cette notification dépend d’une condition de progression.
    \item \textit{Diagramme 16 :} Ajouter clairement la condition de déblocage du badge, par exemple \textit{si premier cours ajouté au Kanban}.
    \item \textit{Diagramme 16 :} Vérifier si \textit{calculerPointsInitiative(userId)} et \textit{updateLeaderboardScore(userId, points)} relèvent du processus principal ; sinon, les déplacer vers un diagramme dédié à la gamification.
    \item \textit{Diagramme 16 :} Ajouter des barres d’activation cohérentes pour toutes les lignes de vie qui exécutent des traitements, notamment \textit{Microservice IA}, \textit{Modèle LLM} et \textit{Base de données}.
    \item \textit{Diagramme 16 :} Vérifier que toutes les réponses sont représentées par des flèches de retour en pointillés et que les appels principaux restent en lignes pleines.
    \item \textit{Diagramme 16 :} Réduire la quantité de texte dans les messages longs afin de garder le diagramme lisible dans le rapport.
   
\end{itemize}
\color{black}
\subsubsection*{Cas 2: Gouvernance et Supervision RH }
Ce diagramme met en évidence les mécanismes de contrôle et de sécurité du module de gouvernance. Contrairement à un utilisateur standard, l'accès aux décisions RH nécessite un contrôle \textbf{RBAC} (\textit{Role-Based Access Control}) strict, modélisé ici par la vérification du privilège \texttt{ROLE\_ADMIN} par le filtre de sécurité. La complexité métier est gérée via un bloc transactionnel conditionnel (\texttt{alt}). Selon la décision de l'administrateur, le système exécute des actions distinctes (comme le déblocage de l'accès au cours en cas d'approbation). Enfin, pour garantir une expérience utilisateur fluide sans rechargement de page, le système déclenche un événement asynchrone via le service de notification (basé sur les \textit{Server-Sent Events} ou \texttt{WebSockets}), informant instantanément le candidat de la décision RH.

\begin{figure}[H]
    \centering
    % Remplacez 'chemin_image_seq_cas2' par le nom exact de votre image
    \includegraphics[width=1\textwidth]{images/seq2Sprint3.png}
    \vspace{2mm}
    \caption{Diagramme de séquence : Gouvernance et Supervision RH}
    \label{fig:seq_tech_cas2}
\end{figure}

\color{red}
\begin{itemize}
    \item \textit{Diagramme 17 :} Renommer le titre en \textit{Diagramme de séquence : gouvernance et supervision RH} pour garder une formulation plus simple et homogène.
    \item \textit{Diagramme 17 :} Renommer l’acteur \textit{Admin (RH)} en \textit{Administrateur RH} afin d’utiliser la même terminologie dans tout le rapport.
    \item \textit{Diagramme 17 :} Remplacer \textit{Security Filter} par \textit{Filtre de sécurité} ou garder tous les noms techniques en anglais ; il faut éviter un mélange non maîtrisé entre français et anglais.
    \item \textit{Diagramme 17 :} Séparer la ligne de vie \textit{GovernanceService} en deux éléments si nécessaire : un contrôleur \textit{GovernanceController} et un service métier \textit{GovernanceService}.
    \item \textit{Diagramme 17 :} Remplacer \textit{Base de Données (PostgreSQL)} par \textit{Base de données PostgreSQL} pour une formulation plus standard.
    \item \textit{Diagramme 17 :} Ajouter un fragment \texttt{alt} après la validation du jeton JWT pour distinguer le cas \textit{jeton valide avec rôle administrateur} du cas \textit{jeton invalide ou rôle non autorisé}.
    \item \textit{Diagramme 17 :} Ajouter un retour d’erreur possible \textit{401 Unauthorized} ou \textit{403 Forbidden} lorsque le jeton est invalide ou lorsque le rôle n’est pas \texttt{ROLE\_ADMIN}.
    \item \textit{Diagramme 17 :} Remplacer \textit{Valider JWT \& Vérifier ROLE\_ADMIN} par \textit{Valider le jeton JWT et vérifier le rôle administrateur} pour améliorer la lisibilité.
    \item \textit{Diagramme 17 :} Remplacer \textit{fetchPendingRequestsAndKPIs()} par \textit{récupérerDemandesEtIndicateurs()} ou harmoniser tous les messages en anglais selon la convention adoptée dans le rapport.
    \item \textit{Diagramme 17 :} Remplacer \textit{Liste des demandes \& Statistiques d’équipe} par \textit{Liste des demandes et indicateurs d’équipe}.
    \item \textit{Diagramme 17 :} Ajouter une étape explicite de sélection d’une demande par l’administrateur avant l’appel \textit{processDecision(reqId, action, motif)}.
    \item \textit{Diagramme 17 :} Corriger \textit{reqId} en \textit{requestId} ou utiliser \textit{idDemande} afin d’éviter une abréviation peu claire dans un diagramme académique.
    \item \textit{Diagramme 17 :} Ajouter un fragment \texttt{alt} ou \texttt{opt} pour vérifier que le motif est obligatoire uniquement lorsque l’action correspond à un refus.
    \item \textit{Diagramme 17 :} Dans la branche \textit{APPROUVER}, préciser si \textit{unlockCourseAccess(userId, courseId)} est toujours exécuté ; si ce n’est pas systématique, le placer dans un fragment \texttt{opt}.
    \item \textit{Diagramme 17 :} Dans la branche \textit{SUGGÉRER}, préciser comment \textit{newCourseId} est choisi ; ajouter une étape \textit{sélectionner une formation alternative} avant la mise à jour du statut.
    \item \textit{Diagramme 17 :} Dans la branche \textit{REFUSER}, vérifier que le motif est bien transmis et stocké avec le statut \textit{REJECTED}.
    \item \textit{Diagramme 17 :} Ajouter un retour de la base de données après chaque opération \textit{updateStatus} et \textit{unlockCourseAccess} afin de montrer que l’opération a réussi ou échoué.
    \item \textit{Diagramme 17 :} Remplacer les mentions \textit{Transaction Commit} par \textit{Validation de la transaction} ou les déplacer dans le texte technique si elles alourdissent le diagramme.
    \item \textit{Diagramme 17 :} Ajouter un fragment \texttt{alt} autour du traitement de la décision pour distinguer le cas \textit{opération réussie} du cas \textit{erreur lors de la mise à jour}.
    \item \textit{Diagramme 17 :} Vérifier si la notification doit être envoyée dans les trois cas \textit{APPROUVER}, \textit{SUGGÉRER} et \textit{REFUSER} ; si oui, le message \textit{pushDecisionNotification} est bien placé après le fragment \texttt{alt}.
    \item \textit{Diagramme 17 :} Remplacer \textit{NotificationService (SSE / WebSockets)} par \textit{Service de notification temps réel} pour garder une formulation plus claire.
    \item \textit{Diagramme 17 :} Corriger la note \textit{Envoi de l’événement en temps réel au navigateur du Candidat} si le destinataire est un employé ; utiliser le même acteur que dans le reste du rapport.
    \item \textit{Diagramme 17 :} Ajouter une réponse du service de notification vers le service métier indiquant que l’événement a été publié.
    \item \textit{Diagramme 17 :} Ajouter des barres d’activation cohérentes sur toutes les lignes de vie, notamment \textit{Base de données PostgreSQL} et \textit{NotificationService}.
    \item \textit{Diagramme 17 :} Vérifier que les flèches de retour sont bien en pointillés et que les appels principaux restent en lignes pleines.
    \item \textit{Diagramme 17 :} Réduire la taille des messages longs afin de rendre le diagramme plus lisible dans une page de rapport.

\end{itemize}
\color{black}
\subsubsection*{Cas 3 Évaluation et Certification}
Ce diagramme de séquence est le plus riche car il fait collaborer la logique transactionnelle, l'intelligence artificielle et un moteur de rendu documentaire. Le système impose d'abord une vérification d'éligibilité (le cours doit être terminé) avant de demander au microservice IA de générer un test sur mesure. Pour prévenir toute fraude, la correction du test est traitée exclusivement côté serveur (\textit{Backend}). Si le seuil de réussite est atteint, le \textit{Readiness Score} du talent est immédiatement mis à jour en base de données. Enfin, le processus illustre le concept de \textit{Human-in-the-loop} : la génération du diplôme officiel au format PDF n'est jamais automatique. Elle nécessite l'approbation explicite de l'Administrateur RH pour déclencher l'utilisation de la librairie PDF (\texttt{PDFBox}), garantissant l'intégrité de la certification.

\begin{figure}[H]
    \centering
    % Remplacez 'chemin_image_seq_cas3' par le nom exact de votre image
    \includegraphics[width=1\textwidth]{images/seq3Sprint3.png}
    \vspace{2mm}
    \caption{Diagramme de séquence :  Évaluation et Certification }
    \label{fig:seq_tech_cas3}
\end{figure}

\color{red}
\begin{itemize}
    \item \textit{Diagramme 18 :} Renommer le titre en \textit{Diagramme de séquence : évaluation et certification} pour alléger l’intitulé.
    \item \textit{Diagramme 18 :} Renommer l’acteur \textit{Candidat (Employé)} en \textit{Employé} ou \textit{Candidat}, puis utiliser le même terme dans tout le rapport.
    \item \textit{Diagramme 18 :} Renommer l’acteur \textit{Admin (Responsable RH)} en \textit{Administrateur RH} afin d’harmoniser la terminologie avec les autres diagrammes.
    \item \textit{Diagramme 18 :} Remplacer \textit{Security Filter} par \textit{Filtre de sécurité} ou garder tous les noms techniques en anglais ; il faut éviter le mélange non maîtrisé entre français et anglais.
    \item \textit{Diagramme 18 :} Séparer la ligne de vie \textit{Backend Services} en deux éléments si nécessaire : \textit{EvaluationController} et \textit{EvaluationService}, afin de mieux représenter la séparation entre contrôleur et logique métier.
    \item \textit{Diagramme 18 :} Remplacer \textit{Base de Données (PostgreSQL)} par \textit{Base de données PostgreSQL} pour garder une formulation plus standard.
    \item \textit{Diagramme 18 :} Ajouter un fragment \texttt{alt} après la validation du JWT pour distinguer le cas \textit{jeton valide} du cas \textit{jeton invalide ou expiré}.
    \item \textit{Diagramme 18 :} Ajouter un retour d’erreur \textit{401 Unauthorized} lorsque le jeton JWT est invalide ou absent.
    \item \textit{Diagramme 18 :} Ajouter un fragment \texttt{alt} après \textit{verifyEligibility(userId, courseId)} pour distinguer le cas \textit{cours terminé} du cas \textit{cours non terminé}.
    \item \textit{Diagramme 18 :} Ajouter un retour d’erreur ou de refus lorsque le cours n’est pas terminé, car l’évaluation ne doit pas être lancée dans ce cas.
    \item \textit{Diagramme 18 :} Vérifier si le microservice IA est réellement responsable de la génération du test ; si oui, l’appel \textit{POST /ai/test/generate} est cohérent.
    \item \textit{Diagramme 18 :} Ajouter une étape de validation du test structuré retourné par le microservice IA avant de l’envoyer au frontend.
    \item \textit{Diagramme 18 :} Ajouter une étape d’enregistrement de la session de test dans la base de données avant le retour \textit{TestSessionDto}.
    \item \textit{Diagramme 18 :} Remplacer \textit{Test Structuré (JSON)} par \textit{Test structuré au format JSON} pour améliorer la lisibilité.
    \item \textit{Diagramme 18 :} Remplacer la note \textit{L’IA génère un QCM inédit basé sur le cours} par une note plus courte, par exemple \textit{Génération IA d’un QCM adapté au cours}.
    \item \textit{Diagramme 18 :} Dans la phase de soumission des réponses, ajouter une étape de vérification du délai ou du chronomètre si le test est limité dans le temps.
    \item \textit{Diagramme 18 :} Vérifier si la correction est entièrement faite par le backend ou par le microservice IA ; si le microservice intervient, ajouter une ligne de message vers \textit{Microservice IA}.
    \item \textit{Diagramme 18 :} Ajouter un retour de la base de données après \textit{updateReadinessScore(userId, newScore)} afin de montrer que le score a bien été enregistré.
    \item \textit{Diagramme 18 :} Remplacer \textit{Score actualisé} par \textit{Score de préparation actualisé} pour clarifier le sens métier.
    \item \textit{Diagramme 18 :} Dans le fragment \texttt{alt}, remplacer \textit{Score >= Seuil de réussite} et \textit{Score < Seuil de réussite} par des conditions plus lisibles en français.
    \item \textit{Diagramme 18 :} Vérifier si \textit{createPendingCertification()} doit être enregistré en base de données ; si oui, ajouter un message vers la base et un retour de confirmation.
    \item \textit{Diagramme 18 :} Dans le cas d’échec, préciser si \textit{Préparer un plan de révision} est généré par le backend ou par le microservice IA.
    \item \textit{Diagramme 18 :} La partie d’approbation RH commence après l’affichage du résultat candidat ; il serait plus lisible de la placer dans un fragment ou dans un diagramme séparé consacré à la validation RH.
    \item \textit{Diagramme 18 :} Ajouter une validation du rôle administrateur avant \textit{approveCertification(certId)}, avec un cas d’erreur \textit{403 Forbidden} si le rôle n’est pas autorisé.
    \item \textit{Diagramme 18 :} Ajouter un retour de la base de données après \textit{updateStatus(APPROVED)} et \textit{saveCertificatePath(pdfUrl)}.
    \item \textit{Diagramme 18 :} Remplacer \textit{Transaction Commit} par \textit{Validation de la transaction} ou déplacer ce détail dans le texte technique.
    \item \textit{Diagramme 18 :} Vérifier si la génération du certificat PDF est faite par le backend Java ou par un service dédié ; si elle est faite côté backend, la ligne actuelle est acceptable.
    \item \textit{Diagramme 18 :} Ajouter une notification vers l’employé après l’approbation de la certification, car l’utilisateur doit être informé que le certificat est disponible.
    \item \textit{Diagramme 18 :} Ajouter un fragment \texttt{alt} pour la génération du certificat afin de distinguer le cas \textit{PDF généré avec succès} du cas \textit{échec de génération}.
    \item \textit{Diagramme 18 :} Vérifier si le téléchargement du certificat doit passer par le filtre de sécurité ; si oui, la validation JWT doit être représentée avant l’appel \textit{download()}.
    \item \textit{Diagramme 18 :} Ajouter une vérification d’autorisation lors du téléchargement pour s’assurer que l’employé télécharge uniquement son propre certificat.
    \item \textit{Diagramme 18 :} Vérifier que toutes les réponses sont représentées par des flèches de retour en pointillés et que les appels principaux restent en lignes pleines.
    \item \textit{Diagramme 18 :} Ajouter des barres d’activation cohérentes pour toutes les lignes de vie, notamment \textit{Base de données PostgreSQL}, \textit{Microservice IA} et \textit{Administrateur RH}.
    \item \textit{Diagramme 18 :} Réduire certains messages longs afin de rendre la figure plus lisible dans une page de rapport.

\end{itemize}
\color{black}
\subsubsection*{10. Diagramme de Classes}
% Ajustez le niveau de sectionnement selon la structure de votre rapport

\begin{figure}[H]
    \centering
    % Remplacez 'chemin_diagramme_classes_sprint3.png' par le nom exact de votre fichier
    \includegraphics[width=\textwidth]{images/classSprint3.png}
    \vspace{2mm}
    \caption{Diagramme de classes du Sprint 3}
    \label{fig:class_diagram_sprint3}
\end{figure}


\color{red}
\begin{itemize}
    \item \textit{Diagramme 19 :} Renommer le titre en \textit{Diagramme de classes du Sprint 3} ou \textit{Diagramme de classes de l’architecture du Sprint 3}, car l’intitulé actuel \textit{Architectures Complet} est grammaticalement incorrect.
    \item \textit{Diagramme 19 :} Corriger \textit{Couche Repositories} en \textit{Couche des repositories} ou \textit{Couche de persistance}, afin d’utiliser une formulation plus naturelle.
    \item \textit{Diagramme 19 :} Vérifier que le diagramme est bien présenté comme un diagramme de classes d’architecture en couches, et non comme un diagramme de classes métier classique.
    \item \textit{Diagramme 19 :} Supprimer les annotations techniques comme \textit{@Autowired}, car elles relèvent du code d’implémentation et alourdissent le diagramme.
    \item \textit{Diagramme 19 :} Uniformiser la langue du diagramme, car les noms de couches sont en français alors que la majorité des classes et méthodes sont en anglais.
    \item \textit{Diagramme 19 :} Vérifier que chaque repository porte un nom explicite et lisible, par exemple \textit{RoadmapRepository}, \textit{EvaluationRepository} et \textit{TrainingRequestRepository}.
    \item \textit{Diagramme 19 :} Ajouter les noms des relations entre contrôleurs et services uniquement si elles apportent une vraie valeur ; sinon, garder de simples dépendances UML.
    \item \textit{Diagramme 19 :} Remplacer les verbes libres comme \textit{utilise}, \textit{gère}, \textit{génère}, \textit{soumet} et \textit{possède} par des associations UML plus claires, avec cardinalités si nécessaire.
    \item \textit{Diagramme 19 :} Ajouter les multiplicités sur les relations entre les entités, car elles ne sont pas clairement visibles dans la partie ORM.
    \item \textit{Diagramme 19 :} Vérifier la relation entre \textit{User} et \textit{Roadmap} ; si chaque utilisateur possède une seule roadmap active, indiquer clairement une cardinalité \textit{1..1}, sinon utiliser \textit{1..*}.
    \item \textit{Diagramme 19 :} Vérifier la relation entre \textit{Roadmap} et \textit{CourseTask} ; si une roadmap est composée de plusieurs tâches, utiliser une composition avec cardinalité \textit{1..*}.
    \item \textit{Diagramme 19 :} Vérifier la relation entre \textit{Evaluation} et \textit{Certificate} ; si un certificat n’est créé qu’après réussite, préciser une relation optionnelle de type \textit{0..1}.
    \item \textit{Diagramme 19 :} Vérifier la relation entre \textit{User} et \textit{TrainingRequest} ; si un utilisateur peut soumettre plusieurs demandes, indiquer explicitement une cardinalité \textit{1..*}.
    \item \textit{Diagramme 19 :} Ajouter une relation claire entre \textit{Evaluation} et \textit{User}, car l’évaluation dépend d’un utilisateur, mais la cardinalité doit être explicitée.
    \item \textit{Diagramme 19 :} Enrichir la classe \textit{User} avec les attributs utiles au Sprint 3 si nécessaire, par exemple \textit{email}, \textit{role}, \textit{firstName} et \textit{lastName}, car la classe actuelle est très minimale.
    \item \textit{Diagramme 19 :} Vérifier si l’attribut \textit{readinessScore} de \textit{User} doit rester dans cette classe ou être déplacé dans une entité dédiée si plusieurs scores sont historisés.
    \item \textit{Diagramme 19 :} Vérifier si \textit{gamificationPoints} doit être stocké dans \textit{User} ou dans une entité séparée liée à la gamification, afin de mieux séparer les responsabilités métier.
    \item \textit{Diagramme 19 :} Ajouter les énumérations utilisées par les entités, comme \textit{KanbanStatus}, \textit{CertStatus} et \textit{RequestStatus}, ou au moins les mentionner dans une note ou une légende.
    \item \textit{Diagramme 19 :} Vérifier que les méthodes des contrôleurs sont cohérentes avec les cas d’utilisation du Sprint 3, notamment pour la roadmap, la gouvernance RH et l’évaluation/certification.
    \item \textit{Diagramme 19 :} Vérifier que \textit{RoadmapController} expose aussi les opérations de consultation si l’utilisateur doit lire sa roadmap, pas seulement la générer et mettre à jour le Kanban.
    \item \textit{Diagramme 19 :} Vérifier que \textit{GovernanceController} couvre bien les trois scénarios du sprint : consultation des demandes, prise de décision et suggestion éventuelle d’une alternative.
    \item \textit{Diagramme 19 :} Vérifier que \textit{EvaluationController} couvre bien tout le flux métier attendu : démarrage du test, soumission, génération de certificat et téléchargement si ces fonctionnalités font partie du sprint.
    \item \textit{Diagramme 19 :} Ajouter une dépendance plus explicite entre \textit{GovernanceService} et \textit{PdfExportService} seulement si le certificat ou les documents RH sont réellement générés dans ce module.
    \item \textit{Diagramme 19 :} Vérifier si \textit{SseNotificationService} doit être relié uniquement à \textit{GovernanceService}, ou aussi à d’autres services si des notifications temps réel sont envoyées ailleurs dans le sprint.
    \item \textit{Diagramme 19 :} Vérifier si \textit{AiMicroserviceClient} doit contenir uniquement des appels liés à la roadmap et au test dynamique, ou s’il faut le scinder en plusieurs clients plus spécialisés.
    \item \textit{Diagramme 19 :} Réduire les longues flèches horizontales entre les couches, car elles diminuent la lisibilité globale du diagramme.
    \item \textit{Diagramme 19 :} Rapprocher visuellement chaque repository des entités qu’il manipule afin de rendre les correspondances plus immédiates.
    \item \textit{Diagramme 19 :} Ajouter une légende expliquant la signification des icônes \textit{C} et \textit{I}, ainsi que les différents types de relations utilisés.

\end{itemize}
\color{black}
Ce diagramme illustre l'implémentation concrète de l'architecture logicielle du \textit{Backend} (\texttt{Spring Boot}) pour le Sprint 3. Il met en évidence le respect strict du patron de conception en couches (\textit{N-Tier Architecture} / \textit{Clean Architecture}), garantissant la modularité, la testabilité et la maintenabilité du système :

\begin{itemize}
    \item \textbf{La Couche d'Exposition (Controllers) :} Les classes annotées avec \texttt{@RestController} (ex. : \texttt{GovernanceController}) constituent le point d'entrée de l'API REST. Dépourvues de toute logique métier, elles se limitent à l'interception des requêtes HTTP, au routage, et à l'encapsulation des données de sortie dans des objets standardisés \texttt{ResponseEntity<T>}.

    \item \textbf{La Couche d'Orchestration (Services) :} C'est le cœur névralgique de l'application (annoté avec \texttt{@Service}). Elle encapsule la logique métier complexe (calcul des scores, \textit{workflows} d'approbation). Le diagramme révèle également l'utilisation de composants hautement spécialisés qui délèguent des responsabilités spécifiques pour éviter de surcharger les services principaux :
    \begin{itemize}
        \item \textbf{\texttt{AiMicroserviceClient} :} Abstrait la communication réseau vers le microservice \texttt{Python} (\texttt{FastAPI}) via un \texttt{RestTemplate}.
        \item \textbf{\texttt{PdfExportService} :} Isole la complexité du moteur de rendu documentaire (\texttt{Thymeleaf} et \texttt{OpenHTMLtoPDF}).
        \item \textbf{\texttt{SseNotificationService} :} Encapsule la logique de communication asynchrone (\textit{Server-Sent Events}) pour "pousser" les alertes en temps réel au \textit{Frontend}.
    \end{itemize}
    \vspace{2mm}

    \item \textbf{L'Inversion de Contrôle (IoC) :} Les flèches descendantes annotées \texttt{@Autowired} démontrent l'application du principe d'injection de dépendances natif de \texttt{Spring}. Les couches supérieures dépendent d'abstractions (interfaces) plutôt que d'implémentations concrètes, assurant un couplage lâche (\textit{Loose Coupling}).

    \item \textbf{La Couche de Persistance (Repositories \& Entités) :} Basée sur \texttt{Spring Data JPA}, elle assure le \textit{mapping} objet-relationnel (ORM) avec la base \texttt{PostgreSQL}. Les flèches de cardinalité au niveau des entités montrent la rigueur du modèle de données, notamment la composition forte (\texttt{*-- cascade ALL}) entre \texttt{Roadmap} et \texttt{CourseTask}, garantissant que la suppression d'un parcours entraîne la purge automatique de ses tâches associées pour préserver l'intégrité référentielle.
\end{itemize}

\subsection{Codage }
Consacrée à la phase d'implémentation spécifique au Sprint 3, cette section met en lumière les aspects les plus complexes du développement. Les fondations infrastructurelles de la plateforme étant désormais opérationnelles, notre travail de codage s'est concentré sur des enjeux à haute valeur ajoutée : le pilotage des modèles d'IA générative, la gestion des échanges entre les différents microservices, ainsi que le rendu dynamique et sécurisé des documents de certification.

\paragraph*{1. Orchestration de l'Intelligence Artificielle (Architecture Distribuée)}
Pour répondre au besoin de génération de contenu dynamique (\textit{Roadmaps} et Mini-tests), nous avons mis en place un écosystème \texttt{Python} découplé du cœur transactionnel Java :

\begin{itemize}
    \item \textbf{\texttt{FastAPI} \& \texttt{Uvicorn} :} Sélectionné comme socle de notre API cognitive pour ses hautes performances basées sur l'\textit{Event Loop} ASGI (\texttt{Starlette}). Il expose les \textit{webhooks} de génération tout en générant nativement les contrats d'interface (OpenAPI/Swagger).
    \item \textbf{\texttt{Ollama Python SDK } (\texttt{AsyncIO}) :}Librairie utilisée pour instancier des clients non-bloquants interagissant avec le modèle Llama 3 hébergé localement. Elle orchestre les échanges asynchrones avec le moteur d'IA en exploitant le mode de forçage JSON (JSON Schema). Cela garantit que les analyses de compétences et les prédictions produites par l'IA soient déterministes, structurées et directement parsables par les services du backend Java.
    \item \textbf{\texttt{Pydantic} :} Déployée de manière systématique pour imposer un typage fort en Python. Elle assure la validation sémantique et syntaxique des \textit{Data Models} à la volée, agissant comme un bouclier contre les erreurs de désérialisation avant l'envoi des données vers le \textit{Backend} Java.
    \item \textbf{\texttt{Spring RestTemplate} :} Côté Java, implémentation d'un client HTTP \textit{thread-safe} configuré avec des stratégies de \textit{Timeout} agressives et des gestionnaires d'exceptions (\texttt{ResponseErrorHandler}) pour assurer une résilience réseau (prévention de la famine des \textit{threads}) lors des appels au microservice IA.
\end{itemize}

\paragraph*{2. Workflow de Gouvernance et Rendu Documentaire (Full-Stack)}
Pour transformer la plateforme en un outil certifiant et interactif, nous avons implémenté des mécanismes avancés de manipulation du DOM et de génération à la volée :

\begin{itemize}
    \item \textbf{\texttt{OpenHTMLToPDF} \& \texttt{Thymeleaf} :} Ce duo constitue notre moteur de \textit{Server-Side Rendering} (SSR) documentaire. \texttt{Thymeleaf} résout dynamiquement l'arbre DOM en y injectant le contexte métier (identifiants cryptographiques, \textit{Readiness Score}). \texttt{OpenHTMLtoPDF} compile ensuite ce DOM XHTML en un flux d'octets PDF vectoriel, inviolable et prêt au téléchargement.
    \item \textbf{\texttt{Spring SseEmitter} \& \texttt{RxJS} :} Implémentation du \textit{pattern Publisher/Subscriber}. Le \textit{backend} Java maintient une connexion HTTP unidirectionnelle (\textit{Server-Sent Events}) pour publier les décisions de l'Admin. Côté Angular, \texttt{RxJS} écoute ce flux via l'objet \texttt{EventSource}. Les flux sont gérés de manière réactive avec un nettoyage strict des souscriptions (\textit{unsubscribes}) pour éviter toute fuite de mémoire (\textit{Memory Leak}).
    \item \textbf{\texttt{Angular CDK} (\textit{Drag and Drop}) :} Intégration des composants de bas niveau de Google pour manipuler l'arbre DOM virtuel. Cela permet au candidat de gérer les statuts de sa formation sur un tableau \textit{Kanban} interactif garantissant une fluidité à 60 FPS.
    \item \textbf{\texttt{Ngx-Toastr} :} Couplée aux \textit{Observables} \texttt{RxJS}, cette librairie consomme les flux asynchrones SSE pour afficher des alertes visuelles élégantes et éphémères (ex. "Certification approuvée"), améliorant considérablement l'UX globale.
\end{itemize}

\paragraph*{3. Outillage DevOps et Qualité du Code (\textit{Clean Code})}
Afin de garantir la robustesse de ces nouvelles fonctionnalités critiques, nous avons maintenu un niveau d'exigence industriel sur l'outillage et les tests :

\begin{itemize}
    \item \textbf{\texttt{Postman} :} Utilisé intensivement pour la création de collections d'API. Il a permis de simuler (\textit{mocker}) les retours du modèle LLM et de tester les flux binaires (\textit{Headers} PDF) de manière isolée avant l'intégration du client Angular.
    \item \textbf{\texttt{Java Bean Validation} (\texttt{Hibernate Validator}) :} Mécanisme déployé systématiquement par programmation orientée aspect (AOP). Les annotations (\texttt{@Valid}, \texttt{@NotNull}) appliquées aux objets de transfert (DTOs) assurent l'assainissement (\textit{Sanitization}) des requêtes entrantes, bloquant les \textit{payloads} malformés dès l'entrée du contrôleur.
    \item \textbf{\texttt{GitLab CI/CD} \& Stratégie GitFlow :} L'historique du code source est protégé par une stratégie de branches stricte. Les \textit{Merge Requests} ciblent systématiquement la branche de pré-production (\texttt{main}), déclenchant les pipelines d'intégration continue pour auditer le code et prévenir toute régression architecturale.
\end{itemize}










\subsection{Réalisation}
Cette partie détaille les résultats techniques du sprint, notamment les interfaces utilisateur développées avec \textbf{Angular}. L’objectif était d’offrir une expérience immersive, gamifiée et intuitive, tout en reflétant visuellement la complexité des algorithmes d’IA et des workflows de validation en arrière-plan.
\Needspace{20\baselineskip}
\subsection*{1. Tableau de Bord Global pour l'Upskilling}

Le Tableau de Bord Global pour l’Upskilling débute par l’identification des manques de compétences. Il mesure l’employabilité d’un collaborateur grâce à un indicateur \textit{« Prêt pour le poste »} (\textit{Readiness Score}) qqui permet d’analyser précisément les écarts de compétences techniques et comportementales, mettant ainsi en évidence le besoin de formation.
L’interface montre aussi en temps réel le niveau de l’utilisateur (\textbf{LVL}) et ses points d'expérience (\textbf{XP}),  intégrant ainsi une dimension gamifiée pour inciter l’apprenant à s’y engager davantage.
\begin{figure}[H]
    \centering
    \includegraphics[width=\textwidth]{images/dashboredFormation.png} 
    \caption{Tableau de Bord Global pour l'Upskilling}
    \label{fig:dashboard_upskilling}
\end{figure}
\Needspace{20\baselineskip}
\subsection*{2. La Sélection (Catalogue de Recommandations Intelligentes)}

Afin de mettre en œuvre la \textit{Roadmap}, le système suggère une sélection de cours adaptés proposés par plusieurs prestataires extérieurs (\textbf{Coursera}, \textbf{Udemy}, \textbf{LinkedIn Learning}). 

L’innovation principale de cette interface est son aspect prédictif : chaque carte de formation affiche l’impact estimé sur l’employabilité du candidat (ex: \textit{+5\% Readiness}).L’utilisateur se voit ainsi permettre de filtrer les cours par priorité et, en cliquant sur Ajouter au plan, de formuler une demande d’inscription.
\begin{figure}[H]
    \centering
    \includegraphics[width=\textwidth]{images/recommandation.png}
    \caption{ La Sélection (Catalogue de Recommandations Intelligentes)}
    \label{fig:course_recommendations}
\end{figure}
\subsection*{3. La Stratégie (Roadmap générée par l'IA)}

À partir du diagnostic, le système propose une solution macroscopique. Page d’apparence Parcours Stratégique transcrit les résultats du microservice IA dans une feuille de route actionnable.

Cette dernière propose des phases d’apprentissage précises  (ex: \textit{Phase 1 — Foundations \& Time Blocking}) sur mesure pour combler directement les faiblesses identifiées lors de l’évaluation initiale.
\begin{figure}[H]
    \centering
    \includegraphics[width=\textwidth]{images/roadmap.png}
    \caption{La Stratégie (Roadmap)}
    \label{fig:ai_roadmap}
\end{figure}
\subsection* {4. L'Exécution (Gestion Kanban des Formations)}
Lorsque les cours choisis dans le catalogue, le collaborateur passe à l’action. L’interface sous forme de tableau  \textbf{Kanban}à l’usager de piloter visuellement le cycle de vie de ses formations (de l'état \textit{Proposées} à \textit{Terminées}). 

C'est l'espace de travail quotidien de l'apprenant pour suivre, d'une manière fluide et interactive, sa progression, ayant une visibilité claire sur les tâches en cours et les objectifs atteints.
\begin{figure}[H]
    \centering
    \includegraphics[width=\textwidth]{images/mescourse.png}
    \caption{L'Exécution (Gestion Kanban des Formations)}
    \label{fig:kanban_execution}
\end{figure}
\subsection*{4. La Validation (Interface du Mini-Test IA)}

Une fois la formation complétée à 100\%, ll'utilisateur peut passer un  \textbf{Mini-Test de certification} généré dynamiquement par l'IA,  Ce test apparaît directement sur le tableau \textbf{Kanban} dans une fenêtre contextuelle. . Ce questionnaire à choix multiples de quatre options est créé en fonction du contenu pédagogique (\textit{ex: Angular Material}), Cela garantit que l'évaluation soit pertinente et unique.es réponses sont vérifiées de manière déterministe par  \textbf{Spring Boot} de manière déterministe pour maximiser la performance . l'obtention d'un score $\geq 70\%$ automatise le passage au statut \texttt{TERMINEE}, l'attribution des XP et la mise à jour du \textit{Readiness Score} .
\begin{figure}[H]
    \centering
    \includegraphics[width=\textwidth]{images/minitest.png}
    \caption{Interface du Mini-Test}
    \label{fig:minitest_interface}
\end{figure}
\subsection*{5. La Gouvernance (Espace Administration RH)}

L'exécution de la formation nécessite un encadrement managérial rigoureux. L'interface d'administration centralise les \textbf{KPI} et permet aux responsables RH de superviser les demandes issues du tableau Kanban des collaborateurs. 

LLes administrateurs peuvent valider ou rejeter les inscriptions. Ils vérifient les certificats soumis et ajoutent des notes, ce que l'on appelle des (\textit{adminNote}), Cela permet d'assurer que les formations individuelles sont alignées sur les objectifs globaux de l'entreprise.

\begin{figure}[H]
    \centering
    \includegraphics[width=\textwidth]{images/adminFormation.png}
    \caption{La Gouvernance (Espace Administration RH)}
    \label{fig:rh_governance}
\end{figure}

\subsection*{6. La Rétention (Leaderboard et Gamification)}

Le parcours se termine par la reconnaissance des efforts fournis. Une fois que l'administration a validé les formations, le collaborateur obtient des badges spéciaux, comme  \textit{Premier Certif} / \textit{Compétences comportementales Pro} et accumule de l' \textbf{XP} .

L'interface de classement (\textit{Leaderboard}) stimule l'engagement continu et la saine compétition en permettant aux apprenants de se situer par rapport à leurs pairs, renforçant ainsi la culture de l'apprentissage au sein de l'organisation.
D'un point de vue technique, le calcul des points XP est effectué côté backend à chaque transition d'état dans le Kanban (\texttt{TO\_DO} $\rightarrow$ \texttt{IN\_PROGRESS} $\rightarrow$ \texttt{DONE}), garantissant ainsi l'intégrité du classement. Le Leaderboard est mis à jour en temps réel via le même canal SSE utilisé pour les notifications de gouvernance.
\begin{figure}[H]
    \centering
    \includegraphics[width=\textwidth]{images/classement.png}
    \caption{ La Rétention (Leaderboard et Gamification)}
    \label{fig:gamification_leaderboard}
\end{figure}

\subsection{Phase de Test et Déploiement }

Cette section décrit la méthodologie appliquée pour valider la robustesse de la solution et garantir que l'intégration des nouvelles bibliothèques complexes (moteurs PDF, connecteurs IA, flux asynchrones) ne compromet pas la stabilité du cœur transactionnel de l'application.

\subsubsection*{1. Stratégie d'Assurance Qualité (QA)}

Afin de garantir une couverture optimale, la validation du code a été découpée en deux niveaux de tests automatisés :

\begin{itemize}
    \item \textbf{Tests Unitaires (\texttt{JUnit 5} \& \texttt{Mockito}) :} L'objectif principal a été d'isoler la logique métier sans dépendre de la base de données. Des tests ont été rédigés pour le \texttt{PdfExportService} en simulant (\textit{mocking}) le contexte pour vérifier que la conversion HTML-vers-PDF ne lève pas d'exception. De même, l'algorithme de calcul du \textit{Readiness Score} a été testé avec différents jeux de réponses pour s'assurer de sa fiabilité mathématique.
    \item \textbf{Tests d'Intégration d'API (\texttt{Postman}) :} L'effort s'est concentré sur l'architecture distribuée. Des collections \texttt{Postman} ont été créées pour vérifier le dialogue HTTP entre le \textit{Backend} Java et le microservice \texttt{Python} (\texttt{FastAPI}). Nous avons volontairement simulé des latences réseau pour valider que le \texttt{RestTemplate} Java appliquait correctement ses \textit{Timeouts} sans saturer la mémoire du serveur.
\end{itemize}

\subsubsection*{2. Plan de Validation Fonctionnelle (Tests \textit{End-to-End})}

Pour valider l'expérience utilisateur, un plan de tests de bout en bout (E2E) a été exécuté depuis l'interface \texttt{Angular}, ciblant les trois cas d'utilisation majeurs du sprint :

\begin{table}[H]
    \centering
    \renewcommand{\arraystretch}{1.5}
    \small
    \begin{tabularx}{\textwidth}{|c|>{\raggedright\arraybackslash}p{3cm}|>{\raggedright\arraybackslash}X|>{\raggedright\arraybackslash}X|}
        \hline
        \rowcolor[HTML]{F2F2F2}
        \textbf{ID} & \textbf{Scénario Testé} & \textbf{Pré-conditions \& Étapes} & \textbf{Résultat Attendu} \\
        \hline
        TEST-01 & Génération de la \textit{Roadmap} par l'IA & \textbf{Pré-cond :} Profil utilisateur connecté avec un \textit{Skill Gap} existant. \newline \textbf{Étapes :} 1. Accès espace \textit{Upskilling}. 2. Choix poste cible. 3. Clic sur "Générer". & Le \textit{Backend} interroge \texttt{FastAPI} avec succès. Le JSON est désérialisé sans erreur. Le \textit{Kanban} est alimenté dynamiquement. \\
        \hline
        TEST-02 & Soumission du test et calcul du Score & \textbf{Pré-cond :} Test lancé et chronomètre actif. \newline \textbf{Étapes :} 1. Sélectionner les réponses. 2. Soumission. & Le \textit{Backend} corrige le test. Le \textit{Readiness Score} est actualisé. Une demande de certification \texttt{PENDING\_APPROVAL} est créée. \\
        \hline
        TEST-03 & Approbation RH et Notifications (SSE) & \textbf{Pré-cond :} Admin et utilisateur connectés simultanément. \newline \textbf{Étapes :} 1. Choix de la demande par l'Admin. 2. Clic sur "Approuver". & La transaction SQL passe à \texttt{APPROVED}. L'utilisateur reçoit une alerte visuelle (\texttt{Ngx-Toastr}) instantanément via \texttt{SSE}. \\
        \hline
        TEST-04 & Génération documentaire (Moteur PDF) & \textbf{Pré-cond :} Certification approuvée. \newline \textbf{Étapes :} 1. Clic sur "Télécharger le diplôme". & Pas d'erreur 500. Téléchargement d'un \texttt{.pdf} valide incluant dynamiquement le nom et le score via \texttt{Thymeleaf}. \\
        \hline
    \end{tabularx}
    \caption{Plan de validation fonctionnelle du Sprint 3}
    \label{tab:validation_sprint3}
\end{table}

\subsubsection*{3. Résultats obtenus et Corrections apportées (Débogage)}

Durant cette phase d'intégration, nous avons été confrontés à des erreurs techniques complexes :

\begin{itemize}
    \item \textbf{Erreur 422 (\textit{Unprocessable Content}) :} 
    \textit{Problème :} \texttt{FastAPI} rejetait le \textit{payload} JSON du Java.
    \textit{Correction :} Alignement des DTOs Java avec le modèle \texttt{Pydantic} (\texttt{CareerPredictionBody}) côté \texttt{Python}.
    \item \textbf{Erreur 500 (Conflit de dépendances) :} 
    \textit{Problème :} Conflit entre \texttt{PDFBox 3.x} et \texttt{OpenHTMLtoPDF}.
    \textit{Correction :} Ajustement des versions dans le \texttt{pom.xml} pour restaurer la compatibilité d'architecture.
    \item \textbf{Erreur 500 (Rendu XHTML) :} 
    \textit{Problème :} Erreur fatale lors du \textit{parsing} du certificat.
    \textit{Correction :} Réécriture du \textit{template} en XHTML strict (balises fermées) et correction du CSS pour le moteur de rendu.
\end{itemize}

\subsubsection*{3. Processus de Déploiement et Exécution}

L'intégration du microservice IA a nécessité une évolution de l'infrastructure via la conteneurisation :

\begin{itemize}
    \item \textbf{Conteneurisation (\texttt{Docker}) :} Création de \texttt{Dockerfiles} isolés pour le \textit{Backend} (JDK 21), le \textit{Frontend} (\texttt{Nginx}) et le microservice IA (\texttt{Python 3.11}).
    \item \textbf{Orchestration Locale (\texttt{Docker Compose}) :} Le fichier \texttt{docker-compose.yml} gère le démarrage ordonné : \texttt{PostgreSQL} $\rightarrow$ \texttt{FastAPI} $\rightarrow$ \texttt{Spring Boot} $\rightarrow$ \texttt{Angular}.
    \item \textbf{Sécurisation des Secrets :} LLes variables sensibles (Clé API inter-services, secrets \texttt{JWT}, identifiants \texttt{PostgreSQL}) sont injectées au \textit{runtime} via un fichier \texttt{.env} protégé. Cela permet de sécuriser la communication entre le \textit{backend} \texttt{Java} et le microservice IA sans exposer de clés privées dans le code source.
\end{itemize}
\section{Suivi et Pilotage du Sprint}
\subsection{ Scrum Board}
Afin de garantir la transparence et l’inspection continue, qui sont essentielles dans la méthodologie Scrum, nous avons utilisé deux outils visuels importants pour gérer cette itération. Le sprint a duré \textbf{3 semaines} (soit 15 jours ouvrés), avec une charge de travail estimée à \textbf{34 Story Points}.

Le suivi quotidien des tâches s'est matérialisé par un tableau de bord reflétant l'état d'avancement des \textit{User Stories}. Ce management visuel a permis à l'équipe de développement d'identifier rapidement les goulots d'étranglement et de fluidifier le passage des tâches de la phase de conception vers la phase de test.

\begin{figure}[H]
    \centering
    % Assurez-vous que l'image scrum_board.png est bien dans votre dossier d'images
    \includegraphics[width=0.9\textwidth]{images/scrum_bored.png}
    \caption{Scrum board du sprint 3}
    \label{fig:scrum_board}
\end{figure}

\subsection{ Le Burndown Chart}

Pour évaluer notre vélocité et s'assurer de l'atteinte de l'objectif dans le temps imparti, nous avons maintenu à jour le \textit{Burndown Chart} lors de nos \textit{Daily Stand-ups}. 

Le graphique montre la courbe de l’effort restant en  \textit{Story Points} au fil des 15 jours du sprint. Nous avons rencontré quelques difficultés techniques entre\textbf{Jour 9 et le Jour 11} s'expliquent par les défis techniques rencontrés lors de la configuration du réseau interne Docker pour le microservice IA (FastAPI). La résolution de ces obstacles a permis une accélération de la validation finale, clôturant le sprint dans les délais.

\begin{figure}[H]
    \centering
    \begin{tikzpicture}
        \begin{axis}[
            width=0.9\textwidth,
            height=8cm,
            title={\textbf{Burndown Chart - Sprint 3 (Upskilling \& Formation)}},
            xlabel={Jours ouvrés du Sprint (3 Semaines)},
            ylabel={Story Points restants},
            xmin=0, xmax=15,
            ymin=0, ymax=40,
            xtick={0,1,2,3,4,5,6,7,8,9,10,11,12,13,14,15},
            ytick={0,10,20,30,34,40},
            legend pos=north east,
            ymajorgrids=true,
            xmajorgrids=true,
            grid style=dashed,
            axis background/.style={fill=gray!5},
            thick
        ]

        % La ligne idéale (l'objectif linéaire)
        \addplot[
            color=gray,
            dashed,
            very thick,
            mark=none
            ]
            coordinates {
            (0,34)(15,0)
            };
            \addlegendentry{Tendance Idéale}

        % La ligne réelle (l'avancement de l'équipe avec les blocages)
        \addplot[
            color=blue!70!black,
            very thick,
            mark=square*,
            mark options={fill=blue!70!black, scale=0.9}
            ]
            coordinates {
            (0,34)     % Jour 0 : Lancement
            (1,34)     % Jour 1 : Analyse
            (2,32)     % Jour 2 : Premières validations
            (3,29)     % Jour 3
            (4,29)     % Jour 4 : Focus sur le code long
            (5,26)     % Jour 5 : Fin de semaine 1
            (6,23)     % Jour 6 : Début semaine 2
            (7,21)     % Jour 7
            (8,21)     % Jour 8
            (9,18)     % Jour 9 : Début des problèmes Docker/IA
            (10,18)    % Jour 10 : Recherche de solution
            (11,18)    % Jour 11 : Blocage résolu
            (12,10)    % Jour 12 : Validation massive des tickets
            (13,6)     % Jour 13 : Gamification terminée
            (14,3)     % Jour 14 : Tests finaux
            (15,0)     % Jour 15 : Démo et clôture
            };
            \addlegendentry{Avancement Réel}

        \end{axis}
    \end{tikzpicture}
    \caption{Burndown Chart du sprint 3}
    \label{fig:burndown_chart}
\end{figure}

\subsection{Indicateurs de Performance (KPIs) du Sprint 3}

Afin d'évaluer la réussite de cette itération focalisée sur l'intelligence RH et l'\textit{upskilling}, nous avons mis en place un système de suivi basé sur des indicateurs précis.

\paragraph*{1. Indicateurs de Pilotage (Agilité)}
Le pilotage de ce sprint a été marqué par une forte densité technique (intégration de l'IA et des \texttt{WebSockets}).

\begin{itemize}
    \item \textbf{Vélocité Réalisée :} Sur les 34 \textit{Story Points} planifiés, la totalité a été livrée dans les délais, démontrant une maturité croissante de l'équipe sur la \textit{stack} technologique.
    \item \textbf{Densité de Développement :} La livraison de 12 nouveaux \textit{endpoints} \texttt{API} complexes souligne l'efficacité de la collaboration entre les modules \textit{Backend}, \textit{Frontend} et IA.
    \item \textbf{Qualité des \textit{User Stories} :} Le taux de validation en fin de sprint par le \textit{Product Owner} a atteint 100\%, sans aucun report de tâche (\textit{Carry-over}).
\end{itemize}

\paragraph*{2. Indicateurs Techniques et Qualité}
Le tableau ci-dessous synthétise les performances de la plateforme à l'issue du Sprint 3, couvrant les aspects logiciels, l'intelligence artificielle et l'infrastructure \texttt{Docker}.

\begin{table}[H]
    \centering
    \renewcommand{\arraystretch}{1.5}
    \small
    \begin{tabularx}{0.95\textwidth}{|X|c|c|c|}
        \hline
        \rowcolor[HTML]{F2F2F2}
        \textbf{Indicateur} & \textbf{Cible} & \textbf{Résultat Obtenu} & \textbf{État} \\
        \hline
        \textit{Story Points} complétés & 34 & 34 & $\checkmark$ Validé \\
        \hline
        Disponibilité (\textit{Uptime} \texttt{Docker}) & $100\%$ & $100\%$ & $\checkmark$ Validé \\
        \hline
        Poids des images \texttt{Docker} & $< 600$ Mo & $\sim 410$ Mo & $\checkmark$ Validé \\
        \hline
        Précision du \textit{Gap Analysis} (IA) & $> 90\%$ & $94\%$ & $\checkmark$ Validé \\
        \hline
        Fiabilité du \textit{Matching} Formation & $> 85\%$ & $89\%$ & $\checkmark$ Validé \\
        \hline
        Taux de réussite des tests E2E & $> 90\%$ & $100\%$ & $\checkmark$ Validé \\
        \hline
        Temps de génération \textit{Roadmap} IA & $< 5$ s & $\sim 3.2$ s & $\checkmark$ Validé \\
        \hline
        Temps de génération PDF (Passeport) & $< 3$ s & $\sim 1.8$ s & $\checkmark$ Validé \\
        \hline
        Latence \texttt{SSE} (Notifications temps réel) & $< 500$ ms & $\sim 180$ ms & $\checkmark$ Validé \\
        \hline
    \end{tabularx}
    \caption{Indicateurs de Performance du Sprint 3}
    \label{tab:kpi_tech_s3}
\end{table}

\paragraph*{3. Analyse du Bilan}
Le bilan du Sprint 3 est très positif, tant sur le plan fonctionnel que technique. L'intégration réussie de l'analyse des écarts (\textit{Skill-Gap Analysis}) sur les deux dimensions (\textit{compétences techniques et comportementales}) constitue la plus grande valeur ajoutée de cette itération.

La stabilité des infrastructures \texttt{Docker}, combinée à une latence de notification très faible (180 ms), prouve que la plateforme est capable de supporter des interactions en temps réel complexes. Enfin, la rapidité de génération des documents officiels (\texttt{PDF}) et des \textit{roadmaps} d'apprentissage confirme que l'architecture est optimisée pour une utilisation intensive par les départements RH.


\section{Conclusion }

Dans ce chapitre, nous avons présenté la conception, la réalisation et la validation du module  de\textit{Formation} et d'\textit{Upskilling},qui est au cœur de la plateforme \textbf{TalentPredict}. Nous avons utilisé une modélisation UML rigoureuse pour structurer les flux de données complexes entre les collaborateurs, l’administration RH et l’Intelligence Artificielle.
Le passage de la théorie à la pratique s'est matérialisé par le développement d'interfaces \textbf{Angular} dynamiques et ergonomiques, soutenues par une architecture \textit{backend} hybride (\textbf{Spring Boot} et \textbf{FastAPI}). Nous avons réussi à concrétiser la promesse initiale du projet : transformer les lacunes de compétences (\textit{Skill Gaps}) en opportunités d'apprentissage personnalisées grâce aux prédictions génératives du LLM.

Par ailleurs, l'intégration de mécanismes de gamification (niveaux, badges) et d'un pilotage visuel (\textbf{Kanban}) favorise l'engagement de l'apprenant, tandis que le \textit{workflow} d'approbation et les notifications en temps réel garantissent une gouvernance RH maîtrisée. Enfin, l'application stricte de la méthodologie Agile \textbf{Scrum} et l'exécution des tests fonctionnels ont prouvé la fiabilité et la robustesse de ce système.

Avec la validation de ce troisième sprint, le cycle de vie de la gestion des talents est désormais opérationnel de bout en bout. Le chapitre suivant s'attachera à présenter le quatrième sprint.
\chapter{ Sprint 4 :Pilotage, Management \& Mise en Production}
\section*{Introduction}
\phantomsection
\addcontentsline{toc}{section}{Introduction}
\phantomsection 
Ce chapitre est la dernière étape du projet \textbf{TalentPredict}. Nous avons développé les moteurs d'évaluation et l'intelligence artificielle lors des étapes précédentes. Maintenant, nous nous concentrons sur la façon dont nous allons utiliser les données pour aider les ressources humaines. Nous voulons créer des outils pour les aider à prendre des décisions et à communiquer facilement. Nous allons également consolider tous les modules que nous avons créés, générer le \texttt{Talent Passport} et mettre la plateforme en ligne grâce à \textbf{Docker} et l'implémentation d'une chaîne
 \textbf{CI/CD} (Intégration et Déploiement Continus). Cette étape finale assure une qualité de code constante, une automatisation des tests et une mise en production sécurisée, livrant ainsi une plateforme solide, adaptable et prête à l’emploi.
\section{Objectifs du Sprint 4}
% Remarque : Utilisez \section{} au lieu de \subsection{} si "4.1" correspond au premier niveau de titre dans votre chapitre.

L’objectif principal de ce quatrième et dernier sprint est de renforcer les différents modules développés précédemment pour obtenir une version finale et stable de la plateforme \textbf{TalentPredict} . Ce sprint se concentre sur deux axes : la finition des interfaces de pilotage stratégique et l'industrialisation de la plateforme. 

    \subsubsection*{2. Objectifs fonctionnels et interactifs :}
    \begin{itemize}
        \item \textbf{Conception d'un tableau de bord décisionnel (\textit{Executive Dashboard}) :} Fournir aux administrateurs RH une vue centralisée et en temps réel des indicateurs clés de performance (KPIs) stratégiques pour le pilotage des talents.
        \item \textbf{Développement de l'espace collaborateur (\textit{User Dashboard}) :} Créer une interface personnalisée et ergonomique, en adoptant le design \textit{Bento Grid} . Cela inclut un radar de compétences dynamique et la présentation du persona IA.
        \item \textbf{Mise en place de la gestion administrative (Module RH) :} Développer un module complet (CRUD) pour gérer les comptes employés. Cela permettra de consulter leurs profils publics (liens sociaux, résultats détaillés, informations biographiques) tout en sécurisant les accès grâce à un contrôle strict basé sur les rôles (RBAC).
        \item \textbf{Création du \textit{Communication Hub} :} Centraliser les échanges avec un système de notifications en temps réel SSE(\textit{Server-Sent Events}) et permettre aux employés d’interagir avec les recommandations de l’IA via une boucle de rétroaction.
        \item \textbf{Génération de rapports professionnels :} 
        Produire, côté serveur, des documents structurés et fidèles. Cela inclut le \textit{Talent Passport} pour le bilan de compétences individuel des employés et le \textit{Rapport Global RH }  pour la synthèse analytique et stratégique de l'ensemble du capital humain..
    \end{itemize}

    \subsubsection*{2. Objectifs techniques et opérationnels (DevOps) :}
    \begin{itemize}
        \item \textbf{Conteneurisation de l'architecture :} Mettre en place un déploiement complet et isolé des différents services (Frontend, Backend, IA, Base de données) via \texttt{Docker} et \texttt{Docker Compose}.
        \item \textbf{Automatisation CI/CD :} Configurer un pipeline d'intégration et de déploiement continus sous \texttt{GitLab CI} pour garantir que les tests automatisés soient exécutés à chaque évolution du code.
        \item \textbf{Stabilisation et Performance :} Optimiser les temps de réponse des services asynchrones et assurer une cohérence totale des données entre les modules d'évaluation, de formation et de dashboarding.
    \end{itemize}

\subsubsection*{Critères d'acceptation (\textit{Definition of Done})}
Pour considérer ce sprint comme achevé, l'équipe a défini les critères de validation suivants :
\begin{itemize}
    \item Tous les tests unitaires du backend doivent s'exécuter avec succès.
    \item La génération du génération de rapports professionnels doit s'effectuer correctement, en respectant la mise en page spécifiée pour chaque utilisateur.
    \item Les notifications asynchrones doivent être réceptionnées en temps réel sans perte de connexion.
    \item  Le système doit montrer une stabilité satisfaisante dans l’environnement de pré-production  (\textit{staging}).
\end{itemize}
\section{Séance de Sprint Planning}

La séance de Sprint Planning a réuni l'équipe de développement, composée de \textit{Mohamed Ben Abdeladhim} et \textit{Rahma Barouni}, ainsi que le Product Owner,Mr \textit{Med Amine Gharbi}, et notre encadrante pédagogique, Mme \textit{Ben Aïcha Takwa}, afin de définir le périmètre de cette dernière itération de développement. Les points clés suivants ont été déterminés lors de cette réunion.

\begin{itemize}
    \item \textbf{Durée :} La durée de cette itération a été fixée à un mois, ce qui nous permet de disposer du temps nécessaire pour achever les développements complexes, réaliser les tests finaux et assurer le déploiement.
    \item \textbf{Priorités absolues :} Le développement des tableaux de bord pour les ressources humaines et les utilisateurs, le centre de communication, la génération automatisée des rapports professionnels et la mise en place du module de gestion des employés permettent aux ressources humaines d'avoir un contrôle total.
    \item \textbf{Priorités secondaires :} Optimisation du moteur d’IA pour améliorer ses performances.
 et la configuration du pipeline CI/CD via \texttt{Docker},et \texttt{GitLab}.
    \item \textbf{Prérequis :} Un prérequis essentiel pour cette itération est la stabilité confirmée des API livrées lors des sprints précédents, notamment le moteur d’IA, le module d’évaluation et le module de recommendations.
\end{itemize}

À la fin de cette réunion, les besoins ont été décomposés en \textit{User Stories} (US) et estimés en jours-homme, en donnant la priorité maximale aux livrables critiques pour la démonstration finale de \textbf{TalentPredict}.

\section{Backlog du Sprint}
% Utilisez \subsubsection{} selon la hiérarchie de vos titres.
\Needspace{18\baselineskip}

Pour mettre en œuvre les priorités définies lors de la planification du sprint et faciliter la répartition des tâches entre les membres de l'équipe (Mohamed et Rahma), chaque \textit{User Story} a été décomposée en tâches techniques spécifiques. Le tableau présente ce backlog enrichi, en incluant les responsabilités et les critères de validation associés.

\vspace{-0.3em}
\begin{table}[H]
    \centering
    \renewcommand{\arraystretch}{1}
    \scriptsize
    \begin{tabularx}{\textwidth}{|c|>{\raggedright\arraybackslash}p{3cm}|>{\raggedright\arraybackslash}X|c|>{\raggedright\arraybackslash}X|c|}
        \hline
        \textbf{ID} & \textbf{Épic} & \textbf{User Story} & \textbf{Priorité} & \textbf{Critères d'acceptation} & \textbf{Estim. (SP)} \\
        \hline
        
        % US 4.1
        US4.1 & Système de Reporting 
        & En tant qu'employé, je veux exporter mon Passeport Talent en PDF. 
        & Moyenne 
        & \textbullet{} PDF fidèle au design, incluant scores IA et radar. 
        & 8 \\
        \hline
        
        % US 4.2
        US4.2 & \textit{Executive Analytics} 
        & En tant que RH, je veux visualiser les KPIs stratégiques sur un \textit{Dashboard} Admin. 
        & Haute 
        & \textbullet{} \textit{Dashboard Bento} fonctionnel avec graphiques \texttt{Chart.js}. 
        & 8 \\
        \hline
        
        % US 4.3
        US4.3 & Expérience Utilisateur 
        & En tant qu'employé, je veux voir mon radar de compétences et mon \textit{persona} IA. 
        & Haute
        & \textbullet{} Visualisation dynamique des compétences et résultats PCM. 
        & 5 \\
        \hline
        
        % US 4.4
        US4.4 & \textit{Communication Hub} 
        & En tant qu'utilisateur, je veux recevoir des notifications système en temps réel. 
        & Haute 
        & \textbullet{} Réception via WebSockets sans rafraîchissement de page. 
        & 5 \\
        \hline
        
        % US 4.5
        US4.5 & Pilotage Stratégique 
        & En tant qu'admin, je veux télécharger un rapport analytique global de l'entreprise. 
        & Moyenne 
        & \textbullet{} PDF consolidé avec statistiques par département. 
        & 5 \\
        \hline
        
        % US 4.6
        US4.6 & Gouvernance RH 
        & En tant qu'admin, je veux gérer (CRUD) et activer/désactiver les comptes employés. 
        & Haute 
        & \textbullet{} Changement de statut instantané et accès aux profils publics. 
        & 5 \\
        \hline
        
        % US 4.7
        US4.7 & Infrastructure \& DevOps 
        & En tant que Dev, je veux conteneuriser l'appli et automatiser le déploiement. 
        & Haute 
        & \textbullet{} Projet déployable via \texttt{Docker Compose} et pipeline \texttt{GitLab CI}. 
        & 13 \\
        \hline
        
    \end{tabularx}
    \vspace{1mm}
    \caption{Baklog du Sprint 4}
    \label{tab:backlog_sprint4}
\end{table}
\section{Démarrage du Sprint 4}
% Vous pouvez adapter le niveau de titre (\section, \subsection) selon la structure exacte de votre chapitre.
Une fois que l’équipe a validé le périmètre du sprint et a réparti les tâches entre ses membres, la phase de réalisation a réellement commencé. Cette phase a débuté par une analyse approfondie et détaillée des spécifications techniques. Ensuite, l’équipe a conçu l’architecture complète des nouveaux modules, ce qui est une étape cruciale. L’objectif principal de cette étape est de garantir que les solutions implémentées s’intègrent parfaitement avec le code produit lors des itérations précédentes. Les solutions doivent s’intégrer de manière harmonieuse avec le code existant pour éviter tout problème.

\subsubsection*{1. Activités réalisées}% La numérotation 3.4.1 se fera automatiquement selon la structure de votre document.

Durant cette itération, nous avons mené de front plusieurs activités critiques :

\begin{itemize}
    \item \textbf{Intégration du Moteur de Reporting :} Développement de la logique \textit{Backend} pour transformer les \textit{templates} HTML en documents PDF officiels.
    \item \textbf{Conception des Dashboards Bento :} Développement de l'interface réactive et intégration de la bibliothèque \texttt{Chart.js} pour la visualisation des compétences.
    \item \textbf{Configuration du Hub de Communication :}Implémentation du flux Server-Sent Events (SSE) pour les notifications en temps réel.
    \item \textbf{Travaux DevOps :} Écriture des fichiers de configuration Docker et mise en place du pipeline d'intégration continue sous \texttt{GitLab CI}.
    \item \textbf{Synchronisation des Données :} \textit{Mapping} des prédictions issues du service Python vers les objets Java (DTO) pour une restitution fluide.
\end{itemize}

\subsubsection*{2. Livrables produits}

À l'issue du sprint, les livrables suivants ont été finalisés :

\begin{itemize}
    \item \textbf{Modules Frontend :} Dashboards Admin/User, Centre de Notifications et Gestion des employés.
    \item \textbf{Services Backend :} Reporting Service,SSE Notification Hub et API de gestion RH.
    \item \textbf{Documents de Restitution :} \textit{Templates} PDF du Passeport Talent et du Rapport Global RH.
    \item \textbf{Artefacts DevOps :} Fichiers \texttt{Dockerfile} et \texttt{docker-compose.yml} opérationnels.
    \item \textbf{Suite de Tests :} Rapports de tests unitaires et d'intégration validés par la CI/CD.
\end{itemize}

\subsubsection{3. Difficultés rencontrées et traitement}

Le développement a présenté certains défis techniques que nous avons dû résoudre :

\begin{itemize}
    \item \textbf{Difficulté 1 : Rendu des polices dans le PDF}
    \begin{itemize}
        \item \textbf{Problème :} Les caractères spéciaux et les icônes ne s'affichaient pas correctement dans les rapports exportés.
        \item \textbf{Traitement :} Intégration de polices vectorielles TrueType (TTF) directement dans le moteur de rendu \texttt{OpenHTMLToPDF}.
    \end{itemize}
    \vspace{2mm} % Espacement léger pour aérer la liste
    
    \item \textbf{Difficulté 2 : Latence des notifications via Docker}
    \begin{itemize}
        \item \textbf{Problème :} Latence et stabilité du flux SSE via Docker.
        \item \textbf{Traitement :} Implémentation d'un mécanisme de 'Keep-Alive' (envoi périodique de commentaires vides) côté SseEmitter pour maintenir la connexion active.
    \end{itemize}
    \vspace{2mm}
    
    \item \textbf{Difficulté 3 : Cohérence des données IA}
    \begin{itemize}
        \item \textbf{Problème :} Décalage entre les types de données envoyés par \texttt{FastAPI} (Python) et ceux attendus par \texttt{Spring Boot} (Java).
        \item \textbf{Traitement :} Mise en place d'une couche de validation stricte via des DTOs communs pour garantir l'intégrité des analyses.
    \end{itemize}
\end{itemize}

\subsection{Analyse}


Cette phase a pour objectif de définir précisément le périmètre fonctionnel et technique du sprint, en identifiant les interactions entre les utilisateurs et les nouveaux modules de pilotage et de reporting.
\subsubsection*{1. Identification des Acteurs}
% J'utilise \subsubsection pour correspondre au niveau 4.4.2

Dans le cadre de ce sprint, les rôles se spécialisent autour de l'exploitation et de la valorisation des données :

\begin{itemize}
     
     \item \textbf{L'Administrateur RH (Acteur Principal) :} Responsable de la gestion stratégique, il utilise les outils décisionnels, gère le vivier de talents et valide les parcours de formation.
    
    \item \textbf{L'Employé (\textit{Acteur Principal}) :}tUtilisateur final qui consulte ses indicateurs de performance, reçoit des notifications et exporte son bilan de compétences officiel.

    \item \textbf{Le Système IA (Acteur Système) :} Service externe automatisé qui fournit les analyses prédictives nécessaires au fonctionnement des dashboards.
    \item \textbf{Le Service SSE (Acteur Système) :} Composant technique responsable du maintien de la connexion persistante et de la diffusion instantanée des notifications asynchrones vers les utilisateurs.
\end{itemize}

\subsubsection*{2. Clarification du besoin}

L'enjeu de ce sprint est de transformer la donnée brute en information décisionnelle. Le besoin se décline en trois axes :

\begin{itemize}
    \item \textbf{Visualisation :} Rendre les résultats de l'IA et des tests intelligibles via des graphiques dynamiques.
    \item \textbf{Interactivité :} Assurer une communication fluide et instantanée entre la direction et les employés.
    \item \textbf{Formalisation :} Produire des documents officiels (PDF) garantissant la traçabilité des compétences acquises.
\end{itemize}

\subsubsection*{3. Scénarios de validation}

Pour valider la réussite des développements, nous avons défini des scénarios clés :

\begin{itemize}
    \item \textbf{Scénario de Pilotage :} L'Admin RH consulte le \textit{dashboard} global ; le système doit afficher les KPIs agrégés après avoir sollicité le service IA.
    \item \textbf{Scénario de Communication :} L'Admin envoie un "\textit{Broadcast}" ; tous les employés connectés doivent recevoir l'alerte instantanément via SSE (Server-Sent Events).
    \item \textbf{Scénario d'Export :} L'Employé génère son "\textit{Talent Passport}" ; le système doit produire un PDF fidèle au design \textit{Bento}, incluant ses dernières analyses IA.
\end{itemize}

\subsubsection*{4. Contraintes techniques}

L'implémentation de ces besoins doit respecter plusieurs contraintes :

\begin{itemize}
    \item \textbf{Temps réel :} Utilisation du protocole \texttt{ SSE (Server-Sent Events)} pour garantir une latence minimale des notifications.
    \item \textbf{Performance :} La génération de rapports PDF complexes est gérée en dehors du cycle de vie principal (via des flux de sortie optimisés) pour ne pas bloquer les requêtes utilisateurs.
    \item \textbf{Sécurité (RBAC) :} Chaque action (export, consultation, gestion) doit être validée par le module de sécurité \texttt{JWT} selon le rôle de l'utilisateur.
\end{itemize}

\subsubsection{5. Diagramme de Cas d’Utilisation du Sprint 4}

La figure illustre les interactions fonctionnelles spécifiques à ce dernier sprint, mettant en lumière les nouveaux services de pilotage et de communication.

\begin{figure}[H]
    \centering
    % Remplacez 'edited-image.png' par le chemin exact si l'image est dans un sous-dossier (ex: images/edited-image.png)
    \includegraphics[width=0.85\textwidth]{images/usecasesprint4.png} 
    \vspace{2mm}
    \caption{Diagramme de cas d'utilisation globale du sprint 4}
    \label{fig:use_case_sprint4}
\end{figure}

\color{red}
\begin{itemize}
    \item \textit{Diagramme 20 :} Renommer le titre en \textit{Diagramme de cas d’utilisation global du Sprint 4 : pilotage et reporting} pour garder une formulation académique et homogène.
    \item \textit{Diagramme 20 :} Remplacer l’acteur \textit{Admin RH} par \textit{Administrateur RH}, afin d’utiliser la même terminologie dans tout le rapport.
    \item \textit{Diagramme 20 :} Remplacer \textit{Système IA (FastAPI)} par \textit{Service IA}, puis expliquer dans le texte que ce service est implémenté avec FastAPI.
    \item \textit{Diagramme 20 :} Remplacer \textit{Système SSE (Notifications)} par \textit{Service de notification temps réel}, pour éviter de mettre une technologie directement comme acteur métier.
    \item \textit{Diagramme 20 :} Supprimer les flèches sur les associations entre les acteurs et les cas d’utilisation ; en UML, une association acteur--cas d’utilisation est représentée par une ligne simple sans flèche.
    \item \textit{Diagramme 20 :} Remplacer \textit{Consulter le Profil Public d’un employé} par \textit{Consulter le profil public d’un employé}, avec une casse plus cohérente.
    \item \textit{Diagramme 20 :} Remplacer \textit{Gérer le vivier d’employés (CRUD / Statuts)} par \textit{Gérer le vivier d’employés}, puis détailler les opérations CRUD dans le texte ou dans un diagramme de classes.
    \item \textit{Diagramme 20 :} Remplacer \textit{Consulter Executive Dashboard (Analytics)} par \textit{Consulter le tableau de bord exécutif}, afin d’éviter le mélange français--anglais.
    \item \textit{Diagramme 20 :} Remplacer \textit{Envoyer des relances groupées (Broadcast)} par \textit{Envoyer des relances groupées}.
    \item \textit{Diagramme 20 :} Remplacer \textit{Consulter Dashboard Personnel} par \textit{Consulter le tableau de bord personnel}.
    \item \textit{Diagramme 20 :} Remplacer \textit{Générer Insights \& Simulations} par \textit{Générer des analyses et des simulations}, avec une formulation entièrement en français.
    \item \textit{Diagramme 20 :} Remplacer \textit{Exporter Rapport RH Global (PDF)} par \textit{Exporter le rapport RH global}.
    \item \textit{Diagramme 20 :} Remplacer \textit{Exporter Passeport Talent (PDF)} par \textit{Exporter le passeport de compétences}.
    \item \textit{Diagramme 20 :} Remplacer \textit{Recevoir des notifications} par \textit{Consulter les notifications} si l’action est réalisée par l’utilisateur dans l’interface.
    \item \textit{Diagramme 20 :} Vérifier les relations \texttt{<<include>>} vers \textit{S’authentifier} ; elles sont correctes si l’authentification est obligatoire, mais elles peuvent être supprimées pour alléger le diagramme global.
    \item \textit{Diagramme 20 :} Corriger la relation \texttt{<<extend>>} entre \textit{Exporter Rapport RH Global} et \textit{Consulter Executive Dashboard} ; elle est cohérente si l’export PDF est une option disponible depuis le tableau de bord exécutif.
    \item \textit{Diagramme 20 :} Corriger la relation \texttt{<<extend>>} entre \textit{Exporter Passeport Talent} et \textit{Consulter Dashboard Personnel} ; elle est cohérente si l’export du passeport est une action optionnelle depuis le tableau de bord personnel.
    \item \textit{Diagramme 20 :} Vérifier la relation \texttt{<<include>>} entre \textit{Consulter le tableau de bord exécutif} et \textit{Générer des analyses et des simulations} ; elle est correcte seulement si ces analyses sont toujours générées lors de la consultation.
    \item \textit{Diagramme 20 :} Si les analyses et simulations sont déclenchées seulement à la demande, remplacer la relation \texttt{<<include>>} par une relation \texttt{<<extend>>}.
    \item \textit{Diagramme 20 :} Relier le \textit{Service IA} uniquement au cas \textit{Générer des analyses et des simulations}, car c’est le seul cas qui dépend explicitement d’un traitement intelligent.
    \item \textit{Diagramme 20 :} Relier le \textit{Service de notification temps réel} uniquement aux cas liés aux notifications ou aux relances, afin d’éviter les associations trop générales.
    \item \textit{Diagramme 20 :} Vérifier si \textit{Approuver les formations} appartient réellement au Sprint 4 ; si cette fonctionnalité a déjà été traitée dans le Sprint 3, il faut l’enlever ou la justifier comme fonctionnalité de pilotage final.
    \item \textit{Diagramme 20 :} Ajouter un cas \textit{Consulter les indicateurs RH globaux} si le Sprint 4 porte sur le reporting et le pilotage stratégique.
    \item \textit{Diagramme 20 :} Ajouter un cas \textit{Visualiser la cartographie des compétences} si cette fonctionnalité est présente dans le tableau de bord RH.
    \item \textit{Diagramme 20 :} Ajouter un cas \textit{Télécharger le rapport RH} ou garder \textit{Exporter le rapport RH global}, mais éviter d’utiliser deux formulations différentes pour la même action.
    \item \textit{Diagramme 20 :} Vérifier si l’employé peut consulter un tableau de bord personnel dans le Sprint 4 ; si oui, ajouter éventuellement les sous-cas \textit{Consulter ses scores}, \textit{Consulter sa progression} et \textit{Télécharger son passeport}.
    \item \textit{Diagramme 20 :} Réorganiser les cas d’utilisation pour réduire les croisements de lignes, en plaçant les cas RH à gauche ou au centre, les cas employé en bas, les exports à droite et les services externes à l’extrémité droite.
    \item \textit{Diagramme 20 :} Ajouter une légende expliquant les associations simples, les relations \texttt{<<include>>} et les relations \texttt{<<extend>>}.
  
\end{itemize}
\color{black}

% --- MODULE 1 ---
\subsubsection*{Cas 1 : Pilotage Collaborateur (User Dashboard)}

Ce module offre à l'employé une interface de pilotage personnel. L'objectif est de centraliser les résultats des évaluations diagnostiques et de fournir des outils de développement de carrière basés sur l'IA, aboutissant à la génération d'un "Passeport Talent" certifié.

\begin{figure}[H]
    \centering
    \includegraphics[width=1\textwidth]{images/case1Sprint4.png}
    \caption{Diagramme de cas d'utilisation : User Dashboard \& Passeport}
    \label{fig:usecase_dash_user}
\end{figure}

\color{red}
\begin{itemize}
    \item Renommer le titre du module en français homogène, par exemple : \textit{Module : Tableau de bord utilisateur et passeport de compétences}.
    \item Remplacer le cas d’utilisation \textit{Consulter Dashboard Personnel} par \textit{Consulter le tableau de bord personnel}.
    \item Remplacer \textit{Authentication} par \textit{S'authentifier} ou \textit{Se connecter}, afin d’éviter le mélange français--anglais.
    \item Remplacer \textit{Voir Recommandations de Carrière} par \textit{Consulter les recommandations de carrière}.
    \item Remplacer \textit{Suivre ses Formations \& Progrès} par \textit{Suivre ses formations et sa progression}.
    \item Remplacer \textit{Consulter Profil Psychologique IA} par \textit{Consulter le profil comportemental généré par l'IA} ou \textit{Consulter le profil psychologique}.
    \item Remplacer \textit{Visualiser Radar de Compétences} par \textit{Visualiser le radar des compétences}.
    \item Remplacer \textit{Exporter Passeport Talent (Certification PDF)} par \textit{Exporter le passeport talent en PDF} ou \textit{Télécharger le passeport de compétences}.
    \item Renommer \textit{Agrégation \& Analyse IA} par un libellé plus métier, par exemple \textit{Générer l’analyse synthétique du profil} ou \textit{Calculer le passeport talent}.
    \item Renommer l’acteur secondaire \textit{«Service» Moteur IA} par \textit{Service IA} pour garder une notation plus simple et plus cohérente.
    \item Vérifier la relation \texttt{<<include>>} entre \textit{Consulter le tableau de bord personnel} et \textit{S'authentifier} ; elle est correcte seulement si l’authentification est obligatoire.
    \item Vérifier les relations \texttt{<<include>>} entre \textit{Consulter le tableau de bord personnel} et les cas \textit{Consulter l’historique des tests}, \textit{Suivre ses formations}, \textit{Consulter les recommandations}, \textit{Consulter le profil psychologique} et \textit{Visualiser le radar des compétences}.
    \item Si ces fonctionnalités ne sont pas toujours exécutées automatiquement lors de l’ouverture du tableau de bord, remplacer les \texttt{<<include>>} par de simples associations ou par des cas séparés.
    \item Vérifier la relation \texttt{<<extend>>} entre \textit{Exporter le passeport talent en PDF} et \textit{Consulter le tableau de bord personnel} ; l’export doit être modélisé comme une action optionnelle déclenchée depuis le tableau de bord.
    \item Corriger le sens de la flèche \texttt{<<extend>>} si nécessaire : la flèche doit partir du cas optionnel \textit{Exporter le passeport talent en PDF} vers le cas de base \textit{Consulter le tableau de bord personnel}.
    \item Vérifier la relation entre \textit{Exporter le passeport talent en PDF} et \textit{Agrégation / Analyse IA} ; elle n’est pertinente que si l’export déclenche réellement un recalcul automatique du profil.
    \item Si l’analyse IA est aussi utilisée pour afficher le tableau de bord, relier directement le \textit{Service IA} au cas métier principal correspondant, sans surcharger le diagramme avec des détails trop techniques.
    \item Ajouter éventuellement un cas d’utilisation \textit{Consulter ses scores globaux} si cette information est réellement visible dans le tableau de bord.
    \item Ajouter éventuellement un cas d’utilisation \textit{Consulter ses badges ou certifications} si le passeport talent contient aussi ces éléments.
    \item Supprimer les détails trop techniques si l’objectif est un diagramme fonctionnel ; le diagramme doit rester centré sur les actions de l’employé.
    \item Réorganiser les cas d’utilisation pour réduire les croisements et améliorer la lisibilité, en plaçant les fonctionnalités de consultation au centre et l’export en bas à droite.
    \item Ajouter une légende expliquant les relations \texttt{<<include>>} et \texttt{<<extend>>}.

\end{itemize}
\color{black}
\begin{table}[H]
    \centering
    \renewcommand{\arraystretch}{1.5}
    \small
    \begin{tabularx}{0.95\textwidth}{|>{\raggedright\arraybackslash}p{4.5cm}|X|}
        \hline
        \rowcolor[HTML]{F2F2F2}
        \textbf{Élément} & \textbf{Description} \\
        \hline
        \textbf{Nom du cas d'utilisation} & Consulter le \textit{Dashboard} Personnel et Exporter son \textit{Passeport Talent} \\
        \hline
        \textbf{Acteurs} & Employé (Principal), Moteur IA (Secondaire) \\
        \hline
        \textbf{Objectif} & Permettre au collaborateur de piloter sa carrière et de certifier ses acquis. \\
        \hline
        \textbf{Description \newline (Scénario Nominal)} & 
        L'employé consulte son \textit{dashboard} (\textit{Design Bento}) incluant son radar de compétences et son profil psychologique IA. Il peut suivre ses formations et exporter son \textit{Passeport Talent}, un document PDF certifié résumant son parcours et son score de préparation (\textit{Readiness Score}). \\
        \hline
        \textbf{Résultat \newline (Post-conditions)} & Un collaborateur autonome dans son développement professionnel. \\
        \hline
    \end{tabularx}
    \caption{Fiche de cas d'utilisation : \textit{User Dashboard} \& Passeport (Module Employé)}
    \label{tab:uc_user_dashboard_passeport}
\end{table}

\vspace{1cm}

% --- MODULE 2 ---
\subsubsection*{CAS 2: Pilotage Stratégique (Executive Analytics)}

Destiné à la direction RH, ce module transforme les données individuelles en indicateurs de performance globaux (KPIs). Il permet d'identifier les écarts de compétences et de détecter les "Hauts Potentiels" (HiPo).

\begin{figure}[H]
    \centering
    \includegraphics[width=1\textwidth]{images/case2Sprint4.png}
    \caption{Diagramme de cas d'utilisation : Pilotage Stratégique \& Analytics}
    \label{fig:usecase_analytics}
\end{figure}
\color{red}
\begin{itemize}
    \item \textit{Diagramme 21 :} Renommer le titre en \textit{Diagramme de cas d’utilisation : pilotage stratégique et analytics RH} pour garder une formulation académique homogène.
    \item \textit{Diagramme 21 :} Remplacer le nom du module \textit{Pilotage Stratégique \& Analytics} par \textit{Pilotage stratégique et analyses RH}, afin d’éviter le mélange français--anglais.
    \item \textit{Diagramme 21 :} Renommer l’acteur \textit{Admin / Responsable RH} en \textit{Administrateur RH} ou \textit{Responsable RH}, puis utiliser le même nom dans tout le rapport.
    \item \textit{Diagramme 21 :} Remplacer l’acteur secondaire \textit{«Service» Moteur IA} par \textit{Service IA}, avec une notation plus simple et cohérente avec les autres diagrammes.
    \item \textit{Diagramme 21 :} Supprimer les flèches sur les associations entre l’acteur principal et les cas d’utilisation ; en UML, une association acteur--cas d’utilisation est une ligne simple sans flèche.
    \item \textit{Diagramme 21 :} Remplacer \textit{Consulter Executive Dashboard (Analytics RH)} par \textit{Consulter le tableau de bord exécutif RH}.
    \item \textit{Diagramme 21 :} Remplacer \textit{Filtrer \& Rechercher des Talents} par \textit{Filtrer et rechercher des talents}.
    \item \textit{Diagramme 21 :} Remplacer \textit{Suivre les KPIs de Performance Globale} par \textit{Suivre les indicateurs de performance globale}.
    \item \textit{Diagramme 21 :} Remplacer \textit{Identifier les Talents Haut Potentiel} par \textit{Identifier les talents à haut potentiel}.
    \item \textit{Diagramme 21 :} Remplacer \textit{Analyser Distribution des Compétences} par \textit{Analyser la distribution des compétences}.
    \item \textit{Diagramme 21 :} Remplacer \textit{Agrégation \& Analyse IA Globale} par \textit{Générer une analyse globale par IA} ou \textit{Produire une synthèse analytique globale}.
    \item \textit{Diagramme 21 :} Remplacer \textit{Exporter Rapport RH Global (Audit PDF)} par \textit{Exporter le rapport RH global}.
    \item \textit{Diagramme 21 :} Remplacer \textit{Authentication} par \textit{S’authentifier} afin d’uniformiser la langue des cas d’utilisation.
    \item \textit{Diagramme 21 :} Vérifier les relations \texttt{<<include>>} entre \textit{Consulter le tableau de bord exécutif RH} et les sous-cas de visualisation ; elles sont correctes uniquement si ces éléments sont toujours affichés dans le tableau de bord.
    \item \textit{Diagramme 21 :} Si les fonctionnalités \textit{Filtrer et rechercher des talents}, \textit{Identifier les talents à haut potentiel} ou \textit{Analyser la distribution des compétences} sont déclenchées à la demande, remplacer les relations \texttt{<<include>>} par des relations \texttt{<<extend>>} ou par de simples associations.
    \item \textit{Diagramme 21 :} Corriger la relation \texttt{<<extend>>} entre \textit{Exporter le rapport RH global} et \textit{Consulter le tableau de bord exécutif RH} ; elle est correcte si l’export est une action optionnelle disponible depuis le tableau de bord.
    \item \textit{Diagramme 21 :} Vérifier le sens de la relation \texttt{<<extend>>} : la flèche doit partir du cas optionnel \textit{Exporter le rapport RH global} vers le cas de base \textit{Consulter le tableau de bord exécutif RH}.
    \item \textit{Diagramme 21 :} Vérifier la relation \texttt{<<include>>} entre \textit{Exporter le rapport RH global} et \textit{Générer une analyse globale par IA} ; elle est pertinente seulement si l’export déclenche toujours une analyse IA.
    \item \textit{Diagramme 21 :} Si le rapport PDF exporte uniquement les données déjà affichées, supprimer la relation avec l’analyse IA et relier l’export uniquement au tableau de bord exécutif.
    \item \textit{Diagramme 21 :} Relier le \textit{Service IA} uniquement au cas \textit{Générer une analyse globale par IA}, afin d’éviter des liaisons trop générales.
    \item \textit{Diagramme 21 :} Ajouter éventuellement un cas \textit{Visualiser la cartographie des compétences} si le module propose une vue stratégique des compétences par équipe, métier ou département.
    \item \textit{Diagramme 21 :} Ajouter éventuellement un cas \textit{Comparer les profils et les postes cibles} si l’identification des talents repose sur une analyse d’adéquation profil--poste.
    \item \textit{Diagramme 21 :} Ajouter éventuellement un cas \textit{Consulter les alertes RH} si le tableau de bord signale automatiquement des écarts, risques ou besoins de formation.
    \item \textit{Diagramme 21 :} Réduire les grandes courbes \texttt{<<include>>}, car elles rendent le diagramme difficile à lire.
    \item \textit{Diagramme 21 :} Réorganiser le diagramme en plaçant le cas principal à gauche, les sous-fonctionnalités d’analyse au centre, l’export en bas à droite et le service IA à droite.
    \item \textit{Diagramme 21 :} Ajouter une légende expliquant les associations simples, les relations \texttt{<<include>>} et les relations \texttt{<<extend>>}.

\end{itemize}
\color{black}

\begin{table}[H]
    \centering
    \renewcommand{\arraystretch}{1.5}
    \small
    \begin{tabularx}{0.95\textwidth}{|>{\raggedright\arraybackslash}p{4.5cm}|X|}
        \hline
        \rowcolor[HTML]{F2F2F2}
        \textbf{Élément} & \textbf{Description} \\
        \hline
        \textbf{Nom du cas d'utilisation} & Consulter l'\textit{Executive Dashboard} \\
        \hline
        \textbf{Acteurs} & Administrateur RH (Principal), Moteur IA (Secondaire) \\
        \hline
        \textbf{Objectif} & Offrir une vue d'ensemble du capital humain pour le pilotage stratégique. \\
        \hline
        \textbf{Pré-conditions} & L'Admin est authentifié et possède le rôle \texttt{ROLE\_ADMIN}. \\
        \hline
        \textbf{Description \newline (Scénario Nominal)} & 
        1. L'Admin accède à l'espace \textit{Dashboard}. \newline
        2. Le système sollicite le service IA pour agréger les données de compétences. \newline
        3. Le \textit{dashboard} affiche les KPIs : taux de complétion, distribution des niveaux, et identification des Hauts Potentiels (\textit{HiPo}). \newline
        4. L'Admin peut filtrer les résultats par département ou par technologie. \\
        \hline
        \textbf{Post-conditions} & L'Admin possède une vision claire du vivier de talents pour la prise de décision. \\
        \hline
    \end{tabularx}
    \caption{Fiche de cas d'utilisation : Consulter l'\textit{Executive Dashboard}}
    \label{tab:uc_executive_dashboard}
\end{table}

\vspace{1cm}

% --- MODULE 3 ---
\subsubsection*{Cas 3 : Hub de Communication (Real-time Hub)}

Ce module assure la réactivité du système via les Server-Sent Events (SSE). Il permet une communication instantanée entre l'administration et les employés sans rafraîchissement de page.

\begin{figure}[H]
    \centering
    \includegraphics[width=1\textwidth]{images/case3Sprint4.png}
    \caption{Diagramme de cas d'utilisation : Communication Hub \& Server-Sent Events (SSE)}
    \label{fig:usecase_comm}
\end{figure}
\color{red}
\begin{itemize}
    \item \textit{Diagramme 22 :} Renommer le titre en \textit{Diagramme de cas d’utilisation : communication et notifications temps réel} pour garder une formulation plus académique.
    \item \textit{Diagramme 22 :} Remplacer le nom du module \textit{Communication Hub \& Server-Sent Events (SSE)} par \textit{Module de communication et notifications temps réel}.
    \item \textit{Diagramme 22 :} Renommer l’acteur \textit{Admin / Responsable RH} en \textit{Administrateur RH} afin d’utiliser la même terminologie dans tout le rapport.
    \item \textit{Diagramme 22 :} Renommer l’acteur \textit{Service SSE (Notifications Temps Réel)} en \textit{Service de notification temps réel}.
    \item \textit{Diagramme 22 :} Supprimer les flèches sur les associations entre les acteurs et les cas d’utilisation ; en UML, une association acteur--cas d’utilisation est une ligne simple sans flèche.
    \item \textit{Diagramme 22 :} Remplacer \textit{Diffuser une annonce globale (Broadcast)} par \textit{Diffuser une annonce globale}.
    \item \textit{Diagramme 22 :} Remplacer \textit{Consulter le Centre de Notifications (Historique)} par \textit{Consulter le centre de notifications}.
    \item \textit{Diagramme 22 :} Remplacer \textit{Recevoir des alertes instantanées (Toasts)} par \textit{Recevoir des alertes instantanées}.
    \item \textit{Diagramme 22 :} Remplacer \textit{Maintenir la connexion Server-Sent Events (SSE)} par \textit{Maintenir la connexion temps réel}.
    \item \textit{Diagramme 22 :} Vérifier si \textit{Maintenir la connexion temps réel} doit apparaître comme un cas d’utilisation ; il s’agit plutôt d’un mécanisme technique interne à décrire dans un diagramme d’architecture ou de séquence.
    \item \textit{Diagramme 22 :} Si le diagramme doit rester fonctionnel, supprimer \textit{Maintenir la connexion temps réel} et garder seulement les cas liés aux notifications visibles par l’utilisateur.
    \item \textit{Diagramme 22 :} Vérifier la relation \texttt{<<include>>} entre \textit{Notifier une validation de formation} et \textit{Recevoir des alertes instantanées} ; elle est correcte si une alerte est toujours envoyée après validation.
    \item \textit{Diagramme 22 :} Vérifier la relation \texttt{<<include>>} entre \textit{Diffuser une annonce globale} et \textit{Recevoir des alertes instantanées} ; elle est correcte si l’annonce apparaît toujours sous forme d’alerte instantanée.
    \item \textit{Diagramme 22 :} Vérifier la relation \texttt{<<include>>} entre \textit{Consulter le centre de notifications} et \textit{S’authentifier} ; elle est acceptable si l’accès à l’historique nécessite toujours une authentification.
    \item \textit{Diagramme 22 :} Corriger la relation \texttt{<<extend>>} entre \textit{Marquer une notification comme lue} et \textit{Consulter le centre de notifications} ; elle est cohérente si le marquage est une action optionnelle réalisée depuis l’historique.
    \item \textit{Diagramme 22 :} Vérifier le sens de la relation \texttt{<<extend>>} : la flèche doit partir de \textit{Marquer une notification comme lue} vers \textit{Consulter le centre de notifications}.
    \item \textit{Diagramme 22 :} Vérifier si \textit{Envoyer un message de relance individuel} appartient à ce module ; si oui, le relier clairement à l’acteur \textit{Administrateur RH}.
    \item \textit{Diagramme 22 :} Ajouter une relation entre \textit{Envoyer un message de relance individuel} et \textit{Recevoir des alertes instantanées} uniquement si la relance produit réellement une notification instantanée.
    \item \textit{Diagramme 22 :} Ajouter éventuellement un cas \textit{Consulter les notifications non lues} si cette fonctionnalité existe dans l’interface.
    \item \textit{Diagramme 22 :} Ajouter éventuellement un cas \textit{Filtrer les notifications} si le centre de notifications permet de trier par type, date ou statut.
    \item \textit{Diagramme 22 :} Relier le \textit{Service de notification temps réel} uniquement aux cas d’envoi ou de réception de notifications, afin d’éviter une liaison trop générale.
    \item \textit{Diagramme 22 :} Réduire les grandes courbes et les croisements de lignes, car ils diminuent la lisibilité du diagramme.
    \item \textit{Diagramme 22 :} Réorganiser le diagramme en plaçant les actions de l’administrateur RH à gauche, les actions de l’employé au centre et le service de notification à droite.
    \item \textit{Diagramme 22 :} Ajouter une légende expliquant les associations simples, les relations \texttt{<<include>>} et les relations \texttt{<<extend>>}.

\end{itemize}
\color{black}

\begin{table}[H]
    \centering
    \renewcommand{\arraystretch}{1.5}
    \small
    \begin{tabularx}{0.95\textwidth}{|>{\raggedright\arraybackslash}p{4.5cm}|X|}
        \hline
        \rowcolor[HTML]{F2F2F2}
        \textbf{Élément} & \textbf{Description} \\
        \hline
        \textbf{Nom du cas d'utilisation} & Recevoir des alertes instantanées (\textit{Toasts}) \\
        \hline
        \textbf{Acteurs} & Employé, Admin, Service \texttt{SSE} \\
        \hline
        \textbf{Objectif} & Informer instantanément l'utilisateur d'un événement critique sans rechargement de page. \\
        \hline
        \textbf{Pré-conditions} & Une connexion \texttt{SSE} (\textit{Server-Sent Events}) est active et maintenue entre le client et le serveur. \\
        \hline
        \textbf{Description \newline (Scénario Nominal)} & 
        1. Un événement survient sur le serveur (ex : l'Admin valide une formation). \newline
        2. Le service \texttt{NotificationSseService} identifie l'utilisateur concerné. \newline
        3. Le serveur pousse un flux \texttt{JSON} via le canal \texttt{SSE} ouvert. \newline
        4. Le \textit{frontend} \texttt{Angular} intercepte l'événement et affiche un "\textit{Toast}" visuel via \texttt{Ngx-Toastr}. \\
        \hline
        \textbf{Post-conditions} & L'utilisateur est informé en temps réel, améliorant significativement la réactivité du système. \\
        \hline
    \end{tabularx}
    \caption{Fiche de cas d'utilisation : Recevoir des Alertes Temps Réel}
    \label{tab:uc_alerts_sse}
\end{table}

\vspace{1cm}

% --- MODULE 4 ---
\subsubsection{Cas 4: Gouvernance \& Administration RH}

Module de gestion du cycle de vie des collaborateurs, permettant de gérer les habilitations de sécurité (RBAC) et de contrôler l'accès au système.
\begin{table}[H]
    \centering
    \renewcommand{\arraystretch}{1.5}
    \small
    \begin{tabularx}{0.95\textwidth}{|>{\raggedright\arraybackslash}p{4.5cm}|X|}
        \hline
        \rowcolor[HTML]{F2F2F2}
        \textbf{Élément} & \textbf{Description} \\
        \hline
        \textbf{Nom du cas d'utilisation} & Gérer le Vivier d'Employés \\
        \hline
        \textbf{Acteurs} & Administrateur RH \\
        \hline
        \textbf{Objectif} & Assurer la gestion administrative et sécurisée des comptes utilisateurs. \\
        \hline
        \textbf{Description \newline (Scénario Nominal)} & 
        L'Admin possède un \textit{CRUD} complet sur les employés (Création, Modification, Activation/Désactivation). Il peut consulter le profil public détaillé de chaque membre (liens \texttt{GitHub}/\texttt{LinkedIn}, historique de performance) pour un suivi personnalisé. \\
        \hline
        \textbf{Résultat \newline (Post-conditions)} & Une base de données RH toujours à jour et conforme aux politiques de sécurité (\textit{RBAC}). \\
        \hline
    \end{tabularx}
    \caption{Fiche de cas d'utilisation : Gouvernance \& Administration RH (Module Gestion)}
    \label{tab:uc_gouvernance_rh}
\end{table}

\subsection{Conception}

% La numérotation 3.2 s'adaptera automatiquement

Cette phase définit les choix techniques et architecturaux permettant de répondre aux besoins analysés précédemment. L'objectif est de garantir la robustesse du système et la fluidité des échanges de données.

\subsubsection*{1. Objectif de conception}

L'objectif principal de la conception pour ce sprint est la \textit{scalabilité} et l'intégration. Il s'agit de s'assurer que :

\begin{itemize}
    \item Le moteur de \textit{reporting} peut générer des fichiers lourds sans impacter les performances.
    \item Le système de notifications peut gérer des centaines de connexions simultanées via WebSockets.
    \item L'architecture permet une séparation nette entre la logique métier (Java) et l'intelligence artificielle (Python).
\end{itemize}
\subsubsection*{2. Architecture Technique}
Le système repose sur une architecture N-Tier (Multicouche), renforcée par une approche Micro-services hybride.
\begin{itemize}
    \item \textbf{Frontend :} Framework \texttt{Angular 17}, gérant l'état et la visualisation \textit{Bento Grid}.
    \item \textbf{Backend (Cœur) :} \texttt{Spring Boot 3}, agissant comme un chef d'orchestre (\textit{Gateway}) pour la sécurité et les données.
    \item \textbf{Service IA :} \texttt{FastAPI} (Python), dédié exclusivement aux calculs intensifs et aux modèles prédictifs.
    \item \textbf{Communication Temps Réel :} Protocole \texttt{SEE} via HTTP pour la gestion des notifications.
\end{itemize}

\subsubsection*{3. Algorithme de Détection \textit{HiPo} (\textit{High Potential})}

Pour automatiser l'identification des talents clés, nous avons implémenté un algorithme \textit{"Data-Driven"} basé sur la \textbf{Matrice 9-Box}, croisant les deux axes fondamentaux de l'évaluation RH :

\begin{itemize}
    \item \textbf{Performance ($S_{perf}$) :} Calculée via la moyenne arithmétique globale (sur 100) de l'ensemble des tests techniques passés par le candidat sur la plateforme.
    \item \textbf{Potentiel ($S_{pot}$) :} Évalué via la normalisation des scores comportementaux issus du test de personnalité \textit{PCM} (\textit{Process Communication Model}) ou, à défaut, des prédictions de l'IA.
\end{itemize}

L'algorithme génère un score \textit{HiPo} global pondéré selon la formule suivante :

\begin{equation}
    Score_{HiPo} = (S_{perf} \times 0.6) + (S_{pot} \times 0.4)
\end{equation}

Enfin, des seuils mathématiques placent automatiquement le collaborateur dans l'un des 9 quadrants de la matrice. Par exemple, un profil respectant la condition :

\begin{equation}
    (S_{perf} \ge 80\%) \cap (S_{pot} \ge 80\%)
\end{equation}

est immédiatement labellisé \textbf{"High Potential"} (\textit{Top Right}), permettant aux décideurs d'identifier les futurs leaders de manière objective et sans biais cognitifs. Cette classification automatisée facilite le pilotage stratégique du capital humain au sein de \textbf{TalentPredict}.

\subsubsection*{4. Modèle de données }

Pour ce sprint, nous avons enrichi le schéma de données avec :

\begin{itemize}
    \item \textbf{Entité \texttt{Notification} :} Stockage de l'historique des alertes (Type, Message, Date, Statut de lecture).
    \item \textbf{Entité \texttt{Audit/Historique} :} Pour tracer les évaluations et alimenter les graphiques d'évolution du \textit{dashboard}.
    \item \textbf{Entité \texttt{TalentProfile} :} Extension des données utilisateurs pour inclure les liens sociaux et la biographie.
\end{itemize}

\subsubsection*{5. Contrats API (Exemple : Dashboard)}

Pour garantir l'interopérabilité, nous utilisons des DTO (\textit{Data Transfer Objects}). Exemple de réponse JSON pour le \textit{Dashboard} Admin :

\begin{lstlisting}[language=json, caption={Exemple de structure JSON pour le Dashboard Admin}, label={lst:json_dashboard}]
{
  "totalEmployees": 150,
  "highPotentialCount": 12,
  "skillsDistribution": { 
      "Java": 85, 
      "Python": 70, 
      "BehavioralSkills": 90 
  },
  "recentActivities": [...]
}
\end{lstlisting}

\subsubsection{6. Diagrammes de Séquence}

Cette section détaille la chorégraphie des interactions entre les différents composants de la plateforme à travers quatre scénarios clés du Sprint 4.

\subsubsection{Cas n°1 : Pilotage Collaborateur \& Passeport Talent}

Ce diagramme illustre le flux complexe d'agrégation de données pour l'employé.
\begin{figure}[H]
    \centering
    \IfFileExists{images/seq1Spring4.png}{%
      \includegraphics[width=0.9\textwidth]{images/seq1Spring4.png}%
    }{%
      \fbox{\parbox[c][6cm][c]{0.82\textwidth}{\centering Figure non disponible : images/seq1Spring4.png}}%
    }
    \caption{Diagramme de séquence : Génération du Passeport Talent}
    \label{fig:seq_cas1}
\end{figure}

\color{red}
\begin{itemize}
    \item \textit{Diagramme 24 :} Renommer le titre en \textit{Diagramme de séquence : génération du passeport talent} pour garder une formulation plus simple et homogène.
    \item \textit{Diagramme 24 :} Vérifier la cohérence du titre interne \textit{User Dashboard \& Passport} ; si le rapport est en français, le remplacer par \textit{Tableau de bord utilisateur et passeport talent}.
    \item \textit{Diagramme 24 :} Harmoniser la langue de toutes les lignes de vie, car le diagramme mélange français et anglais : \textit{Frontend}, \textit{Security Layer}, \textit{Business Logic}, \textit{AI Engine}, \textit{Persistence Layer}.
    \item \textit{Diagramme 24 :} Renommer \textit{Rest Controller} en \textit{Contrôleur REST} ou garder toute la terminologie technique en anglais, mais de manière uniforme.
    \item \textit{Diagramme 24 :} Renommer \textit{Business Logic (Service Layer)} en \textit{Service métier} ou \textit{Couche métier}.
    \item \textit{Diagramme 24 :} Remplacer \textit{AI Engine (FastAPI / Python)} par \textit{Microservice IA} et préciser FastAPI/Python dans une note technique si nécessaire.
    \item \textit{Diagramme 24 :} Remplacer \textit{Persistence Layer (PostgreSQL)} par \textit{Base de données PostgreSQL} pour une formulation plus claire.
    \item \textit{Diagramme 24 :} Ajouter un fragment \texttt{alt} ou \texttt{break} dans la phase 2 pour gérer le cas où l’utilisateur n’est pas autorisé à exporter le passeport.
    \item \textit{Diagramme 24 :} Dans la phase 1, le fragment \texttt{alt} d’authentification est pertinent, mais il faut expliciter les deux branches avec des intitulés plus lisibles : \textit{échec d’authentification} et \textit{authentification réussie}.
    \item \textit{Diagramme 24 :} Ajouter un retour \textit{401 Unauthorized} ou \textit{403 Forbidden} cohérent selon le cas, puis conserver une seule formulation d’erreur dans le diagramme.
    \item \textit{Diagramme 24 :} Vérifier si \textit{Validation Signature JWT} doit être renommé en \textit{Validation du jeton JWT}, pour rester plus clair pour les étudiants.
    \item \textit{Diagramme 24 :} Clarifier le message \textit{Dispatcher vers DashboardController}, qui relève d’un détail interne ; il peut être simplifié en un appel direct vers le contrôleur.
    \item \textit{Diagramme 24 :} Vérifier si \textit{loadUserDashboardData(userId)} doit être scindé en plusieurs appels explicites si le dashboard récupère plusieurs sources distinctes.
    \item \textit{Diagramme 24 :} Le message \textit{Fetch Skills, Tests \& History} doit être reformulé en français, par exemple \textit{Récupérer les compétences, les tests et l’historique}.
    \item \textit{Diagramme 24 :} Le message \textit{Data ResultSet} est trop technique ; le remplacer par \textit{Données utilisateur consolidées}.
    \item \textit{Diagramme 24 :} Vérifier si l’appel au microservice IA dans la phase 1 est toujours exécuté ; si ce n’est pas systématique, l’encadrer par un fragment \texttt{opt}.
    \item \textit{Diagramme 24 :} Remplacer \textit{Analyse Prédictive \& Profiling Persona} par \textit{Analyse prédictive et génération du profil} pour une formulation plus claire.
    \item \textit{Diagramme 24 :} Remplacer \textit{Insights JSON (PCM, Career Path)} par \textit{Résultats d’analyse IA} ou \textit{Profil IA et recommandations}, afin d’éviter un libellé trop technique.
    \item \textit{Diagramme 24 :} Le message \textit{DTO Consolidé (Bento Grid Ready)} doit être simplifié, par exemple en \textit{DTO consolidé du tableau de bord}.
    \item \textit{Diagramme 24 :} Remplacer \textit{Rendering Bento Grid (Chart.js)} par \textit{Rendu du tableau de bord} si l’objectif est fonctionnel ; les technologies comme Chart.js peuvent être déplacées dans le texte.
    \item \textit{Diagramme 24 :} Dans la phase 2, préciser si l’export du passeport passe aussi par une validation JWT ; si oui, montrer explicitement cette vérification.
    \item \textit{Diagramme 24 :} Vérifier si l’appel \textit{/api/v1/reporting/export} doit contenir un identifiant utilisateur ou un type de document dans les paramètres.
    \item \textit{Diagramme 24 :} Renommer \textit{orchestratePdfGeneration(userId)} par \textit{Générer le passeport PDF} ou \textit{Orchestrer la génération du passeport}, afin d’être plus explicite.
    \item \textit{Diagramme 24 :} Clarifier le rôle du moteur \textit{OpenHTMLToPDF} ; si c’est un composant interne, il peut être représenté comme un service technique plutôt qu’une ligne de vie métier complète.
    \item \textit{Diagramme 24 :} Le message \textit{Récupération Profil RH final} doit être renommé si le passeport concerne l’employé ; utiliser par exemple \textit{Récupération du profil final de l’employé}.
    \item \textit{Diagramme 24 :} Le message \textit{User Details} doit être reformulé en français, par exemple \textit{Détails de l’utilisateur}.
    \item \textit{Diagramme 24 :} Remplacer \textit{Conversion HTML -> PDF Stream} par \textit{Conversion du modèle HTML en flux PDF}.
    \item \textit{Diagramme 24 :} Le message \textit{Resource (InputStream)} est trop technique ; le remplacer par \textit{Document PDF généré}.
    \item \textit{Diagramme 24 :} Remplacer \textit{Transfer-Encoding: chunked (PDF File)} par \textit{Transmission du fichier PDF} pour rester centré sur la logique fonctionnelle.
    \item \textit{Diagramme 24 :} Ajouter un fragment \texttt{alt} dans la phase 2 pour gérer le cas où la génération PDF échoue ou si les données nécessaires sont incomplètes.
    \item \textit{Diagramme 24 :} Vérifier que toutes les réponses sont bien représentées par des flèches de retour en pointillés, et que les appels principaux restent en lignes pleines.
    \item \textit{Diagramme 24 :} Ajouter des barres d’activation cohérentes sur toutes les lignes de vie, notamment le moteur PDF, le microservice IA et la couche de persistance.
    \item \textit{Diagramme 24 :} Réduire le nombre de notes internes dans la figure si elles alourdissent le diagramme ; certaines précisions techniques peuvent être déplacées dans le texte.

\end{itemize}
\color{black}
\paragraph*{ \textit{User Dashboard} \& Passeport Talent}
Ce flux illustre le processus d'agrégation de données et de certification individuelle au sein de \textbf{TalentPredict}.
\begin{itemize}
    \item \textbf{Phase d'Analyse :} Lors de la connexion, le \textit{backend} \texttt{Spring Boot} orchestre la collecte des scores de tests et des compétences depuis \texttt{PostgreSQL}. Un appel asynchrone est envoyé au microservice IA (\texttt{FastAPI}) pour générer le "\textit{Profiling Persona}" (profil psychologique et recommandations).
    \item \textbf{Rendu et Export :} Les données consolidées sont renvoyées au format \texttt{JSON} pour alimenter la \textit{Bento Grid} du \textit{frontend}. Pour la certification, le moteur \texttt{OpenHTMLToPDF} transforme dynamiquement un \textit{template} \texttt{Thymeleaf} en un document PDF sécurisé, intégrant les résultats officiels de l'employé.
\end{itemize}

\subsubsection{Cas n°2 : Pilotage Stratégique \& Audit RH}

Ce flux met en évidence les capacités de \textit{Business Intelligence} (BI) de la plateforme pour les administrateurs.

\begin{figure}[H]
    \centering
    \includegraphics[width=0.9\textwidth]{images/seq2Spring4.png}
    \caption{Diagramme de séquence : Pilotage stratégique et Audit RH}
    \label{fig:seq_cas2}
\end{figure}

\color{red}
\begin{itemize}
    \item \textit{Diagramme 25 :} Renommer le titre en \textit{Diagramme de séquence : pilotage stratégique et audit RH} pour garder une formulation plus simple et homogène.
    \item \textit{Diagramme 25 :} Remplacer l’acteur \textit{Admin / Responsable RH} par \textit{Administrateur RH} afin d’utiliser la même terminologie dans tout le rapport.
    \item \textit{Diagramme 25 :} Harmoniser la langue des lignes de vie, car le diagramme mélange français et anglais : \textit{Frontend}, \textit{RBAC Guard}, \textit{Admin Controller}, \textit{Analytics Service}, \textit{Global Export Service}, \textit{IA Service}.
    \item \textit{Diagramme 25 :} Remplacer \textit{Frontend (Angular 17)} par \textit{Interface Angular} ou conserver le nom technique, mais appliquer la même logique à toutes les lignes de vie.
    \item \textit{Diagramme 25 :} Remplacer \textit{RBAC Guard} par \textit{Garde d’accès RBAC} ou \textit{Filtre de sécurité RBAC} pour améliorer la lisibilité.
    \item \textit{Diagramme 25 :} Remplacer \textit{Backend Analytics Engine} par \textit{Moteur backend d’analyse RH} si le rapport est rédigé en français.
    \item \textit{Diagramme 25 :} Remplacer \textit{IA Service} par \textit{Service IA}, formulation plus correcte en français.
    \item \textit{Diagramme 25 :} Remplacer \textit{PostgreSQL} par \textit{Base de données PostgreSQL} pour garder une désignation cohérente avec les autres diagrammes.
    \item \textit{Diagramme 25 :} Le fragment \texttt{alt} sur le contrôle d’accès est pertinent ; il faut toutefois renommer les branches en \textit{accès refusé} et \textit{accès autorisé} avec une casse homogène.
    \item \textit{Diagramme 25 :} Vérifier si le retour \textit{403 Forbidden} correspond bien au cas d’un rôle insuffisant ; pour un jeton invalide ou expiré, utiliser plutôt \textit{401 Unauthorized}.
    \item \textit{Diagramme 25 :} Ajouter explicitement la validation du jeton JWT avant la vérification du rôle administrateur, car le diagramme mentionne seulement le contrôle RBAC.
    \item \textit{Diagramme 25 :} Remplacer \textit{Vérifier Rôle (ROLE\_ADMIN)} par \textit{Vérifier le rôle administrateur}.
    \item \textit{Diagramme 25 :} Remplacer \textit{Contrôle d’accès granulaire (RBAC)} par une note plus courte, par exemple \textit{Contrôle d’accès selon le rôle}.
    \item \textit{Diagramme 25 :} Remplacer \textit{fetchGlobalKPIs()} par \textit{récupérer les indicateurs globaux} ou harmoniser tous les messages en anglais.
    \item \textit{Diagramme 25 :} Remplacer \textit{Requête agrégée (GROUP BY, AVG)} par \textit{Requête d’agrégation des indicateurs}, car les détails SQL peuvent être expliqués dans le texte.
    \item \textit{Diagramme 25 :} Remplacer \textit{Données brutes (Skills, Complétion)} par \textit{Données de compétences et de complétion}.
    \item \textit{Diagramme 25 :} Remplacer \textit{Détection des Hauts Potentiels} par \textit{Détection des talents à haut potentiel}.
    \item \textit{Diagramme 25 :} Remplacer \textit{POST /analyze/macro-trends} par une formulation plus fonctionnelle, par exemple \textit{Analyser les tendances globales}.
    \item \textit{Diagramme 25 :} Remplacer \textit{Liste des ``HiPo'' + Tendances} par \textit{Liste des talents à haut potentiel et tendances détectées}.
    \item \textit{Diagramme 25 :} Vérifier si l’appel au service IA est toujours exécuté lors de la consultation du tableau de bord ; sinon, encadrer cette partie avec un fragment \texttt{opt}.
    \item \textit{Diagramme 25 :} Remplacer \textit{AdminOverviewDTO (Data Ready)} par \textit{DTO de synthèse RH prêt}.
    \item \textit{Diagramme 25 :} Remplacer \textit{Response JSON (200 OK)} par \textit{Réponse JSON avec les indicateurs RH}.
    \item \textit{Diagramme 25 :} Remplacer \textit{Rendu Graphiques (Distribution, KPIs)} par \textit{Affichage des graphiques et indicateurs}.
    \item \textit{Diagramme 25 :} Dans la phase 2, ajouter une vérification d’autorisation avant la génération du rapport d’audit, comme cela est fait dans la phase 1.
    \item \textit{Diagramme 25 :} Vérifier si l’export du rapport doit utiliser le même contrôle RBAC que le tableau de bord exécutif ; si oui, représenter clairement cette étape.
    \item \textit{Diagramme 25 :} Remplacer \textit{generateGlobalAuditPdf()} par \textit{Générer le rapport d’audit RH global}.
    \item \textit{Diagramme 25 :} Remplacer \textit{Extraire statistiques consolidées} par \textit{Extraire les statistiques consolidées}.
    \item \textit{Diagramme 25 :} Remplacer \textit{Global Data Set} par \textit{Jeu de données global consolidé}.
    \item \textit{Diagramme 25 :} Remplacer \textit{Compilation PDF (Audit Format)} par \textit{Compilation du rapport d’audit au format PDF}.
    \item \textit{Diagramme 25 :} Remplacer \textit{Resource (PDF)} par \textit{Rapport PDF généré}.
    \item \textit{Diagramme 25 :} Remplacer \textit{application/pdf stream} par \textit{Flux du fichier PDF}.
    \item \textit{Diagramme 25 :} Remplacer \textit{Téléchargement du Rapport\_Global.pdf} par \textit{Téléchargement du rapport d’audit RH global}.
    \item \textit{Diagramme 25 :} Ajouter un fragment \texttt{alt} dans la phase 2 pour distinguer le cas \textit{rapport généré avec succès} du cas \textit{erreur lors de la génération du PDF}.
    \item \textit{Diagramme 25 :} Ajouter une étape de journalisation ou d’audit de l’export si le téléchargement du rapport RH global doit être traçable.
    \item \textit{Diagramme 25 :} Vérifier que toutes les réponses sont représentées par des flèches de retour en pointillés, notamment les retours de la base de données et du service IA.
    \item \textit{Diagramme 25 :} Ajouter des barres d’activation cohérentes sur toutes les lignes de vie qui exécutent des traitements, notamment \textit{Analytics Service}, \textit{Global Export Service}, \textit{Service IA} et \textit{Base de données PostgreSQL}.
    \item \textit{Diagramme 25 :} Réduire les libellés trop longs pour améliorer la lisibilité dans le rapport, surtout dans les messages liés aux requêtes et aux réponses.

\end{itemize}
\color{black}
\paragraph*{ \textit{Executive Dashboard} \& Audit RH Global}
Ce flux détaille l'intelligence décisionnelle destinée à la direction RH.
\begin{itemize}
    \item \textbf{\textit{Business Intelligence} :} Le système vérifie d'abord les droits d'accès via le \texttt{RBAC Guard}. Une fois autorisé, le \texttt{AnalyticsService} exécute des requêtes \texttt{SQL} complexes (agrégations et moyennes) pour obtenir l'état global du capital humain.
    \item \textbf{Analyse de Tendances :} Les données brutes sont transmises au service IA pour identifier les Hauts Potentiels (\textit{HiPo}) et les trajectoires de croissance. Le résultat est une vue stratégique interactive complétée par la possibilité de générer un rapport d'audit PDF consolidé.
\end{itemize}

\subsubsection{Cas n°3 : Communication Hub \& Server-Sent Events (SSE)}

Ce diagramme détaille l'implémentation de la communication bidirectionnelle en temps réel.
\begin{figure}[H]
    \centering
    \includegraphics[width=0.9\textwidth]{images/seq3Sprint4.png}
    \caption{Diagramme de séquence : Flux de notification temps réel}
    \label{fig:seq_cas3}
\end{figure}

\color{red}
\begin{itemize}
    \item \textit{Diagramme 26 :} Renommer le titre en \textit{Diagramme de séquence : flux de notification temps réel} pour garder une formulation académique et homogène.
    \item \textit{Diagramme 26 :} Remplacer l’acteur \textit{Employé / Client} par \textit{Employé}, \textit{Client} ou \textit{Utilisateur}, selon le vocabulaire adopté dans le rapport.
    \item \textit{Diagramme 26 :} Remplacer \textit{Admin / System Action} par un nom plus clair, par exemple \textit{Administrateur RH} ou \textit{Action système}, car le rôle actuel mélange un acteur humain et un déclencheur automatique.
    \item \textit{Diagramme 26 :} Harmoniser la langue des lignes de vie, car le diagramme mélange français et anglais : \textit{Frontend}, \textit{NotificationService}, \textit{NotificationController}, \textit{NotificationSseService}, \textit{Active Session}.
    \item \textit{Diagramme 26 :} Remplacer \textit{Frontend (Angular)} par \textit{Interface Angular} ou conserver la terminologie anglaise partout, mais de manière uniforme.
    \item \textit{Diagramme 26 :} Remplacer \textit{Backend (Spring Boot)} par \textit{Backend Spring Boot} pour une formulation plus standard.
    \item \textit{Diagramme 26 :} Remplacer \textit{Évènement Métier} par \textit{Événement métier} pour corriger l’orthographe.
    \item \textit{Diagramme 26 :} Corriger \textit{NotificationService (RxJS / EventSource)} si cette ligne représente un service Angular côté client ; il serait plus clair d’écrire \textit{Service de notification Angular}.
    \item \textit{Diagramme 26 :} Corriger \textit{Ngx-Toastr UI Component} en \textit{Composant d’affichage des notifications} ou préciser \textit{Ngx-Toastr} dans une note technique.
    \item \textit{Diagramme 26 :} Ajouter un fragment \texttt{alt} dans la phase 1 pour distinguer le cas \textit{jeton JWT valide} du cas \textit{jeton invalide ou expiré}.
    \item \textit{Diagramme 26 :} Ajouter un retour d’erreur \textit{401 Unauthorized} si le jeton JWT est invalide lors de l’abonnement au flux SSE.
    \item \textit{Diagramme 26 :} Remplacer \textit{GET /api/notifications/subscribe (JWT)} par une formulation plus lisible, par exemple \textit{S’abonner au flux de notifications}.
    \item \textit{Diagramme 26 :} Vérifier si le JWT est transmis dans l’en-tête ou dans un paramètre de connexion SSE ; ce détail doit être cohérent avec l’implémentation réelle.
    \item \textit{Diagramme 26 :} Remplacer \textit{createEmitter(userId)} par \textit{Créer une session SSE pour l’utilisateur} ou harmoniser tous les messages en anglais.
    \item \textit{Diagramme 26 :} Remplacer \textit{new SseEmitter(timeout=30min)} par \textit{Créer un émetteur SSE avec délai d’expiration}, car le détail technique peut être expliqué dans le texte.
    \item \textit{Diagramme 26 :} Remplacer \textit{Store Emitter in Map} par \textit{Enregistrer la session SSE active}.
    \item \textit{Diagramme 26 :} Remplacer \textit{Connection Established} par \textit{Connexion SSE établie}.
    \item \textit{Diagramme 26 :} Remplacer \textit{HTTP 200 (text/event-stream)} par \textit{Ouverture du flux SSE}, afin de garder le diagramme centré sur le comportement fonctionnel.
    \item \textit{Diagramme 26 :} Dans la phase 2, conserver le fragment \texttt{loop}, car il représente correctement l’envoi périodique du signal de maintien de connexion.
    \item \textit{Diagramme 26 :} Remplacer \textit{send(comment: ``heartbeat'')} par \textit{Envoyer un signal de maintien de connexion}.
    \item \textit{Diagramme 26 :} Corriger la note \textit{Prévient la coupure Proxy/Docker} en \textit{Prévenir la coupure de connexion par le proxy ou Docker}.
    \item \textit{Diagramme 26 :} Dans la phase 3, clarifier le déclencheur métier : préciser si l’alerte est déclenchée par l’administrateur RH, par une validation de formation ou par une action système.
    \item \textit{Diagramme 26 :} Remplacer \textit{sendToUser(userId, ``Formation Approuvée'')} par \textit{Envoyer une notification à l’utilisateur}.
    \item \textit{Diagramme 26 :} Remplacer \textit{send(event: ``message'', data: JSON)} par \textit{Transmettre l’événement SSE au client}.
    \item \textit{Diagramme 26 :} Remplacer \textit{Data: \{title: ``Succès''\}} par \textit{Données de notification reçues}.
    \item \textit{Diagramme 26 :} Remplacer \textit{showSuccess()} par \textit{Afficher une alerte visuelle}.
    \item \textit{Diagramme 26 :} Le fragment \texttt{alt} \textit{émetteur trouvé / émetteur non trouvé} est pertinent ; il faut toutefois expliquer clairement que l’émetteur correspond à une session SSE active de l’utilisateur.
    \item \textit{Diagramme 26 :} Dans la branche \textit{émetteur non trouvé}, remplacer \textit{Log for next login} par \textit{Enregistrer la notification pour consultation ultérieure}, si la notification est réellement conservée.
    \item \textit{Diagramme 26 :} Ajouter la base de données ou le dépôt de notifications dans le diagramme si les notifications non reçues sont stockées pour être affichées plus tard.
    \item \textit{Diagramme 26 :} Dans la phase 4, remplacer \textit{Close Connection} par \textit{Fermer la connexion SSE}.
    \item \textit{Diagramme 26 :} Remplacer \textit{onCompletion()} par \textit{Détecter la fermeture de la connexion}.
    \item \textit{Diagramme 26 :} Remplacer \textit{Remove userId from Map} par \textit{Supprimer la session SSE active de l’utilisateur}.
    \item \textit{Diagramme 26 :} Ajouter un scénario de reconnexion automatique si l’application frontend tente de rétablir le flux SSE après une perte de connexion.
    \item \textit{Diagramme 26 :} Ajouter un fragment \texttt{alt} ou \texttt{opt} pour traiter le cas d’expiration du \textit{SseEmitter} après le délai configuré.
    \item \textit{Diagramme 26 :} Vérifier que toutes les réponses sont représentées par des flèches de retour en pointillés, notamment les confirmations du backend vers le frontend.
    \item \textit{Diagramme 26 :} Ajouter des barres d’activation cohérentes pour toutes les lignes de vie qui exécutent un traitement, notamment \textit{NotificationController}, \textit{NotificationSseService} et \textit{SseEmitter}.
    \item \textit{Diagramme 26 :} Réduire les libellés techniques trop longs afin d’améliorer la lisibilité dans le rapport.

\end{itemize}
\color{black}
\paragraph*{\textit{Communication Hub} \& \textit{Server-Sent Events} (\texttt{SSE})}
Ce mécanisme expose la mécanique du temps réel unidirectionnel pour une réactivité maximale.
\begin{itemize}
    \item \textbf{Maintien du Flux :} Contrairement au \texttt{HTTP} classique, une connexion persistante (\texttt{text/event-stream}) est établie entre \texttt{Angular} et \texttt{Spring Boot}. Un mécanisme de "\textit{Heartbeat}" (pulsation) est envoyé toutes les 30 secondes pour empêcher la coupure de la connexion par les \textit{proxies} réseau ou \texttt{Docker}.
    \item \textbf{\textit{Push} de Notification :} Lorsqu'un événement métier survient, le \texttt{NotificationSseService} identifie l'émetteur actif dans une \texttt{ConcurrentHashMap} et "pousse" instantanément l'objet \texttt{JSON}. Le \textit{frontend} intercepte ce flux via un \texttt{Observable RxJS} pour déclencher une alerte visuelle (\textit{Toast}).
\end{itemize}

\subsubsection{Cas n°4 : Gouvernance \& Administration RH}

Ce dernier flux illustre la gestion du cycle de vie des collaborateurs au sein de \textbf{TalentPredict}.

\begin{figure}[H]
    \centering
    \includegraphics[width=0.9\textwidth]{images/seq4Spring4.png}
    \caption{Diagramme de séquence : Administration des comptes et accès}
    \label{fig:seq_cas4}
\end{figure}
\color{red}
\begin{itemize}
    \item \textit{Diagramme 27 :} Corriger le titre annoncé sous la figure : il ne s’agit pas d’un flux de notification temps réel, mais d’un flux de gestion du vivier et de consultation des profils publics.
    \item \textit{Diagramme 27 :} Renommer le titre en \textit{Diagramme de séquence : gestion du vivier et consultation des profils publics}.
    \item \textit{Diagramme 27 :} Remplacer l’acteur \textit{Admin / Responsable RH} par \textit{Administrateur RH} afin d’utiliser la même terminologie dans tout le rapport.
    \item \textit{Diagramme 27 :} Harmoniser la langue des lignes de vie, car le diagramme mélange français et anglais : \textit{Frontend}, \textit{Security Layer}, \textit{Employee Controller}, \textit{Employee Service}.
    \item \textit{Diagramme 27 :} Remplacer \textit{Frontend (Angular 17)} par \textit{Interface Angular} ou garder tous les intitulés techniques en anglais, mais de manière uniforme.
    \item \textit{Diagramme 27 :} Remplacer \textit{Security Layer (RBAC Filter)} par \textit{Filtre de sécurité RBAC} pour améliorer la lisibilité.
    \item \textit{Diagramme 27 :} Remplacer \textit{HR Governance Module} par \textit{Module de gouvernance RH}.
    \item \textit{Diagramme 27 :} Remplacer \textit{Employee Controller} par \textit{Contrôleur des employés} ou conserver \textit{EmployeeController} si le diagramme reste technique.
    \item \textit{Diagramme 27 :} Remplacer \textit{Employee Service} par \textit{Service des employés} ou conserver \textit{EmployeeService} de manière homogène.
    \item \textit{Diagramme 27 :} Ajouter explicitement la validation du jeton JWT avant la vérification du rôle administrateur.
    \item \textit{Diagramme 27 :} Ajouter un fragment \texttt{alt} après la vérification du rôle pour distinguer le cas \textit{accès autorisé} du cas \textit{accès refusé}.
    \item \textit{Diagramme 27 :} Ajouter un retour \textit{401 Unauthorized} si le jeton JWT est invalide ou expiré.
    \item \textit{Diagramme 27 :} Ajouter un retour \textit{403 Forbidden} si le jeton est valide mais que l’utilisateur ne possède pas le rôle administrateur.
    \item \textit{Diagramme 27 :} Remplacer \textit{Vérifier Rôle (ADMIN)} par \textit{Vérifier le rôle administrateur}.
    \item \textit{Diagramme 27 :} Remplacer \textit{findAllEmployees()} par \textit{récupérer la liste des employés} ou harmoniser tous les messages en anglais.
    \item \textit{Diagramme 27 :} Remplacer \textit{SELECT * FROM employees} par \textit{Requête de récupération des employés}, car le détail SQL peut être expliqué dans le texte technique.
    \item \textit{Diagramme 27 :} Remplacer \textit{List<Employee>} par \textit{Liste des employés}.
    \item \textit{Diagramme 27 :} Remplacer \textit{EmployeeSummaryDTOs} par \textit{DTO de synthèse des employés} ou \textit{Liste synthétique des employés}.
    \item \textit{Diagramme 27 :} Remplacer \textit{Sélectionne un profil pour voir les détails} par \textit{Sélectionne un employé pour consulter son profil public}.
    \item \textit{Diagramme 27 :} Remplacer \textit{GET /api/v1/employees/\{id\}/public-profile} par \textit{Demander le profil public de l’employé}.
    \item \textit{Diagramme 27 :} Remplacer \textit{getDetailedPublicProfile(id)} par \textit{récupérer le profil public détaillé}.
    \item \textit{Diagramme 27 :} Remplacer \textit{Fetch SocialLinks, Bio, Performance} par \textit{Récupérer les liens sociaux, la biographie et les indicateurs de performance}.
    \item \textit{Diagramme 27 :} Remplacer \textit{Profile Data} par \textit{Données du profil public}.
    \item \textit{Diagramme 27 :} Remplacer \textit{PublicProfileDTO} par \textit{DTO du profil public}.
    \item \textit{Diagramme 27 :} Dans la phase 2, ajouter une demande de confirmation avant la désactivation du compte si cette confirmation existe dans l’interface.
    \item \textit{Diagramme 27 :} Remplacer \textit{PATCH /api/v1/employees/\{id\}/status?active=false} par \textit{Demander la désactivation du compte employé}.
    \item \textit{Diagramme 27 :} Remplacer \textit{updateEmployeeStatus(id, false)} par \textit{mettre à jour le statut de l’employé}.
    \item \textit{Diagramme 27 :} Remplacer \textit{UPDATE employees SET is\_active = false} par \textit{Mise à jour du statut du compte employé}.
    \item \textit{Diagramme 27 :} Remplacer \textit{Row Updated} par \textit{Statut mis à jour}.
    \item \textit{Diagramme 27 :} Remplacer \textit{StatusUpdateDTO} par \textit{DTO de mise à jour du statut}.
    \item \textit{Diagramme 27 :} Ajouter un fragment \texttt{alt} pour distinguer le cas \textit{mise à jour réussie} du cas \textit{échec de mise à jour}.
    \item \textit{Diagramme 27 :} Ajouter un retour d’erreur si l’identifiant de l’employé est introuvable ou si le compte est déjà désactivé.
    \item \textit{Diagramme 27 :} La note \textit{Invalidation immédiate des droits d’accès futurs} est pertinente ; il faut toutefois préciser dans le texte si les sessions déjà ouvertes sont aussi invalidées.
    \item \textit{Diagramme 27 :} Vérifier si la désactivation doit déclencher une notification à l’employé ; si oui, ajouter une ligne de vie \textit{Service de notification}.
    \item \textit{Diagramme 27 :} Ajouter des barres d’activation cohérentes sur toutes les lignes de vie qui exécutent un traitement, notamment \textit{Employee Controller}, \textit{Employee Service} et \textit{PostgreSQL}.
    \item \textit{Diagramme 27 :} Vérifier que toutes les réponses sont représentées par des flèches de retour en pointillés, notamment les réponses de la base de données et du service métier.
    \item \textit{Diagramme 27 :} Réduire certains libellés techniques trop longs afin d’améliorer la lisibilité dans le rapport.

\end{itemize}
\color{black}
\paragraph*{Flux : Gestion du Vivier \& Profils Publics}
Ce flux décrit la gouvernance administrative et le contrôle rigoureux des comptes utilisateurs.
\begin{itemize}
    \item \textbf{Supervision RH :} L'administrateur peut lister et rechercher des collaborateurs via une interface sécurisée. La consultation d'un profil public déclenche la récupération de données hétérogènes (\textit{bio}, historique de performance, liens \texttt{GitHub}/\texttt{LinkedIn}) pour une évaluation à 360°.
    \item \textbf{Gestion du Cycle de Vie :} Le processus inclut l'action critique de désactivation d'un compte. Cette opération entraîne une mise à jour immédiate en base de données et l'invalidation des droits d'accès dans le \texttt{SecurityContext}, garantissant une réactivité de sécurité immédiate (\textit{Kill Switch}).
\end{itemize}

\subsubsection{7. Diagramme de Conception }
\begin{figure}[H]
    \centering
    \includegraphics[width=\textwidth]{images/classSprint4.png}
    \vspace{2mm}
    \caption{Diagramme de classes technique du système TalentPredict}
    \label{fig:diagramme_classe_technique}
\end{figure}

\color{red}
\begin{itemize}
    \item \textit{Diagramme 28 :} Renommer le titre en \textit{Diagramme de classes technique du système TalentPredict AI} pour garder une formulation claire et académique.
    \item \textit{Diagramme 28 :} Réduire le grand espace vide à gauche de la figure afin d’améliorer l’équilibre visuel et la lisibilité.
    \item \textit{Diagramme 28 :} Réorganiser les blocs pour obtenir une structure plus compacte : modèle de données au centre, modules fonctionnels autour, services externes à droite.
    \item \textit{Diagramme 28 :} Harmoniser la langue du diagramme, car certains intitulés sont en anglais et d’autres en français.
    \item \textit{Diagramme 28 :} Remplacer \textit{Core Data Models} par \textit{Modèle de données central} si le rapport est rédigé en français.
    \item \textit{Diagramme 28 :} Remplacer \textit{Module Dashboard \& Analytics} par \textit{Module tableau de bord et analyses}.
    \item \textit{Diagramme 28 :} Remplacer \textit{Module RH \& Learning} par \textit{Module RH et apprentissage}.
    \item \textit{Diagramme 28 :} Remplacer \textit{Module Reporting} par \textit{Module de reporting} ou \textit{Module d’export des rapports}.
    \item \textit{Diagramme 28 :} Ajouter une légende expliquant les stéréotypes utilisés, notamment \texttt{<<Entity>>}, \texttt{<<Service>>}, \texttt{<<Controller>>}, \texttt{<<DTO>>}, \texttt{<<ValueObject>>} et \texttt{<<System>>}.
    \item \textit{Diagramme 28 :} Clarifier la différence entre les entités persistées et les objets de valeur ; \textit{Skill} et \textit{Prediction} semblent contenir des données métier importantes.
    \item \textit{Diagramme 28 :} Vérifier si \textit{Skill} doit rester un \texttt{<<ValueObject>>} ou devenir une entité persistée, car les compétences sont généralement associées et sauvegardées pour chaque utilisateur.
    \item \textit{Diagramme 28 :} Vérifier si \textit{Prediction} doit rester un \texttt{<<ValueObject>>} ou devenir une entité si les prédictions et scores sont historisés.
    \item \textit{Diagramme 28 :} Ajouter les multiplicités sur les relations entre \textit{User}, \textit{Formation}, \textit{TestResult}, \textit{Skill}, \textit{Prediction} et \textit{Notification}.
    \item \textit{Diagramme 28 :} Corriger la relation entre \textit{User} et \textit{Formation} si plusieurs utilisateurs peuvent suivre plusieurs formations ; dans ce cas, ajouter une classe d’association comme \textit{UserFormation} ou \textit{Enrollment}.
    \item \textit{Diagramme 28 :} Ajouter une classe \textit{TrainingRequest} si le système gère les demandes de formation, leur approbation, leur refus ou leur remplacement.
    \item \textit{Diagramme 28 :} Ajouter une classe \textit{TalentPassport} ou \textit{Report} si les passeports et rapports PDF sont générés, stockés ou téléchargés.
    \item \textit{Diagramme 28 :} Vérifier si les DTOs \textit{DashboardDto} et \textit{AdminOverviewDto} doivent être affichés dans le même diagramme que les entités métier ; sinon, les déplacer vers un diagramme technique séparé.
    \item \textit{Diagramme 28 :} Ajouter les repositories ou interfaces d’accès aux données si l’objectif est de représenter une architecture backend complète.
    \item \textit{Diagramme 28 :} Ajouter les contrôleurs manquants, par exemple \textit{DashboardController}, \textit{FormationController} ou \textit{NotificationController}, si ces endpoints existent réellement dans le backend.
    \item \textit{Diagramme 28 :} Clarifier la relation entre \textit{DashboardService} et le service IA, car le lien actuel n’indique pas clairement quelles données sont envoyées ni quels résultats sont retournés.
    \item \textit{Diagramme 28 :} Renommer \textit{Service IA (FastAPI)} en \textit{Microservice IA}, puis préciser FastAPI dans une note technique.
    \item \textit{Diagramme 28 :} Supprimer les endpoints REST à l’intérieur du bloc \textit{Service IA} si le schéma est présenté comme un diagramme de classes ; ces détails relèvent plutôt d’un diagramme d’architecture ou d’API.
    \item \textit{Diagramme 28 :} Si les endpoints REST sont conservés, les représenter plutôt via une interface \textit{AiMicroserviceClient} appelée par le backend.
    \item \textit{Diagramme 28 :} Vérifier si l’endpoint \textit{predict-mobility} appartient réellement au projet ; sinon, le remplacer par un endpoint cohérent comme \textit{predict-career}, \textit{recommend-training} ou \textit{analyze-candidate}.
    \item \textit{Diagramme 28 :} Vérifier la cohérence de la classe \textit{FormationService}, car la méthode \textit{approveEnrollment} suppose l’existence d’une entité d’inscription ou de demande.
    \item \textit{Diagramme 28 :} Ajouter une classe \textit{PdfExportService} ou renommer \textit{ExportService} si ce service est dédié à la génération des fichiers PDF.
    \item \textit{Diagramme 28 :} Clarifier la relation entre \textit{NotificationService} et l’entité \textit{Notification}, notamment si le service crée, stocke et diffuse les notifications.
    \item \textit{Diagramme 28 :} Ajouter éventuellement une classe \textit{SseEmitterManager} ou \textit{RealTimeNotificationService} si les notifications temps réel sont importantes dans le Sprint 4.
    \item \textit{Diagramme 28 :} Ajouter un champ \textit{type} ou \textit{category} dans \textit{Notification} si plusieurs types de notifications sont gérés.
    \item \textit{Diagramme 28 :} Ajouter les champs \textit{createdAt} et \textit{updatedAt} dans les entités principales pour améliorer la traçabilité.
    \item \textit{Diagramme 28 :} Vérifier si la classe \textit{User} doit contenir uniquement les informations d’identité, tandis que les scores, compétences et prédictions doivent être placés dans des entités séparées.

\end{itemize}
\color{black}
\color{red}
\begin{itemize}
    \item \textit{Diagramme 23 :} Renommer le titre en \textit{Diagramme de classes technique du système TalentPredict AI} pour garder une formulation plus académique.
    \item \textit{Diagramme 23 :} Réduire le vide important à gauche de la figure, car il diminue la lisibilité et rend le diagramme difficile à exploiter dans le rapport.
    \item \textit{Diagramme 23 :} Réorganiser les modules pour obtenir une structure plus compacte : données centrales au centre, modules fonctionnels autour, services externes à droite.
    \item \textit{Diagramme 23 :} Remplacer \textit{Core Data Models} par \textit{Modèle de données central} ou harmoniser tous les intitulés en anglais.
    \item \textit{Diagramme 23 :} Remplacer \textit{Module Dashboard \& Analytics} par \textit{Module tableau de bord et analyses} si le rapport est rédigé en français.
    \item \textit{Diagramme 23 :} Remplacer \textit{Module Reporting} par \textit{Module de reporting} ou \textit{Module d’export des rapports}.
    \item \textit{Diagramme 23 :} Remplacer \textit{Module RH \& Learning} par \textit{Module RH et apprentissage}.
    \item \textit{Diagramme 23 :} Remplacer \textit{Services Externes} par \textit{Services externes}, avec une casse homogène.
    \item \textit{Diagramme 23 :} Éviter le mélange entre français et anglais dans les relations, par exemple \textit{suit}, \textit{réalise}, \textit{reçoit}, \textit{produit}, et les noms de méthodes en anglais.
    \item \textit{Diagramme 23 :} Ajouter une légende expliquant les stéréotypes utilisés, notamment \texttt{<<Entity>>}, \texttt{<<Service>>}, \texttt{<<Controller>>}, \texttt{<<DTO>>}, \texttt{<<ValueObject>>} et \texttt{<<System>>}.
    \item \textit{Diagramme 23 :} Clarifier la différence entre les entités persistées et les objets de valeur ; \textit{Skill} et \textit{Prediction} sont marqués comme \texttt{<<ValueObject>>}, mais ils semblent contenir des informations métier importantes qui peuvent être persistées.
    \item \textit{Diagramme 23 :} Vérifier si \textit{Skill} doit être une entité indépendante, car les compétences sont généralement stockées et associées à l’utilisateur.
    \item \textit{Diagramme 23 :} Vérifier si \textit{Prediction} doit être une entité, surtout si les scores, le type de carrière et les résultats IA sont historisés.
    \item \textit{Diagramme 23 :} Ajouter les multiplicités sur toutes les relations importantes, notamment entre \textit{User}, \textit{Skill}, \textit{Prediction}, \textit{TestResult}, \textit{Formation} et \textit{Notification}.
    \item \textit{Diagramme 23 :} Corriger la relation entre \textit{User} et \textit{Formation} : si plusieurs utilisateurs peuvent suivre plusieurs formations, il faut ajouter une entité d’association, par exemple \textit{UserFormation} ou \textit{Enrollment}.
    \item \textit{Diagramme 23 :} Ajouter une classe ou entité \textit{TrainingRequest} si le système gère l’approbation, le refus ou la suggestion de formations par l’administrateur RH.
    \item \textit{Diagramme 23 :} Ajouter une classe \textit{Report} ou \textit{TalentPassport} si les rapports et passeports PDF sont générés, stockés ou téléchargés.
    \item \textit{Diagramme 23 :} Vérifier si \textit{DashboardDto} et \textit{AdminOverviewDto} doivent apparaître dans un diagramme de classes global ; si le diagramme vise le modèle métier, ces DTOs peuvent être déplacés vers un diagramme technique séparé.
    \item \textit{Diagramme 23 :} Ajouter les contrôleurs manquants si le diagramme se veut technique complet, par exemple \textit{DashboardController}, \textit{NotificationController} ou \textit{FormationController}.
    \item \textit{Diagramme 23 :} Ajouter les repositories ou ports d’accès aux données si l’objectif est de représenter une architecture backend complète.
    \item \textit{Diagramme 23 :} Clarifier la relation entre \textit{DashboardService} et \textit{Service IA} ; actuellement, le lien passe indirectement par plusieurs éléments et son rôle n’est pas immédiatement lisible.
    \item \textit{Diagramme 23 :} Renommer \textit{Service IA (FastAPI)} en \textit{Microservice IA} ou \textit{Service IA}, puis préciser FastAPI dans une note technique.
    \item \textit{Diagramme 23 :} Supprimer les endpoints REST à l’intérieur de la classe \textit{Service IA} si le diagramme doit rester un diagramme de classes ; ces endpoints relèvent plutôt d’un diagramme d’architecture ou d’API.
    \item \textit{Diagramme 23 :} Si les endpoints REST sont conservés, les présenter dans une interface séparée, par exemple \textit{AiMicroserviceClient}, afin de distinguer le service externe de son client côté backend.
    \item \textit{Diagramme 23 :} Remplacer la note \textit{Microservice Python pour le traitement du langage naturel et la simulation} par une formulation plus précise : \textit{Microservice Python dédié à l’analyse IA, aux prédictions et aux simulations}.
    \item \textit{Diagramme 23 :} Vérifier si \textit{predict-mobility} appartient réellement au périmètre de TalentPredict AI ; sinon, remplacer par un endpoint cohérent avec le projet, comme \textit{predict-career} ou \textit{recommend-training}.
    \item \textit{Diagramme 23 :} Harmoniser les noms des méthodes : utiliser soit l’anglais technique partout, soit le français, mais éviter un mélange dans le même diagramme.
    \item \textit{Diagramme 23 :} Ajouter les types de retour des méthodes principales lorsque cela est utile, mais éviter de surcharger le diagramme avec trop de détails techniques.
    \item \textit{Diagramme 23 :} Vérifier la classe \textit{FormationService} : la méthode \textit{approveEnrollment} semble liée à une demande ou inscription à une formation, ce qui nécessite une classe métier dédiée.
    \item \textit{Diagramme 23 :} Vérifier la cohérence de \textit{ReportingController} : les méthodes d’export PDF doivent être reliées clairement à \textit{ExportService}.
    \item \textit{Diagramme 23 :} Ajouter une classe \textit{PdfExportService} ou renommer \textit{ExportService} de manière plus explicite si ce service génère uniquement des PDF.
    \item \textit{Diagramme 23 :} Clarifier la relation entre \textit{NotificationService} et l’entité \textit{Notification} ; si le service crée et envoie des notifications, la relation doit être nommée plus précisément.
    \item \textit{Diagramme 23 :} Ajouter éventuellement une classe \textit{SseEmitterManager} ou \textit{RealTimeNotificationService} si la communication temps réel par SSE est un élément important du Sprint 4.
    \item \textit{Diagramme 23 :} Vérifier que \textit{Notification} contient un champ \textit{type} ou \textit{category} si plusieurs types de notifications sont gérés.
    \item \textit{Diagramme 23 :} Ajouter les champs \textit{createdAt} et \textit{updatedAt} dans les entités principales si la traçabilité est nécessaire.
    \item \textit{Diagramme 23 :} Vérifier si la classe \textit{User} doit contenir uniquement les informations d’identité, tandis que les scores, compétences et prédictions doivent être placés dans des entités séparées.
    \item \textit{Diagramme 23 :} Ajouter une relation claire entre \textit{User} et \textit{DashboardDto}, ou éviter de représenter les DTOs comme des éléments centraux du modèle métier.
    \item \textit{Diagramme 23 :} Réduire les longues lignes diagonales, car elles rendent les dépendances difficiles à suivre.

\end{itemize}
\color{black}
Le diagramme de conception illustre une architecture logicielle robuste et modulaire, structurée selon le modèle \textit{N-Tier} (Multicouche). Cette organisation garantit une séparation nette des responsabilités entre la gestion des données, la logique métier et les interfaces de communication.

\paragraph{ Cœur du Système (\textit{Core Data Models})}
Au centre de l'application se trouve l'entité \texttt{User}, pivot central autour duquel gravitent les données d'évaluation. On observe une relation de composition avec les résultats de tests (\texttt{TestResult}) et les compétences (\texttt{Skill}). L'innovation réside dans l'intégration de l'objet \texttt{Prediction}, qui stocke les résultats issus de l'intelligence artificielle (\textit{Persona} PCM, \textit{Career Path}), permettant ainsi un accès rapide aux \textit{insights} sans recalcul systématique.

\paragraph{ Module \textit{Dashboard \& Analytics} (Agrégation de Données)}
Ce module est le cerveau visuel de l'application. Le \texttt{DashboardService} joue un rôle d'orchestrateur : il récupère les informations brutes des différents modules pour les transformer en DTO (\textit{Data Transfer Objects}).
\begin{itemize}
    \item Le \texttt{DashboardDto} est optimisé pour la vue "\textit{Bento}" du collaborateur.
    \item L'\texttt{AdminOverviewDto} agrège les données de l'ensemble du vivier pour offrir une vision macroscopique à la direction.
\end{itemize}

\paragraph{Module \textit{Reporting} (Moteur d'Exportation)}
Le système de \textit{reporting} est conçu pour être indépendant de la logique de calcul. Le \texttt{ReportingController} sollicite l'\texttt{ExportService} qui, à son tour, utilise les données consolidées du \textit{Dashboard} pour générer des documents PDF fidèles et certifiés via un moteur de rendu HTML vers PDF.

\paragraph{Intégration Micro-services (Services Externes)}
L'architecture se distingue par son ouverture vers un micro-service externe développé en \texttt{Python} (\texttt{FastAPI}). Ce service est spécialisé dans le traitement du langage naturel (NLP) et les simulations prédictives basées sur des modèles \texttt{Llama3} / \texttt{Ollama}. La communication est assurée par des appels REST sécurisés, permettant une évolution indépendante des moteurs d'IA.

\paragraph{Modules de Soutien (Communication \& Learning)}
\begin{itemize}
    \item \textbf{Module Communication :} Gère le cycle de vie des notifications. Le service dédié assure la persistance en base de données et la diffusion en temps réel.
    \item \textbf{Module RH \& Learning :} Assure la gestion administrative du vivier (\texttt{AdminController}) et la recommandation intelligente de formations (\texttt{FormationService}), créant ainsi un écosystème d'\textit{upskilling} complet.
\end{itemize}

\subsubsection*{8. Architecture Temps Réel via \texttt{SSE} (\textit{Server-Sent Events})}
Afin de garantir une communication fluide et instantanée, \textbf{TalentPredict} intègre une architecture temps réel basée sur le protocole \texttt{SSE}. Ce choix technologique, plus léger et robuste que les \texttt{WebSockets} pour un flux unidirectionnel (\textit{Serveur vers Client}), est parfaitement adapté à la diffusion de notifications et d'alertes sans rechargement de page.

Le \textit{backend} \texttt{Spring Boot} agit comme un émetteur d’événements grâce à l’\texttt{API} \texttt{SseEmitter}. Après authentification, les utilisateurs établissent une connexion persistante pour recevoir des mises à jour asynchrones.

\paragraph*{Composants de l'architecture :}
L’architecture repose sur les modules suivants :
\begin{itemize}
    \item \textbf{\texttt{NotificationSseService} :} Gère le cycle de vie des connexions (\texttt{SseEmitter}) et le registre des utilisateurs connectés via une structure \texttt{ConcurrentHashMap} sécurisée.
    \item \textbf{\texttt{NotificationController} :} Expose l’\textit{endpoint} de souscription \texttt{api}/\texttt{notifications}/\texttt{subscribe} permettant au client d'ouvrir le canal de communication.
    \item \textbf{Système Événementiel :} Centralise les notifications métiers (nouvelle recommandation IA, succès d’un test, validation RH, génération du \textit{Talent Passport}).
    \item \textbf{Angular \textit{EventSource Client} :} Utilise l’\texttt{API} native du navigateur pour maintenir la connexion ouverte et rafraîchir dynamiquement les composants \texttt{Angular} via \texttt{RxJS} à la réception d’un événement.
\end{itemize}

\begin{figure}[H]
    \centering
    % Remplacez 'nom_du_fichier' par le nom réel de votre image (ex: flux_sse.png)
    \includegraphics[width=0.85\textwidth]{images/DiagrammedeComposants.png}
    \caption{Architecture événementielle et gestion du registre des connexions  \texttt{SSE}}
    \label{fig:flux_communication_sse}
\end{figure}

Le schéma ci-dessus illustre le flux de communication temps réel au sein de \textbf{TalentPredict}. Le \texttt{NotificationController} réceptionne les demandes de souscription et délègue au \texttt{NotificationSseService} la création d'un canal persistant. 

Ce dernier stocke l'instance de la connexion dans une \texttt{ConcurrentHashMap} (mémoire vive), garantissant un accès \textit{thread-safe} et performant. Lors du déclenchement d'un événement métier, le système identifie le destinataire et "pousse" l'information instantanément vers le client \texttt{Angular}.

\paragraph*{Avantages techniques retenus :}
\begin{itemize}
    \item \textbf{Reconnexion Automatique :} Le protocole \texttt{SSE} gère nativement la reprise de connexion en cas de coupure réseau.
    \item \textbf{Efficacité énergétique :} Consomme moins de ressources côté client et serveur qu’une connexion \texttt{WebSocket} bidirectionnelle.
    \item \textbf{Compatibilité Proxy/Firewall :} Contrairement aux \texttt{WebSockets}, le \texttt{SSE} utilise le protocole \texttt{HTTP} standard, évitant ainsi les problèmes de blocage par les pare-feux d'entreprise ou les \textit{proxys} \texttt{Docker}.
\end{itemize}






\subsection{ Codage}
% Adaptez le niveau (section, subsection, subsubsection) selon votre plan.

Cette sous-section détaille l'étape de codage spécifique au quatrième sprint. L'infrastructure globale de l'application étant déjà en place, nos efforts de développement se sont concentrés sur l'intégration de bibliothèques spécialisées pour le rendu documentaire et la communication asynchrone.

\paragraph*{1. Implémentation du Moteur de Reporting (Backend)}
Pour répondre au besoin de dématérialisation (\textit{Passeport Talent} et Rapports RH), nous avons mis en place un service de génération de documents côté serveur (\textit{Server-Side Rendering}) :

\begin{itemize}
    \item \textbf{\texttt{OpenHTMLToPDF} (v1.0.10) :} Il s'agit de la bibliothèque centrale de ce sprint. Elle a été intégrée pour \textit{parser} nos gabarits dynamiques (HTML5/CSS3) et les transformer en documents PDF professionnels. Ce choix a été motivé par sa grande fidélité de rendu des standards Web récents.
    \item \textbf{\texttt{Apache PDFBox} :} Utilisée en complément pour la manipulation de bas niveau des flux générés (gestion des métadonnées du fichier et sécurisation du PDF).
    \item \textbf{\texttt{Lombok} :} Déployée systématiquement sur nos nouveaux objets de transfert (DTOs) de \textit{reporting} pour réduire le code \textit{boilerplate} et garantir la clarté des classes métiers.
    \item \textbf{Thymeleaf :} Choisi comme moteur de gabarits (Templating Engine) côté serveur pour injecter dynamiquement les métadonnées (nom de l'employé, scores IA, radar) dans les templates HTML avant leur transformation finale en documents PDF.
\end{itemize}
\subsubsection*{2. Génération Dynamique des Rapports PDF}

La génération documentaire constitue l’une des fonctionnalités majeures du Sprint 4. Le \textit{backend} \texttt{Spring Boot} génère dynamiquement des documents PDF professionnels à partir des données analytiques produites par les modules IA.

Deux types de documents clés ont été implémentés :
\begin{itemize}
    \item \textbf{\textit{Talent Passport} :} Une synthèse individuelle exhaustive regroupant les compétences extraites, les scores techniques, le profil comportemental PCM et les recommandations de carrière de l'IA.
    \item \textbf{Rapport RH Global :} Une vue consolidée destinée à la direction, présentant les statistiques et les tendances macroscopiques du capital humain.
\end{itemize}

\paragraph*{Mécanisme de production et persistance}
Le moteur de génération repose sur une architecture orientée \textit{Templates} (\texttt{Thymeleaf}) permettant de produire des documents cohérents, structurés et respectant la charte graphique de la plateforme. Une attention particulière a été portée à la persistance : les fichiers générés sont automatiquement stockés dans un volume \texttt{Docker} dédié (\texttt{app\_uploads}), assurant leur disponibilité et leur archivage même après le redémarrage des services applicatifs.


\subsubsection*{3. Hub de Communication et Tableaux de Bord }
Pour transformer la plateforme en un outil de pilotage interactif, nous avons implémenté des mécanismes de visualisation et de mise à jour en temps réel :

\begin{itemize}
    \item \textbf{\texttt{Spring SseEmitter} \& \texttt{RxJS} :} Ce duo technique a permis l'intégration du protocole SSE (Server-Sent Events). Côté serveur, \texttt{Spring Boot} gère la diffusion des messages. Côté client (\texttt{Angular}), \texttt{RxJS} orchestre les flux de données asynchrones, permettant de "pousser" les notifications aux utilisateurs sans aucun rafraîchissement manuel de l'interface.
    \item \textbf{\texttt{Ngx-Toastr} :} Consomme les flux SSE pour afficher des alertes visuelles élégantes et non intrusives (ex. "Nouvelle formation approuvée").
    \item \textbf{\texttt{Chart.js} :} Bibliothèque de \textit{data-visualisation} intégrée pour générer dynamiquement les radars de compétences de l'employé et les graphiques de tendances macroscopiques de l'\textit{Executive Dashboard}.
    \item \textbf{\texttt{Lucide-Angular} \& \texttt{FontAwesome} :} Packs vectoriels utilisés pour standardiser l'iconographie des nouveaux modules (boutons d'export, indicateurs d'état).
\end{itemize}

\subsubsection*{4. Outillage DevOps et Validation API}
Afin de garantir la stabilité de ces nouvelles fonctionnalités complexes avant leur déploiement, nous avons maintenu une rigueur stricte sur l'outillage :

\begin{itemize}
    \item \textbf{\texttt{Postman} :} Utilisé pour documenter et tester unitairement les nouveaux points d'accès REST complexes, notamment les flux binaires (téléchargement de PDF) et les retours de l'IA, avant leur intégration dans le \textit{frontend}.
    \item \textbf{Communication Inter-services :} Mise en œuvre de clients REST robustes (Spring RestTemplate) configurés avec des mécanismes de Timeout pour assurer une communication fluide et résiliente entre le Backend Java et le micro-service d'Intelligence Artificielle (FastAPI).
    \item \textbf{\texttt{Docker} \& \texttt{Docker Compose} :}la configuration du flux persistant pour le SSE a nécessité la mise à jour de nos fichiers de conteneurisation pour garantir un déploiement fluide de ce nouvel écosystème.
    \item \textbf{\texttt{GitLab CI/CD} :} Le pipeline d'intégration continue a été maintenu pour exécuter automatiquement les tests sur ces nouveaux modules à chaque validation de code (\textit{Commit}).
\end{itemize}
\subsection{Realisation}
% Ajustez le numéro de section au besoin (ex: remplacez \subsection par \section si c'est le titre principal 4.X)

L'aboutissement de ce quatrième sprint se matérialise par la livraison d'interfaces interactives et de modules de restitution documentaire. Le \textit{Frontend}, développé en \texttt{Angular 17}, communique de manière asynchrone avec nos différents microservices pour offrir une expérience utilisateur fluide et en temps réel.

\subsubsection*{1. Pilotage Stratégique : Executive Dashboard et Gestion RH}
Le cœur de ce sprint consistait à fournir aux administrateurs et décideurs RH une vue macroscopique et exploitable du capital humain de l'entreprise.


    \subsubsection*{ L'\textit{Executive Dashboard} (Vue d'ensemble) :} Cette interface centralise les indicateurs de performance (KPIs) en temps réel. Grâce à l'intégration de la bibliothèque \texttt{Chart.js}, l'administrateur peut visualiser les tendances d'évolution des compétences, le suivi de l'\textit{Onboarding}, ainsi que les alertes critiques (ex. : taux de couverture des tests insuffisant). L'IA intervient également ici pour mettre en évidence les "Talents à haut potentiel" (HiPo).


\begin{figure}[H]
    \centering
    \includegraphics[width=0.9\textwidth]{images/dashboardAdmin.png} % Remplacez par le nom de l'Image 1
    \vspace{2mm}
    \caption{\textit{Executive Dashboard} : Vue macroscopique et indicateurs clés}
    \label{fig:executive_dashboard}
\end{figure}

    \subsubsection*{Gestion des Employés et Profil Public :} Ce module permet une administration complète (CRUD) des collaborateurs. La liste interactive intègre des filtres dynamiques (statuts, rôles, départements). Un panneau latéral (\textit{Slide-over}) offre un aperçu rapide du profil sélectionné, affichant son volume de formations, ses prédictions IA et permettant de vérifier la validité de ses liens sociaux (LinkedIn, GitHub) sans quitter la page principale.


% --- Figure 1 : Aperçu du profil public ---
\begin{figure}[H]
    \centering
    \includegraphics[width=1\textwidth]{images/publicprofileUser.png}
    \caption{Interface de l'aperçu rapide du profil collaborateur}
    \label{fig:profil_public}
\end{figure}

% --- Figure 2 : Gestion des collaborateurs ---
\begin{figure}[H]
    \centering
    \includegraphics[width=1\textwidth]{images/gestiondeEmployee.png}
    \caption{Interface de gestion globale des effectifs (Vue Administrateur)}
    \label{fig:gestion_employes}
\end{figure}

\subsubsection*{2. Espace Collaborateur et Profiling IA}
Du côté de l'employé, l'interface a été enrichie pour restituer les analyses prédictives générées par notre moteur \texttt{FastAPI}.

    \subsubsection*{Tableau de bord de l'Employé :} Le collaborateur accède à son espace personnel ("\textit{Welcome back}") où figure son "\textit{AI Persona}" (ex. : Profil Stratégique). L'interface affiche une matrice de distribution des compétences (\textit{Soft \& Compétences techniques}), son inertie d'apprentissage (\textit{Momentum}), et surtout, les parcours de carrière recommandés (\textit{Career Pathways}) avec un pourcentage de correspondance calculé par l'intelligence artificielle.


\begin{figure}[H]
    \centering
    \includegraphics[width=0.9\textwidth]{images/dashboardUser.png} % Remplacez par le nom de l'Image 3
    \vspace{2mm}
    \caption{Tableau de bord collaborateur avec restitution des analyses IA}
    \label{fig:dashboard_employe}
\end{figure}

\subsubsection*{3. Interaction Temps Réel : Communication Hub et Notifications}
Pour transformer la plateforme en un écosystème vivant, nous avons implémenté des mécanismes de communication instantanée basés sur les \texttt{SEE}.


    \subsubsection*{Le \textit{Hub} de Communication :} Réservé aux administrateurs, cet espace permet de composer et d'envoyer des messages ciblés (Informations, Succès, Alertes) à des collaborateurs spécifiques, avec une option de relai par e-mail.

\begin{figure}[H]
    \centering
    \includegraphics[width=0.85\textwidth]{images/communicationHUB.png}
    \caption{Interface du Hub de communication pour l'envoi de messages}
    \label{fig:communication_hub}
\end{figure}
    \subsubsection*{Centre de Notifications (SSE) :} Côté utilisateur, un volet de notifications interactif capte les événements poussés par le serveur via le protocole \texttt{SSE} (ex. : "Demande de formation approuvée", "Nouveau test assigné"). L'état non lu/lu est géré dynamiquement sans aucun rafraîchissement de la page.


\begin{figure}[H]
    \centering
    \includegraphics[width=0.6\textwidth]{images/NoificationUser.png}
    \caption{Volet de notifications en temps réel pour l'utilisateur}
    \label{fig:notifications_user}
\end{figure}

\subsubsection*{4. Restitution Documentaire : Le Moteur de Reporting (PDF)}
La dématérialisation et l'exportation sécurisée des données ont été assurées par la bibliothèque \texttt{OpenHTMLToPDF} côté serveur.

    \subsubsection*{Le Passeport Talent:} Destiné à l'employé, ce document généré à la volée compile ses scores, son analyse IA détaillée, les recommandations de développement (\textit{Soft \& Tech}) et son historique de formation. La mise en page respecte une charte graphique professionnelle, idéale pour des entretiens annuels.

    \begin{figure}[H]
    \centering
    \includegraphics[width=0.85\textwidth]{images/talentPassport.png}
    \caption{Passeport Talent : Document PDF généré dynamiquement pour l'employé}
    \label{fig:talent_passport}
    \end{figure}
    \subsubsection*{Le Rapport Analytique RH :} Export destiné à la direction, synthétisant les données globales du personnel (département, nombre de tests, statut d'activité) sous forme de tableau paginé.


-
\begin{figure}[H]
    \centering
    \includegraphics[width=0.65\textwidth]{images/rapportRH.png}
    \caption{Rapport RH : Synthèse décisionnelle générée par le backend}
    \label{fig:rapport_rh}
\end{figure}
\subsection{Test et Déploiement }
% Ajustez la numérotation si nécessaire

Cette phase ultime constitue le point d'orgue de notre cycle de développement. Elle garantit que la plateforme \textbf{TalentPredict} est non seulement fonctionnelle d'un point de vue métier, mais aussi stable, sécurisée et prête à être exploitée dans un environnement de production à haute disponibilité.

\subsubsection*{1. Stratégie de Test et Assurance Qualité (QA)}

Afin de garantir une qualité logicielle irréprochable et de prévenir les régressions lors des mises à jour, nous avons mis en place une pyramide de tests rigoureuse :

\begin{itemize}
    \item \textbf{Tests Unitaires (Backend) :} En nous appuyant sur \texttt{JUnit 5} et \texttt{Mockito}, nous avons isolé et testé les composants critiques. Un effort particulier a été porté sur l'\texttt{ExportService} (vérification de la bonne agrégation des données avant injection dans le PDF) et le \texttt{DashboardService} (validation algorithmique du calcul des indicateurs globaux).
    \item \textbf{Tests d'Intégration et Validation API :} Nous avons utilisé \texttt{Postman} pour valider les contrats d'interface. Nous avons systématiquement testé la communication inter-services, notamment les requêtes asynchrones entre le \textit{Backend} \texttt{Spring Boot} et le moteur IA \texttt{FastAPI}, en simulant des scénarios de réponses complexes.
    \item \textbf{Tests de Recette Fonctionnelle et Temps Réel :} Les \textit{User Stories} ont fait l'objet d'une validation manuelle de bout en bout. Nous avons soumis le \textit{Communication Hub} (SSE) à des tests de charge simulée pour vérifier qu'aucune perte de notification n'intervenait lors de sessions concurrentes.
\end{itemize}

Afin de matérialiser notre démarche d'assurance qualité, le tableau ci-dessous présente un extrait des scénarios de tests majeurs exécutés lors de la phase de validation de ce quatrième sprint :

\begin{table}[H]
    \centering
    \renewcommand{\arraystretch}{1}
    \small
    \begin{tabularx}{\textwidth}{|c|>{\raggedright\arraybackslash}p{3cm}|>{\raggedright\arraybackslash}X|>{\raggedright\arraybackslash}X|}
        \hline
        \textbf{ID} & \textbf{Scénario Testé} & \textbf{Pré-conditions \& Étapes} & \textbf{Résultat Attendu} \\
        \hline
        
        % TC-01
        TC-01 & \textbf{Génération PDF (Passeport Talent)} \newline \textit{(Test Unitaire / Backend)} 
        & \textbf{Pré-requis :} Employé avec des compétences évaluées en base. \newline
          \textbf{Action :} Appel de la méthode \texttt{generatePassport(userId)} de l'\texttt{ExportService}.
        & Génération réussie du flux binaire (\texttt{byte[]}). Le document s'ouvre sans erreur et les scores correspondent aux données de la base PostgreSQL. \\
        \hline
        
        % TC-02
        TC-02 & \textbf{Prédictions IA (API FastAPI)} \newline \textit{(Test d'Intégration)} 
        & \textbf{Pré-requis :} Microservice IA démarré. \newline
          \textbf{Action :} Envoi d'une requête \texttt{POST /predict/persona} avec un \textit{payload} JSON depuis le \textit{backend}.
        & Retour HTTP 200 (OK). Le \textit{payload} de réponse contient un "Persona" valide et des scores prédictifs correctement formatés. \\
        \hline
        
        % TC-03
        TC-03 & \textbf{Notification Temps Réel (SSE)} \newline \textit{(Test Fonctionnel de bout en bout)} 
        & \textbf{Pré-requis :} Employé et Manager RH connectés . \newline
          \textbf{Action :} Le manager clique sur "Approuver la formation".
        & L'employé reçoit instantanément une notification visuelle sur son interface Angular, sans aucun rafraîchissement manuel. \\
        \hline
        
        % TC-04
        TC-04 & \textbf{Restitution Visuelle (\textit{Executive Dashboard})} \newline \textit{(Test de Recette UI)} 
        & \textbf{Pré-requis :} Authentification réussie avec le rôle ADMIN. \newline
          \textbf{Action :} Navigation vers la vue macroscopique RH.
        & Les graphiques (\texttt{Chart.js}) se chargent sans erreur console. Les KPI globaux correspondent aux agrégations de la base de données. \\
        \hline
        
        % TC-05
        TC-05 & \textbf{Contrôle d'Accès (RBAC)} \newline \textit{(Test de Sécurité)} 
        & \textbf{Pré-requis :} Utilisateur authentifié avec le rôle EMPLOYEE. \newline
          \textbf{Action :} Tentative de téléchargement du rapport global RH (\texttt{GET /api/reporting/hr-global}).
        & Le filtre JWT intercepte l'appel. La requête est rejetée avec un code HTTP 403 (Forbidden). L'accès est bloqué. \\
        \hline
        
    \end{tabularx}
    \vspace{2mm}
    \caption{Synthèse des cas de test majeurs (Sprint 4)}
    \label{tab:test_cases_sprint4}
\end{table}

\subsubsection*{2.Tests des Flux Temps Réel (SSE)}

Afin de garantir la fiabilité des communications asynchrones au sein de \textbf{TalentPredict}, plusieurs scénarios de validation ont été réalisés sur l’architecture temps réel basée sur le protocole \texttt{SSE} (\textit{Server-Sent Events}).

Ces tests ont permis de vérifier plusieurs aspects critiques du système :

\begin{itemize}
    \item \textbf{Réception instantanée :} Les événements générés côté serveur, tels que les recommandations IA ou les validations RH, sont affichés sur le \textit{Dashboard} utilisateur avec une latence moyenne inférieure à 200 ms.
    
    \item \textbf{Synchronisation multi-session :} Un même utilisateur connecté simultanément sur plusieurs onglets ou appareils reçoit les notifications de manière cohérente et synchronisée en temps réel.
    
    \item \textbf{Résilience réseau :} Des tests ont été effectués afin d’évaluer la capacité du client \texttt{Angular} à rétablir automatiquement la connexion \texttt{SSE} après une interruption volontaire du réseau. Les résultats ont confirmé la continuité du flux événementiel sans perte de données ni désynchronisation des interfaces.
    
    \item \textbf{Sécurité et Confidentialité :} Des validations ont été menées afin de garantir que chaque flux \texttt{SSE} reste strictement associé à l’identité authentifiée de l’utilisateur via \texttt{JWT}. Cette approche empêche tout accès non autorisé aux notifications privées et renforce la confidentialité des échanges temps réel.
\end{itemize}

L’ensemble des expérimentations réalisées a confirmé la stabilité, la légèreté et la robustesse de l’architecture \texttt{SSE}, même dans des scénarios simulant plusieurs dizaines de connexions simultanées.

\subsubsection*{3. Architecture de Conteneurisation (Docker)}

Pour résoudre les problématiques de disparité d'environnements et faciliter le déploiement, nous avons intégralement conteneurisé l'écosystème :

\begin{itemize}
    \item \textbf{Dockerisation \textit{Multi-Stage} :} L'écriture de fichiers \texttt{Dockerfile} optimisés a permis de séparer les environnements de compilation des environnements d'exécution, réduisant drastiquement la taille des images finales.
    \item \textbf{Orchestration via Docker Compose :} La configuration du fichier \texttt{docker-compose.yml} orchestre le démarrage simultané de nos quatre micro-environnements au sein d'un réseau privé isolé (\texttt{talentpredict-net}) :
    \begin{itemize}
        \item \texttt{tp-frontend} : Application \texttt{Angular} servie par \texttt{Nginx}.
        \item \texttt{tp-backend} : Serveur \texttt{Spring Boot} exposant l'API REST.
        \item \texttt{tp-ai-service} : Moteur analytique \texttt{FastAPI}.
        \item \texttt{tp-database} : Instance \texttt{PostgreSQL} couplée à des volumes persistants (\texttt{tp-db-data}).
    \end{itemize}
\end{itemize}
\subsubsection*{4.Architecture Globale de Déploiement}

\begin{figure}[H]
    \centering
    % Remplacez 'chemin_architecture_globale.png' par le nom exact de votre image
    \includegraphics[width=0.95\textwidth]{images/ArchitectureGlobalede Déploiement.png}
    \vspace{2mm}
    \caption{Architecture globale de déploiement conteneurisé de la plateforme TalentPredict}
    \label{fig:archi_deploiement_globale}
\end{figure}

L'infrastructure de \textit{TalentPredict} est entièrement conteneurisée via \textbf{Docker Compose}, garantissant isolation et portabilité. L'écosystème repose sur un réseau privé \texttt{talentpredict-net} orchestrant cinq services clés :
\begin{itemize}
    \item \textbf{Frontend} : Application Angular servie par \textbf{Nginx} (Port 80).
    \item \textbf{Backend} : Serveur \textbf{Spring Boot 3.4} gérant la logique métier et les flux \textbf{SSE}.
    \item \textbf{IA Service} : Micro-service \textbf{FastAPI} pilotant l'intelligence artificielle.
    \item \textbf{Ollama Engine} : Inférence locale du modèle \textbf{Llama 3.2} pour la souveraineté des données.
    \item \textbf{Base de données} : Instance \textbf{PostgreSQL 16} avec volumes persistants (\texttt{pg\_data}).
\end{itemize}
Cette architecture modulaire assure une maintenance simplifiée et une \textbf{scalabilité horizontale} optimale, permettant de faire évoluer chaque composant indépendamment selon les besoins de charge.

\subsubsection*{5. Persistance des Données avec Docker Volumes}

Afin de garantir la conservation durable des données applicatives indépendamment du cycle de vie des conteneurs \texttt{Docker}, l’infrastructure de \textbf{TalentPredict} repose sur l’utilisation de \textit{Docker Volumes} dédiés.

Deux volumes principaux ont été configurés dans l’architecture de déploiement :

\begin{itemize}
    \item \textbf{\texttt{pg\_data}} : Volume associé au conteneur \texttt{PostgreSQL} permettant de persister l’ensemble des données relationnelles de la plateforme, notamment les comptes utilisateurs, les profils collaborateurs, les résultats d’évaluation, les notifications et les données analytiques RH.
    
    \item \textbf{\texttt{app\_uploads}} : Volume utilisé pour le stockage des fichiers générés ou téléversés par les utilisateurs, tels que les CV, les photos de profil, les documents PDF et les certificats produits automatiquement par la plateforme.
\end{itemize}

Cette approche permet de dissocier les données persistantes des conteneurs applicatifs eux-mêmes. Ainsi, même en cas de suppression, redémarrage ou mise à jour des conteneurs \texttt{Docker}, les données restent conservées et réutilisables lors du redéploiement du système.

L’utilisation des \textit{Docker Volumes} améliore également :
\begin{itemize}
    \item la fiabilité de l’infrastructure ;
    \item la sécurité des données ;
    \item la maintenabilité des environnements ;
    \item ainsi que la préparation à une future migration vers une infrastructure \textit{Cloud} ou \texttt{Kubernetes}.
\end{itemize}

\subsubsection{6.Stratégie de Supervision et Observabilité}

La pérennité de \textbf{TalentPredict} repose sur une surveillance proactive de ses composants conteneurisés. Pour assurer cette stabilité, nous avons déployé un dispositif d'observabilité articulé autour de deux axes : les métriques applicatives et la santé des conteneurs.

\paragraph*{Mise en œuvre du \textit{Monitoring} Applicatif}
Grâce à l’intégration de \texttt{Spring Boot Actuator}, le \textit{backend} de l'application devient "transparent". Il expose en continu des indicateurs critiques tels que l’utilisation de la mémoire \texttt{JVM}, la charge \texttt{CPU} et la disponibilité des ressources de persistance. Ces données permettent d'anticiper d'éventuelles dégradations de performance et facilitent le diagnostic technique à travers une centralisation des \textit{logs} applicatifs.

\paragraph*{Résilience et Auto-réparation \texttt{Docker}}
La robustesse de l'infrastructure est garantie par des mécanismes de \texttt{Healthcheck} intégrés à l'orchestration \texttt{Docker}. Chaque service (\textit{Backend}, IA, \texttt{PostgreSQL}) est soumis à des tests de santé périodiques. 

En cas d’anomalie ou d’indisponibilité d’un composant, \texttt{Docker} déclenche automatiquement un cycle de redémarrage. Cette capacité d’auto-réparation assure une continuité de service optimale et réduit drastiquement le temps d’indisponibilité (\textit{Downtime}) sans intervention manuelle, renforçant ainsi la fiabilité globale de la plateforme.


\subsubsection*{7. Pipeline CI/CD (Intégration et Déploiement Continus)}

Pour \textbf{TalentPredict}, nous avons mis en œuvre une architecture de \textit{pipeline} sophistiquée permettant de gérer indépendamment le cycle de vie du \textit{Backend} (\texttt{Spring Boot}) et du \textit{Frontend} (\texttt{Angular}).

\paragraph*{ Architecture \textit{Parent-Child} (Monorepo)}
Le fichier racine \texttt{.gitlab-ci.yml} agit comme un chef d'orchestre. Il ne déclenche les \textit{pipelines} spécifiques que si des modifications sont détectées dans les répertoires respectifs (\texttt{BackEnd/**} ou \texttt{FrontEnd/**}). Cette approche optimise les ressources et réduit le temps global d'intégration.

\paragraph*{Flux de Travail et Automatisation}
Le \textit{pipeline} du \textit{backend} automatise l'intégralité de la chaîne de production, de la validation syntaxique (\textit{Validate}) à la compilation (\textit{Build}). Une étape de test cruciale est intégrée, instanciant dynamiquement un service \texttt{PostgreSQL 16-alpine} pour valider la persistance des données. S'ensuivent une analyse de qualité via \texttt{Checkstyle}, la génération du \textit{package} exécutable, et la construction d'une image \texttt{Docker} poussée sur le registre privé de \texttt{GitLab}. Le déploiement final est protégé par une action manuelle, garantissant un contrôle humain avant la mise à jour sur le serveur.

\vspace{2mm}
\begin{lstlisting}[language=yaml, caption=Extrait du fichier de configuration (\texttt{BackEnd/.gitlab-ci.yml})]
test:
  stage: test
  services:
    - name: postgres:16-alpine
      alias: postgres
  variables:
    SPRING_DATASOURCE_URL: "jdbc:postgresql://postgres:5432/talentpredict_test"
  script:
    - cd BackEnd
    - mvn test -B
\end{lstlisting}
\vspace{2mm}

\paragraph*{ Bénéfices pour le Projet}
Cette industrialisation offre une fiabilité totale grâce au cache \texttt{Maven} et à l'archivage des rapports \texttt{JUnit}. La sécurité est renforcée par l'utilisation de variables d'environnement masquées, assurant une conformité aux standards \textbf{DevOps} professionnels.


\subsubsection*{8. Sécurisation du Pipeline DevOps (DevSecOps)}

Le déploiement automatisé via \texttt{GitLab CI/CD} introduit des vecteurs de risques (fuite de clés, librairies vulnérables). \textbf{TalentPredict} intègre des mesures de sécurité à chaque étape du \textit{pipeline} pour garantir l'intégrité de la plateforme.

\paragraph*{Gestion des Secrets et Variables d'Environnement}
Le principe fondamental de "\textit{No Secrets in Git}" est strictement appliqué.
\begin{itemize}
    \item \textbf{Variables \texttt{GitLab CI/CD} :} Les identifiants sensibles (mot de passe DB, clé secrète \texttt{JWT}, identifiants du registre \texttt{Docker}) sont stockés de manière chiffrée dans les paramètres \texttt{GitLab}. Ils ne sont jamais présents dans le code source.
    \item \textbf{Injection Dynamique :} Ces variables sont injectées sous forme de variables d'environnement au moment du \textit{build}, assurant qu'elles ne sont accessibles qu'au processus autorisé.
\end{itemize}

\paragraph*{Analyse de Sécurité des Dépendances}
Le \textit{pipeline} intègre des vérifications automatiques pour prévenir l'utilisation de librairies obsolètes ou malveillantes :
\begin{itemize}
    \item \textbf{\texttt{Maven Dependency Check} :} Lors de l'étape de Test, \texttt{Maven} vérifie si les dépendances \texttt{Java} (\texttt{pom.xml}) contiennent des vulnérabilités connues (CVE).
    \item \textbf{Audit \texttt{npm} :} Une vérification similaire est effectuée sur le \textit{frontend} \texttt{Angular} pour s'assurer de l'intégrité du catalogue de modules \texttt{Node.js}.
\end{itemize}

\paragraph*{Isolation et Sécurité du Registre Docker}
\begin{itemize}
    \item \textbf{Registre Privé :} Nous utilisons le registre privé de \texttt{GitLab} pour stocker les images de nos conteneurs. Seul le \textit{pipeline} de déploiement et le serveur de production (via des identifiants \texttt{deploy\_token}) peuvent y accéder.
    \item \textbf{Principe du Moindre Privilège :} Les conteneurs en production sont configurés pour ne pas fonctionner avec des privilèges "\texttt{root}", limitant ainsi l'impact potentiel en cas de compromission d'un service.
\end{itemize}

\paragraph*{Contrôle d'Accès et Traçabilité}
\begin{itemize}
    \item \textbf{\textit{Protected Branches} :} Seuls les "\textit{Maintainers}" du projet peuvent déclencher un déploiement sur la branche \texttt{main} (Production).
    \item \textbf{\textit{Audit Log} :} Chaque étape du \textit{pipeline} est tracée dans \texttt{GitLab}, permettant de savoir exactement qui a déclenché un déploiement, avec quel code et à quel moment.
\end{itemize}

\subsubsection{Bilan du Déploiement}

Le déploiement final a permis de confirmer l'extrême résilience de l'architecture choisie. Grâce au couplage de l'orchestration Docker et des pipelines CI/CD, la plateforme \textbf{TalentPredict} peut désormais être mise à jour en quelques minutes de manière transparente. De plus, cette architecture modulaire garantit la \textit{scalabilité} future du système, permettant par exemple de dupliquer indépendamment le moteur IA en cas de forte montée en charge.


\section{Suivi et Pilotage du Sprint}
\subsection{Le Scrum Board}

Le \textit{Scrum Board} a constitué le pilier de notre gestion visuelle et opérationnelle tout au long de ce sprint. Structuré sous forme de tableau \textit{Kanban}, il a agi comme un véritable radiateur d'information, offrant une transparence totale sur le cycle de vie de chaque \textit{User Story}, de sa phase de conception jusqu'à sa validation finale. Cet outil a grandement facilité la coordination quotidienne des développements au sein de l'équipe technique, en synergie avec Mohamed Ben Abdeladhim et Rahma Barouni, permettant de répartir efficacement la charge de travail et d'identifier instantanément les points de blocage. Par ailleurs, ce formalisme a offert à notre \textit{Product Owner}, Med Amine Gharbi, un cadre clair et structuré pour effectuer la recette fonctionnelle continue et certifier que chaque module respectait scrupuleusement la \textit{Definition of Done} (DoD).

% Insertion de l'image du Scrum Board
\begin{figure}[H]
    \centering
    % Remplacez 'nom_image_scrum_board.png' par le nom exact de votre fichier
    \includegraphics[width=0.9\textwidth]{images/scrumboardSpring.png} 
    \vspace{2mm}
    \caption{Tableau de suivi visuel (Scrum Board) du Sprint 4}
    \label{fig:scrum_board_sprint4}
\end{figure}

\subsection{Le Burndown Chart}

En complément du suivi quotidien, le \textit{Burndown Chart} a représenté notre principal outil de pilotage prédictif. Ce graphique nous a permis de mesurer avec précision la vélocité de l'équipe en confrontant l'effort global restant à fournir avec la trajectoire de complétion idéale sur les quatre semaines de ce mois de mai 2026. L'analyse itérative de cette courbe s'est révélée cruciale pour la réussite de cette dernière phase : elle nous a permis de détecter de manière proactive les variations de rythme, notamment lors des défis techniques liés à l'intégration asynchrone, d'ajuster notre capacité d'exécution en temps réel, et de sécuriser ainsi la livraison de l'intégralité du périmètre fonctionnel avant la clôture du projet.

\begin{figure}[H]
    \centering
    \begin{tikzpicture}
    \begin{axis}[
        title={Burndown Chart - Sprint 4 (TalentPredict)},
        xlabel={Jours du Sprint},
        ylabel={Effort Restant (Story Points)},
        xmin=1, xmax=20,
        ymin=0, ymax=45,
        xtick={1,4,8,12,16,20},
        ytick={0,10,20,30,40},
        legend pos=north east,
        xmajorgrids=true,
        ymajorgrids=true,
        grid style=dashed,
        width=12cm,
        height=7cm
    ]

    % Ligne de l'effort IDÉAL (Ligne droite de 40 à 0)
    \addplot[
        color=blue,
        domain=1:20,
        dashed,
        line width=1pt
    ] {40 - (40/19)*(x-1)};
    \addlegendentry{Effort Idéal}

    % Ligne de l'effort RÉEL (Données fictives réalistes pour le Sprint 4)
    \addplot[
        color=red,
        mark=square*,
        line width=1.2pt
    ] coordinates {
        (1,40) (3,38) (5,35) (8,30) (10,28) (12,22) (15,12) (18,5) (20,0)
    };
    \addlegendentry{Effort Réel}

    \end{axis}
    \end{tikzpicture}
    \caption{Suivi de l'avancement du Sprint 4}
\end{figure} 

\subsection{Indicateurs de Performance du Sprint 4}

Ce dernier sprint a été le garant de l'industrialisation et de la mise en service. Nous avons mesuré la performance à travers trois prismes : l'agilité, la robustesse \textbf{DevOps} et la valeur métier des tableaux de bord.

\paragraph*{1. Indicateurs de Pilotage (Agilité)}
Le sprint s'est achevé avec une vélocité de 28 \textit{Story Points}. Bien que le volume soit inférieur aux sprints précédents, la complexité résidait dans l'automatisation et l'interfaçage des systèmes.
\begin{itemize}
    \item \textbf{Taux de complétion :} 100\%. Toutes les tâches de déploiement et de pilotage ont été terminées.
    \item \textbf{Qualité des \textit{User Stories} :} Les \textit{dashboards} RH ont été validés avec un taux d'acceptation immédiat, grâce à la précision des données agrégées venant du moteur d'IA.
\end{itemize}

\paragraph*{2. Indicateurs DevOps et Infrastructure}
L'industrialisation via \texttt{GitLab CI/CD} a permis d'atteindre des niveaux de performance remarquables :
\begin{itemize}
    \item \textbf{Vitesse d'Intégration :} Le \textit{build} du \textit{backend} s'effectue en moyenne en 1 min 30 s, tandis que le \textit{bundle} \textit{frontend} nécessite 2 min 45 s.
    \item \textbf{Fiabilité du Déploiement :} L'utilisation de \texttt{docker-compose} en mode \textit{detached} assure des mises à jour transparentes avec une interruption de service quasi-nulle.
    \item \textbf{Optimisation des Images :} La taille finale des images \texttt{Docker} a été réduite de 25\% par rapport aux premiers tests, optimisant ainsi l'espace disque sur le serveur.
\end{itemize}

\paragraph*{3. Indicateurs de Pilotage RH (\textit{Dashboards})}
\begin{itemize}
    \item \textbf{Temps de réponse du \textit{Dashboard} :} L'agrégation des données statistiques (plusieurs milliers d'entrées potentielles) se fait en moins de 300 ms, garantissant une fluidité de navigation pour l'administrateur.
    \item \textbf{Cohérence des données :} 100\% de concordance entre les résultats des tests individuels et les graphiques de tendances globales affichés.
\end{itemize}

\subsubsection*{4. Tableau de Synthèse Final (KPIs Sprint 4)}
\begin{table}[H]
    \centering
    \renewcommand{\arraystretch}{1.5}
    \small
    \begin{tabularx}{0.95\textwidth}{|X|c|c|c|}
        \hline
        \rowcolor[HTML]{F2F2F2}
        \textbf{Indicateur} & \textbf{Cible} & \textbf{Résultat Obtenu} & \textbf{État} \\
        \hline
        \textit{Story Points} complétés & 28 & 28 & $\checkmark$ Validé \\
        \hline
        Succès des \textit{Pipelines} \texttt{GitLab} & $100\%$ & $100\%$ & $\checkmark$ Validé \\
        \hline
        Temps moyen de \textit{Build} (Global) & $< 6$ min & 4 min 15 s & $\checkmark$ Validé \\
        \hline
        \textit{Uptime} Production (\texttt{Docker}) & $100\%$ & $100\%$ & $\checkmark$ Validé \\
        \hline
        Latence agrégation \textit{Dashboards} & $< 500$ ms & $\sim 280$ ms & $\checkmark$ Validé \\
        \hline
        Audit Sécurité (CORS/\texttt{JWT}/SSL) & Conforme & $100\%$ & $\checkmark$ Validé \\
        \hline
    \end{tabularx}
    \caption{Tableau de bord des métriques globales du Sprint 4}
    \label{tab:kpi_tech_s4}
\end{table}

\paragraph*{5. Analyse du Bilan}
Le bilan du Sprint 4 confirme la viabilité industrielle de \textbf{TalentPredict}. L'automatisation du cycle de vie logiciel (CI/CD) couplée à une infrastructure conteneurisée robuste permet d'envisager sereinement une mise en production réelle. Les tableaux de bord fournissent enfin une vision macroscopique des talents, transformant les données IA individuelles en un véritable outil de pilotage stratégique pour l'entreprise.

\section*{Conclusion}
\phantomsection
\addcontentsline{toc}{section}{Conclusion}
% Utilisez \section*{Conclusion} si vous ne souhaitez pas qu'elle soit numérotée

Ce quatrième et dernier chapitre s'est attaché à décrire avec précision l'ultime itération de notre cycle de développement, marquant l'aboutissement technique de la plateforme \textbf{TalentPredict}. En nous appuyant sur une conception architecturale solide et modulaire, nous avons exposé l'implémentation de fonctionnalités avancées et critiques, telles que le moteur de \textit{reporting} documentaire et le système de notifications asynchrones en temps réel.

L'intégration d'une démarche \textit{DevOps} rigoureuse a constitué un axe majeur de cette phase finale. La mise en place d'une stratégie de tests multiniveau, couplée à la conteneurisation via \texttt{Docker} et à l'automatisation des déploiements (CI/CD), a permis de transformer un code en développement en un produit logiciel fiable, sécurisé et prêt pour la production. De plus, le pilotage agile de ce sprint, soutenu par des artefacts \textit{Scrum} pertinents, a garanti une gestion optimale de l'effort d'équipe et le respect scrupuleux des délais impartis.

La réussite de cette dernière itération clôture ainsi la phase de réalisation technique de notre système, ouvrant naturellement la voie au bilan final et à la conclusion générale de ce projet de fin d'études.


\clearpage
\chapter{Conclusion Générale}
\addcontentsline{toc}{chapter}{Conclusion Générale}

Le présent mémoire de fin d’études, réalisé au sein de la société AgileX Management dans le cadre de la formation dispensée à l’Institut Supérieur d’Informatique de Mahdia, a porté sur la conception et le développement de \textit{TalentPredict AI}, une plateforme intelligente dédiée à l’évaluation et au développement des compétences en entreprise.

Ce projet s’inscrit dans un contexte où les organisations cherchent à mieux identifier, mesurer et valoriser les compétences de leurs collaborateurs. Les méthodes classiques d’évaluation, souvent fondées sur des CV déclaratifs, des entretiens ponctuels ou des appréciations subjectives, présentent plusieurs limites. Elles ne permettent pas toujours de croiser les informations disponibles, de valider objectivement les compétences techniques ni d’analyser de manière structurée les compétences comportementales. Face à ces limites, \textit{TalentPredict AI} propose une approche intégrée, multi-sources et orientée aide à la décision RH.

La réalisation du projet a été conduite selon une démarche Agile Scrum, organisée en quatre sprints complémentaires. Le premier sprint a permis de mettre en place les fondations techniques de la plateforme, notamment l’architecture applicative, le système d’authentification sécurisé par JWT, la gestion différenciée des profils employé et administrateur RH, ainsi que l’environnement conteneurisé avec Docker. Le deuxième sprint a constitué le cœur fonctionnel de l’évaluation des compétences, à travers l’analyse du CV, l’exploitation du profil GitHub, l’évaluation technique par QCM et challenge de code, ainsi que l’évaluation comportementale fondée sur le modèle PCM et enrichie par des traitements d’analyse sémantique. Le troisième sprint a transformé les résultats d’évaluation en recommandations concrètes, en intégrant un moteur d’upskilling, des parcours personnalisés, un suivi de progression et des mécanismes de validation. Enfin, le quatrième sprint a permis de finaliser la plateforme à travers les tableaux de bord analytiques, la génération du Talent Passport, la communication en temps réel et la préparation du déploiement.

Sur le plan technique, le projet a permis de mobiliser plusieurs technologies complémentaires. Le backend métier a été développé avec Spring Boot, tandis que les services d’intelligence artificielle ont été organisés autour de FastAPI, Ollama et n8n. L’interface utilisateur a été réalisée avec Angular, et la persistance des données a été assurée par PostgreSQL. Cette architecture modulaire facilite la séparation des responsabilités, améliore la maintenabilité du système et prépare l’intégration future de nouveaux modules d’analyse ou de recommandation.

La valeur ajoutée de TalentPredict AI réside principalement dans sa capacité à dépasser une évaluation purement déclarative des compétences. La plateforme confronte les informations fournies par le collaborateur à des preuves techniques observables, telles que le profil GitHub, les tests et les challenges de code. Elle prend également en compte les compétences comportementales à travers un cadre psychométrique structuré et une analyse sémantique locale. Cette combinaison permet de produire une vision plus complète du profil collaborateur, tout en générant des recommandations de formation adaptées aux écarts identifiés.

Le projet met également l’accent sur la souveraineté et la confidentialité des données. Le recours à une inférence locale via Ollama limite l’externalisation des données sensibles et renforce la conformité avec les exigences de protection des données personnelles. Cette orientation est particulièrement importante dans un contexte RH, où les informations traitées peuvent concerner les profils professionnels, les résultats d’évaluation, les préférences d’apprentissage et certaines dimensions comportementales.

Malgré les résultats obtenus, certaines limites doivent être prises en considération. L’évaluation sémantique pourrait être améliorée par l’entraînement ou l’adaptation de modèles spécialisés sur des corpus RH francophones. Les performances de la plateforme devraient également être testées sur un volume plus important d’utilisateurs et de profils afin de mieux évaluer sa robustesse en conditions réelles. Par ailleurs, l’intégration avec des référentiels de compétences reconnus, tels que ROME ou ESCO, permettrait d’améliorer la normalisation des résultats et la pertinence des recommandations.

Plusieurs perspectives peuvent être envisagées pour faire évoluer la solution. Il serait pertinent d’intégrer un module de feedback périodique entre collaborateurs afin d’enrichir l’évaluation comportementale par des observations issues du contexte de travail réel. La plateforme pourrait également être connectée à des systèmes d’information RH existants afin d’automatiser davantage la gestion des parcours de formation. À moyen terme, l’évolution vers une architecture plus scalable, basée sur une orchestration avancée des services, pourrait faciliter le passage d’un prototype académique à une solution exploitable à plus grande échelle.

En définitive, ce projet de fin d’études a permis de concevoir une plateforme complète, cohérente et évolutive répondant à une problématique réelle de gestion des talents. Il a également constitué une expérience formatrice, mobilisant des compétences en génie logiciel, intelligence artificielle, DevOps, conception d’interfaces et gestion de projet Agile. TalentPredict AI représente ainsi une base solide pour le développement futur d’un outil RH intelligent, capable d’accompagner les entreprises dans l’évaluation, la valorisation et le développement continu des compétences.

\color{red}
% =========================================================
% Suggestions de figures à ajouter pour améliorer le rapport
% =========================================================

\section*{Suggestions de figures à intégrer dans le rapport}
\phantomsection
\addcontentsline{toc}{section}{Suggestions de figures à intégrer dans le rapport}

Afin d’améliorer la lisibilité du rapport et de renforcer la valeur explicative des chapitres, plusieurs figures peuvent être ajoutées. Ces figures ont pour objectif de synthétiser les choix méthodologiques, fonctionnels, techniques et expérimentaux du projet \textit{TalentPredict AI}. Le tableau~\ref{tab:suggestions_figures_rapport} présente les principales figures recommandées, leur emplacement conseillé et leur objectif.

\begin{longtable}{|P{0.24\textwidth}|P{0.28\textwidth}|P{0.40\textwidth}|}
\hline
\rowcolor{TableHeaderBlue}
\color{white}\textbf{Emplacement recommandé} &
\color{white}\textbf{Titre proposé de la figure} &
\color{white}\textbf{Objectif de la figure} \\
\hline
\endfirsthead

\hline
\rowcolor{TableHeaderBlue}
\color{white}\textbf{Emplacement recommandé} &
\color{white}\textbf{Titre proposé de la figure} &
\color{white}\textbf{Objectif de la figure} \\
\hline
\endhead

\hline
\endfoot

Chapitre 1 -- Méthodologie Scrum &
Organisation Scrum du projet &
Présenter les rôles Scrum, les artefacts, les cérémonies et leur application dans le projet. \\
\hline

Chapitre 1 -- Solution proposée &
Architecture en couches &
Présenter les couches de la solution : présentation, API métier, orchestration IA, inférence locale, persistance et déploiement. \\
\hline


Chapitre 2 -- Concepts et technologies clés &
Pipeline NLP appliqué à l’analyse des profils &
Expliquer le passage du CV ou du profil professionnel vers l’extraction des compétences, la normalisation et le scoring. \\
\hline

Chapitre 2 -- Analyse GitHub &
Exploitation des données GitHub comme preuve technique &
Illustrer les signaux extraits de GitHub : langages, dépôts, commits, régularité, contribution et activité. \\
\hline

Chapitre 2 -- LLM et souveraineté des données &
Inférence locale avec Ollama &
Montrer la différence entre une approche cloud et une approche locale, en insistant sur la confidentialité et la conformité RGPD. \\
\hline

Chapitre 2 -- Analyse comparative &
Matrice de positionnement des solutions existantes &
Comparer \textit{TalentPredict AI} aux solutions existantes selon deux axes : objectivité de l’évaluation et souveraineté des données. \\
\hline

Chapitre 2 -- Contributions &
Positionnement scientifique et technique de \textit{TalentPredict AI} &
Synthétiser les trois contributions : validation multi-sources, évaluation technique et comportementale, inférence locale. \\
\hline

Sprint 1 -- Objectif du sprint &
Objectifs techniques et fonctionnels du Sprint 1 &
Présenter les objectifs du sprint : architecture, sécurité, authentification, profils et environnement Docker. \\
\hline

Sprint 1 -- Conception &
Architecture de sécurité du Sprint 1 &
Illustrer le fonctionnement de l’authentification JWT, du hachage BCrypt, des rôles et des routes protégées. \\
\hline


Sprint 1 -- Conception &
Modèle de données utilisateur et profil &
Montrer les entités principales liées à l’utilisateur, au rôle, au profil, au token et aux informations professionnelles. \\
\hline

Sprint 1 -- Déploiement &
Infrastructure Docker du Sprint 1 &
Présenter les conteneurs Angular, Spring Boot, PostgreSQL et leurs interactions réseau. \\
\hline

Sprint 1 -- Tests &
Stratégie de test de l’authentification &
Synthétiser les tests réalisés : inscription, connexion, récupération de mot de passe, routes protégées et accès par rôle. \\
\hline

Sprint 2 -- Objectif du sprint &
Pipeline global d’analyse multi-sources &
Montrer le flux complet depuis le CV, GitHub, QCM et challenge de code vers le scoring technique et comportemental. \\
\hline

Sprint 2 -- Analyse technique &
Pipeline d’analyse du CV &
Illustrer les étapes : chargement du fichier, extraction du texte, nettoyage, détection des compétences, normalisation et stockage. \\
\hline

Sprint 2 -- Analyse GitHub &
Workflow d’analyse GitHub &
Présenter les étapes d’appel à l’API GitHub, récupération des données, extraction des signaux techniques et calcul des scores. \\
\hline

Sprint 2 -- Analyse comportementale &
Pipeline d’évaluation PCM &
Montrer les étapes : questionnaire PCM, calcul heuristique, analyse sémantique, profil dominant et scores comportementaux. \\
\hline

Sprint 2 -- Orchestration &
Workflow n8n d’analyse des compétences &
Présenter comment n8n orchestre les appels entre les services, les prompts IA, FastAPI et le backend métier. \\
\hline

Sprint 2 -- Scoring &
Fusion multi-sources des scores &
Illustrer la triangulation entre CV, GitHub, QCM, challenge de code et analyse IA pour produire un score de confiance. \\
\hline

Sprint 3 -- Objectif du sprint &
Architecture du moteur de recommandation &
Présenter comment les résultats d’évaluation alimentent le moteur de recommandation de formations. \\
\hline

Sprint 2 -- IA générative &
Génération de QCM adaptatifs par IA &
Montrer le processus : profil de l’utilisateur, compétence cible, prompt, génération du QCM et validation des réponses. \\
\hline

Sprint 2 -- Challenge de code &
Évaluation automatique des challenges de code &
Illustrer le flux : génération du problème, soumission du code, évaluation, feedback et scoring. \\
\hline

Sprint 3 -- Upskilling &
Roadmap personnalisée de formation &
Présenter la construction du parcours de formation à partir des lacunes détectées et du profil comportemental. \\
\hline


Sprint 4 -- Objectif du sprint &
Pilotage stratégique des compétences &
Présenter la transformation des résultats individuels en indicateurs RH globaux. \\
\hline

Sprint 4 -- Tableaux de bord &
Tableau de bord employé &
Montrer les principales zones du dashboard employé : scores, recommandations, progression et Talent Passport. \\
\hline

Sprint 4 -- Tableaux de bord &
Tableau de bord administrateur RH &
Illustrer les vues RH : cartographie des compétences, écarts, formations, indicateurs globaux et rapports. \\
\hline

Sprint 4 -- Reporting &
Génération du Talent Passport &
Présenter le processus de génération du document synthétique à partir des résultats d’évaluation. \\
\hline

Sprint 4 -- DevOps &
Chaîne CI/CD et mise en production &
Montrer les étapes : commit, build, tests, image Docker, déploiement et supervision. \\
\hline

Sprint 4 -- Tests &
Stratégie globale de tests et validation &
Synthétiser les tests unitaires, tests d’intégration, tests API, tests UI et validation fonctionnelle. \\
\hline
\caption{Suggestions de figures pour améliorer le rapport}
\label{tab:suggestions_figures_rapport}\\
\end{longtable}
\color{black}

%\chapter*{Bibliographie}
\begin{thebibliography}{99}
\bibitem{ch2-ulrich1997} Ulrich, D. (1997). \emph{Human Resource Champions: The Next Agenda for Adding Value and Delivering Results}. Boston, MA: Harvard Business School Press.
\bibitem{ch2-salgado1997} Salgado, J. F. (1997). The Five Factor Model of Personality and Job Performance in the European Community. \emph{Journal of Applied Psychology}, 82(1), 30--43.
\bibitem{ch2-donaldson2002} Donaldson, S. I., \& Grant-Vallone, E. J. (2002). Understanding Self-Report Bias in Organizational Behavior Research. \emph{Journal of Business and Psychology}, 17(2), 245--260.
\bibitem{ch2-kruger1999} Kruger, J., \& Dunning, D. (1999). Unskilled and Unaware of It: How Difficulties in Recognizing One's Own Incompetence Lead to Inflated Self-Assessments. \emph{Journal of Personality and Social Psychology}, 77(6), 1121--1134.
\bibitem{ch2-rgpd2016} Règlement (UE) 2016/679 du Parlement européen et du Conseil du 27 avril 2016 (RGPD), Articles 9 et 25. \emph{Journal officiel de l'Union européenne}, L 119, 4 mai 2016.
\bibitem{ch2-mcdaniel1994} McDaniel, M. A., Whetzel, D. L., Schmidt, F. L., \& Maurer, S. D. (1994). The Validity of Employment Interviews: A Comprehensive Review and Meta-Analysis. \emph{Journal of Applied Psychology}, 79(4), 599--616.
\bibitem{ch2-costa1992} Costa, P. T., \& McCrae, R. R. (1992). \emph{Revised NEO Personality Inventory (NEO-PI-R) and NEO Five-Factor Inventory (NEO-FFI): Professional Manual}. Odessa, FL: Psychological Assessment Resources.
\bibitem{ch2-pittenger1993} Pittenger, D. J. (1993). The Utility of the Myers-Briggs Type Indicator. \emph{Review of Educational Research}, 63(4), 467--488.
\bibitem{ch2-kahler1996} Kahler, T. (1996). \emph{Process Communication Model}. Little Rock, AR: Kahler Communications.
\bibitem{ch2-thornton2006} Thornton, G. C., \& Rupp, D. E. (2006). \emph{Assessment Centers in Human Resource Management: Strategies for Prediction, Diagnosis, and Development}. Mahwah, NJ: Lawrence Erlbaum Associates.
\bibitem{ch2-celik2013} Celik, D., \& Elçi, A. (2013). A SWRL-Based Ontology for Resume Information Extraction. In \emph{Proceedings of the 2013 International Conference on Information Technology Based Higher Education and Training (ITHET)}. IEEE.
\bibitem{ch2-chien2020} Chien, H.-J., \& Chia-Hui, C. (2020). Resume Information Extraction with Cascaded Hybrid Model. In \emph{Proceedings of the 58th Annual Meeting of the Association for Computational Linguistics (ACL 2020)}, pp. 5559--5569.
\bibitem{ch2-shi2023} Shi, W., et al. (2023). Large Language Models Are Not Robust Multiple Choice Selectors. \emph{arXiv preprint arXiv:2309.03882}.
\bibitem{ch2-kalliamvakou2014} Kalliamvakou, E., Gousios, G., Blincoe, K., Singer, L., German, D. M., \& Damian, D. (2014). The Promises and Perils of Mining GitHub. In \emph{Proceedings of the 11th Working Conference on Mining Software Repositories (MSR 2014)}, pp. 92--101. ACM.
\bibitem{ch2-marlow2013} Marlow, J., Dabbish, L., \& Herbsleb, J. (2013). Impression Formation in Online Peer Production: Activity Traces and Personal Profiles in GitHub. In \emph{Proceedings of the 2013 Conference on Computer Supported Cooperative Work (CSCW 2013)}, pp. 117--128. ACM.
\bibitem{ch2-vaswani2017} Vaswani, A., Shazeer, N., Parmar, N., Uszkoreit, J., Jones, L., Gomez, A. N., Kaiser, L., \& Polosukhin, I. (2017). Attention Is All You Need. In \emph{Advances in Neural Information Processing Systems (NeurIPS 2017)}, vol. 30.
\bibitem{ch2-meta2024} Meta AI. (2024). Llama 3 Model Card. Meta Platforms. Disponible à : \url{https://llama.meta.com/llama3}
\bibitem{ch2-mistral2023} Mistral AI. (2023). Mistral 7B Technical Report. Disponible à : \url{https://arxiv.org/abs/2310.06825}
\bibitem{ch2-abdin2024} Abdin, M., et al. (2024). Phi-3 Technical Report: A Highly Capable Language Model Locally on Your Phone. \emph{arXiv preprint arXiv:2404.14219}.
\bibitem{ch2-touvron2023} Touvron, H., et al. (2023). LLaMA: Open and Efficient Foundation Language Models. \emph{arXiv preprint arXiv:2302.13971}.
\bibitem{ch2-xu2024} Xu, Y., et al. (2024). Benchmarking LLaMA 3 on HR NLP Tasks. \emph{arXiv preprint}.
\bibitem{ch2-vanderaalst2013} van der Aalst, W. M. P. (2013). Business Process Management: A Comprehensive Survey. \emph{ISRN Software Engineering}, vol. 2013, Article ID 507984.
\bibitem{ch2-newman2015} Newman, S. (2015). \emph{Building Microservices: Designing Fine-Grained Systems}. Sebastopol, CA: O'Reilly Media.
\bibitem{ch2-drachsler2009} Drachsler, H., Hummel, H., \& Koper, R. (2009). Factors Affecting the Robustness of Recommender Systems in Technology Enhanced Learning Environments. \emph{Journal of Digital Information}, 10(2).

\bibitem{kahler1996} Kahler, T. (1996). \emph{Process Communication Model}. Little Rock, AR: Kahler Communications.
\bibitem{touvron2023llama} Touvron, H., et al. (2023). LLaMA: Open and Efficient Foundation Language Models. \emph{arXiv:2302.13971}.
\bibitem{roziere2023codellama} Rozière, B., et al. (2023). Code Llama: Open Foundation Models for Code. \emph{arXiv:2308.12950}.
\bibitem{chien2020} Chien, H.-J., \& Chang, C.-H. (2020). Resume Information Extraction with Cascaded Hybrid Model. In \emph{Proceedings of ACL 2020}, pp. 5559--5569.
\bibitem{kalliamvakou2014} Kalliamvakou, E., Gousios, G., Blincoe, K., Singer, L., German, D. M., \& Damian, D. (2014). The Promises and Perils of Mining GitHub. In \emph{MSR 2014}, pp. 92--101. ACM.
\bibitem{celik2013} Celik, I., \& Elçi, A. (2013). An ontology-based information extraction system for resume parsing.
\bibitem{rgpd2016} Règlement (UE) 2016/679 du Parlement européen et du Conseil (RGPD), Article 9.
\bibitem{donaldson2002} Donaldson, S. I., \& Grant-Vallone, E. J. (2002). Understanding Self-Report Bias in Organizational Behavior Research. \emph{Journal of Business and Psychology}, 17(2), 245--260.
\bibitem{kruger1999dunning} Kruger, J., \& Dunning, D. (1999). Unskilled and Unaware of It. \emph{Journal of Personality and Social Psychology}, 77(6), 1121--1134.
\bibitem{yao2023react} Yao, S., et al. (2023). ReAct: Synergizing Reasoning and Acting in Language Models. \emph{ICLR 2023}.
\bibitem{drachsler2009} Drachsler, H., Hummel, H., \& Koper, R. (2009). Factors Affecting the Robustness of Recommender Systems. \emph{Journal of Digital Information}, 10(2).
\bibitem{angular} Google. (2024). Angular Documentation. \url{https://angular.dev}
\bibitem{springboot} VMware Tanzu. (2025). Spring Boot Reference Guide. \url{https://spring.io/projects/spring-boot}
\bibitem{fastapi} Ramírez, S. (2024). FastAPI Documentation. \url{https://fastapi.tiangolo.com}
\bibitem{n8n} n8n GmbH. (2025). n8n Documentation. \url{https://docs.n8n.io}
\bibitem{ollama} Ollama. (2024). Ollama Documentation. \url{https://ollama.com}
\bibitem{docker} Docker Inc. (2024). Docker Documentation. \url{https://docs.docker.com}
\bibitem{eks} Amazon Web Services. (2025). Amazon Elastic Kubernetes Service User Guide. \url{https://docs.aws.amazon.com/eks/}
\bibitem{postgresql} The PostgreSQL Global Development Group. (2024). PostgreSQL Documentation. \url{https://www.postgresql.org/docs/}
\bibitem{pdfjs} Mozilla Foundation. (2025). PDF.js Documentation. \url{https://mozilla.github.io/pdf.js/}
\bibitem{chartjs} Chart.js Contributors. (2025). Chart.js Documentation. \url{https://www.chartjs.org}
\bibitem{pdfplumber} jsvine. (2024). pdfplumber Documentation. \url{https://github.com/jsvine/pdfplumber}
\bibitem{scrum2020} Schwaber, K., \& Sutherland, J. (2020). \emph{The Scrum Guide}. \url{https://scrumguides.org}
\bibitem{devlin2019bert} Devlin, J., Chang, M.-W., Lee, K., \& Toutanova, K. (2019). BERT: Pre-training of Deep Bidirectional Transformers for Language Understanding. In \emph{NAACL-HLT 2019}.
\bibitem{salgado1997} Salgado, J. F. (1997). The Five Factor Model of Personality and Job Performance. \emph{Journal of Applied Psychology}, 82(1), 30--43.
\bibitem{ulrich1997} Ulrich, D. (1997). \emph{Human Resource Champions}. Harvard Business School Press.
\end{thebibliography}
\end{document}
